\documentclass[11pt]{amsart}
\usepackage{amsmath,amssymb,amsthm,booktabs,enumitem}
\usepackage[margin=1.15in]{geometry}
\usepackage[colorlinks,linkcolor=blue,citecolor=blue,urlcolor=blue]{hyperref}

\newtheorem{theorem}{Theorem}
\newtheorem{lemma}[theorem]{Lemma}
\newtheorem{proposition}[theorem]{Proposition}
\newtheorem{corollary}[theorem]{Corollary}
\theoremstyle{definition}
\newtheorem{definition}[theorem]{Definition}
\theoremstyle{remark}
\newtheorem{remark}[theorem]{Remark}

\newcommand{\eps}{\varepsilon}
\newcommand{\jt}{\tilde{\jmath}}
\newcommand{\NOR}{\mathrm{NOR}}
\newcommand{\MAJ}{\mathrm{MAJ}}
\newcommand{\val}{\mathrm{val}}
\newcommand{\PP}{\mathbb{P}}
\newcommand{\EE}{\mathbb{E}}
\newcommand{\cF}{\mathcal{F}}
\newcommand{\ghat}{\hat g}
\newcommand{\Comp}{\mathrm{Comp}}

\title[Sequential NOR on random and smoothed instances]{Sequential NOR on random and smoothed instances}
\author{Generated Computers}
\date{24 November 2026}

\setlength{\emergencystretch}{2.5em}
\begin{document}

\begin{abstract}
Sequential NOR is the P-complete recurrence $x_0=1$, $x_i=\NOR(x_{i-1},x_{j_i})$ with
$0\le j_i<i$ (Okhotin, 2011). We study it on random instances and on smoothed instances, in which
an adversary chooses the pointers and each pointer is then resampled with probability $\eps$.

(1) \emph{Upper bounds.} On uniformly random instances Newton's method over GF(2) is exact after
$O(\log n)$ rounds with probability $1-o(1)$, and the density of ones in the trace tends to
$(3-\sqrt5)/2$. For every $\eps>3/4$ an elementary iteration whose rounds are prefix
computations finishes in $O(\log n)$ rounds on the $\eps$-smoothing of every base, and there is a
deterministic logspace algorithm that is never wrong and gives up with probability at most $3/n$.
With relevance pruning, $O(\log n)$ rounds suffice on every base for every $\eps>0.7288$.

(2) \emph{Hardness.} For every $\eps\le0.4468$, a randomized NC algorithm that outputs the trace
with probability only $n^{-K}$ on the $\eps$-smoothing of every base gives
$\mathrm P\subseteq\mathrm{BPNC}$. Four embeddings of circuits reach noise $1/64$, $1/4$, every
level below $4/9$, and $0.4468$. The third keeps every value on two rails restored by chains of
set--reset latches; the fourth replaces the latch by chains that run a program of period $24$,
whose stages oscillate between two states, and rests on a certificate computed in exact
arithmetic. So, unless $\mathrm P\subseteq\mathrm{BPNC}$, the critical noise level lies in
$[0.4468,0.7288]$; in particular it is neither $4/9$ nor $3/4$. The same holds at each level
$\mathrm{NC}^k$,
$k\ge2$. Below $0.4468$ a polynomial speedup with near-linear work for smoothed instances would
give one for circuit evaluation.

(3) \emph{The two natural iterations.} Against adaptive adversaries both need linearly many rounds.
On random block bases, fixed in advance, they need $\Omega(n/\log n)$ rounds below explicit noise
levels, and $4/9$ is the exact level of the relevance-pruned iteration there. A first-moment count
of dependency chains cannot reach below about $0.72$ for all bases. On fixed-distance bases with
a large distance, a count of positions gives $O(\log n)$ rounds from noise $0.455$ on. A table iteration gives
$O(\sqrt{\eps n}+\log n)$ rounds on the block bases; a sound variant keeps sets of windows and is
fast, in simulation, on a base where the two iterations are slow. In the ideal model of the latch, layered
designs with independent reads forget their input from noise $0.455$ on, while the design of
period $24$ behind the fourth embedding keeps a bit up to about $0.447$. From noise $0.455$ on, an
errorless NC algorithm solves in $O(\log n)$ rounds the smoothing of every layered base with wide
bundles (logarithmically many chains, reading random copies or one or two positions back) over a long preamble; the circuit
embeddings can be built in this form, so this method of proving hardness stops between $0.4468$
and $0.455$.

None of this says anything about $\mathrm P$ versus $\mathrm{NC}$ in the worst case.
\end{abstract}

\maketitle

\noindent\textbf{Status.} Not refereed by a human. The proofs were developed with the assistance
of an AI system. One law-of-large-numbers step in Section~\ref{sec:newton} is only sketched; it is
marked there. All other proofs are complete; several rest on computations in exact or interval
arithmetic, and Section~\ref{sec:checks} lists them together with the passes of automated
referees and the deterministic lemmas that are checked in Lean.

\section{Introduction}\label{sec:intro}

Okhotin~\cite{okhotin} introduced Sequential NOR: an instance is a parent array
$j=(j_1,\dots,j_n)$ with $0\le j_i<i$, and its trace is $x_0=1$, $x_i=\NOR(x_{i-1},x_{j_i})$.
Evaluating the trace is P-complete under logspace reductions, so no algorithm of polylogarithmic
parallel time is expected to work on all instances~\cite{ghr}. This paper asks what happens on
typical instances: on uniformly random parent arrays, and on smoothed instances in the sense of
Spielman and Teng~\cite{gb-spielmanteng}, where an adversary fixes a base $j$ and every pointer is
replaced, with probability $\eps$, by a uniformly random earlier position.

\emph{Organisation.} Section~\ref{sec:setting} fixes the notation. Sections~\ref{sec:iterations}
and~\ref{sec:dense} introduce the two natural parallel iterations, the Kleene--prefix iteration
and its relevance-pruned variant, identify their round counts with two dependency depths, and
show that both finish in $O(\log n)$ rounds on smoothed instances for every $\eps>3/4$, using
that ones are dense in every smoothed trace. Section~\ref{sec:logspace} turns the chain count
behind this into an errorless logspace algorithm. Section~\ref{sec:newton} studies Newton's
method over GF(2) on uniformly random instances.

Sections~\ref{sec:gapbase}--\ref{sec:s24} are about hardness. Section~\ref{sec:gapbase}
compiles NOR circuits into parent arrays that survive noise $1/64$; Section~\ref{sec:gap} shows
that success probability $n^{-K}$ already suffices for the conclusion
$\mathrm P\subseteq\mathrm{BPNC}$, and derives a zero--one gap for the critical noise level.
Sections~\ref{sec:tree} and~\ref{sec:tree2} replace the restoring organ by trees of single NOR
gates and reach $3/20$ and $1/4$; Section~\ref{sec:latch} uses set--reset latches and reaches every
level below $4/9$. Section~\ref{sec:s24} replaces the latch by chains with a program of period
$24$ and reaches $0.4468$; so $4/9$ is not the critical noise level, unless
$\mathrm P\subseteq\mathrm{BPNC}$.

Sections~\ref{sec:adaptive}--\ref{sec:table} are about the upper end and the two
iterations: a barrier against adaptive adversaries, a two-valued iteration that is never slower,
bounds below $3/4$ for fixed-distance and ordered bases, a bound from noise $0.455$ on for
fixed-distance bases with a large distance (Section~\ref{sec:rows}), a base on which the expected number of
dependency chains grows exponentially at noise $0.72$, short pointers, lower bounds on random
block bases with the exact level $4/9$ for the pruned iteration, a proof that the relevance-pruned
iteration is fast on every base above noise $0.7288$ (Section~\ref{sec:allbases}), and a table
iteration with a sound variant that keeps sets of windows (Section~\ref{sec:setval}).
Section~\ref{sec:designs} compares the latch with other layered designs in its ideal model, and
shows that from noise $0.455$ on every layered base with wide bundles over a long preamble, the
circuit embeddings built in this form included, is solved by the relevance-pruned iteration in
$O(\log n)$ rounds; and
Section~\ref{sec:levels} restates the hardness and the algorithms for the levels of NC and for
polynomial speedups. Section~\ref{sec:checks} records the computations and checks.

\section{Setting}\label{sec:setting}

An \emph{instance} is a parent array $j=(j_1,\dots,j_n)$ with $0\le j_i<i$. Its \emph{trace} $x=x^*(j)$
is $x_0=1$ and $x_i=\NOR(x_{i-1},x_{j_i})$ for $1\le i\le n$. Gate $i$ is a \emph{NOT gate} if
$j_i=i-1$. A gate $k$ is a \emph{reset} if it is not a NOT gate and $x_{j_k}=1$; then $x_k=0$. Every
other gate satisfies $x_k=1\oplus x_{k-1}$. Let $\varrho(i)$ be the largest $r\le i$ that is $0$ or
a reset. Then
\begin{equation}\label{eq:reset}
x_i=x_{\varrho(i)}\oplus\bigl((i-\varrho(i))\bmod 2\bigr),\qquad x_0=1,\quad x_r=0\text{ at a reset }r .
\end{equation}

\emph{Smoothing.} For $\eps\in[0,1]$ and a parent array $j$ (the \emph{base}), the
$\eps$-smoothing $S_\eps(j)$ is the random array $\jt$ with $\jt_i=j_i$ if $\pi_i=0$ and
$\jt_i=\rho_i$ if $\pi_i=1$, where $\pi_i$ is Bernoulli$(\eps)$, $\rho_i$ is uniform on
$\{0,\dots,i-1\}$, and all these variables are independent. $S_1(j)$ is the uniform distribution
$U_n$ on parent arrays, whatever the base. We write $\cF_g=\sigma(\pi_i,\rho_i:i\le g)$, and for
the trace $x$ of $\jt$ we write $O_g=|\{p\le g:x_p=1\}|$. The problem is due to
Okhotin~\cite{okhotin}; the smoothing model is that of smoothed analysis~\cite{gb-spielmanteng}.

\emph{Chains.} A \emph{dependency chain} of length $\ell\ge 0$ from $P_0$ in an instance consists
of positions $P_0,\dots,P_\ell$ and gates $k_1,\dots,k_\ell$ such that, for $1\le s\le \ell$,
\[
k_s\le P_{s-1},\qquad k_s\text{ is not a NOT gate},\qquad\text{no gate in }(k_s,P_{s-1}]\text{ is a reset},\qquad P_s=j_{k_s}.
\]
This is Definition~\ref{kp:chain}; the middle two conditions say
$\max(\varrho(P_{s-1}),1)\le k_s$. Since $P_s\le k_s-2$, the positions decrease and the intervals
$[k_s,P_{s-1}]$ are pairwise disjoint. The \emph{weight} of a chain is
\[
\Omega=\prod_{s=1}^{\ell}\bigl(1+P_{s-1}-k_s\bigr).
\]

% ---------------------------------------------------------------------
% Fragment: two parallel iterations, their depths, and the threshold 3/4.
% Requires: amsmath, amssymb, amsthm, booktabs, enumitem, hyperref;
% environments theorem, lemma, proposition, corollary, definition, remark;
% macros \eps, \NOR, \PP, \EE, \cF, \val, \jt.
% All labels start with "kp:".
% ---------------------------------------------------------------------
\providecommand{\ind}{\mathbf{1}}
\providecommand{\epsK}{\eps_{\mathrm K}}
\providecommand{\epsB}{\eps_{\mathrm B}}
\providecommand{\gmin}{g_{\min}}
\providecommand{\AC}{\mathrm{AC}}
\providecommand{\NC}{\mathrm{NC}}
\providecommand{\BPNC}{\mathrm{BPNC}}

\section{Two parallel iterations and their depths}\label{sec:iterations}

This section introduces the Kleene--prefix iteration, which records only values
it knows for certain and resolves chains of NOT gates by a prefix computation,
and its relevance-pruned variant, which in addition skips gates that follow a
one. Each round of either iteration is in $\AC^0$. We prove that the number of
rounds of the first iteration is exactly a dependency depth $d$
(Lemma~\ref{kp:rounds-d}), that the number of rounds of the second is a depth
$D$ with $D\le d$ (Lemma~\ref{kp:pruned}), and we put both depths in a common
recursive form (Lemma~\ref{kp:switch}).

\subsection{Resets and the Kleene--prefix iteration}\label{kp:sub-resets}

Throughout this section $j=(j_1,\dots,j_n)$ is a fixed parent array,
$0\le j_k<k$, and $x^*=x^*(j)$ is its trace: $x^*_0=1$ and
$x^*_k=\NOR(x^*_{k-1},x^*_{j_k})$ for $1\le k\le n$. Recall that gate $k$ is a
\emph{NOT gate} if $j_k=k-1$; gate $1$ is always one. Gate $k$ is a
\emph{reset} if it is not a NOT gate and $x^*_{j_k}=1$; then $x^*_k=0$. Every
other gate satisfies $x^*_k=1\oplus x^*_{k-1}$: a NOT gate by definition, and a
gate that is neither a NOT gate nor a reset because $x^*_{j_k}=0$. For
$0\le i\le n$ let $\varrho(i)$ be the largest $r\le i$ that is $0$ or a reset.
(We write $\varrho$ rather than $\rho$, since $\rho_i$ will denote the
resampled pointers of the smoothing.) Then
\begin{equation}\label{kp:eq-alternate}
  x^*_i=x^*_{\varrho(i)}\oplus\bigl((i-\varrho(i))\bmod 2\bigr),
\end{equation}
where $x^*_0=1$ and $x^*_r=0$ for a reset $r$.
A gate that is not a NOT gate satisfies $j_k\le k-2$. Intervals are of
integers: $(a,b]=\{a+1,\dots,b\}$ and so on.

\begin{definition}[Kleene--prefix iteration]\label{kp:kleene}
The iteration maintains a set $K_t\subseteq\{0,\dots,n\}$ of \emph{known}
positions and values $v_p$ for $p\in K_t$. Start with $K_0=\{0\}$ and $v_0=1$.
Round $t+1$ is as follows.
\begin{itemize}[leftmargin=2em]
\item Gate $k$ is \emph{resolved} if it is a NOT gate or $j_k\in K_t$.
\item A resolved gate that is not a NOT gate and has $v_{j_k}=1$ is a
  \emph{resolved reset}.
\item For each $i$, let $r_i$ be the largest $r\le i$ that is $0$ or a resolved
  reset.
\item Then $i\in K_{t+1}$ if and only if every gate in $(r_i,i]$ is resolved,
  and in that case $v_i=v_{r_i}\oplus((i-r_i)\bmod 2)$, where $v_r=0$ for a
  resolved reset $r$.
\end{itemize}
\end{definition}

Each round is computed by a logspace-uniform circuit of polynomial size and
constant depth with unbounded fan-in. Resolution is a lookup. The index $r_i$
is the maximum of a set of polynomially many indices defined by $\AC^0$
predicates. ``Every gate in $(r_i,i]$ is resolved'' is one AND. The parity of
$i-r_i$ is the exclusive-or of the low bits of $i$ and $r_i$.

Define the \emph{dependency depth} by $d(0)=0$ and, for $1\le i\le n$,
\begin{equation}\label{kp:eq-d}
  d(i)=1+\max\bigl\{d(j_k): k \text{ not a NOT gate},\
  \max(\varrho(i),1)\le k\le i\bigr\},
\end{equation}
with $\max\emptyset=0$. Gates that are not NOT gates have $j_k\le k-2$, so $d$
is well defined, and induction gives $1\le d(i)\le i$ for $i\ge 1$.

\begin{lemma}[Rounds equal dependency depth]\label{kp:rounds-d}
For every $t\ge 0$, $K_t=\{i: d(i)\le t\}$, and $v_i=x^*_i$ for every
$i\in K_t$.
\end{lemma}

\begin{proof}
Induction on $t$; the case $t=0$ is clear, since $d(i)\ge 1$ for $i\ge 1$.
Assume the claim for $t$. Then gate $k$ is resolved in round $t+1$ exactly when
it is a NOT gate or $d(j_k)\le t$. A resolved gate $k$ that is not a NOT gate
has $v_{j_k}=x^*_{j_k}$; so it is a resolved reset if and only if it is a
reset. In particular every resolved reset is a reset, and hence
$r_i\le\varrho(i)$.

Suppose $d(i)\le t+1$, that is, every gate in $[\max(\varrho(i),1),i]$ that is
not a NOT gate has $d(j_k)\le t$. Then every gate in $[\max(\varrho(i),1),i]$
is resolved. If $\varrho(i)\ge 1$, then $\varrho(i)$ is a resolved reset, so
$r_i\ge\varrho(i)$ and therefore $r_i=\varrho(i)$. If $\varrho(i)=0$, then
$(0,i]$ contains no reset, hence no resolved reset, so $r_i=0$. Either way every
gate in $(r_i,i]$ is resolved, and $i\in K_{t+1}$.

Conversely, suppose $i\in K_{t+1}$. Every gate in $(r_i,i]$ is resolved, and
$r_i$ itself is $0$ or a resolved reset. Since $r_i\le\varrho(i)$, every gate
in $[\max(\varrho(i),1),i]$ is resolved, so every such gate that is not a NOT
gate has $d(j_k)\le t$, and $d(i)\le t+1$.

For the values, let $i\in K_{t+1}$. By maximality of $r_i$, no gate in
$(r_i,i]$ is a resolved reset. All of them are resolved, so by the first
paragraph none is a reset. Hence $x^*_k=1\oplus x^*_{k-1}$ for $k\in(r_i,i]$.
Moreover $x^*_{r_i}=v_{r_i}$: either $r_i=0$, or $r_i$ is a reset and
$x^*_{r_i}=0$. So $x^*_i=v_{r_i}\oplus((i-r_i)\bmod 2)=v_i$.
\end{proof}

In particular $K_t\subseteq K_{t+1}$, and a value, once set, never changes.
Running $R$ rounds and outputting $v$ if $K_R=\{0,\dots,n\}$, and $\bot$
otherwise, is therefore errorless, and it succeeds exactly when
$\max_i d(i)\le R$.

\begin{definition}[Dependency chains]\label{kp:chain}
A \emph{dependency chain} of length $L$ from $P_0$ consists of positions
$P_0,P_1,\dots,P_L$ and gates $k_1,\dots,k_L$, none of them a NOT gate, with
\[
  \max(\varrho(P_{s-1}),1)\le k_s\le P_{s-1}
  \quad\text{and}\quad P_s=j_{k_s}\qquad(1\le s\le L).
\]
\end{definition}

Because $P_s\le k_s-2$, the positions decrease and the intervals
$[k_s,P_{s-1}]$ are pairwise disjoint. Note that
$\max(\varrho(P),1)\le k\le P$ holds if and only if $1\le k\le P$ and no gate in
$(k,P]$ is a reset. If $d(i)=D\ge 1$, choose $k_1$ attaining the maximum
in~\eqref{kp:eq-d} for $i$, then $k_2$ attaining it for $P_1=j_{k_1}$, and so
on. As long as $d(P_{s-1})\ge 2$ the set in~\eqref{kp:eq-d} is nonempty, so this
gives a chain of length $D-1$ from $i$ with $d(P_s)=D-s$ for $0\le s\le D-1$.

\subsection{The relevance-pruned iteration}\label{kp:sub-pruned}

Between two resets the trace alternates, and a gate $k$ with $x^*_{k-1}=1$
outputs $0$ whatever its pointer reads. The iteration of
Definition~\ref{kp:kleene} nevertheless waits for such gates. Call gate $k$
\emph{relevant} if $x^*_{k-1}=0$, and call a reset \emph{relevant} or
\emph{irrelevant} according to its gate.

\begin{definition}[Relevance-pruned iteration]\label{kp:pruned-def}
As Definition~\ref{kp:kleene}, except for the last item. With $r_i$ as there,
and with $v_0=1$ and $v_r=0$ for a resolved reset $r$, put
$\alpha^{(i)}_p=v_{r_i}\oplus((p-r_i)\bmod 2)$ for $r_i\le p\le i$. Then
$i\in K_{t+1}$ if and only if every gate $k\in(r_i,i]$ with
$\alpha^{(i)}_{k-1}=0$ is resolved, and in that case $v_i=\alpha^{(i)}_i$.
\end{definition}

Each round is still in $\AC^0$: the new condition reads $v_{r_i}$ and the low
bits of $k-1$ and $r_i$. When $i$ is fixed we write $\alpha_p$ for
$\alpha^{(i)}_p$.

For the analysis, let $\varrho'(i)$ be the largest $r\le i$ that is $0$ or a
relevant reset. The \emph{anchors} of $i$ are the resets in $[\varrho'(i),i]$,
together with $0$ if $\varrho'(i)=0$; in particular $\varrho'(i)$ is an anchor.
For an anchor $r$ of $i$ put
\begin{equation}\label{kp:eq-V}
  V_r(i)=\max\Bigl(\{D(j_r)\}\cup\{D(j_k): k\in R_r(i)\}\Bigr),
\end{equation}
where $R_r(i)$ is the set of relevant gates in $(r,i]$ that are not NOT gates,
the first set $\{D(j_r)\}$ is omitted when $r=0$, and $\max\emptyset=0$. Define
$D(0)=0$ and, for $1\le i\le n$,
\begin{equation}\label{kp:eq-D}
  D(i)=1+\min_{r}V_r(i),
\end{equation}
where $r$ ranges over the anchors of $i$. All pointers involved are below $i$,
so $D$ is well defined.

\begin{lemma}\label{kp:pruned}
Let $K_t$ and $v$ refer to the relevance-pruned iteration.
\begin{enumerate}[label=\textup{(\alph*)},leftmargin=2em]
\item For every $t\ge 0$, $K_t=\{i: D(i)\le t\}$, and $v_i=x^*_i$ for every
  $i\in K_t$.
\item After every round, every position known to the Kleene--prefix iteration
  is known to the relevance-pruned iteration. Consequently $D(i)\le d(i)$ for
  every $i$.
\end{enumerate}
\end{lemma}

\begin{proof}
(a) Induction on $t$; the case $t=0$ is clear. Assume the claim for $t$. As in
the proof of Lemma~\ref{kp:rounds-d}, gate $k$ is resolved in round $t+1$
exactly when it is a NOT gate or $D(j_k)\le t$, and a resolved gate that is not
a NOT gate is a resolved reset if and only if it is a reset. Fix $i$ and write
$\alpha_p=\alpha^{(i)}_p$.

\emph{Values.} Let $i\in K_{t+1}$. We show $x^*_p=\alpha_p$ for
$r_i\le p\le i$, by induction on $p$. At $p=r_i$ this holds because
$v_{r_i}=x^*_{r_i}$ ($r_i$ is $0$ or a reset). Let $r_i<p\le i$. If
$\alpha_{p-1}=1$, then $x^*_{p-1}=1$ and $x^*_p=0=\alpha_p$. If
$\alpha_{p-1}=0$, then gate $p$ is resolved. A NOT gate gives
$x^*_p=1=\alpha_p$. Otherwise $x^*_{j_p}=v_{j_p}=0$, since gate $p$ would else be
a resolved reset in $(r_i,i]$, contradicting the maximality of $r_i$; so
$x^*_p=\NOR(0,0)=1=\alpha_p$. In particular $v_i=\alpha_i=x^*_i$.

\emph{$D(i)\le t+1$ implies $i\in K_{t+1}$.} Let $r$ attain the minimum
in~\eqref{kp:eq-D}. If $r\ge 1$, then $r$ is a reset with $D(j_r)\le t$, hence a
resolved reset. So $r_i\ge r\ge\varrho'(i)$. There is no relevant reset in
$(r_i,i]$, so on $[r_i,i]$ the trace alternates starting from
$x^*_{r_i}=v_{r_i}$: if $x^*_{p-1}=1$ then $x^*_p=0$, and if $x^*_{p-1}=0$ then
gate $p$ is relevant, hence not a reset, and $x^*_p=1$. Hence $\alpha_p=x^*_p$
on $[r_i,i]$, and the gates $k\in(r_i,i]$ with $\alpha_{k-1}=0$ are exactly the
relevant gates in $(r_i,i]$. Each of them is a NOT gate or, since $k>r_i\ge r$,
satisfies $D(j_k)\le V_r(i)\le t$. So each is resolved, and $i\in K_{t+1}$.

\emph{$i\in K_{t+1}$ implies $D(i)\le t+1$.} By the values paragraph, the gates
$k\in(r_i,i]$ with $\alpha_{k-1}=0$ are the relevant gates in $(r_i,i]$, and all
of them are resolved. None of them is a reset, since it would be a resolved
reset after $r_i$. So $\varrho'(i)\le r_i$, and $r_i$ is an anchor of $i$: if
$r_i\ge 1$ it is a reset in $[\varrho'(i),i]$, and if $r_i=0$ then
$\varrho'(i)=0$. Its own pointer, if $r_i\ge 1$, and the pointers of the
relevant gates after it that are not NOT gates lie in $K_t$, because these
gates are resolved. So $V_{r_i}(i)\le t$ and $D(i)\le t+1$.

(b) Write $K_t$, $r_i$ for the relevance-pruned iteration and $K^{\mathrm{KP}}_t$,
$r^{\mathrm{KP}}_i$ for the Kleene--prefix iteration. We show
$K^{\mathrm{KP}}_t\subseteq K_t$ by induction on $t$; the case $t=0$ is clear.
If $K^{\mathrm{KP}}_t\subseteq K_t$, then every gate resolved by the
Kleene--prefix iteration in round $t+1$ is resolved here, and, since both
iterations record correct values (Lemma~\ref{kp:rounds-d} and part~(a)), every
reset it resolves is a resolved reset here. So $r_i\ge r^{\mathrm{KP}}_i$. If it
puts $i$ into $K^{\mathrm{KP}}_{t+1}$, then every gate in
$(r^{\mathrm{KP}}_i,i]$ is resolved, hence so is every gate in the shorter
interval $(r_i,i]$, and $i\in K_{t+1}$. Finally, if $d(i)=t$ then
$i\in K^{\mathrm{KP}}_t\subseteq K_t$ by Lemma~\ref{kp:rounds-d}, so $D(i)\le t$
by part~(a).
\end{proof}

So every round bound proved for the Kleene--prefix iteration holds verbatim for
the relevance-pruned iteration. The next lemma puts both depths in a common
form.

\begin{lemma}[Switch characterisation]\label{kp:switch}
Let $j$ be a parent array of length $n$.
\begin{enumerate}[label=\textup{(\alph*)},leftmargin=2em]
\item $D(1)=1$, and for $2\le i\le n$,
  \[
    D(i)=\begin{cases}
      1+D(j_i) & \text{$i$ a relevant reset,}\\
      \min\bigl(D(i-1),\,1+D(j_i)\bigr) & \text{$i$ an irrelevant reset,}\\
      \max\bigl(D(i-1),\,1+D(j_i)\bigr) & \text{$i$ relevant, neither a
        reset nor a NOT gate,}\\
      D(i-1) & \text{otherwise.}
    \end{cases}
  \]
  The depth $d$ of Lemma~\ref{kp:rounds-d} satisfies $d(1)=1$ and, for
  $2\le i\le n$, $d(i)=1+d(j_i)$ if $i$ is a reset,
  $d(i)=\max(d(i-1),1+d(j_i))$ if $i$ is neither a reset nor a NOT gate, and
  $d(i)=d(i-1)$ otherwise.
\item Let $t\ge 1$ and $1\le i\le n$. Then $D(i)\ge t+1$ if and only if there
  is a relevant gate $a\le i$, not a NOT gate, with $D(j_a)\ge t$ and with
  $D(j_g)\ge t$ for every reset $g\in(a,i]$. The same holds for $d$ with the
  word ``relevant'' deleted.
\end{enumerate}
\end{lemma}

\begin{proof}
(a) By~\eqref{kp:eq-D}, $D(i)=1+\min_r V_r(i)$ over the anchors $r$ of $i$.
Gate $1$ is a NOT gate, hence not a reset, so $\varrho'(1)=0$, the only anchor
of $1$ is $0$, $V_0(1)=\max\emptyset=0$ and $D(1)=1$.

Let $i\ge 2$. If $i$ is a relevant reset, then $\varrho'(i)=i$, the only anchor
of $i$ is $i$, and $V_i(i)=D(j_i)$; so $D(i)=1+D(j_i)$. Otherwise
$\varrho'(i)=\varrho'(i-1)$, and the anchors of $i$ are those of $i-1$, together
with $i$ if $i$ is an irrelevant reset, in which case $V_i(i)=D(j_i)$. For an
anchor $r<i$ we have $V_r(i)=\max(V_r(i-1),D(j_i))$ if $i$ is relevant and not
a NOT gate, and $V_r(i)=V_r(i-1)$ otherwise. A relevant gate that is not a
relevant reset is not a reset at all. Since $\max(\cdot,c)$ commutes with
$\min$, taking the minimum over $r$ gives the four cases: for an irrelevant
reset $\min_r V_r(i)=\min(D(i-1)-1,D(j_i))$; for a relevant gate that is
neither a reset nor a NOT gate $\min_r V_r(i)=\max(D(i-1)-1,D(j_i))$; and
otherwise $\min_r V_r(i)=D(i-1)-1$.

For $d$: gate $1$ is a NOT gate, so the set in~\eqref{kp:eq-d} is empty for
$i=1$ and $d(1)=1$. Let $i\ge 2$. If $i$ is a reset then $\varrho(i)=i$, the
set in~\eqref{kp:eq-d} is $\{d(j_i)\}$, and $d(i)=1+d(j_i)$. Otherwise
$\varrho(i)=\varrho(i-1)$, and the set for $i$ is the set for $i-1$, together
with $d(j_i)$ exactly when $i$ is not a NOT gate. This gives the other two
cases (with $\max\emptyset=0$ and $d\ge 0$).

(b) Fix $t$, and write $I_i=\ind[D(i)\ge t+1]$ and $J_g=\ind[D(j_g)\ge t]$. By
(a), $I_1=0$ (as $t\ge 1$), and for $i\ge 2$:
$I_i=J_i$ at a relevant reset; $I_i=I_{i-1}\wedge J_i$ at an irrelevant reset;
$I_i=I_{i-1}\vee J_i$ at a relevant gate that is neither a reset nor a NOT
gate; and $I_i=I_{i-1}$ otherwise.

Suppose $a$ is as in the statement. Then $a\ge 2$, because gate $1$ is a NOT
gate, and $I_a=1$: $a$ is relevant and not a NOT gate, so either it is a
relevant reset and $I_a=J_a=1$, or $I_a=I_{a-1}\vee J_a=1$. For $g\in(a,i]$,
every case of the recursion gives $I_g\ge I_{g-1}$ unless $g$ is a reset with
$J_g=0$, which is excluded. So $I_i=1$.

Conversely, suppose $I_i=1$. Then $i\ge 2$, and we let $a\in[2,i]$ be largest
with $I_{a-1}=0$; it exists because $I_1=0$. By maximality $I_g=1$ for
$a\le g\le i$. Since $I_{a-1}=0$ and $I_a=1$, the recursion shows that $a$ is a
relevant reset with $J_a=1$ or a relevant gate, neither a reset nor a NOT gate,
with $J_a=1$; in both cases $a$ is relevant, not a NOT gate, and
$D(j_a)\ge t$. For $g\in(a,i]$ we have $I_{g-1}=I_g=1$, which a reset with
$J_g=0$ would contradict (it gives $I_g=0$).

The argument for $d$ is the same, with $I_i=\ind[d(i)\ge t+1]$,
$J_g=\ind[d(j_g)\ge t]$, $I_i=J_i$ at every reset, $I_i=I_{i-1}\vee J_i$ at
every gate that is neither a reset nor a NOT gate, and $I_i=I_{i-1}$ at NOT
gates.
\end{proof}

\begin{remark}\label{kp:rem-switch}
For $d$, replacing $a$ in part~(b) by the last reset in $(a,i]$, when there is
one, gives a gate $k\in[\max(\varrho(i),1),i]$, not a NOT gate, with
$d(j_k)\ge t$: one step of a dependency chain (Definition~\ref{kp:chain}). For
$D$ the last reset may be irrelevant, and then it cannot serve as $a$. Part~(b)
asks instead that it point to a deep position \emph{and} that a relevant gate
before it already make the level deep. This conjunction is all that the
minimum in~\eqref{kp:eq-D} adds.
\end{remark}

% =====================================================================

\section{Ones are dense, and the threshold \texorpdfstring{$3/4$}{3/4}}\label{sec:dense}

This section proves that on the $\eps$-smoothing of every base the
Kleene--prefix iteration, and hence the relevance-pruned one, finishes in
$O_\eps(\log n)$ rounds with probability at least $1-2/n$ whenever
$\eps>3/4$. The proof bounds the density of ones in the trace from below
(Lemmas~\ref{kp:dense}, \ref{kp:twosided} and~\ref{kp:limdens}) and then counts
dependency chains (Theorem~\ref{kp:thm34}). The three density bounds give the
thresholds $\epsK\approx 0.7808$ with explicit constants, $\epsB\approx 0.7676$,
and $3/4$ (Theorems~\ref{kp:thm34}, \ref{kp:sharper} and~\ref{kp:thm39}); by
Lemma~\ref{kp:limdens}(b), density bounds cannot go further. We deduce
errorless average-case $\AC^1$ and $\mathrm{SC}^2$ algorithms
(Corollaries~\ref{kp:cor315} and~\ref{kp:cor316}).

\subsection{Setting and a martingale inequality}\label{kp:sub-setting}

From now on $j$ is an arbitrary base of length $n$, $\eps\in(0,1]$, and
$\jt=S_\eps(j)$: $\jt_i=j_i$ if $\pi_i=0$ and $\jt_i=\rho_i$ if $\pi_i=1$, where
$\pi_i$ is Bernoulli$(\eps)$, $\rho_i$ is uniform on $\{0,\dots,i-1\}$, and all
these variables are independent. Let $x=x^*(\jt)$; resets, $\varrho$, the depths
$d$ and $D$, and dependency chains all refer to $\jt$. Let
\[
  \cF_g=\sigma(\pi_i,\rho_i: i\le g),\quad
  O_g=\bigl|\{p\le g: x_p=1\}\bigr|,\quad
  \psi_g=\frac{O_g}{g+1},\quad H_g=\sum_{i=1}^{g}\frac1i .
\]
So $\psi_g$ is the fraction of ones among the positions $0,\dots,g$, and
$O_g\ge 1$ since $x_0=1$. The prefix $x_0,\dots,x_{g}$ is
$\cF_{g}$-measurable, and $(\pi_i,\rho_i)$ is independent of $\cF_{i-1}$.
Ones are never adjacent, because $x_p=1$ forces $x_{p+1}=\NOR(1,\cdot)=0$;
hence
\begin{equation}\label{kp:eq-nonadj}
  O_g\le\frac{g+2}{2}\qquad(0\le g\le n).
\end{equation}

We use the Hoeffding--Azuma inequality with conditional ranges
\cite{hoeffding,kp-mcdiarmid}.

\begin{lemma}\label{kp:azuma}
Let $(\mathcal G_i)_{i\ge 0}$ be a filtration and let $Y_1,\dots,Y_m$ be random
variables with $Y_i$ $\mathcal G_i$-measurable, $\EE[Y_i\mid\mathcal G_{i-1}]=0$,
and $A_i\le Y_i\le A_i+1$ for some $\mathcal G_{i-1}$-measurable $A_i$. Then for
every $u\ge 0$,
$\PP\bigl[\sum_{i=1}^m Y_i\le -u\bigr]\le e^{-2u^2/m}$ and
$\PP\bigl[\sum_{i=1}^m Y_i\ge u\bigr]\le e^{-2u^2/m}$.
\end{lemma}

We apply it with $Y_i=x_i-a_i$, where $a_i=\EE[x_i\mid\cF_{i-1}]$; then
$-a_i\le Y_i\le 1-a_i$.

\subsection{Ones are dense in every smoothed instance}\label{kp:sub-dense}

\begin{lemma}[Ones are dense]\label{kp:dense}
Let $n\ge 2$. With probability at least $1-1/n$, for every $g\in[1,n]$,
\[
  O_g\ \ge\ \frac{\eps g+2-\eps H_g-2\sqrt{g\ln n}}{2+\eps}.
\]
\end{lemma}

\begin{proof}
By~\eqref{kp:eq-nonadj}, $O_{i-1}\le(i+1)/2$. Let $a_i=\EE[x_i\mid\cF_{i-1}]$.
If $x_{i-1}=0$ and gate $i$ reads a zero, then $x_i=1$. With probability $\eps$,
independently of $\cF_{i-1}$, gate $i$ reads the uniform position
$\rho_i\in\{0,\dots,i-1\}$, which holds a zero with conditional probability
$1-O_{i-1}/i\ge\frac12-\frac1{2i}$. So
\[
  a_i\ \ge\ \eps\,\ind[x_{i-1}=0]\Bigl(\frac12-\frac1{2i}\Bigr)
  \ \ge\ \frac\eps2\,\ind[x_{i-1}=0]-\frac{\eps}{2i}.
\]
Summing over $i\le g$, and using that $[0,g-1]$ contains
$g-O_{g-1}\ge g-O_g$ zeros, we get
$\sum_{i\le g}a_i\ge\frac\eps2(g-O_g)-\frac\eps2H_g$. By
Lemma~\ref{kp:azuma},
$\PP\bigl[\sum_{i\le g}(x_i-a_i)\le-\sqrt{g\ln n}\bigr]\le e^{-2\ln n}=n^{-2}$,
and a union bound over $g\le n$ costs $1/n$. Outside these events,
\[
  O_g-1=\sum_{i\le g}x_i\ \ge\ \frac\eps2(g-O_g)-\frac\eps2H_g-\sqrt{g\ln n}.
\]
Multiplying by $2$ and collecting the terms in $O_g$ gives
$(2+\eps)O_g\ge\eps g+2-\eps H_g-2\sqrt{g\ln n}$.
\end{proof}

\subsection{The round bound}\label{kp:sub-round}

Let
\[
  \epsK=\frac{\sqrt{17}-1}{4}\approx 0.7808 .
\]
Since $\eps^2/(2+\eps)>1-\eps$ is equivalent to $2\eps^2+\eps-2>0$, it holds
for $\eps\in(0,1]$ exactly when $\eps>\epsK$.

\begin{theorem}[Kleene--prefix on smoothed instances]\label{kp:thm34}
Let $\eps\in(\epsK,1]$, let $\eta\in(0,\eps/(2+\eps))$ and $\lambda>0$, and
put
\[
  c=\frac{\eps}{2+\eps}-\eta,\qquad
  \theta=\frac{\eps/(\lambda+1)+1-\eps}{\eps c}.
\]
Assume $\theta<1$; such $\eta$ and $\lambda$ exist exactly because
$\eps^2/(2+\eps)>1-\eps$. Let
\[
  \kappa=\Bigl(\frac{2}{0.9\,\eta\,(2+\eps)}\Bigr)^2,
\]
and let $\gmin$ be the least integer $g\ge 3$ with
$(\ln g+4)/(2g)\le\eta/10$. Then, for every $n\ge 2$ and every base $j$ of
length $n$, the following holds for $S_\eps(j)$. Put
$g_0=\lceil\kappa\ln n\rceil+\gmin$ and
\[
  L=\Bigl\lceil\frac{(\lambda+1)\ln(n+1)+\ln n}{\ln(1/\theta)}\Bigr\rceil .
\]
Let $\Gamma$ be the event that no dependency chain of length $L$ has all its
positions at least $g_0$. Then $\PP[\Gamma^c]\le 2/n$, and on $\Gamma$ the
Kleene--prefix iteration finishes within
\[
  A\ln(n+1)+\gmin+2\ \text{rounds},\qquad A=\frac{\lambda+2}{\ln(1/\theta)}+\kappa .
\]
\end{theorem}

\begin{proof}
On the existence of $\eta$ and $\lambda$: $\theta<1$ means
$\eps/(\lambda+1)+1-\eps<\eps c$. As $\eta\to 0$ and $\lambda\to\infty$ the two
sides tend to $1-\eps$ and $\eps^2/(2+\eps)$; conversely
$\theta>(1-\eps)/(\eps c)>(1-\eps)(2+\eps)/\eps^2$ always.

First, $\Gamma$ implies $\max_i d(i)<L+g_0$. Suppose instead that
$\max_i d(i)=D\ge L+g_0$, and take a chain of length $D-1$ from a maximiser
with $d(P_s)=D-s$ (Definition~\ref{kp:chain} and the paragraph after it).
Since $d(q)\le q$, we have $P_s\ge D-s\ge g_0$ for $s\le L$, and
$L\le D-g_0\le D-1$. So the first $L$ steps form a chain of length $L$ with all
positions at least $g_0$, contradicting $\Gamma$. By Lemma~\ref{kp:rounds-d} the
iteration then finishes within
\[
  L+g_0-1\ \le\ \frac{(\lambda+1)\ln(n+1)+\ln n}{\ln(1/\theta)}+1+\kappa\ln n+1+\gmin-1
  \ \le\ A\ln(n+1)+\gmin+2
\]
rounds. It remains to show $\PP[\Gamma^c]\le 2/n$.

\emph{Step 1: the good event.} Let $G$ be the event of Lemma~\ref{kp:dense}.
Let $g_0<g\le n$. On $G$, the lemma at $g-1$, together with $\eps\le 1$,
$\eps/(2+\eps)\le\frac12$ and $H_{g-1}\le\ln g+1$, gives
\begin{align*}
  \frac{O_{g-1}-1}{g}
  &\ \ge\ \frac{\eps}{2+\eps}-\frac{\eps(2+H_{g-1})}{(2+\eps)g}
       -\frac{2\sqrt{(g-1)\ln n}}{(2+\eps)g}\\
  &\ \ge\ \frac{\eps}{2+\eps}-\frac{2}{2+\eps}\sqrt{\frac{\ln n}{g}}
       -\frac{\ln g+4}{2g}
  \ \ge\ \frac{\eps}{2+\eps}-0.9\eta-0.1\eta\ =\ c.
\end{align*}
The last inequality uses $g\ge\kappa\ln n$, so that
$\frac{2}{2+\eps}\sqrt{\ln n/g}\le\frac{2}{(2+\eps)\sqrt\kappa}=0.9\eta$, and
$g\ge\gmin$, together with the fact that $(\ln g+4)/(2g)$ is decreasing in
$g\ge 1$. Let
\[
  G_g=\Bigl\{\frac{O_{g-1}-1}{g}\ge c\Bigr\}.
\]
Then $G_g\in\cF_{g-1}$, and $G\subseteq\bigcap_{g_0<g\le n}G_g$.

\emph{Step 2: one gate.} Let $g_0<g\le n$. On $G_g$, gate $g$ is a reset with
conditional probability at least $\eps c$ given $\cF_{g-1}$. Indeed, with
probability $\eps$ it is resampled, and then it is a reset when
$\rho_g\le g-2$ and $x_{\rho_g}=1$. Given $\cF_{g-1}$, that has probability
$(O_{g-1}-x_{g-1})/g\ge(O_{g-1}-1)/g\ge c$. Also, for $q\le k-2$,
\[
  \PP\bigl[\text{gate $k$ is not a NOT gate and $\jt_k=q$}\bigm|\cF_{k-1}\bigr]
  \ =\ w(k,q):=\frac{\eps}{k}+(1-\eps)\,\ind[j_k=q].
\]

\emph{Step 3: chains.} A \emph{skeleton} is a tuple
$\sigma=(P_0,k_1,P_1,\dots,k_L,P_L)$ of integers with
$P_s\le k_s-2$, $k_s\le P_{s-1}$ and $P_L\ge g_0$. Its event $E_\sigma$ says,
for each $s$: gate $k_s$ is not a NOT gate and $\jt_{k_s}=P_s$, and no gate in
$(k_s,P_{s-1}]$ is a reset. Every dependency chain of length $L$ with all
positions at least $g_0$ realises the event of its skeleton. The gates involved
in $E_\sigma$, namely those in the disjoint intervals $[k_s,P_{s-1}]$, are
distinct and all exceed $g_0$. List them as $g_1<\dots<g_M$, and let $B_m$ be
the event that $E_\sigma$ imposes at gate $g_m$; then $B_m\in\cF_{g_m}$. Its
conditional probability given $\cF_{g_m-1}$ is $w(k_s,P_s)$ if $g_m=k_s$, and
at most $1-\eps c$ on $G_{g_m}$ if $g_m\in(k_s,P_{s-1}]$, by Step~2. Since
$G\subseteq\bigcap_m G_{g_m}$, and since all factors but the last are
$\cF_{g_M-1}$-measurable,
\[
  \PP[E_\sigma\cap G]\ \le\
  \EE\Bigl[\prod_{m<M}\ind_{B_m}\ind_{G_{g_m}}\cdot
  \ind_{G_{g_M}}\PP[B_M\mid\cF_{g_M-1}]\Bigr],
\]
and peeling off the gates in decreasing order gives
\[
  \PP[E_\sigma\cap G]\ \le\ \prod_{s=1}^{L}(1-\eps c)^{P_{s-1}-k_s}\,w(k_s,P_s).
\]

\emph{Step 4: summation.} For functions $f\ge 0$ on $\{g_0,\dots,n\}$ let
\[
  (Wf)(P)=\sum_{k=g_0+1}^{P}(1-\eps c)^{P-k}\sum_{q=g_0}^{k-2}w(k,q)\,f(q),
\]
and let $h(q)=(q+1)^\lambda$. Summing the last display over the skeletons with
first position $P_0$ gives $(W^L\ind)(P_0)$, and $(W^L\ind)(P_0)\le(W^Lh)(P_0)$
because $h\ge 1$ and $W$ is monotone. Since
$\sum_{q\le k-2}(q+1)^\lambda\le\int_1^k x^\lambda\,dx\le k^{\lambda+1}/(\lambda+1)$
and $j_k+1\le k$,
\begin{align*}
  (Wh)(P)&\ \le\ \Bigl(\frac{\eps}{\lambda+1}+1-\eps\Bigr)
  \sum_{k\le P}(1-\eps c)^{P-k}k^\lambda\\
  &\ \le\ \Bigl(\frac{\eps}{\lambda+1}+1-\eps\Bigr)\frac{(P+1)^\lambda}{\eps c}
  \ =\ \theta\,h(P).
\end{align*}
Hence $(W^Lh)(P_0)\le\theta^Lh(P_0)\le\theta^L(n+1)^\lambda$. (The weight
$h(q)=(q+1)^\lambda$ is the standard device for heights of random trees
\cite{kp-devroye,kp-pittel}.)

\emph{Step 5.} The first position of a skeleton lies in $[1,n]$. Summing over
$P_0$ and adding $\PP[G^c]\le 1/n$,
\[
  \PP[\Gamma^c]\ \le\ \PP[G^c]+\sum_{\sigma}\PP[E_\sigma\cap G]
  \ \le\ \frac1n+n(n+1)^\lambda\theta^L\ \le\ \frac1n+\frac1{n+1}\ \le\ \frac2n,
\]
because $\theta^L\le(n+1)^{-(\lambda+1)}/n$.
\end{proof}

\begin{remark}\label{kp:reuse}
The first paragraph after the existence claim and Steps~2--5 of this proof use
only the following: $c>0$, $\lambda>0$, $\theta<1$, $g_0\ge 1$ an integer, and
an event $G$ with $\PP[G^c]\le 1/n$ and $G\subseteq\bigcap_{g_0<g\le n}G_g$.
Under these assumptions, with $L$ as in the theorem, they give
$\PP[\Gamma^c]\le 2/n$, and on $\Gamma$ the iteration finishes within
$L+g_0-1$ rounds. Moreover, Steps~3 and~4 bound
$\sum_\sigma\PP[E_\sigma\cap G]$, over the skeletons of any fixed length
$\ell\ge 0$ from a fixed $P_0$, by $\theta^\ell(n+1)^\lambda$.
\end{remark}

A grid search over $\eta$ and $\lambda$ gives the following constants.
\begin{center}
\begin{tabular}{lrrrrrr}
\toprule
$\eps$ & $0.79$ & $0.80$ & $0.85$ & $0.90$ & $0.95$ & $1$\\
\midrule
$A$ & $22495$ & $5205$ & $407$ & $137$ & $65.4$ & $35.1$\\
$\gmin+2$ & $8209$ & $3758$ & $939$ & $518$ & $348$ & $244$\\
\bottomrule
\end{tabular}
\end{center}
For uniform instances ($\eps=1$) the bound is $35.1\ln(n+1)+244$ rounds.

\subsection{A sharper density bound}\label{kp:sub-sharper}

Lemma~\ref{kp:dense} uses only that a resampled gate after a zero reads a zero
with probability at least about $1/2$. Feeding upper and lower bounds on the
density into each other gives a smaller threshold. Put
\[
  L(\mu)=\frac{\eps(1-\mu)}{1+\eps(1-\mu)},\qquad
  U(\ell)=\frac{1-\eps\ell}{2-\eps\ell},
\]
$\mu_0=1/2$, and $\ell_{s+1}=L(\mu_s)$, $\mu_{s+1}=U(\ell_{s+1})$ for
$s\ge 0$. Note that $\ell_1=\eps/(2+\eps)$ is the constant of
Lemma~\ref{kp:dense}.

\begin{lemma}[Two-sided density]\label{kp:twosided}
The sequence $(\ell_s)$ increases and $(\mu_s)$ decreases. Their limits $\ell^*$
and $\mu^*$ satisfy
\[
  \eps\ell^*=\frac12\Bigl(2+\eps-\sqrt{4+4\eps-3\eps^2}\Bigr).
\]
Let $s\ge 1$ and $\tau\in(0,1/(8s+2)]$. There is a constant $K$, depending only
on $s$ and $\tau$, such that for every $n\ge 2$ and every base, with
probability at least $1-1/n$,
\[
  \ell_s-(4s+1)\tau\ \le\ \psi_g\ \le\ \mu_s+(4s+1)\tau
  \qquad\text{for all } g\in[K\ln(n+1),n].
\]
\end{lemma}

\begin{proof}
\emph{The sequences.} For $\mu\le 1$ the map $\mu\mapsto L(\mu)$ is
decreasing, and for $\eps\ell<2$ so is $\ell\mapsto U(\ell)=1-1/(2-\eps\ell)$;
hence $U\circ L$ is increasing. Since $\mu_1=U(\ell_1)\le U(0)=\mu_0$, induction
gives $\mu_{s+1}=U(L(\mu_s))\le U(L(\mu_{s-1}))=\mu_s$, and then
$\ell_{s+2}=L(\mu_{s+1})\ge L(\mu_s)=\ell_{s+1}$. Both sequences lie in
$[0,1]$, and by continuity their limits satisfy $\ell^*=L(\mu^*)$ and
$\mu^*=U(\ell^*)$. With $w=\eps\ell^*$ we get $1-\mu^*=1/(2-w)$ and
$\ell^*=\eps/(2-w+\eps)$, that is, $w^2-(2+\eps)w+\eps^2=0$. The larger root is
at least $1+\eps/2$, while $w\le\eps\ell^*<1$ since $L<1$; so $w$ is the smaller
root, which is the stated value. Finally,
$L'(\mu)=-\eps/(1+\eps(1-\mu))^2$ and $U'(\ell)=-\eps/(2-\eps\ell)^2$, so
$|L'|\le\eps\le 1$ on $[0,1]$ and $|U'|\le\eps\le 1$ on $[-1,1]$.

\emph{Conditional means.} Let $a_i=\EE[x_i\mid\cF_{i-1}]$ and
$z_{i-1}=1-x_{i-1}$. If $x_{i-1}=1$ then $x_i=0$. Otherwise $x_i=1-x_{\jt_i}$,
and with probability $\eps$ the pointer is uniform on $\{0,\dots,i-1\}$, where
a fraction $1-\psi_{i-1}$ of the positions hold zeros; with probability
$1-\eps$ it is $j_i$. Hence
$a_i=z_{i-1}\bigl(\eps(1-\psi_{i-1})+(1-\eps)(1-x_{j_i})\bigr)$, and
\begin{equation}\label{kp:eq-condmean}
  \eps z_{i-1}(1-\psi_{i-1})\ \le\ a_i\ \le\ z_{i-1}(1-\eps\psi_{i-1}).
\end{equation}

\emph{The good event.} Let $G^*$ be the event that
\[
  \Bigl|\sum_{i=a+1}^{b}(x_i-a_i)\Bigr|\ \le\ \sqrt{2(b-a)\ln(n+1)}
  \qquad\text{for all } 0\le a<b\le n.
\]
By Lemma~\ref{kp:azuma}, applied to the filtration $(\cF_{a+i})_{i\ge 0}$, each
pair $(a,b)$ fails with probability at most $2(n+1)^{-4}$. There are at most
$(n+1)^2/2$ pairs, so $\PP[(G^*)^c]\le(n+1)^{-2}\le 1/n$.

\emph{Intervals.} On $G^*$, let $0\le a<b\le n$, $m=b-a$, $X=O_b-O_a$ and
\[
  \Delta(m)=1+\sqrt{2m\ln(n+1)} .
\]
Positions $a,\dots,b-1$ contain $X-x_b+x_a\in[X-1,X+1]$ ones, hence between
$m-X-1$ and $m-X+1$ zeros, and
$\sum_{i=a+1}^{b}z_{i-1}$ is this number of zeros. If $\psi_{i-1}\le\mu\le 1$
for all $i\in(a,b]$, the lower bound in~\eqref{kp:eq-condmean} gives, with
$u=\eps(1-\mu)\ge 0$,
$X\ge u(m-X-1)-\sqrt{2m\ln(n+1)}$, hence
\begin{equation}\label{kp:eq-int-lower}
  X\ \ge\ \frac{u}{1+u}\,m-\frac{u+\sqrt{2m\ln(n+1)}}{1+u}
  \ \ge\ L(\mu)\,m-\Delta(m).
\end{equation}
If $\psi_{i-1}\ge\ell\ge-1$ for all $i\in(a,b]$ (necessarily $\ell\le 1$, so
$u'=1-\eps\ell\in[0,2]$), the upper bound gives
$X\le u'(m-X+1)+\sqrt{2m\ln(n+1)}$, hence
\begin{equation}\label{kp:eq-int-upper}
  X\ \le\ \frac{u'}{1+u'}\,m+\frac{u'+\sqrt{2m\ln(n+1)}}{1+u'}
  \ \le\ U(\ell)\,m+\Delta(m).
\end{equation}

\emph{From intervals to prefixes.} For an integer $G\ge 0$ let
$G'=\lceil(G+1)/\tau\rceil+\lceil 8\ln(n+1)/\tau^2\rceil$. For $g\ge G'$ we
have $(G+1)/(g+1)\le\tau$ and, since $\ln(n+1)>1$ and $\tau<1$, also
$1/(g+1)\le\tau/8$ and $\sqrt{2g\ln(n+1)}/(g+1)\le\tau/2$, so
$\Delta(g)/(g+1)<\tau$. Let $g\in[G',n]$; note $g>G$. Suppose $\psi_p\le\mu\le1$
for all $p\in[G,n]$. Then~\eqref{kp:eq-int-lower} on $(G,g]$ gives
$O_g\ge L(\mu)(g-G)-\Delta(g)$, so $\psi_g\ge L(\mu)-2\tau$. Suppose instead
$\psi_p\ge\ell\ge-1$ for all $p\in[G,n]$. By~\eqref{kp:eq-nonadj},
$O_G\le(G+2)/2\le G+1$, and~\eqref{kp:eq-int-upper} gives
$O_g\le G+1+U(\ell)(g+1)+\Delta(g)$. So $\psi_g\le U(\ell)+2\tau$.

\emph{Stages.} Let $G_0=\lceil 1/(2\tau)\rceil$ and $G_{r+1}=(G'_r)'$ for
$r\ge 0$. By~\eqref{kp:eq-nonadj}, $\psi_g\le\frac12+\frac1{2(g+1)}\le\mu_0+\tau$
for $g\ge G_0$. Let $r<s$ and suppose $\psi\le\mu_r+e$ on $[G_r,n]$ with
$e=(4r+1)\tau$. Since $\tau\le 1/(8s+2)$, we have $(4r+3)\tau\le 1/2$, so
$\mu_r+e\le 1$ and $\ell_{r+1}-e-2\tau\ge-1/2$, and the Lipschitz bounds apply.
By the previous paragraph,
$\psi\ge L(\mu_r+e)-2\tau\ge\ell_{r+1}-e-2\tau$ on $[G'_r,n]$, and then
$\psi\le U(\ell_{r+1}-e-2\tau)+2\tau\le\mu_{r+1}+e+4\tau$ on $[G_{r+1},n]$.
Starting from $e=\tau$ at $r=0$, this gives
$\psi\le\mu_s+(4s+1)\tau$ on $[G_s,n]$ and, from the stage $r=s-1$,
$\psi\ge\ell_s-(4s-1)\tau$ on $[G'_{s-1},n]\supseteq[G_s,n]$.

Finally, $G\le k\ln(n+1)$ implies
$G'\le\bigl((k+1)/\tau+2+8/\tau^2\bigr)\ln(n+1)$, and
$G_0\le(1/(2\tau)+1)\ln(n+1)$. So $G_s\le K\ln(n+1)$ with $K$ depending only
on $s$ and $\tau$.
\end{proof}

\begin{theorem}[Sharper threshold]\label{kp:sharper}
Let $\epsB=(1+\sqrt{13})/6\approx 0.7676$ and $\eps\in(\epsB,1]$. There are
constants $A$ and $B$, depending only on $\eps$, such that for every $n\ge 2$
and every base $j$, with probability at least $1-2/n$, the Kleene--prefix
iteration on $S_\eps(j)$ finishes within $A\ln(n+1)+B$ rounds.
\end{theorem}

\begin{proof}
Let $f(\eps)=\eps\ell^*(\eps)-(1-\eps)$. By Lemma~\ref{kp:twosided},
\[
  f'(\eps)=\frac32-\frac{2-3\eps}{2\sqrt{4+4\eps-3\eps^2}}\ \ge\ 1
\]
on $[0,1]$, because $4+4\eps-3\eps^2\ge 4$ and $2-3\eps\le 2$ there. At
$\eps=\epsB$, the number $w=1-\eps$ satisfies
$w^2-(2+\eps)w+\eps^2=3\eps^2-\eps-1=0$ and $w<1+\eps/2$, so it is the smaller
root. So $f(\epsB)=0$, and $f(\eps)>0$ for $\eps>\epsB$.

Choose $c$ with $(1-\eps)/\eps<c<\ell^*$, and $\lambda>0$ with
$\theta=(\eps/(\lambda+1)+1-\eps)/(\eps c)<1$. Choose $s\ge 1$ with $\ell_s>c$
and then $\tau\in(0,1/(8s+2)]$ with $\ell_s-(4s+1)\tau\ge c+\tau$. Let $K$ be
the constant of Lemma~\ref{kp:twosided}, let $G$ be its event, and put
$g_0=\lceil K\ln(n+1)\rceil+\lceil 1/\tau\rceil$. For $g_0<g\le n$, on $G$,
\[
  \frac{O_{g-1}-1}{g}=\psi_{g-1}-\frac1g\ \ge\ c+\tau-\tau=c .
\]
So $G\subseteq\bigcap_{g_0<g\le n}G_g$, with $G_g$ as in Step~1 of the proof of
Theorem~\ref{kp:thm34}, and $\PP[G^c]\le 1/n$. By Remark~\ref{kp:reuse}, with
$L=\lceil((\lambda+1)\ln(n+1)+\ln n)/\ln(1/\theta)\rceil$, we get
$\PP[\Gamma^c]\le 2/n$, and on $\Gamma$ the iteration finishes within
$L+g_0-1\le A\ln(n+1)+B$ rounds, where $A=(\lambda+2)/\ln(1/\theta)+K$ and
$B=\lceil 1/\tau\rceil+1$.
\end{proof}

\subsection{The limiting density and the threshold \texorpdfstring{$3/4$}{3/4}}\label{kp:sub-limit}

The upper bounds $\mu_s$ are needed in Lemma~\ref{kp:twosided} only because
each stage uses one bound on the density over a whole range of prefixes.
Prefix densities change slowly, so the lower bound can be fed back locally
instead. This gives the density of the base $j_k=0$ as a lower bound for every
base.

\begin{lemma}[The limiting density]\label{kp:limdens}
Let $\varpi=\varpi(\eps)$ be the root in $(0,1/2)$ of
$\varpi=\eps(1-\varpi)^2$, that is,
$\varpi=\bigl(2\eps+1-\sqrt{4\eps+1}\bigr)/(2\eps)$. Let $\tau\in(0,1/4]$ and
$K=2^{16}\tau^{-7}$, and let $n\ge 2$.
\begin{enumerate}[label=\textup{(\alph*)},leftmargin=2em]
\item For every base, with probability at least $1-1/n$,
  $\psi_g\ge\varpi-\tau$ for all $g\in[K\ln(n+1),n]$.
\item For the base $j_k=0$, with probability at least $1-1/n$, also
  $\psi_g\le\varpi+\tau$ for all $g\in[K\ln(n+1),n]$.
\end{enumerate}
\end{lemma}

\begin{proof}
The function $\varpi\mapsto\varpi-\eps(1-\varpi)^2$ is increasing on $[0,1]$,
negative at $0$ and equal to $\frac12-\frac\eps4>0$ at $\frac12$, so the root
in $(0,1/2)$ exists and is unique; the formula solves the quadratic. Work on
the event $G^*$ of the proof of Lemma~\ref{kp:twosided}, which has probability
at least $1-1/n$. The function $L$ of that lemma is decreasing with
$|L'|\le\eps\le 1$ on $[0,1]$, and $L(\varpi)=\varpi$, because
$\eps(1-\varpi)=\varpi/(1-\varpi)$.

\emph{Windows.} Put $\delta=\tau/8$, $w_0=\lceil 1024\ln(n+1)/\tau^3\rceil$ and
$w_{k+1}=w_k+\lceil\delta w_k\rceil$. Then $w_k\ge w_0\ge 1/\delta$, so
$w_{k+1}\le(1+2\delta)w_k$. Fix $k$ with $w_{k+1}\le n$, and write $a=w_k$,
$b=w_{k+1}$, $m=b-a=\lceil\delta a\rceil$ and $\varphi_k=m/(b+1)$. Since
$\delta\le 1/32$ and $a\ge 32$, we have $b+1\le(1+2\delta)a+1\le 2a$, so
$\delta/2\le\varphi_k\le m/a\le 2\delta\le 1/16$. Also
$m\ge\delta w_0\ge 128\ln(n+1)/\tau^2$, so
$\Delta(m)/m\le\tau^2/128+\tau/8\le\tau/4$. For $i\in(a,b]$,
\[
  \psi_a-2\delta\ \le\ \frac{O_a}{b}\ \le\ \psi_{i-1}=\frac{O_{i-1}}{i}
  \ \le\ \frac{O_a+m}{a+1}\ \le\ \psi_a+2\delta,
\]
where the first inequality uses $O_a/b\ge\psi_a(a+1)/(b+1)=(1-\varphi_k)\psi_a$.
Moreover $\psi_a+2\delta\le 1$, because $\psi_a\le(a+2)/(2a+2)$
by~\eqref{kp:eq-nonadj}. Below we divide $O_b=\psi_a(a+1)+(O_b-O_a)$ by $b+1$
and use $(a+1)/(b+1)=1-\varphi_k$.

(a) Inequality~\eqref{kp:eq-int-lower} with $\mu=\psi_a+2\delta$ gives
$O_b-O_a\ge L(\psi_a+2\delta)m-\Delta(m)\ge(L(\psi_a)-2\delta)m-\Delta(m)$.
Hence, with $\varphi=\varphi_k$ and $2\delta=\tau/4$,
\[
  \psi_b\ \ge\ (1-\varphi)\psi_a+\varphi L(\psi_a)-\varphi\tau/2 .
\]
Let $\eta_k=\max(0,\varpi-\psi_{w_k})$. If $\psi_a\le\varpi$, then
$L(\psi_a)\ge\varpi$, so
$\varpi-(1-\varphi)\psi_a-\varphi L(\psi_a)\le(1-\varphi)(\varpi-\psi_a)$. If
$\psi_a>\varpi$, then $L(\psi_a)\ge\varpi-\eps(\psi_a-\varpi)$, so
$(1-\varphi)\psi_a+\varphi L(\psi_a)-\varpi\ge(1-2\varphi)(\psi_a-\varpi)\ge0$.
In both cases
\begin{equation}\label{kp:eq-contract}
  \eta_{k+1}\ \le\ (1-\varphi_k)\eta_k+\varphi_k\tau/2 .
\end{equation}
Hence $\eta_{k+1}\le\tau/2$ if $\eta_k\le\tau/2$, and otherwise
$\eta_{k+1}-\tau/2\le(1-\varphi_k)(\eta_k-\tau/2)$. Since $\varphi_k\ge\tau/16$,
we get $\eta_k\le\tau/2+e^{-k\tau/16}\eta_0$ for every $k$ with $w_k\le n$. Put
\[
  k^*=\Bigl\lceil\frac{16}{\tau}\ln\frac{5}{2\tau}\Bigr\rceil .
\]
Since $\eta_0\le\varpi<1/2\le 5/8$, we get
$\eta_k\le\tau/2+\frac{2\tau}{5}\cdot\frac58=3\tau/4$ for all $k\ge k^*$ with
$w_k\le n$.

Now let $w_{k^*}\le g\le n$, and let $k\ge k^*$ be largest with $w_k\le g$.
Then $g<w_{k+1}$, so
\[
  \psi_g\ \ge\ \frac{O_{w_k}}{g+1}\ \ge\ \psi_{w_k}\frac{w_k+1}{w_{k+1}}
  \ \ge\ \frac{\psi_{w_k}}{1+2\delta}\ \ge\ \psi_{w_k}-2\delta
  \ \ge\ \varpi-\frac{3\tau}4-\frac\tau4 .
\]
Finally, $2\delta k^*\le 4\ln(5/(2\tau))+\tau/4$, so
\[
  w_{k^*}\ \le\ w_0(1+2\delta)^{k^*}\ \le\ w_0e^{2\delta k^*}
  \ \le\ w_0e^{1/16}\Bigl(\frac{5}{2\tau}\Bigr)^4,
\]
and $w_0\le 1025\ln(n+1)/\tau^3$. So
$w_{k^*}\le 42700\,\tau^{-7}\ln(n+1)\le K\ln(n+1)$.

(b) On the base $j_k=0$, every gate that is not resampled reads $x_0=1$ and
outputs $0$. So the lower bound in~\eqref{kp:eq-condmean} is an equality:
$a_i=\eps z_{i-1}(1-\psi_{i-1})$. If $\psi_{i-1}\ge\ell\in[0,1]$ for all
$i\in(a,b]$, this gives, with $u=\eps(1-\ell)$,
$O_b-O_a\le u\bigl(m-(O_b-O_a)+1\bigr)+\sqrt{2m\ln(n+1)}$, hence
$O_b-O_a\le L(\ell)m+\Delta(m)$ as in~\eqref{kp:eq-int-upper}. Take
$\ell=\max(0,\psi_a-2\delta)$; then $L(\ell)\le L(\psi_a)+2\delta$ (if $\ell=0$,
because $\psi_a<2\delta$ and $|L'|\le 1$). The computation of~(a) gives
\[
  \psi_b\ \le\ (1-\varphi)\psi_a+\varphi L(\psi_a)+\varphi\tau/2 .
\]
The contraction is symmetric: if $\psi_a\ge\varpi$ then $L(\psi_a)\le\varpi$,
and if $\psi_a<\varpi$ then $L(\psi_a)\le\varpi+\eps(\varpi-\psi_a)$. So
$\eta'_k=\max(0,\psi_{w_k}-\varpi)$ obeys~\eqref{kp:eq-contract}, and
$\eta'_0\le\psi_{w_0}\le\frac12+\frac1{2(w_0+1)}\le\frac58$. Hence
$\eta'_k\le 3\tau/4$ for all $k\ge k^*$ with $w_k\le n$. Between windows,
$\psi_g\le(O_{w_k}+g-w_k)/(g+1)\le\psi_{w_k}+(g-w_k)/(g+1)\le\psi_{w_k}+2\delta$
for $w_k\le g<w_{k+1}$. The rest is as in~(a).
\end{proof}

\begin{theorem}[Threshold $3/4$]\label{kp:thm39}
Let $\eps\in(3/4,1]$. There are constants $A$ and $B$, depending only on
$\eps$, such that for every $n\ge 2$ and every base $j$, with probability at
least $1-2/n$, the Kleene--prefix iteration on $S_\eps(j)$ finishes within
$A\ln(n+1)+B$ rounds.
\end{theorem}

\begin{proof}
Differentiating $\varpi=\eps(1-\varpi)^2$ gives
$\varpi'=(1-\varpi)^2/(1+2\eps(1-\varpi))>0$, and $\varpi(3/4)=1/3$. So
$\eps\varpi(\eps)-(1-\eps)$ is increasing and vanishes at $3/4$, and
$\eps\varpi>1-\eps$ for $\eps>3/4$. Choose $c$ with $(1-\eps)/\eps<c<\varpi$,
and $\lambda>0$ with $\theta=(\eps/(\lambda+1)+1-\eps)/(\eps c)<1$. Let
$\tau=(\varpi-c)/2<1/4$, let $K$ and the event $G$ be those of
Lemma~\ref{kp:limdens}(a), and put
$g_0=\lceil K\ln(n+1)\rceil+\lceil 1/\tau\rceil$. For $g_0<g\le n$, on $G$,
\[
  \frac{O_{g-1}-1}{g}=\psi_{g-1}-\frac1g\ \ge\ \varpi-2\tau=c .
\]
So $G\subseteq\bigcap_{g_0<g\le n}G_g$, with $G_g\in\cF_{g-1}$ as in Step~1 of
the proof of Theorem~\ref{kp:thm34}, and $\PP[G^c]\le 1/n$. By
Remark~\ref{kp:reuse}, with
$L=\lceil((\lambda+1)\ln(n+1)+\ln n)/\ln(1/\theta)\rceil$, we get
$\PP[\Gamma^c]\le 2/n$, and on $\Gamma$ the iteration finishes within
$L+g_0-1\le A\ln(n+1)+B$ rounds, where $A=(\lambda+2)/\ln(1/\theta)+K$ and
$B=\lceil 1/\tau\rceil+1$.
\end{proof}

\begin{remark}\label{kp:rem-limits}
The constants of Theorem~\ref{kp:thm39} are explicit but astronomical: already
$K=2^{30}$ at $\tau=1/4$. Theorem~\ref{kp:thm34}, with the table after it, is
the quantitative statement for $\eps>\epsK$. The method cannot go below $3/4$
with density bounds alone. Steps~2--4 need $1-\eps<\eps c$, where $c$ bounds the
density of ones from below on every base. By Lemma~\ref{kp:limdens}(b), on the
base $j_k=0$ the density is at most $\varpi(\eps)+\tau$ for large $g$, with high
probability, and $\eps\varpi(\eps)\le 1-\eps$ for $\eps\le 3/4$. By
Lemma~\ref{kp:pruned}(b), Theorems~\ref{kp:thm34}, \ref{kp:sharper}
and~\ref{kp:thm39}, and Corollary~\ref{kp:cor315} below, hold verbatim for the
relevance-pruned iteration.
\end{remark}

\subsection{Consequences}\label{kp:sub-cons}

Recall that an errorless smoothed algorithm at noise $\eps$ with success $p$ is
a family $(\mathcal A_n)$ of circuits that receive a parent array of length $n$
and a random string $r$, such that $\mathcal A_n(j,r)\in\{x^*(j),\bot\}$ always
and $\PP_{\jt\sim S_\eps(j),\,r}[\mathcal A_n(\jt,r)\ne\bot]\ge p$ for every
base $j$; and that $\eps_c$ is the infimum of the noise levels at which such an
algorithm in randomized $\NC$ with success $1/2$ exists.

\begin{corollary}[Upper bound on the noise threshold]\label{kp:cor315}
For every $\eps\in(3/4,1]$ there is an errorless, deterministic,
logspace-uniform $\AC^1$ algorithm that outputs $\bot$ with probability at most
$2/n$ on $\eps$-smoothed instances of every base. Hence $\eps_c\le 3/4$.
\end{corollary}

\begin{proof}
Take $A$ and $B$ from Theorem~\ref{kp:thm39}. Run
$R=\lceil A\rceil\lceil\log_2(n+1)\rceil+\lceil B\rceil\ge A\ln(n+1)+B$ rounds
of the iteration, each in $\AC^0$; $R$ is computable in logspace. Output $v$ if
every position is known, and $\bot$ otherwise. By Lemma~\ref{kp:rounds-d} the
output is never wrong, and by Theorem~\ref{kp:thm39} it is $\bot$ with
probability at most $2/n$, which is at most $1/2$ for $n\ge 4$. Smaller $n$ are
handled by table lookup. A deterministic algorithm is a randomized one that
ignores its random bits, so every $\eps\in(3/4,1]$ is a noise level of the kind
in the definition of $\eps_c$, and $\eps_c\le 3/4$.
\end{proof}

\noindent Together with Corollary~\ref{cor:gaps24}, this places the critical noise level in
$[0.4468,3/4]$ unless $\mathrm P\subseteq\mathrm{BPNC}$. Section~\ref{sec:allbases} lowers the upper
end to $0.7288$, for the relevance-pruned iteration.

\begin{corollary}[Uniform instances]\label{kp:cor316}
Let $\eps\in(3/4,1]$. On $\eps$-smoothed instances of every base, and in
particular on uniformly random parent arrays (the law of $S_1(j)$, which
does not depend on $j$):
\begin{enumerate}[label=\textup{(\alph*)},leftmargin=2em]
\item there is an errorless, logspace-uniform $\AC^1$ algorithm, hence also
  one in $\NC^2$ and one in space $O(\log^2 n)$, that outputs $\bot$ with
  probability at most $2/n$;
\item there is a deterministic algorithm that runs in polynomial time and space
  $O(\log^2 n)$, is never wrong, and outputs $\bot$ with probability $O(1/n)$.
\end{enumerate}
\end{corollary}

\begin{proof}
(a) is Corollary~\ref{kp:cor315} together with $\AC^1\subseteq\NC^2$ and
Borodin's theorem $\NC^2\subseteq\mathrm{DSPACE}(\log^2 n)$
\cite{kp-borodin}.

(b) Use the notation of the proof of Theorem~\ref{kp:thm34}, or of
Theorem~\ref{kp:thm39} when $\eps\le\epsK$: the constants $c,\lambda,\theta$,
the integers $g_0=O(\log n)$ and $L=O(\log n)$, and the events $G$ and
$\Gamma$. The argument below uses only these and Remark~\ref{kp:reuse}.

First compute $x_0,\dots,x_{g_0}$ by forward evaluation into a table of
$g_0+1$ bits. To evaluate $x_q$ for $q>g_0$, scan $k=q,q-1,\dots$ down to
$g_0+1$. Skip NOT gates. For each other $k$, obtain $x_{\jt_k}$ from the table if
$\jt_k\le g_0$, and by a recursive call otherwise. Stop at the first $k$ with
$x_{\jt_k}=1$, which is $\varrho(q)$, and return $(q-k)\bmod 2$. If no such $k$
exists, there is no reset in $(g_0,q]$, and we return
$x_{g_0}\oplus((q-g_0)\bmod 2)$. Both answers are correct
by~\eqref{kp:eq-alternate}.

A recursive call from $q$ is made on $\jt_k$ for a gate $k$ that is not a NOT
gate, with $\max(\varrho(q),1)\le k\le q$ and $\jt_k>g_0$. So a stack of nested
recursive calls $q=P_0,P_1,\dots,P_\ell$ is a dependency chain of length $\ell$
with all positions above $g_0$, and on $\Gamma$ the recursion depth is below
$L$. Each frame stores two indices, so the space is
$O(L\log n+g_0)=O(\log^2 n)$. The number of calls made while evaluating $x_q$
is at most the number $N_q$ of such chains from $q$, of all lengths. Each chain
of length $\ell$ realises the event $E_\sigma$ of its skeleton, so by
Remark~\ref{kp:reuse}
\[
  \EE[N_q\ind_G]\ \le\ \sum_{\ell\ge0}\theta^\ell(n+1)^\lambda
  =\frac{(n+1)^\lambda}{1-\theta}.
\]
Each call scans at most $n$ gates. Steps~3 and~4 hold for every $\lambda'>0$ with
$\theta(\lambda')=(\eps/(\lambda'+1)+1-\eps)/(\eps c)<1$, with the same $c$ and
$g_0$; so the exponent $\lambda'$ in this call bound may be chosen separately
from the $\lambda$ that fixes $L$ ($\lambda'=\lambda$ is allowed). For example,
at $\eps=1$ Theorem~\ref{kp:thm34} with $\eta=0.01$ allows every
$\lambda'>2.1$.

Evaluate $q=1,\dots,n$ in turn. Abort with $\bot$ if the recursion depth
reaches $L$ or the total number of calls exceeds $(n+1)^{\lambda'+2}$. The
running time is then polynomial; with $\eps=1$ and $\lambda'$ close to $2.1$ it
is $O(n^{5.2})$. On $\Gamma\cap G$ the depth stays below $L$, and Step~5 in
fact shows $\PP[(\Gamma\cap G)^c]\le\PP[G^c]+\sum_\sigma\PP[E_\sigma\cap G]\le
2/n$. By Markov's inequality,
$\PP\bigl[G\cap\{\sum_qN_q>(n+1)^{\lambda'+2}\}\bigr]\le
n(n+1)^{\lambda'}/((1-\theta(\lambda'))(n+1)^{\lambda'+2})$. So the abort
probability is at most $2/n+1/((1-\theta(\lambda'))(n+1))=O(1/n)$. Otherwise
the output is the trace.
\end{proof}

\section{Average-case logspace}\label{sec:logspace}

\subsection{The density of ones}

\begin{lemma}[density of ones]\label{fact:density}
Let $\eps\in(3/4,1]$. There are a constant $c$ with $\eps c>1-\eps$, constants $K_0,K_1$, and for
every $n\ge 2$ an integer $g_0=g_0(n)$ with $1\le g_0\le K_0\ln(n+1)+K_1$, such that for every base $j$
of length $n$ there is an event $G$ with
\[
\PP[G^c]\le \frac1n\qquad\text{and}\qquad G\subseteq\bigcap_{g_0<g\le n}G_g,\quad G_g:=\Bigl\{\frac{O_{g-1}-1}{g}\ge c\Bigr\}\in\cF_{g-1}.
\]
\end{lemma}

\begin{proof}
For $\eps>(\sqrt{17}-1)/4$ this is Step~1 of the proof of Theorem~\ref{kp:thm34}, with
$c=\eps/(2+\eps)-\eta$ and $g_0=\lceil\kappa\ln n\rceil+g_{\min}$. For every $\eps>3/4$ it is the
first paragraph of the proof of Theorem~\ref{kp:thm39}, which rests on
Lemma~\ref{kp:limdens}(a), with $c$ below the limiting density $\varpi(\eps)$ and
$g_0=\lceil K\ln(n+1)\rceil+\lceil1/\tau\rceil$.
\end{proof}

\noindent The chain bound (Lemma~\ref{lem:onechain}) is Steps~2--3 of the proof of
Theorem~\ref{kp:thm34}; we repeat it in the form needed here.

\smallskip
Fix $\eps$, $c$, $g_0$ and $G$ as in Lemma~\ref{fact:density}. Since $(1-\eps)/(\eps c)<1$, we can
fix $\lambda>0$ with
\[
\theta:=\frac{\eps/(\lambda+1)+1-\eps}{\eps c}<1 .
\]
Put
\[
\vartheta:=\frac{\theta}{\eps c},\quad \Theta:=\max(1,\vartheta),\quad
L:=\Bigl\lceil\frac{(\lambda+1)\ln(n+1)+\ln n}{\ln(1/\theta)}\Bigr\rceil,\quad
M:=n\,(n+1)^{\lambda+1}\,L\,\Theta^{L}.
\]
Also let $w(k,q)=\eps/k+(1-\eps)\,\mathbf 1[j_k=q]$ and $h(q)=(q+1)^\lambda$. From now on
$\jt=S_\eps(j)$ and $x$ is its trace.

\subsection{Two first-moment bounds}

A \emph{skeleton} of length $\ell$ is a tuple $\sigma=(P_0,k_1,P_1,\dots,k_\ell,P_\ell)$ of integers
with $P_0\le n$, $k_s\le P_{s-1}$, $P_s\le k_s-2$ and $P_\ell\ge g_0$. Let $E_\sigma$ be the event
that $\sigma$ is a dependency chain of $\jt$: for every $s$, $\jt_{k_s}=P_s$ and no gate in
$(k_s,P_{s-1}]$ is a reset. (Since $P_s\le k_s-2$, the gate $k_s$ is then not a NOT gate.)

\begin{lemma}[one chain]\label{lem:onechain}
For every skeleton $\sigma$,
\[
\PP[E_\sigma\cap G]\le\prod_{s=1}^{\ell}(1-\eps c)^{P_{s-1}-k_s}\,w(k_s,P_s).
\]
\end{lemma}

\begin{proof}
Two facts about a single gate. (i) Let $g_0<g\le n$. Gate $g$ is a reset whenever $\pi_g=1$,
$\rho_g\le g-2$ and $x_{\rho_g}=1$. The pair $(\pi_g,\rho_g)$ is independent of $\cF_{g-1}$, and
$x_p$ is $\cF_{g-1}$-measurable for $p\le g-1$. So
$\PP[g\text{ is a reset}\mid\cF_{g-1}]\ge \eps\,(O_{g-1}-x_{g-1})/g\ge\eps c$ on $G_g$.
(ii) For $q\le k-2$, the event $\{\jt_k=q\}$ is independent of $\cF_{k-1}$ and has probability
$w(k,q)$.

The gates that $E_\sigma$ constrains are the $k_s$ and the gates in the intervals $(k_s,P_{s-1}]$.
They are distinct, and all exceed $g_0$, because $k_s\ge P_s+2\ge g_0+2$. List them as
$g_1<\dots<g_N$. Let $C_i$ be the event attached to $g_i$: either ``$g_i$ is not a reset'' or
``$\jt_{g_i}=P_s$''. Let $p_i$ be $1-\eps c$ or $w(k_s,P_s)$ accordingly, and let $G_i'$ be
$G_{g_i}$ in the first case and the sure event in the second. Then $C_i\in\cF_{g_i}$,
$G_i'\in\cF_{g_i-1}$, and $\PP[C_i\mid\cF_{g_i-1}]\le p_i$ on $G_i'$. Since
$E_\sigma\cap G\subseteq\bigcap_i(C_i\cap G_i')$ and the events $C_i,G_i'$ with $i<N$ lie in
$\cF_{g_N-1}$,
\[
\PP\Bigl[\bigcap_{i\le N}(C_i\cap G_i')\Bigr]
=\EE\Bigl[\mathbf 1_{\bigcap_{i<N}(C_i\cap G_i')}\,\mathbf 1_{G_N'}\,\PP[C_N\mid\cF_{g_N-1}]\Bigr]
\le p_N\,\PP\Bigl[\bigcap_{i<N}(C_i\cap G_i')\Bigr],
\]
and induction on $N$ gives $\prod_i p_i$. The interval $(k_s,P_{s-1}]$ contains $P_{s-1}-k_s$ gates.
\end{proof}

For $f\ge0$ define two kernels on positions $P\le n$:
\begin{align*}
(Wf)(P)&=\sum_{k=g_0+1}^{P}(1-\eps c)^{P-k}\sum_{q=g_0}^{k-2}w(k,q)f(q),\\
(Vf)(P)&=\sum_{k=g_0+1}^{P}(1+P-k)(1-\eps c)^{P-k}\sum_{q=g_0}^{k-2}w(k,q)f(q).
\end{align*}

\begin{lemma}[weighted count]\label{lem:kernel}
$Wh\le\theta h$ and $Vh\le\vartheta h$. Consequently, for every $P_0\le n$ and $\ell\ge1$, summing
over the skeletons $\sigma$ of length $\ell$ from $P_0$,
\[
\sum_\sigma\PP[E_\sigma\cap G]\le\theta^\ell(n+1)^\lambda,\qquad
\sum_\sigma\Omega(\sigma)\,\PP[E_\sigma\cap G]\le\vartheta^\ell(n+1)^\lambda .
\]
\end{lemma}

\begin{proof}
For $k\ge1$ we have $\sum_{q=0}^{k-2}(q+1)^\lambda\le\int_1^kt^\lambda\,dt\le k^{\lambda+1}/(\lambda+1)$,
and at most one $q$ has $j_k=q$, with $h(j_k)\le k^\lambda$. Hence
\[
\sum_{q=g_0}^{k-2}w(k,q)h(q)\le\Bigl(\frac{\eps}{\lambda+1}+1-\eps\Bigr)k^\lambda .
\]
Using $k^\lambda\le(P+1)^\lambda$, $\sum_{d\ge0}y^d=1/(1-y)$ and $\sum_{d\ge0}(1+d)y^d=1/(1-y)^2$
with $y=1-\eps c$,
\begin{align*}
(Wh)(P)&\le\Bigl(\frac{\eps}{\lambda+1}+1-\eps\Bigr)\frac{(P+1)^\lambda}{\eps c}=\theta h(P),\\
(Vh)(P)&\le\Bigl(\frac{\eps}{\lambda+1}+1-\eps\Bigr)\frac{(P+1)^\lambda}{(\eps c)^2}=\vartheta h(P).
\end{align*}
By Lemma~\ref{lem:onechain}, the first sum over skeletons is at most $(W^\ell\mathbf 1)(P_0)$ and
the second at most $(V^\ell\mathbf 1)(P_0)$: every skeleton has $k_s\ge g_0+1$ and $P_s\ge g_0$, and
$\Omega(\sigma)$ contributes the factor $1+P_{s-1}-k_s$ at step $s$. The kernels are positive and
$\mathbf 1\le h$, so $W^\ell\mathbf 1\le W^\ell h\le\theta^\ell h$ and
$V^\ell\mathbf 1\le\vartheta^\ell h$, and $h(P_0)\le(n+1)^\lambda$.
\end{proof}

Define two events:
\begin{itemize}[leftmargin=2.2em]
\item[$\Gamma$:] no dependency chain of $\jt$ of length $L$ has all its positions at least $g_0$;
\item[$\Lambda$:] every dependency chain of $\jt$ of length $\ell\in[1,L-1]$ with all its positions
at least $g_0$ has weight $\Omega\le M$.
\end{itemize}
The event $\Gamma$ is that of Theorem~\ref{kp:thm34}; the event $\Lambda$ is new.

\begin{proposition}\label{prop:events}
$\PP[\Gamma^c\cap G]\le 1/(n+1)$ and $\PP[\Lambda^c\cap G]\le1/(n+1)$. Hence
$\PP[(G\cap\Gamma\cap\Lambda)^c]\le 1/n+2/(n+1)\le 3/n$.
\end{proposition}

\begin{proof}
There are at most $n$ starting positions $P_0$. By Lemma~\ref{lem:kernel} and a union bound over
skeletons, $\PP[\Gamma^c\cap G]\le n(n+1)^\lambda\theta^L\le1/(n+1)$, because
$\theta^L\le(n+1)^{-(\lambda+1)}/n$ by the choice of $L$. For $\Lambda$, a skeleton with
$\Omega(\sigma)>M$ satisfies $\mathbf 1\le\Omega(\sigma)/M$, so
\[
\PP[\Lambda^c\cap G]\le\frac1M\sum_{\ell=1}^{L-1}\sum_{P_0}\sum_{\sigma}\Omega(\sigma)\PP[E_\sigma\cap G]
\le\frac{n(n+1)^\lambda}{M}\sum_{\ell=1}^{L-1}\vartheta^\ell\le\frac{n(n+1)^\lambda L\,\Theta^{L}}{M}=\frac1{n+1}.
\]
The last claim adds $\PP[G^c]\le1/n$.
\end{proof}

\subsection{The evaluator}

The machine is a deterministic Turing machine with a read-only input tape holding $n$ and the
array, a write-only output tape, and work tapes. Fix integers $\ghat=\ghat(n)\ge g_0(n)$ and
$B=B(n)$; they are chosen in the proof of Theorem~\ref{thm:logspace}. Offsets are stored in the
Elias gamma code~\cite{elias}: the integer $o\ge0$ is written as the code of $o+1$, which has
$2\lfloor\log_2(o+1)\rfloor+1$ bits and is self-delimiting, so a sequence of codewords can be
parsed from its beginning.

\medskip\noindent\textbf{Table.} Evaluate $x_0,\dots,x_{\ghat}$ forwards and store them as bits
$T[0],\dots,T[\ghat]$. (If $n\le\ghat$, output this table and stop.)

\medskip\noindent\textbf{Procedure} $\mathrm{EVAL}(q)$ for $\ghat<q\le n$. The work memory is a
\emph{stack} $(o_1,\dots,o_s)$ of offsets, $s\ge1$, starting as $(0)$. It determines positions and
gates by
\[
P_0=q,\qquad k_r=P_{r-1}-o_r\ \ (1\le r\le s),\qquad P_r=\jt_{k_r}\ \ (1\le r<s).
\]
The positions are not stored. Whenever the machine needs $P_{s-1}$ it recomputes it from $q$ and
the stack, by one pass over the stack and one lookup of the input per frame. Repeat the following
with $P=P_{s-1}$, $o=o_s$, $k=P-o$:
\begin{enumerate}[label=(\alph*),leftmargin=2.4em]
\item If $k=\ghat$: \emph{return} $T[\ghat]\oplus((P-\ghat)\bmod2)$.
\item If $k>\ghat$ and gate $k$ is a NOT gate: replace $o_s$ by $o_s+1$.
\item If $k>\ghat$ and gate $k$ is not a NOT gate, let $t=\jt_k$. If $t\le\ghat$, \emph{resolve}
with $v=T[t]$. Otherwise push a new frame with offset $0$; its position is $t$.
\end{enumerate}
To \emph{resolve} with $v$ (the value $x_t$ read by the gate $k$ of the top frame): if $v=1$,
\emph{return} $o_s\bmod2$; if $v=0$, replace $o_s$ by $o_s+1$. To \emph{return} $v$: pop the top
frame; if the stack is now empty, $\mathrm{EVAL}(q)$ ends with result $v$; otherwise resolve with
$v$.

\medskip\noindent\textbf{Main program.} Pass~1: run $\mathrm{EVAL}(q)$ for $q=\ghat+1,\dots,n$ and
discard the results; if the stack ever needs more than $B$ bits, output $\bot$ and halt. Pass~2:
output $T[0],\dots,T[\ghat]$ and then the result of $\mathrm{EVAL}(q)$ for $q=\ghat+1,\dots,n$.

\begin{lemma}[correctness]\label{lem:correct}
On every parent array, every frame created at a position $P$ returns $x_P$, and
$\mathrm{EVAL}(q)$ halts with result $x_q$. While a frame at position $P$ has offset $o$, no gate
in $(P-o,P]$ is a reset.
\end{lemma}

\begin{proof}
Induction on $P$. The invariant holds at $o=0$. Let $k=P-o>\ghat$. A NOT gate is not a reset, so
step (b) keeps the invariant. In step (c), $t\le k-2<P$, so the value $v=x_t$ is correct, by the
table or by the induction hypothesis. If $v=0$, gate $k$ is not a reset and the invariant is
kept. If $v=1$, gate $k$ is a reset and, by the invariant, the last one up to $P$; so
$\varrho(P)=k$ and \eqref{eq:reset} gives $x_P=(P-k)\bmod2=o\bmod2$. If $k=\ghat$ is reached, no
gate in $(\ghat,P]$ is a reset, so $x_g=1\oplus x_{g-1}$ on $(\ghat,P]$ and
$x_P=x_{\ghat}\oplus((P-\ghat)\bmod2)$. The offset of a frame is at most $P-\ghat$ and the
positions of nested frames decrease, so the procedure halts.
\end{proof}

\begin{lemma}[the stack is a chain]\label{lem:stackchain}
At every moment of $\mathrm{EVAL}(q)$, the tuple
\[
(P_0,k_1,P_1,\dots,k_{s-1},P_{s-1})
\]
is a dependency chain of length $s-1$ from $q$ whose positions are all greater than $\ghat$. Its
weight is $\prod_{r<s}(1+o_r)$, and $0\le o_s\le n$.
\end{lemma}

\begin{proof}
Frame $r+1$ was pushed in step (c), when the frame at position $P_{r-1}$ had offset $o_r$, the gate
$k_r=P_{r-1}-o_r$ was not a NOT gate, and $\jt_{k_r}=P_r>\ghat$. By Lemma~\ref{lem:correct} no gate in
$(k_r,P_{r-1}]$ is a reset. Offsets below the top do not change.
\end{proof}

\begin{lemma}[space]\label{lem:space}
On $\Gamma\cap\Lambda$, during every $\mathrm{EVAL}(q)$ the stack has at most $L$ frames and occupies
at most
\[
B_0:=2\log_2M+2\log_2(n+1)+L
\]
bits. Moreover $B_0\le C_1\log_2(n+1)+C_2$ with
$C_1=2(\lambda+3)+a\,(3+2\log_2\Theta)$, $C_2=3+2\log_2\Theta$ and $a=(\lambda+2)/\ln(1/\theta)$.
\end{lemma}

\begin{proof}
By Lemma~\ref{lem:stackchain} the frames below the top form a chain of length $s-1$ with positions
greater than $\ghat\ge g_0$. A chain of length at least $L$ has a prefix of length $L$, so on
$\Gamma$ we get $s-1\le L-1$. On $\Lambda$ its weight is at most $M$ (trivially if $s=1$). The
frames below the top take $\sum_{r<s}(2\lfloor\log_2(1+o_r)\rfloor+1)\le2\log_2M+(s-1)$ bits, and
the top frame at most $2\log_2(n+1)+1$.

For the second claim, $\log_2M=\log_2n+(\lambda+1)\log_2(n+1)+\log_2L+L\log_2\Theta$, and
$\log_2L\le L-1$, so $B_0\le2(\lambda+3)\log_2(n+1)+L(3+2\log_2\Theta)$. Finally
$L\le a\ln(n+1)+1\le a\log_2(n+1)+1$.
\end{proof}

\subsection{The theorem}

\begin{theorem}[average-case logspace]\label{thm:logspace}
Let $\eps\in(3/4,1]$. There is a deterministic Turing machine $A_\eps$ with work space
$O_\eps(\log n)$, and therefore polynomial running time, such that
\begin{enumerate}[label=(\roman*),leftmargin=2.4em]
\item on every parent array $j$, the output is $x^*(j)$ or $\bot$;
\item for every $n\ge2$ and every base $j$ of length $n$,
$\ \PP\bigl[A_\eps(S_\eps(j))=\bot\bigr]\le 3/n$.
\end{enumerate}
In particular, with $\eps=1$: on a uniformly random instance of length $n$ the machine outputs
the trace with probability at least $1-3/n$.
\end{theorem}

\begin{proof}
Take the evaluator with
$\ghat(n)=\lceil K_0\rceil\lceil\log_2(n+1)\rceil+\lceil K_1\rceil$ and
$B(n)=\lceil C_1\rceil\lceil\log_2(n+1)\rceil+\lceil C_2\rceil$. Both are computed from the bit
length of $n+1$, and $\ghat\ge K_0\ln(n+1)+K_1\ge g_0$, $B\ge B_0$.

\emph{(i).} If the machine does not output $\bot$ in Pass~1, then Pass~2 repeats the same
computation and outputs, by Lemma~\ref{lem:correct}, the trace. Nothing is written before Pass~1
has finished.

\emph{(ii).} By Lemma~\ref{lem:space}, on $\Gamma\cap\Lambda$ the stack never exceeds $B_0\le B$
bits, so the machine does not output $\bot$. Proposition~\ref{prop:events} bounds the
probability of the complement by $3/n$.

\emph{Resources.} The table has $\ghat+1=O(\log n)$ bits, the stack at most $B+O(1)$ bits, and
the registers ($q$, the pass, the current $P$, $k$, $t$, a position in the stack) $O(\log n)$
bits. The machine halts on every input, by Lemma~\ref{lem:correct} and the space check. A halting
deterministic machine never repeats a configuration, and there are $n^{O(1)}$ configurations.
\end{proof}

\begin{corollary}\label{cor:classes}
The distributional problem (Sequential NOR, $(U_n)$) has an errorless average-case algorithm in
deterministic logspace with failure probability at most $3/n$, and so has the $\eps$-smoothing of
every base for every $\eps>3/4$. Since
$\mathrm{L}\subseteq\mathrm{NL}\subseteq\mathrm{AC}^1\subseteq\mathrm{NC}^2$ and
$\mathrm{L}=\mathrm{SC}^1$, this contains both parts of Corollary~\ref{kp:cor316} (errorless
$\mathrm{AC}^1$; polynomial time with $O(\log^2n)$ space), with $3/n$ in place of $2/n$.
\end{corollary}

\begin{remark}[what changed]
Corollary~\ref{kp:cor316}(b) runs the same recursion and stores two indices per frame, which costs
$O(\log n)$ bits for each of $O(\log n)$ frames. Here a frame is an offset, and the weighted
kernel $V$ shows that the product of all offsets on the stack is polynomial in $n$. The sum of
their code lengths is therefore $O(\log n)$. The price is the factor $1/(\eps c)$ between
$\theta$ and $\vartheta$, which only affects constants, because $\vartheta$ need not be below $1$.
\end{remark}

\begin{remark}[constants for uniform instances]\label{rem:constants}
For $\eps=1$, with $\eta=0.126$ and $\lambda=13.1$ in the notation of Theorem~\ref{kp:thm34}, one
gets $c\approx0.2073$, $\theta\approx0.342$, $\vartheta\approx1.650$, and the bound $B_0$ of
Lemma~\ref{lem:space}, evaluated without the crude step $\log_2L\le L-1$, is at most
$56.1\log_2(n+1)+2\log_2\log_2(n+1)+10$ bits. The table has
$\lceil34.6\ln n\rceil+398$ bits. In simulation the stack stayed below $6.2\log_2n$ bits on uniform
instances up to $n=4096$ (Section~\ref{sec:checks}).
\end{remark}

\begin{remark}[scope]
The threshold $3/4$ is inherited from Lemma~\ref{fact:density} and from the condition
$\theta<1$, exactly as in Theorem~\ref{kp:thm34}. The argument of this section applies wherever the first-moment
chain bound of Lemma~\ref{lem:onechain} holds with $\theta<1$. The table cannot simply be
dropped: below $g_0$ the density bound is not available, and chains inside $[0,g_0]$ would cost
$O(\log n\log\log n)$ bits.
\end{remark}


\section{Random instances: Newton's method over GF(2)}\label{sec:newton}

\providecommand{\nwone}{\mathbf{1}}
\providecommand{\nwcl}{\mathrm{cl}}
\providecommand{\nwGF}{\mathrm{GF}(2)}
\providecommand{\nwsucc}{\mathrm{succ}}

This section studies Sequential NOR on uniformly random instances, in which
each pointer $j_i$ is uniform on $\{0,\dots,i-1\}$ and the pointers are
independent. Over $\nwGF$ we have $\NOR(a,b)=(1+a)(1+b)$, so the trace is the
unique solution of a triangular polynomial system. Newton's method linearises
that system at the current iterate and solves the resulting unit lower
triangular linear system; one round is one $\mathrm{NC}^2$ computation. We
prove the following, with explicit constants.
\begin{itemize}[leftmargin=*]
\item \emph{Theorem~A} (Theorem~\ref{nw:thm-A}). From the start
  $x^{(0)}=(1,0,\dots,0)$, and also from a uniformly random start, Newton's
  method is exact after $\tfrac12 y_0+\lceil 4.1\,e^{100}\ln n\rceil=O(\log n)$
  rounds with probability $1-o(1)$, where $y_0=\max(e^{200},\ln n)$ (for the
  random start, with $\frac12(y_0+1)$ in place of $\frac12y_0$).
\item \emph{Theorem~B} (Theorem~\ref{nw:thm-B}). From the start
  $(1,0,\dots,0)$, at least $\log_2\ln n-3$ rounds are needed with
  probability $1-o(1)$.
\item \emph{Theorem~C} (Theorem~\ref{nw:thm-C}). The Horton--Strahler rank,
  which bounds the number of rounds syntactically
  (Proposition~\ref{nw:prop-strahler-bound}), is $(\sqrt{8/\pi}+o(1))\sqrt n$
  in probability. One law-of-large-numbers step of this proof is only
  sketched.
\item \emph{Corollary~D} (Corollary~\ref{nw:cor-D}). There is a
  logspace-uniform family of circuits of polynomial size and depth
  $O(\log^3 n)$ that outputs the correct trace or $\bot$ on every instance,
  and outputs $\bot$ with probability $o(1)$ on a uniformly random instance.
\end{itemize}
We also determine Newton's worst case exactly ($\lfloor (n-1)/2\rfloor$ rounds
for $n\ge 3$, Proposition~\ref{nw:prop-worst}), show that random instances have
binary core of size $n-O(\log n)$ and linear live width
(Proposition~\ref{nw:prop-structure}), prove that the density of ones in the
trace tends to $\varphi^{-2}=(3-\sqrt5)/2$ almost surely, where $\varphi$ is
the golden ratio (Theorem~\ref{nw:thm-density}), and derive errorless
average-case $\mathrm{NC}^{k+2}$ and $O(\log^3 n)$-space algorithms
(Corollary~\ref{nw:cor-classes}).

\emph{Outline of the proof of Theorem~A.} Write $E_t$ for the set of wrong
positions after $t$ rounds and $f_t=\min E_t$ for the front. A gate becomes a
new error in round $t+1$ only if both its parents were wrong in round $t$; we
call such gates sources. The first source of round $t$ is exactly $f_{t+1}$
(Lemma~\ref{nw:lem-front}). The second parent $x_{j_i}$ is a fresh uniform
pointer, so the conditional probability that $i$ is a source is (errors so
far)$/i$ whenever $i-1$ is wrong (Lemma~\ref{nw:lem-adapted}). In a short
window past the front, the next round's error mass is at most a constant
times the square of this round's mass divided by $f$, because each source
spreads to $O(1)$ further errors (Lemma~\ref{nw:lem-closure}). In the other
direction, every source of round $r-1$ is witnessed by an error of round $r$
(Lemma~\ref{nw:lem-witness}), so if round $r$ has $m$ errors before position
$y$, round $r-1$ has $O(\sqrt{ym})$ errors there. Before its first source a
round has only $O(\sqrt{f\log f})$ errors (Lemma~\ref{nw:lem-prefront}).
Together these bounds show that, over a window of fixed relative length
$\eps$, the previous round's mass stays at most $f^{1/2+o(1)}$; the current
round then has no source in the window, and the front grows by the factor
$1+\eps$ with probability at least $0.74$ in every round
(Lemma~\ref{nw:lem-window}). A martingale argument over rounds gives
Theorem~A.

The algebraic analogue of the syntactic bound is the analysis of Newton's
method on commutative $\omega$-continuous semirings by Esparza, Kiefer and
Luttenberger~\cite{nw-ekl10}, which relates the iterates to derivation trees
of bounded dimension. Sequential NOR is P-complete under logspace
reductions~\cite{okhotin}; see~\cite{ghr,gb-ladner} for the circuit value
problem in general.

\subsection{Setting and notation}\label{nw:sec-setting}

Throughout this section $x^*$ denotes the trace of the instance, since the
letter $x$ is used for Newton iterates. Thus $x^*_0=1$ and
$x^*_i=(1+x^*_{i-1})(1+x^*_{j_i})$ over $\nwGF$. Gate $i\ge1$ is a NOT gate if
$j_i=i-1$ (then $x^*_i=1+x^*_{i-1}$) and a \emph{binary} gate otherwise.

\emph{Random instances.} The pointers $j_1,j_2,\dots$ are independent, with
$j_i$ uniform on $\{0,\dots,i-1\}$.

\emph{The Newton map.} For $i\ge1$ let
\[
  G_i(y)=(1+y_{i-1})(1+y_{j_i})\ \text{if $i$ is binary},\qquad
  G_i(y)=1+y_{i-1}\ \text{if $i$ is a NOT gate.}
\]
Given $x$ with $x_0=1$, the Newton iterate $x^+$ is the unique solution of
\begin{equation}\label{nw:eq-newton}
  x^+_0=1,\qquad x^+_i=G_i(x)+\sum_k \partial_k G_i(x)\,(x^+_k+x_k)\quad(i\ge1),
\end{equation}
where $\partial_k$ is the formal partial derivative over $\nwGF$. For a binary
gate, $\partial_{i-1}G_i(x)=1+x_{j_i}$ and $\partial_{j_i}G_i(x)=1+x_{i-1}$.
For a NOT gate, $\partial_{i-1}G_i=1$.

With $\delta=x^++x$, equation~\eqref{nw:eq-newton} reads
$(I+J_G(x))\,\delta=x+G(x)$, where $J_G(x)$ is strictly lower triangular. This
is the classical Newton step for $F(y)=y+G(y)=0$. It is computable by forward
substitution, or in depth $O(\log^2 n)$ via
$(I+J)^{-1}=\prod_k(I+J^{2^k})$.

\begin{remark}[NOT-gate convention]\label{nw:rem-not}
A NOT gate is the affine map $x_i=1+x_{i-1}$, with derivative $1$, so Newton
reproduces it exactly. If one wrote it instead as the polynomial
$(1+x_{i-1})^2$, whose formal $\nwGF$ derivative is $0$,
Lemma~\ref{nw:lem-error} would change and the front inequality
$f_{t+1}\ge f_t+2$ of Lemma~\ref{nw:lem-front} would fail.
\end{remark}

\emph{Iterates, errors, fronts.} The start $x^{(0)}$ is either the
\emph{standard start} $(1,0,\dots,0)$ or a random start as in
Remark~\ref{nw:rem-random-start}. We put $x^{(t+1)}=(x^{(t)})^+$ and
$e^{(t)}=x^{(t)}\oplus x^*$, and we use the following notation.
\begin{itemize}[leftmargin=*]
\item the error set $E_t=\{i:e^{(t)}_i=1\}$;
\item the mass $m_t(y)=|E_t\cap[0,y)|$;
\item the front $f_t=\min E_t$, with $f_t=\infty$ if $E_t=\emptyset$;
\item the sources $S_t=\{\,i\ge1: j_i\ne i-1,\ i-1\in E_t,\ j_i\in E_t\,\}$;
\item the filtration $\cF_m=\sigma(j_1,\dots,j_m)$, with any start bits placed
  in $\cF_0$;
\item the compensator $q^{(t)}_i=\PP(i\in S_t\mid\cF_{i-1})$.
\end{itemize}
Intervals are of integers: $[a,b)=\{a,\dots,b-1\}$.

\emph{Infinite embedding and the round count.} We work with the infinite
pointer sequence. By Lemma~\ref{nw:lem-adapted}, the restriction of
$x^{(t)}$ to $[0,n]$ depends only on $j_1,\dots,j_n$, so it coincides with the
Newton iterate of the length-$n$ instance. Define
\[
  T_n=\min\{t:f_t>n\},
\]
the first round that is exact on $[0,n]$.

\emph{Front growth.} For the standard start, $f_0\ge2$, because
$x^*_1=0=x^{(0)}_1$. Lemma~\ref{nw:lem-front} gives $f_{t+1}\ge f_t+2$, hence
$f_t\ge 2t+2$.

\emph{Scale threshold.} Throughout, $y_0\ge e^{200}$ is a scale threshold. In
Theorem~\ref{nw:thm-A} we take $y_0=\max(e^{200},\ln n)$, so $y_0\to\infty$.

\subsection{The syntactic bound and its failure on random instances}
\label{nw:sec-strahler}

Define the Horton--Strahler rank $\rho$ by $\rho(0)=0$, $\rho(i)=\rho(i-1)$
for a NOT gate, and, for a binary gate whose parents have ranks
$r_1\ge r_2$, $\rho(i)=r_1$ if $r_1>r_2$ and $\rho(i)=r_1+1$ if $r_1=r_2$.
Every parent has rank at most that of its child, so $\rho$ is non-decreasing
in $i$.

\begin{proposition}[Newton--Strahler bound]\label{nw:prop-strahler-bound}
For any start, after $t$ Newton rounds every gate $i$ with $\rho(i)<t$ is
exact. In particular $T_n\le\rho(n)+1$.
\end{proposition}

\begin{proof}
Induct on $t$; the case $t=0$ is vacuous. Assume the claim after $t$ rounds.
We show that every gate with $\rho(i)\le t$ is exact after $t+1$ rounds. Let
$D=\{i:\rho(i)\le t\}$, which is downward closed under taking parents.

Take a binary gate $i\in D$. If both of its parents had rank at least $t$,
then $\rho(i)\ge t+1$. So one parent has rank less than $t$, and by induction
it is exact in round $t$. Hence the last term of the recursion of
Lemma~\ref{nw:lem-error} vanishes at every gate of $D$.

On $D$ the recursion of Lemma~\ref{nw:lem-error} is therefore homogeneous
and triangular, with $e^{(t+1)}_0=0$. So $e^{(t+1)}=0$ on $D$.
\end{proof}

\begin{theorem}[Theorem~C: Strahler obstruction]\label{nw:thm-C}
For a uniformly random instance, $\rho(n)/\sqrt n\to\sqrt{8/\pi}$ in
probability.
\end{theorem}

\begin{proof}[Proof (step~(5) is only sketched)]
The law-of-large-numbers step~(5) below is only sketched, not proved; the
theorem therefore rests on that sketch.
\begin{enumerate}[label=(\arabic*),leftmargin=*]
\item Since $\rho$ is non-decreasing, its levels are intervals
  $[a_r,a_{r+1})$.
\item Consider step $i$ in level $r$. A NOT gate keeps the rank, and so does
  a binary gate whose other parent has a lower rank. A binary gate whose other
  parent lies in the current level raises the rank. So the rank increments
  exactly when $j_i\in[a_r,i-2]$.
\item The pointers are fresh. Hence, conditionally on $\cF_{a_r}$ and
  $a_r=s$,
  \[
    \PP(a_{r+1}-a_r>K)=\prod_{k=1}^{K}\frac{s+1}{s+k},
  \]
  since $j_{s+k}$ avoids the $k-1$ values of $[s,s+k-2]$ with probability
  $(s+1)/(s+k)$.
\item As $s\to\infty$, $(a_{r+1}-a_r)/\sqrt s$ converges to the Rayleigh law
  with tail $e^{-x^2/2}$, and
  $\EE[a_{r+1}-a_r\mid a_r=s]=\sqrt{\pi s/2}\,(1+o(1))$.
\item \emph{(Sketch.)} The increments of $\sqrt{a_r}$ have conditional mean
  $\sqrt{\pi/8}+o(1)$ and bounded conditional variance. A martingale law of
  large numbers gives $a_r\sim(\pi/8)r^2$, hence $\rho(n)\sim\sqrt{8n/\pi}$.
  \qedhere
\end{enumerate}
\end{proof}

So the syntactic bound of Proposition~\ref{nw:prop-strahler-bound} gives only
$T_n=O(\sqrt n)$ on random instances. Numerically, $\rho(n)/\sqrt n$ equals
$1.534$, $1.573$ and $1.583$ at $n=10^4,10^5,10^6$; the limit is
$\sqrt{8/\pi}=1.5958$.

\subsection{Tools and basic lemmas}\label{nw:sec-tools}

\begin{lemma}[Counting bounds]\label{nw:lem-counting}
Let $\xi_i\in\{0,1\}$ be adapted to $(\cF_i)$, with
$p_i=\PP(\xi_i=1\mid\cF_{i-1})<1$. Let $I$ be a bounded set of indices such
that each $\nwone\{i\in I\}$ is $\cF_{i-1}$-measurable. Put
$N=\sum_{i\in I}\xi_i$ and $\Lambda=\sum_{i\in I}p_i$. For $x\ge0$ and
$\lambda>0$:
\begin{enumerate}[label=\textup{(T\arabic*)},leftmargin=*]
\item $\PP(N\ge 2\Lambda+2x)\le e^{-x}$.
\item $\PP(N\le\Lambda/2-x)\le e^{-x}$.
\item \textup{(Stopped form.)} Let
  $I'=\{\,i\in I:\sum_{i'\in I,\,i'<i}p_{i'}<\lambda\,\}$, the indices up to
  and including the one where the partial sum of $p$ first reaches $\lambda$.
  Then $\PP(\text{no event in $I'$ and }\Lambda\ge\lambda)\le e^{-\lambda}$.
\end{enumerate}
\end{lemma}

\begin{proof}
For every real $s$, $\EE[e^{s\xi_i}\mid\cF_{i-1}]\le\exp(p_i(e^s-1))$. Hence
$\exp(sN-(e^s-1)\Lambda)$, accumulated along $I$, is a supermartingale with
mean at most $1$.

(T1) Take $s=\ln2$. On $\{N\ge2\Lambda+2x\}$ we have
$N\ln2-\Lambda\ge2x\ln2\ge x$. Apply Markov's inequality.

(T2) Take $s=-1$. On $\{N\le\Lambda/2-x\}$ we have $(1-e^{-1})\Lambda-N\ge x$.
Apply Markov's inequality.

(T3) The process
$\nwone\{\text{no event so far in }I'\}\big/\prod_{i'\in I',\,i'\le i}(1-p_{i'})$
is a martingale with mean $1$. On $\{\Lambda\ge\lambda\}$ the partial sum
reaches $\lambda$ inside $I'$, so $\prod_{I'}(1-p_i)\le e^{-\lambda}$.
Optional stopping at the end of $I'$ gives the bound.
\end{proof}

\begin{lemma}[Error recursion]\label{nw:lem-error}
For every round $t$ and every binary gate $i$,
\[
  e^{(t+1)}_i=(1\oplus x^{(t)}_{j_i})\,e^{(t+1)}_{i-1}\ \oplus\
  (1\oplus x^{(t)}_{i-1})\,e^{(t+1)}_{j_i}\ \oplus\
  e^{(t)}_{i-1}e^{(t)}_{j_i}.
\]
For a NOT gate, $e^{(t+1)}_i=e^{(t+1)}_{i-1}$. Always $e^{(t)}_0=0$.
\end{lemma}

\begin{proof}
Write $x=x^{(t)}$ and $x^+=x^{(t+1)}$, and let $i$ be a binary gate with
$j=j_i$. For any $y$, expanding the bilinear form exactly gives
\[
  G_i(y)=G_i(x)+(1+x_j)(y_{i-1}+x_{i-1})+(1+x_{i-1})(y_j+x_j)
  +(y_{i-1}+x_{i-1})(y_j+x_j).
\]
Take $y=x^*$, so the left side is $x^*_i$, and subtract the defining
equation~\eqref{nw:eq-newton} for $x^+_i$. The NOT gate is affine, so its
expansion has no quadratic term. Position $0$ is fixed to $1$ in every
iterate.
\end{proof}

\begin{lemma}[Front identity]\label{nw:lem-front}
$f_{t+1}=\min S_t$ and $f_{t+1}\ge f_t+2$.
\end{lemma}

\begin{proof}
By Lemma~\ref{nw:lem-error} and induction on $i$, $e^{(t+1)}$ vanishes below
$\min S_t$. At $s=\min S_t$ both propagated terms vanish and the source term
is $1$.

A source $i$ needs $j_i\ge f_t$ and $j_i\le i-2$, so $i\ge f_t+2$.
\end{proof}

\begin{lemma}[Adaptedness and compensator]\label{nw:lem-adapted}
For every $t$ and $i$, the values $x^{(t)}_i$ and $x^*_i$ are
$\cF_i$-measurable. Consequently:
\begin{itemize}[leftmargin=*]
\item $m_t(i)$ and $E_t\cap[0,i)$ are $\cF_{i-1}$-measurable;
\item each $f_t$ is a stopping time;
\item the compensator is
  \[
    q^{(t)}_i=\nwone\{i-1\in E_t\}\,\frac{|E_t\cap[0,i-2]|}{i}.
  \]
\end{itemize}
\end{lemma}

\begin{proof}
Measurability follows by induction on $t$ and then on $i$, since
$x^{(t+1)}_i$ depends only on $x^{(t)}_{\le i}$, $x^{(t+1)}_{<i}$ and $j_i$.

For the compensator: $j_i$ is uniform on $\{0,\dots,i-1\}$ and independent of
$\cF_{i-1}$. A source at $i$ needs $j_i\ne i-1$ and $j_i\in E_t$.
\end{proof}

\begin{lemma}[Stopped closure bound]\label{nw:lem-closure}
Fix the following data:
\begin{itemize}[leftmargin=*]
\item a round $u\ge1$ and positions $f<g$;
\item $\theta\in(0,1)$;
\item a random time $\tau$ with $\{i<\tau\}\in\cF_{i-1}$ such that, for every
  $i\in(f,g)$ with $i<\tau$, $x^{(u-1)}$ has at least $\theta i$ ones in
  $[0,i)$.
\end{itemize}
Define $\nwcl\subseteq\mathbb N$ by $\nwcl\cap[0,f]=E_u\cap[0,f]$. For
$i\in(f,g)$, let $i\in\nwcl$ if and only if
\[
  i\in S_{u-1},\quad\text{or}\quad(i-1\in\nwcl\text{ and }x^{(u-1)}_{j_i}=0),
  \quad\text{or}\quad j_i\in\nwcl.
\]
Then $E_u\cap[0,g)\subseteq\nwcl$. Moreover, with
$G=\prod_{i=f+1}^{g-1}(1+1/(\theta i))\le\exp\bigl((g-f)/(\theta f)\bigr)$,
\[
  \EE\Bigl[\,|\nwcl\cap[f,\min(g,\tau))|\Bigm|\cF_f\Bigr]
  \le(G-1)\,m_u(f)+\frac G\theta
  +\frac G\theta\,\EE\Bigl[\sum_{f<i<g,\ i<\tau}q^{(u-1)}_i\Bigm|\cF_f\Bigr].
\]
\end{lemma}

\begin{proof}
\emph{Inclusion.} By Lemma~\ref{nw:lem-error} for round $u$, if
$e^{(u)}_i=1$ at a binary gate, then $i\in S_{u-1}$, or
($x^{(u-1)}_{j_i}=0$ and $e^{(u)}_{i-1}=1$), or $e^{(u)}_{j_i}=1$. A NOT gate
is covered by $j_i=i-1\in\nwcl$. Induct on $i$.

\emph{Expectation.} Let $a_i=|\nwcl\cap[0,i]|$, $b_i=\nwone\{i\in\nwcl\}$,
$\kappa=(1-\theta)/\theta$ and $\psi_i=a_i+\kappa b_i$. For $f<i<g$ with
$i<\tau$, a union bound over the three ways of joining gives, conditionally on
$\cF_{i-1}$,
\[
  \PP(i\in\nwcl\mid\cF_{i-1})\le q^{(u-1)}_i+(1-\theta)\,b_{i-1}+a_{i-1}/i.
\]
Here we used that $\PP(x^{(u-1)}_{j_i}=0\mid\cF_{i-1})\le1-\theta$ by the
density hypothesis, and that $\PP(j_i\in\nwcl\mid\cF_{i-1})=a_{i-1}/i$. Since
$1+\kappa=1/\theta$,
\[
  \EE[\psi_i\mid\cF_{i-1}]\le\Bigl(1+\frac1{\theta i}\Bigr)\psi_{i-1}
  +\frac{q^{(u-1)}_i}{\theta}.
\]
The start satisfies $\psi_f\le m_u(f)+1+\kappa=m_u(f)+1/\theta$, which counts
position $f$ itself. Iterate up to $\min(g,\tau)-1$ (the factors are
deterministic and at most $G$), then subtract $m_u(f)$.
\end{proof}

\begin{lemma}[Density of ones]\label{nw:lem-density}
Let
$\Omega_6=\{\,\forall m\ge y_0:\ |\{i\in[1,m):x^*_i=1\}|\ge m/3-m^{2/3}\,\}$.
Then
\[
  \PP(\Omega_6^c)\le\epsilon_6(y_0):=\sum_{m\ge y_0}2e^{-m^{1/3}/2}\to0.
\]
On $\Omega_6$ the count is at least $0.33m$, since $y_0\ge3\cdot10^7$.
\end{lemma}

\begin{proof}
A $1$ is always followed by a $0$, so $[0,i)$ contains at least $(i-1)/2$
zeros. Hence
\[
  \PP(x^*_i=1\mid\cF_{i-1},\ x^*_{i-1}=0)\ge\tfrac12-\tfrac1{2i}.
\]
The centred count $D_m$ of $0\to1$ transitions is a martingale with
increments in $[-1,1]$. By Azuma's inequality~\cite{nw-azuma67},
$|D_m|\le m^{2/3}$ except with probability $2e^{-m^{1/3}/2}$.

On that event, with $U_m=|\{i\in[1,m):x^*_i=1\}|$,
\[
  U_m\ge\tfrac12(m-2-U_m)-\tfrac12(1+\ln m)-m^{2/3}.
\]
So $U_m\ge m/3-1-(\ln m)/3-\tfrac23m^{2/3}\ge m/3-m^{2/3}$ for $m\ge30$.
\end{proof}

\subsection{Global good events}\label{nw:sec-global}

\begin{lemma}[Pre-front mass]\label{nw:lem-prefront}
Let
\[
  \Omega_5=\bigl\{\text{for every $t\ge1$ with $f_t\ge y_0$: }
  m_{t-1}(f_t)\le4\sqrt{f_t\ln f_t}+2\bigr\}.
\]
Then $\PP(\Omega_5^c)\le10\,y_0^{-3}$.
\end{lemma}

\begin{proof}
\emph{Setup.} Fix $t$ and $k$ with $2^k>y_0/2$, and put $\lambda_k=4k\ln2$.
Consider the event
\[
  f_t\in[2^k,2^{k+1})\quad\text{and}\quad
  m:=m_{t-1}(f_t)>2+\sqrt{2^{k+2}\lambda_k}.
\]

\emph{Large compensator.} By Lemma~\ref{nw:lem-front}, round $t-1$ has no
source in $[1,f_t-1]$. The interval $[0,f_t-2]$ contains at least $m-1$
errors of round $t-1$. Each pair $a<b$ of them contributes
$1/(b+1)\ge1/f_t$ to $q^{(t-1)}_{b+1}$. Hence
\[
  \sum_{i<f_t}q^{(t-1)}_i\ge\frac{(m-2)^2}{2f_t}\ge\lambda_k.
\]

\emph{Applying \textup{(T3)}.} Take $I=[1,2^{k+1})$ and the stopped set $I'$
of (T3) with $\lambda=\lambda_k$. On the event, the partial sum reaches
$\lambda_k$ at an index at most $f_t-1$. So $I'\subseteq[1,f_t-1]$ contains
no source, while $\Lambda\ge\lambda_k$. By (T3), the event has probability at
most $e^{-\lambda_k}=2^{-4k}$. The unstopped statement would not suffice
here, because $f_t\in I$ is itself a source.

\emph{Union bound.} Rounds with $f_t<2^{k+1}$ satisfy $t<2^k$, since
$f_t\ge2t+2$. So
\[
  \PP(\Omega_5^c)\le\sum_{k:\,2^k>y_0/2}2^k\,2^{-4k}
  \le\tfrac87\,(y_0/2)^{-3}\le10\,y_0^{-3}.
\]

\emph{The stated bound.} Since $2^{k+2}\le4f_t$ and $k\ln2\le\ln f_t$, the
threshold satisfies $2+\sqrt{2^{k+2}\lambda_k}\le2+4\sqrt{f_t\ln f_t}$.
\end{proof}

\begin{lemma}[Witness lemma]\label{nw:lem-witness}
Let $r\ge1$.
\begin{enumerate}[label=\textup{(\roman*)},leftmargin=*]
\item \textup{(Pathwise.)} Every $s\in S_{r-1}$ satisfies $s\in E_r$, or
  $s-1\in E_r$, or $j_s\in E_r$.
\item \textup{(One block.)} For $b\ge8$ let $a=\lceil b/2\rceil$, $B_b=[a,b)$
  and $K=|E_{r-1}\cap[a-1,b-1)|$. With probability at least $1-2b^{-4}$,
  \[
    K\le1+4\sqrt{b\,(m_r(b)+3\ln b)}.
  \]
\item \textup{(Global event and prefix bound.)} Let
  \[
    \Omega_7=\bigcap_{b\ge\sqrt{y_0}}\ \bigcap_{1\le r\le b/2+1}
    \{\text{\textup{(ii)} holds for $(r,B_b)$}\}.
  \]
  Then $\PP(\Omega_7^c)\le3/y_0$. On $\Omega_7$, for every $r\ge1$ and every
  $y\ge y_0$,
  \[
    m_{r-1}(y)\le15\sqrt{y\,(m_r(y)+3\ln y)}.
  \]
\end{enumerate}
\end{lemma}

\begin{proof}
(i) Apply Lemma~\ref{nw:lem-error} with $t=r-1$ at the binary gate $s$. The
source term equals $1$. So if $e^{(r)}_s=0$, one of the two propagated terms
equals $1$.

(ii) Put $x=4\ln b$.
\begin{itemize}[leftmargin=*]
\item \emph{Pointer hits.} Let $H=|\{s\in B_b:j_s\in E_r\}|$. Its compensator
  is $\sum_{s\in B_b}m_r(s)/s\le m_r(b)\ln\frac{b-1}{a-1}\le m_r(b)$, because
  $(b-1)/(a-1)\le7/3$ for $b\ge8$. By (T1), $H\le2m_r(b)+2x$.
\item \emph{Sources.} Let $N=|S_{r-1}\cap B_b|$, with compensator
  $\Lambda=\sum_{s\in B_b}q^{(r-1)}_s$. Each pair $a'<b'$ in
  $E_{r-1}\cap[a-1,b-1)$ has $b'+1\in B_b$ and $a'\le(b'+1)-2$, so it adds at
  least $1/b$ to $q^{(r-1)}_{b'+1}$. Hence $\Lambda\ge K(K-1)/(2b)$, and by
  (T2), $N\ge\Lambda/2-x$.
\item \emph{Witnessing.} By (i),
  $N\le|E_r\cap B_b|+|E_r\cap(B_b-1)|+H\le2m_r(b)+H$.
\end{itemize}
Combining, $K(K-1)/(4b)\le N+x\le4m_r(b)+3x$, so
$(K-1)^2\le16b\,(m_r(b)+3\ln b)$. The two failure events each have
probability $e^{-x}=b^{-4}$.

(iii), \emph{the global event.} By (ii),
$\PP(\Omega_7^c)\le\sum_{b\ge\sqrt{y_0}}(b/2+1)\,2b^{-4}\le3/y_0$. For
$r>b/2+1$ we have $f_{r-1}\ge2(r-1)\ge b$, so $K=0$ and (ii) holds trivially.

(iii), \emph{the prefix bound.} Put $b_0=y$ and $b_{k+1}=\lceil b_k/2\rceil$,
and stop at the first $K'$ with $b_{K'}<\sqrt{y_0}$. The shifted blocks
$[b_{k+1}-1,b_k-1)$ for $k<K'$ tile $[b_{K'}-1,y-1)$. The rest is
$[0,b_{K'}-1)$ together with $\{y-1\}$. Hence, writing $K_k$ for the quantity
$K$ of (ii) for the block $B_{b_k}$,
\[
  m_{r-1}(y)\le\sqrt{y_0}+1+\sum_{k<K'}K_k.
\]
By (ii), $K_k\le1+4\sqrt{b_k}\sqrt{m_r(y)+3\ln y}$. Also
$\sqrt{b_k}\le\sqrt y\,2^{-k/2}+1$, and there are at most $\log_2y+1$ blocks.
So
\[
  \sum_kK_k\le(\log_2y+1)+4\,(3.415\sqrt y+\log_2y+1)\sqrt{m_r(y)+3\ln y}.
\]
For $y\ge2^{20}$ we have $\log_2y+1\le0.021\sqrt y$. Using
$m_r+3\ln y\ge1$ and $\sqrt{y_0}\le\sqrt y$, the total is at most
$(1.03+13.8)\sqrt{y(m_r(y)+3\ln y)}\le15\sqrt{y(m_r(y)+3\ln y)}$.
\end{proof}

\begin{theorem}[Theorem~B: lower bound]\label{nw:thm-B}
For the standard start, on $\Omega_6\cap\Omega_7$ with $n+1\ge y_0$, we have
$T_n\ge\log_2\ln n-3$. With $y_0=\max(e^{200},\ln n)$, this holds with
probability $1-o(1)$.
\end{theorem}

\begin{proof}
Let $y=n+1$, $T=T_n$ and $\mu_k=m_{T-k}(y)/y$. Then $\mu_0=0$.

By Lemma~\ref{nw:lem-witness}(iii), $\mu_1\le\eta:=15\sqrt{3\ln y/y}$.
Define
\[
  U_1=\eta,\qquad U_k=15\sqrt{U_{k-1}+3\ln y/y}\quad(k\ge2).
\]
Then $\mu_k\le U_k$. Since $U_{k-1}\ge\eta\ge3\ln y/y$, we get
$U_k\le\sqrt{450\,U_{k-1}}$, and induction gives
$U_k\le450\,\eta^{2^{1-k}}$.

By Lemma~\ref{nw:lem-density}, $\mu_T=m_0(y)/y\ge0.33$. So
$0.33\le450\,\eta^{2^{1-T}}$, that is, $2^{T-1}\ge\ln(1/\eta)/\ln1364$. For
$\ln y\ge200$,
\[
  \ln(1/\eta)=\tfrac12\ln y-\tfrac12\ln\ln y-\tfrac12\ln3-\ln15\ge0.47\ln y.
\]
Hence $2^{T-1}\ge0.0651\ln y$, i.e.\ $T\ge\log_2\ln y-2.94$.

For the probability statement: $\PP(\Omega_6^c\cup\Omega_7^c)\le
\epsilon_6(y_0)+3/y_0\to0$ by Lemmas~\ref{nw:lem-density}
and~\ref{nw:lem-witness}, and $n+1\ge y_0$ for all large $n$.
\end{proof}

\subsection{The window lemma}\label{nw:sec-window}

The constants of this subsection are collected in the following table. To
avoid a clash with the block endpoint $b$ and block start $a$ of
Lemma~\ref{nw:lem-witness}, the quadratic coefficient and the linear-phase
growth factor are written $\bar b$ and $\bar a$.

\begin{center}
\small
\begin{tabular}{@{}lll@{}}
\toprule
Symbol & Value & Role\\
\midrule
$\theta$ & $1/5$ & density of ones of every linearisation point used\\
$\beta$ & $1/20$ & threshold defining the bottom level\\
$\lambda$ & $450$ & constant in $\Pi_{\ell+1}^2\le\lambda\Pi_\ell$\\
$\eps$ & $e^{-100}$ & relative window length\\
$\bar b$ & $2500$ & quadratic coefficient of the $c$-recursion\\
$\bar a$ & $2\bar b+12=5012$ & linear-phase growth factor\\
$\gamma_0$ & $\lambda/\beta^2=180000$ & bound on the bottom ratio $c_L/\eps$\\
$M_1$ & $800$ & Markov factor in the linear phase\\
$k_1$ & $5$ & number of linear-phase levels\\
\bottomrule
\end{tabular}
\end{center}

\begin{lemma}[Window lemma]\label{nw:lem-window}
Fix a round $t\ge1$, and write $f=f_t$ and $g=f+\lfloor\eps f\rfloor$. Let
$\mathcal G_t\in\cF_f$ be the event that all of the following hold:
\begin{itemize}[leftmargin=*]
\item $y_0\le f<\infty$;
\item the inequalities defining $\Omega_5$ hold at $f$;
\item those defining $\Omega_6$ hold for all $m\in[y_0,f]$;
\item those defining $\Omega_7$ hold for all blocks with $b\le f$;
\item \textup{(S)} $m_0(f)\ge f/4$.
\end{itemize}
For the standard start, \textup{(S)} follows from $\Omega_6$. On
$\mathcal G_t$,
\[
  \PP\bigl(S_t\cap(f,g)\ne\emptyset\bigm|\cF_f\bigr)\le0.26,\qquad\text{hence}
  \qquad\PP\bigl(f_{t+1}\ge f+\lfloor\eps f\rfloor\bigm|\cF_f\bigr)\ge0.74.
\]
Moreover, with conditional probability at least $0.74$ we have
\[
  |E_{t-1}\cap[f,g)|\le5f^{1/2+1/16}\sqrt{\ln f}.
\]
\end{lemma}

The proof occupies the rest of this subsection. The second inequality
follows from the first by Lemma~\ref{nw:lem-front}, since sources are at
least $f+2$ and $f_{t+1}=\min S_t$.

\emph{Formal setup at a random front.} $f=f_t$ is a stopping time. On
$\{f_t=\varphi\}$, conditional expectations given $\cF_{f_t}$ coincide with
those given $\cF_\varphi$. So it suffices to prove every bound below for a
deterministic $\varphi$ on $\{f_t=\varphi\}\cap\mathcal G_t$. All the
quantities $X_\ell$, $L$, $g$ and $M_\ell$ defined below are deterministic
functions of $f$, and $\mathcal G_t\in\cF_f$.

Lemma~\ref{nw:lem-closure} requires its density hypothesis to hold surely. So
off $\mathcal G_t$ we replace every stopping time $\tau_\ell$ below by $f+1$,
an $\cF_f$-measurable choice. This changes nothing on $\mathcal G_t$, which
is the only place where the conclusion is claimed.

\subsubsection{Step 1: canonical pre-window bounds}
Define
\[
  \Pi_1=5\sqrt{\ln f/f},\qquad\Pi_{\ell+1}=15\sqrt{\Pi_\ell+3\ln f/f},
\]
and let $p_\ell=m_{t-\ell}(f)$. On $\mathcal G_t$ the following hold.
\begin{enumerate}[label=(P\arabic*),leftmargin=*]
\item $p_\ell\le\Pi_\ell f$ for $1\le\ell\le t$.
  \begin{itemize}
  \item For $\ell=1$: $\Omega_5$ gives
    $p_1\le4\sqrt{f\ln f}+2\le5\sqrt{f\ln f}$.
  \item For the induction step: Lemma~\ref{nw:lem-witness}(iii) with
    $r=t-\ell$ and $y=f$ gives
    $p_{\ell+1}\le15\sqrt{f(p_\ell+3\ln f)}\le\Pi_{\ell+1}f$.
  \end{itemize}
\item Properties of the canonical sequence.
  \begin{itemize}
  \item \emph{Increasing.} We have
    $\Pi_{\ell+1}\ge15\sqrt{\Pi_\ell}\ge\Pi_\ell$ whenever
    $\Pi_\ell\le225$. The sequence never exceeds $226$, because
    $\Pi_\ell\le226$ implies
    \[
      \Pi_{\ell+1}\le15\sqrt{226+3\ln f/f}<226.
    \]
    So $\Pi$ is increasing, and $\Pi_\ell\ge\Pi_1$ for all $\ell$.
  \item \emph{Quadratic relation.} Since $3\ln f/f\le\Pi_1\le\Pi_\ell$, we get
    $\Pi_{\ell+1}^2=225(\Pi_\ell+3\ln f/f)\le\lambda\Pi_\ell$.
  \item \emph{Two-sided bounds.} Since $\Pi_{\ell+1}\ge\sqrt{\Pi_\ell}$
    always, induction gives
    $\Pi_1^{2^{1-\ell}}\le\Pi_\ell\le\lambda\,\Pi_1^{2^{1-\ell}}$.
  \item \emph{Sharper lower bound.} From
    $\Pi_{\ell+1}\ge15\sqrt{\Pi_\ell}$, induction gives
    \[
      \Pi_\ell\ge225\,(\Pi_1/225)^{2^{1-\ell}}.
    \]
  \end{itemize}
\item \emph{The bottom level.} Let
  $\ell_b=\max\{\ell\le t-1:\Pi_{\ell+1}\le\beta\}$.
  \begin{itemize}
  \item \emph{Not too deep.} By (S) and (P1), $\Pi_t\ge1/4>\beta$. Since $\Pi$
    is increasing, every $\ell$ in the set has $\ell+1<t$. Hence
    $\ell_b\le t-2$ and $\Pi_{\ell_b+2}>\beta$, so
    $\Pi_{\ell_b+1}\ge\Pi_{\ell_b+2}^2/\lambda>\beta^2/\lambda=1/\gamma_0$.
  \item \emph{Not empty.} $\Pi_2\le\sqrt{\lambda\Pi_1}\le\beta$ for
    $f\ge e^{200}$. Hence $t\ge3$ and $\ell_b\ge1$.
  \item \emph{Depth bound (a).} From
    $\Pi_1^{2^{-\ell_b}}\le\Pi_{\ell_b+1}\le\beta$ we get
    $2^{-\ell_b}\ln(1/\Pi_1)\ge\ln20$.
  \item \emph{Depth bound (b).} From
    \[
      225(\Pi_1/225)^{2^{-\ell_b}}\le\Pi_{\ell_b+1}\le1/20
    \]
    we get $2^{-\ell_b}\ln(225/\Pi_1)\ge\ln4500$. Since
    \[
      \ln(225/\Pi_1)=\ln45+\tfrac12\ln f-\tfrac12\ln\ln f,
    \]
    we obtain
    \[
      2^{\ell_b}\le\frac{\ln(225/\Pi_1)}{\ln4500}\le\frac{\ln f}{16}
      \qquad(\ln f\ge151).
    \]
  \end{itemize}
\item The rounds $t-\ell$ with $0\le\ell\le\ell_b+1$ are at least $1$. Also
  $|\{i<f:x^*_i=1\}|\ge0.33f$ by $\Omega_6$.
\end{enumerate}

\subsubsection{Step 2: the stopped levels}
Let $L=\ell_b+1$. For $0\le\ell\le L$ and $i\in(f,g]$ let
$w_\ell(i)=|E_{t-\ell}\cap[f,i)|$, and let
\[
  \tau_\ell=\min\{\,i\in(f,g]:w_\ell(i)>X_\ell f\,\}\quad
  (\infty\text{ if there is none}),
\]
with thresholds $X_\ell$ defined below. Since $w_\ell(i)$ is
$\cF_{i-1}$-measurable, so is $\{i<\tau_\ell\}$. Set $X_L=\eps$, so that
$\tau_L=\infty$ surely, because $w_L(g)\le g-f\le\eps f$. Set
$X_0=f^{1/3}/f$.

\begin{lemma}[Level inequality]\label{nw:lem-level}
Let $0\le\ell\le L-1$ and suppose $X_{\ell+1}\le\beta$. Put $\Pi_0=0$,
$G=e^{\eps/\theta}$, and
\[
  R_\ell=(G-1)\Pi_\ell f+\frac G\theta
  +\frac G\theta\,(X_{\ell+1}f+1)(\Pi_{\ell+1}+X_{\ell+1}).
\]
Then $\EE\bigl[w_\ell(\min(g,\tau_{\ell+1}))\bigm|\cF_f\bigr]\le R_\ell$.
\end{lemma}

\begin{proof}
Apply Lemma~\ref{nw:lem-closure} with $u=t-\ell$, the window $(f,g)$,
$\theta=1/5$ and $\tau=\tau_{\ell+1}$.

\emph{Density.} For $f<i<\tau_{\ell+1}$, the vector $x^{(t-\ell-1)}$ differs
from $x^*$ on $[0,i)$ in at most $p_{\ell+1}+w_{\ell+1}(i)\le2\beta f$
positions. So it has at least $0.33f-0.1f=0.23f\ge0.2(1+\eps)f\ge\theta i$
ones in $[0,i)$.

\emph{Compensator.} For $f<i<\min(g,\tau_{\ell+1})$,
\[
  q^{(t-\ell-1)}_i\le\nwone\{i-1\in E_{t-\ell-1}\}\,
  \frac{p_{\ell+1}+X_{\ell+1}f}{f}.
\]
At most $X_{\ell+1}f+1$ such $i$ have $i-1\in E_{t-\ell-1}$. So the
compensator sum is at most $(X_{\ell+1}f+1)(\Pi_{\ell+1}+X_{\ell+1})$.

\emph{Conclusion.} Finally $w_\ell\le|\nwcl\cap[f,\cdot)|$ and
$m_{t-\ell}(f)=p_\ell\le\Pi_\ell f$.
\end{proof}

\emph{Markov step.} For $1\le\ell\le L-1$ set $X_\ell=M_\ell R_\ell/f$, with
$M_\ell$ fixed in Step~3. Then
\[
  \PP\bigl(w_\ell(\min(g,\tau_{\ell+1}))>X_\ell f\bigm|\cF_f\bigr)
  \le1/M_\ell.
\]
Since $\{\tau_\ell<\infty\}\subseteq\{\tau_{\ell+1}<\infty\}\cup
\{w_\ell(\min(g,\tau_{\ell+1}))>X_\ell f\}$, the following holds provided
$X_k\le\beta$ for all $k\in[\ell+1,L]$:
\begin{equation}\label{nw:eq-tau}
  \PP(\tau_\ell<\infty\mid\cF_f)\le\sum_{k=\ell}^{L-1}1/M_k.
\end{equation}

\subsubsection{Step 3: the $c$-recursion}
For $1\le\ell\le L$ let $c_\ell=X_\ell/\Pi_\ell$, and let $\nu=1/(f\Pi_1)$.

\emph{Constants from $G$.} Since $\eps/\theta\le0.1$, we have
$G-1\le1.06\,\eps/\theta=5.3\eps$ and $G/\theta\le5.53$.

\emph{Dividing $R_\ell$ by $f\Pi_\ell$.} The middle term is at most
$5.53\nu$. Moreover $X_{\ell+1}+1/f\le\Pi_{\ell+1}(c_{\ell+1}+\nu)$ and
$\Pi_{\ell+1}+X_{\ell+1}=\Pi_{\ell+1}(1+c_{\ell+1})$. Using (P2)
($\Pi_{\ell+1}^2\le\lambda\Pi_\ell$) and $5.53\cdot450\le\bar b$,
\[
  c_\ell\le M_\ell\bigl[5.3\eps+5.53\nu+\bar b\,(c_{\ell+1}+\nu)(1+c_{\ell+1})
  \bigr].
\]
For $f\ge e^{200}$, $f\Pi_1=5\sqrt{f\ln f}\ge e^{100}$, so $\nu\le\eps$ and
\begin{equation}\label{nw:eq-crec}
  c_\ell\le M_\ell\bigl[11\eps+\bar b\,(c_{\ell+1}+\eps)(1+c_{\ell+1})\bigr].
\end{equation}
At the bottom, $c_L=\eps/\Pi_L\le\gamma_0\eps$ by (P3).

\emph{Allocation of failure probabilities.} Index levels by $k=L-\ell$ for
$k=1,\dots,L-1$. Put $M_\ell=M_1=800$ if $k\le k_1=5$, and
$M_\ell=2^{(k-k_1)+2}$ if $k>k_1$. Then
\[
  \sum1/M\le5/800+\sum_{s\ge1}2^{-s-2}\le0.257.
\]

\emph{Phase~1 ($k\le k_1$): $c\le1$.} If $c_{\ell+1}\le1$, then
\eqref{nw:eq-crec} gives
$c_\ell+\eps\le M_1(2\bar b+12)(c_{\ell+1}+\eps)=\bar aM_1(c_{\ell+1}+\eps)$.
Hence $c_{L-k}+\eps\le(\bar aM_1)^k(\gamma_0+1)\eps$. At $k=k_1$,
\[
  5\ln(5012\cdot800)+\ln180001-100=76.0+12.1-100<0,
\]
so $c\le1$ throughout Phase~1.

\emph{Phase~2 ($k=k_1+s$, $s\ge1$): $\log c$ at most doubles per level.} Let
$D=\max(1,c)$. For $D\ge1$ we have $(D+\eps)(1+D)\le2.01D^2$ and
$11\eps\le0.01\,\bar bD^2$, so $D_\ell\le4\bar bM_\ell D_{\ell+1}^2$. Starting
from $D=1$ at the end of Phase~1,
\[
  \ln D\text{ at step }s\le\sum_{s'=1}^{s}2^{s-s'}\bigl(\ln(16\bar b)
  +s'\ln2\bigr)\le2^s\ln(64\bar b)\le12\cdot2^s.
\]
The number $Q$ of Phase-2 levels satisfies
$2^Q\le2^{\ell_b-k_1}\le(\ln f/16)/32$ by (P3)(b). Therefore
\begin{equation}\label{nw:eq-c1}
  \ln c_1\le12\cdot2^Q\le\tfrac{3}{128}\ln f,\qquad\text{i.e.}\qquad
  c_1\le f^{3/128}\le f^{1/16}.
\end{equation}

\emph{The proviso $X_\ell\le\beta$ at every level.}
\begin{itemize}[leftmargin=*]
\item At the bottom, $X_L=\eps$.
\item In Phase~1, $X_\ell=c_\ell\Pi_\ell\le\Pi_\ell\le\Pi_L\le\beta$.
\item In Phase~2, at level $\ell=L-k_1-s$, we have
  $2^{1-\ell}=2^{k_1+s-\ell_b}$. By (P3)(a),
  $2^{-\ell_b}\ln(1/\Pi_1)\ge\ln20$, so, using
  $X_\ell=c_\ell\Pi_\ell\le c_\ell\lambda\Pi_1^{2^{1-\ell}}$,
  \[
    \ln X_\ell\le12\cdot2^s+\ln\lambda-2^{k_1+s}\ln20
    \le2^s(12-95.8)+6.11<\ln\beta.
  \]
\end{itemize}
The $X_\ell$ are deterministic functions of $f$, defined bottom-up, and the
argument above is pure algebra, so it involves no circularity. The
verification scripts recheck this arithmetic numerically at $11{,}202$ values
of $\ln f$ between $200$ and $10^{20}$; this is a check, not part of the
proof.

\subsubsection{Step 4: the top level and the failure budget}
\emph{Top closure.} We have $E_t\cap[0,f)=\emptyset$, $\Pi_1\le\beta$ and
$X_1\le\beta$. Lemma~\ref{nw:lem-level} with $\ell=0$ (so $p_0=0$) gives
\[
  \EE[w_0(\min(g,\tau_1))\mid\cF_f]\le R_0
  =\frac G\theta+\frac G\theta(X_1f+1)(\Pi_1+X_1).
\]
With $X_1=c_1\Pi_1\le5f^{1/16}\sqrt{\ln f/f}$ we get
$X_1f+1\le6f^{1/16}\sqrt{f\ln f}$ and
$\Pi_1+X_1\le10f^{1/16}\sqrt{\ln f/f}$. Hence
\[
  R_0\le5.53\,(1+60f^{1/8}\ln f)\le340f^{1/8}\ln f.
\]

\emph{Sources of round $t$ in the window.} Let $B=f^{1/3}$ and
$\tau_0=\min\{i\in(f,g]:w_0(i)>B\}$. Then
\[
  \PP(S_t\cap(f,g)\ne\emptyset\mid\cF_f)\le\PP(\tau_0<\infty\mid\cF_f)
  +\EE\Bigl[\sum_{f<i<g,\ i<\tau_0}q^{(t)}_i\Bigm|\cF_f\Bigr].
\]
For $i<\tau_0$, at most $B$ positions $i$ have $i-1\in E_t$, and each has
$|E_t\cap[0,i-2]|\le B$. So the second term is at most $B^2/f=f^{-1/3}$. By
\eqref{nw:eq-tau} and Markov's inequality,
\[
  \PP(\tau_0<\infty\mid\cF_f)\le\PP(\tau_1<\infty\mid\cF_f)+R_0/B
  \le0.257+340f^{-5/24}\ln f.
\]
For $f\ge e^{200}$ the total is at most $0.26$. On the event
$\{\tau_1=\infty\}$, which has conditional probability at least $0.74$, we
also get $|E_{t-1}\cap[f,g)|\le X_1f\le5f^{1/2+1/16}\sqrt{\ln f}$. This
proves Lemma~\ref{nw:lem-window}.\qed

\subsection{Proof of Theorem A}\label{nw:sec-main}

\begin{theorem}[Theorem~A: upper bound]\label{nw:thm-A}
For the standard start, let $y_0=\max(e^{200},\ln n)$. Then, with probability
$1-o(1)$,
\[
  T_n\le\tfrac12y_0+\lceil4.1\,e^{100}\ln n\rceil=O(\log n).
\]
\end{theorem}

\begin{proof}
\emph{Good event.} Let $\mathfrak G=\Omega_5\cap\Omega_6\cap\Omega_7$. Since
$y_0\to\infty$,
\[
  \PP(\mathfrak G^c)\le10y_0^{-3}+\epsilon_6(y_0)+3/y_0\to0.
\]
On $\mathfrak G$, (S) holds and $\mathcal G_t$ holds for every $t$ with
$y_0\le f_t<\infty$.

\emph{Burn-in.} Since $f_t\ge2t+2$, $T_0:=\min\{t:f_t\ge y_0\}\le y_0/2$.
Also $\{T_0\le t\}=\{f_t\ge y_0\}\in\cF_{f_t}$.

\emph{Successes.} Let $\nwsucc_t=\{f_{t+1}\ge f_t+\lfloor\eps f_t\rfloor\}$
if $f_t<\infty$, and the sure event if $f_t=\infty$. It is
$\cF_{f_{t+1}}$-measurable. By Lemma~\ref{nw:lem-window},
$\PP(\nwsucc_t\mid\cF_{f_t})\ge0.74\,\nwone_{\mathcal G_t}$.

\emph{Azuma over a predictable window.} Let $T'=\lceil4.1e^{100}\ln n\rceil$
and $w_t=\nwone\{T_0\le t<T_0+T'\}$. It is $\cF_{f_t}$-measurable, since
$\{t<T_0+T'\}=\{f_{t-T'}<y_0\}\in\cF_{f_t}$. Put
$\zeta_t=\nwone_{\nwsucc_t}-0.74\,\nwone_{\mathcal G_t}$. Then:
\begin{itemize}[leftmargin=*]
\item the increments $w_t\zeta_t$ have nonnegative conditional mean given
  $\cF_{f_t}$;
\item they lie in $[-1,1]$, with predictable bounds $w_t$ satisfying
  $\sum_tw_t^2\le T'$.
\end{itemize}
By the Azuma--Hoeffding inequality~\cite{nw-azuma67,hoeffding,nw-freedman75},
$\sum_tw_t\zeta_t\ge-T'/4$ except with probability $e^{-T'/32}$.

\emph{Counting.} On $\mathfrak G$, $\nwone_{\mathcal G_t}=1$ at every window
round with $f_t<\infty$. So either the iteration terminates inside the
window, or the window contains at least $0.74T'-T'/4=0.49T'$ successes.

\emph{Growth.} Each success multiplies the front by at least $1+\eps/2$, since
$\lfloor\eps f\rfloor\ge\eps f/2$ when $\eps f\ge2$. Since
$\ln(1+\eps/2)\ge0.4999\eps$,
\[
  0.49\,T'\ln(1+\eps/2)\ge0.49\cdot4.1\cdot0.4999\,\ln n>\ln n.
\]
So $f_{T_0+T'}>n$, i.e.\ $T_n\le T_0+T'$. The failure probability is at most
$\PP(\mathfrak G^c)+e^{-T'/32}=o(1)$.
\end{proof}

\begin{remark}[Random start]\label{nw:rem-random-start}
Let $x^{(0)}_0=1$, and let $x^{(0)}_i$ ($i\ge1$) be i.i.d.\ uniform bits
independent of the pointers, placed in $\cF_0$. Theorem~\ref{nw:thm-A} holds
with the following changes.
\begin{itemize}[leftmargin=*]
\item \emph{Front growth.} Only $f_0\ge1$ is guaranteed, so $f_t\ge2t+1$. The
  union bounds still hold. In Lemma~\ref{nw:lem-prefront}, $f_t<2^{k+1}$
  forces $t<2^k$. In $\Omega_7$, $r>b/2+1$ forces $f_{r-1}\ge2r-1>b$, so
  $K=0$. The burn-in becomes $T_0\le(y_0+1)/2$.
\item \emph{Condition \textup{(S)}.} Given the pointers, the indicators
  $\nwone\{x^{(0)}_i\ne x^*_i\}$ are independent Bernoulli$(1/2)$ variables.
  By Chernoff's bound, $\PP(m_0(f)<f/4)\le e^{-f/16}$, so (S) holds for all $f\ge y_0$
  simultaneously with probability $1-o(1)$.
\item \emph{Unchanged parts.} Lemmas~\ref{nw:lem-error}, \ref{nw:lem-front},
  \ref{nw:lem-closure} and~\ref{nw:lem-witness}(i) are pathwise and do not
  involve the start.
\end{itemize}
\end{remark}

\begin{corollary}[Corollary~D: errorless average-case NC]\label{nw:cor-D}
Consider the logspace-uniform circuit family that:
\begin{enumerate}[label=(\arabic*),leftmargin=*]
\item runs $\lceil y_0/2+4.1e^{100}\ln n\rceil$ Newton rounds, each a
  unit-lower-triangular $\nwGF$ solve of depth $O(\log^2n)$;
\item checks every gate equation $x_i=\NOR(x_{i-1},x_{j_i})$;
\item outputs the trace if all checks pass, and $\bot$ otherwise.
\end{enumerate}
It has polynomial size and depth $O(\log^3n)$. It never outputs a wrong trace,
and by Theorem~\ref{nw:thm-A} it outputs $\bot$ with probability $o(1)$ on a
uniformly random instance.
\end{corollary}

\begin{proof}
Each round has depth $O(\log^2 n)$ and there are $O(\log n)$ rounds; the
check has depth $O(\log n)$. A trace that passes every gate check is the
trace, so the output is never wrong. Since $F(x^*)=0$, the Newton step at
$x^*$ is $\delta=0$. So $x^*$ is a fixed point of the Newton map, and any
number of rounds $R\ge T_n$ returns $x^*$; by Theorem~\ref{nw:thm-A} this
happens with probability $1-o(1)$. In particular the round count may be
replaced by the logspace-computable bound
$R=\lceil e^{200}\rceil+\lceil4.1e^{100}+1\rceil\lceil\log_2n\rceil$.
\end{proof}

\subsection{Newton's worst case}\label{nw:sec-worst}

In this subsection the instance is fixed and of length $n$. The sets $E_t$,
the fronts $f_t$ (with $\min\emptyset=\infty$) and $T_n=\min\{t:f_t>n\}$ are
defined for it as in Subsection~\ref{nw:sec-setting}.

\begin{proposition}[Newton's worst case]\label{nw:prop-worst}
\leavevmode
\begin{enumerate}[label=\textup{(\alph*)},leftmargin=*]
\item For every start with $x^{(0)}_0=1$, $f_t\ge f_0+2t$. From the standard
  start moreover $f_1\ge5$, so $f_t\ge2t+3$ for $t\ge1$, and
  $T_n\le\lfloor(n-1)/2\rfloor$ for every $n\ge3$.
\item For $n\ge5$ let $W_n$ be the instance with
  $(j_1,\dots,j_5)=(0,1,0,1,2)$, $j_i=1$ for even $i\ge6$, and $j_i=i-2$ for
  odd $i\ge7$. From the standard start, $T_n(W_n)=\lfloor(n-1)/2\rfloor$.
\item Hence, for every $n\ge3$, the maximum of $T_n$ over all $n!$ instances,
  from the standard start, is exactly $\lfloor(n-1)/2\rfloor$.
\end{enumerate}
\end{proposition}

\begin{proof}
(a) The first claim is Lemma~\ref{nw:lem-front}. From the standard start,
$E_0=\{i\ge2:x^*_i=1\}$, since $x^*_0=1$ and $x^*_1=0$. A source $i\in S_0$ is
a binary gate with $x^*_{i-1}=x^*_{j_i}=1$ and $2\le j_i\le i-2$. Two ones are
never adjacent in the trace, so $j_i\ne i-2$ and $i\ge j_i+3\ge5$. By
Lemma~\ref{nw:lem-front}, $f_1=\min S_0\ge5$, and then $f_t\ge2t+3$ for
$t\ge1$. If $n\ge3$, then $t=\lfloor(n-1)/2\rfloor\ge1$ and
$f_t\ge2\lfloor(n-1)/2\rfloor+3>n$, so $T_n\le t$.

(b) By induction, $x^*_i=1$ for even $i$ and $x^*_i=0$ for odd $i$. Each
Newton round is triangular, so the iterates of $W_n$ are the restrictions to
$[0,n]$ of those of the infinite instance $W_\infty$ defined by the same rule;
we work with $W_\infty$. By (a), $f_t\ge2$, so $e^{(t)}_1=0$ and
$x^{(t)}_1=0$ in every round.

\emph{Copy rule.} For even $i\ge4$ the pointer is $j_i=1$, so
Lemma~\ref{nw:lem-error} gives $e^{(t+1)}_i=e^{(t+1)}_{i-1}$ in every round.
The same holds at the NOT gate $i=2$.

\emph{Odd rule.} Let $i\ge7$ be odd. By Lemma~\ref{nw:lem-error} with
$j_i=i-2$ and the copy rule $e^{(t+1)}_{i-1}=e^{(t+1)}_{i-2}$,
\[
  e^{(t+1)}_i=(x^{(t)}_{i-2}\oplus x^{(t)}_{i-1})\,e^{(t+1)}_{i-2}
  \oplus e^{(t)}_{i-1}e^{(t)}_{i-2}.
\]
For $t\ge1$ the copy rule in round $t$ gives $e^{(t)}_{i-1}=e^{(t)}_{i-2}$, so
$x^{(t)}_{i-2}\oplus x^{(t)}_{i-1}=x^*_{i-2}\oplus x^*_{i-1}=1$ and
\[
  e^{(t+1)}_i=e^{(t+1)}_{i-2}\oplus e^{(t)}_{i-2}\qquad(t\ge1).
\]

\emph{Round 1.} Here $x^{(0)}_{i-2}=x^{(0)}_{i-1}=0$ and $e^{(0)}_{i-2}=0$ for
odd $i\ge7$ (odd positions are correct at the start), so the display gives
$e^{(1)}_i=0$, and by the copy rule $e^{(1)}_{i+1}=0$. Position $5$ is a
source of round $0$, since $x^*_4=x^*_2=1$; by (a), $f_1=5$. With the copy
rule at $6$, $E_1=\{5,6\}$.

\emph{Generating function.} For $t\ge1$ put
$A^{(t)}(z)=\sum_{k\ge2}e^{(t)}_{2k+1}z^k$, a formal power series over
$\nwGF$. For $t\ge2$, (a) gives $f_t\ge7$, so $e^{(t)}_5=0$. Multiplying the
odd rule by $z^k$ and summing over $k\ge3$ gives, for $t\ge1$,
\[
  A^{(t+1)}=\frac{z}{1+z}A^{(t)},\qquad A^{(1)}=z^2,\qquad\text{hence}\qquad
  A^{(t)}=z^{t+1}(1+z)^{-(t-1)}.
\]
The lowest term is $z^{t+1}$, and even positions copy their odd predecessors.
So $f_t=2t+3$ for $t\ge1$, while $f_0=2$. Hence
$T_n(W_n)=\min\{t\ge1:2t+3>n\}=\lfloor(n-1)/2\rfloor$ for $n\ge5$.

(c) For $n\ge5$ combine (a) and (b). For $n\in\{3,4\}$ the bound is $1$, and
it is attained by any instance with $j_2=1$, since then
$x^*_2=1\ne x^{(0)}_2$.
\end{proof}

An exhaustive search over all instances with $n\le13$ in the verification
scripts agrees with (c). So the number of Newton rounds is $\Theta(n)$ in the
worst case, exactly $\lfloor(n-1)/2\rfloor$, and $O(\log n)$ on random
instances. The family $W_n$ is itself trivial, since its trace simply
alternates; the number of Newton rounds measures the iteration, not the
difficulty of the instance.

\subsection{Structure of random instances and the density of ones}
\label{nw:sec-structure}

For a parent array $j$ define the \emph{binary core}
\[
  B(j)=\{\,i\ge2:1\le j_i\le i-2,\ j_{i-1}\ne0\,\}
\]
and the \emph{live width}
\[
  L(j)=\max_c\,|\{\,p\le c:\exists\,i>c,\ j_i=p\,\}|,
\]
where $i$ ranges over $[1,n]$.

\begin{proposition}[Structural parameters]\label{nw:prop-structure}
For a uniformly random instance, the following hold with probability
$1-o(1)$.
\begin{enumerate}[label=\textup{(\alph*)},leftmargin=*]
\item The Horton--Strahler rank is $(\sqrt{8/\pi}+o(1))\sqrt n$
  \textup{(}Theorem~\ref{nw:thm-C}, whose law-of-large-numbers step is only
  sketched\textup{)}.
\item $|B(j)|\ge n-7\ln n$.
\item $L(j)\ge n/4-n^{2/3}$.
\end{enumerate}
\end{proposition}

\begin{proof}
(a) This is Theorem~\ref{nw:thm-C}.

(b) A position $i\in[1,n]$ lies outside $B(j)$ if and only if $i=1$, or
$j_i\in\{0,i-1\}$, or $j_{i-1}=0$. So $n-|B(j)|\le1+Y+Y'$, where
$Y=|\{2\le i\le n:j_i\in\{0,i-1\}\}|$ and
$Y'=|\{1\le i\le n-1:j_i=0\}|$. These are sums of independent
Bernoulli$(2/i)$ and Bernoulli$(1/i)$ variables respectively (the terms $i=2$
in $Y$ and $i=1$ in $Y'$ are certain), with means at most $2\ln n$ and
$1+\ln n$. By Chernoff's bound, $Y\le4\ln n$ and $Y'\le2+2\ln n$ except with
probability $o(1)$, and $3+6\ln n\le7\ln n$ for $n\ge21$. The exact mean is
$\EE[n-|B(j)|]=3H_n-3+1/n$.

(c) Take $c=\lfloor n/2\rfloor$ and let $\Lambda$ be the set on the
right-hand side of the definition of $L(j)$. A position $p\le c$ is never
referenced from $(c,n]$ with probability $\prod_{i>c}(1-1/i)=c/n$, so
$\EE|\Lambda|=(c+1)(1-c/n)\ge n/4$. Moreover $|\Lambda|$ is a function of the
$N=\lceil n/2\rceil$ independent pointers $j_{c+1},\dots,j_n$, and changing
any one of them changes $|\Lambda|$ by at most $1$. McDiarmid's
inequality~\cite{kp-mcdiarmid} gives $|\Lambda|\ge n/4-n^{2/3}$ except with
probability $e^{-2n^{4/3}/N}\le e^{-2n^{1/3}}$.
\end{proof}

Consequently, methods whose cost is polylogarithmic only when the binary core
has polylogarithmic size, methods using space $O(L(j)\log n)$ (here
$\Theta(n\log n)$), and the Newton--Strahler bound
(Proposition~\ref{nw:prop-strahler-bound}, here $\Theta(\sqrt n)$ rounds) give
no polylogarithmic bound on a typical instance, while Newton's actual round
count is logarithmic by Theorem~\ref{nw:thm-A}.

\begin{theorem}[Golden-ratio density]\label{nw:thm-density}
Let $(j_i)_{i\ge1}$ be the infinite sequence of independent uniform pointers
of Subsection~\ref{nw:sec-setting}, and $x^*$ its trace. Almost surely, as
$m\to\infty$,
\[
  \frac1m\,\bigl|\{\,i<m:x^*_i=1\,\}\bigr|\longrightarrow\frac{3-\sqrt5}{2},
\]
the inverse square of the golden ratio.
\end{theorem}

\begin{proof}
Let $\cF_i=\sigma(j_1,\dots,j_i)$, let $O_m$ be the number of ones in $[0,m)$
and $\sigma_m=O_m/m$. For $i\ge1$, $x^*_i=1$ if and only if $x^*_{i-1}=0$ and
$x^*_{j_i}=0$; given $\cF_{i-1}$ this has probability
$\nwone\{x^*_{i-1}=0\}(1-\sigma_i)$. So, with $M_m$ the sum of the centred
indicators,
\[
  O_m=1+\sum_{1\le i<m}\nwone\{x^*_{i-1}=0\}(1-\sigma_i)+M_m.
\]
$M_m$ is a martingale with increments in $[-1,1]$. By Azuma's inequality and
the Borel--Cantelli lemma, $M_m/m\to0$ almost surely.

Let $L=\liminf\sigma_m$ and $U=\limsup\sigma_m$. Both lie in $[0,1/2]$, since
no two ones are adjacent. Put $g(u)=(1-u)/(2-u)$. Fix $\epsilon>0$ and choose
$i_0$ with $\sigma_i\le U+\epsilon$ for $i\ge i_0$. The number of zeros in
$[0,m-1)$ is at least $m-1-O_m$, so
\[
  O_m\ge(1-U-\epsilon)(m-1-O_m-i_0)+M_m,
\]
that is, $\sigma_m\ge(1-U-\epsilon)(1-\sigma_m)-o(1)$, or
$\sigma_m\ge g(U+\epsilon)-o(1)$. Symmetrically,
$\sigma_m\le g(L-\epsilon)+o(1)$. Letting $m\to\infty$ and then
$\epsilon\to0$ gives $L\ge g(U)$ and $U\le g(L)$.

The map $g$ is decreasing, with $|g'(u)|=(2-u)^{-2}\le4/9$ on $[0,1/2]$.
Hence $U-L\le g(L)-g(U)\le\frac49(U-L)$, so $U=L=:\sigma$ and $\sigma=g(\sigma)$.
That is, $\sigma^2-3\sigma+1=0$, so $\sigma=(3-\sqrt5)/2$.
\end{proof}

Numerically the density is within $0.0005$ of $(3-\sqrt5)/2=0.38197$ at
$n=10^3,10^4,10^5$.

\subsection{Consequences for complexity classes}\label{nw:sec-classes}

\begin{corollary}[Errorless average-case algorithms]\label{nw:cor-classes}
Let $\mathcal U_n$ denote the uniform distribution on parent arrays of
length $n$.
\begin{enumerate}[label=\textup{(\alph*)},leftmargin=*]
\item For every fixed $k\ge1$, the distributional problem
  $(\text{Sequential NOR},(\mathcal U_n)_{n\ge1})$ has an errorless
  average-case algorithm in logspace-uniform $\mathrm{NC}^{k+2}$ with failure
  probability $O((\log n)^{-k})$. On every instance the algorithm outputs
  either the correct trace or $\bot$.
\item There is a deterministic algorithm using $O(\log^3n)$ space that outputs
  the trace or $\bot$, is never wrong, and outputs $\bot$ with probability
  $O(1/\log n)$ under $\mathcal U_n$.
\end{enumerate}
\end{corollary}

\begin{proof}
(a) Run the proof of Theorem~\ref{nw:thm-A} with
$y_0=\max(e^{200},(\ln n)^k)$. The burn-in costs at most $y_0/2$ rounds. The
failure probability is at most
$10y_0^{-3}+\epsilon_6(y_0)+3/y_0+e^{-T'/32}=O((\ln n)^{-k})$. Since $x^*$ is
a fixed point of the Newton map, the circuit may run the logspace-computable
number
\[
  R=\lceil e^{200}\rceil+\lceil\log_2n\rceil^k
  +\lceil4.1e^{100}+1\rceil\lceil\log_2n\rceil
\]
of rounds, which is at least $\frac12y_0+\lceil4.1e^{100}\ln n\rceil$ for
this $y_0$. Each round has depth $O(\log^2n)$, and the final check is in
$\mathrm{NC}^1$.

(b) Take $k=1$ in (a). Logspace-uniform circuits of depth $d\ge\log n$ can be
evaluated in space $O(d)$~\cite{kp-borodin}.
\end{proof}

The simulation of~\cite{kp-borodin} takes time $n^{O(\log^2 n)}$, so (b)
does not give polynomial time and polylogarithmic space simultaneously. None of
this bears on worst-case inclusions such as $\mathrm P\subseteq\mathrm{NC}$:
an errorless average-case algorithm may output $\bot$ on a set of instances of
$\mathcal U_n$-probability $o(1)$, and nothing prevents the instances produced
by the P-completeness reduction from lying in that set. Newton itself needs
linearly many rounds on some instances (Proposition~\ref{nw:prop-worst}).

\begin{corollary}[Greedy independent sets]\label{nw:cor-lfmis}
Let $\Gamma_n$ be the random graph on the vertex set $\{0,\dots,n\}$ with
the edges $\{i-1,i\}$ and $\{j_i,i\}$, where $j_i$ is uniform on $\{0,\dots,i-1\}$. This is a
Hamiltonian path together with a uniform-attachment random recursive tree. The
lexicographically first maximal independent set of $\Gamma_n$ in the natural
order equals $\{i:x^*_i=1\}$. Hence it is computed by the errorless
average-case circuits of Corollary~\ref{nw:cor-classes}.
\end{corollary}

\begin{proof}
The neighbours of $i$ that come earlier in the order are exactly $i-1$ and
$j_i$. So the greedy rule includes $i$ if and only if neither of them is
included, that is, $x_i=\NOR(x_{i-1},x_{j_i})$, with vertex $0$ always
included.
\end{proof}

For the Erd\H{o}s--R\'enyi graph $G_{n,p}$ in the natural order, the
lexicographically first maximal independent set is computable in
polylogarithmic expected parallel time~\cite{nw-crt87,nw-cf90}. For an
arbitrary graph in a uniformly random order, the greedy dependence depth is
$O(\log n)$~\cite{nw-bfs12,nw-fn18}. The graph $\Gamma_n$ of
Corollary~\ref{nw:cor-lfmis} falls in neither regime. It is sparse, and its
edges are tied to the processing order, so unlike $G_{n,p}$ it is not
exchangeable, and processing it in a random order would be a different
problem.

\begin{remark}[Open problems]\label{nw:rem-open}
(1) The window length $\eps=e^{-100}$, and hence the constant in
Theorem~\ref{nw:thm-A}, comes from the crude Phase-1 growth factor
$\bar aM_1\approx4\cdot10^6$ and from Markov truncation; sharper concentration
should give a reasonable constant. (2) Only $\log_2\ln n-3$ is proved as a
lower bound (Theorem~\ref{nw:thm-B}); an $\Omega(\log n)$ bound needs a lower
bound on error growth past the front, which has to rule out cancellation.
(3) In experiments with $2^{16}\le n\le2^{25}$, Newton takes about
$0.87\log_2n$ rounds; we conjecture that $T_n/\log_2n\to c$ in probability
for some $c\approx0.87$.
\end{remark}

% ---------------------------------------------------------------------
% Section: a noise-tolerant embedding at noise 1/64
% Labels: all prefixed gb:.  External references: okhotin, ghr, hoeffding,
% gb-vonneumann, gb-pippenger, gb-spielmanteng, gb-blasermanthey, gb-ladner.
% ---------------------------------------------------------------------
\providecommand{\cG}{\mathcal{G}}
\providecommand{\Bin}{\operatorname{Bin}}

\section{A noise-tolerant embedding at noise \texorpdfstring{$1/64$}{1/64}}\label{sec:gapbase}

This section shows that Sequential NOR stays hard when an adversary chooses
the parent array and each pointer is then resampled with probability
$\eps\le 1/64$. We compile an arbitrary NOR circuit into a parent array by von
Neumann's construction: every wire becomes a bundle of copies, majority
gadgets restore the bundles, and the wiring between copies is random. After
smoothing, a strict majority of the output bundle still holds the value of
the circuit, except with probability $s^{-a}+(n+1)^{-(a+1)}$
(Theorem~\ref{gb:embed}). A simpler version of the construction works up to
noise $1/1000$ (Theorem~\ref{gb:thm1000}). Since smoothings compose
(Lemma~\ref{gb:comp}), an errorless randomized $\mathrm{NC}$ algorithm that
succeeds with probability $1/20$ on the $\eps$-smoothing of every base, for
some $\eps\le 1/64$, would give $\mathrm{P}\subseteq\mathrm{BPNC}$.

\subsection{Smoothing and smoothed algorithms}\label{gb:sec-smooth}

\begin{definition}[Smoothing]\label{gb:smooth}
Let $\eps\in[0,1]$ and let $j=(j_1,\dots,j_n)$ be a parent array,
$0\le j_i<i$. The \emph{$\eps$-smoothing} $S_\eps(j)$ is the random parent
array $\jt$ with $\jt_i=j_i$ if $\pi_i=0$ and $\jt_i=\rho_i$ if $\pi_i=1$.
Here $\pi_i$ is Bernoulli$(\eps)$, $\rho_i$ is uniform on $\{0,\dots,i-1\}$,
and all of these variables are independent. We call $\pi_i$ the
\emph{resampling indicator} of gate $i$, and $j$ the \emph{base}.
\end{definition}

Thus $S_0(j)=j$, and $S_1(j)$ is the uniform distribution $U_n$ on parent
arrays of length $n$, for every base $j$. Since $j_1=0$ is forced, every
instance has $x^*_0=1$ and $x^*_1=\NOR(1,1)=0$; we use this throughout.

\begin{definition}[Smoothed algorithms]\label{gb:smalg}
Let $\eps\in[0,1]$ and $p\in(0,1]$. An \emph{errorless smoothed algorithm at
noise $\eps$ with success $p$} is a family $(A_n)$ of circuits, where $A_n$
receives a parent array of length $n$ and a string $r$ of random bits of
length polynomial in $n$, such that
\begin{enumerate}[label=(\roman*)]
\item $A_n(j,r)\in\{x^*(j),\bot\}$ for every array $j$ and every $r$, and
\item $\PP_{\jt\sim S_\eps(j),\,r}\bigl[A_n(\jt,r)\ne\bot\bigr]\ge p$ for
every $n$ and every base $j$ of length $n$.
\end{enumerate}
It is \emph{in randomized $\mathrm{NC}$} if the family is logspace-uniform,
of polynomial size and of polylogarithmic depth. $\mathrm{BPNC}$ is the class
of languages decided by such randomized families with two-sided error at
most $1/3$.
\end{definition}

This is the smoothed analysis of Spielman and
Teng~\cite{gb-spielmanteng} applied to instances: the adversary chooses the
base, and nature perturbs it. It is not the smoothed complexity theory of
Bl\"aser and Manthey~\cite{gb-blasermanthey}, which has its own notion of
reduction; no notion of reduction to a smoothed problem is needed below.

\begin{lemma}[Composition of smoothings]\label{gb:comp}
Let $\eps,\eps'\in[0,1]$ and let $j$ be a parent array. If $S_\eps$ and
$S_{\eps'}$ are applied independently, then $S_\eps(S_{\eps'}(j))$ has the
law of $S_{\eps''}(j)$, where $1-\eps''=(1-\eps)(1-\eps')$. Consequently, an
errorless smoothed algorithm at noise $\eps$ with success $p$ is also an
errorless smoothed algorithm at every noise $\eps''\in[\eps,1]$, with the
same success $p$.
\end{lemma}

\begin{proof}
Let $(\pi'_i,\rho'_i)$ be the variables of the inner smoothing and
$(\pi_i,\rho_i)$ those of the outer one, so that $j'=S_{\eps'}(j)$ and
$\jt=S_\eps(j')$. Coordinate $i$ of $\jt$ is $\rho_i$ if $\pi_i=1$, it is
$\rho'_i$ if $\pi_i=0$ and $\pi'_i=1$, and it is $j_i$ if $\pi_i=\pi'_i=0$.
Put $\pi''_i=\max(\pi_i,\pi'_i)$, and $\rho''_i=\rho_i$ if $\pi_i=1$,
$\rho''_i=\rho'_i$ otherwise. Then $\jt_i=j_i$ if $\pi''_i=0$ and
$\jt_i=\rho''_i$ if $\pi''_i=1$. The indicator $\pi''_i$ is
Bernoulli$(\eps'')$ with $1-\eps''=(1-\eps)(1-\eps')$. Conditionally on
$(\pi_i,\pi'_i)$, the variable $\rho''_i$ is uniform on $\{0,\dots,i-1\}$, so
$\rho''_i$ is uniform and independent of $\pi''_i$. The pairs
$(\pi''_i,\rho''_i)$ for different $i$ are functions of disjoint families of
independent variables, hence independent. So $\jt$ has the law of
$S_{\eps''}(j)$.

For the second claim, let $(A_n)$ be an errorless smoothed algorithm at
noise $\eps$ with success $p$, and let $\eps''\in[\eps,1]$. Condition~(i) of
Definition~\ref{gb:smalg} does not involve the noise. If $\eps=1$ then
$\eps''=1$ and there is nothing to prove. Otherwise put
$\eps'=(\eps''-\eps)/(1-\eps)\in[0,1]$; then $(1-\eps)(1-\eps')=1-\eps''$.
By the first claim, for every base $j$,
\[
\PP_{\jt\sim S_{\eps''}(j),\,r}[A_n(\jt,r)\ne\bot]
=\EE_{j'\sim S_{\eps'}(j)}\,
\PP_{\jt\sim S_\eps(j'),\,r}[A_n(\jt,r)\ne\bot]\ge p,
\]
because the inner probability is at least $p$ for every base $j'$.
\end{proof}

\subsection{NOR circuits and gadgets}\label{gb:sec-gadgets}

\begin{definition}[NOR circuits]\label{gb:nor}
A \emph{NOR circuit} $C$ with $s$ gates is a sequence of gates
$g_1,\dots,g_s$ with $g_t=\NOR(u_t,v_t)$, where each of $u_t,v_t$ is an
earlier gate $g_{t'}$, $t'<t$, or one of the constants $c_0,c_1$. Values are
$\val(c_b)=b$ and $\val(g_t)=\NOR(\val(u_t),\val(v_t))$, and the output is
$\val(C)=\val(g_s)$.
\end{definition}

Evaluating NOR circuits is $\mathrm{P}$-complete under logspace reductions.
For NAND circuits this is \cite[Problem~A.1.5]{ghr} (see also
\cite{gb-ladner}). Complementing every value turns a NAND circuit into a NOR
circuit with complemented constants and complemented output, since
$\lnot\,\mathrm{NAND}(u,v)=u\wedge v=\NOR(\bar u,\bar v)$.

To \emph{append a gate with pointer $q$} to a parent array means to create
the next index $i$ and set $j_i=q$. A \emph{gadget} is a block of
consecutive gates, given by the list of their pointers. Let $k$ be the index
of the first gate of a gadget, and let $a,b,c<k$ be positions already
present. We use three gadgets:
\begin{align*}
N(a,b)&:\ (0,\ a,\ 1,\ b), \\
M(a,b,c)&:\ (0,\ a,\ 1,\ b,\ b,\ k+1,\ c,\ k+3), \\
M'(a,b,c)&:\ (a,\ a,\ c,\ k+1,\ b,\ k), \qquad\text{used only when } c=k-1.
\end{align*}
The gadget $N(a,b)$ has $4$ gates, $M(a,b,c)$ has $8$ and $M'(a,b,c)$ has
$6$. The \emph{output} of a gadget is its last gate.

Let $q_i$ be the pointer that a gadget lists for its gate $i$. We say that
gate $i$ \emph{reads its listed pointer} in a $0/1$ sequence $x$ if
$x_i=\NOR(x_{i-1},x_{q_i})$. This holds when $\jt_i=q_i$, but also when $\jt_i$ is any
position with $x_{\jt_i}=x_{q_i}$.

\begin{lemma}[Gadgets]\label{gb:gadgets}
Let $x$ be the trace of a parent array, so $x_0=1$ and $x_1=0$, and write
$a,b,c$ also for the values $x_a,x_b,x_c$.
\begin{enumerate}[label=(\roman*)]
\item If the gates $k,\dots,k+3$ of an $N(a,b)$ gadget all read their listed
pointers, then $x_{k+3}=\NOR(a,b)$, and $x_{k+2}=1-x_{k+1}$.
\item If the gates $k,\dots,k+7$ of an $M(a,b,c)$ gadget all read their
listed pointers, then $x_{k+7}=\MAJ(a,b,c)$.
\item If $c=k-1$ and the gates $k,\dots,k+5$ of an $M'(a,b,c)$ gadget all
read their listed pointers, then $x_{k+5}=\MAJ(a,b,c)$, and at least two of
$x_k,\dots,x_{k+5}$ equal $1$.
\end{enumerate}
Parts (i) and (ii) hold whatever the value of $x_{k-1}$.
\end{lemma}

\begin{proof}
(i) and (ii). Both gadgets begin with $x_k=\NOR(x_{k-1},x_0)=\NOR(\cdot,1)=0$,
then $x_{k+1}=\NOR(0,a)=\bar a$, $x_{k+2}=\NOR(\bar a,x_1)=\NOR(\bar a,0)=a$,
and $x_{k+3}=\NOR(a,b)$. This settles (i). For $M$, continue:
\begin{align*}
x_{k+4}&=\NOR(\NOR(a,b),\,b)=(a\vee b)\,\bar b=a\bar b,\\
x_{k+5}&=\NOR(a\bar b,\,x_{k+1})=\NOR(a\bar b,\,\bar a)=(\bar a\vee b)\,a=ab,\\
x_{k+6}&=\NOR(ab,\,c),\\
x_{k+7}&=\NOR(x_{k+6},\,x_{k+3})=(ab\vee c)(a\vee b)=ab\vee ac\vee bc
=\MAJ(a,b,c).
\end{align*}
(iii) Here $x_{k-1}=x_c=c$. Then
\begin{align*}
x_k&=\NOR(c,a), & x_{k+1}&=\NOR(\NOR(c,a),a)=c\bar a,\\
x_{k+2}&=\NOR(c\bar a,c)=\bar c, & x_{k+3}&=\NOR(\bar c,c\bar a)=ca,\\
x_{k+4}&=\NOR(ca,b)=\overline{ca}\,\bar b, &
x_{k+5}&=\NOR(x_{k+4},x_k)=(ca\vee b)(c\vee a),
\end{align*}
and $(ca\vee b)(c\vee a)=ca\vee ab\vee bc=\MAJ(a,b,c)$. For the count of ones:
if $c=0$, the six values are $\bar a,0,1,0,\bar b,ab$, and
$\bar a+\bar b+ab\ge1$; if $c=1$, they are $0,\bar a,0,a,\bar a\bar b,a\vee
b$, and $\bar a+a=1$ while $\bar a\bar b+(a\vee b)=1$. In both cases at least
two values are $1$.
\end{proof}

Call a gadget \emph{oblivious} if it works for every value of $x_{k-1}$ and
its pointers go only to $0$, $1$, its inputs and its own earlier gates. The
gadgets $N$ and $M$ are oblivious; $M'$ is not, and is used only where the
layout puts its third input at position $k-1$.

\subsection{The embedding at noise \texorpdfstring{$1/1000$}{1/1000}}
\label{gb:sec-1000}

A \emph{cell of depth $h$} is a complete ternary tree of depth $h$ whose
leaves are $N$ gadgets and whose internal nodes are majority gadgets, laid
out in post-order: the gadgets of the three subtrees of a node come first,
in order, then the gadget of the node itself, and each internal node takes
as $a,b,c$ the outputs of its first, second and third child. In this layout
the third input $c$ of every internal node is the output of its third
child, which is the position $k-1$ just before the node's first gate. The
output of a cell is the output of its root.

\begin{definition}[The compiler $\Comp$]\label{gb:compiler}
Let $C$ be a NOR circuit with $s\ge2$ gates, and let $m$ be the least power
of two with $m\ge 20\log_2 s$. The cells of $\Comp$ have depth $3$, with
$27$ leaves $N$ and $13$ internal nodes $M$, hence $27\cdot4+13\cdot8=212$
gates. The compiler appends gate $1$ (pointer $0$) and then, for
$t=1,\dots,s$ and $\ell=1,\dots,m$ in this order, a cell $C(t,\ell)$. Write
$o(t,\ell)$ for the output position of $C(t,\ell)$, and put $o(c_1,\ell)=0$
and $o(c_0,\ell)=1$ for all $\ell$. Each leaf of $C(t,\ell)$ is
$N\bigl(o(u_t,r),o(v_t,r')\bigr)$, where $r,r'$ are fresh independent
uniform elements of $\{1,\dots,m\}$; these are the \emph{coins} of the
compiler. The base $j=\Comp(C)$ has length $n=1+212sm$, and its last
position $n$ is $o(s,m)$.
\end{definition}

Since $m$ is a power of two, each coin is drawn exactly from $\log_2 m$
random bits. All positions are explicit arithmetic expressions in $t$,
$\ell$ and the position inside a cell, so $\Comp$ is computable in logspace
from $C$ and the coins. In the unsmoothed instance every gate reads its
listed pointer, and since $x_0=1=\val(c_1)$ and $x_1=0=\val(c_0)$,
Lemma~\ref{gb:gadgets} and induction along the layout give
$x^*_{o(t,\ell)}(\Comp(C))=\val(g_t)$ for all $t,\ell$, whatever the coins.

\begin{theorem}[Noise-tolerant embedding up to $1/1000$]\label{gb:thm1000}
Let $s\ge2$ and $0\le\eps\le1/1000$. For every NOR circuit $C$ with $s$
gates,
\[
\PP\bigl[x^*_n\bigl(S_\eps(\Comp(C))\bigr)\ne\val(C)\bigr]\le 0.0140+1/s,
\]
where the probability is over the coins of the compiler and the smoothing.
\end{theorem}

\begin{proof}
Let $j=\Comp(C)$, let $\jt=S_\eps(j)$ with variables $(\pi_i,\rho_i)$, and
let $x=x^*(\jt)$. Call the cell $C(t,\ell)$ \emph{bad} if
$x_{o(t,\ell)}\ne\val(g_t)$; the constants $c_0,c_1$ (positions $1$ and $0$)
are never bad. Let $\beta_t$ be the fraction of bad cells among
$C(t,1),\dots,C(t,m)$, and put $\beta=0$ for the constants. A gadget is
\emph{hit} if $\pi_i=1$ for one of its gates.

The positions $0,1$ together with the cells $C(t',\ell)$, $t'\le t$, form an
initial segment $I_t$ of positions. Let $\cG_t$ be the $\sigma$-field
generated by the coins of these cells and by $(\pi_i,\rho_i)_{i\in I_t}$.
Every pointer points backwards and the listed pointers in $I_t$ depend only
on $C$ and the coins of cells in $I_t$, so $x$ restricted to $I_t$ is
$\cG_t$-measurable. In particular $\beta_1,\dots,\beta_t$ and the sets of
bad cells of bundles $1,\dots,t$ are $\cG_t$-measurable.

\emph{Step 1: dominating events.} Fix $t$ and $\ell$. For a node $\nu$ of
$C(t,\ell)$ define an event $D(\nu)$ by induction on height. If $\nu$ is
the leaf $N(o(u_t,r),o(v_t,r'))$, then
$D(\nu)=\{\nu\text{ hit}\}\cup\{C(u_t,r)\text{ bad}\}\cup\{C(v_t,r')\text{
bad}\}$, where a constant is never bad. If $\nu$ is internal, then
$D(\nu)=\{\nu\text{ hit}\}\cup\{D(\mu)\text{ holds for at least two children
}\mu\text{ of }\nu\}$. We claim: \emph{if the output of $\nu$ differs from
$\val(g_t)$, then $D(\nu)$ holds.} A leaf that is not hit carries its listed
pointers, so by Lemma~\ref{gb:gadgets}(i) its output is
$\NOR(x_{o(u_t,r)},x_{o(v_t,r')})$, which equals $\val(g_t)$ when neither
input cell is bad. An internal node that is not hit outputs the majority of
the outputs of its children, by Lemma~\ref{gb:gadgets}(ii). If $D(\mu)$
fails for at least two children, then by induction at least two children
output $\val(g_t)$, and so does the majority.

\emph{Step 2: conditional law.} The event $D(\nu)$ is determined by three
things: the indicators $\pi_i$ of the gates in the subtree of $\nu$; the
coins of the leaves of that subtree; and the sets of bad cells of the
bundles $u_t,v_t$, which are $\cG_{t-1}$-measurable because $u_t,v_t$ are
constants or gates before $g_t$. The first two are independent of
$\cG_{t-1}$, and they are independent across disjoint subtrees. Hence,
conditionally on $\cG_{t-1}$, the events $D(\mathrm{root}\,C(t,\ell))$ for
$\ell=1,\dots,m$ are independent, and so are the events $D(\mu)$ for the
three children $\mu$ of an internal node, together with the event that the
node itself is hit. For a leaf, the three events defining $D$ are
conditionally independent, with probabilities $1-(1-\eps)^4$,
$\beta_{u_t}$ and $\beta_{v_t}$, because $r$ and $r'$ are uniform and
independent of $\cG_{t-1}$. So every node of height $l$ in bundle $t$ has
the same conditional probability $q_l$ of $D$, depending on $t$ only
through $\beta_{u_t},\beta_{v_t}$:
\begin{equation}\label{gb:eq-q1000}
q_0=1-(1-\eps)^4(1-\beta_{u_t})(1-\beta_{v_t}),\qquad
q_{l+1}=1-(1-\eps)^8\bigl(1-\phi(q_l)\bigr),
\end{equation}
where $\phi(q)=3q^2-2q^3$ is the probability that at least two of three
independent events of probability $q$ occur.

\emph{Step 3: numbers.} The function $\phi$ is increasing on $[0,1]$, so
$q_3$ is increasing in $\beta_{u_t}$, $\beta_{v_t}$ and $\eps$. Let
$A_t=\{\beta_{u_t}\le1/10,\ \beta_{v_t}\le1/10\}$, an event in
$\cG_{t-1}$. On $A_t$, using $1-(1-\eps)^k\le k\eps$ and
$\eps\le1/1000$, we get $q_0\le 4\eps+2/10\le0.204$ and
$q_{l+1}\le8\eps+\phi(q_l)\le 0.008+\phi(q_l)$. Hence $q_1\le0.11587$,
$q_2\le0.04517$ and $q_3\le0.01394=:\bar q$.

\emph{Step 4: concentration.} On $A_t$, by Step~1 every bad cell of bundle
$t$ has $D(\mathrm{root})$, and conditionally on $\cG_{t-1}$ the number of
roots with $D$ is $\Bin(m,q_3)$ with $q_3\le\bar q$. Let
$k=\lfloor m/10\rfloor+1$, so that $\beta_t>1/10$ means at least $k$ bad
cells. Then on $A_t$
\[
\PP[\beta_t>1/10\mid\cG_{t-1}]\le\binom mk\bar q^{\,k}
\le\Bigl(\frac{em\bar q}{k}\Bigr)^k\le(10e\bar q)^k\le2^{-k}
\le2^{-m/10}\le s^{-2},
\]
since $10e\bar q<0.38$ and $m\ge20\log_2 s$.

\emph{Step 5: all gates.} Let $H_t=\{\beta_1,\dots,\beta_t\le1/10\}$. The
inputs $u_t,v_t$ are constants or gates before $g_t$, so $H_{t-1}\subseteq
A_t$, and $H_{t-1}\in\cG_{t-1}$. By Step~4,
$\PP[H_{t-1}\setminus H_t]\le s^{-2}$, hence $\PP[H_{s-1}^c]\le
(s-1)/s^2<1/s$. Finally $x_n$ is the output of $C(s,m)$, so by Steps~1--3
\[
\PP[x_n\ne\val(g_s)]\le\PP[H_{s-1}^c]
+\EE\bigl[\mathbf 1_{H_{s-1}}\,\PP\bigl(D(\mathrm{root}\,C(s,m))\bigm|
\cG_{s-1}\bigr)\bigr]\le\frac1s+\bar q.\qedhere
\]
\end{proof}

The compiler does not depend on $\eps$. Theorem~\ref{gb:thm1000} stops at
$1/1000$ because every hit gadget counts as failed, the answer is read from
a single position, and the tail bound of Step~4 with $m\ge20\log_2 s$ needs
$\bar q\le1/(20e)$. The next construction removes all three losses.

\subsection{The embedding at noise \texorpdfstring{$1/64$}{1/64}}
\label{gb:sec-64}

Four changes raise the admissible noise to $1/64$. Internal nodes use the
$6$-gate majority gadget $M'$; cells are deeper; the bundles are longer and
are preceded by padding, which makes the density of ones bounded away from
$0$ and $1$ (Lemma~\ref{gb:density}); and the answer is read from the whole
output bundle. The density is what makes a resampled pointer harmless most
of the time: it does no harm when it happens to read the value its listed
pointer would have read.

\begin{definition}[The compiler $\Comp'_a$]\label{gb:compiler64}
Let $a>0$ be rational, and let $C$ be a NOR circuit with $s\ge2$ gates.
\begin{enumerate}[label=(\roman*)]
\item A cell is a complete ternary tree of depth $10$, laid out in
post-order as above. Its $3^{10}=59049$ leaves are $N$ gadgets and its
$(3^{10}-1)/2=29524$ internal nodes are $M'$ gadgets, so a cell has
$59049\cdot4+29524\cdot6=413340$ gates. In the post-order layout the third
input of every internal node is the position $k-1$, as $M'$ requires.
\item $m$ is the least power of two with $m\ge65(a+1)\log_2 s$.
\item Let $L=\lceil\log_2(10^7(a+2)sm)\rceil$ and
$P=\lceil23(a+2)L\rceil+125$. After gate $1$ (pointer $0$) the compiler
appends $P$ padding gadgets $N(0,0)$, occupying positions $2,\dots,4P+1$.
\item Then, for $t=1,\dots,s$ and $\ell=1,\dots,m$ in this order, it appends
the cell $C(t,\ell)$. With $o(t,\ell)$, $o(c_1,\ell)=0$, $o(c_0,\ell)=1$ and
the coins as in Definition~\ref{gb:compiler}, each leaf of $C(t,\ell)$ is
$N\bigl(o(u_t,r),o(v_t,r')\bigr)$ with fresh coins $r,r'$.
\end{enumerate}
The exponent $a$ is unrelated to the gadget inputs $a,b,c$ of
Lemma~\ref{gb:gadgets}. The base $j=\Comp'_a(C)$ has length $n=4P+1+413340\,sm$. Its \emph{output
bundle} is the set of positions $o(s,1),\dots,o(s,m)$, and the first gate of
the first cell is $g_0=4P+2$.
\end{definition}

As before, $\Comp'_a$ is computable in logspace from $C$ and the coins, the
coins are exactly uniform, and the unsmoothed instance satisfies
$x^*_{o(t,\ell)}=\val(g_t)$ for all $t,\ell$ by Lemma~\ref{gb:gadgets}(i)
and (iii). We record two inequalities on the length. First,
$a+2\le m\le sm$ since $m\ge65(a+1)$, so
$L\le\log_2(10^7)+1+2\log_2(sm)\le25+2sm$, and
\begin{equation}\label{gb:eq-nbound}
n+1=4P+2+413340\,sm\le92(a+2)L+506+413340\,sm\le10^7(a+2)sm\le2^L.
\end{equation}
Second, since $\ln(n+1)\le L\ln2\le0.6932\,L$ by~\eqref{gb:eq-nbound},
\begin{equation}\label{gb:eq-g0}
g_0=4P+2\ge92(a+2)L+502\ge130(a+2)\ln(n+1)+500=:g^*.
\end{equation}
So every gate of every cell lies at a position $g$ with $g^*\le g\le n$.

\begin{lemma}[Density of compiled instances]\label{gb:density}
Let $a>0$ and $0\le\eps\le1/64$. Let $j$ be a parent array of length $n$
consisting of gate $1$ (pointer $0$) followed by consecutive gadgets, each
of the form $N(a',b')$ or $M'(a',b',c')$ with $c'=k-1$. Let $x$ be the trace
of $S_\eps(j)$, let $O_{g-1}$ be the number of ones among
$x_0,\dots,x_{g-1}$, and let $\psi_{g-1}=O_{g-1}/g$. With probability at
least $1-(n+1)^{-(a+1)}$,
\[
\psi_{g-1}\ge0.19\qquad\text{for all $g$ with }
130(a+2)\ln(n+1)+500\le g\le n.
\]
\end{lemma}

\begin{proof}
A gadget none of whose gates is resampled has all its gates reading their
listed pointers, so at least a quarter of its positions are ones, whatever
its input values: for $N$, Lemma~\ref{gb:gadgets}(i) gives
$x_{k+2}=1-x_{k+1}$, one $1$ among $4$ positions; for $M'$,
Lemma~\ref{gb:gadgets}(iii) gives two $1$s among $6$.

Fix $g$, and let $R$ be the number of resampled gates among positions
$2,\dots,g-1$; it has law $\Bin(g-2,\eps)$. A position in $\{2,\dots,g-1\}$
that does not lie in a complete gadget inside $\{2,\dots,g-1\}$ free of
resampled gates either lies in a gadget containing a resampled gate, which
accounts for at most $6R$ positions, or in the last, incomplete gadget, at
most $5$ positions. Hence $O_{g-1}\ge(g-2-6R-5)/4$. If $R\le(g-2)/30$, this
gives $O_{g-1}\ge0.2g-1.65\ge0.19g$, as $g\ge165$.

Let $D(\alpha\|\beta)=\alpha\ln(\alpha/\beta)+(1-\alpha)\ln((1-\alpha)/(1-\beta))$
be the relative entropy. It decreases in $\beta$ on $(0,\alpha)$, and
$D(1/30\|1/64)\ge0.0077$. By the Chernoff--Hoeffding
bound~\cite{hoeffding}, for $0<\eps\le1/64$,
\[
\PP[R>(g-2)/30]\le\exp\bigl(-(g-2)D(1/30\|\eps)\bigr)\le e^{-0.0077(g-2)}
\le(n+1)^{-(a+2)},
\]
since $g-2\ge130(a+2)\ln(n+1)$ and $0.0077\cdot130>1$; for $\eps=0$ the
probability is $0$. A union bound over the at most $n$ values of $g$ gives
failure probability at most $n(n+1)^{-(a+2)}<(n+1)^{-(a+1)}$.
\end{proof}

\begin{theorem}[Embedding at noise $1/64$]\label{gb:embed}
Let $a>0$ be rational, $s\ge2$ and $0\le\eps\le1/64$. For every NOR circuit
$C$ with $s$ gates, let $j=\Comp'_a(C)$, of length $n$. With probability at
least $1-s^{-a}-(n+1)^{-(a+1)}$ over the coins and the smoothing, at least
five sixths of the positions $o(s,1),\dots,o(s,m)$ of $x^*(S_\eps(j))$ hold
$\val(C)$. In particular, the probability that the majority of the output
bundle is wrong, that is, that at most $m/2$ of these positions hold
$\val(C)$, is at most $s^{-a}+(n+1)^{-(a+1)}$.
\end{theorem}

\begin{proof}
Let $\jt=S_\eps(j)$ with variables $(\pi_i,\rho_i)$, let $x=x^*(\jt)$, and
let $G_d$ be the event of Lemma~\ref{gb:density}, which applies to $j$ and
has $\PP[G_d^c]\le(n+1)^{-(a+1)}$. Let $(U_g)$ be independent uniform
variables on $[0,1]$, independent of everything else; they serve only the
analysis. Let $\cF_g$ be the $\sigma$-field generated by all coins, by
$(\pi_i,\rho_i)_{i\le g}$ and by $(U_i)_{i\le g}$. Then
$x_0,\dots,x_g$ are $\cF_g$-measurable.

\emph{Flags.} For a gate $g\ge g_0$ of a cell, with listed pointer $j_g$,
let $f_g$ be the fraction of positions $p<g$ with $x_p\ne x_{j_g}$; it is
$\cF_{g-1}$-measurable, and $f_g<1$ since $p=j_g$ itself agrees. Put
$h_g=0.501$ if $j_g=1$ and $h_g=0.81$ otherwise. The listed pointer $1$
occurs only as the third pointer of a leaf and as $o(c_0,r)=1$, and the
other listed pointers of cells are $0$ or positions $\ge g_0>1$; so whether
$j_g=1$ depends only on $C$, and $h_g$ is a constant, not depending on the
coins. Define the \emph{flag} $B_g\in\{0,1\}$ as follows. If $f_g\le h_g$,
let $B_g=1$ when $\pi_g=1$ and either $x_{\rho_g}\ne x_{j_g}$ or
$U_g<(h_g-f_g)/(1-f_g)$. If $f_g>h_g$, let $B_g=1$ when $\pi_g=1$ and
$U_g<h_g$. Given $\cF_{g-1}$, the variables $\pi_g$, $\rho_g$, $U_g$ are
independent, $\pi_g$ is Bernoulli$(\eps)$ and $\rho_g$ is uniform on
$\{0,\dots,g-1\}$, so $x_{\rho_g}\ne x_{j_g}$ with conditional probability
$f_g$. In both cases
\[
\PP[B_g=1\mid\cF_{g-1}]=\eps h_g,
\]
a constant. Since $B_g$ is $\cF_g$-measurable, induction on $g$ shows that
the flags $(B_g)_{g\ge h}$ are independent Bernoulli$(\eps h_g)$ variables,
jointly independent of $\cF_{h-1}$, for every $h\ge g_0$. Note that
$\cF_{h-1}$ contains all coins.

On $G_d$ the flags dominate harmful resamplings. Ones are never adjacent in
a trace, so $\psi_{g-1}\le(g+1)/(2g)\le0.501$ for $g\ge500$. If $j_g=1$,
then $x_{j_g}=0$ and $f_g=\psi_{g-1}\le0.501=h_g$. Otherwise
$f_g\in\{\psi_{g-1},1-\psi_{g-1}\}$, and on $G_d$, by~\eqref{gb:eq-g0},
$f_g\le\max(\psi_{g-1},1-\psi_{g-1})\le0.81=h_g$. So on $G_d$ every cell gate
with $\pi_g=1$ and $x_{\rho_g}\ne x_{j_g}$ has $B_g=1$. A gate with
$\pi_g=0$, or with $x_{\rho_g}=x_{j_g}$, reads its listed pointer. Call a
gadget \emph{flagged} if $B_g=1$ for one of its gates. On $G_d$, every gate
of an unflagged gadget reads its listed pointer, so by
Lemma~\ref{gb:gadgets} an unflagged leaf outputs $\NOR$ of the values at
its inputs and an unflagged internal node outputs their majority.

\emph{Filtration.} Let $I_0=\{0,\dots,g_0-1\}$ (gate $1$ and the padding)
and $I_t=I_0\cup\bigcup_{t'\le t,\,\ell}C(t',\ell)$, an initial segment of
positions. Let $\cG_t$ be generated by the coins of the cells in $I_t$ and
by $(\pi_i,\rho_i,U_i)_{i\in I_t}$. As before, $x$ restricted to $I_t$, the
fractions $\beta_1,\dots,\beta_t$ of bad cells and the sets of bad cells of
bundles $1,\dots,t$ are $\cG_t$-measurable, with bad cells and $\beta$
defined as in the proof of Theorem~\ref{gb:thm1000}. Let $h_t$ be the first
position of bundle $t$. Then $\cG_{t-1}$ and the coins of bundle $t$ lie in
$\cF_{h_t-1}$, so the flags of bundle $t$ are independent of them jointly.
The coins of bundle $t$ are independent of $\cG_{t-1}$, being primitive
variables not among its generators. Hence the flags and the coins of bundle
$t$ are jointly independent of $\cG_{t-1}$, and all of them are mutually
independent.

\emph{Steps 1 and 2.} Define $D(\nu)$ as in the proof of
Theorem~\ref{gb:thm1000}, with ``flagged'' in place of ``hit''; the
definition does not refer to $G_d$. The argument of Step~1 there, using the
previous paragraph, shows: \emph{on $G_d$, if the output of $\nu$ differs
from $\val(g_t)$, then $D(\nu)$ holds.} The event $D(\nu)$ is determined by
the flags of the gates of the subtree of $\nu$, the coins of its leaves,
and the bad sets of the bundles $u_t,v_t$, which are
$\cG_{t-1}$-measurable. By the independence just shown, conditionally on
$\cG_{t-1}$ the events $D(\mathrm{root}\,C(t,\ell))$, $\ell=1,\dots,m$, are
independent, and for each internal node the events $D(\mu)$ of its three
children and the event that the node is flagged are independent. A leaf has
three gates with $h_g\le0.81$ and one, its third, with $h_g=0.501$, so it is
flagged with probability at most $\eps(3\cdot0.81+0.501)=2.931\,\eps$; an
internal node has six gates with $h_g=0.81$ and is flagged with probability
at most $6\cdot0.81\,\eps=4.86\,\eps$. The probability that at least two of
three independent events occur is increasing in each of their
probabilities. Hence, on
$A_t=\{\beta_{u_t}\le1/6,\ \beta_{v_t}\le1/6\}\in\cG_{t-1}$, every node of
height $l$ in bundle $t$ has $\PP[D(\nu)\mid\cG_{t-1}]\le q_l$, where
\begin{equation}\label{gb:eq-q64}
q_0=1-(1-2.931\,\eps)\Bigl(\frac56\Bigr)^2,\qquad
q_{l+1}=1-(1-4.86\,\eps)\bigl(1-\phi(q_l)\bigr).
\end{equation}

\emph{Step 3: numbers.} The right-hand sides of~\eqref{gb:eq-q64} increase
in $\eps$ and $q_l$, so it suffices to take $\eps=1/64$. Exact rational
arithmetic, rounding up to five decimals after every step, gives the
following upper bounds:
\begin{center}
\begin{tabular}{@{}lcccccc@{}}
\toprule
$l$ & 0 & 1 & 2 & 3 & 4 & 5\\
$q_l\le$ & 0.33736 & 0.32049 & 0.29985 & 0.27537 & 0.24756 & 0.21780\\
\midrule
$l$ & 6 & 7 & 8 & 9 & 10 & \\
$q_l\le$ & 0.18835 & 0.16194 & 0.14079 & 0.12573 & 0.11609 & \\
\bottomrule
\end{tabular}
\end{center}
So the root of every cell of bundle $t$ has
$\PP[D(\mathrm{root})\mid\cG_{t-1}]\le\bar q:=0.1161$ on $A_t$.

\emph{Step 4: concentration.} On $A_t\cap G_d$, every bad cell of bundle $t$
has $D(\mathrm{root})$. Conditionally on $\cG_{t-1}$, on $A_t$, the number
of roots with $D$ is a sum of $m$ independent indicators of probability at
most $\bar q<1/6$, so it is stochastically dominated by $\Bin(m,\bar q)$.
By the Chernoff--Hoeffding bound~\cite{hoeffding}, on $A_t$,
\begin{align*}
\PP\bigl[\{\beta_t>1/6\}\cap G_d\bigm|\cG_{t-1}\bigr]
&\le\exp\bigl(-m\,D(1/6\|\bar q)\bigr)\\
&\le\exp\bigl(-65(a+1)\cdot0.0111\cdot\log_2 s\bigr)\le s^{-(a+1)},
\end{align*}
since $D(1/6\|0.1161)\ge0.0111$ and $65\cdot0.0111/\ln2>1$.

\emph{Step 5: all gates.} Let $H_t=\{\beta_1,\dots,\beta_t\le1/6\}$, so
$H_0$ is the sure event. As before $H_{t-1}\subseteq A_t$ and
$H_{t-1}\in\cG_{t-1}$, so Step~4 gives
$\PP[H_{t-1}\cap\{\beta_t>1/6\}\cap G_d]\le s^{-(a+1)}$. Since
\[
H_s^c\subseteq G_d^c\cup\bigcup_{t\le s}\bigl(H_{t-1}\cap\{\beta_t>1/6\}\cap
G_d\bigr),
\]
we get $\PP[H_s^c]\le(n+1)^{-(a+1)}+s\cdot s^{-(a+1)}
=(n+1)^{-(a+1)}+s^{-a}$. On $H_s$ at most $m/6$ cells of bundle $s$ are
bad, so at least $5m/6>m/2$ positions of the output bundle hold
$\val(g_s)=\val(C)$.
\end{proof}

For $a=1$ the bound is $1/s+1/(n+1)^2$, with $m\ge130\log_2 s$. The
probability bound of Theorem~\ref{gb:embed} does not depend on the noise
level, and the construction does not depend on $\eps$.

\subsection{Consequences}\label{gb:sec-cons}

The following consequences are stated in their simplest form; stronger
versions appear elsewhere in this paper.

\begin{corollary}[Smoothed P-hardness at noise $1/1000$]\label{gb:cor1000}
Suppose that for some $\eps\le1/1000$ there is an errorless smoothed
algorithm in randomized $\mathrm{NC}$ at noise $\eps$ with success at least
$1/20$. Then $\mathrm{P}\subseteq\mathrm{BPNC}$.
\end{corollary}

\begin{proof}
By Lemma~\ref{gb:comp}, the same family $(A_n)$ is an errorless smoothed
algorithm at noise $\eps_*=1/1000$ with success at least $1/20$. Let $C$ be
a NOR circuit with $s$ gates; prefixing dummy gates $\NOR(c_0,c_0)$, we may
assume $s\ge1000$. A \emph{trial} computes $j=\Comp(C)$ with fresh coins,
draws $\jt$ from a distribution within total variation $n^{-2}$ of
$S_{\eps_*}(j)$, runs $A_n(\jt,r)$ and reads bit $n$ of its output. The
sampling is in $\mathrm{NC}^1$: reduce random integers of
$4\lceil\log_2 n\rceil+10$ bits modulo $1000$ and modulo $i$; since
$n\ge1000$, each of the $2n$ draws is within total variation
$n/(1024\,n^4)$ of its target.
Let $Y=+1$ if the bit read is $\val(C)$, $Y=-1$ if it is $1-\val(C)$, and
$Y=0$ on $\bot$. Since $A_n$ is errorless, $Y=-1$ forces
$x^*_n(\jt)\ne\val(C)$, so $\PP[Y=-1]\le\delta+n^{-2}$ with
$\delta=0.0140+1/s\le0.015$ by Theorem~\ref{gb:thm1000}. Also
$\PP[Y=+1]=\PP[Y\ne0]-\PP[Y=-1]\ge1/20-\delta-2n^{-2}$. Hence
$\EE Y\ge1/20-2\delta-3n^{-2}\ge0.019$.

Run $T=10^4s$ independent trials in parallel and output the bit with more
votes among the trials that did not return $\bot$. This is wrong only if
$\sum Y\le0$, which by Hoeffding's inequality~\cite{hoeffding} has
probability at most $\exp(-T(\EE Y)^2/2)\le e^{-1.8s}$. Every step is in
logspace-uniform randomized $\mathrm{NC}$. Evaluating NOR circuits is
$\mathrm{P}$-complete under logspace reductions, and $\mathrm{BPNC}$ is
closed under them. Hence $\mathrm{P}\subseteq\mathrm{BPNC}$.
\end{proof}

\begin{corollary}[Smoothed P-hardness at noise $1/64$]\label{gb:cor64}
Suppose that for some $\eps\le1/64$ there is an errorless smoothed algorithm
in randomized $\mathrm{NC}$ at noise $\eps$ with success at least $1/20$.
Then $\mathrm{P}\subseteq\mathrm{BPNC}$.
\end{corollary}

\begin{proof}
Repeat the proof of Corollary~\ref{gb:cor1000} with $\eps_*=1/64$ and the
compiler $\Comp'_1$, changing only how a trial reads its answer. Sampling
$S_{1/64}$ is now exact up to the residues modulo $i$. When $A_n$ does not
return $\bot$, its output is the trace, and the trial takes the majority of
its bits at the output bundle $o(s,1),\dots,o(s,m)$. By
Theorem~\ref{gb:embed} with $a=1$ this majority is wrong with probability
at most $\delta=1/s+(n+1)^{-2}\le0.0011$ for $s\ge1000$, up to the sampling
error $n^{-2}$. Then $\EE Y\ge1/20-2\delta-3n^{-2}\ge0.047$, and with
$T=10^4s$ trials Hoeffding's inequality bounds the error by
$\exp(-T(\EE Y)^2/2)\le e^{-11s}$.
\end{proof}

\begin{remark}[Algorithms that may err]\label{gb:rem-err}
Suppose a randomized $\mathrm{NC}$ algorithm, at some noise $\eps\le1/1000$,
outputs $x^*_n(\jt)$ correctly with probability at least $1/2+1/20$ for
every base. The averaging argument of Lemma~\ref{gb:comp} shows that it
does so at noise $1/1000$ as well. In a trial it then outputs $\val(C)$ with
probability at least $0.55-0.015-n^{-2}\ge0.53$, and majority amplification
again gives $\mathrm{P}\subseteq\mathrm{BPNC}$. At noise $1/64$ the same
holds, by the majority readout, for algorithms that output the whole trace
correctly with probability at least $1/2+1/20$; for algorithms that output a
single bit we do not claim it.
\end{remark}

Recall that Newton's method over $\mathrm{GF}(2)$ for an instance $j$ solves
$F(y)=y+G(y)=0$, where $G_i(y)=(1+y_{i-1})(1+y_{j_i})$ if $j_i\ne i-1$ and
$G_i(y)=1+y_{i-1}$ if $j_i=i-1$ (a NOT gate), by the step
$(I+J_G(x))(x^++x)=x+G(x)$, with $x^+_0=1$, from the standard start
$x^{(0)}=(1,0,\dots,0)$. Here $J_G(x)$ is the matrix of formal partial
derivatives $\partial_kG_i(x)$; it is strictly lower triangular,
so a round is computable in $\mathrm{NC}^2$ via
$(I+J)^{-1}=\prod_{k<\lceil\log_2 n\rceil+1}(I+J^{2^k})$. The trace is the
unique solution of $F(x)=0$, by forward induction on $i$, and a fixed point
of the Newton map, since the step at $x^*$ is $0$.

\begin{corollary}[Newton on smoothed instances]\label{gb:newton}
Suppose $\mathrm{P}\not\subseteq\mathrm{BPNC}$. Let $T(j)$ be the number of
Newton rounds from the standard start until the iterate equals $x^*(j)$.
Then for every $k\ge1$ and every $\eps\le1/64$ there are infinitely many $n$
and bases $j$ of length $n$ with
$\PP\bigl[T(S_\eps(j))>\lceil\log_2 n\rceil^k\bigr]>19/20$.
\end{corollary}

\begin{proof}
Run $R=\lceil\log_2 n\rceil^k$ Newton rounds, and output the iterate if it
satisfies $F(x)=0$ and $\bot$ otherwise. By the facts just recalled this is
an errorless, logspace-uniform $\mathrm{NC}^{k+2}$ algorithm whose success on
$\jt$ is $\PP[T(\jt)\le R]$. If the conclusion failed, this success would be
at least $1/20$ for every base of every large $n$; finitely many small $n$
are handled by table lookup. This contradicts
Corollary~\ref{gb:cor64}.
\end{proof}

\begin{theorem}[A gap in the noise threshold]\label{gb:gap}
Let $E\subseteq[0,1]$ be the set of noise levels $\eps$ at which an
errorless smoothed algorithm in randomized $\mathrm{NC}$ with success $1/2$
exists, and let $\eps_c=\inf(E\cup\{1\})$. Then:
\begin{enumerate}[label=(\alph*)]
\item $E$ is an up-set: if $\eps\in E$ and $\eps''\in[\eps,1]$, then
$\eps''\in E$;
\item if $\mathrm{P}\subseteq\mathrm{BPNC}$ then $E=[0,1]$, and otherwise
$E\subseteq(1/64,1]$.
\end{enumerate}
In particular $\eps_c\in\{0\}\cup[1/64,1]$, and $\eps_c=0$ if and only if
$\mathrm{P}\subseteq\mathrm{BPNC}$.
\end{theorem}

\begin{proof}
Part (a) is Lemma~\ref{gb:comp}. For (b), suppose first that
$\mathrm{P}\subseteq\mathrm{BPNC}$. The language of pairs $(j,i)$ with
$x^*_i(j)=1$ is in $\mathrm{P}$, hence in $\mathrm{BPNC}$. Computing all bits
in parallel, each with error at most $1/(3n)$ after majority amplification,
and outputting the result only if it satisfies $x_0=1$ and
$x_i=\NOR(x_{i-1},x_{j_i})$ for all $i$, gives an errorless algorithm with
success at least $2/3$ on every instance, hence at every noise level. If
instead $\mathrm{P}\not\subseteq\mathrm{BPNC}$, then
Corollary~\ref{gb:cor64} excludes every $\eps\le1/64$, since success $1/2$
exceeds $1/20$. The final sentence follows from (b).
\end{proof}

At $\eps=1$ the smoothed instance has law $U_n$ whatever the base, so any
errorless average-case algorithm under $U_n$ with success $1/2$ in
randomized $\mathrm{NC}$ puts $1$ in $E$, and then $\eps_c=\inf E$.

\begin{remark}[Scope]\label{gb:rem-scope}
(i) The hardness is conditional, and it concerns worst-case bases; we
exhibit no explicit base on which smoothed instances are provably hard.
(ii) Reliable computation with gates that fail independently goes back to
von Neumann~\cite{gb-vonneumann}; see also Pippenger~\cite{gb-pippenger}.
We have given the argument in full because in Sequential
NOR~\cite{okhotin} one input of every gate is the fixed predecessor, and the
noise acts by resampling a pointer rather than by flipping a gate's output.
(iii) The constants are not optimized. With deeper cells and longer padding
the proof of Theorem~\ref{gb:embed} extends somewhat beyond $1/64$; going
much further would require using that errors are often masked downstream.
\end{remark}

\section{Hardness at inverse-polynomial success}\label{sec:gap}

This section uses the embedding of Section~\ref{sec:gapbase} and only strengthens the conclusion
drawn from it: inverse-polynomial success suffices, and it stays at noise $1/64$. Section~\ref{sec:tree} replaces the embedding and raises the noise level to $3/20$,
Section~\ref{sec:tree2} to $1/4$, Section~\ref{sec:latch}, with a different restoring organ,
to every level below $4/9$, and Section~\ref{sec:s24} to $0.4468$.

\subsection{Ingredients}

A \emph{randomized NC family} is a logspace-uniform family $(A_n)$ of circuits of polynomial size
and polylogarithmic depth; $A_n$ receives a parent array of length $n$ and a string $r$ of
polynomially many random bits. BPNC is the class of languages decided by such families with
two-sided error at most $1/3$. The family is \emph{errorless} if $A_n(j,r)\in\{x^*(j),\bot\}$ for
every $j$ and $r$. We use the composition of smoothings (Lemma~\ref{gb:comp}), the gadgets
(Lemma~\ref{gb:gadgets}), the compiler $\Comp'_a$ (Definition~\ref{gb:compiler64}) and the
embedding at noise $1/64$ with exponent $a$ (Theorem~\ref{gb:embed}): for every NOR circuit $C$
with $s\ge2$ gates and $0\le\eps\le1/64$, with probability at least $1-s^{-a}-(n+1)^{-(a+1)}$ over
the coins and the smoothing, at least five sixths of the output bundle $o(s,1),\dots,o(s,m)$ of
$S_\eps(\Comp'_a(C))$ hold $\val(C)$. Evaluating NOR circuits is P-complete under logspace
reductions~\cite{gb-ladner,ghr}.

\begin{lemma}[length]\label{lem:param}
For $\Comp'_a(C)$ with an integer $a\ge1$: $m<130(a+1)\log_2s$, and $n+1\le s^2$ whenever
$s\ge s_a:=10^{11}(a+2)^4$.
\end{lemma}

\begin{proof}
The least power of two that is at least $y\ge1$ is less than $2y$, which gives the bound on $m$.
By~\eqref{gb:eq-nbound}, $n+1\le10^7(a+2)sm<1.3\cdot10^9(a+2)^2\,s\log_2s$. The function
$s/\log_2s$ increases for $s\ge3$, and $\log_2s_a\le36.6+4\log_2(a+2)\le14.4\,(a+2)$, so
$s_a/\log_2s_a\ge6.9\cdot10^9(a+2)^3$. Hence $n+1\le s^2$ for $s\ge s_a$.
\end{proof}

\subsection{The theorem}

\begin{lemma}[errorlessness is free]\label{lem:errorless}
Let $(A_n)$ be a randomized NC family. Let $A'_n(j,r)$ compute $y=A_n(j,r)$ and output $y$ if
$y_0=1$ and $y_i=\NOR(y_{i-1},y_{j_i})$ for all $1\le i\le n$, and $\bot$ otherwise. Then $(A'_n)$
is an errorless randomized NC family, and $A'_n(j,r)\ne\bot$ exactly when $A_n(j,r)=x^*(j)$.
\end{lemma}

\begin{proof}
The trace is the unique solution of the gate equations, by induction on $i$. The $n+1$ checks
are done in parallel in constant depth.
\end{proof}

\begin{theorem}[hardness at inverse-polynomial success]\label{thm:hard}
Let $K\ge1$ be an integer and $\eps\le1/64$. Suppose there are a randomized NC family $(A_n)$
and an $n_0$ such that, for every $n\ge n_0$ and every base $j$ of length $n$,
\[
\PP_{\jt\sim S_\eps(j),\,r}\bigl[A_n(\jt,r)=x^*(\jt)\bigr]\ge n^{-K}.
\]
Then $\mathrm P\subseteq\mathrm{BPNC}$.
\end{theorem}

\begin{proof}
By Lemma~\ref{lem:errorless} we may assume the family is errorless, with
$\PP[A_n(\jt,r)\ne\bot]\ge n^{-K}$. By Lemma~\ref{gb:comp} the same holds at noise
$\eps^*=1/64$. Put $a=2K+4$.

Let $C$ be a NOR circuit. Prefixing dummy gates $\NOR(c_0,c_0)$, we may assume that its number of
gates $s$ satisfies $s\ge\max(s_a,n_0,3)$. Let $n$ be the length of $\Comp'_a(C)$; then
$n\ge s\ge n_0$, and $n+1\le s^2$ by Lemma~\ref{lem:param}(iii).

\emph{A trial.} Compute $j=\Comp'_a(C)$ with fresh coins; they are exactly uniform because $m$ is
a power of two. Draw $\jt$ from a distribution within total variation $\nu:=(n+1)^{-(K+2)}$ of
$S_{1/64}(j)$: the indicators $\pi_i$ use $6$ random bits each and are exact, and $\rho_i$ is the
residue modulo $i$ of a random integer with $\lceil(K+4)\log_2(n+1)\rceil$ bits, which is within
$i\,(n+1)^{-(K+4)}$ of uniform, so the total error is at most $n^2(n+1)^{-(K+4)}\le\nu$. Run
$A_n(\jt,r)$ with fresh $r$. If the output is $\bot$, the trial abstains, $Y=0$. Otherwise the
output is the trace of $\jt$; the trial votes for the majority bit among its positions
$o(s,1),\dots,o(s,m)$, and we put $Y=+1$ if the vote equals $\val(C)$ and $Y=-1$ if not.

\emph{Its bias.} The success guarantee holds for every base, hence on average over the coins:
$\PP[Y\ne0]\ge n^{-K}-\nu$. A wrong vote forces the majority of the trace of $\jt$ on the output
bundle to be wrong, so by Theorem~\ref{gb:embed}, $\PP[Y=-1]\le\delta+\nu$ with
\[
\delta:=s^{-a}+(n+1)^{-(a+1)}\le2(n+1)^{-a/2}=2(n+1)^{-(K+2)},
\]
using $s\ge(n+1)^{1/2}$. Therefore
\[
\EE Y=\PP[Y\ne0]-2\,\PP[Y=-1]\ge n^{-K}-2\delta-3\nu\ge n^{-K}-7(n+1)^{-(K+2)}\ge\tfrac12\,n^{-K},
\]
since $(n+1)^2\ge14$.

\emph{Amplification.} Run $T=8n^{2K+1}$ independent trials in parallel and output the bit with
more votes. The output is wrong only if $\sum Y\le0$. The $Y$ are independent with values in
$[-1,1]$, so by Hoeffding's inequality~\cite{hoeffding} this has probability at most
$\exp(-T(\EE Y)^2/2)\le e^{-n}\le1/3$.

\emph{Complexity.} The compiler is computable in logspace, the sampling and the votes are in
$\mathrm{NC}^1$, the trials are independent copies of a polylogarithmic-depth circuit, and $T$
is polynomial. So evaluating NOR circuits is in BPNC. That problem is P-complete under logspace
reductions, and BPNC is closed under them.
\end{proof}

\subsection{A zero--one gap}

For a randomized NC family $A$ let
\[
\sigma_A(\eps,n)=\min_{j\text{ of length }n}\ \PP_{\jt\sim S_\eps(j),\,r}\bigl[A_n(\jt,r)=x^*(\jt)\bigr]
\]
be its success probability at noise $\eps$ on the worst base.

\begin{corollary}[zero--one gap]\label{cor:gap}
\begin{enumerate}[label=(\alph*),leftmargin=2.4em]
\item If $\mathrm P\subseteq\mathrm{BPNC}$, there is an errorless randomized NC family $A$ with
$\sigma_A(\eps,n)\ge1-2^{-n}$ for all $\eps\in[0,1]$ and all $n$.
\item If $\mathrm P\not\subseteq\mathrm{BPNC}$, then for every randomized NC family $A$, every
$\eps\le1/64$ and every $K$, we have $\sigma_A(\eps,n)<n^{-K}$ for infinitely many $n$.
\item For every $\eps>3/4$ there is a deterministic logspace algorithm $A$ with
$\sigma_A(\eps,n)\ge1-3/n$.
\end{enumerate}
Consequently, let $E_K$ be the set of $\eps\in[0,1]$ at which some randomized NC family has
$\sigma_A(\eps,n)\ge n^{-K}$ for all large $n$. Then every $E_K$ is an up-set containing
$(3/4,1]$, and $\inf E_K\in\{0\}\cup[1/64,3/4]$, with $\inf E_K=0$ if and only if
$\mathrm P\subseteq\mathrm{BPNC}$. So whether the threshold is $0$ does not depend on $K$.
\end{corollary}

\begin{proof}
(a) The language $\{(j,i):x^*_i(j)=1\}$ is in P, hence in BPNC. The majority of $36(n+1)$
independent runs errs with probability at most $e^{-2(n+1)}\le2^{-n}/(n+1)$, by Hoeffding's
inequality. Compute all $n+1$ bits in this way in parallel and apply the check of
Lemma~\ref{lem:errorless}. This works on every instance, hence at every noise level and for
every base.
(b) is Theorem~\ref{thm:hard}. (c) is Theorem~\ref{thm:logspace}.
For the last claims: $E_K$ is an up-set by Lemma~\ref{lem:errorless} and Lemma~\ref{gb:comp}; it
contains $(3/4,1]$ by (c); under $\mathrm P\subseteq\mathrm{BPNC}$ it is $[0,1]$ by (a), and
otherwise it is contained in $(1/64,1]$ by (b).
\end{proof}

\begin{remark}[relation to Section~\ref{sec:gapbase}]
Theorem~\ref{gb:gap} is the case of errorless families with success $1/2$, and
Corollary~\ref{gb:cor64} needs success $1/20$. The constant $1/20$ comes from a failure
probability of the embedding that a trial cannot tell from a correct answer. With the exponent $a$
of Theorem~\ref{gb:embed} that failure is $O(n^{-(K+2)})$, which is below any prescribed inverse
polynomial success. Lemma~\ref{lem:errorless} shows that for algorithms that output the whole
trace, ``errorless'' is not a restriction.
\end{remark}

\begin{remark}[what is not covered]
\begin{enumerate}[label=(\roman*),leftmargin=2.4em]
\item Algorithms that output a single bit of the trace are not covered at noise $1/64$, as in
Section~\ref{sec:gapbase}: one position of a compiled instance is wrong with probability about $0.12$.
\item The conclusion is $\mathrm P\subseteq\mathrm{BPNC}$, not $\mathrm P\subseteq\mathrm{NC}$,
because the reduction samples the noise. By part (a), $\mathrm P\subseteq\mathrm{BPNC}$ already
gives an errorless randomized NC algorithm for Sequential NOR on every instance, and by Adleman's
argument~\cite{adleman} it implies that every problem in P has non-uniform circuits of polynomial
size and polylogarithmic depth.
\item Part (b) is conditional. It is not evidence for or against
$\mathrm P\subseteq\mathrm{BPNC}$.
\end{enumerate}
\end{remark}

\section{A NOR-tree embedding: hardness up to noise \texorpdfstring{$3/20$}{3/20}}\label{sec:tree}

Section~\ref{sec:gap} used the embedding of Section~\ref{sec:gapbase}. This section
replaces the embedding. The new one has no majority gadget: its restoring organ is a binary tree
of single NOR gates. It is analysed from scratch, and outside this section only Lemma~\ref{gb:comp} is used.
The result is Theorem~\ref{thm:hard2}: the conclusion of Theorem~\ref{thm:hard} holds for every
$\eps\le3/20$, and already for $\eps\le1/10$ by an embedding that needs no density argument.

\emph{Where the gain comes from.} The predecessor input $x_{i-1}$ of a gate is never resampled,
and a gate with $x_{i-1}=1$ outputs $0$ whatever its pointer reads. Lay a tree of NOR gates out in
post-order. Each internal node is then one gate; its predecessor is one child and its pointer goes
to the other. A node whose children should both hold $1$ is wrong only if the predecessor child
is wrong \emph{and} the other child is wrong or misread, so there the noise is multiplied by an
error probability. A node whose children should both hold $0$ is wrong if a child is wrong or if
its pointer misreads a $0$, which needs the resampled pointer to hit a one. The analysis of Section~\ref{sec:gapbase}
charges every resampled gate of a gadget in full, which limits that analysis to about $0.02$.

\subsection{Misreads}

In this section a base $j$ of length $n$ may depend on random \emph{coins}. Fix $\eps\in[0,1]$, let
$\jt=S_\eps(j)$ be built from $(\pi_i,\rho_i)$ as in Section~\ref{sec:setting}, and let $x$ be its
trace. Let $(U_i)_{i\ge1}$ be independent uniform variables on $[0,1]$, independent of everything
else. Let $\cF_g$ be generated by the coins and by $(\pi_i,\rho_i)_{i\le g}$, and let
$\mathcal H_g=\cF_g\vee\sigma(U_i:i\le g)$. For a gate $g\ge1$ define:
\begin{itemize}[leftmargin=2em]
\item the value it reads, $y_g=x_{\jt_g}$, so that $x_g=\NOR(x_{g-1},y_g)$;
\item gate $g$ \emph{misreads} if $\pi_g=1$ and $x_{\rho_g}\ne x_{j_g}$; then $y_g=1\oplus x_{j_g}$, and
otherwise $y_g=x_{j_g}$, the value at its listed pointer;
\item $f_g$, the fraction of positions $p<g$ with $x_p\ne x_{j_g}$, which is $\cF_{g-1}$-measurable;
\item $V_g=\eps f_g\,U_g$ if $g$ misreads, and $V_g=\eps f_g+(1-\eps f_g)\,U_g$ otherwise.
\end{itemize}
For a constant $\bar a\in(0,1]$, fixed below, the \emph{flags} of gate $g$ are
\[
A_g=\{V_g<\eps\bar a\},\qquad B_g=\{V_g<\eps\}.
\]

\begin{lemma}[misreads]\label{lem:misread}
\begin{enumerate}[label=(\alph*),leftmargin=2.4em]
\item The variables $V_g$, $g\ge1$, are independent and uniform on $[0,1]$, and for every $g_0$ the
family $(V_g)_{g\ge g_0}$ is independent of $\mathcal H_{g_0-1}$. In particular the flags of
different gates are independent, they are independent of the coins, and $\PP[A_g]=\eps\bar a$,
$\PP[B_g]=\eps$.
\item Almost surely, gate $g$ misreads if and only if $V_g<\eps f_g$. Hence a gate that misreads has
$B_g$; and if moreover $x_{j_g}=0$ and $O_{g-1}\le\bar a\,g$, it has $A_g$.
\item $O_{g-1}\le(g+1)/2$ for every $g$, so $O_{g-1}\le0.501\,g$ for $g\ge500$.
\end{enumerate}
\end{lemma}

\begin{proof}
(a) The triple $(\pi_g,\rho_g,U_g)$ is independent of $\mathcal H_{g-1}$, and $x_p$ for $p<g$ and
$j_g$ are $\mathcal H_{g-1}$-measurable. Given $\mathcal H_{g-1}$, gate $g$ misreads with
probability $\eps f_g$, so for $v\in[0,1]$
\[
\PP[V_g\le v\mid\mathcal H_{g-1}]=\eps f_g\min\Bigl(1,\frac{v}{\eps f_g}\Bigr)+(1-\eps f_g)\max\Bigl(0,\frac{v-\eps f_g}{1-\eps f_g}\Bigr)=v
\]
(with the first term read as $0$ when $\eps f_g=0$). So $V_g$ is uniform and independent of
$\mathcal H_{g-1}$, and it is $\mathcal H_g$-measurable. Induction on the number of gates gives the
joint statement.
(b) If $g$ misreads then $V_g=\eps f_gU_g<\eps f_g$ unless $U_g=1$; otherwise $V_g\ge\eps f_g$.
Always $f_g<1$, since position $j_g$ agrees with itself. If $x_{j_g}=0$, then $f_g=O_{g-1}/g$.
(c) A one is always followed by a zero.
\end{proof}

So the smoothed instance behaves like the base with read errors. Gate $g$ computes the NOR of
$x_{g-1}$ and $x_{j_g}$ unless it misreads. A misread always raises the flag $B_g$, and the misread
of a $0$ raises the smaller flag $A_g$, as long as ones are not too dense.

\subsection{The compiler}

Let $C$ be a NOR circuit with gates $g_t=\NOR(u_t,v_t)$, $t\le s$. Fix an odd depth $d$, a bundle
size $m$, a padding ratio $\rho\ge0$ and a preamble length $\mathrm{pre}\ge499$. The base is made
of \emph{zero gates}, which are gates with pointer $0$, and of blocks defined as follows. For an
input $w$ (an earlier gate or a constant):
\begin{itemize}[leftmargin=2em]
\item $T_0(w)$ is a \emph{leaf}: a zero gate $z$ followed by a gate $\ell$ with pointer $o(w,r)$,
where $r$ is a fresh uniform element of $\{1,\dots,m\}$ (a coin);
\item $T_{h+1}(w)$ consists of two blocks $T_h(w)$, with independent coins, followed by one gate
$\kappa$ whose pointer is the last gate of the first of them;
\item the \emph{cell} $C(t,l)$ consists of $T_d(u_t)$, then $T_d(v_t)$, then one gate $\tau$ whose
pointer is the last gate of $T_d(u_t)$.
\end{itemize}
Here $o(t',l')$ is the position of the gate $\tau$ of $C(t',l')$, and $o(c_1,\cdot)=0$,
$o(c_0,\cdot)=1$. The base $\mathrm{Tree}(C)$ is: gate $1$ and further zero gates up to position
$\mathrm{pre}$; then, for $t=1,\dots,s$ and $l=1,\dots,m$ in this order, $\rho c$ zero gates
followed by the cell $C(t,l)$. A block $T_h$ has $3\cdot2^h-1$ gates, a cell has $c=6\cdot2^d-1$,
and the length is $n=\mathrm{pre}+(1+\rho)\,c\,sm$. All positions are explicit arithmetic
expressions, so the base is computable in logspace from $C$ and the coins.

\begin{lemma}[exactness]\label{lem:exact}
In the trace of $\mathrm{Tree}(C)$ itself, the last gate of every block $T_h(w)$ holds
$1\oplus\val(w)$ if $h$ is even and $\val(w)$ if $h$ is odd, and $x_{o(t,l)}=\val(g_t)$ for all
$t,l$ and all coins.
\end{lemma}

\begin{proof}
Induction along the layout. A zero gate holds $\NOR(\cdot,x_0)=0$, so the gate $\ell$ of a leaf
holds $\NOR(0,\val(w))$. The gate $\kappa$ holds the NOR of the last gates of its two blocks,
which agree. As $d$ is odd, $\tau$ holds $\NOR(\val(v_t),\val(u_t))$.
\end{proof}

The value of Lemma~\ref{lem:exact} is the \emph{intended value} of a block or cell; it depends on
$C$ only. In the smoothed instance a block or cell is \emph{wrong} if its last gate differs from
its intended value. A wrong cell is called \emph{bad}; positions $0$ and $1$ are never bad, because
$x_0=1$ and $x_1=0$ in every instance. Let $\beta_t$ be the fraction of bad cells among
$C(t,1),\dots,C(t,m)$, and $\beta=0$ for the constants.

\subsection{Dominating events}

For each block define an event $D$ by recursion. Let $b=\val(w)$.
\begin{itemize}[leftmargin=2em]
\item \emph{Leaf} with gates $z,\ell$ and copy $p=o(w,r)$: if $b=0$,
$D=\{p\text{ bad}\}\cup B_z\cup A_\ell$; if $b=1$, $D=\{p\text{ bad}\}\cup B_\ell$.
\item \emph{Node} $\kappa$ over blocks $T'$ (pointer) and $T''$ (predecessor), both of intended
value $e$: if $e=0$, $D=D(T')\cup D(T'')\cup A_\kappa$; if $e=1$,
$D=D(T'')\cap\bigl(D(T')\cup B_\kappa\bigr)$.
\item \emph{Cell} with top gate $\tau$ and blocks $T_u=T_d(u_t)$, $T_v=T_d(v_t)$: according to
$(\val(u_t),\val(v_t))$,
\[
\begin{array}{ll}
(0,0):\ D=D(T_u)\cup D(T_v)\cup A_\tau, & (1,1):\ D=D(T_v)\cap\bigl(D(T_u)\cup B_\tau\bigr),\\[2pt]
(1,0):\ D=D(T_u)\cup B_\tau, & (0,1):\ D=D(T_v).
\end{array}
\]
\end{itemize}
Let $\Delta$ be the event that $O_{g-1}\le\bar a\,g$ for every gate $g$ of every cell.

\begin{lemma}[domination]\label{lem:dom}
On $\Delta$, every wrong block has its event $D$, and every bad cell has its event $D$.
\end{lemma}

\begin{proof}
Induction along the layout, using $x_g=\NOR(x_{g-1},y_g)$ and Lemma~\ref{lem:misread}(b).

\emph{Leaf.} The gate $z$ reads $x_0=1$ unless it misreads, so $x_z=1$ implies $B_z$. If $b=0$ the
intended value is $1$. If $x_\ell=0$, then $x_z=1$, or $y_\ell=1$; in the second case $x_p=1$, so
$p$ is bad, or $x_p=0$ and $\ell$ misreads a $0$, which on $\Delta$ gives $A_\ell$. If $b=1$ the
intended value is $0$. If $x_\ell=1$ then $y_\ell=0$, so $x_p=0$ and $p$ is bad, or $\ell$ misreads,
which gives $B_\ell$.

\emph{Node, $e=0$} (intended value $1$). If $x_\kappa=0$, then $x_{\kappa-1}=1$, so $T''$ is wrong,
or $y_\kappa=1$; then the last gate of $T'$ holds $1$, so $T'$ is wrong, or it holds $0$ and
$\kappa$ misreads a $0$, which gives $A_\kappa$.

\emph{Node, $e=1$} (intended value $0$). If $x_\kappa=1$, then $x_{\kappa-1}=0$, so $T''$ is wrong,
and $y_\kappa=0$; then $T'$ is wrong or $\kappa$ misreads, which gives $B_\kappa$.

\emph{Cell.} The cases $(0,0)$ and $(1,1)$ are the two node cases. In the other two the intended
value is $0$, and $x_\tau=1$ forces $x_{\tau-1}=0$ and $y_\tau=0$. For $(1,0)$, $y_\tau=0$ means
that $T_u$ is wrong or $\tau$ misreads. For $(0,1)$, $x_{\tau-1}=0$ means that $T_v$ is wrong.
\end{proof}

\noindent(The script \texttt{domination\_test.py} checks this implication pathwise on random small
instances, with each flag replaced by the smallest event it is known to contain;
Section~\ref{sec:checks}.)

\subsection{Probabilities}

Let $\mathcal G_t$ be generated by the coins of the cells of bundles $1,\dots,t$ and by
$(\pi_i,\rho_i,U_i)$ for the positions $i$ up to the end of bundle $t$. The trace up to that
position is $\mathcal G_t$-measurable, and so are the sets of bad cells of bundles $1,\dots,t$.

Fix \emph{tolerances} $\omega_0,\omega_1\in(0,1/2)$, one for gates of value $0$ and one for gates
of value $1$. For $p\in[0,1]$ put
\[
\mathrm{or}(p)=1-(1-p)^2(1-\eps\bar a),\qquad\mathrm{and}(p)=p\,\bigl(1-(1-p)(1-\eps)\bigr),
\]
which are nondecreasing in $p$, $\eps$ and $\bar a$. Define
\begin{align*}
p^{(1)}_0&=1-(1-\omega_1)(1-\eps), & p^{(1)}_h&=\mathrm{or}\bigl(p^{(1)}_{h-1}\bigr)\ (h\text{ odd}),\quad \mathrm{and}\bigl(p^{(1)}_{h-1}\bigr)\ (h\text{ even}),\\
p^{(0)}_0&=1-(1-\omega_0)(1-\eps)(1-\eps\bar a), & p^{(0)}_h&=\mathrm{and}\bigl(p^{(0)}_{h-1}\bigr)\ (h\text{ odd}),\quad \mathrm{or}\bigl(p^{(0)}_{h-1}\bigr)\ (h\text{ even}),
\end{align*}
and
\[
r_1=p^{(1)}_d,\quad r_0=p^{(0)}_d,\qquad Q_1=1-(1-r_0)^2(1-\eps\bar a),\quad Q_0=1-(1-r_1)(1-\eps).
\]
Let $I_t=\{\beta_{u_t}\le\omega_{\val(u_t)},\ \beta_{v_t}\le\omega_{\val(v_t)}\}$, an event in
$\mathcal G_{t-1}$.

\begin{lemma}[conditional law]\label{lem:law}
Conditionally on $\mathcal G_{t-1}$, the events $D(C(t,l))$, $l=1,\dots,m$, are independent, and on
$I_t$ each has conditional probability at most $Q_{\val(g_t)}$.
\end{lemma}

\begin{proof}
The event $D$ of a block is a function of the flags of its gates, of its coins, and of the set of
bad cells of the bundle it reads, which is $\mathcal G_{t-1}$-measurable. By
Lemma~\ref{lem:misread}(a) the flags of the gates of bundle $t$ are independent of each other, of
all coins and of $\mathcal G_{t-1}$, and the coins of bundle $t$ are independent of
$\mathcal G_{t-1}$. Disjoint blocks have disjoint gates and coins. So, given $\mathcal G_{t-1}$, the
events $D$ of disjoint blocks are independent, and the flag of a node is independent of the events
of its two blocks.

For a leaf over $w$, the copy is bad with conditional probability $\beta_w$, because $r$ is
uniform. With $b=\val(w)$, on $\{\beta_w\le\omega_b\}$ the event $D$ of the leaf has conditional
probability at most $p^{(b)}_0$. A block $T_{h-1}(w)$ has intended value $1\oplus b$ for odd $h$
and $b$ for even $h$. So the node of $T_h(w)$ is of the kind $e=0$ exactly when ($b=1$ and $h$ is
odd) or ($b=0$ and $h$ is even), and then its event has probability
$1-(1-p')(1-p'')(1-\eps\bar a)$; otherwise it is of the kind $e=1$, with probability
$p''\,(1-(1-p')(1-\eps))$. Here $p',p''$ are the probabilities for its two blocks. Both expressions
are nondecreasing in $p',p''$, so induction on $h$ bounds the event of $T_h(w)$ by $p^{(b)}_h$, and
those of $T_d(u_t)$, $T_d(v_t)$ by $r_{\val(u_t)}$, $r_{\val(v_t)}$. For the cell, the four cases
give at most $1-(1-r_0)^2(1-\eps\bar a)=Q_1$ for $(0,0)$, where $\val(g_t)=1$, and at most
$r_1(1-(1-r_1)(1-\eps))$, $1-(1-r_1)(1-\eps)$ and $r_1$ for $(1,1)$, $(1,0)$ and $(0,1)$, where
$\val(g_t)=0$; all three are at most $Q_0$. A constant input has $\beta=0$, which only lowers the
bounds.
\end{proof}

\begin{theorem}[embedding]\label{thm:embed}
Let $\eps_0\in(0,1]$ and suppose that, with $\eps_0$ in place of $\eps$, the numbers above satisfy
$Q_0<\omega_0$ and $Q_1<\omega_1$. Let
$\varkappa=\min\bigl(D(\omega_0\,\|\,Q_0),D(\omega_1\,\|\,Q_1)\bigr)$, where $D(p\|q)$ is the
relative entropy of Lemma~\ref{gb:density} (not a dominating event). Then for every
$\eps\le\eps_0$ and every NOR circuit $C$ with $s$ gates, over the coins and the smoothing,
\[
\PP\Bigl[\Delta\cap\bigl\{\exists\,t\le s:\ \beta_t>\omega_{\val(g_t)}\bigr\}\Bigr]\le s\,e^{-\varkappa m}.
\]
\end{theorem}

\begin{proof}
The quantities $Q_0,Q_1$ are nondecreasing in $\eps$, so their values at $\eps$ are at most those at
$\eps_0$. Let $H_t=\{\beta_{t'}\le\omega_{\val(g_{t'})}\text{ for all }t'\le t\}$; then
$H_{t-1}\subseteq I_t$. Fix $t$ and let $b=\val(g_t)$. On $\Delta$, by Lemma~\ref{lem:dom}, the
number of bad cells of bundle $t$ is at most the number of $l$ with $D(C(t,l))$. By
Lemma~\ref{lem:law}, on $I_t$ and given $\mathcal G_{t-1}$ that number is stochastically dominated
by a binomial variable with parameters $m$ and $Q_b<\omega_b$. The Chernoff bound gives
\[
\PP\bigl[\{\beta_t>\omega_b\}\cap\Delta\mid\mathcal G_{t-1}\bigr]\le e^{-m\,D(\omega_b\|Q_b)}\le e^{-\varkappa m}\qquad\text{on }I_t .
\]
Since $H_s^c\cap\Delta\subseteq\bigcup_{t\le s}\bigl(H_{t-1}\cap\{\beta_t>\omega_{\val(g_t)}\}\cap\Delta\bigr)$,
a union bound finishes the proof.
\end{proof}

\subsection{Two choices of parameters}

\begin{corollary}[noise $1/10$, no padding]\label{cor:tenth}
Take $\rho=0$, $\mathrm{pre}=499$, $d=15$, $\bar a=0.501$, $\omega_0=3/10$ and $\omega_1=7/50$, and
let $m$ be a power of two with $m\ge56(a+1)\log_2s$, where $a\ge1$. Then for every $\eps\le1/10$
and every NOR circuit $C$ with $s\ge2$ gates,
\[
\PP\bigl[\exists\,t\le s:\ \beta_t>\omega_{\val(g_t)}\bigr]\le s^{-a}.
\]
In particular, with at least the complementary probability, at least seven tenths of the
output positions $o(s,1),\dots,o(s,m)$ hold $\val(C)$. A cell has $196{,}607$ gates.
\end{corollary}

\begin{proof}
Every cell gate is at a position $g\ge500$, so $\Delta$ is the sure event by
Lemma~\ref{lem:misread}(c). At $\eps_0=1/10$ the recursion gives $Q_0\le0.2311$ and
$Q_1\le0.0900$ (computed in exact arithmetic with upward rounding; Section~\ref{sec:checks}), and
$D(0.3\,\|\,0.2311)\ge0.0125$, $D(0.14\,\|\,0.09)\ge0.0125$. So $\varkappa\ge0.0125$, and
$s\,e^{-\varkappa m}\le s\cdot s^{-0.0125\cdot56\,(a+1)/\ln2}\le s^{-a}$.
\end{proof}

For higher noise we lower the density of ones. Zero gates hold $0$ unless they misread, so padding
with zero gates brings the density down to about $\eps$, and then the flags $A_g$ become rare.

\begin{lemma}[density under zero padding]\label{lem:zeropad}
Let $\rho=99$, $\eps\le3/20$, $a\ge1$, and suppose that $c\ge102$ and
$\mathrm{pre}\ge g^\circ:=1200(a+2)\ln(n+1)+10^4$. Then, with $\bar a=0.17$,
$\ \PP[\Delta^c]\le(n+1)^{-(a+1)}$.
\end{lemma}

\begin{proof}
Fix $g>\mathrm{pre}$ and consider the positions below $g$. A zero gate $p$ reads $x_0=1$ unless it
misreads, so $x_p=1$ implies that $p$ misreads, hence $V_p<\eps\le3/20$. Let $N_g$ be the number of
zero gates $p<g$ with $V_p<3/20$; it is dominated by a binomial variable with parameters $g$ and
$3/20$. Every cell that meets $[0,g)$ is preceded by its $\rho c$ zero gates, so if $k$ cells meet
$[0,g)$, there are at most $kc$ cell gates and at least $k\rho c$ zero gates below $g$. Hence
there are at most $g/100$ cell gates below $g$, and $k\le g/(99c)$. Ones are never adjacent, so the
cell gates below $g$ hold at most $(g/100+k)/2$ ones. Counting position $0$,
\[
O_{g-1}\le1+N_g+\frac{g}{200}+\frac{g}{198\,c}\le N_g+0.0052\,g ,
\]
since $g\ge10^4$ and $c\ge102$. By the Chernoff bound and $D(0.1648\,\|\,0.15)\ge0.000836$,
\[
\PP[N_g>0.1648\,g]\le e^{-0.000836\,g}\le(n+1)^{-(a+2)}\qquad\text{for }g\ge g^\circ ,
\]
because $1200\cdot0.000836\ge1$. A union bound over $g\le n$ gives the claim, as
$0.1648+0.0052=0.17$.
\end{proof}

\begin{definition}[$\mathrm{Tree}_a$]\label{def:treea}
For an integer $a\ge1$ and a circuit with $s\ge2$ gates, $\mathrm{Tree}_a(C)$ is the base
$\mathrm{Tree}(C)$ with $d=19$ (so $c=3{,}145{,}727$), $\rho=99$,
\begin{align*}
m&=\text{the least power of two with }m\ge33(a+1)\log_2s,\\
Z&=4\cdot10^8(a+2)^2sm,\qquad b=\lceil\log_2Z\rceil,\qquad\mathrm{pre}=840(a+2)b+10^4 .
\end{align*}
\end{definition}

\begin{lemma}[parameters]\label{lem:param2}
For $\mathrm{Tree}_a(C)$: (i) $n+1\le Z$; (ii) $\mathrm{pre}\ge g^\circ$; (iii)
$m<66(a+1)\log_2s$, and $n+1\le s^2$ whenever $s\ge s^\circ_a:=10^{13}(a+2)^4$.
\end{lemma}

\begin{proof}
(i) $b\le29.6+2\log_2(a+2)+\log_2(sm)$, so
$\mathrm{pre}+1\le(24864+1680+840+10001)(a+2)^2sm$, and $100c\le3.15\cdot10^8$. Hence
$n+1=\mathrm{pre}+1+100\,c\,sm\le4\cdot10^8(a+2)^2sm$.
(ii) $840\,b\ge840\log_2Z\ge1200\ln Z\ge1200\ln(n+1)$, because $840/\ln2\ge1200$.
(iii) As in Lemma~\ref{lem:param}: $n+1\le Z<2.64\cdot10^{10}(a+2)^3s\log_2s$, and
$\log_2s^\circ_a\le16.6\,(a+2)$, so $s^\circ_a/\log_2s^\circ_a\ge6\cdot10^{11}(a+2)^3$.
\end{proof}

\begin{corollary}[noise $3/20$]\label{cor:threetwentieths}
Let $a\ge1$, $\eps\le3/20$, and let $C$ be a NOR circuit with $s\ge2$ gates. For the trace of
$S_\eps(\mathrm{Tree}_a(C))$, with $\omega_0=3/10$ and $\omega_1=1/10$,
\[
\PP\bigl[\exists\,t\le s:\ \beta_t>\omega_{\val(g_t)}\bigr]\le s^{-a}+(n+1)^{-(a+1)} .
\]
In particular, with at least the complementary probability, at least seven tenths of the
positions $o(s,1),\dots,o(s,m)$ hold $\val(C)$.
\end{corollary}

\begin{proof}
Take $\bar a=0.17$. By Lemmas~\ref{lem:zeropad} and~\ref{lem:param2}(ii),
$\PP[\Delta^c]\le(n+1)^{-(a+1)}$. At $\eps_0=3/20$ the recursion gives $Q_0\le0.2099$ and
$Q_1\le0.0493$, and $D(0.3\,\|\,0.2099)\ge0.0214$, $D(0.1\,\|\,0.0493)\ge0.0214$. So
$\varkappa\ge0.0214$, and Theorem~\ref{thm:embed} bounds the rest by
$s\cdot s^{-0.0214\cdot33\,(a+1)/\ln2}\le s^{-a}$.
\end{proof}

\subsection{Hardness}

\begin{theorem}[hardness up to noise $3/20$]\label{thm:hard2}
Theorem~\ref{thm:hard} holds with $\eps\le3/20$ in place of $\eps\le1/64$: if for some
$\eps\le3/20$ a randomized NC family outputs the trace with probability at least $n^{-K}$ on the
$\eps$-smoothing of every base of every length $n\ge n_0$, then $\mathrm P\subseteq\mathrm{BPNC}$.
\end{theorem}

\begin{proof}
Repeat the proof of Theorem~\ref{thm:hard} with $\eps^*=3/20$, with $\mathrm{Tree}_a$ in place of
$\Comp'_a$, with $s\ge\max(s^\circ_a,n_0,3)$, and with Corollary~\ref{cor:threetwentieths} and
Lemma~\ref{lem:param2}(iii) in place of Theorem~\ref{gb:embed} and
Lemma~\ref{lem:param}(iii). Three points change. The coins of a cell are $2^{d+1}$ elements of
$\{1,\dots,m\}$, exactly uniform as $m$ is a power of two. The indicators $\pi_i$ are not dyadic
now; they are drawn by comparing a random integer of $\lceil(K+4)\log_2(n+1)\rceil$ bits with a
threshold, with error at most $(n+1)^{-(K+4)}$ each, so the total sampling error is at most
$(n+n^2)(n+1)^{-(K+4)}\le\nu$. And the vote of a trial is the majority bit of the output bundle,
which is $\val(C)$ whenever at most $3/10$ of its cells are bad. With
$\delta=s^{-a}+(n+1)^{-(a+1)}$ and $a=2K+4$ the rest is unchanged.
\end{proof}

\begin{corollary}[the gap, with the new lower end]\label{cor:gap2}
Corollary~\ref{cor:gap} holds with $3/20$ in place of $1/64$. In particular
$\inf E_K\in\{0\}\cup[3/20,3/4]$ for every $K$, and if
$\mathrm P\not\subseteq\mathrm{BPNC}$, the critical noise level $\eps_c$ satisfies
$3/20\le\eps_c\le3/4$.
\end{corollary}

\begin{proof}
Use Theorem~\ref{thm:hard2} in place of Theorem~\ref{thm:hard} in the proof of
Corollary~\ref{cor:gap}. For $\eps_c$: an errorless family with success $1/2$ at some noise
$\eps\le3/20$ would give $\mathrm P\subseteq\mathrm{BPNC}$ by Theorem~\ref{thm:hard2}.
\end{proof}

\begin{remark}[comparison of the embeddings]
The six embeddings compare as follows.
\begin{center}
\begin{tabular}{lccccc}
\toprule
 & noise & gates per cell & padding & $m$ ($a=1$) & density\\
\midrule
Thm.~\ref{gb:embed} & $1/64$ & $413{,}340$ & none & $130\log_2s$ & $\ge0.19$ needed\\
Cor.~\ref{cor:tenth} & $1/10$ & $196{,}607$ & none & $112\log_2s$ & not used\\
Cor.~\ref{cor:threetwentieths} & $3/20$ & $3{,}145{,}727$ & $\times99$ & $66\log_2s$ & $\le0.17$ needed\\
Cor.~\ref{cor:twentythree} & $1/4$ & $3.2\cdot10^{16}$ & $\times999$ & $750\log_2s$ & $\varpi\pm\frac1{400}$ needed\\
Cor.~\ref{cor:twofifths} & $2/5$ & $12{,}160$ & $\times1999$ & $2\cdot10^6\log_2s$ & $\varpi\pm\frac1{500}$ needed\\
Prop.~\ref{prop:s24cert} & $0.4468$ & $576{,}096$ & $\times2\cdot10^6$ & $1.3\cdot10^{13}\log_2s$ & $\varpi\pm10^{-6}$ needed\\
\bottomrule
\end{tabular}
\end{center}
With deeper cells the proof of Theorem~\ref{gb:embed} extends somewhat beyond $1/64$
(Remark~\ref{gb:rem-scope}). The last row is the latch
embedding of Section~\ref{sec:latch}, which has no cells; the entry counts the $152$ chains of
$80$ gates that carry one copy of a gate through $76$ stages on two rails. The row below it is the
embedding of Section~\ref{sec:s24} at noise $3661/8192$: $6001$ chains of $96$ gates, one per
stage, each feeding five rails; its padding is roughly the least $\rho$ for which
Lemma~\ref{lem:dens2any} gives density error $10^{-6}$.
\end{remark}

\begin{remark}[how far this goes]\label{rem:limits}
The recursion of this section stops working, for any depth and any tolerances, at about
$\eps=0.13$ without padding and at about $0.18$ with zero padding (evaluated numerically up to
$d=31$). The construction itself does better than the proof: in simulation, cells of depth $11$
without padding kept a correct majority in every bundle up to $\eps=0.20$ and lost some at $0.25$
(Section~\ref{sec:checks}). The proof loses in two places. It bounds a misread by $\eps$ or
$\eps\bar a$, while the true probabilities are $\eps$ times the density of zeros or of ones, about
$0.65$ and $0.35$ without padding. And it uses union bounds throughout.

Refined cells (a zero gate that survives one misread, leaves that read two copies, a top gate that
reads its pointer twice), together with two-sided density bounds, are analysed in
Section~\ref{sec:tree2}, which proves noise $1/4$. For the cells of this section, with the exact
misread probabilities in place of the flags, the recursion stops at $0.214$ at depth $31$ and at
about $0.2195$ for large depth, because the stable pair of error levels of the NOR tree (one for
blocks that should hold $0$, one for blocks that should hold $1$) disappears there. So these cells
do not reach noise $1/4$ with this kind of recursion, however the flags are bounded; the refined
cells of Section~\ref{sec:tree2} reach $1/4$, and their recursion stops near $0.259$
(script \texttt{tree2\_explore.py}). The upper end $3/4$ of the interval is the subject of Section~\ref{sec:adaptive}.
\end{remark}

\section{A refined embedding: hardness up to noise \texorpdfstring{$1/4$}{1/4}}\label{sec:tree2}

This section raises the noise level of Theorem~\ref{thm:hard2} from $3/20$ to $1/4$. It keeps
the NOR tree of Section~\ref{sec:tree} and changes three things.
\begin{itemize}[leftmargin=2em]
\item \emph{Cells.} The zero gate of a leaf is protected by a gate that holds $1$; a leaf reads
two copies of its input; and the top of a cell reads its pointer twice.
\item \emph{Flags.} A misread of a $1$ is charged $\eps$ times the density of zeros, not $\eps$.
\item \emph{Density.} With $999$ zero gates per cell gate, the density of ones stays within
$1/400$ of the limiting density $\varpi(\eps)$ of Lemma~\ref{kp:limdens}, from below and from above
(Lemma~\ref{lem:dens2}).
\end{itemize}
The deterministic part, exactness and domination for the new cells, is checked in Lean.

\subsection{The refined compiler}

As in Section~\ref{sec:tree}, a \emph{zero gate} is a gate with pointer $0$. A gate with pointer
$1$ reads $x_1$, which is $0$ in every instance, because gate $1$ is a NOT gate with $j_1=0$.

\begin{definition}[$\mathrm{Tree}'$]\label{def:treeprime}
Let $C$, $d$ (odd), $m$, $\rho$ and $\mathrm{pre}$ be as in Section~\ref{sec:tree}. For an input $w$:
\begin{itemize}[leftmargin=2em]
\item $T'_0(w)$ is a \emph{leaf} of six gates $z_0,\gamma,z,\ell_1,\ell_2,\ell_3$, with pointers $0$,
$1$, $0$, $o(w,r_2)$, $1$, $o(w,r_1)$, where $r_1,r_2$ are independent uniform elements of
$\{1,\dots,m\}$ (coins);
\item $T'_{h+1}(w)$ consists of two blocks $T'_h(w)$, with independent coins, followed by one gate
$\kappa$ whose pointer is the last gate of the first of them;
\item the cell $C'(t,l)$ consists of $T'_d(u_t)$, then $T'_d(v_t)$, then three gates
$\tau_1,\tau_2,\tau_3$ with pointers: the last gate of $T'_d(u_t)$, then $1$, then the last gate of
$T'_d(u_t)$ again.
\end{itemize}
Here $o(t',l')$ is the position of the gate $\tau_3$ of $C'(t',l')$, and $o(c_1,\cdot)=0$,
$o(c_0,\cdot)=1$. The base $\mathrm{Tree}'(C)$ is: zero gates at positions $1,\dots,\mathrm{pre}$
(gate $1$ is then the NOT gate); then, for $t=1,\dots,s$ and $l=1,\dots,m$ in this order, $\rho c$
zero gates followed by the cell $C'(t,l)$. A block $T'_h$ has $7\cdot2^h-1$ gates, a cell has
$c=7\cdot2^{d+1}+1$, and $n=\mathrm{pre}+(1+\rho)\,c\,sm$. The gates of the cells are the
\emph{cell gates}; all others are zero gates.
\end{definition}

\begin{lemma}[exactness]\label{lem:exact2}
In the trace of $\mathrm{Tree}'(C)$ itself, the last gate of every block $T'_h(w)$ holds
$1\oplus\val(w)$ if $h$ is even and $\val(w)$ if $h$ is odd, and $x_{o(t,l)}=\val(g_t)$, for all
$t,l$ and all coins.
\end{lemma}

\begin{proof}
In a leaf, $z_0$ holds $0$, $\gamma$ holds $\NOR(0,x_1)=1$, $z$ holds $0$, $\ell_1$ holds
$\NOR(0,b)=1\oplus b$ with $b=\val(w)$, $\ell_2$ holds $\NOR(1\oplus b,0)=b$, and $\ell_3$ holds
$\NOR(b,b)=1\oplus b$. Nodes are as in Lemma~\ref{lem:exact}. In a cell with $u=\val(u_t)$ and
$v=\val(v_t)$, $\tau_1=\NOR(v,u)$, $\tau_2=1\oplus\tau_1=u\vee v$ and
$\tau_3=\NOR(u\vee v,u)=\NOR(u,v)$.
\end{proof}

\noindent For this section, fix constants $0<\bar a\le\bar b$ with $\eps\bar b\le1$ and redefine the flags of a gate
$g$ as $A_g=\{V_g<\eps\bar a\}$ and $B_g=\{V_g<\eps\bar b\}$, with $V_g$ as in
Section~\ref{sec:tree}. Let
\[
\Delta'=\bigl\{(1-\bar b)\,g\le O_{g-1}\le\bar a\,g\ \text{ for every cell gate }g\bigr\}.
\]

\begin{lemma}[misreads, two-sided]\label{lem:misread2}
The flags of different gates are independent, independent of the coins, and
$\PP[A_g]=\eps\bar a$, $\PP[B_g]=\eps\bar b$. On $\Delta'$, a cell gate that misreads has $B_g$, and
if it misreads a $0$ it has $A_g$.
\end{lemma}

\begin{proof}
The first sentence is Lemma~\ref{lem:misread}(a). By Lemma~\ref{lem:misread}(b), gate $g$ misreads
if and only if $V_g<\eps f_g$, where $f_g$ is the fraction of positions below $g$ that hold
$1\oplus x_{j_g}$. If $x_{j_g}=0$ then $f_g=O_{g-1}/g\le\bar a\le\bar b$ on $\Delta'$; if
$x_{j_g}=1$ then $f_g=(g-O_{g-1})/g\le\bar b$ on $\Delta'$.
\end{proof}

\noindent Dominating events are defined by recursion, with $D$ of a node as in
Section~\ref{sec:tree}. For a leaf over $w$ with $b=\val(w)$ and copies $p_1=o(w,r_1)$,
$p_2=o(w,r_2)$:
\begin{align*}
b=0:&\quad D=\{p_1\text{ bad}\}\cup\{p_2\text{ bad}\}\cup\bigl((B_{z_0}\cup A_\gamma)\cap B_z\bigr)\cup A_{\ell_1}\cup A_{\ell_3},\\
b=1:&\quad D=\bigl(\{p_1\text{ bad}\}\cup B_{\ell_3}\bigr)\cap\bigl(\{p_2\text{ bad}\}\cup B_{\ell_1}\cup A_{\ell_2}\bigr).
\end{align*}
For a cell with blocks $T_u=T'_d(u_t)$, $T_v=T'_d(v_t)$, according to $(\val(u_t),\val(v_t))$:
\[
\begin{array}{ll}
(0,0):\ D=D(T_u)\cup D(T_v)\cup A_{\tau_1}\cup A_{\tau_3}, & (1,0),(1,1):\ D=D(T_u)\cup\bigl(B_{\tau_3}\cap(B_{\tau_1}\cup A_{\tau_2})\bigr),\\[2pt]
(0,1):\ D=D(T_v)\cup A_{\tau_2}.
\end{array}
\]

\begin{lemma}[domination]\label{lem:dom2}
On $\Delta'$, every wrong block of $\mathrm{Tree}'(C)$ has its event $D$, and every bad cell has
its event $D$.
\end{lemma}

\begin{proof}
Induction along the layout, as for Lemma~\ref{lem:dom}, with Lemma~\ref{lem:misread2} in place of
Lemma~\ref{lem:misread}(b). Nodes are as there. Write $y$ for the value a gate reads.

\emph{Leaf, $b=0$} (intended value $1$). If $\ell_3$ holds $0$, then $\ell_2$ holds $1$ or
$y_{\ell_3}=1$. The second means that $p_1$ is bad or that $\ell_3$ misreads a $0$. The first
needs $\ell_1$ to hold $0$, so $z$ holds $1$ or $y_{\ell_1}=1$; the latter means that $p_2$ is bad
or that $\ell_1$ misreads a $0$. And $z$ holds $1$ only if $\gamma$ holds $0$ and $z$ misreads,
while $\gamma$ holds $0$ only if $z_0$ holds $1$, which needs a misread of $z_0$, or $\gamma$
misreads the $0$ at position $1$.

\emph{Leaf, $b=1$} (intended value $0$). If $\ell_3$ holds $1$, then $y_{\ell_3}=0$, so $p_1$ is bad
or $\ell_3$ misreads; and $\ell_2$ holds $0$, so $\ell_2$ misreads the $0$ at position $1$, or
$\ell_1$ holds $1$, which needs $y_{\ell_1}=0$, so that $p_2$ is bad or $\ell_1$ misreads.

\emph{Cell.} Let $L_u,L_v$ be the last gates of $T_u,T_v$. For $(0,0)$ (intended $1$): if $\tau_3$
holds $0$, then $y_{\tau_3}=1$, so $T_u$ is wrong or $\tau_3$ misreads a $0$; or $\tau_2$ holds $1$,
so $\tau_1$ holds $0$, so $L_v$ holds $1$, or $y_{\tau_1}=1$, so $T_u$ is wrong or $\tau_1$ misreads
a $0$. In the other cases the intended value is $0$, and $\tau_3$ holds $1$ only if $\tau_2$ holds
$0$ and $y_{\tau_3}=0$. For $(1,\cdot)$: if $T_u$ is right, $y_{\tau_3}=0$ needs a misread of
$\tau_3$; and $\tau_2$ holds $0$ only if $\tau_2$ misreads the $0$ at position $1$ or $\tau_1$
holds $1$, which needs $y_{\tau_1}=0$, a misread of $\tau_1$. For $(0,1)$: $\tau_2$ holds $0$ only
if $\tau_2$ misreads or $\tau_1$ holds $1$; the latter needs $L_v$ to hold $0$, so $T_v$ is wrong.
\end{proof}

\noindent The file \texttt{lean/SeqNORTree2.lean} checks Lemma~\ref{lem:exact2} and
Lemma~\ref{lem:dom2} in the abstract form of Section~\ref{sec:checks}: every flag replaced by the
smallest event it is known to contain, positions forgotten, and the leaf, node and cell cases
proved for all values (theorems \texttt{leaf\_domination}, \texttt{domination},
\texttt{top\_domination}, \texttt{exact}, \texttt{cell\_exact}).

\subsection{Probabilities}

With $A=\eps\bar a$ and $B=\eps\bar b$, put $\mathrm{or}(p)=1-(1-p)^2(1-A)$,
$\mathrm{and}(p)=p\,(1-(1-p)(1-B))$, $\zeta=(1-(1-B)(1-A))\,B$, and
\begin{align*}
p^{(0)}_0&=1-(1-\omega_0)^2(1-\zeta)(1-A)^2,\\
p^{(1)}_0&=\bigl(1-(1-\omega_1)(1-B)\bigr)\bigl(1-(1-\omega_1)(1-B)(1-A)\bigr),
\end{align*}
then $p^{(b)}_h$ for $1\le h\le d$ by the alternation of Section~\ref{sec:tree}, $r_b=p^{(b)}_d$, and
\begin{align*}
Q_1&=1-(1-r_0)^2(1-A)^2,\\
Q_0&=\max\Bigl(1-(1-r_1)\bigl(1-B\,(1-(1-B)(1-A))\bigr),\ 1-(1-r_1)(1-A)\Bigr).
\end{align*}
All of these are nondecreasing in $A$, $B$, $\omega_0$, $\omega_1$.

\begin{lemma}[conditional law]\label{lem:law2}
Lemma~\ref{lem:law} holds for $\mathrm{Tree}'(C)$ with these $Q_0,Q_1$.
\end{lemma}

\begin{proof}
As for Lemma~\ref{lem:law}. The two copies of a leaf use independent coins, so given
$\mathcal G_{t-1}$ each is bad with probability $\beta_w$, independently of each other and of the
flags; this gives $p^{(b)}_0$ for the two leaf events. In the cell events the flags of
$\tau_1,\tau_2,\tau_3$ are independent of each other and of the events of the blocks, which gives
$Q_1$ for $(0,0)$, the first term of $Q_0$ for $(1,0)$ and $(1,1)$, and the second for $(0,1)$.
\end{proof}

\begin{theorem}[embedding]\label{thm:embed2}
Let $\bar a,\bar b$ be functions of $\eps$ such that $A(\eps)=\eps\bar a(\eps)$ and
$B(\eps)=\eps\bar b(\eps)$ are nondecreasing on $[0,\eps_0]$, and suppose that at $\eps_0$ the
numbers of this section satisfy $Q_0<\omega_0$ and $Q_1<\omega_1$. Then Theorem~\ref{thm:embed}
holds for $\mathrm{Tree}'(C)$, with the flags defined by $\bar a(\eps),\bar b(\eps)$ and with
$\Delta'$ in place of $\Delta$.
\end{theorem}

\begin{proof}
The proof of Theorem~\ref{thm:embed}, with Lemmas~\ref{lem:dom2} and~\ref{lem:law2}. The proviso
replaces the monotonicity in $\eps$ used there.
\end{proof}

\subsection{Density under heavy zero padding}

Recall $\varpi(\eps)$, the root in $(0,1/2)$ of $\varpi=\eps(1-\varpi)^2$, and
$\psi_g=O_g/(g+1)$. For $\mu\le1$ put
\[
\Lambda(\mu)=\frac{\eps(1-\mu)}{1+\eps(1-\mu)} .
\]
$\Lambda$ is decreasing, $|\Lambda'|\le\eps$ on $(-\infty,1]$, and $\Lambda(\varpi)=\varpi$, since
$\varpi(1+\eps(1-\varpi))=\eps(1-\varpi)$ is the defining equation.

\begin{lemma}[two-sided density]\label{lem:dens2}
Let $0<\eps\le1/4$, $\rho\ge1$, $\tau\in(0,1/8]$, $K\ge0$, $c\ge1$, and let $j$ be a base of length
$n\ge2$ whose gates are zero gates and \emph{cell gates}, such that positions $1,\dots,\mathrm{pre}$
are zero gates and every interval $(a,b]$ of positions contains at most $(b-a)/(1+\rho)+c$ cell
gates. Let $W=c+3+(K+2)\ln(n+1)$ and suppose $\mathrm{pre}\ge(5/\tau^2)^7\,W$. Then with
probability at least $1-(n+1)^{-K}$, for every $g\in[\mathrm{pre},n]$,
\[
\bigl|\psi_g-\varpi(\eps)\bigr|\ \le\ \frac{1/(1+\rho)+\tau}{1-\eps}+\frac{\eps^5}2 .
\]
\end{lemma}

\begin{proof}
This follows the proof of Lemma~\ref{kp:twosided} and of Lemma~\ref{lem:fddens}, with exact conditional
means for the zero gates.

\emph{Conditional means.} Let $a_i=\PP[x_i=1\mid\cF_{i-1}]$ and $z_{i-1}=1-x_{i-1}$. A zero gate
$i\ge2$ reads $x_0=1$ unless it is resampled, and a resampled pointer finds a zero with probability
$1-\psi_{i-1}$; so $a_i=\eps\,z_{i-1}(1-\psi_{i-1})$. For gate $1$ both sides are $0$. For a cell
gate, $0\le a_i\le z_{i-1}$.

\emph{The good event.} Let $\ell n=\ln(n+1)$, $\lambda(m)=\sqrt{(K+2)\,m\,\ell n}$, and let $G^*$
be the event that $\bigl|\sum_{i=a+1}^b(x_i-a_i)\bigr|\le\lambda(b-a)$ for all $0\le a<b\le n$. The
increments lie in intervals of length $1$ that are $\cF_{i-1}$-measurable, so by the
Azuma--Hoeffding inequality each pair fails with probability at most $2(n+1)^{-2(K+2)}$, and
$\PP[(G^*)^c]\le(n+1)^{2-2K-4}\le(n+1)^{-K}$.

\emph{Intervals.} Work on $G^*$. Let $0\le a<b\le n$, $m=b-a$, $X=O_b-O_a$, let $k$ be the number of
cell gates in $(a,b]$, and $\Delta(m)=1+c+m/(1+\rho)+\lambda(m)\ge1+k+\lambda(m)$. Positions
$a,\dots,b-1$ hold between $m-X-1$ and $m-X+1$ zeros, so the sum of $z_{i-1}$ over the zero gates
$i\in(a,b]$ lies between $m-X-1-k$ and $m-X+1$. If $\psi_{i-1}\le\mu\le1$ for all $i\in(a,b]$, then
\[
X\ \ge\ \eps(1-\mu)(m-X-1-k)-\lambda(m),\qquad\text{hence}\qquad X\ge\Lambda(\mu)\,m-\Delta(m).
\]
If $\psi_{i-1}\ge\ell$ for all $i\in(a,b]$, with $\ell\le1$, then
$X\le\eps(1-\ell)(m-X+1)+k+\lambda(m)$, hence $X\le\Lambda(\ell)\,m+\Delta(m)$.

\emph{Prefixes.} For an integer $Q\ge0$ let
$Q'=\lceil2(Q+2+c)/\tau\rceil+\lceil4(K+2)\ell n/\tau^2\rceil$. For $g\ge Q'$ we have
$(Q+2+c)/(g+1)\le\tau/2$ and $\lambda(g)/(g+1)\le\sqrt{(K+2)\ell n/g}\le\tau/2$, so
$(Q+1+\Delta(g))/(g+1)\le\eta:=1/(1+\rho)+\tau$. If $\psi_p\le\mu$ for all $p\in[Q,n]$, the
interval bound on $(Q,g]$ gives $O_g\ge\Lambda(\mu)(g-Q)-\Delta(g)$, so $\psi_g\ge\Lambda(\mu)-\eta$
for $g\ge Q'$, as $\Lambda\le1$. If $\psi_p\ge\ell$ for all $p\in[Q,n]$, then, as $O_Q\le Q+1$,
$\psi_g\le\Lambda(\ell)+\eta$ for $g\ge Q'$.

\emph{Stages.} Let $Q_0=\lceil1/(2\tau)\rceil$ and $Q_{k+1}=Q_k'$. Ones are never adjacent, so
$\psi_g\le\frac12+\frac1{2(g+1)}\le u_0:=\frac12+\tau$ for $g\ge Q_0$. Define $\ell_1=\Lambda(u_0)-\eta$,
$u_2=\Lambda(\ell_1)+\eta$, $\ell_3=\Lambda(u_2)-\eta$, and so on up to $u_6$. By the previous
paragraph, induction gives $\psi\ge\ell_k$ on $[Q_k,n]$ for odd $k$ and $\psi\le u_k$ on $[Q_k,n]$
for even $k$. Compare with $u^*_0=\frac12$, $\ell^*_1=\Lambda(u^*_0)$, $u^*_2=\Lambda(\ell^*_1)$, and so on.
Since $|\Lambda'|\le\eps$ (all arguments are at most $1$), the differences satisfy $e_{k+1}\le\eps e_k+\eta$
with $e_0=\tau\le\eta$, so $e_k\le\eta/(1-\eps)$; and the starred sequence approaches the fixed point
$\varpi$ by a factor $\eps$ per step, so $|u^*_k-\varpi|,|\ell^*_k-\varpi|\le\eps^k/2$. Hence, on
$[Q_6,n]$, $\psi\ge\ell_5\ge\varpi-\eps^5/2-\eta/(1-\eps)$ and
$\psi\le u_6\le\varpi+\eps^6/2+\eta/(1-\eps)$.

\emph{Size of $Q_6$.} By induction, $Q_k\le(5/\tau^2)^{k+1}W$: this holds for $k=0$, and
$Q_{k+1}\le(2/\tau)Q_k+(4/\tau^2)(c+2+(K+2)\ell n)+2\le(2/\tau)Q_k+(5/\tau^2)W$, while
$2\tau/5+(\tau^2/5)^{k+1}\le1$. So $Q_6\le(5/\tau^2)^7W\le\mathrm{pre}$.
\end{proof}

\subsection{Parameters and hardness}

\begin{definition}[$\mathrm{Tree}'_a$]\label{def:treeprimea}
For an integer $a\ge1$ and a circuit with $s\ge2$ gates, $\mathrm{Tree}'_a(C)$ is the base
$\mathrm{Tree}'(C)$ with $d=51$ (so $c=7\cdot2^{52}+1\le3.2\cdot10^{16}$), $\rho=999$,
\begin{align*}
m&=\text{the least power of two with }m\ge375(a+1)\log_2s,\\
Z&=10^{68}(a+2)^2sm,\qquad b=\lceil\log_2Z\rceil,\qquad
\mathrm{pre}=\bigl\lceil1.3\cdot10^{51}\bigl(c+3+(a+3)\,b\ln2\bigr)\bigr\rceil .
\end{align*}
\end{definition}

\begin{lemma}[parameters]\label{lem:param3}
For $\mathrm{Tree}'_a(C)$: (i) $n+1\le Z$; (ii) with $\tau=1/2000$ and $K=a+1$, the hypothesis
$\mathrm{pre}\ge(5/\tau^2)^7W$ of Lemma~\ref{lem:dens2} holds; (iii) $m<750(a+1)\log_2s$, and
$n+1\le s^2$ whenever $s\ge s^\bullet_a:=10^{76}(a+2)^4$.
\end{lemma}

\begin{proof}
(i) $b\ln2\le\ln Z+1\le158+2\ln(a+2)+\ln(sm)$. With $\ln x\le x$, $a+3\le\frac43(a+2)$, $a+2\ge3$
and $sm\ge1500$, this gives $(a+3)\,b\ln2\le(a+2)^2sm$, so
$\mathrm{pre}\le1.3\cdot10^{51}(3.2\cdot10^{16}+1)(a+2)^2sm+1\le4.2\cdot10^{67}(a+2)^2sm$, and
$n+1=\mathrm{pre}+1+1000\,c\,sm\le10^{68}(a+2)^2sm$.
(ii) $(5/\tau^2)^7=(2\cdot10^7)^7\le1.3\cdot10^{51}$, and $b\ln2\ge\ln Z\ge\ln(n+1)$ by (i), so
$\mathrm{pre}\ge1.3\cdot10^{51}(c+3+(K+2)\ln(n+1))$.
(iii) As in Lemma~\ref{lem:param}: $n+1\le Z<7.5\cdot10^{70}(a+2)^3s\log_2s$, and for
$s\ge s^\bullet_a$ we have $\log_2s\le253(a+2)$ at $s=s^\bullet_a$ and $s/\log_2s$ increasing, so
$s/\log_2s\ge3.9\cdot10^{73}(a+2)^3$.
\end{proof}

\begin{corollary}[noise $1/4$]\label{cor:twentythree}
Let $a\ge1$, $\eps\le1/4$, and let $C$ be a NOR circuit with $s\ge2$ gates. For the trace of
$S_\eps(\mathrm{Tree}'_a(C))$, with $\omega_0=4/25$ and $\omega_1=91/500$,
\[
\PP\bigl[\exists\,t\le s:\ \beta_t>\omega_{\val(g_t)}\bigr]\le s^{-a}+(n+1)^{-(a+1)} .
\]
In particular, with at least the complementary probability, at least $409/500$ of the positions
$o(s,1),\dots,o(s,m)$ hold $\val(C)$.
\end{corollary}

\begin{proof}
Take $\bar a=\varpi(\eps)+1/400$ and $\bar b=1-\varpi(\eps)+1/400$. Lemma~\ref{lem:dens2} with
$\rho=999$, $\tau=1/2000$ and $K=a+1$ applies by Lemma~\ref{lem:param3}(ii), and its error is at
most $0.0015/0.75+4^{-5}/2<1/400$. Every cell gate $g$ has $g-1\ge\mathrm{pre}$, so
$\PP[(\Delta')^c]\le(n+1)^{-(a+1)}$. The functions $A(\eps)=\eps(\varpi(\eps)+1/400)$ and
$B(\eps)=\eps(1-\varpi(\eps)+1/400)$ are nondecreasing on $[0,1/4]$: the first is a product of
nondecreasing positive functions, and the derivative of the second is at least
$1-\varpi-\eps\varpi'\ge1-0.18-0.25>0$, because $\varpi'=(1-\varpi)^2/(1+2\eps(1-\varpi))\le1$.
At $\eps_0=1/4$, $\varpi\in[0.17157287,0.17157288]$, so $A\le0.043519$ and $B\le0.207732$, and
the recursion with $d=51$, evaluated at these two rounded values, gives $Q_0\le0.13848$ and
$Q_1\le0.15904$ (exact arithmetic with upward rounding; Section~\ref{sec:checks}). Then
$D(0.16\,\|\,0.13848)\ge0.00185$ and $D(0.182\,\|\,0.15904)\ge0.00185$, so $\varkappa\ge0.00185$, and Theorem~\ref{thm:embed2} bounds the
rest by $s\cdot s^{-0.00185\cdot375\,(a+1)/\ln2}\le s^{-a}$.
\end{proof}

\begin{theorem}[hardness up to noise $1/4$]\label{thm:hard3}
Theorem~\ref{thm:hard} holds with $\eps\le1/4$ in place of $\eps\le1/64$: if for some
$\eps\le1/4$ a randomized NC family outputs the trace with probability at least $n^{-K}$ on the
$\eps$-smoothing of every base of every length $n\ge n_0$, then $\mathrm P\subseteq\mathrm{BPNC}$.
\end{theorem}

\begin{proof}
Repeat the proof of Theorem~\ref{thm:hard2} with these changes: $\eps^*=1/4$;
$\mathrm{Tree}'_a$ in place of $\mathrm{Tree}_a$; $s\ge\max(s^\bullet_a,n_0,3)$;
Corollary~\ref{cor:twentythree} in place of Corollary~\ref{cor:threetwentieths}; and
Lemma~\ref{lem:param3}(iii) in place of Lemma~\ref{lem:param2}(iii). The coins of a cell are
$2^{d+2}$ elements of $\{1,\dots,m\}$, exactly uniform as $m$ is a power of two; the vote of a trial is the majority bit of the output bundle,
which is $\val(C)$ whenever at most $91/500$ of its cells are bad.
\end{proof}

\begin{corollary}[the gap, with the new lower end]\label{cor:gap3}
Corollary~\ref{cor:gap} holds with $1/4$ in place of $1/64$. In particular
$\inf E_K\in\{0\}\cup[1/4,3/4]$ for every $K$, and if $\mathrm P\not\subseteq\mathrm{BPNC}$,
the critical noise level satisfies $1/4\le\eps_c\le3/4$.
\end{corollary}

\begin{proof}
As for Corollary~\ref{cor:gap2}, with Theorem~\ref{thm:hard3}.
\end{proof}

\begin{remark}[how far this goes]\label{rem:limits2}
\begin{enumerate}[label=(\roman*),leftmargin=2.4em]
\item The gain over Section~\ref{sec:tree} comes from two sources: the density bounds (misread
probabilities $\eps(\varpi\pm1/400)$ and $\eps(1-\varpi\pm1/400)$ instead of $0.17\,\eps$ and
$\eps$) and the three refinements of the cells. At depth $25$ with a margin of $20\%$ in the
tolerances, the first raises the threshold from $0.173$ by about $0.028$ and the second by a
further $0.041$, to $0.241$.
\item Noise $1/4$ is close to the limit of this construction. With a margin of $20\%$ in the
tolerances and the density error $1/400$ of the proof, the threshold is $0.237$ at depth $21$,
$0.241$ at depth $25$ and $0.245$ at depth $31$, and with no margin it is $0.256$ at depth $41$. At
noise $23/100$, depth $25$ and tolerances $7/40$, $9/40$ already give $\varkappa\ge0.0148$
(script \texttt{check\_tree2.py}). With the exact misread probabilities and no margin the recursion
stops at $0.255$ at depth $31$ and $0.259$ at depth $61$. In that setting, three reading gates
(like $\tau_1,\tau_2,\tau_3$) at the nodes that should hold $0$ would raise it to $0.29$ at depth
$61$, but which nodes these are depends on the values of the circuit, which the compiler does not
know; three gates at every node give only $0.246$, less than one gate. All these thresholds come
from \texttt{tree2\_explore.py}. A higher lower end needs a different restoring organ;
Section~\ref{sec:latch} gives one.
\item The constants are astronomical: the preamble has about $10^{68}$ gates, and
Theorem~\ref{thm:hard3} concerns circuits with more than $10^{76}$ gates. This comes from the
restarts in the proof of Lemma~\ref{lem:dens2} and does not affect the statements, which are
asymptotic, like those of Section~\ref{sec:gapbase}.
\item In simulation (script \texttt{embed2.c}) cells of depth $7$ with padding ratio $99$ kept a
correct output majority at noise $0.23$, $0.25$, $0.28$ and $0.32$, with the density of ones
within $0.003$ of $\varpi(\eps)$. At $0.25$ (bundles of $31$) and at $0.32$ (bundles of $15$) some
bundles were more than half bad, as expected at this depth (Section~\ref{sec:checks}; the run at
$0.32$ is the last row of the table there).
\end{enumerate}
\end{remark}

\section{A latch embedding: hardness at every noise level below \texorpdfstring{$4/9$}{4/9}}\label{sec:latch}

Sections~\ref{sec:tree} and~\ref{sec:tree2} restore values with a tree of NOR gates, and
Remark~\ref{rem:limits2} explains why trees of that kind stop near noise $0.26$. This section uses
a different restoring organ and proves Theorem~\ref{thm:hard4}: the conclusion of
Theorem~\ref{thm:hard} holds for every $\eps<4/9$. Of the earlier sections it uses the composition
of smoothings (Lemma~\ref{gb:comp}), the coupling of Lemma~\ref{lem:misread}(a), the two-sided
density bound (Lemmas~\ref{lem:dens2} and~\ref{lem:dens2any}; the proof of the second does not
depend on the sections between) and the reduction in the proof of Theorem~\ref{thm:hard}, with
Lemma~\ref{lem:errorless}.

\emph{Where the gain comes from.}
\begin{itemize}[leftmargin=2em]
\item \emph{Evidence.} A gate that reads a one outputs zero, whatever its predecessor holds. With
heavy zero padding a resampled pointer finds a one with probability about $\varpi(\eps)$ only. So
``the gate read a one'' is good evidence that its listed pointer holds a one: a listed zero is
read as a one with probability about $\eps\varpi$ only. ``The gate read a zero'' is poor evidence,
because a listed one is read as a zero with probability about $\eps(1-\varpi)$.
\item \emph{Latches.} Two consecutive gates form a latch. If they read $y$ and $y'$, the second
one holds $(x\vee y)\wedge\neg y'$, where $x$ is the value before the pair: a one read by the
first gate sets the state, a one read by the second gate resets it, and otherwise the state is
kept. It is kept without noise, because the predecessor input of a gate is never resampled. A
chain of such pairs ends in the state given by the last evidence it saw. So errors do not add up
along a chain; later evidence overrides earlier evidence.
\item \emph{Two rails.} Every value is kept twice, in a bundle of copies that should hold the
value and in a bundle that should hold its negation. A new copy of the first bundle is a chain
whose first gates of the pairs read random copies of the first bundle and whose second gates read
random copies of the second; for the second bundle the roles are exchanged. Then true evidence
arrives at a rate close to the fraction of correct copies times $1-\eps$, and false evidence at
a rate close to $\eps\varpi$ plus the fraction of wrong copies times $1-\eps$.
\end{itemize}
In the trees of Section~\ref{sec:tree} a node has to be right in two roles, and in one of them
the errors of its two children add. Here there is one kind of chain. The analysis is a monotone
coupling with exact read probabilities; there are no dominating events.

The level $4/9$ is exact for this analysis: the map that describes one restoring stage has the
fixed point that is needed exactly below $4/9$, and from $4/9$ on no parameters satisfy the
conditions of the proof (Propositions~\ref{prop:ideal} and~\ref{prop:adm}). It is also the level
below which the relevance-pruned iteration needs at least $cn/\ln n$ rounds on random block
bases over zero gates (Theorem~\ref{thm:blockzero}), and the two conditions are the same
inequality (Remark~\ref{rem:latch}).

\emph{Notation.} In this section $s$ and $r$ also denote the probabilities that a set check and a
reset check read a one, and $\tau=(1-s)(1-r)$; in Proposition~\ref{prop:ideal}, $u$ and $v$ are
coordinates. The number of gates of the circuit is $s$, and the inputs of a gate are $u,v$, only
where a circuit is in sight.

\subsection{Latches}

For bits $x_0,y_1,\dots,y_{2h}$ let $x_i=\NOR(x_{i-1},y_i)$ for $1\le i\le2h$, and put
\[
\mathrm{Lat}_h(x_0;y_1,\dots,y_{2h})=x_{2h}.
\]
For $s,r\in[0,1]$, not both $0$, put
\[
\tau=(1-s)(1-r),\qquad F(s,r)=\frac{s\,(1-r)}{1-\tau}=\frac{s\,(1-r)}{s+r-sr},
\]
and for $h\ge1$
\[
F^-_h(s,r)=F(s,r)\,(1-\tau^h),\qquad F^+_h(s,r)=F(s,r)+(1-F(s,r))\,\tau^h .
\]

\begin{lemma}[latch]\label{lem:latch}
\begin{enumerate}[label=(\alph*),leftmargin=2.4em]
\item $x_{2i}=(x_{2i-2}\vee y_{2i-1})\wedge\neg y_{2i}$ for $1\le i\le h$. Hence $\mathrm{Lat}_h$ is
nondecreasing in $x_0$ and in $y_i$ for odd $i$, and nonincreasing in $y_i$ for even $i$.
\item Let $x_0$ be fixed and let the $y_i$ be independent, with $\PP[y_i=1]=s$ for odd $i$ and
$\PP[y_i=1]=r$ for even $i$. Then $\PP[x_{2h}=1]=F(s,r)(1-\tau^h)+x_0\,\tau^h$. So this probability
is $F^-_h(s,r)$ for $x_0=0$ and $F^+_h(s,r)$ for $x_0=1$.
\item $F$, $F^-_h$ and $F^+_h$ are nondecreasing in $s$ and nonincreasing in $r$, and
$F^-_h\le F\le F^+_h\le F^-_h+\tau^h$. On $(0,1)^2$, $F$ is strictly increasing in $s$ and
strictly decreasing in $r$.
\end{enumerate}
\end{lemma}

\begin{proof}
(a) $x_{2i-1}=\neg x_{2i-2}\wedge\neg y_{2i-1}$, so
$x_{2i}=\neg x_{2i-1}\wedge\neg y_{2i}=(x_{2i-2}\vee y_{2i-1})\wedge\neg y_{2i}$. Induction on $i$
gives the monotonicity.
(b) Let $p_i=\PP[x_{2i}=1]$. By (a), $p_i=(1-r)\bigl(s+(1-s)p_{i-1}\bigr)=(1-\tau)F(s,r)+\tau p_{i-1}$,
so $p_i-F=\tau\,(p_{i-1}-F)$ and $p_h=F+\tau^h(x_0-F)$.
(c) Realise $y_i$ as the indicator of $\{\xi_i<s\}$ or $\{\xi_i<r\}$ with independent uniform
variables $\xi_i$. Raising $s$ raises the odd $y_i$ and raising $r$ raises the even ones, so by
(a) and (b) $F^-_h$ and $F^+_h$ are monotone as stated, and so is their common limit $F$ when
$\tau<1$. The inequalities follow from the formulas. With $D=s+r-sr$,
$\partial F/\partial s=r(1-r)/D^2$ and $\partial F/\partial r=-s/D^2$, because $D+(1-s)(1-r)=1$.
\end{proof}

\subsection{Reads}

We are in the setting of Section~\ref{sec:tree}: the base $j$ may depend on coins,
$\jt=S_\eps(j)$ is built from $(\pi_i,\rho_i)$, the $U_i$ are independent uniform variables,
$\cF_g$ is generated by all coins and by $(\pi_i,\rho_i)_{i\le g}$, and
$\mathcal H_g=\cF_g\vee\sigma(U_i:i\le g)$. Gate $g$ reads the value $y_g=x_{\jt_g}$, and
$x_g=\NOR(x_{g-1},y_g)$. Given $\mathcal H_{g-1}$, it reads a one with probability
\[
\lambda_g=(1-\eps)\,x_{j_g}+\eps\,\psi_{g-1},\qquad\psi_{g-1}=O_{g-1}/g ,
\]
because $\rho_g$ is uniform on $\{0,\dots,g-1\}$. Put $\Xi_g=\lambda_gU_g$ if $y_g=1$ and
$\Xi_g=\lambda_g+(1-\lambda_g)U_g$ if $y_g=0$. For a constant $\eta\in(0,\varpi(\eps))$ let
\[
y^-_g=\mathbf 1\bigl[\Xi_g<(1-\eps)x_{j_g}+\eps(\varpi(\eps)-\eta)\bigr],\qquad
y^+_g=\mathbf 1\bigl[\Xi_g<(1-\eps)x_{j_g}+\eps(\varpi(\eps)+\eta)\bigr].
\]

\begin{lemma}[reads]\label{lem:reads}
\begin{enumerate}[label=(\alph*),leftmargin=2.4em]
\item The variables $\Xi_g$, $g\ge1$, are independent and uniform on $[0,1]$, and for every $g_0$
the family $(\Xi_g)_{g\ge g_0}$ is independent of $\mathcal H_{g_0-1}$; in particular it is
independent of all coins.
\item Almost surely, $y_g=1$ if and only if $\Xi_g<\lambda_g$. Hence, if
$|\psi_{g-1}-\varpi(\eps)|\le\eta$, then $y^-_g\le y_g\le y^+_g$.
\end{enumerate}
\end{lemma}

\begin{proof}
As for Lemma~\ref{lem:misread}(a),(b), with ``reads a one'' in place of ``misreads'': given
$\mathcal H_{g-1}$, the event $\{y_g=1\}$ has probability $\lambda_g$, so $\Xi_g$ is uniform and
independent of $\mathcal H_{g-1}$, and it is $\mathcal H_g$-measurable. If $y_g=1$ then
$\Xi_g<\lambda_g$ unless $U_g=1$, and if $y_g=0$ then $\Xi_g\ge\lambda_g$.
\end{proof}

\subsection{The compiler}

Position $0$ holds $1$ and position $1$ holds $0$ in every instance. A \emph{source} is a
probability distribution on positions; a gate \emph{reads from} a source if its listed pointer is
drawn from it with fresh coins. A \emph{rail} is a list of $m$ positions, its \emph{copies}; as a
source it is uniform on the list. The rail $\mathbf 1$ has all copies equal to position $0$, and
the rail $\mathbf 0$ has all copies equal to position $1$. For sources $S,S'$ and
$\theta\in(0,1]$, the source $\frac12S+\frac12S'$ draws from $S$ or from $S'$ with probability
$\frac12$ each, and $\theta S$ draws from $S$ with probability $\theta$ and is position $1$
otherwise.

\begin{definition}[$\mathrm{Latch}(C)$]\label{def:latch}
Let $C$ be a NOR circuit with gates $g_t=\NOR(u_t,v_t)$, $t\le s$. Fix integers $h,J,\rho\ge1$, a
bundle size $m$ that is a power of two, a dyadic rational $\theta\in(0,1]$ and a preamble length
$\mathrm{pre}$.
\begin{itemize}[leftmargin=2em]
\item A \emph{chain} with set source $S$ and reset source $S'$ is a sequence of $2h$ consecutive
gates; the first, third, \dots\ gate read from $S$ and the second, fourth, \dots\ from $S'$, all with
independent coins. The last gate is the \emph{copy} of the chain. A \emph{block} is $2h\rho$ zero
gates followed by a chain. The rail \emph{built from} $(S,S')$ is a sequence of $m$ consecutive
blocks whose chains have set source $S$ and reset source $S'$; its copies are the copies of the
chains.
\item The \emph{wires} are the two constants and the gates of $C$. Every wire $w$ has a positive
rail $P_k(w)$ and a negative rail $N_k(w)$ for every \emph{stage} $k=0,\dots,J$.
\item \emph{Constants.} $P_0(c_1)=N_0(c_0)=\mathbf 1$ and $N_0(c_1)=P_0(c_0)=\mathbf 0$.
\item \emph{Gate stage.} For $g=\NOR(u,v)$ let $\Sigma^+=\frac12P_J(u)+\frac12P_J(v)$ and
$\Sigma^-=\frac12N_J(u)+\frac12N_J(v)$. Then $P_0(g)$ is built from $(\theta\Sigma^-,\Sigma^+)$ and
$N_0(g)$ is built from $(\Sigma^+,\theta\Sigma^-)$.
\item \emph{Restoring stages.} For every wire $w$ and $1\le k\le J$, $P_k(w)$ is built from
$(P_{k-1}(w),N_{k-1}(w))$ and $N_k(w)$ from $(N_{k-1}(w),P_{k-1}(w))$.
\end{itemize}
The base $\mathrm{Latch}(C)$ consists of zero gates at positions $1,\dots,\mathrm{pre}$, followed
by the rails of stages $1,\dots,J$ of $c_0$ and then of $c_1$, and then, for $t=1,\dots,s$ in
this order, by the rails of stages $0,\dots,J$ of $g_t$; within a wire the order is
$P_0,N_0,P_1,N_1,\dots$. Its length is
\[
n=\mathrm{pre}+2h(1+\rho)\,m\,\bigl(2(J+1)s+4J\bigr).
\]
The gates of the chains are the \emph{chain gates}.
\end{definition}

All positions are explicit arithmetic expressions, a chain gate uses a bounded number of fair
coins plus $\log_2m$ coins for the index of a copy, and every chain reads only positions $0$, $1$
and copies of rails that lie before it. So the base is computable in logspace from $C$ and the
coins. Evidence for $\val(g)=1$ comes from the negative rails of both inputs, and the factor
$\theta$ makes one such piece of evidence count less than evidence for $\val(g)=0$, which a
single input can give.

Put $\val(c_1)=1$ and $\val(c_0)=0$. At stage $k$, the \emph{high rail} of a wire $w$ is $P_k(w)$
if $\val(w)=1$ and $N_k(w)$ if $\val(w)=0$; the other one is its \emph{low rail}. For a rail
$\mathcal R$, in the trace of the smoothed instance, $\mathrm{dens}(\mathcal R)$ is the fraction of
its copies that hold $1$; so $\mathrm{dens}(\mathbf 1)=1$ and $\mathrm{dens}(\mathbf 0)=0$.

\subsection{Admissible parameters and the embedding theorem}

Fix $\eps\in(0,1)$ and write $\varpi=\varpi(\eps)$. For $\eta\in(0,\varpi)$, $\gamma>0$, $h\ge1$ and
$d,d'\in[0,1]$ put
\begin{align*}
\ell^-(d)&=\eps(\varpi-\eta)+(1-\eps)d, & \mathrm{lo}(d,d')&=F^-_h\bigl(\ell^-(d),\ell^+(d')\bigr)-\gamma,\\
\ell^+(d)&=\eps(\varpi+\eta)+(1-\eps)d, & \mathrm{up}(d,d')&=F^+_h\bigl(\ell^+(d),\ell^-(d')\bigr)+\gamma .
\end{align*}
By Lemma~\ref{lem:latch}(c), $\mathrm{lo}$ and $\mathrm{up}$ are nondecreasing in $d$ and
nonincreasing in $d'$. The number $\mathrm{lo}(d,d')$ will bound from below the density of a rail
whose set source holds ones with density at least $d$ and whose reset source with density at most
$d'$; $\mathrm{up}(d,d')$ bounds it from above when the set density is at most $d$ and the reset
density at least $d'$.

\begin{definition}[admissible parameters]\label{def:adm}
Parameters $(\theta,h,J,\eta,\gamma)$ are \emph{admissible for} $\eps$ if there are numbers
$0\le a_k\le a'_k\le1$ and $0\le b'_k\le b_k\le1$, $k=0,\dots,J$, with $a'_0=1$ and $b'_0=0$, such
that:
\begin{itemize}[leftmargin=3em]
\item[(R)] for $1\le k\le J$,
\[
a_k\le\mathrm{lo}(a_{k-1},b_{k-1}),\quad b_k\ge\mathrm{up}(b_{k-1},a_{k-1}),\quad
a'_k\ge\mathrm{up}(a'_{k-1},b'_{k-1}),\quad b'_k\le\mathrm{lo}(b'_{k-1},a'_{k-1});
\]
\item[(G)] with $\mu^-=(a_J+b'_J)/2$ and $\mu^+=(a'_J+b_J)/2$, for each of the three pairs
\[
(d,d')=(\theta a_J,\,b_J),\qquad(\mu^-,\,\theta\mu^+),\qquad(a_J,\,\theta b_J)
\]
we have $a_0\le\mathrm{lo}(d,d')$ and $b_0\ge\mathrm{up}(d',d)$;
\item[(V)] $a_J>b_J$.
\end{itemize}
\end{definition}

\noindent The numbers $a_k,a'_k$ will bound the density of a high rail at stage $k$ and
$b'_k,b_k$ that of a low rail. Let $E_\eta$ be the event that $|\psi_g-\varpi(\eps)|\le\eta$ for
every $g$ with $\mathrm{pre}\le g\le n$.

\begin{theorem}[latch embedding]\label{thm:latch}
Let $\eps\in(0,1)$, let $(\theta,h,J,\eta,\gamma)$ be admissible for $\eps$ with the numbers
$a_k,a'_k,b'_k,b_k$, and let $C$ be a NOR circuit with $s$ gates. In the trace of
$S_\eps(\mathrm{Latch}(C))$ call a rail of stage $k$ \emph{good} if its density lies in
$[a_k,a'_k]$ when it is a high rail and in $[b'_k,b_k]$ when it is a low rail. Then, over the
coins and the smoothing, for all $m$, $\rho$ and $\mathrm{pre}$,
\[
\PP\bigl[E_\eta\cap\{\text{some rail is not good}\}\bigr]\ \le\ 4(J+1)(s+2)\,e^{-2\gamma^2m}.
\]
\end{theorem}

\begin{proof}
The rails of stage $0$ of the constants have density $1$ (high) and $0$ (low) and are good,
because $a_0\le a'_0=1$ and $b_0\ge b'_0=0$. Let $\mathcal R_1,\mathcal R_2,\dots$ be the other
rails in the order of the layout; there are at most $2(J+1)(s+2)$ of them. Let $\mathcal G_i$ be
generated by the coins of the chains of $\mathcal R_1,\dots,\mathcal R_{i-1}$ and by
$(\pi_g,\rho_g,U_g)$ for the positions $g$ below the first position of $\mathcal R_i$. The trace
below $\mathcal R_i$ is $\mathcal G_i$-measurable. Let $H_i\in\mathcal G_i$ be the event that
$\mathcal R_1,\dots,\mathcal R_{i-1}$ are good.

\emph{Sandwich.} Fix $\mathcal R=\mathcal R_i$, built from $(S,S')$. For a chain $c$ of
$\mathcal R$ with gates $g_1<\dots<g_{2h}$ put
\[
X^-(c)=\mathrm{Lat}_h\bigl(0;\,y^-_{g_1},y^+_{g_2},y^-_{g_3},\dots,y^+_{g_{2h}}\bigr),\qquad
X^+(c)=\mathrm{Lat}_h\bigl(1;\,y^+_{g_1},y^-_{g_2},y^+_{g_3},\dots,y^-_{g_{2h}}\bigr).
\]
Every chain gate $g$ has $g-1\ge\mathrm{pre}$. So on $E_\eta$, by Lemma~\ref{lem:reads}(b) and
Lemma~\ref{lem:latch}(a), the copy of $c$ holds a value between $X^-(c)$ and $X^+(c)$, whatever the
gate before the chain holds.

\emph{Conditional law.} The listed pointers of the gates of $\mathcal R$ are positions below
$\mathcal R$ chosen by the coins of $\mathcal R$. So $X^\pm(c)$ is a function of the variables
$\Xi_g$ of the gates of $c$, of the coins of $c$, and of the trace below $\mathcal R$. By
Lemma~\ref{lem:reads}(a) the $\Xi_g$ of the gates of $\mathcal R$ are independent, uniform, and
independent of $\mathcal G_i$ and of all coins; and the coins of $\mathcal R$ are independent of
$\mathcal G_i$. Let $d_S$ be the probability that a position drawn from $S$ holds a one in the
trace, and $d_{S'}$ the same for $S'$; they are $\mathcal G_i$-measurable. Given $\mathcal G_i$,
the indicators $y^\pm_g$ of the gates of $\mathcal R$ are therefore independent across gates, with
\[
\PP[y^\pm_g=1\mid\mathcal G_i]=(1-\eps)d+\eps(\varpi\pm\eta)=\ell^\pm(d),\qquad
d=d_S\ \text{or}\ d_{S'},
\]
according to the source of $g$; the thresholds in the definition of $y^\pm_g$ lie in $[0,1]$
because $0<\eta<\varpi$ and $\varpi+\eta<1$. Hence, given $\mathcal G_i$, the pairs $(X^-(c),X^+(c))$ of the
$m$ chains are independent, and by Lemma~\ref{lem:latch}(b)
\[
\PP[X^-(c)=1\mid\mathcal G_i]=F^-_h\bigl(\ell^-(d_S),\ell^+(d_{S'})\bigr),\qquad
\PP[X^+(c)=1\mid\mathcal G_i]=F^+_h\bigl(\ell^+(d_S),\ell^-(d_{S'})\bigr).
\]

\emph{The bounds.} We claim that on $H_i$: if $\mathcal R$ is a high rail of stage $k$, then
$\PP[X^-(c)=1\mid\mathcal G_i]\ge a_k+\gamma$ and, for $k\ge1$,
$\PP[X^+(c)=1\mid\mathcal G_i]\le a'_k-\gamma$; if it is a low rail, then
$\PP[X^+(c)=1\mid\mathcal G_i]\le b_k-\gamma$ and, for $k\ge1$,
$\PP[X^-(c)=1\mid\mathcal G_i]\ge b'_k+\gamma$. For $k=0$ the other two bounds of a good rail are
trivial, because $a'_0=1$ and $b'_0=0$.

For $k\ge1$ the high rail of $w$ is built from the high and the low rail of stage $k-1$, in this
order, and the low rail from the low and the high rail. On $H_i$ these are good, so
$d_S\in[a_{k-1},a'_{k-1}]$ and $d_{S'}\in[b'_{k-1},b_{k-1}]$ for the high rail, and the other way
round for the low rail. The claim follows from (R) and the monotonicity of $\mathrm{lo}$ and
$\mathrm{up}$.

For $k=0$ and a gate $g=\NOR(u,v)$, the rails of stage $J$ of $u$ and $v$ are good on $H_i$, so
the high ones have density in $[a_J,a'_J]$ and the low ones in $[b'_J,b_J]$; this includes the
case that $u$ or $v$ is a constant, and the case $u=v$. Write $d^+,d^-$ for the densities of
$\Sigma^+,\Sigma^-$; the source $\theta\Sigma^-$ has density $\theta d^-$, because position $1$
holds $0$. There are three cases.
\begin{itemize}[leftmargin=2em]
\item $\val(u)=\val(v)=0$, so $\val(g)=1$ and the high rail is $P_0(g)$. The positive rails of the
inputs are low and the negative ones high: $d^+\le b_J$ and $d^-\ge a_J$. The high rail has set
density $\theta d^-\ge\theta a_J$ and reset density $d^+\le b_J$; the low rail has set density at
most $b_J$ and reset density at least $\theta a_J$. This is the pair $(\theta a_J,b_J)$ of (G).
\item $\val(u)\ne\val(v)$, so $\val(g)=0$ and the high rail is $N_0(g)$. Each of $\Sigma^+$,
$\Sigma^-$ is the mean of a high and a low rail, so $d^\pm\in[\mu^-,\mu^+]$. The high rail has set
density $d^+\ge\mu^-$ and reset density $\theta d^-\le\theta\mu^+$, and the low rail the reverse.
This is the pair $(\mu^-,\theta\mu^+)$.
\item $\val(u)=\val(v)=1$, so $\val(g)=0$ and the high rail is $N_0(g)$, with set density
$d^+\ge a_J$ and reset density $\theta d^-\le\theta b_J$, and the low rail the reverse. This is the
pair $(a_J,\theta b_J)$.
\end{itemize}
In each case (G) and the monotonicity give $\PP[X^-(c)=1\mid\mathcal G_i]\ge a_0+\gamma$ for the
high rail and $\PP[X^+(c)=1\mid\mathcal G_i]\le b_0-\gamma$ for the low rail.

\emph{Concentration.} Let $\mathcal R$ be a high rail of stage $k$. Given $\mathcal G_i$, on $H_i$,
the number of chains with $X^-(c)=1$ is a sum of $m$ independent indicators with mean at least
$m(a_k+\gamma)$, so by Hoeffding's inequality~\cite{hoeffding} it is less than $ma_k$ with
conditional probability at most $e^{-2\gamma^2m}$. On $E_\eta$ the density of $\mathcal R$ is at
least that number divided by $m$. So
$\PP[E_\eta\cap H_i\cap\{\mathrm{dens}(\mathcal R)<a_k\}]\le e^{-2\gamma^2m}$. The same argument
with $X^+$ bounds the probability of $\mathrm{dens}(\mathcal R)>a'_k$ for $k\ge1$, and the two
bounds for a low rail. If some rail is not good, let $\mathcal R_i$ be the first one; then $H_i$
holds and one of the two bounds for $\mathcal R_i$ fails. A union bound over $i$ and over the two
bounds gives the theorem.
\end{proof}

\subsection{The ideal map and the level \texorpdfstring{$4/9$}{4/9}}

Admissible parameters exist exactly when a simple map of the unit square has a fixed point off the
diagonal. Let $A=\eps\varpi$ and $B=\eps(1-\varpi)$, so $A+B=\eps$; these are the probabilities
that a listed zero is read as a one and a listed one as a zero when the density of ones is
$\varpi$. Put $\ell(d)=A+(1-\eps)d$ and, for $d,d'\in[0,1]$,
\[
\Phi(d,d')=\Bigl(F\bigl(\ell(d),\ell(d')\bigr),\,F\bigl(\ell(d'),\ell(d)\bigr)\Bigr),\qquad T(x,y)=\Phi(x,y).
\]
A \emph{state} is a point $z=(x,y)\in[0,1]^2$: the densities of a high and of a low rail. $T$
describes one restoring stage with $\eta=\gamma=0$ and $h=\infty$. Write $z\preceq z'$ if
$x\le x'$ and $y\ge y'$, and $z\prec\!\!\prec z'$ if both inequalities are strict.

\begin{proposition}[the ideal map]\label{prop:ideal}
Let $0<\eps<1$.
\begin{enumerate}[label=(\alph*),leftmargin=2.4em]
\item $\Phi$ is continuous, and $\Phi(d_1,d_1')\prec\!\!\prec\Phi(d_2,d_2')$ whenever $d_1\le d_2$,
$d_1'\ge d_2'$ and at least one of the two inequalities is strict. In particular $T$ is monotone:
$z\preceq z'$ implies $Tz\preceq Tz'$.
\item $T$ has exactly one fixed point on the diagonal, $z_0=(\pi_0,\pi_0)$ with
$\pi_0=(1-A-q)/(1-\eps)$, where $q\in(0,1)$ is the positive root of $q^2+(1-B)q-(1-A)=0$. This
$q$ is the number $q(\eps)$ of~\eqref{eq:qeps}.
\item If $\eps\ge4/9$, $T$ has no other fixed point. If $\eps<4/9$, it has exactly two others,
$L=(x_*,y_*)$ and $(y_*,x_*)$, with $y_*<\pi_0<x_*$ and $x_*+y_*>2\pi_0$. Here
$1-\ell(x_*)=u_*$ and $1-\ell(y_*)=v_*$ are the two roots of $z^2-(1+B)z+\eps=0$; with
$w=1-\varpi$,
\[
u_*=\frac{1-\sqrt{4w-3}}{2w},\qquad v_*=\frac{1+\sqrt{4w-3}}{2w}.
\]
Moreover $q^2>\eps$ for $\eps<4/9$.
\item Let $\eps<4/9$, $z_e=\bigl(\pi_0+e/(1-\eps),\,\pi_0-e/(1-\eps)\bigr)$ and
$e_1=\min\bigl((q^2-\eps)/2,\ q-B,\ 1-A-q\bigr)>0$. For $0<e\le e_1$, $z_e$ is a state,
$z_e\prec\!\!\prec Tz_e$, and $T^nz_e\to L$ as $n\to\infty$. Also $T^n(1,0)\to L$.
\end{enumerate}
\end{proposition}

\begin{proof}
(a) $\ell(d)\in[A,1-B]\subset(0,1)$, so this is Lemma~\ref{lem:latch}(c).

\emph{Coordinates.} For a state $(x,y)$ put $u=1-\ell(x)$ and $v=1-\ell(y)$, the probabilities
that a check of each kind reads a zero. Then $F(\ell(x),\ell(y))=(1-u)v/(1-uv)$, and $T$ becomes
\[
u'=1-A-(1-\eps)\frac{v(1-u)}{1-uv},\qquad v'=1-A-(1-\eps)\frac{u(1-v)}{1-uv},
\]
so that
\begin{equation}\label{eq:uv}
v'-u'=\frac{(1-\eps)(v-u)}{1-uv},\qquad u'+v'=2(1-A)-(1-\eps)\,\frac{u+v-2uv}{1-uv}.
\end{equation}
States correspond to $(u,v)\in[B,1-A]^2$, and a state is larger in the order $\preceq$ when $u$
is smaller and $v$ larger.

(b) For $u=v=q$ the fixed point equation is $q=1-A-(1-\eps)q/(1+q)$, which is~\eqref{eq:qeps}
and, since $\eps-A=B$, is equivalent to $Q(q)=0$ with $Q(z)=z^2+(1-B)z-(1-A)$. $Q$ is increasing
on $[0,\infty)$, $Q(B)=\eps-1<0$ and $Q(1-A)=(1-A)(1-\eps)>0$, so there is exactly one root, and
$B<q<1-A$.

(c) By~\eqref{eq:uv}, a fixed point with $u\ne v$ has $uv=\eps$, and then the second equation
reads $u+v=2(1-A)-(u+v)+2\eps$, that is $u+v=1+B$. Conversely such a pair is a fixed point. So
$u,v$ are the roots of $P(z)=z^2-(1+B)z+\eps$. From $\varpi=\eps w^2$ we get $\eps=(1-w)/w^2$,
$B=(1-w)/w$ and $1+B=1/w$, so the discriminant is
\[
(1+B)^2-4\eps=\frac{4w-3}{w^2}.
\]
Since $\eps=\varpi/(1-\varpi)^2$ increases with $\varpi$ and $\varpi(4/9)=1/4$, it is positive
exactly for $\eps<4/9$; for $\eps=4/9$ the double root is $u=v=2/3$, which is $q$. For $\eps<4/9$
the roots are the $u_*<v_*$ of the statement. They belong to states: $P(B)=A>0$,
$P(1-A)=\eps A>0$, and $B<(1+B)/2<1-A$ because $1-\eps-A>0$; hence $B<u_*<v_*<1-A$. For the
position of $q$: $Q((1+B)/2)=\bigl(4A-(1-B)^2\bigr)/4=-\bigl((1+B)^2-4\eps\bigr)/4<0$, so
$q>(1+B)/2$, which is $x_*+y_*>2\pi_0$ and gives $q^2>(1+B)^2/4>\eps$. And, using $Q(q)=0$,
$P(q)=(1-A)-(1-B)q-(1+B)q+\eps=1+B-2q<0$, so $u_*<q<v_*$, which is $y_*<\pi_0<x_*$.

(d) In the coordinates $(u,v)$, $z_e$ is $(q-e,q+e)$, which lies in $[B,1-A]^2$ for $e\le e_1$. Put $D=1-q^2+e^2$. By~\eqref{eq:uv} and the
equation for $q$,
\[
v'-u'=\frac{2(1-\eps)e}{D},\qquad u'+v'=2q-\frac{2(1-\eps)e^2}{(1+q)D} .
\]
(For the second: $2q=2(1-A)-2(1-\eps)q/(1+q)$, and
$2qD-(1+q)\bigl(2q(1-q)+2e^2\bigr)=-2e^2$.) So $u'<q-e$ as soon as $(1-\eps)/D\ge1$, that is
$e^2\le q^2-\eps$; and $v'>q+e$ if and only if $q^2-\eps>e^2+(1-\eps)e/(1+q)$. Both hold for
$0<e\le(q^2-\eps)/2$, because then $e<\frac12$ and $e^2+(1-\eps)e/(1+q)<\frac32e<q^2-\eps$. So
$z_e\prec\!\!\prec Tz_e$. By (a) the sequence $T^nz_e$ is nondecreasing for $\preceq$, hence
convergent, and its limit is a fixed point that is $\succeq z_e$; it has $u<q<v$, so by (b) and (c)
it is $L$. Finally $(1,0)$ is the largest state, so $T^n(1,0)$ is nonincreasing and stays
$\succeq L$; its limit is a fixed point $\succeq L$, which has $u\le u_*<v_*\le v$ and so is $L$.
\end{proof}

\begin{proposition}[admissible parameters]\label{prop:adm}
For every $\eps\in(0,4/9)$ there are admissible parameters $(\theta,h,J,\eta,\gamma)$, with
$\theta$ a dyadic rational. For $\eps\in[4/9,1)$ there are none.
\end{proposition}

\begin{proof}
Use the notation of Proposition~\ref{prop:ideal}. Let first $\eps<4/9$.

\emph{Thinning.} Since $y_*<x_*$ and $x_*+y_*>2\pi_0$, the interval
$\bigl(\pi_0/x_*,\,2\pi_0/(x_*+y_*)\bigr)$ is nonempty and lies in $(0,1)$. Let $\theta$ be a
dyadic rational in it. Then each of the three pairs
\[
(\theta x_*,\,y_*),\qquad\bigl(\tfrac12(x_*+y_*),\,\tfrac\theta2(x_*+y_*)\bigr),\qquad(x_*,\,\theta y_*)
\]
has first entry above $\pi_0$ and second entry below $\pi_0$, so by
Proposition~\ref{prop:ideal}(a) its image under $\Phi$ is $\succ\!\!\succ\Phi(\pi_0,\pi_0)=z_0$.

\emph{A neighbourhood.} For two states $z=(x,y)$ and $\tilde z=(\tilde x,\tilde y)$, the
densities of the inputs $u$ and $v$ of a gate, let
\[
\Gamma_1=\Phi\bigl(\tfrac\theta2(x+\tilde x),\tfrac12(y+\tilde y)\bigr),\quad
\Gamma_2=\Phi\bigl(\tfrac12(y+\tilde x),\tfrac\theta2(x+\tilde y)\bigr),\quad
\Gamma_3=\Phi\bigl(\tfrac12(x+\tilde x),\tfrac\theta2(y+\tilde y)\bigr).
\]
These are the ideal densities of the high and the low rail of the gate stage when
$\val(u)=\val(v)=0$, when $\val(u)=0\ne\val(v)$, and when $\val(u)=\val(v)=1$; the fourth case
is $\Gamma_2$ with the inputs exchanged. They are continuous in $(z,\tilde z)$, and at
$z=\tilde z=L$ they are the three images of the previous paragraph. So there are $e_0\in(0,e_1]$, $\xi\in(0,(x_*-y_*)/4)$ and $\zeta>0$
such that, whenever $z$ and $\tilde z$ are within distance $\xi$ of $L$ in the maximum norm, each
$\Gamma_i(z,\tilde z)$ has first coordinate at least $\pi_0+e_0/(1-\eps)+\zeta$ and second
coordinate at most $\pi_0-e_0/(1-\eps)-\zeta$.

\emph{Stages.} By Proposition~\ref{prop:ideal}(d) there is $J$ such that $T^Jz_{e_0}$ and
$T^J(1,0)$ are within $\xi/2$ of $L$.

\emph{Perturbation.} For $(\eta,\gamma,h)$ consider the maps
$\Phi^-(d,d')=(\mathrm{lo}(d,d'),\mathrm{up}(d',d))$ and
$\Phi^+(d,d')=(\mathrm{up}(d,d'),\mathrm{lo}(d',d))$ of $[0,1]^2$. For $\eta\le\varpi/2$ the
arguments of $F$ stay in $[A/2,1)$, where $F$ is Lipschitz, $|\ell^\pm-\ell|\le\eps\eta$, and
$|F^\pm_h-F|\le\tau^h\le(1-A/2)^{2h}$ by Lemma~\ref{lem:latch}(c). So $\Phi^-$ and $\Phi^+$
converge to $\Phi$ uniformly on $[0,1]^2$ as $\eta\to0$, $\gamma\to0$ and $h\to\infty$. $\Phi$
maps $[0,1]^2$ into a compact subset of $(0,1)^2$, so for small perturbations $\Phi^\pm$ map
$[0,1]^2$ into $[0,1]^2$. Define
\[
(a_0,b_0)=z_{e_0},\quad(a_k,b_k)=\Phi^-(a_{k-1},b_{k-1}),\qquad
(a'_0,b'_0)=(1,0),\quad(a'_k,b'_k)=\Phi^+(a'_{k-1},b'_{k-1}).
\]
Then (R) holds with equality. As $\mathrm{lo}\le\mathrm{up}$ and both are monotone, induction
gives $a_k\le a'_k$ and $b'_k\le b_k$. The $J$-fold compositions of $\Phi^-$ and of $\Phi^+$ converge uniformly to $T^J$, because $T$ is
uniformly continuous. So for small perturbations $(a_J,b_J)$ and $(a'_J,b'_J)$ are within $\xi$ of
$L$, which gives (V), and $\Phi^-$ is within $\zeta$ of $\Phi$. The three pairs of (G) are the
arguments of $\Phi$ in $\Gamma_1(z,z)$, $\Gamma_2(z',z)$ and $\Gamma_3(z,z)$ for $z=(a_J,b_J)$ and
$z'=(a'_J,b'_J)$. Hence for each of them $\mathrm{lo}(d,d')\ge\pi_0+e_0/(1-\eps)=a_0$ and
$\mathrm{up}(d',d)\le\pi_0-e_0/(1-\eps)=b_0$, which is (G).

\emph{No parameters for $\eps\ge4/9$.} Suppose $(\theta,h,J,\eta,\gamma)$ is admissible, with
numbers $a_k,b_k$, and put $\sigma_k=(a_k,b_k)$. By Lemma~\ref{lem:latch}(c), for all $d,d'$ the
first coordinate of $\Phi(d,d')$ is at least $\mathrm{lo}(d,d')$ and the second is at most
$\mathrm{up}(d',d)$. So (R) gives $\sigma_k\preceq T\sigma_{k-1}$, hence
$\sigma_J\preceq T^J\sigma_0$ by monotonicity, and the first pair of (G), with $\theta\le1$,
gives $\sigma_0\preceq\Phi(\theta a_J,b_J)\preceq T\sigma_J$. Therefore
$\sigma_0\preceq T^{n(J+1)}\sigma_0$ for every $n$. For $\eps\ge4/9$ the only fixed point of $T$
is $(\pi_0,\pi_0)$ (Proposition~\ref{prop:ideal}(c)). The sequences $T^n(1,0)$ and $T^n(0,1)$ are
monotone and converge to fixed points, so both converge to it, and so does $T^nz$ for every state
$z$, which lies between $(0,1)$ and $(1,0)$. Hence $\sigma_0\preceq(\pi_0,\pi_0)$, then
$\sigma_J\preceq T^J\sigma_0\preceq(\pi_0,\pi_0)$, that is $a_J\le\pi_0\le b_J$, which
contradicts (V).
\end{proof}

\begin{corollary}[noise $2/5$, explicit]\label{cor:twofifths}
For $\eps=2/5$ the parameters $\theta=1/2$, $h=40$, $J=75$, $\eta=\gamma=1/500$ are admissible,
with
\[
a_0=0.5269,\quad b_0=0.3270,\qquad a_J=0.663,\quad a'_J=0.717,\quad b'_J=0.133,\quad b_J=0.173 .
\]
So in the trace of $S_{2/5}(\mathrm{Latch}(C))$ with these parameters, on $E_{1/500}$ and outside
an event of probability at most $304(s+2)e^{-m/125000}$, at least $66.3\%$ of the copies of
$P_J(g_s)$ hold $1$ if $\val(C)=1$, and at most $17.3\%$ if $\val(C)=0$.
\end{corollary}

\begin{proof}
The script \texttt{latch\_check.py} computes the numbers of Definition~\ref{def:adm} in exact
rational arithmetic, rounding every lower bound down and every upper bound up to a multiple of
$10^{-18}$ and enclosing $\varpi(2/5)=(9-\sqrt{65})/4$ between rationals. With the box above for
stage $J$ it finds $\mathrm{lo}\ge0.5269$ and $\mathrm{up}\le0.3270$ for the three pairs of (G)
(the least and the largest values are $0.52694$ and $0.32691$), and then, starting from
$a_0=0.5269$, $b_0=0.3270$, $a'_0=1$, $b'_0=0$ and taking equality in (R), it arrives at
$a_{75}\ge0.6637$, $b_{75}\le0.1721$, $a'_{75}\le0.7167$ and $b'_{75}\ge0.1336$, which lie inside
the box (Section~\ref{sec:checks}). The value $\theta=1/2$ lies outside the interval used in the
proof of Proposition~\ref{prop:adm}, which is $(0.577,0.948)$ here; that interval is sufficient,
not necessary. The last sentence is Theorem~\ref{thm:latch} with $4(J+1)=304$ and
$2\gamma^2=1/125000$.
\end{proof}

\subsection{Hardness}

Fix $\eps\in(0,4/9)$ and admissible parameters $(\theta,h,J,\eta,\gamma)$ for it
(Proposition~\ref{prop:adm}). Fix integers $\rho\ge1$ and $k_0\ge3$ and a number
$\tau_0\in(0,1/8]$, to be used as $\rho$, $k$ and $\tau$ in Lemma~\ref{lem:dens2any}, with
\[
\frac1{1+\rho}+\tau_0\le\frac14\qquad\text{and}\qquad
\frac{1/(1+\rho)+\tau_0}{1-\eps}+\frac{\eps^{2k_0-1}}2\le\eta ,
\]
and let
\[
\Lambda_0=(5/\tau_0^2)^{2k_0+1},\qquad M=\Bigl\lceil\frac{1+\ln(12(J+1))}{2\gamma^2}\Bigr\rceil,\qquad
N_1=2h(1+\rho)(6J+2),
\]
and let $Z_0$ be the least integer with $Z_0\ge\Lambda_0\bigl(2h+5+\tfrac43(\ln Z_0+4)\bigr)+N_1$.
All of these depend on $\eps$ only.

\begin{definition}[$\mathrm{Latch}_a$]\label{def:latcha}
For an integer $a\ge1$ and a circuit with $s\ge2$ gates, $\mathrm{Latch}_a(C)$ is the base
$\mathrm{Latch}(C)$ with the parameters $\theta,h,J,\rho$ fixed above and
\begin{align*}
m&=\text{the least power of two with }m\ge M(a+1)\log_2s,\\
Z&=Z_0(a+2)^2sm,\qquad b=\lceil\log_2Z\rceil,\qquad
\mathrm{pre}=\bigl\lceil\Lambda_0\bigl(2h+3+(a+3)\,b\ln2\bigr)\bigr\rceil .
\end{align*}
\end{definition}

\begin{lemma}[parameters]\label{lem:param4}
For $\mathrm{Latch}_a(C)$: (i) $n+1\le Z$; (ii)
$\mathrm{pre}\ge\Lambda_0\bigl(2h+3+(a+3)\ln(n+1)\bigr)$; (iii) there is $s^\ast_a$, depending on
$\eps$ and $a$ only, such that $n+1\le s^2$ whenever $s\ge s^\ast_a$.
\end{lemma}

\begin{proof}
(i) $b\ln2\le\ln Z+1=\ln Z_0+1+2\ln(a+2)+\ln(sm)$. With $\ln x\le x$ we have
$2\ln(a+2)\le2(a+2)sm$ and $\ln(sm)\le(a+2)sm$, so $b\ln2\le(\ln Z_0+4)(a+2)sm$. With
$a+3\le\frac43(a+2)$ and $\Lambda_0\ge1$ this gives
$\mathrm{pre}+1\le\Lambda_0\bigl(2h+5+\frac43(\ln Z_0+4)\bigr)(a+2)^2sm$. The chains and their
padding occupy $2h(1+\rho)m(2(J+1)s+4J)\le N_1sm$ positions. So $n+1\le Z$ by the choice of $Z_0$.
(ii) $b\ln2\ge\ln Z\ge\ln(n+1)$ by (i).
(iii) $m<2M(a+1)\log_2s$, so $n+1\le Z<2Z_0M(a+2)^3s\log_2s$, which is at most $s^2$ for large $s$.
\end{proof}

\begin{corollary}[embedding below $4/9$]\label{cor:latch}
Let $\eps\in(0,4/9)$, $a\ge1$, and let $C$ be a NOR circuit with $s\ge2$ gates. In the trace of
$S_\eps(\mathrm{Latch}_a(C))$,
\[
\PP[\text{some rail is not good}]\ \le\ s^{-a}+(n+1)^{-(a+1)} .
\]
In particular, with at least the complementary probability, at least a fraction $a_J$ of the
copies of $P_J(g_s)$ hold $1$ if $\val(C)=1$, and at most a fraction $b_J<a_J$ if $\val(C)=0$.
\end{corollary}

\begin{proof}
The base consists of zero gates and chain gates, positions $1,\dots,\mathrm{pre}$ are zero gates,
and every chain of $2h$ gates is preceded by $2h\rho$ zero gates, so every interval $(p,p']$ of
positions contains at most $(p'-p)/(1+\rho)+2h$ chain gates. Which positions are zero gates does
not depend on the coins. So Lemma~\ref{lem:dens2any} can be applied conditionally on the coins,
with the chain gates as its cell gates, $c=2h$, $K=a+1$ and the $\rho,\tau_0,k_0$ fixed above; its
condition on $\mathrm{pre}$ is Lemma~\ref{lem:param4}(ii), and its error is at most $\eta$. So
$\PP[E_\eta^c]\le(n+1)^{-(a+1)}$. By Theorem~\ref{thm:latch} the rest is at most
$4(J+1)(s+2)e^{-2\gamma^2m}\le12(J+1)\,s\,e^{-2\gamma^2m}$. Now
$2\gamma^2m\ge(1+\ln(12(J+1)))(a+1)\log_2s\ge(a+1)\ln s+\ln(12(J+1))$, since
$(a+1)\log_2s\ge\max(1,(a+1)\ln s)$. So the rest is at most $s^{-a}$. The last sentence is the definition of a good rail, and
$b_J<a_J$ is (V).
\end{proof}

\begin{theorem}[hardness below noise $4/9$]\label{thm:hard4}
Theorem~\ref{thm:hard} holds for every $\eps<4/9$: if for some $\eps<4/9$ a randomized NC family
outputs the trace with probability at least $n^{-K}$ on the $\eps$-smoothing of every base of
every length $n\ge n_0$, then $\mathrm P\subseteq\mathrm{BPNC}$.
\end{theorem}

\begin{proof}
Let $\eps^*$ be a rational number with $\eps\le\eps^*<4/9$ and $\eps^*>0$. By Lemma~\ref{lem:errorless} and
Lemma~\ref{gb:comp}, as in the proof of Theorem~\ref{thm:hard}, we may assume that the family is
errorless and has success probability at least $n^{-K}$ at noise $\eps^*$. Repeat the proof of
Theorem~\ref{thm:hard2} at the noise level $\eps^*$, with $\mathrm{Latch}_a$ (for $\eps^*$) in
place of $\mathrm{Tree}_a$, with $s\ge\max(s^\ast_a,n_0,3)$, and with
Corollary~\ref{cor:latch} and Lemma~\ref{lem:param4}(iii) in place of
Corollary~\ref{cor:threetwentieths} and Lemma~\ref{lem:param2}(iii). Three points change. The
coins of a chain gate are a bounded number of fair bits for the choice within $\frac12S+\frac12S'$
and within $\theta S$, which are exact because $\theta$ is dyadic, and $\log_2m$ bits for the
copy. The indicators $\pi_i$ are drawn by comparing a random integer with a threshold, as there;
the threshold is computable because $\eps^*$ is rational. And a trial votes $1$ if more than a
fraction $\beta$ of the copies of $P_J(g_s)$ hold $1$ in the trace it received, and $0$ otherwise,
where $\beta$ is a fixed rational number with $b_J<\beta<a_J$; by Corollary~\ref{cor:latch} this
vote differs from $\val(C)$ with probability at most $\delta=s^{-a}+(n+1)^{-(a+1)}$. With $a=2K+4$ the rest is unchanged.
\end{proof}

\begin{corollary}[the gap, with the lower end $4/9$]\label{cor:gap4}
Corollary~\ref{cor:gap} holds with ``$\eps<4/9$'' in place of ``$\eps\le1/64$'' in part (b). In
particular $\inf E_K\in\{0\}\cup[4/9,3/4]$ for every $K$, and if
$\mathrm P\not\subseteq\mathrm{BPNC}$, the critical noise level satisfies
$4/9\le\eps_c\le3/4$.
\end{corollary}

\begin{proof}
As for Corollary~\ref{cor:gap2}, with Theorem~\ref{thm:hard4}.
\end{proof}

\begin{remark}[the level $4/9$]\label{rem:latch}
\begin{enumerate}[label=(\roman*),leftmargin=2.4em]
\item \emph{The analysis stops at $4/9$.} By~\eqref{eq:uv}, one ideal stage multiplies the
difference of the two rails by $(1-\eps)/(1-uv)$, where $u$ and $v$ are the probabilities that a
set check and a reset check read a zero. For $\eps\ge4/9$ the only fixed point of the ideal map is
the diagonal one, where the two rails look alike, and no parameters are admissible
(Proposition~\ref{prop:adm}). That the construction itself fails there for every choice of
parameters is not proved; in simulation the two rails could not be told apart at noise $0.46$ and
$0.50$ (Section~\ref{sec:checks}). For $\eps<4/9$ the fixed point $L$ is the following, with $\lambda=(1-\eps)/(1-q^2)$ the multiplier at
the diagonal fixed point (script \texttt{latch\_check.py table}):
\begin{center}
\begin{tabular}{lcccccccc}
\toprule
$\eps$ & $0.10$ & $0.25$ & $0.30$ & $0.35$ & $0.40$ & $0.42$ & $0.43$ & $0.44$\\
\midrule
$x_*$ & $0.990$ & $0.922$ & $0.876$ & $0.807$ & $0.693$ & $0.621$ & $0.571$ & $0.495$\\
$y_*$ & $0.001$ & $0.021$ & $0.041$ & $0.077$ & $0.150$ & $0.204$ & $0.244$ & $0.310$\\
$\lambda$ & $1.525$ & $1.326$ & $1.248$ & $1.165$ & $1.079$ & $1.044$ & $1.026$ & $1.008$\\
\bottomrule
\end{tabular}
\end{center}
The diagonal fixed point has $\pi_0$ between $0.390$ and $0.400$ throughout.
\item \emph{The same inequality as in Section~\ref{sec:block}.} At the diagonal fixed point a
chain reads uniformly random copies of the stage below, whatever the rail. That is the situation
of a random block base over zero gates (Definition~\ref{def:block}), with $q=q(\eps)$ the
probability that a gate does not read a one. The condition for $L$ to exist, $q^2>\eps$, is the
condition of Theorem~\ref{thm:blockzero}(b), under which the relevance-pruned iteration needs
$n/\log n$ rounds. Both say that $(1-\eps)/(1-q^2)>1$: a copy depends, through relevant gates after
its last reset, on more than one copy of the stage below on average, counting only listed
pointers that were not resampled. In Section~\ref{sec:block} this makes long dependency chains;
here it makes a difference between the rails grow. Also $\pi_0=q/(1+q)$ is the density of ones
of such blocks. This is an explanation of the coincidence, not a theorem that relates the two
statements.
\item \emph{What is open.} For every $\eps<4/9$ there is no randomized NC algorithm with
inverse-polynomial success on all bases unless $\mathrm P\subseteq\mathrm{BPNC}$, and for
$\eps>3/4$ there is a logspace algorithm. Section~\ref{sec:s24} extends the first statement to
every $\eps\le\eps_S=3661/8192$, and Section~\ref{sec:allbases} extends the second to every
$\eps>\eps_1$, with $\eps_1<0.7288$ and depth $O(\log n)$ in place of logspace; between $\eps_S$ and
$\eps_1$ nothing is proved for all bases.
Among all periodic patterns of at most six pairs in which some checks read position $1$ instead
of a rail, the plain latch has the highest level in the ideal model (script
\texttt{latch\_patterns.py}). Longer designs do better: Section~\ref{sec:designs} gives one, of
period $24$, whose pairs of layers have two attracting states in the ideal model up to noise about
$0.447$ (numerically), and it proves that designs with independent reads cannot keep one from noise $0.455$ on. That
design gives hardness up to $0.4468$ (Section~\ref{sec:s24}), so $4/9$ is not the critical noise
level unless $\mathrm P\subseteq\mathrm{BPNC}$; by Section~\ref{sec:allbases} that level is below
$0.7288$, and where it lies in between is open. On random block bases the relevance-pruned
iteration is fast above $4/9$ (Theorem~\ref{thm:blockabove}); on the mixed base of
Remark~\ref{rem:mixed} it is slow up to about $0.55$ in simulation.
\item \emph{The gate stage is soft.} Unlike the embeddings of Sections~\ref{sec:tree}
and~\ref{sec:tree2}, this one is not exact without noise: a chain of $P_0(g)$ whose set checks are
all thinned away keeps its initial state, and when the inputs of a gate differ, both rails of its
stage $0$ receive sets and resets. Without noise the rails separate only over the stages, and only
for most choices of the coins. The results of this section are stated for $\eps>0$.
\item \emph{Constants.} For $\eps=2/5$ and the parameters of Corollary~\ref{cor:twofifths} one can
take $\rho=1999$, $\tau_0=1/2000$ and $k_0=5$ (density error at most $0.0018$), and then
$M=976{,}955$ and $\Lambda_0\approx2\cdot10^{80}$. So a copy costs $80$ chain gates and
$159{,}920$ zero gates, a bundle has about $10^6(a+1)\log_2s$ copies, and the preamble is
astronomical, as in Section~\ref{sec:tree2}. The statements are asymptotic. As $\eps\to4/9$ the
multiplier $\lambda$ tends to $1$ and all constants grow without bound.
\item \emph{Simulation.} The script \texttt{latch\_embed.c} builds $\mathrm{Latch}(C)$ for a random
circuit, with small parameters, and evaluates its smoothing; the results are in
Section~\ref{sec:checks}. Every gate of the circuit came out right in all runs at noise $0.30$,
$0.35$, $0.40$ and $0.42$, the densities of the rails were close to the fixed point of the ideal
map for the measured density of ones, and at noise $0.46$ and $0.50$ the two rails could not be
told apart.
\end{enumerate}
\end{remark}

\section{Above \texorpdfstring{$4/9$}{4/9}: an embedding with a design of period \texorpdfstring{$24$}{24}}\label{sec:s24}

Section~\ref{sec:latch} stops at $4/9$ because the ideal map of a latch stage has a fixed point
off the diagonal only below $4/9$ (Proposition~\ref{prop:ideal}). This section replaces the latch
by chains that run a fixed program of period $24$, the design $S_{24}$ of
Remark~\ref{rem:beyond}, and proves Theorem~\ref{thm:hards24}: the conclusion of
Theorem~\ref{thm:hard} holds for every
\[
\eps\le\eps_S:=\frac{3661}{8192}=0.44689\ldots
\]
So, unless $\mathrm P\subseteq\mathrm{BPNC}$, the critical noise level is at least $\eps_S>0.4468$, and
$4/9$ is not its value (Corollary~\ref{cor:gaps24}). Of the earlier sections it uses what
Section~\ref{sec:latch} uses, and Lemma~\ref{lem:reads} and the proof of
Theorem~\ref{thm:latch}.

\emph{What changes.}
\begin{itemize}[leftmargin=2em]
\item \emph{A bit is a phase.} A stage is one bundle of chains, and it feeds five rails, one for
each phase of the program that some gate reads. In the ideal model the map that sends the
densities of these rails to those of the next stage reverses an order (even phases up, odd phases
down), so the map of two stages preserves it, and that map has two attracting fixed points up to
noise about $0.447$ (Remark~\ref{rem:beyond}(ii)). The stages of a wire oscillate between two
states, and the bit is which of the two the even stages are near.
\item \emph{The sandwich works stage by stage.} The value of a chain gate is a unate function of
what the chain reads: whether it rises or falls with a read depends only on the parity of the
distance (Lemma~\ref{lem:progchain}). Every gate of the program reads a phase of its own parity.
So the bounds of the proof of Theorem~\ref{thm:latch} carry over, with order intervals in place of
pairs of numbers; a map that reverses the order sends an order interval into the order interval
between the images of its two ends, swapped.
\item \emph{The gate reads both inputs at once,} with a small bias toward the top state. Two low
inputs give a high output, two high inputs a low one, and a mixed pair, which lands near the
middle, is pushed to the low side by the bias.
\item \emph{The margins are thin.} The certificate (Proposition~\ref{prop:s24cert}) uses $6000$
restoring stages and slacks of $10^{-6}$, computed in integer arithmetic with outward rounding. All constants are
astronomical, more so than in Section~\ref{sec:latch}; the statements are asymptotic.
\end{itemize}

\subsection{The program and its chains}

\begin{definition}[the program $S_{24}$]\label{def:s24}
For a phase $\varphi\in\{0,\dots,23\}$ let $\kappa(\varphi)$ be the entry of number $\varphi$
(counting from $0$) in
\[
22,\,9,\,22,\,9,\,22,\,9,\,22,\,9,\,22,\,9,\,\ast,\,3,\,8,\,9,\,8,\,9,\,8,\,9,\,8,\,11,\,8,\,11,\,8,\,9 .
\]
The \emph{read phases} are $\mathcal L=\{3,8,9,11,22\}$. Every entry $\kappa(\varphi)\ne\ast$ has the
parity of $\varphi$. For $h\ge1$ and sources $(S_c)_{c\in\mathcal L}$, a \emph{program chain} is a
sequence of $24h$ consecutive gates $g_1<\dots<g_{24h}$: gate $g_r$ has phase
$\varphi=(r-1)\bmod24$, and it lists position $1$ if $\kappa(\varphi)=\ast$ and otherwise reads
from the source $S_{\kappa(\varphi)}$, with independent coins. For $c\in\mathcal L$ its \emph{copy of
phase $c$} is the gate $g_{o(c)}$, $o(c)=24(h-1)+c+1$, the gate of phase $c$ of its last period.
\end{definition}

\noindent This is the design $S_{24}$ of Remark~\ref{rem:beyond}, with $\mathrm Z_1$ written $\ast$;
in both, the gate lists position $1$, which holds $0$.

\begin{lemma}[program chains]\label{lem:progchain}
Let $x_0,y_1,\dots,y_N$ be bits and $x_r=\NOR(x_{r-1},y_r)$ for $1\le r\le N$.
\begin{enumerate}[label=(\alph*),leftmargin=2.4em]
\item For $1\le r\le o\le N$, $x_o$ is nondecreasing in $y_r$ if $o-r$ is odd and nonincreasing if
$o-r$ is even. It is nondecreasing in $x_0$ if $o$ is even and nonincreasing if $o$ is odd.
\item If $x_0$ is fixed and $y_1,\dots,y_N$ are independent with $\PP[y_r=1]=\ell_r$, then
$p_r=\PP[x_r=1]$ satisfies $p_0=x_0$ and $p_r=(1-p_{r-1})(1-\ell_r)$. For $r\le o$, $p_o$ is
nondecreasing in $\ell_r$ if $o-r$ is odd and nonincreasing if $o-r$ is even.
\end{enumerate}
\end{lemma}

\begin{proof}
(a) Induction on $o$. $x_o$ is nonincreasing in $y_o$ and in $x_{o-1}$, and every other variable
enters only through $x_{o-1}$, in which, by induction, its direction is the opposite of the one
claimed for $x_o$. (b) $x_r=1$ exactly when $x_{r-1}=0$ and $y_r=0$, and $y_r$ is independent of
$x_{r-1}$, a function of $x_0,y_1,\dots,y_{r-1}$. For the monotonicity write $y_r$ as the
indicator of $\{\xi_r<\ell_r\}$ with independent uniform $\xi_r$ and use (a).
\end{proof}

For a vector $d=(d_c)_{c\in\mathcal L}\in[0,1]^\mathcal L$ let $t(d)$ be the vector with $t(d)_c=d_c$ for
even $c$ and $t(d)_c=1-d_c$ for odd $c$. Vectors in $[0,1]^\mathcal L$ are compared coordinatewise, and
$\mathbf 1$ and $\mathbf 0$ are the vectors with all entries $1$ and $0$. So $t(d)\le t(d')$ means
that $d'$ is higher than $d$ on the even phases and lower on the odd ones.

\subsection{The compiler}

We use the sources and rails of Section~\ref{sec:latch}: position $0$ holds $1$ and position $1$
holds $0$, the rail $\mathbf 1$ is position $0$ repeated and the rail $\mathbf 0$ is position $1$
repeated. For $c\in\mathcal L$ let $K_c$ be the rail $\mathbf 1$ if $c$ is even and $\mathbf 0$ if $c$
is odd. For a dyadic rational $\beta\in(0,1)$ and sources $S,S'$, the source
$\frac{1-\beta}2S+\frac{1-\beta}2S'+\beta K_c$ draws from $S$, from $S'$ and from $K_c$ with these
probabilities, using fair coins.

\begin{definition}[$\mathrm{Prog}(C)$]\label{def:prog}
Let $C$ be a NOR circuit with gates $g_t=\NOR(u_t,v_t)$, $t\le s$. Fix integers $h,\rho\ge1$, an
even $J\ge2$, a bundle size $m$ that is a power of two, a dyadic rational $\beta\in(0,1)$ and a
preamble length $\mathrm{pre}$.
\begin{itemize}[leftmargin=2em]
\item A \emph{block} is $24h\rho$ zero gates followed by a program chain. A \emph{bundle} with
sources $(S_c)_{c\in\mathcal L}$ is a sequence of $m$ consecutive blocks whose chains have these
sources. For $c\in\mathcal L$, its \emph{rail of phase $c$} is the list of the copies of phase $c$ of
its $m$ chains.
\item Every gate $w$ of $C$ has a bundle $B_k(w)$ for every stage $k=0,\dots,J$, and each constant
$w\in\{c_0,c_1\}$ one for every $k=1,\dots,J$. The rails of $B_k(w)$ are written $R_k(w,c)$. For
the constants put $R_0(c_1,c)=K_c$ and let $R_0(c_0,c)$ be the other constant rail.
\item \emph{Gate stage.} For $g=\NOR(u,v)$ the bundle $B_0(g)$ has the sources
$S_c=\frac{1-\beta}2R_J(u,c)+\frac{1-\beta}2R_J(v,c)+\beta K_c$.
\item \emph{Restoring stages.} For $1\le k\le J$ the bundle $B_k(w)$ has the sources
$S_c=R_{k-1}(w,c)$.
\end{itemize}
The base $\mathrm{Prog}(C)$ consists of zero gates at positions $1,\dots,\mathrm{pre}$, followed by
the bundles of stages $1,\dots,J$ of $c_0$ and then of $c_1$, and then, for $t=1,\dots,s$ in this
order, by the bundles of stages $0,\dots,J$ of $g_t$. Its length is
\[
n=\mathrm{pre}+24h(1+\rho)\,m\,\bigl((J+1)s+2J\bigr).
\]
\end{definition}

As for $\mathrm{Latch}(C)$, all positions are explicit arithmetic expressions, a chain gate uses
a bounded number of fair coins and $\log_2m$ coins for the copy, and every chain reads only
positions $0$, $1$ and copies that lie before it; so the base is computable in logspace from $C$
and the coins. With $\val(c_1)=1$, $\val(c_0)=0$ and the vector
$\mathrm{dens}_k(w)=\bigl(\mathrm{dens}(R_k(w,c))\bigr)_{c\in\mathcal L}$, the constants have
$t(\mathrm{dens}_0(c_1))=\mathbf 1$ and $t(\mathrm{dens}_0(c_0))=\mathbf 0$.

\subsection{Envelopes, admissible parameters and the embedding theorem}

Fix $\eps\in(0,1)$, write $\varpi=\varpi(\eps)$, and let $\eta\in(0,\varpi)$, $\gamma\ge0$ and
$h\ge1$. As in Section~\ref{sec:latch}, $\ell^\pm(d)=(1-\eps)d+\eps(\varpi\pm\eta)$. For
$c\in\mathcal L$, $\sigma\in\{-,+\}$ and $d\in[0,1]^\mathcal L$, let $P^\sigma_c(d)$ be the number $p_o$ of
Lemma~\ref{lem:progchain}(b) for $o=o(c)$ and the following choices: $x_0$ is the bit for which
$p_o$ is smallest ($\sigma=-$) or largest ($\sigma=+$), that is, $x_0=0$ for $\sigma=-$ and $x_0=1$
for $\sigma=+$ if $o$ is even, and the other way round if $o$ is odd; and
$\ell_r=\ell^{s_r}(d_{\kappa(\varphi)})$ for the phase $\varphi$ of $r$, with $d_\ast=0$, where
$s_r=\sigma$ if $o-r$ is odd and $s_r=-\sigma$ if $o-r$ is even. In words, every read is given the
probability, among $\ell^-$ and $\ell^+$, that pushes $p_o$ down (for $\sigma=-$) or up. Define
$\Psi^-,\Psi^+:[0,1]^\mathcal L\to\mathbb R^\mathcal L$ by
\[
\Psi^-(d)_c=\begin{cases}P^-_c(d)-\gamma,\\ 1-P^+_c(d)-\gamma,\end{cases}\qquad
\Psi^+(d)_c=\begin{cases}P^+_c(d)+\gamma&c\text{ even},\\ 1-P^-_c(d)+\gamma&c\text{ odd}.\end{cases}
\]
With $\eta=\gamma=0$ and $h\to\infty$ both tend to $t\circ\Phi$, where $\Phi$ is the density map of
Definition~\ref{def:design} for $S_{24}$, restricted to the read phases.

\begin{lemma}[the envelopes reverse the order]\label{lem:envelope}
If $t(d)\le t(d')$ then $\Psi^\pm(d')\le\Psi^\pm(d)$. Also $\Psi^-\le\Psi^+$.
\end{lemma}

\begin{proof}
Let $r\le o=o(c)$ and let $\varphi=(r-1)\bmod24$ be the phase of $r$. Then
$o-r\equiv c-\varphi\pmod2$, and $\kappa(\varphi)\equiv\varphi$ if $\kappa(\varphi)\ne\ast$. So,
by Lemma~\ref{lem:progchain}(b) and since $\ell^\pm$ is increasing, $P^\pm_c$ is nondecreasing in
$d_{c'}$ if $c$ and $c'$ have different parities and nonincreasing if they have the same parity.
Raising $t(d)_{c'}$ raises $d_{c'}$ for even $c'$ and lowers it for odd $c'$. For even $c$ this
lowers $P^\pm_c=\Psi^\pm(d)_c\mp\gamma$ in both cases; for odd $c$ it raises $P^\pm_c$, and so lowers
$1-P^\pm_c$. The last claim holds because $P^-_c\le P^+_c$: the choices for $\sigma=-$ give, read by
read and at the start, the smaller of the two probabilities.
\end{proof}

We write $\Psi^\pm(t)$ for $\Psi^\pm(d)$ when $t=t(d)$, and $\mathrm{Mix}(x,y)=\frac{1-\beta}2(x+y)+\beta\mathbf 1$
for $x,y\in[0,1]^\mathcal L$. The source $S_c$ of a gate stage has density $d_c$ with
$t(d)=\mathrm{Mix}\bigl(t(\mathrm{dens}_J(u)),t(\mathrm{dens}_J(v))\bigr)$, because $t$ is affine in each
coordinate with slope $\pm1$, the weights sum to $1$, and $t$ of the densities of the rails $K_c$
is $\mathbf 1$.

\begin{definition}[admissible parameters for $S_{24}$]\label{def:adms24}
Parameters $(\beta,h,J,\eta,\gamma)$, with $\beta$ a dyadic rational in $(0,1)$ and $J\ge2$ even, are
\emph{admissible for} $\eps$ if there are vectors $L^v_k\le U^v_k$ in $[0,1]^\mathcal L$, for
$k=0,\dots,J$ and $v\in\{0,1\}$, with $U^1_0=\mathbf 1$ and $L^0_0=\mathbf 0$, such that:
\begin{itemize}[leftmargin=3em]
\item[(R)] for $1\le k\le J$ and $v\in\{0,1\}$,
\ $L^v_k\le\Psi^-(U^v_{k-1})$ and $U^v_k\ge\Psi^+(L^v_{k-1})$;
\item[(G)] $L^1_0\le\Psi^-\bigl(\mathrm{Mix}(U^0_J,U^0_J)\bigr)$,
\ $U^0_0\ge\Psi^+\bigl(\mathrm{Mix}(L^0_J,L^1_J)\bigr)$ and
\ $U^0_0\ge\Psi^+\bigl(\mathrm{Mix}(L^1_J,L^1_J)\bigr)$;
\item[(V)] $(L^1_J)_8>(U^0_J)_8$.
\end{itemize}
\end{definition}

\noindent The vectors $L^v_k,U^v_k$ will bound $t(\mathrm{dens}_k(w))$ for a wire of value $v$. Let
$E_\eta$ be as in Section~\ref{sec:latch}: $|\psi_g-\varpi|\le\eta$ for all $g$ with
$\mathrm{pre}\le g\le n$.

\begin{theorem}[embedding with $S_{24}$]\label{thm:progembed}
Let $\eps\in(0,1)$, let $(\beta,h,J,\eta,\gamma)$ be admissible for $\eps$ with the vectors
$L^v_k,U^v_k$, and let $C$ be a NOR circuit with $s$ gates. In the trace of $S_\eps(\mathrm{Prog}(C))$
call the bundle $B_k(w)$ \emph{good} if $L^v_k\le t(\mathrm{dens}_k(w))\le U^v_k$ for
$v=\val(w)$. Then, over the coins and the smoothing, for all $m$, $\rho$ and $\mathrm{pre}$,
\[
\PP\bigl[E_\eta\cap\{\text{some bundle is not good}\}\bigr]\ \le\ 10(J+1)(s+2)\,e^{-2\gamma^2m}.
\]
\end{theorem}

\begin{proof}
We follow the proof of Theorem~\ref{thm:latch}. The stage-$0$ vectors of the constants are
$\mathbf 1\in[L^1_0,\mathbf 1]$ and $\mathbf 0\in[\mathbf 0,U^0_0]$, so these are good. Let
$\mathcal B_1,\mathcal B_2,\dots$ be the other bundles in the order of the layout; there are at
most $(J+1)(s+2)$ of them. Define $\mathcal G_i$ and $H_i\in\mathcal G_i$ as there, with bundles in
place of rails: $\mathcal G_i$ is generated by the coins of the chains of
$\mathcal B_1,\dots,\mathcal B_{i-1}$ and by $(\pi_g,\rho_g,U_g)$ for the positions $g$ below the
first position of $\mathcal B_i$, and $H_i$ is the event that $\mathcal B_1,\dots,\mathcal B_{i-1}$
are good.

\emph{Sandwich.} Let $\mathcal B_i$ have sources $(S_c)$, let $c$ be a read phase, and let a chain
of $\mathcal B_i$ have the gates $g_1<\dots<g_{24h}$; put $o=o(c)$. For $\sigma\in\{-,+\}$ let
$X^\sigma(c)$ be the value $x_o$ of Lemma~\ref{lem:progchain} with the start chosen as in the
definition of $P^\sigma_c$ and with the reads $y^{s_r}_{g_r}$ of Lemma~\ref{lem:reads}, $s_r$ as
there. Every chain gate $g$ has $g-1\ge\mathrm{pre}$, so on $E_\eta$ we have
$y^-_{g_r}\le y_{g_r}\le y^+_{g_r}$, and by Lemma~\ref{lem:progchain}(a) the copy of phase $c$ of
the chain holds a value between $X^-(c)$ and $X^+(c)$, whatever the gate before the chain holds.

\emph{Conditional law.} As in the proof of Theorem~\ref{thm:latch}, given $\mathcal G_i$ the
indicators $y^\pm_g$ of the gates of $\mathcal B_i$ are independent across gates, with
$\PP[y^\pm_g=1\mid\mathcal G_i]=\ell^\pm(d)$, where $d$ is the density of the source of $g$, which is
$\mathcal G_i$-measurable, and $d=0$ for a gate that lists position $1$. Let $d=(d_c)_{c\in\mathcal L}$
be the vector of the densities of the sources $S_c$. So for each $c$ and $\sigma$ the variables
$X^\sigma(c)$ of the $m$ chains are independent given $\mathcal G_i$, and by
Lemma~\ref{lem:progchain}(b), $\PP[X^\sigma(c)=1\mid\mathcal G_i]=P^\sigma_c(d)$.

\emph{The bounds.} We claim that on $H_i$, if $\mathcal B_i=B_k(w)$ and $v=\val(w)$, then
\begin{equation}\label{eq:s24bounds}
L^v_k\le\Psi^-(d)\quad\text{unless }k=0\text{ and }v=0,\qquad
\Psi^+(d)\le U^v_k\quad\text{unless }k=0\text{ and }v=1 .
\end{equation}
For $k\ge1$ the sources are the rails of $B_{k-1}(w)$, which is good on $H_i$, so
$L^v_{k-1}\le t(d)\le U^v_{k-1}$. By Lemma~\ref{lem:envelope},
$\Psi^-(d)\ge\Psi^-(U^v_{k-1})\ge L^v_k$ and $\Psi^+(d)\le\Psi^+(L^v_{k-1})\le U^v_k$, by (R). For
$k=0$ and $w=g=\NOR(u,v')$, the bundles $B_J(u)$ and $B_J(v')$ are good on $H_i$ (this includes
constants, and $u=v'$), and $t(d)=\mathrm{Mix}(t(\mathrm{dens}_J(u)),t(\mathrm{dens}_J(v')))$, where $\mathrm{Mix}$ is
nondecreasing in both arguments. If $\val(u)=\val(v')=0$, then $\val(g)=1$ and
$t(d)\le \mathrm{Mix}(U^0_J,U^0_J)$, so $\Psi^-(d)\ge L^1_0$ by Lemma~\ref{lem:envelope} and (G). If the values
differ, then $\val(g)=0$ and $t(d)\ge \mathrm{Mix}(L^0_J,L^1_J)$, in either order of the inputs; if both are
$1$, then $\val(g)=0$ and $t(d)\ge \mathrm{Mix}(L^1_J,L^1_J)$. In both cases $\Psi^+(d)\le U^0_0$. The bounds
excluded in~\eqref{eq:s24bounds} are those against $U^1_0=\mathbf 1$ and $L^0_0=\mathbf 0$, which
hold for every vector of densities.

\emph{Concentration.} Let $c$ be even. By~\eqref{eq:s24bounds}, given $\mathcal G_i$ and on $H_i$, the
number of chains with $X^-(c)=1$ is a sum of $m$ independent indicators with mean
$mP^-_c(d)\ge m((L^v_k)_c+\gamma)$; by Hoeffding's inequality it is less than $m(L^v_k)_c$ with
conditional probability at most $e^{-2\gamma^2m}$, and on $E_\eta$ the density of $R_k(w,c)$ is at
least that number divided by $m$. In the same way the number of chains with $X^+(c)=1$ bounds the
density from above, with mean $mP^+_c(d)\le m((U^v_k)_c-\gamma)$. For odd $c$ the two roles are
exchanged: $t(\mathrm{dens})_c=1-\mathrm{dens}(R_k(w,c))$ is at least $1$ minus the fraction of
chains with $X^+(c)=1$, whose mean is $P^+_c(d)\le1-(L^v_k)_c-\gamma$, and at most $1$ minus the
fraction with $X^-(c)=1$, whose mean is $P^-_c(d)\ge1-(U^v_k)_c+\gamma$. So a bundle fails one of
its at most ten bounds with probability at most $e^{-2\gamma^2m}$ each, on $E_\eta\cap H_i$. If some
bundle is not good, let $\mathcal B_i$ be the first one; then $H_i$ holds and one of its bounds
fails. A union bound gives the theorem.
\end{proof}

\begin{proposition}[a certificate at noise $3661/8192$]\label{prop:s24cert}
For $\eps=\eps_S=3661/8192$ the parameters
\[
\beta=3/128,\qquad h=4,\qquad J=6000,\qquad \eta=\gamma=10^{-6}
\]
are admissible, with vectors that satisfy
\[
(L^1_J)_8\ge0.4426,\qquad(U^0_J)_8\le0.3982 .
\]
So in the trace of $S_{\eps_S}(\mathrm{Prog}(C))$ with these parameters, on $E_{10^{-6}}$ and outside
an event of probability at most $60010(s+2)e^{-m/(5\cdot10^{11})}$, at least $44.26\%$ of the copies
in $R_J(g_s,8)$ hold $1$ if $\val(C)=1$, and at most $39.82\%$ if $\val(C)=0$.
\end{proposition}

\begin{proof}
The script \texttt{s24\_check.py} represents every number as an integer multiple of $2^{-96}$,
rounds every lower bound down and every upper bound up, and encloses
$\varpi(\eps_S)=((2\eps_S+1)-\sqrt{4\eps_S+1})/(2\eps_S)$ between rationals; it evaluates
$P^\pm_c$ by the recursion of Lemma~\ref{lem:progchain}(b) with every product rounded in the
direction that, by the alternation of Lemma~\ref{lem:progchain}(b), lowers (for $P^-_c$) or raises
(for $P^+_c$) the final value, so its values of $\Psi^-$ are at most and its values of $\Psi^+$ at
least the true ones. It fixes two boxes
for stage $J$, $Q^0=[\underline b^0,\overline b^0]$ and $Q^1=[\underline b^1,\overline b^1]$, of
radius $0.004$ around the two points of the ideal cycle of $S_{24}$ (computed in floating point;
this only places the boxes). It sets $L^1_0$ to its value of
$\Psi^-(\mathrm{Mix}(\overline b^0,\overline b^0))$ and $U^0_0$ to the coordinatewise maximum of its values of
$\Psi^+(\mathrm{Mix}(\underline b^0,\underline b^1))$ and $\Psi^+(\mathrm{Mix}(\underline b^1,\underline b^1))$, and
$U^1_0=\mathbf 1$, $L^0_0=\mathbf 0$; and for $k=1,\dots,J$ it sets $L^v_k$ and $U^v_k$ to its
values of $\Psi^-(U^v_{k-1})$ and $\Psi^+(L^v_{k-1})$. Then (R) holds, the script asserts that all
these vectors lie in $[0,1]^{\mathcal L}$, and $L^v_k\le U^v_k$ by induction on $k$: if
$L^v_{k-1}\le U^v_{k-1}$, then by Lemma~\ref{lem:envelope}
$L^v_k\le\Psi^-(U^v_{k-1})\le\Psi^+(U^v_{k-1})\le\Psi^+(L^v_{k-1})\le U^v_k$. It checks that
$[L^v_J,U^v_J]\subseteq Q^v$ for $v=0,1$; with Lemma~\ref{lem:envelope} this gives (G), because $\mathrm{Mix}$ is
nondecreasing. It finds $(L^1_J)_8\ge0.44265$ and $(U^0_J)_8\le0.39816$, which is (V). The other
coordinates and the boxes are listed in Section~\ref{sec:checks}. The last sentence is
Theorem~\ref{thm:progembed} with $10(J+1)=60010$ and $2\gamma^2=2\cdot10^{-12}$.
\end{proof}

\noindent The script needs under a minute. The same check succeeds at noise $0.445$ with
$\beta=1/16$, $J=1000$ and $\eta=\gamma=10^{-5}$, and at $1117/2500=0.4468$ with $J=4000$. In
floating point the construction does not close at $0.447$ (with $\beta\in\{1/64,3/128\}$, $J$ up to
$12000$ and slacks $10^{-6}$), where the cycle of the ideal map is about to disappear
(Remark~\ref{rem:beyond}(ii)). The level $\eps_S$ is dyadic, which Theorem~\ref{thm:work} uses.

\subsection{Hardness up to noise \texorpdfstring{$0.4468$}{0.4468}}

Fix $\eps=\eps_S$ and the parameters of Proposition~\ref{prop:s24cert}. Fix $\rho$, $k_0$ and
$\tau_0$ as in Section~\ref{sec:latch}, with $\eta=10^{-6}$ there, and let $\Lambda_0$ be as there,
\[
M=\Bigl\lceil\frac{1+\ln(30(J+1))}{2\gamma^2}\Bigr\rceil,\qquad N_1=24h(1+\rho)(2J+1),
\]
and let $Z_0$ be the least integer with
$Z_0\ge\Lambda_0\bigl(24h+5+\tfrac43(\ln Z_0+4)\bigr)+N_1$.

\begin{definition}[$\mathrm{Prog}_a$]\label{def:proga}
For an integer $a\ge1$ and a circuit with $s\ge2$ gates, $\mathrm{Prog}_a(C)$ is the base
$\mathrm{Prog}(C)$ with the parameters $\beta,h,J,\rho$ fixed above and
\begin{align*}
m&=\text{the least power of two with }m\ge M(a+1)\log_2s,\\
Z&=Z_0(a+2)^2sm,\qquad b=\lceil\log_2Z\rceil,\qquad
\mathrm{pre}=\bigl\lceil\Lambda_0\bigl(24h+3+(a+3)\,b\ln2\bigr)\bigr\rceil .
\end{align*}
\end{definition}

\begin{lemma}[parameters]\label{lem:paramprog}
For $\mathrm{Prog}_a(C)$: (i) $n+1\le Z$; (ii)
$\mathrm{pre}\ge\Lambda_0\bigl(24h+3+(a+3)\ln(n+1)\bigr)$; (iii) there is $s^\ast_a$, depending on
$a$ only, such that $n+1\le s^2$ whenever $s\ge s^\ast_a$.
\end{lemma}

\begin{proof}
As for Lemma~\ref{lem:param4}, with $24h$ in place of $2h$; the chains and their padding occupy
$24h(1+\rho)m((J+1)s+2J)\le N_1sm$ positions, because $s\ge2$.
\end{proof}

\begin{corollary}[embedding up to noise $0.4468$]\label{cor:prog}
Let $a\ge1$ and let $C$ be a NOR circuit with $s\ge2$ gates. In the trace of
$S_{\eps_S}(\mathrm{Prog}_a(C))$,
\[
\PP[\text{some bundle is not good}]\ \le\ s^{-a}+(n+1)^{-(a+1)} .
\]
In particular, with at least the complementary probability, at least $44.26\%$ of the copies in
$R_J(g_s,8)$ hold $1$ if $\val(C)=1$, and at most $39.82\%$ if $\val(C)=0$.
\end{corollary}

\begin{proof}
As for Corollary~\ref{cor:latch}: every chain of $24h$ gates is preceded by $24h\rho$ zero gates,
so Lemma~\ref{lem:dens2any} applies with $c=24h$, and its condition on $\mathrm{pre}$ is
Lemma~\ref{lem:paramprog}(ii); so $\PP[E_\eta^c]\le(n+1)^{-(a+1)}$. By
Theorem~\ref{thm:progembed} the rest is at most $10(J+1)(s+2)e^{-2\gamma^2m}\le30(J+1)\,s\,e^{-2\gamma^2m}$,
which is at most $s^{-a}$ by the choice of $M$. The last sentence is
Proposition~\ref{prop:s24cert}.
\end{proof}

\begin{theorem}[hardness up to noise $0.4468$]\label{thm:hards24}
Theorem~\ref{thm:hard} holds for every $\eps\le\eps_S=3661/8192$: if for some $\eps\le\eps_S$ a
randomized NC family outputs the trace with probability at least $n^{-K}$ on the
$\eps$-smoothing of every base of every length $n\ge n_0$, then $\mathrm P\subseteq\mathrm{BPNC}$.
\end{theorem}

\begin{proof}
As for Theorem~\ref{thm:hard4}, with $\eps^*=\eps_S$, which is dyadic: by
Lemma~\ref{lem:errorless} and Lemma~\ref{gb:comp} we may assume that the family is errorless and
has success probability at least $n^{-K}$ at noise $\eps_S$. Repeat the proof of
Theorem~\ref{thm:hard2} at that level with $\mathrm{Prog}_a$ in place of $\mathrm{Tree}_a$, with
Corollary~\ref{cor:prog} and Lemma~\ref{lem:paramprog}(iii) in place of
Corollary~\ref{cor:threetwentieths} and Lemma~\ref{lem:param2}(iii), and with the three changes
listed in the proof of Theorem~\ref{thm:hard4}: the coins of a chain gate are a bounded number of
fair bits for the mixture, which is exact because $\beta$ is dyadic, and $\log_2m$ bits for the
copy; the noise is drawn with a computable threshold; and a trial votes $1$ if more than $42\%$ of
the copies in $R_J(g_s,8)$ hold $1$ in the trace it received, and $0$ otherwise. By
Corollary~\ref{cor:prog} the vote differs from $\val(C)$ with probability at most
$s^{-a}+(n+1)^{-(a+1)}$. With $a=2K+4$ the rest is unchanged.
\end{proof}

\begin{corollary}[the gap, with the lower end $0.4468$]\label{cor:gaps24}
Corollary~\ref{cor:gap} holds with ``$\eps\le\eps_S$'' in place of ``$\eps\le1/64$'' in part
(b). In particular $\inf E_K\in\{0\}\cup[\eps_S,3/4]$ for every $K$, and if
$\mathrm P\not\subseteq\mathrm{BPNC}$, the critical noise level satisfies
$0.4468<\eps_S\le\eps_c\le3/4$. Corollary~\ref{cor:upperend} lowers the upper end to $\eps_1<0.7288$.
\end{corollary}

\begin{proof}
As for Corollary~\ref{cor:gap2}, with Theorem~\ref{thm:hards24}.
\end{proof}

\begin{remark}[what this changes]\label{rem:s24}
\begin{enumerate}[label=(\roman*),leftmargin=2.4em]
\item \emph{The level $4/9$.} Below $4/9$ the latch embedding suffices and has simpler constants;
Theorem~\ref{thm:hards24} adds the noise levels in $[4/9,\eps_S]$. The level $4/9$ remains exact
for two other things: the analysis of the latch (Proposition~\ref{prop:adm}) and the
relevance-pruned iteration on random block bases over zero gates (Corollary~\ref{cor:exact}).
So between $4/9$ and $\eps_S$ that iteration is fast on random block bases, while no randomized
NC algorithm succeeds on the smoothing of every base unless $\mathrm P\subseteq\mathrm{BPNC}$. The
coincidence explained in Remark~\ref{rem:latch}(ii) belongs to the latch.
\item \emph{How far this route goes.} The cycle of the ideal map of $S_{24}$ disappears near
$0.447$ (Remark~\ref{rem:beyond}(ii)), and the floating-point version of the certificate fails at
$0.447$; the exact one succeeds at $\eps_S=0.44689\ldots$. By Theorem~\ref{thm:indep}, no design with independent reads has two attracting sets in
the ideal model from noise $0.455$ on, and the present proof uses designs with independent reads
whose gates read phases of their own parity. Corollary~\ref{cor:layered} makes this rigorous for
the bases built with long chains, many copies and a long preamble, which still prove
Theorem~\ref{thm:hards24} (Lemma~\ref{lem:longer}): from noise $0.455$ on, their smoothings are solved
by the relevance-pruned iteration in $O(\log n)$ rounds. So designs of this kind cannot lift the
lower end to $0.455$; better ones may exist between $\eps_S$ and $0.455$.
\item \emph{Constants.} A chain has $96$ gates, after $96\rho$ zero gates, and gives one copy to each of
five rails; a bundle has $m\approx6.6\cdot10^{12}(a+1)\log_2s$ chains (from $M$), there are $J+1=6001$ stages, and with
$\eta=10^{-6}$ the preamble is far larger than in Section~\ref{sec:latch}. Only the asymptotic
statement has content.
\item \emph{The gate stage is soft,} as for the latch (Remark~\ref{rem:latch}(iv)): without noise
the rails of a stage need not be close to the two states.
\item \emph{Simulation, far from the level only.} The parameters of the certificate are far beyond
what can be simulated. The script \texttt{s24\_embed.c} builds $\mathrm{Prog}(C)$ for a random
circuit with small parameters and decodes every gate from its rail of phase $8$ at stage $J$. At
noise $0.40$, with $h=1$, $J=32$, $\rho=20$, $m=1000$ and $\beta=1/4$, all $64$ gates of eight
random circuits with $8$ gates came out right; the decoded densities were at least $0.51$ for value
$1$ and at most $0.27$ for value $0$. With $m=300$, $\rho=10$ and $\beta=1/8$ one gate of $12$ was
wrong in two of three runs, each time a gate with one input of each value, whose gate stage lands
near the boundary between the two basins. At noise $0.43$ the sizes that fit in memory
($\rho=20$, $m=600$, $J=104$) fail: the density of ones is $0.251$ instead of $\varpi=0.245$ and the
sampling error of a rail is about $0.02$, while in the ideal model the construction needs
perturbations below $10^{-3}$ there. So the simulation tests the compiler, not the margins near
$\eps_S$. The smoothed simulations of layers of $S_{24}$ in Remark~\ref{rem:beyond}(ii) show the
oscillation persisting over $60$ layers at noise $0.445$.
\end{enumerate}
\end{remark}

\section{Why the upper end resists: adaptive adversaries}\label{sec:adaptive}

The upper end $3/4$ of the interval comes from the Kleene--prefix iteration (Theorem~\ref{kp:thm39});
Section~\ref{sec:allbases} lowers it to $0.7288$ with the relevance-pruned iteration. This
section records a barrier for improving it. In the smoothed model the adversary is
\emph{oblivious}: the base is fixed before the noise is drawn. Call a base \emph{adaptive} if the
pointer $j_k$ may depend on the noise $(\pi_i,\rho_i)_{i<k}$ of the earlier gates; the smoothing
then sets $\jt_k=j_k$ if $\pi_k=0$ and $\jt_k=\rho_k$ otherwise, as before. Against adaptive bases
both iterations fail at every noise level below $1$.

\begin{lemma}[depths]\label{fact:depths}
For every parent array, the Kleene--prefix iteration needs exactly $\max_id(i)$ rounds and its
relevance-pruned variant exactly $\max_iD(i)$ rounds, where $d$ and $D$ are the depths of
Section~\ref{sec:iterations}. The depth $d$ is given by $d(0)=0$ and
$d(i)=1+\max\{d(j_k):k\text{ not a NOT gate},\ \max(\varrho(i),1)\le k\le i\}$, so
$d(k)\ge1+d(j_k)$ whenever gate $k$ is not a NOT gate. Call gate $k$ \emph{relevant} if
$x_{k-1}=0$. If gate $k$ is relevant, not a reset and not a NOT gate, then
$D(k)=\max(D(k-1),1+D(j_k))$; if gate $k\ge2$ is not relevant and not a reset, then $D(k)=D(k-1)$.
\end{lemma}

\begin{proof}
Lemmas~\ref{kp:rounds-d}, \ref{kp:pruned}(a) and~\ref{kp:switch}(a), and the definition of $d$
in Section~\ref{sec:iterations}.
\end{proof}

\begin{proposition}[adaptive adversaries]\label{prop:adaptive}
Let $\eps\in(0,1)$ and $n\ge3$.
\begin{enumerate}[label=(\alph*),leftmargin=2.4em]
\item There is an adaptive base on whose $\eps$-smoothing the Kleene--prefix iteration needs at
least
\[
\frac{(1-\eps)\,\eps}{2+\eps}\,n-2\sqrt{n\ln n}-\ln n-3
\]
rounds, with probability at least $1-2/n^2$.
\item There is an adaptive base on whose $\eps$-smoothing the relevance-pruned iteration needs at
least $\frac{(1-\eps)(1-\eps/2)}{6}\,n-\sqrt{n\ln n}-1$ rounds, with probability at least $1-2/n^2$.
\end{enumerate}
\end{proposition}

\begin{proof}
In both parts $j_1=j_2=0$, and for $k\ge3$ the pointer $j_k$ is a position $p\le k-2$ with
$x_p=0$ of largest depth ($d$ in (a), $D$ in (b)). Such a position exists because $x_1=0$, and
the choice depends only on the gates below $k-1$. A gate that keeps this pointer reads a $0$ and
is not a NOT gate. The pair $(\pi_k,\rho_k)$ is independent of $\cF_{k-1}$, as in the oblivious
case.

(a) Let $M_k=\max\{d(p):p\le k-1,\ x_p=0\}$, which is nondecreasing in $k$. Suppose $\pi_k=0$ and
$x_{k-1}=1$. Then $x_k=0$, the zeros below $k$ are the zeros up to $k-2$, so $d(j_k)=M_k$, and
$d(k)\ge1+M_k$ by Lemma~\ref{fact:depths}. Hence $M_{k+1}\ge M_k+1$. So the number of rounds is at
least $N=|\{3\le k\le n:\pi_k=0,\ x_{k-1}=1\}|$.

Let $S$ be the number of ones among positions $1,\dots,n-1$. Ones are never adjacent, so
positions $0,\dots,i-1$ hold at least $(i-1)/2$ zeros, and a resampled gate after a zero that
reads a zero outputs a one. Hence
$a_i:=\PP[x_i=1\mid\cF_{i-1}]\ge\eps\,\mathbf 1[x_{i-1}=0]\,(\tfrac12-\tfrac1{2i})$. By the
Azuma--Hoeffding inequality, $\sum_{i<n}(x_i-a_i)\ge-\sqrt{n\ln n}$ except with probability
$n^{-2}$. At least $n-2-S$ of the positions $0,\dots,n-2$ hold a zero, so on this event
\[
S\ge\frac\eps2\,(n-2-S)-\frac\eps2\,(1+\ln n)-\sqrt{n\ln n},\qquad\text{that is,}\qquad
S\ge\frac{\eps\,(n-3-\ln n)-2\sqrt{n\ln n}}{2+\eps}.
\]
As $x_1=0$, exactly $S$ of the indices $3\le k\le n$ have $x_{k-1}=1$, and $\pi_k$ is independent
of $\cF_{k-1}$. A second application of the inequality gives $N\ge(1-\eps)S-\sqrt{n\ln n}$ except
with probability $n^{-2}$. Combining the two bounds gives the claim.

(b) Let $L_k=\max\{D(p):p\le k-2,\ x_p=0\}$, nondecreasing in $k$, and for $3\le k\le n-1$ let
$E_k$ be the event that $\pi_k=0$, $x_{k-1}=0$, and gate $k+1$ is not a reset. On $E_k$ gate $k$
reads a $0$ after a $0$, so $x_k=1$, and it is relevant, not a reset and not a NOT gate; by
Lemma~\ref{fact:depths}, $D(k)\ge1+D(j_k)=1+L_k$. Gate $k+1$ is then not relevant and not a reset, so
$x_{k+1}=0$ and $D(k+1)=D(k)$. Thus $L_{k+3}\ge1+L_k$ on $E_k$. Selecting events greedily at
distance at least $3$ shows that the number of rounds is at least a third of the number of $k$
with $E_k$.

Given $\cF_k$, $\pi_k=0$ and $x_{k-1}=0$, gate $k+1$ is not a reset if it keeps its pointer, and
if it is resampled it is a reset only if $\rho_{k+1}\ne k$ and $x_{\rho_{k+1}}=1$, which has
probability at most $1/2$ because at most $k/2$ of the positions $0,\dots,k-1$ hold a one. So
$\PP[E_k\mid\cF_{k-1}]\ge\mathbf 1[x_{k-1}=0]\,(1-\eps)(1-\eps/2)$. At least $n/2-2$ of the indices
$3\le k\le n-1$ have $x_{k-1}=0$. The indicator of $E_k$ is $\cF_{k+1}$-measurable, so the even and
the odd $k$ each give a martingale; by the Azuma--Hoeffding inequality the number of $k$ with $E_k$
is at least $(1-\eps)(1-\eps/2)(n/2-2)-2\sqrt{n\ln n}$ except with probability $2n^{-2}$.
\end{proof}

\begin{remark}[what this says about the upper end]\label{rem:upper}
\begin{enumerate}[label=(\roman*),leftmargin=2.4em]
\item For oblivious bases and $\eps>3/4$, both iterations finish in $O(\log n)$ rounds (Theorem~\ref{kp:thm39}).
Proposition~\ref{prop:adaptive} shows that this uses obliviousness: no argument that remains valid
when each pointer may depend on the earlier noise can prove a bound below linear, at any noise
level. In particular a chain count that grants the adversary the values its pointers read is false if
the grant is adaptive.
\item So a proof that lowers $3/4$ must use that every target is fixed before the noise is drawn.
The proposition does not show that the oblivious threshold of either iteration is $3/4$; it only
locates the difficulty.
\item In simulation the adaptive bases above give $\max_id(i)\approx0.069\,n$ and
$\max_iD(i)\approx0.037\,n$ at $\eps=0.9$, against $98$ rounds for the oblivious base $j_i=i-2$ at
$n=10^6$ (Section~\ref{sec:checks}).
\end{enumerate}
\end{remark}

\section{The upper end: a two-valued iteration, and where a proof below \texorpdfstring{$3/4$}{3/4} breaks}\label{sec:upper}

This section does not lower $3/4$; Section~\ref{sec:allbases} does, to $0.7288$, by a different
route. This section records a small rigorous fact, measurements on bases that are fixed before the
noise, and the step at which the first attempts made for this paper failed.

\begin{definition}[prefix--Jacobi iteration]\label{def:jacobi}
Let $j$ be a parent array of length $n$. For $y\in\{0,1\}^{n+1}$ with $y_0=1$ let $G(y)=z$, where
$z_0=1$ and, for $i\ge1$, $z_i=\NOR(z_{i-1},r_i)$ with $r_i=z_{i-1}$ if gate $i$ is a NOT gate and
$r_i=y_{j_i}$ otherwise. Put $y^{(0)}=(1,0,\dots,0)$ and $y^{(s+1)}=G(y^{(s)})$. The iteration
stops at the first $s$ with $y^{(s+1)}=y^{(s)}$ and outputs $y^{(s)}$.
\end{definition}

One application of $G$ is a prefix computation: mark gate $i$ if it is not a NOT gate and
$y_{j_i}=1$; then $z_i=0$ at a marked gate and $z_i=1-z_{i-1}$ at every other gate, so
$z_i=(i-m)\bmod 2$ if $m$ is the last marked gate up to $i$, and $z_i=(i+1)\bmod2$ if there is
none. This is in $\mathrm{AC}^0$, like a round of the iterations of Section~\ref{sec:iterations}. No value is ever
``unknown'': the iteration carries a full guess and repairs it.

For the next statement recall the depth $D$ of Section~\ref{sec:iterations}. Gate $k$ is relevant if
$x_{k-1}=0$. Let $\varrho'(i)$ be the largest $r\le i$ that is $0$ or a reset at a relevant gate.
The anchors of $i$ are the resets in $[\varrho'(i),i]$, together with $0$ if $\varrho'(i)=0$. Then
$D(0)=0$ and $D(i)=1+\min_r\max\bigl(\{D(j_r):r\ge1\}\cup\{D(j_k):r<k\le i,\ k\text{ relevant and
not a NOT gate}\}\bigr)$, the minimum over the anchors $r$ of $i$.

\begin{proposition}[the two-valued iteration is never slower]\label{prop:jacobi}
Let $j$ be a parent array with trace $x$.
\begin{enumerate}[label=(\alph*),leftmargin=2.4em]
\item If $G(y)=y$ then $y=x$.
\item If $D(i)\le t\le s$ then $y^{(s)}_i=x_i$.
\end{enumerate}
Hence $y^{(s)}=x$ for every $s\ge\max_iD(i)$: the iteration outputs the trace, is never wrong,
and applies $G$ at most $\max_iD(i)+1$ times. By Lemma~\ref{fact:depths} it reaches the trace no
later than the relevance-pruned iteration, hence no later than the Kleene--prefix iteration
(Lemma~\ref{kp:pruned}(b)), on every instance.
\end{proposition}

\begin{proof}
(a) If $G(y)=y$ then $y_0=1$ and $y_i=\NOR(y_{i-1},y_{j_i})$ for every $i\ge1$, also at NOT
gates, where $j_i=i-1$. The trace is the only solution of these equations.

(b) Induction on $t$. For $t=0$ only $i=0$ has $D(i)=0$, and $y^{(s)}_0=1=x_0$. Let $D(i)\le t+1\le
s$ and let $r$ be an anchor of $i$ attaining the minimum, so that $D(j_r)\le t$ if $r\ge1$, and
$D(j_k)\le t$ for every relevant gate $k\in(r,i]$ that is not a NOT gate. Write $z=y^{(s)}=G(y^{(s-1)})$;
by the induction hypothesis $y^{(s-1)}$ agrees with $x$ at all those positions $j_r$ and $j_k$.
We show $z_p=x_p$ for $r\le p\le i$. At $p=r$: if $r=0$ this is $z_0=1$; otherwise $r$ is a reset,
so it is not a NOT gate and reads $y^{(s-1)}_{j_r}=x_{j_r}=1$, whence $z_r=0=x_r$ whatever
$z_{r-1}$ is. Let $r<p\le i$ and $z_{p-1}=x_{p-1}$. If $x_{p-1}=1$ then $z_p=0=x_p$. If
$x_{p-1}=0$ then gate $p$ is relevant, and it is not a reset, because $(r,i]$ contains no reset at
a relevant gate. If $p$ is a NOT gate, $z_p=1-z_{p-1}=1=x_p$. Otherwise $x_{j_p}=0$ and
$D(j_p)\le t$, so gate $p$ reads $y^{(s-1)}_{j_p}=0$ and $z_p=1=x_p$.
\end{proof}

\noindent So every round bound of Section~\ref{sec:dense} for either iteration, in particular $O(\log n)$ rounds for
$\eps>3/4$, holds for the prefix--Jacobi iteration. Its errors obey a sharper rule than the
unknowns of the pruned iteration: an error that reaches a reset at a relevant gate always stops
there, whether or not that reset's own read is correct. The script \texttt{jacobi\_check.py}
checks (a) and (b) on all $46{,}233$ parent arrays with $n\le8$; the largest saving there is two
rounds. On the smoothed instances of Section~\ref{sec:checks} the saving is between $1$ and $43$
rounds. We do not know whether its threshold is lower.

\begin{lemma}[fresh reads]\label{lem:fresh}
Let $j$ be a base, $0<\eps\le1$, $x$ the trace of the $\eps$-smoothed instance, $\cF_g$ the
$\sigma$-field of the noise of gates $1,\dots,g$, and $\psi_g$ the fraction of ones among
$x_0,\dots,x_g$. For every $\tau\ge2$,
\[
\PP[x_\tau=0\mid\cF_{\tau-1}]\ \ge\ \eps\,\psi_{\tau-1},
\qquad
\PP[x_\tau=1\mid\cF_{\tau-2}]\ \ge\ \eps^2\,\psi_{\tau-2}\,\frac{\tau-1}{2\tau}.
\]
\end{lemma}

\begin{proof}
With probability $\eps$ gate $\tau$ is resampled, and then it reads a one with probability
$\psi_{\tau-1}$ and outputs $0$. For the second bound, first $\PP[x_{\tau-1}=0\mid\cF_{\tau-2}]\ge
\eps\psi_{\tau-2}$: if $x_{\tau-2}=1$ the probability is $1$, and otherwise apply the first bound
at $\tau-1$. Given $x_{\tau-1}=0$, gate $\tau$ is resampled with probability $\eps$ and then reads
a zero with probability at least $\lfloor\tau/2\rfloor/\tau\ge(\tau-1)/(2\tau)$, because ones are
never adjacent; in that case $x_\tau=1$.
\end{proof}

\begin{remark}[where a proof below $3/4$ breaks]\label{rem:needs}
Write $s=1-\eps$ and $r=1-\eps c$, with $c$ the density bound of Theorem~\ref{kp:thm34}, so that the chain count of
Theorem~\ref{kp:thm34} grows at rate $s/(1-r)$, up to the resampled pointers.
\begin{enumerate}[label=(\roman*),leftmargin=2.4em]
\item \emph{The slack.} That count charges probability $1$ to a gate of a window that keeps its
base pointer: such a gate is not a reset only if its target holds a zero, and the count does not
use this. Lemma~\ref{lem:fresh} says that no target is a zero for certain, and this is a
statement about bases fixed in advance; against the adaptive bases of
Proposition~\ref{prop:adaptive} it is false. If the target $\tau$ lies at least two positions above
everything the rest of the chain depends on, the lemma applies conditionally and the gate can be
charged $1-\eps^2c(\tau-1)/(2\tau)$. For bases in which the gates of a window read increasing
targets at least two apart this would lower the threshold slightly, to roughly $0.735$ if every
such gate could be charged in this way. We did not write that special case out: the gain is
small, and the general case is not reached.
\item \emph{The obstruction.} The base may aim the reads of a window at or below the positions
that the chain visits later. Then ``the target holds a zero'' is not fresh: it can be implied by
the event that the chain below exists, and one unlikely value can pay for many windows. Types do
not exclude this. Let $N_t(q)$ be the number of chains of $t$ base pointers from $q$, and
$A_t(q)=\EE[N_t(q);x_q=1]$, $B_t(q)=\EE[N_t(q);x_q=0]$. Conditioning on $\cF_{q-1}$ gives,
where the density of ones is at least $c$,
\[
A_t(q)\le r\,B_t(q-1)+s\,\sup_p B_{t-1}(p),\qquad
B_t(q)\le r\,A_t(q-1)+s\,\sup_p\EE N_{t-1}(p),
\]
and these inequalities are satisfied with growth rate exactly $s/(1-r)$ by the assignment in
which the even positions are deep only when they hold a one and the odd positions only when they
hold a zero. No instance behaves like that, since a resampled reset shifts the phase of the run
above it; but to exclude it one must show that the phase of a run is nearly independent of the
depth below it. We could not bound this correlation. Section~\ref{sec:allbases} goes around it: it
takes the gain from the gate of a step, whose relevance is decided by fresh noise, and not from
the reads of the window.
\item \emph{Measurements.} The script \texttt{theta.c} estimates the growth rate of the number of
chains of base pointers, the quantity a sharp first-moment proof would bound, on an active region
(the last tenth of $n=1.6\cdot10^7$ positions) inside filler $j_i=0$, which keeps the density of
ones near its minimum. The rate is $(m_{12}/m_4)^{1/8}$, where $m_t$ is the mean number of chains
of length $t$ over the upper half of the region; means of three smoothings.

\begin{center}
\begin{tabular}{llrrrrrr}
\toprule
base in the region & count for & $0.62$ & $0.64$ & $0.66$ & $0.68$ & $0.70$ & $0.74$\\
\midrule
distance in $[16,32)$ & $d$ & 1.10 & 1.04 & 0.97 & 0.92 & 0.86 & 0.74\\
 & $D_{\mathrm{last}}$ & 0.88 & 0.83 & 0.78 & 0.73 & 0.68 & 0.56\\
the same, even targets & $d$ & 1.20 & 1.12 & 1.04 & 0.95 & 0.89 & 0.76\\
 & $D_{\mathrm{last}}$ & 0.88 & 0.83 & 0.78 & 0.72 & 0.67 & 0.57\\
$j_i=i-3$ & $d$ & 1.01 & 0.97 & 0.91 & 0.86 & 0.82 & 0.70\\
 & $D_{\mathrm{last}}$ & 0.96 & 0.92 & 0.87 & 0.81 & 0.78 & 0.63\\
\bottomrule
\end{tabular}
\end{center}

\noindent On these bases the Kleene--prefix count stops growing between $0.62$ and $0.68$ and
the count for $D_{\mathrm{last}}$, the chains that step only through relevant gates and through the last reset, is below $1$ throughout. Two families of
designs with pinned resets, meant to make targets reliably zero, did worse. So these measurements
found no base fixed before the noise on which chains multiply at noise $0.70$. This is not evidence about the worst base:
the designs are ours, the rates come from three smoothings, and a count below $1$ on a few bases
proves nothing. Indeed Section~\ref{sec:barrier} gives a base fixed in advance on which the
expected number of chains provably grows at noise $0.70$ and $0.72$; on that base sampled rates
of this kind underestimate the exact ones.
\end{enumerate}
\end{remark}

\section{Below \texorpdfstring{$3/4$}{3/4} for fixed-distance bases}\label{sec:fixed}

For one natural family of bases the obstruction of Remark~\ref{rem:needs}(ii) does not arise, and
the fresh reads of Lemma~\ref{lem:fresh} can be charged. For an integer $H\ge1$ the
\emph{fixed-distance base} $j^H$ has $j_i=\max(i-H,0)$; for $H=1$ every gate of the base is a NOT gate. These are the bases $j_i=i-C$ of the experiments of
Section~\ref{sec:checks}. The result is about this family only; it says nothing about
the worst base, hence nothing about $\eps_c$. For fixed $H$, Theorem~\ref{thm:short} later gives
$O(\log^3n)$ rounds at every positive noise level by a different argument, with a constant that
grows like $\eps^{-2H}$; see Remark~\ref{rem:short}(ii). For large $H$ and the relevance-pruned
iteration, Subsection~\ref{sec:rows} gives $O(\log n)$ rounds from noise $0.455$ on, by counting
positions instead of chains.

Throughout this section $\psi_g=O_g/(g+1)$ is the fraction of ones among $x_0,\dots,x_g$.

\begin{definition}[admissible pairs]\label{def:admissible}
Let $0<\eps\le1$, and fix one base for each length $n$. A pair $(c,\beta)$ with $c>0$ and
$\beta\ge0$ is \emph{admissible} if there are constants $K_0,K_1$ and integers
$g_0=g_0(n)\in[1,K_0\ln(n+1)+K_1]$ such that for every $n\ge2$ there is an event $G$ with
$\PP[G^c]\le1/n$ and $G\subseteq\bigcap_{g_0<g\le n}G_g$, where
\[
G_g=\Bigl\{\frac{O_{g-1}-1}{g}\ge c\ \text{ and }\ \frac{g-O_{g-1}}{g}\ge\beta\Bigr\}\in\cF_{g-1}.
\]
\end{definition}

\begin{lemma}[every base]\label{fact:densall}
For every choice of bases, every pair with $0<c<\varpi(\eps)$ and $0\le\beta<1/2$ is admissible,
where $\varpi(\eps)$ is the root in $(0,1/2)$ of $\varpi=\eps(1-\varpi)^2$.
\end{lemma}

\begin{proof}
The first condition is the first paragraph of the proof of Theorem~\ref{kp:thm39}, which does
not use $\eps>3/4$: with $\tau=\min((\varpi-c)/2,1/4)$ and $K$ from Lemma~\ref{kp:limdens}(a), take
$g_0=\lceil K\ln(n+1)\rceil+\lceil1/\tau\rceil+\lceil1/(1-2\beta)\rceil$. The second condition
holds for every trace: ones are never adjacent, so $O_{g-1}\le(g+1)/2$ and
$(g-O_{g-1})/g\ge1/2-1/(2g)\ge\beta$ for $g>g_0$.
\end{proof}

For the fixed-distance bases both constants can be improved. Put
\[
\Lambda(\mu)=\frac{1-\eps\mu}{3-\eps-\eps\mu},\qquad U(\ell)=\frac{1-\eps\ell}{2-\eps\ell},
\]
and $\mu_0=1/2$, $\ell_{s+1}=\Lambda(\mu_s)$, $\mu_{s+1}=U(\ell_{s+1})$ for $s\ge0$.

\begin{lemma}[density of fixed-distance bases]\label{lem:fddens}
Let $0<\eps\le1$ and $H\ge1$. The sequence $(\ell_s)_{s\ge1}$ increases, $(\mu_s)_{s\ge0}$
decreases, all terms lie in $[0,1/2]$, and $\ell_1=1/3$. Let $s\ge1$ and $\tau\in(0,1/(8s+2)]$.
There is a constant $K$, depending only on $s$, $\tau$ and $H$, such that for every $n\ge2$, on
the $\eps$-smoothing of $j^H$, with probability at least $1-1/n$,
\[
\ell_s-(4s+1)\tau\ \le\ \psi_g\ \le\ \mu_s+(4s+1)\tau\qquad\text{for all }g\in[K\ln(n+1),n].
\]
More precisely, there is a constant $K'$, depending only on $s$ and $\tau$, such that this holds
for all $g\in[K'(H+\ln(n+1)),n]$.
\end{lemma}

\begin{proof}
This is the proof of Lemma~\ref{kp:twosided} with one change: the lower bound on the conditional means
uses the base. We give it in full.

\emph{The sequences.} For $\mu\le1$ and $\ell\le1$ the denominators $3-\eps-\eps\mu$ and
$2-\eps\ell$ are at least $1$, and $\Lambda'(\mu)=-\eps(2-\eps)/(3-\eps-\eps\mu)^2$,
$U'(\ell)=-\eps/(2-\eps\ell)^2$. So $\Lambda$ and $U$ are decreasing there, $|\Lambda'|\le1$ and
$|U'|\le1$, and $U\circ\Lambda$ is increasing. We have $\ell_1=\Lambda(1/2)=(1-\eps/2)/(3-3\eps/2)=1/3$
and $\mu_1=U(\ell_1)\le U(0)=\mu_0$. Induction gives $\mu_{s+1}\le\mu_s$ and then
$\ell_{s+2}=\Lambda(\mu_{s+1})\ge\Lambda(\mu_s)=\ell_{s+1}$. Since $0\le\mu_s\le1/2$, also
$1/3\le\ell_{s+1}\le\Lambda(0)\le1/2$, and $0\le U(\ell)\le1/2$ for $0\le\ell\le1$.

\emph{Conditional means.} Let $a_i=\PP[x_i=1\mid\cF_{i-1}]$ and $z_{i-1}=1-x_{i-1}$. If $x_{i-1}=1$
then $x_i=0$; otherwise $x_i=1$ exactly when gate $i$ reads a zero. Let $i\ge H$. With probability
$\eps$ the pointer is uniform on $\{0,\dots,i-1\}$, where a fraction $1-\psi_{i-1}$ of the
positions hold zeros, and otherwise it is $i-H$. Hence
$a_i=z_{i-1}\bigl(\eps(1-\psi_{i-1})+(1-\eps)(1-x_{i-H})\bigr)$, and since
$z_{i-1}(1-x_{i-H})\ge z_{i-1}-x_{i-H}$,
\[
\eps z_{i-1}(1-\psi_{i-1})+(1-\eps)(z_{i-1}-x_{i-H})\ \le\ a_i\ \le\ z_{i-1}(1-\eps\psi_{i-1}).
\]

\emph{The good event.} Let $G^*$ be the event that
$\bigl|\sum_{i=a+1}^{b}(x_i-a_i)\bigr|\le\sqrt{2(b-a)\ln(n+1)}$ for all $0\le a<b\le n$. The
summands are martingale differences, and $x_i-a_i$ lies in an interval of length $1$ that is
$\cF_{i-1}$-measurable. By the Azuma--Hoeffding inequality in that form, each pair $(a,b)$ fails
with probability at most $2(n+1)^{-4}$, so $\PP[(G^*)^c]\le(n+1)^{-2}\le1/n$.

\emph{Intervals.} Work on $G^*$. Let $H\le a<b\le n$, $m=b-a$, $X=O_b-O_a$ and
$\Delta(m)=H+1+\sqrt{2m\ln(n+1)}$. Positions $a,\dots,b-1$ hold between $m-X-1$ and $m-X+1$ zeros,
and positions $a+1-H,\dots,b-H$ hold at most $O_b-O_{a-H}\le X+H$ ones. If $\psi_{i-1}\le\mu$ for
all $i\in(a,b]$, where $0\le\mu\le1$, summing the lower bound on $a_i$ gives
\[
X\ \ge\ \eps(1-\mu)(m-X-1)+(1-\eps)\bigl((m-X-1)-(X+H)\bigr)-\sqrt{2m\ln(n+1)},
\]
that is, $(3-\eps-\eps\mu)X\ge(1-\eps\mu)m-(H+1)-\sqrt{2m\ln(n+1)}$, hence
$X\ge\Lambda(\mu)m-\Delta(m)$. If $\psi_{i-1}\ge\ell$ for all $i\in(a,b]$, where $-1\le\ell\le1$,
the upper bound gives $X\le(1-\eps\ell)(m-X+1)+\sqrt{2m\ln(n+1)}$, hence
$X\le U(\ell)m+\Delta(m)$.

\emph{From intervals to prefixes.} For an integer $Q\ge H$ let
$Q'=\lceil(Q+H+1)/\tau\rceil+\lceil8\ln(n+1)/\tau^2\rceil$. For $g\ge Q'$ we have
$(Q+H+1)/(g+1)\le\tau$ and, since $\ln(n+1)>1$,
$(1+\sqrt{2g\ln(n+1)})/(g+1)\le\tau/8+\tau/2<\tau$. Suppose $\psi_p\le\mu$ for all $p\in[Q,n]$,
with $0\le\mu\le1$. The interval bound on $(Q,g]$ gives $O_g\ge\Lambda(\mu)(g-Q)-\Delta(g)$, and
$\Lambda(\mu)\le1$, so $\psi_g\ge\Lambda(\mu)-2\tau$ for $g\ge Q'$. Suppose instead $\psi_p\ge\ell$
for all $p\in[Q,n]$, with $-1\le\ell\le1$. Since $O_Q\le(Q+2)/2\le Q+1$, we get
$O_g\le Q+1+U(\ell)(g+1)+\Delta(g)$, so $\psi_g\le U(\ell)+2\tau$ for $g\ge Q'$.

\emph{Stages.} Let $Q_0=\max(H,\lceil1/(2\tau)\rceil)$ and $Q_{r+1}=(Q_r')'$. Since
$O_g\le(g+2)/2$, we have $\psi_g\le\mu_0+\tau$ for $g\ge Q_0$. Let $r<s$ and suppose
$\psi\le\mu_r+e$ on $[Q_r,n]$ with $e=(4r+1)\tau$. Since $\tau\le1/(8s+2)$ we have
$e+2\tau\le1/2$, so $\mu_r+e\le1$ and $\ell_{r+1}-e-2\tau\ge-1$, and the Lipschitz bounds apply:
$\psi\ge\Lambda(\mu_r+e)-2\tau\ge\ell_{r+1}-e-2\tau$ on $[Q_r',n]$, and then
$\psi\le U(\ell_{r+1}-e-2\tau)+2\tau\le\mu_{r+1}+e+4\tau$ on $[Q_{r+1},n]$. Starting from $e=\tau$
at $r=0$, this gives both bounds with error at most $(4s+1)\tau$ on $[Q_s,n]$. Finally put
$L_n=H+\ln(n+1)\ge2$. If $Q\le kL_n$ then $Q'\le((k+1)/\tau+8/\tau^2+1)L_n$, because
$H+1\le L_n$ and $2\le L_n$; and $Q_0\le(2+1/(2\tau))L_n$. So $Q_s\le K'L_n$ with $K'$ depending
only on $s$ and $\tau$, and $K'L_n\le K'(H+1)\ln(n+1)$, which gives the first statement with
$K=K'(H+1)$.
\end{proof}

\begin{corollary}[admissible pairs for $j^H$]\label{cor:fdadm}
Let $\ell^*=\lim_s\ell_s$ and $\mu^*=\lim_s\mu_s$. For the bases $j^H$, every pair with
$0<c<\ell^*$ and $0\le\beta<1-\mu^*$ is admissible.
\end{corollary}

\begin{proof}
Choose $s$ with $\ell_s>c$ and $\mu_s<1-\beta$, and put $\eta=\min(\ell_s-c,\,1-\beta-\mu_s)/2$
and $\tau=\min(\eta/(4s+1),\,1/(8s+2))$. Let $K$ and the event $G$ be those of
Lemma~\ref{lem:fddens}, and $g_0=\lceil K\ln(n+1)\rceil+\lceil1/\eta\rceil$. Let $g_0<g\le n$.
On $G$, since $\psi_{g-1}=O_{g-1}/g$, we get $(O_{g-1}-1)/g\ge\ell_s-\eta-1/g\ge c$ and
$(g-O_{g-1})/g=1-\psi_{g-1}\ge1-\mu_s-\eta\ge\beta$.
\end{proof}

\noindent At $\eps=0.7$, for example, $\varpi=0.322$, while $\ell^*=0.3496$ and
$1-\mu^*=0.5697$. For $\eps=1$ both sequences converge to $(3-\sqrt5)/2$, the density of uniform
instances (Section~\ref{sec:newton}).

Fix $\eps<1$ and a pair $(c,\beta)$ that is admissible for the bases $j^H$, with $g_0$ and $G$ as
in Definition~\ref{def:admissible}, and write
\[
r=1-\eps c,\qquad e_0=\eps(1-c),\qquad\text{so that } r=e_0+(1-\eps).
\]
The two conditions in $G_g$ give $c+\beta\le(g-1)/g<1$, so $\eps\beta<e_0<r$. The first one
also gives $c<1/2$, because $O_{g-1}\le(g+1)/2$.

\begin{lemma}[one gate]\label{lem:onegate}
Let $g_0<g\le n$. On $G_g$:
\begin{enumerate}[label=(\roman*),leftmargin=2.4em]
\item $\PP[\pi_g=1,\ g\text{ is not a reset}\mid\cF_{g-1}]\le e_0$;
\item if $x_{g-1}=1$ then $x_g=0$, and if $x_{g-1}=0$ then
$\eps\beta\le\PP[x_g=1\mid\cF_{g-1}]\le r$.
\end{enumerate}
\end{lemma}

\begin{proof}
(i) If $\pi_g=1$, $\rho_g\le g-2$ and $x_{\rho_g}=1$, then $g$ is a reset; given $\cF_{g-1}$ this has
probability $\eps(O_{g-1}-x_{g-1})/g\ge\eps c$. (ii) A reset outputs $0$, which gives the upper
bound. If $\pi_g=1$ and $x_{\rho_g}=0$, then $x_g=1-x_{g-1}$; given $\cF_{g-1}$ this has probability
$\eps(g-O_{g-1})/g\ge\eps\beta$.
\end{proof}

\begin{definition}[the weights $T_m$]\label{def:Tm}
Let $H\ge2$ and $m\ge0$, and put $q=\min(m,H-1)$. Let $t_q(0)=t_q(1)=r^{m-q}$ and, for
$i=q-1,\dots,0$,
\[
t_i(1)=r\,t_{i+1}(0),\qquad
t_i(0)=\max_{\alpha\in\{\eps\beta,\;r\}}\bigl(\alpha\,e_0\,t_{i+1}(1)+(1-\alpha)\,r\,t_{i+1}(0)\bigr).
\]
Put $T_m=\max(t_0(0),t_0(1))$ and $\Theta_H=\Theta_H(\eps;c,\beta)=(1-\eps)\sum_{m\ge0}T_m$.
\end{definition}

By induction $t_i(v)\le r^{m-i}$, so $T_m\le r^m$ and $\Theta_H\le(1-\eps)/(\eps c)$, the growth
factor of Theorem~\ref{kp:thm34}. The recursion is monotone in $t_{i+1}$, because $0\le\alpha\le1$; so the same
induction shows that $T_m$ does not increase when $H$ is replaced by $H+1$, and $\Theta_H$ is nonincreasing in $H$. For $H=2$ there is no gain: $t_0(1)=r^m$.

\begin{lemma}[one step]\label{lem:onestep}
Let $\jt=S_\eps(j^H)$. Let $P\ge g_0$, $k=P+H$, $m\ge0$ and $k+m\le n$. Let $C$ be the event that
$\pi_k=0$ and no gate in $(k,k+m]$ is a reset, and let $G'=\bigcap_{P<g\le k+m}G_g$. Then
$\PP[C\cap G'\mid\cF_P]\le(1-\eps)\,T_m$.
\end{lemma}

\begin{proof}
Let $q=\min(m,H-1)$ and, for $v\in\{0,1\}$, $\varphi(v)=e_0+(1-\eps)(1-v)$; so $\varphi(0)=r$ and
$\varphi(1)=e_0$. For $1\le i\le q$ the base pointer of gate $k+i$ is $P+i$, which lies in $(P,k)$.
If that gate keeps its pointer and is not a reset, then $x_{P+i}=0$. Since $\pi_{k+i}$ is independent
of $\cF_{k+i-1}$ and $x_{P+i}$ is $\cF_{k-1}$-measurable, Lemma~\ref{lem:onegate}(i) gives, on
$G_{k+i}$,
\[
\PP[k+i\text{ is not a reset}\mid\cF_{k+i-1}]\le e_0+(1-\eps)\mathbf 1[x_{P+i}=0]=\varphi(x_{P+i})\quad(i\le q),
\]
and the bound $r$ for $q<i\le m$. Also $\PP[\pi_k=0\mid\cF_{k-1}]=1-\eps$. Conditioning
successively on $\cF_{k+m-1},\dots,\cF_{k-1}$, as in the proof of Lemma~\ref{lem:onechain},
\[
\PP[C\cap G'\mid\cF_{k-1}]\le(1-\eps)\,r^{m-q}\prod_{i=1}^{q}\varphi(x_{P+i})\cdot\mathbf 1_{G''},
\qquad G''=\bigcap_{P<g<k}G_g .
\]
For $0\le i\le q$ let $U_i=\EE\bigl[r^{m-q}\prod_{i<i'\le q}\varphi(x_{P+i'})\,\mathbf 1_{G''}\mid\cF_{P+i}\bigr]$.
We show $U_i\le t_i(x_{P+i})$ by induction from $i=q$ down; this is clear for $i=q$. For $i<q$,
$U_i\le\EE[\varphi(x_{P+i+1})\,t_{i+1}(x_{P+i+1})\mathbf 1_{G_{P+i+1}}\mid\cF_{P+i}]$, and
$G_{P+i+1}\in\cF_{P+i}$. On it, by Lemma~\ref{lem:onegate}(ii): if $x_{P+i}=1$ then $x_{P+i+1}=0$
and the expectation is $r\,t_{i+1}(0)=t_i(1)$; if $x_{P+i}=0$ then $x_{P+i+1}=1$ with a conditional
probability $\alpha\in[\eps\beta,r]$, and the expectation is
$\alpha e_0t_{i+1}(1)+(1-\alpha)r\,t_{i+1}(0)\le t_i(0)$, since this is linear in $\alpha$. Finally
$\PP[C\cap G'\mid\cF_P]\le(1-\eps)U_0\le(1-\eps)T_m$.
\end{proof}

A \emph{typed skeleton} is a skeleton $\sigma=(P_0,k_1,P_1,\dots,k_\ell,P_\ell)$ as in
Section~\ref{sec:logspace} together with a type for each step, \emph{base} or \emph{resampled},
such that $P_s=k_s-H$ at every base step. Let $E^*_\sigma$ be the event that, for every $s$, no
gate in $(k_s,P_{s-1}]$ is a reset, and $\pi_{k_s}=0$ if step $s$ is a base step, while
$\pi_{k_s}=1$ and $\rho_{k_s}=P_s$ if it is a resampled step.

\begin{lemma}[one chain, typed]\label{lem:typed}
For $\jt=S_\eps(j^H)$ and every typed skeleton,
\[
\PP[E^*_\sigma\cap G]\le\prod_{s\text{ resampled}}r^{P_{s-1}-k_s}\,\frac{\eps}{k_s}\;
\prod_{s\text{ base}}(1-\eps)\,T_{P_{s-1}-k_s}.
\]
\end{lemma}

\begin{proof}
Since $G\subseteq G^\circ:=\bigcap_{P_\ell<g\le P_0}G_g$, it suffices to bound
$\PP[E^*_\sigma\cap G^\circ]$. The condition on step $s$ is $\cF_{P_{s-1}}$-measurable: it involves
the noise of the gates in $(P_s,P_{s-1}]$ and the trace below them. Induction on $\ell$: the
conditions on steps $2,\dots,\ell$ and the events
$G_g$ with $g\le P_1$ are $\cF_{P_1}$-measurable, and the conditional probability of the condition
on step $1$, intersected with $\bigcap_{P_1<g\le P_0}G_g$, given $\cF_{P_1}$, is at most
$(1-\eps)T_{P_0-k_1}$ at a base step (Lemma~\ref{lem:onestep} with $P=P_1$) and at most
$r^{P_0-k_1}\eps/k_1$ at a resampled step (as in Lemma~\ref{lem:onechain}; the gates in $(P_1,k_1)$
are not constrained).
\end{proof}

\begin{theorem}[fixed-distance bases]\label{thm:fixed}
Let $H\ge2$ and $0<\eps<1$, and suppose that $\Theta_H(\eps;c,\beta)<1$ for some pair $(c,\beta)$
that is admissible for the bases $j^H$. Then there are constants $A,B$ such that for every $n\ge2$, with probability at
least $1-2/n$, the Kleene--prefix iteration on $S_\eps(j^H)$ finishes within $A\ln(n+1)+B$ rounds.
The same holds for the relevance-pruned iteration and, with one more round, for the
prefix--Jacobi iteration. With the pairs of Corollary~\ref{cor:fdadm} the hypothesis holds
\[
\text{for }H\ge3\text{ at }\eps=0.73,\qquad\text{for }H\ge8\text{ at }\eps=0.70,\qquad
\text{for }H\ge16\text{ at }\eps=0.69 ,
\]
and already with the pairs of Lemma~\ref{fact:densall}, which are admissible for every base, for
$H\ge3$ at $\eps=0.74$, for $H\ge8$ at $\eps=0.72$ and for $H\ge16$ at $\eps=0.71$.
\end{theorem}

\begin{proof}
Choose $\lambda>0$ with $\theta_H:=1/((\lambda+1)c)+\Theta_H<1$, let $h(q)=(q+1)^\lambda$, and
define, for $f\ge0$,
\[
(W_Hf)(P)=\sum_{k=g_0+1}^{P}\Bigl(r^{P-k}\,\frac{\eps}{k}\sum_{q=g_0}^{k-2}f(q)
+(1-\eps)\,T_{P-k}\,f(k-H)\,\mathbf 1[k-H\ge g_0]\Bigr).
\]
As in Lemma~\ref{lem:kernel}, $\sum_{q\le k-2}h(q)\le k^{\lambda+1}/(\lambda+1)$ and
$h(k-H)\le(P+1)^\lambda$, so
$(W_Hh)(P)\le(P+1)^\lambda\bigl(\eps/((\lambda+1)\eps c)+(1-\eps)\sum_mT_m\bigr)=\theta_Hh(P)$.
Every dependency chain of $\jt$ with positions at least $g_0$ satisfies $E^*_\sigma$ for one
typing: if $\jt_{k_s}=P_s$ and $\pi_{k_s}=0$, then $P_s=j_{k_s}=k_s-H$, because $P_s\ge g_0\ge1$.
By Lemma~\ref{lem:typed} the expected number of dependency chains of length $\ell$ from $P_0$ that
occur together with $G$ is at most $(W_H^\ell\mathbf 1)(P_0)\le\theta_H^\ell(n+1)^\lambda$. With
$\ell^*=\lceil((\lambda+1)\ln(n+1)+\ln n)/\ln(1/\theta_H)\rceil$ and a union bound over $P_0$, the
event $\Gamma$ that no dependency chain of length $\ell^*$ has all its positions at least $g_0$
fails together with $G$ with probability at most $1/(n+1)$. On $\Gamma$ the Kleene--prefix
iteration finishes within $\ell^*+g_0-1$ rounds (proof of Theorem~\ref{kp:thm34}; this step is
deterministic). The other two iterations are covered by Lemma~\ref{kp:pruned}(b) and
Proposition~\ref{prop:jacobi}.

The numerical claims: $\Theta_H$ is nonincreasing in $H$. With $c=\varpi(\eps)-10^{-4}$
and $\beta=0.4999$ the script \texttt{fd\_bound.py} gives $\Theta_3(0.74)\le0.985$,
$\Theta_8(0.72)\le0.972$ and $\Theta_{16}(0.71)\le0.994$. With $c=\ell_{40}-10^{-4}$ and
$\beta=1-\mu_{40}-10^{-4}$, which are admissible by Corollary~\ref{cor:fdadm} because
$\ell_{40}\le\ell^*$ and $\mu_{40}\ge\mu^*$, it gives $\Theta_3(0.73)\le0.964$,
$\Theta_8(0.70)\le0.977$ and $\Theta_{16}(0.69)\le0.995$. All six are computed in double
precision (Section~\ref{sec:checks}).
\end{proof}

\noindent Values of $\Theta_H(\eps;\ell^*,1-\mu^*)$, the constants of
Corollary~\ref{cor:fdadm}, and the noise level at which it equals $1$:

\begin{center}
\begin{tabular}{rrrrrrr}
\toprule
$H$ & $\eps=0.69$ & $0.70$ & $0.71$ & $0.72$ & $0.73$ & $\Theta_H=1$ at\\
\midrule
2 & 1.287 & 1.226 & 1.166 & 1.108 & 1.052 & 0.7395\\
3 & 1.169 & 1.115 & 1.063 & 1.013 & 0.964 & 0.7225\\
4 & 1.114 & 1.065 & 1.017 & 0.971 & 0.925 & 0.7137\\
6 & 1.046 & 1.003 & 0.960 & 0.918 & 0.877 & 0.7006\\
8 & 1.017 & 0.976 & 0.935 & 0.896 & 0.857 & 0.6941\\
16 & 0.995 & 0.956 & 0.917 & 0.880 & 0.842 & 0.6887\\
64 & 0.994 & 0.955 & 0.917 & 0.879 & 0.842 & 0.6884\\
\bottomrule
\end{tabular}
\end{center}

\noindent The same with the constants of Lemma~\ref{fact:densall}, which hold for every base,
$\Theta_H(\eps;\varpi(\eps),1/2)$:

\begin{center}
\begin{tabular}{rrrrrr}
\toprule
$H$ & $\eps=0.71$ & $0.72$ & $0.73$ & $0.74$ & $\Theta_H=1$ at\\
\midrule
2 & 1.260 & 1.191 & 1.125 & 1.061 & 0.7500\\
3 & 1.160 & 1.099 & 1.041 & 0.984 & 0.7371\\
4 & 1.110 & 1.054 & 1.000 & 0.947 & 0.7300\\
6 & 1.047 & 0.997 & 0.948 & 0.901 & 0.7194\\
8 & 1.018 & 0.971 & 0.925 & 0.880 & 0.7138\\
16 & 0.993 & 0.949 & 0.905 & 0.862 & 0.7084\\
64 & 0.992 & 0.947 & 0.904 & 0.862 & 0.7081\\
\bottomrule
\end{tabular}
\end{center}

\begin{remark}[scope]\label{rem:fixedscope}
\begin{enumerate}[label=(\roman*),leftmargin=2.4em]
\item Since $(1-\eps)/(\eps\varpi(\eps))>1$ for $\eps<3/4$, the count of Theorem~\ref{kp:thm34} does not give these
cases. Part of the gain is the density of Lemma~\ref{lem:fddens}, which is special to these
bases; for $H=2$, where the weights $T_m$ gain nothing, it alone moves the point at which the
factor reaches $1$ from $3/4$ to $0.7395$. The rest comes
from two facts about a base fixed in advance: a gate above $k_s$ that keeps
its pointer must read a zero, and the positions it can read are fresh. Two neighbouring positions
are both zero only if the upper one is a reset, which is why the gain is larger than the estimate
of Remark~\ref{rem:needs}(i).
\item What the proof uses about $j^H$ is that the gates of a window above $k_s$ read distinct
positions strictly between $P_s$ and $k_s$, so that nothing else in the chain depends on the
noise of those positions. Section~\ref{sec:ordered} extends the argument to every base whose non-NOT pointers do not
decrease with the gate, with a bound over the positions between the reads in place of
Definition~\ref{def:Tm}. For a base that aims those reads strictly
below $P_s$ the argument gives nothing, for the reason in Remark~\ref{rem:needs}(ii).
\item The bound is still far from the behaviour of these bases. Near $\eps=0.7$
Lemma~\ref{lem:fddens} gives a density of about $0.35$, while the measured density of ones on
$S_\eps(j^H)$ is about $0.39$ for $H\ge3$. (For $H=2$ and small $\eps$ the lower bound $1/3$ is
nearly attained: the trace of $j^2$ has period $3$.) At $\eps=0.6$ the measured number of base
gates in a window, averaged over the value at $k-H$, is $0.92$ for $H=3$ and $0.98$ for $H=8$ (it
is above $1$ when that value is a one), while $\Theta_3=1.73$ and $\Theta_8=1.43$ there
(Section~\ref{sec:checks}).
\item The theorem is qualitative, like Theorem~\ref{kp:thm39}. The constant $K$ of
Lemma~\ref{lem:fddens} grows like $\tau^{-(2s+1)}$, and that of Lemma~\ref{kp:limdens}(a) like
$\tau^{-7}$. For the pairs used in the numerical claims, which stay within $10^{-4}$ of the
limiting constants, $g_0$ is of order $10^{40}\ln n$, so the bound on the number of rounds says
nothing at any size that can be run. The measurements of Section~\ref{sec:checks} are the
evidence at such sizes.
\end{enumerate}
\end{remark}

\subsection{Bases with nondecreasing pointers}\label{sec:ordered}

The argument extends to every base whose non-NOT pointers do not decrease with the gate. The
positions read by the gates of a window are then again fresh, or equal to the next target of the
chain. The distances between them are arbitrary, and the weights $T_m$ are replaced by a bound
that holds for every choice of distances. Several gates may read the same position; of these,
only the highest needs to be counted.

\begin{definition}[ordered bases]\label{def:ordered}
Call gate $i$ of a base \emph{active} if it is not a NOT gate and $j_i\ge1$. Let $D\ge2$ and
$N\ge0$. A base is \emph{$(D,N)$-ordered} if every active gate $g>N$ satisfies $j_g\le g-D$ and
$j_k\le j_g$ for every active gate $k<g$.
\end{definition}

\noindent The base $j^H$ is $(H,H)$-ordered. So is the base with $j_i=i-H$ on an arbitrary set of
gates above $H$ and $j_i=i-1$ elsewhere. The base $j_i=\lfloor i/2\rfloor$ is $(D,2D)$-ordered
for every $D$. A base in which all gates of a block read one position, at least $D$ below the
block, with these positions increasing from block to block, is $(D,0)$-ordered.

Fix $\eps\le3/4$ and a pair $(c,\beta)$ with $\beta\le2/3$ that is admissible for the bases
considered, with $g_0$ and $G$; enlarging $g_0$ by a constant, assume $g_0\ge N$. Keep
$r=1-\eps c$ and $e_0=\eps(1-c)$, and recall $r-e_0=1-\eps$. For $w=(w(0),w(1))$ with nonnegative
entries let
\[
(Mw)(1)=w(0),\qquad (Mw)(0)=\max_{\alpha\in\{\eps\beta,\,r\}}\bigl(\alpha\,w(1)+(1-\alpha)\,w(0)\bigr),
\qquad(\varphi w)(v)=\varphi(v)\,w(v),
\]
with $\varphi(0)=r$ and $\varphi(1)=e_0$ as before, and $B_d=M^d\circ\varphi$ for $d\ge1$. These
maps are monotone and positively homogeneous, $Mw\le\|w\|_\infty\mathbf 1$ and
$B_dw\le r\|w\|_\infty\mathbf 1$. Put
\[
\gamma=\eps^2\beta c,\qquad s(y)=r(1-\gamma)\,y(0)+\gamma e_0\,y(1),\qquad
s_0=s(\mathbf 1)=r-\gamma(1-\eps),\qquad S_\mu=\sum_{a=0}^{\mu-1}e_0^{\,a}.
\]

\begin{lemma}[all distances]\label{lem:spacing}
Let $y\ge0$ with $\Delta:=r\,y(0)-e_0\,y(1)\ge0$. Then $B_1y=\bigl(r\,y(0)-\eps\beta\Delta,\ r\,y(0)\bigr)$, and
$B_dy\le s(y)\mathbf 1$ for every $d\ge2$.
\end{lemma}

\begin{proof}
Let $w=\varphi y$, so $w(0)-w(1)=\Delta\ge0$. We claim $M^dw=(w(0)-c_d\Delta,\ w(0)-c_{d-1}\Delta)$
for $d\ge1$, where $c_0=0$, $c_1=\eps\beta$ and
$c_{d+1}=\min_{\alpha\in\{\eps\beta,r\}}(\alpha c_{d-1}+(1-\alpha)c_d)$. For $d=1$ this is the
definition of $M$, since $w(1)\le w(0)$. If it holds for $d$, then $(M^{d+1}w)(1)=w(0)-c_d\Delta$
and $(M^{d+1}w)(0)=\max_\alpha\bigl(w(0)-(\alpha c_{d-1}+(1-\alpha)c_d)\Delta\bigr)=w(0)-c_{d+1}\Delta$.
Now $c_2=(1-r)\eps\beta=\gamma$, and for $d\ge2$ the number $c_{d+1}$ is a convex combination of
$c_{d-1}$ and $c_d$; so $\gamma\le c_d\le\eps\beta$ for all $d\ge1$. Hence, for $d\ge2$, both
entries of $M^dw$ are at most $w(0)-\gamma\Delta=s(y)$. The case $d=1$ is the first line.
\end{proof}

\begin{lemma}[sum over distances]\label{lem:sumspacing}
Let $d_u\ge1$ and $\mu_u\ge1$ be integers and $\rho_u\in[0,1]$, for $u\ge1$. Let $y_1=\mathbf 1$ and
$y_{u+1}=\rho_u\,r^{\mu_u-1}B_{d_u}y_u$. Then
\[
\sum_{u\ge1}S_{\mu_u}\|y_u\|_\infty\ \le\ \frac1{1-s_0}=\frac1{\eps c+\gamma(1-\eps)} .
\]
\end{lemma}

\begin{proof}
Let $A_1y=(r\,y(0)-\eps\beta(r\,y(0)-e_0y(1)),\ r\,y(0))$, a linear map with nonnegative entries.
Define $\hat y_1=\mathbf 1$ and $\hat y_{u+1}=\rho_ur^{\mu_u-1}A_1\hat y_u$ if $d_u=1$,
$\hat y_{u+1}=\rho_ur^{\mu_u-1}s(\hat y_u)\mathbf 1$ if $d_u\ge2$.

\emph{The cone.} Let $\varrho_1=r/(r-\eps\beta(1-\eps))$ and let $\mathcal K$ be the set of $y$ with
$y(0)>0$ and $1\le y(1)/y(0)\le\varrho_1$. Since $\eps\beta\le1$ we have
$r-\eps\beta(1-\eps)\ge r-(1-\eps)=e_0$, so $\varrho_1\le r/e_0$, and every $y\in\mathcal K$ has
$r\,y(0)\ge e_0y(1)$. Positive multiples of $\mathbf 1$ lie in $\mathcal K$. If $y\in\mathcal K$
has ratio $\varrho=y(1)/y(0)$, then $A_1y$ has positive first entry and ratio
$F(\varrho)=r/(r(1-\eps\beta)+\eps\beta e_0\varrho)$. The function $F$ is nonincreasing,
$F(1)=\varrho_1$, and $F(\varrho_1)\ge1$ because $e_0\varrho_1\le r$. So $A_1\mathcal K\subseteq\mathcal K$,
and by induction $\hat y_u\in\mathcal K$ or $\hat y_u=0$ for every $u$.

\emph{Domination.} $y_u\le\hat y_u$ for every $u$: if $y_u\le\hat y_u$, then
$B_{d_u}y_u\le B_{d_u}\hat y_u$ by monotonicity, and by Lemma~\ref{lem:spacing} the latter
equals $A_1\hat y_u$ if $d_u=1$ and is at most $s(\hat y_u)\mathbf 1$ if $d_u\ge2$.

\emph{A potential.} Let $\Phi(y)=\|y\|_\infty+s(y)/(1-s_0)$; then $\Phi(\mathbf 1)=1/(1-s_0)$, and
$\Phi$ is positively homogeneous. Three facts, for $y\in\mathcal K$.

(a) $\Phi(y')\le\Phi(y)-\|y\|_\infty$ for $y'=s(y)\mathbf 1$ and for $y'=A_1y$. In the first case
$\Phi(y')=s(y)\Phi(\mathbf 1)=s(y)/(1-s_0)$. In the second, write $\Delta=r\,y(0)-e_0y(1)\ge0$; then
$\|A_1y\|_\infty=r\,y(0)$, $s(y)=r\,y(0)-\gamma\Delta$, and, using $r(1-\gamma)+\gamma e_0=s_0$,
\[
s(A_1y)=r(1-\gamma)\bigl(r\,y(0)-\eps\beta\Delta\bigr)+\gamma e_0r\,y(0)=s_0\,r\,y(0)-r(1-\gamma)\eps\beta\Delta .
\]
So $(1-s_0)\Phi(A_1y)=r\,y(0)-r(1-\gamma)\eps\beta\Delta\le r\,y(0)-\gamma\Delta=s(y)$, because
$\gamma=\eps\beta\cdot\eps c$ and $\eps c\le r(1-\gamma)$. The last inequality holds because $c<1/2$
and $c+\beta<1$ for an admissible pair, as noted before Lemma~\ref{lem:onegate}; so
$\beta c\le1/4$, $\eps c<3/8$, $r>5/8$ and $\gamma\le\eps^2/4<1/6$, and $r(1-\gamma)>25/48>\eps c$.

(b) $\Phi(y)\ge\|y\|_\infty/c$. Here $\|y\|_\infty=y(1)$ and $y(0)\ge y(1)/\varrho_1$, so
$s(y)\ge y(1)\bigl((1-\gamma)(r-\eps\beta(1-\eps))+\gamma e_0\bigr)$. Since
$1-s_0=\eps c\,(1+\eps\beta(1-\eps))$, the claim $s(y)\ge(1-s_0)\,y(1)\,(1-c)/c$ follows from
$(1-\gamma)(r-\eps\beta(1-\eps))+\gamma e_0\ge e_0(1+\eps\beta(1-\eps))$. As $r-e_0=1-\eps$, this is
$(1-\eps)\bigl((1-\gamma)(1-\eps\beta)-e_0\eps\beta\bigr)\ge0$, which holds because
$\eps\beta\le1/2$, $\gamma\le9/64$ and $e_0\le3/4$.

(c) $c\,(S_\mu-r^{\mu-1})\le1-r^{\mu-1}$ for every $\mu\ge1$. Let
$f(\mu)=1-(1-c)r^{\mu-1}-cS_\mu$. Then $f(1)=0$ and
$f(\mu+1)-f(\mu)=(1-c)(1-r)r^{\mu-1}-ce_0^{\,\mu}=ce_0\,(r^{\mu-1}-e_0^{\,\mu-1})\ge0$, because
$(1-c)(1-r)=ce_0$.

Now let $\hat y_u\in\mathcal K$, $\mu=\mu_u$, and let $y'$ be $A_1\hat y_u$ or
$s(\hat y_u)\mathbf 1$ according to $d_u$. By (a), then (b) and (c),
\[
\Phi(\hat y_{u+1})\le r^{\mu-1}\Phi(y')\le r^{\mu-1}\bigl(\Phi(\hat y_u)-\|\hat y_u\|_\infty\bigr)
\le\Phi(\hat y_u)-S_\mu\|\hat y_u\|_\infty ,
\]
the last step being
$(1-r^{\mu-1})\Phi(\hat y_u)\ge(1-r^{\mu-1})\|\hat y_u\|_\infty/c\ge(S_\mu-r^{\mu-1})\|\hat y_u\|_\infty$.
Summing gives
$\sum_{u\le U}S_{\mu_u}\|y_u\|_\infty\le\sum_{u\le U}S_{\mu_u}\|\hat y_u\|_\infty\le\Phi(\mathbf 1)$
for every $U$.
\end{proof}

The next lemma holds for an arbitrary base. Among the gates that keep a given pointer, only the
highest one of a window is used as the gate of a chain step; the others must then be resampled.

\begin{lemma}[one step, any base]\label{lem:ordstep}
Let $\jt=S_\eps(j)$ for an arbitrary base $j$. Let $k$ be an active gate with $P:=j_k\ge g_0$ and
let $m\ge0$ with $k+m\le n$. Let $C^{\circ}$ be the event that $\pi_k=0$, that no gate in
$(k,k+m]$ is a reset, and that $\pi_g=1$ for every gate $g\in(k,k+m]$ with $j_g=P$. Let $a$ be the
number of gates $g\in(k,k+m]$ with $j_g=P$. Let $t_1<\dots<t_u$ be the positions in $(P,k)$ that
are the pointer of at least one active gate of $(k,k+m]$, and put $t_0=P$ and $d_i=t_i-t_{i-1}$.
Then, with $G'$ as in Lemma~\ref{lem:onestep},
\[
\PP[C^{\circ}\cap G'\mid\cF_P]\ \le\ (1-\eps)\,e_0^{\,a}\,r^{m-a-u}\,\bigl(B_{d_1}B_{d_2}\cdots B_{d_u}\mathbf 1\bigr)(x_P).
\]
\end{lemma}

\begin{proof}
All gates involved exceed $g_0$. For a gate $g\in(k,k+m]$ with $j_g=P$, the event requires that
it is resampled and not a reset; by Lemma~\ref{lem:onegate}(i) this has conditional probability
at most $e_0$ given $\cF_{g-1}$, on $G_g$. For $1\le i\le u$ let $\mu_i\ge1$ be the number of
active gates of $(k,k+m]$ with pointer $t_i$. If such a gate keeps its pointer and is not a
reset, then $x_{t_i}=0$, and $x_{t_i}$ is $\cF_{k-1}$-measurable; so, as in the proof of
Lemma~\ref{lem:onestep}, it is not a reset with conditional probability at most
$\varphi(x_{t_i})$. Every other gate of $(k,k+m]$ is not a reset with conditional probability at
most $r$, and $\PP[\pi_k=0\mid\cF_{k-1}]=1-\eps$. Since $\varphi\le r$, we have
$\varphi^{\mu_i}\le r^{\mu_i-1}\varphi$, and therefore
\[
\PP[C^{\circ}\cap G'\mid\cF_{k-1}]\le(1-\eps)\,e_0^{\,a}\,r^{m-a-u}\prod_{i=1}^{u}\varphi(x_{t_i})\cdot\mathbf 1_{G''},
\qquad G''=\bigcap_{P<g<k}G_g .
\]
For $w\ge0$, $q\ge g_0$ and $d\ge1$, Lemma~\ref{lem:onegate}(ii) gives
$\EE[w(x_{q+1})\mathbf 1_{G_{q+1}}\mid\cF_q]\le(Mw)(x_q)$, and by iteration
$\EE[w(x_{q+d})\prod_{q<g\le q+d}\mathbf 1_{G_g}\mid\cF_q]\le(M^dw)(x_q)$. Applying this from
$t_{u-1}$ to $t_u$ with $w=\varphi\mathbf 1$, then from $t_{u-2}$ to $t_{u-1}$ with
$w=\varphi\cdot B_{d_u}\mathbf 1$, and so on down to $t_0=P$, gives
$\EE[\prod_i\varphi(x_{t_i})\mathbf 1_{G''}\mid\cF_P]\le(B_{d_1}\cdots B_{d_u}\mathbf 1)(x_P)$.
\end{proof}

For an active gate $k$ with $j_k\ge g_0$ and $m\ge0$, with $a$, $u$ and the $d_i$ as in
Lemma~\ref{lem:ordstep}, let
\[
w_{k,m}=(1-\eps)\,e_0^{\,a}\,r^{m-a-u}\,\|B_{d_1}\cdots B_{d_u}\mathbf 1\|_\infty,
\qquad
\Theta^{\mathrm{ord}}_D=(1-\eps)\Bigl(\frac1{\eps c+\eps^2\beta c\,(1-\eps)}+\frac{r^D}{\eps c}\Bigr).
\]

\begin{lemma}[one window]\label{lem:ordsum}
For a $(D,N)$-ordered base and every $P_0\le n$, the sum of $w_{k,P_0-k}$ over the active gates
$k\le P_0$ with $j_k\ge g_0$ is at most $\Theta^{\mathrm{ord}}_D$.
\end{lemma}

\begin{proof}
Since $e_0\le r$ and $B_dw\le r\|w\|_\infty\mathbf 1$, we have $w_{k,m}\le(1-\eps)r^m$; the gates
with $P_0-k\ge D$ contribute at most $(1-\eps)\sum_{m\ge D}r^m=(1-\eps)r^D/(\eps c)$. Let $\mathcal A$
be the set of active gates $k$ in $(P_0-D,P_0]$ with $j_k\ge g_0$. Pointers of active gates do not
decrease, and an active gate at most $N\le g_0$ has a pointer below $g_0$; so $\mathcal A$
consists of the highest active gates up to $P_0$. Let $T_1>T_2>\dots>T_q$ be the pointers that
occur in $\mathcal A$, and let $\mu_i$ be the number of gates of $\mathcal A$ with pointer $T_i$;
those with pointer $T_i$ lie below those with pointer $T_{i-1}$. Let $k\in\mathcal A$ have pointer
$T_i$, and let $a$ of the gates of $\mathcal A$ with the same pointer lie above $k$. All active
gates of $(k,P_0]$ belong to $\mathcal A$ and are less than $D$ above $k$, so their pointers are
at most $k-1$: in Lemma~\ref{lem:ordstep}, with $m=P_0-k$, the number $a$ is this one, $u=i-1$,
and the positions are $T_{i-1}<\dots<T_1$, so that
$B_{d_1}\cdots B_{d_u}\mathbf 1=B_{\delta_{i-1}}B_{\delta_{i-2}}\cdots B_{\delta_1}\mathbf 1$ with
$\delta_l=T_l-T_{l+1}\ge1$. Let $y_1=\mathbf 1$ and $y_{l+1}=r^{\mu_l-1}B_{\delta_l}y_l$. Then, by homogeneity,
$B_{\delta_{i-1}}\cdots B_{\delta_1}\mathbf 1=r^{-(\mu_1-1)-\dots-(\mu_{i-1}-1)}y_i$, while
$m-a-u\ge(\mu_1-1)+\dots+(\mu_{i-1}-1)$, because $(k,P_0]$ contains the $a$ gates above $k$ with
pointer $T_i$ and all $\mu_1+\dots+\mu_{i-1}$ gates with a larger pointer. Hence
$w_{k,P_0-k}\le(1-\eps)\,e_0^{\,a}\,\|y_i\|_\infty$. Summing over $a=0,\dots,\mu_i-1$ and over $i$,
Lemma~\ref{lem:sumspacing} bounds the total by $(1-\eps)\sum_iS_{\mu_i}\|y_i\|_\infty\le(1-\eps)/(1-s_0)$.
\end{proof}

\begin{theorem}[ordered bases]\label{thm:ordered}
Let $D\ge2$, $N\ge0$ and $0<\eps\le3/4$. For each $n$ fix a $(D,N)$-ordered base of length $n$, and
suppose that $\Theta^{\mathrm{ord}}_D(\eps;c,\beta)<1$ for some pair $(c,\beta)$ with
$\beta\le2/3$ that is admissible for these bases. Then the conclusion of
Theorem~\ref{thm:fixed} holds for them. With the pairs of Lemma~\ref{fact:densall} the hypothesis
holds at $\eps=0.74$ for every $D\ge16$, and at $\eps=0.737$ for every $D\ge32$.
\end{theorem}

\begin{proof}
As for Theorem~\ref{thm:fixed}, with these changes. A base step of a typed skeleton requires
$P_s=j_{k_s}$; then $k_s$ is active, since $P_s\ge g_0\ge1$ and $P_s\le k_s-2$. Call a dependency
chain \emph{top-reading} if at every step with $\pi_{k_s}=0$ no gate $g\in(k_s,P_{s-1}]$ has
$\pi_g=0$ and $j_g=P_s$. If a dependency chain exists, so does a top-reading one with the same
positions $P_0,\dots,P_\ell$: at each step with $\pi_{k_s}=0$ replace $k_s$ by the largest
$k\in[k_s,P_{s-1}]$ with $\pi_k=0$ and $j_k=P_s$; then $\jt_k=P_s$ still, and $(k,P_{s-1}]$ still
contains no reset. So it suffices to count top-reading chains. For a typed skeleton let
$E^{\circ}_\sigma$ be the event $E^*_\sigma$ together with $\pi_g=1$ for every base step $s$ and
every $g\in(k_s,P_{s-1}]$ with $j_g=P_s$. By Lemma~\ref{lem:ordstep} in place of
Lemma~\ref{lem:onestep}, the bound of Lemma~\ref{lem:typed} holds for $E^{\circ}_\sigma$ with the
factor $w_{k_s,P_{s-1}-k_s}$ at each base step. In the kernel the base term becomes
$w_{k,P-k}f(j_k)$, summed over the active gates $k\le P$ with $j_k\ge g_0$, and
Lemma~\ref{lem:ordsum} gives $(W h)(P)\le(1/((\lambda+1)c)+\Theta^{\mathrm{ord}}_D)h(P)$. The
numerical claims: $\Theta^{\mathrm{ord}}_D$ decreases in $D$, and with $c=\varpi(\eps)-10^{-4}$ and
$\beta=0.4999$ it is at most $0.981$ for $(\eps,D)=(0.74,16)$ and at most $0.986$ for
$(0.737,32)$ (script \texttt{ord\_bound.py}).
\end{proof}

\begin{remark}\label{rem:ordscope}
\begin{enumerate}[label=(\roman*),leftmargin=2.4em]
\item With $c=\varpi(\eps)$ and $\beta=1/2$ the factor $\Theta^{\mathrm{ord}}_D$ equals $1$ at
$\eps=0.7365$ for $D=16$ and at $0.7342$ for $D=\infty$; for $D\le8$ it gives nothing below $3/4$.
\item Lemma~\ref{lem:sumspacing} is nearly sharp: when all distances are $2$ and no two gates
read the same position, the sum is within $1.5\%$ of $1/(1-s_0)$ at $\eps=0.74$, and in a
numerical optimization over sequences of distances, at the four noise levels tried, the constant
sequence $2$ was the worst. For $j^H$, where all distances are $1$, Theorem~\ref{thm:fixed} is
stronger.
\item The density of Lemma~\ref{lem:fddens} was not extended to ordered bases, so only the
pairs of Lemma~\ref{fact:densall} are available here. The constants are astronomical, as in
Remark~\ref{rem:fixedscope}(iv).
\item Gates that read the same position cost nothing here because of the top-reading chains:
of the gates of a window that keep a given pointer only the highest is counted. A block of $\mu$
gates reading one position then contributes $S_\mu=\sum_{a<\mu}e_0^{\,a}\le1/(1-e_0)$, where the
count of Theorem~\ref{kp:thm34} has $\sum_{a<\mu}r^a\le1/(\eps c)$.
\item What remains is one case. Lemma~\ref{lem:ordstep} holds for every base. It gives the
factor $r$ of Theorem~\ref{kp:thm34}, no gain, for the gates of a window whose pointer lies strictly below the next
target $P$ of the chain or at or above $k$, and for all but one of the gates that share a pointer
in $(P,k)$. For ordered bases the last two kinds are accounted for in Lemma~\ref{lem:ordsum}.
Reads strictly below the next target are the obstruction of
Remark~\ref{rem:needs}(ii): they can land on positions whose values the chain below already
fixes. A base whose pointers decrease along every window gets no gain from this section.
\end{enumerate}
\end{remark}

\subsection{Reads below the next target: two lemmas, and what is missing}\label{sec:below}

This subsection does not prove a bound for bases whose window gates read below the next target.
It records two facts that such a proof would use, and a calculation that shows how much more it
needs than Sections~\ref{sec:fixed} and~\ref{sec:ordered}. The assumptions and notation are
those of Section~\ref{sec:ordered}; the base is arbitrary.

Consider two consecutive base steps of a chain, with gates $k_s$ and $k'=k_{s+1}$ and targets
$P_s=j_{k_s}$ and $P'=j_{k'}$. A gate $g$ of the upper window that keeps a pointer
$\tau<P_s$ requires $x_\tau=0$. If $\tau$ lies in the lower window, $k'\le\tau<P_s$, then
$x_\tau=x_{k'}\oplus((\tau-k')\bmod2)$, because that window contains no reset. So the requirement
is a condition on the single value $x_{k'}$, and the same holds for $x_{P_s}$ and for every
other read into that window. The first lemma says that this value is not known in advance.

\begin{lemma}[the value at a chain gate]\label{lem:typebit}
Let $k'$ be a gate with $P':=j_{k'}\ge g_0$ and $k'-P'\ge3$, and let $\Psi\ge0$ be a function on
$\{0,1\}$. Put $\Psi^\sharp(1)=\Psi(0)$ and
$\Psi^\sharp(0)=\max\bigl(r\,\Psi(0)+(1-r)\,\Psi(1),\ (1-\gamma)\,\Psi(1)+\gamma\,\Psi(0)\bigr)$. Then
\[
\EE\bigl[\mathbf 1[\pi_{k'}=0]\;\Psi(x_{k'})\;\mathbf 1_{G_{k'-1}\cap G_{k'-2}}\bigm|\cF_{k'-3}\bigr]\ \le\ (1-\eps)\,\Psi^\sharp(x_{P'}) .
\]
\end{lemma}

\begin{proof}
If $\pi_{k'}=0$ then $x_{k'}=\NOR(x_{k'-1},x_{P'})$, and $\pi_{k'}$ is independent of
$\cF_{k'-1}$, while $x_{P'}$ is $\cF_{k'-3}$-measurable. If $x_{P'}=1$ then $x_{k'}=0$, which gives
$\Psi^\sharp(1)$. Let $x_{P'}=0$; then $x_{k'}=1-x_{k'-1}$. Put $w(b)=\Psi(1-b)$. By
Lemma~\ref{lem:onegate}(ii), applied at $k'-1$ and then at $k'-2$, as in the proof of
Lemma~\ref{lem:ordstep},
$\EE[w(x_{k'-1})\mathbf 1_{G_{k'-1}\cap G_{k'-2}}\mid\cF_{k'-3}]\le(M^2w)(x_{k'-3})$. If
$w(0)\ge w(1)$, the formula for $M^dw$ in the proof of Lemma~\ref{lem:spacing} gives both entries
of $M^2w$ at most
$w(0)-\gamma(w(0)-w(1))=(1-\gamma)\Psi(1)+\gamma\Psi(0)$. If $w(0)<w(1)$, then
$Mw=(w(0)+r(w(1)-w(0)),\,w(0))$, so $(M^2w)(1)=(1-r)w(0)+r\,w(1)$ and
$(M^2w)(0)=(Mw)(0)-\eps\beta((Mw)(0)-w(0))\le(Mw)(0)$; both are at most
$(1-r)\Psi(1)+r\,\Psi(0)$.
\end{proof}

\noindent For example, a read into the lower window has $\Psi=(\Psi(0),\Psi(1))$ equal to
$(e_0,r)$ or $(r,e_0)$, according to the value of $x_{k'}$ that suits it, and then
$\Psi^\sharp(0)$ is $s_0$ or $r-\eps c(1-\eps)\le s_0$: when the lower step lands on a zero, such
a read costs at least as much as a fresh read two or more above the target. When the lower step lands on a one,
$x_{k'}=0$ is certain, and the cost appears only through the parity of the window length. All
reads into one window share the one value $x_{k'}$, so they are charged once, not once each.

The second lemma is deterministic. It lets a chain count discard every step that has a sibling
with a higher target in the same run.

\begin{lemma}[highest targets]\label{lem:highest}
Fix a parent array and $g_0\ge1$, and call a dependency chain \emph{admissible} if all its
positions are at least $g_0$. If there is an admissible dependency chain of length $\ell$ from
$P_0$, then there is one, $(P_0,k_1,P_1,\dots,k_\ell,P_\ell)$, with the following property for
every $s$: no gate $k\le P_{s-1}$ with no reset in $(k,P_{s-1}]$ has a pointer $t$ with
$P_s<t\le k-2$ and no reset in $(P_s,t]$. Moreover, of the gates $k\le P_{s-1}$ with pointer
$P_s$ and no reset in $(k,P_{s-1}]$, the gate $k_s$ is the highest.
\end{lemma}

\begin{proof}
For $q\ge g_0$ let $\delta(q)$ be the largest length of an admissible dependency chain from $q$.
If $q<t$ and $(q,t]$ contains no reset, then $\delta(t)\ge\delta(q)$: a chain from $q$ begins with
a gate $k_1\le q$ with no reset in $(k_1,q]$, hence none in $(k_1,t]$, so the same sequence is a
chain from $t$. Now build the chain step by step. Given $P_{s-1}$ with
$\delta(P_{s-1})\ge\ell-s+1$, the set of gates $k\le P_{s-1}$ with no reset in $(k,P_{s-1}]$, whose
pointer $q$ satisfies $g_0\le q\le k-2$ and $\delta(q)\ge\ell-s$, is not empty. Among these choose
$k_s$ with the largest pointer and, for that pointer, the largest gate; put $P_s=\jt_{k_s}$. If a
gate $k$ as in the statement existed, its pointer $t$ would satisfy $t\ge g_0$ and
$\delta(t)\ge\delta(P_s)\ge\ell-s$, against the choice of $k_s$.
\end{proof}

\begin{remark}[what is missing]\label{rem:below}
\begin{enumerate}[label=(\roman*),leftmargin=2.4em]
\item \emph{Where a read below the target can land.} Suppose that every active gate reads at
least $D$ below itself, and that no window gate has a pointer $D$ or more below the pointer of
the step's gate; let $D\ge3$, and let the lower step be a base step. Then a read below $P_s$ lies
above $P'$. Either it is in the lower window or at $k'$, where Lemma~\ref{lem:typebit} applies,
or it is in the lower gap $(P',k')$, where it is fresh as in Lemma~\ref{lem:ordstep}. So each such read has a cost. Without the assumption a read can land
several steps further down, and a proof has to follow it there.
\item \emph{Why this is not yet a theorem.} The weight of the upper steps becomes a function of
the value at the current target, a pair of numbers, and the reads still to be resolved have to
be carried from one step to the next. The potential of Lemma~\ref{lem:sumspacing} handles a pair
within one step; here the pair, together with the pending reads, passes between steps, and the
base chooses for every gate where its read lands. We did not find a potential for this.
\item \emph{How much is there to gain.} Two calculations, neither a proof. In a conservative
model that credits only the gate just above the step's gate, and lets the base choose for every
window length where that gate reads, a value iteration gives the growth factor $1.046$ at
$\eps=0.74$, against $1.061$ for the count of Theorem~\ref{kp:thm34}; that would reach $1$ only near $\eps=0.7475$.
So one gate per window is not enough. If every gate is credited and every read below the target
lands in the next window with the parity that suits a zero, the weights form the recursion
\[
Y'(0)=(1-\eps)\sum_{m\ge0}r^m\,\Psi_m^\sharp(0),\qquad Y'(1)=(1-\eps)\sum_{m\ge0}e_0^{\,m}\,Y(m\bmod2),
\qquad\Psi_m(c)=Y(c\oplus m),
\]
whose growth factor is $0.955$ at $\eps=0.74$ and equals $1$ near $\eps=0.732$; the script is
\texttt{typed\_model.py}. This suggests that the case can be brought below $3/4$ by about as
much as Theorem~\ref{thm:ordered}, but only by an argument that credits every gate.
\item \emph{Lemma~\ref{lem:highest} is not used above.} For a gate of the window that keeps a
pointer $t>P_s$ it turns the requirement $x_t=0$ into: $x_t=0$ and there is a reset in $(P_s,t]$.
For a base whose pointers decrease with the gate it gives a requirement on the gates just below
$k_s$, which have higher targets. In both cases the gain we could extract is small, and we did
not pursue it.
\end{enumerate}
\end{remark}

\subsection{Counting positions: large distances from noise \texorpdfstring{$0.455$}{0.455}}\label{sec:rows}

Theorem~\ref{thm:fixed} counts dependency chains, and for large $H$ its count decays only from
noise $0.69$ on. On a fixed-distance base the chains overlap: a gate that keeps its pointer leads
to the row $H$ below, and many chains pass through the same positions. This subsection counts
positions. For large $H$ it shows that the relevance-pruned iteration needs logarithmically many
rounds from noise $0.455$ on. The proof follows the open positions of two consecutive rows
together, as a Markov chain with $16$ states that is driven by the row below, and uses a potential
that a program finds and checks in interval arithmetic. The result is again about this family of
bases only.

\paragraph{Statuses.} For $t\ge0$ call position $i$ \emph{open at level $t$} if $D(i)>t$, and
\emph{closed at level $t$} otherwise, and put $I^t_i=\ind[D(i)>t]$. (Section~\ref{sec:above} calls
such positions $t$-open and $t$-determined.) So $I^t_0=0$, $I^0_i=1$ for
$i\ge1$, and $I^{t+1}_i\le I^t_i$. By Lemma~\ref{kp:pruned}(a) the relevance-pruned iteration has
finished after $t$ rounds exactly if no position is open at level $t$. For bits $v,o,b,w$ put
\[
f(v,o,b,w)=\begin{cases}
w&\text{if }v=0,\ b=1,\\
o\wedge w&\text{if }v=1,\ b=1,\\
o\vee w&\text{if }v=0,\ b=0,\\
o&\text{if }v=1,\ b=0.
\end{cases}
\]
It is nondecreasing in $o$ and in $w$, and $f(v,0,b,0)=0$.

\begin{lemma}[status of a gate]\label{lem:rowstatus}
Let $j$ be a parent array with trace $x$, let $i\ge1$, $t\ge0$ and $p=j_i$. Then
\[
I^{t+1}_i\le f\bigl(x_{i-1},\,I^{t+1}_{i-1},\,x_p,\,I^t_p\bigr),
\]
with equality if $i$ is not a NOT gate, and also if $I^t_p=0$. Moreover $D(i)\le i$. If
$I^{t+1}_i=1$, then there is a gate $k\le i$, not a NOT gate, such that no gate in $(k,i]$ is a
reset and $I^t_{j_k}=1$.
\end{lemma}

\begin{proof}
Write $v=x_{i-1}$, $o=I^{t+1}_{i-1}$, $b=x_p$, $w=I^t_p$. Gate $1$ has $p=0$ and
$v=b=1$; by Lemma~\ref{kp:switch}(a), $D(1)=1$, so $I^{t+1}_1=0=o\wedge w$. Let $i\ge2$ and use
the four cases of Lemma~\ref{kp:switch}(a). If $i$ is not a NOT gate: for $v=0,b=1$ it is a
relevant reset and $I^{t+1}_i=\ind[1+D(p)>t+1]=w$; for $v=1,b=1$ it is an irrelevant reset and
$I^{t+1}_i=o\wedge w$; for $v=0,b=0$ it is relevant and neither a reset nor a NOT gate, and
$I^{t+1}_i=o\vee w$; for $v=1,b=0$, $I^{t+1}_i=o$. If $i$ is a NOT gate, then $p=i-1$, $b=v$,
$D(i)=D(i-1)$ and $I^{t+1}_i=o$; here $f=o\wedge w$ or $f=o\vee w$, and $w=I^t_{i-1}\ge o$, so
$f\ge o$, with equality if $w=0$. The bound $D(i)\le i$ follows by induction from $D(1)=1$ and
$D(i)\le\max(D(i-1),1+D(j_i))$. For the last claim apply Lemma~\ref{kp:switch}(b) with $t+1$ in
place of $t$: there is a relevant gate $a\le i$, not a NOT gate, with $D(j_a)\ge t+1$, and
$D(j_g)\ge t+1$ for every reset $g\in(a,i]$. Let $k$ be the last reset in $(a,i]$ if there is
one, and $k=a$ otherwise; a reset is not a NOT gate.
\end{proof}

\begin{definition}[the two-row chain]\label{def:tworow}
Let $0<\eps\le1$ and $0\le d_a\le d_b\le1$. A \emph{state} is a quadruple
$s=(x',o',x'',o'')$ of bits; it is \emph{closed} if $o'=o''=0$. A \emph{symbol} is a pair
$y=(b,w)$ of bits, its value and its status. A pair $(v,o)$ of bits that \emph{reads} a pair
$(b,w)$ becomes $(\NOR(v,b),f(v,o,b,w))$. In a \emph{step} with the symbol $y$ and with
densities $d_1,d_2$, first $(x',o')$ reads $y$ with probability $1-\eps$, the pair $(1,0)$ with
probability $\eps d_1$ and the pair $(0,0)$ with probability $\eps(1-d_1)$; then $(x'',o'')$
reads the new pair $(x',o')$ with probability $1-\eps$, the pair $(1,0)$ with probability
$\eps d_2$ and $(0,0)$ with probability $\eps(1-d_2)$. For a real function $g$ on the states,
$(T_yg)(s)$ is the largest expected value of $g$ after one step from $s$ with the symbol $y$,
over $d_1\in[d_a,d_b]$ and over $d_2\in[d_a,d_b]$, where $d_2$ may depend on what $(x',o')$ has
read; $(T^-_yg)(s)$ is the least such value. A word of symbols is \emph{admissible} if no two
consecutive symbols both have the value $1$. Finally $\phi(s)=o''-\theta o'$ for a number
$\theta$.
\end{definition}

\noindent The operators $T_y$ and $T^-_y$ are monotone, map $0$ to $0$, and do not increase the
largest absolute value of a function. A closed state stays closed in a step with a symbol of
status $0$.

\begin{lemma}[potentials]\label{lem:rowscert}
Let $0.455\le\eps<1$, $\delta=2\cdot10^{-5}$, $\theta=0.999$ and $K=8$, and let $\ell_4$ and
$\mu_4$ be the numbers of Lemma~\ref{lem:fddens} for this $\eps$. For every $d$ with
$\ell_4-\delta\le d\le\mu_4+\delta$ there are functions $c_y$ on the states, one for each symbol
$y$, with values in $[-4,4]$ and equal to $0$ on the closed states, such that, with
$[d_a,d_b]=[d-\delta,d+\delta]$, for every admissible word $z_1\dots z_Ky$:
\begin{align}
T_{z_1}\cdots T_{z_K}\bigl(T_y(\phi+c_y)-c_{z_K}\bigr)&\le0,\label{eq:rowsC1}\\
T^-_{z_1}\cdots T^-_{z_K}\,c_{z_K}&\ge0.\label{eq:rowsC2}
\end{align}
\end{lemma}

\begin{proof}
By computation. The program \texttt{rows\_check.c} covers the interval of noise levels by
intervals $[e_a,e_b]$. On each it encloses $\ell_4$ and $\mu_4$ and covers
$[\ell_4-2\delta,\mu_4+2\delta]$ by intervals of $d$ of which consecutive ones overlap by at
least $2\delta$; so every interval $[d-\delta,d+\delta]$ with $\ell_4-\delta\le d\le\mu_4+\delta$
lies in one of them. For each such box it finds four functions by a search in floating point,
which is not part of the proof, and then checks~\eqref{eq:rowsC1} and~\eqref{eq:rowsC2} for all
admissible words, for the operators in which the maximum and the minimum are over the whole box:
over $d_1,d_2$ in the interval of $d$ and over the noise level in $[e_a,e_b]$, separately for the
two reads. These operators lie above $T_y$ and below $T^-_y$, and all of them are monotone, so the
two inequalities for them imply those of the lemma. The expectation is affine in each of the
densities and noise levels, so the maximum and the minimum are attained at the ends of the
intervals. The check uses doubles with every operation rounded in the safe direction, as in
\texttt{potential\_check.c}. It walks the tree of words from the last symbol back, and leaves a
branch when the function is already nonpositive, respectively nonnegative, at all states; this
is allowed because the operators are monotone and map $0$ to $0$. The runs
\texttt{rows\_check 0.455 0.5 0.999 8 0.002 0.004},
\texttt{rows\_check 0.5 0.75 0.999 8 0.008 0.012} and
\texttt{rows\_check 0.75 1.0 0.999 8 0.008 0.012} print \texttt{PROVED}, with $747$, $712$ and $535$ boxes
(Section~\ref{sec:checks}).
\end{proof}

\begin{theorem}[fixed-distance bases, counting positions]\label{thm:rows}
Let $0.455\le\eps<1$. There are constants $H_0$, $A$ and $B$ such that for every $H\ge H_0$ and
every $n\ge2$, with probability at least $1-3/n$, the relevance-pruned iteration on
$S_\eps(j^H)$ finishes within $A\ln(n+1)+B$ rounds, and the prefix--Jacobi iteration within one
round more.
\end{theorem}

\noindent The constants do not depend on $H$. The proof takes the rest of this subsection. Fix
$\eps$, and let $\delta$, $\theta$, $K$ be as in Lemma~\ref{lem:rowscert}. Let $H\ge2K$ and
$n\ge2$, and let $x$ and $D$ be the trace and the depth of $\jt=S_\eps(j^H)$.

\paragraph{Step 0: constants, events, weights.}
The root $\varpi(\eps)$ of Lemma~\ref{fact:densall} increases with $\eps$ and
$\varpi(0.455)>1/4$. By that lemma the pair $(1/4,0)$ is admissible: there are constants and
integers $g_0=g_0(n)\le K_0\ln(n+1)+K_1$, which by its proof depend on $\eps$ only and not on the
base, such that, with probability at least $1-1/n$,
$(O_{g-1}-1)/g\ge1/4$ for all $g_0<g\le n$. Let $\mathsf T$ be the least $g\in(g_0,n]$ with
$(O_{g-1}-1)/g<1/4$, and $\mathsf T=n+1$ if there is none; the event $\{\mathsf T>g\}$ is in
$\cF_{g-1}$. Call a gate $g$ a \emph{stop} if it is a reset, or $g\ge\mathsf T$, or $g>n$. For
$g_0<g\le n$,
\begin{equation}\label{eq:rowsstop}
\PP[g\text{ is not a stop}\mid\cF_{g-1}]\le r_0:=1-\eps/4,
\end{equation}
because on $\{\mathsf T>g\}$ gate $g$ is a reset with conditional probability at least
$\eps(O_{g-1}-x_{g-1})/g\ge\eps/4$, as in Lemma~\ref{lem:onegate}. Put $K_r=r_0/(1-r_0)$.

By Lemma~\ref{lem:fddens} with $s=4$ and $\tau=\delta/40$ there is an integer $K_d\ge13$, which
depends on $\eps$ only, such that with probability at least $1-1/n$
\[
\ell_4-\delta/2\le\psi_g\le\mu_4+\delta/2\quad\text{for all }g\in[g_1,n],\qquad g_1:=K_d\bigl(H+\lceil\ln(n+1)\rceil\bigr).
\]
For an integer $b$ let $\Gamma^\circ_b$ be the event that these two inequalities hold for all
$g\in[g_1,\min(b,n)]$; it is in $\cF_b$, and it is the sure event if $b<g_1$. Let $m_0\ge8/\delta$
be an integer, fixed in Step~4, and put
\[
t_0=\max\bigl(g_0,K_d\lceil\ln(n+1)\rceil\bigr),\quad m_1=K_d+m_0+4,\quad A_0=g_1+(m_0+2)H,\quad A_1=A_0+H-K .
\]
Then $A_1\le t_0+m_1H$. For $g<g'$ we have $|\psi_{g'}-\psi_g|\le(g'-g)/(g'+1)$. Hence:
\begin{equation}\label{eq:rowslate}
\begin{minipage}{0.86\linewidth}
if $b\ge g_1+m_0H$ and $\Gamma^\circ_b\cap\{\mathsf T>g_1\}$ holds, then
$|\psi_g-\psi_{g'}|\le\delta/2$ for all $g,g'\in[b,\min(b+3H,n)]$, $\ \ell_4-\delta\le\psi_g\le\mu_4+\delta$
for all $g\in[g_1,\min(b+3H,n)]$, and $\mathsf T>\min(b+3H,n)$.
\end{minipage}
\end{equation}
For the last claim, $(O_{g-1}-1)/g=\psi_{g-1}-1/g\ge1/3-\delta-1/13>1/4$ for $g_1<g\le\min(b+3H,n)$,
since $\ell_4\ge\ell_1=1/3$.

For $1\le i\le n$ let $\Gamma^*_i=\{\mathsf T>i\}\cap\Gamma^\circ_{i-2H}$. These events decrease
with $i$, and $\Gamma^*_i\in\cF_{i-1}$. For $t\ge t_0$ put
\[
u_t(i)=\PP[I^t_i=1,\ \Gamma^*_i],\qquad \nu_i=i^{-\lambda},\qquad a_t=\sum_{t_0<i\le n}\nu_i\,u_t(i),\qquad
q_t(g)=\frac1g\sum_{p<g}u_t(p),
\]
with a number $\lambda\ge2$ fixed in Step~4, and $u_t(0)=0$. Since $D(i)\le i$, $u_t(i)=0$ for $i\le t_0$; and
$a_t\le\sum_{i\ge2}i^{-2}<1$. Since $\sum_{g>p}g^{-\lambda-1}\le p^{-\lambda}/\lambda$,
\begin{equation}\label{eq:rowsq}
\sum_{g\le n}\nu_g\,q_t(g)\le\frac{a_t}{\lambda}.
\end{equation}
Call gate $g$ a \emph{source at level $t$} if $\pi_g=1$ and $I^t_{\rho_g}=1$. Given $\cF_{g-1}$
the position $\rho_g$ is uniform on $\{0,\dots,g-1\}$, and $\Gamma^*_g\subseteq\Gamma^*_p$ for
$p<g$, so
\begin{equation}\label{eq:rowssource}
\PP[g\text{ is a source at level }t,\ \Gamma^*_g]\le\eps\,q_t(g).
\end{equation}

\paragraph{Step 1: early positions.}
For $k\le n$ let $N(k)$ be the largest $i\le n$ such that no gate in $(k,i]$ is a stop, and let
$N(k)=k$ for $k>n$; $N$ is nondecreasing. For a position $p$ put $E_0(p)=p$ and
$E_j(p)=N(E_{j-1}(p)+H)$, and for a gate $k$ put $\bar E_0(k)=N(k)$ and
$\bar E_j(k)=N(\bar E_{j-1}(k)+H)$.

\emph{Claim.} Let $t\ge0$, $m\ge1$, $1\le i\le n$, $\mathsf T>i$ and $I^{t+m}_i=1$. Then
\begin{enumerate}[label=(\alph*),leftmargin=2.4em]
\item there is a position $p\ge1$ with $I^t_p=1$ and $p+mH\le i\le E_m(p)$, or
\item there are $0\le j<m$ and a gate $k$ that is a source at level $t+m-j-1$, with
$k+jH\le i\le\bar E_j(k)$.
\end{enumerate}
\emph{Proof.} Induction on $m$. By Lemma~\ref{lem:rowstatus} there is a gate $k\le i$ with no reset
in $(k,i]$ and $I^{t+m-1}_{\jt_k}=1$. As $\mathsf T>i$, there is no stop in $(k,i]$, so
$i\le N(k)$. If $\pi_k=1$, this is (b) with $j=0$. Otherwise $\jt_k=\max(k-H,0)$, and this
position is open, so $i':=k-H\ge1$, $I^{t+m-1}_{i'}=1$ and $i'+H\le i\le N(i'+H)$. For $m=1$ this
is (a) with $p=i'$. For $m>1$ apply the claim to $i'$ with $m-1$; then (a) for $i'$ with a
position $p$ gives $p+mH\le i\le N(E_{m-1}(p)+H)=E_m(p)$, and (b) for $i'$ with $j$ and $k'$ gives
$k'+(j+1)H\le i\le\bar E_{j+1}(k')$, where $k'$ is a source at level $t+m-(j+1)-1$.

The number of $i\le n$ with $p+mH\le i\le E_m(p)$ is at most $1+\tau_1+\dots+\tau_m$, where
$\tau_j=N(\kappa_j)-\kappa_j$ and $\kappa_j=E_{j-1}(p)+H$. Here $E_{j-1}(p)+1$ is the first stop
after $\kappa_{j-1}$, so $\kappa_j$ is a stopping time, and by~\eqref{eq:rowsstop}
$\EE[\tau_j\mid\cF_{\kappa_j}]\le\sum_{s\ge1}r_0^s=K_r$ if $p>g_0$; this count is used in
Remark~\ref{rem:rows}(i) only. In the same way the number of
$i\le n$ with $k+jH\le i\le\bar E_j(k)$ has conditional expectation at most $1+(j+1)K_r$ given
$\cF_k$, if $k>g_0$. Now let $t\ge t_0$ and $t_0<i\le t_0+mH$. In case (a) we would have
$p\le t_0$, which is closed at level $t$. So only (b) occurs. The event that $k$ is a source at
level $t'$ and $\Gamma^*_k$ holds is in $\cF_k$, $\Gamma^*_i\subseteq\Gamma^*_k$ for $k\le i$, and
$\nu_i\le\nu_k$; a gate $k\le t_0$ is not a source at a level $t'\ge t_0$. So
by~\eqref{eq:rowssource} and~\eqref{eq:rowsq}
\begin{equation}\label{eq:rowsearly}
\sum_{t_0<i\le t_0+mH}\!\!\nu_i\,u_{t+m}(i)\le\sum_{j=0}^{m-1}\bigl(1+(j+1)K_r\bigr)\frac{\eps}{\lambda}\,a_{t+m-j-1}
\le\frac{\eps\,m(1+mK_r)}{\lambda}\max_{1\le j\le m}a_{t+m-j}
\end{equation}
for $t\ge t_0$.

\paragraph{Step 2: one window, and an ideal upper row.}
Fix $t\ge t_0$ and a base point $a$ with $A_0\le a<n$. Put $s=a-H-K$ and $J=\min(H-K,n-a)$, and
for $-K\le j\le J$ let $g'_j=a-H+j$ and $g''_j=a+j$. The gates $g'_j$ with $j>-K$ are the
\emph{lower row} $(s,a-K]$ and the gates $g''_j$ with $j>-K$ the \emph{upper row} $(a-K,a+J]$;
gate $g''_j$ has the base pointer $g'_j$. Let
\[
z_j=\bigl(x_{g'_j-H},\,I^t_{g'_j-H}\bigr),\qquad (x'_j,o'_j)=\bigl(x_{g'_j},\,I^{t+1}_{g'_j}\bigr),\qquad
m_j=\ind\bigl[\Gamma^\circ_{a+j-2H}\cap\{\mathsf T>g_1\}\bigr].
\]
We use $m_j$ for $1\le j\le J$ only. Since $g'_j-H\le s$ and $a+j-2H\le s$, the symbols $z_j$ and
the numbers $m_j$ are $\cF_s$-measurable; the $m_j$ do not increase with $j$; and every word
$z_i\dots z_j$ is admissible, because the positions $g'_i-H\ge1$ are consecutive and a trace has
no two consecutive ones. For $j\ge1$ we have $a+j-2H\ge g_1+m_0H$, and all positions from $s$ to
$a+J$ lie in $[a+j-2H,\,a+j+H]$. So by~\eqref{eq:rowslate}, on $\{m_j=1\}$,
\begin{equation}\label{eq:rowswin}
\ell_4-\delta\le\psi_s\le\mu_4+\delta,\qquad|\psi_g-\psi_s|\le\delta/2\ \ (s\le g\le a+J),\qquad
\mathsf T>a+j,
\end{equation}
so that $m_j=\ind[\Gamma^*_{a+j}]$ for $j\ge1$, and $m_j\le\ind[\Gamma^*_p]$ for every $p\le a+j$.

\emph{The ideal upper row.} Let $\beta_j$, $-K<j\le J$, be bits that, given $\cF_s$, are
independent of each other and of the noise, and equal to $1$ with probability $\psi_s$. Put
$(y_{-K},\hat o_{-K})=(0,0)$, and for $j>-K$ let $(y_j,\hat o_j)$ be the pair
$(y_{j-1},\hat o_{j-1})$ after it has read
\[
\begin{cases}
(x'_j,o'_j)&\text{if }\pi_{g''_j}=0,\\
(x_\rho,I^{t+1}_\rho)&\text{if }\pi_{g''_j}=1\text{ and }\rho:=\rho_{g''_j}\le g'_j,\\
(\beta_j,0)&\text{if }\pi_{g''_j}=1\text{ and }\rho_{g''_j}>g'_j .
\end{cases}
\]
In the third case call $j$ a \emph{near step}. Let $S_j=(x'_j,o'_j,y_j,\hat o_j)$, and let
$\mathcal G_j$ be generated by $\cF_{g'_j}$ and by the noise of the gates $g''_i$ and the bits
$\beta_i$ with $-K<i\le j$; so $\mathcal G_{-K}=\cF_s$. The upper row lies above the lower row, so $S_j$ is
$\mathcal G_j$-measurable, and given $\mathcal G_{j-1}$ the noise of the gates $g'_j$ and $g''_j$
and the bit $\beta_j$ are independent and have their unconditional laws. Call $j$ a \emph{source
step} if $g'_j$ is a source at level $t$, or if $\pi_{g''_j}=1$, $\rho_{g''_j}\le g'_j$ and
$I^{t+1}_{\rho_{g''_j}}=1$; in the second case $g''_j$ is a source at level $t+1$.

\emph{One step.} Let $-K<j\le J$, work on $\{m_{\max(j,1)}=1\}$, put $d=\psi_s$ and
$[d_a,d_b]=[d-\delta,d+\delta]$, and let $g$ be a function on the states with $|g|\le M$. Write
$e_j=\PP[j\text{ is a source step}\mid\mathcal G_{j-1}]$. Then
\begin{equation}\label{eq:rowsone}
(T^-_{z_j}g)(S_{j-1})-2M\,e_j\le\EE[g(S_j)\mid\mathcal G_{j-1}]\le(T_{z_j}g)(S_{j-1})+2M\,e_j .
\end{equation}
Indeed, gate $g'_j$ keeps its pointer with probability $1-\eps$; it is then not a NOT gate, and by
Lemma~\ref{lem:rowstatus} the pair $(x'_{j-1},o'_{j-1})$ reads $z_j$. With probability $\eps$ it
lists a uniform position below $g'_j$. If that position is closed at level $t$, the pair reads
its value with the status $0$, again by Lemma~\ref{lem:rowstatus}; otherwise $g'_j$ is a source. The
position holds a one with probability $d_1=\psi_{g'_j-1}$. So, up to an event of probability
$\PP[g'_j\text{ is a source at level }t\mid\mathcal G_{j-1}]$ on which a probability mass is moved,
the lower pair makes the first half of a step with the density $d_1$. Given also the noise of
$g'_j$, the ideal upper pair reads $(x'_j,o'_j)$ with probability $1-\eps$; with probability
$\eps(g'_j+1)/g''_j$ it reads a uniform position of $[0,g'_j]$, which holds a one with probability
$\psi_{g'_j}$ and gives a source step if it is open at level $t+1$; and with probability
$\eps(H-1)/g''_j$ it reads $(\beta_j,0)$. So, up to the source steps, it makes the second half of
a step with the density $d_2=\bigl((g'_j+1)\psi_{g'_j}+(H-1)\psi_s\bigr)/g''_j$.
By~\eqref{eq:rowswin}, $d_1$ and $d_2$ lie in $[d_a,d_b]$. Moving a probability mass $e$ changes
an expectation of a function with values in $[-M,M]$ by at most $2Me$, which
gives~\eqref{eq:rowsone}.

\emph{The window inequality.} By~\eqref{eq:rowswin} and Lemma~\ref{lem:rowscert} there are
functions $c_y$ for $d=\psi_s$ when $m_1=1$; they are chosen as a function of $\psi_s$, so they
are $\cF_s$-measurable, and we put $c_y=0$ when $m_1=0$. Let $1\le j\le J$ with $m_j=1$. The functions $\phi+c_{z_j}$ and
$T_{z_j}(\phi+c_{z_j})-c_{z_{j-1}}$ are at most $9$ in absolute value. Apply~\eqref{eq:rowsone}
to the steps $j,j-1,\dots,j-K$, and then~\eqref{eq:rowsC1} for the word $z_{j-K}\dots z_{j-1}z_j$:
\[
\EE\bigl[\phi(S_j)+c_{z_j}(S_j)-c_{z_{j-1}}(S_{j-1})\bigm|\cF_s\bigr]\le18\sum_{i=j-K}^{j}\sigma_i,
\qquad\sigma_i:=\PP[i\text{ is a source step}\mid\cF_s].
\]
In the same way~\eqref{eq:rowsC2} gives
$\Phi_j:=\EE[c_{z_j}(S_j)\mid\cF_s]\ge-8\sum_{i=j-K+1}^{j}\sigma_i$ for $1\le j\le J$ with
$m_j=1$. Let $\omega_j=\nu_{a+j}m_j$ for $1\le j\le J$; these numbers are $\cF_s$-measurable and
do not increase. Summation by parts gives
\[
\sum_{j=1}^{J}\omega_j\bigl(\Phi_{j-1}-\Phi_j\bigr)=\omega_1\Phi_0-\omega_J\Phi_J-\sum_{j=1}^{J-1}(\omega_j-\omega_{j+1})\Phi_j
\le\omega_1\Phi_0+8\sum_{i=1-K}^{J}\omega_{\max(i,1)}\,\sigma_i ,
\]
and so
\begin{equation}\label{eq:rowsideal0}
\sum_{j=1}^{J}\omega_j\,\EE[\phi(S_j)\mid\cF_s]\le\omega_1\Phi_0+\bigl(18(K+1)+8\bigr)\sum_{i=1-K}^{J}\omega_{\max(i,1)}\,\sigma_i .
\end{equation}
The functions $c_y$ vanish on closed states and are at most $4$, so
$\Phi_0\le4\,\EE[o'_0+\hat o_0\mid\cF_s]$. The ideal row starts closed, so $\hat o_0=1$ only if for
some $-K<i\le0$ the pair read a pair of status $1$: then $\pi_{g''_i}=0$ and $o'_i=1$, or $i$ is
a source step. So
\[
\Phi_0\le4\,\EE\Bigl[o'_0+\sum_{i=1-K}^{0}o'_i\Bigm|\cF_s\Bigr]+4\sum_{i=1-K}^{0}\sigma_i .
\]
Take expectations in~\eqref{eq:rowsideal0}. We have $\phi=\hat o-\theta o'$ on $S_j$;
$\omega_{\max(i,1)}\le\nu_{g''_i}\le\nu_{g'_i}$; and $m_j\le\ind[\Gamma^*_p]$ for $p\le a+j$, so
that $\EE[m_jo'_i]\le u_{t+1}(g'_i)$ for $i\le j$, and by~\eqref{eq:rowssource}
$\EE[m_{\max(i,1)}\sigma_i]\le\eps\,q_t(g'_i)+\eps\,q_{t+1}(g''_i)$. This gives
\begin{multline}\label{eq:rowsideal}
\sum_{j=1}^{J}\nu_{a+j}\,\EE[m_j\hat o_j]\le\theta\sum_{j=1}^{J}\nu_{g'_j}u_{t+1}(g'_j)
+8\sum_{i=1-K}^{0}\nu_{g'_i}u_{t+1}(g'_i)\\
+174\,\eps\sum_{i=1-K}^{J}\bigl(\nu_{g'_i}q_t(g'_i)+\nu_{g''_i}q_{t+1}(g''_i)\bigr),
\end{multline}
where $174=18(K+1)+8+4$.

\paragraph{Step 3: the upper row and its ideal copy.}
The true pair of gate $g''_j$ is $(x_{g''_j},I^{t+2}_{g''_j})$. Call a step $i>-K$ a
\emph{regeneration} if $i$ is not a near step, both pairs have the value $0$ before it, that is
$y_{i-1}=x_{g''_i-1}=0$, and the position that gate $g''_i$ lists holds a one. A gate of the upper
row at a step that is not near is not a NOT gate, since $H\ge2$; it reads the same pair for the
true and for the ideal upper row, and the true pair changes as in Definition~\ref{def:tworow}, by
the case of equality in Lemma~\ref{lem:rowstatus}. At a regeneration both pairs read a pair $(1,w)$
and become $(0,w)$. After it they stay equal as long as there is no near step. So if $I^{t+2}_{a+j}\ne\hat o_j$, then
there is no regeneration in $(j^*,j]$, where $j^*$ is the last near step up to $j$, and
$j^*=-K$ if there is none; call this event $\mathcal B_j$.

Let $i-1>-K$. If at both steps $i-1$ and $i$ the gate is resampled to a position at most $g'_{i-1}$,
respectively $g'_i$, that holds a one, then both pairs have the value $0$ after step $i-1$, and
step $i$ is a regeneration. Given $\cF_{g''_i-1}$, gate $g''_i$ is resampled to such a position
with probability $\eps\,O_{g'_i}/g''_i=\eps\,\psi_{g'_i}(g'_i+1)/g''_i$, and on $\{m_j=1\}$,
$i\le j$, this is at least $\eps/4$: by~\eqref{eq:rowswin} $\psi_{g'_i}\ge\ell_4-3\delta/2\ge0.33$,
and $g'_i+1\ge(1-1/m_0)\,g''_i$ because $g''_i\ge m_0H$. An interval $(i,j]$ contains
$\lfloor(j-i)/2\rfloor$ disjoint pairs of consecutive steps. So with $\gamma=(\eps/4)^2$ and
$\kappa=(1-\gamma)^{1/2}$, for $-K\le i\le j$, $j\ge1$,
\begin{equation}\label{eq:rowsregen}
m_j\,\PP[\text{no regeneration in }(i,j]\mid\cF_{g''_i}]\le\kappa^{\,j-i-1}.
\end{equation}
A pair of status $0$ keeps it as long as it reads pairs of status $0$. Gate $g''_i$ reads a pair of
status $1$ only if $\pi_{g''_i}=0$ and $o'_i=1$, or if it is a source at level $t+1$. Hence, for
$1\le j\le J$, on $\mathcal B_j$ with $j^*$ as above,
\begin{equation}\label{eq:rowssplit}
I^{t+2}_{a+j}\le I^{t+2}_{b(j^*)}+\sum_{i=\max(j^*,1-K)}^{j}\bigl(o'_i+\ind[g''_i\text{ is a source at level }t+1]\bigr),
\end{equation}
where $b(j^*)=g''_{j^*}-1$ if $j^*>-K$ and $b(-K)=a-K$. The event that $j^*=i$ is not
$\cF_{g''_i}$-measurable; we use only that $\mathcal B_j\cap\{j^*=i\}$ is contained in the event
that $i$ is a near step and there is no regeneration in $(i,j]$, for $i>-K$, and in the event that
there is no regeneration in $(-K,j]$, for $i=-K$, and we add over all $i\le j$. Multiply~\eqref{eq:rowssplit} by $\nu_{a+j}m_j\ind[\mathcal B_j]$,
sum over $j$, and take expectations, term by term. Let $C_\gamma=\kappa^{-1}(1-\kappa)^{-2}$; note $\nu_{a+j}\le\nu_g$ for every position $g\le a+j$.

\emph{Sources.} For a source $g''_i$ and $j\ge i$, the event $\mathcal B_j$ implies that there is
no regeneration in $(i,j]$. By~\eqref{eq:rowsregen} and~\eqref{eq:rowssource} these terms add up
to at most $C_\gamma\,\eps\sum_{i=1-K}^{J}\nu_{g''_i}q_{t+1}(g''_i)$.

\emph{Near steps.} The event that a step $j^*$ is a near step is independent of $\cF_{g''_{j^*}-1}$ and
has probability at most $\eps H/g''_{j^*}\le\eps/m_0$; the variables $m_j$, $I^{t+2}_{g''_{j^*}-1}$
and $o'_i$ are $\cF_{g''_{j^*}-1}$-measurable, because $g''_{j^*}-1\ge a-K$. So by~\eqref{eq:rowsregen} the terms
of~\eqref{eq:rowssplit} with $j^*>-K$ add up to at most
\[
\frac{\eps\,C_\gamma}{m_0}\Bigl(\sum_{i=-K}^{J-1}\nu_{g''_i}u_{t+2}(g''_i)+\sum_{i=1-K}^{J}\nu_{g'_i}u_{t+1}(g'_i)\Bigr),
\]
where we used $\sum_{j\ge j^*}\kappa^{j-j^*-1}\le C_\gamma$ for the first sum and $\sum_{j^*\le i\le j}\kappa^{j-j^*-1}\le C_\gamma$ for the second.

\emph{The start.} The terms with $j^*=-K$ add up to at most
\[
C_\gamma\,\nu_{a-K}u_{t+2}(a-K)+C_\gamma\sum_{i=1-K}^{J}\kappa^{\max(i,1)}\,\nu_{g'_i}u_{t+1}(g'_i),
\]
because $\sum_{j\ge1}\kappa^{j+K-1}\le C_\gamma$ and $\sum_{j\ge\max(i,1)}\kappa^{j+K-1}\le C_\gamma\kappa^{\max(i,1)}$; here $I^{t+2}_{a-K}$ and the $o'_i$ are $\cF_{a-K}$-measurable.

Since $u_{t+2}(a+j)=\EE[m_jI^{t+2}_{a+j}]\le\EE[m_j\hat o_j]+\EE[m_jI^{t+2}_{a+j}\ind[\mathcal B_j]]$,
these three bounds and~\eqref{eq:rowsideal} bound $\sum_{j=1}^{J}\nu_{a+j}u_{t+2}(a+j)$.

\paragraph{Step 4: the recursion.}
Let $H'=H-K\ge K$. For $0\le\varsigma<H'$ the windows with the base points $a=A_0+\varsigma+kH'$,
$k\ge0$, $a<n$, have upper positions $(a,a+J]$ that partition $(A_0+\varsigma,n]$. Add the bounds
of Steps~2 and~3 over these windows and average over $\varsigma$. On the left this gives at least
$\sum_{A_1<g\le n}\nu_gu_{t+2}(g)$. On the right, every position is $g'_i$ for at most two windows
of one $\varsigma$, and the same holds for $g''_i$; and for a fixed $i$, every position is $g'_i$,
respectively $g''_i$, for at most one pair of $\varsigma$ and a window. So the terms that carry
one position per window, or a sum with the weights $\kappa^{\max(i,1)}$, or the $K$ positions
$g'_i$ with $i\le0$, get a factor $1/H'$ in the average. With~\eqref{eq:rowsq} this gives
\begin{equation}\label{eq:rowslaterec}
\sum_{A_1<g\le n}\nu_gu_{t+2}(g)\le\theta\,a_{t+1}+\eta\,(a_{t+2}+a_{t+1}+a_t),\qquad
\eta=C_*\Bigl(\frac1{H'}+\frac1{m_0}+\frac1\lambda\Bigr),
\end{equation}
with a constant $C_*$ that depends on $\eps$ only. In detail, the right side is at most
\begin{multline*}
\theta a_{t+1}+\frac{\bigl(8K+C_\gamma(K+(1-\kappa)^{-1})\bigr)a_{t+1}+C_\gamma a_{t+2}}{H'}\\
+\frac{2\eps C_\gamma\,(a_{t+2}+a_{t+1})}{m_0}+\frac{(348+2C_\gamma)\,\eps\,(a_t+a_{t+1})}{\lambda},
\end{multline*}
so one can take $C_*=8K+348+2C_\gamma\bigl(K+1+(1-\kappa)^{-1}\bigr)$. By~\eqref{eq:rowsearly} with
$m=m_1$, and $A_1\le t_0+m_1H$, for $t+2\ge t_0+m_1$,
\[
\sum_{t_0<g\le A_1}\nu_gu_{t+2}(g)\le\eta'\max_{1\le j\le m_1}a_{t+2-j},\qquad\eta'=\frac{\eps\,m_1(1+m_1K_r)}{\lambda}.
\]
Adding the two, $(1-\eta)\,a_{t+2}\le(\theta+2\eta+\eta')\max_{1\le j\le m_1}a_{t+2-j}$. Here $m_1\ge2$. Choose
$m_0\ge\max(8/\delta,\,10^5C_*)$, then $\lambda\ge\max(10^5C_*,\,10^5m_1(1+m_1K_r))$, and
$H_0\ge2K$ with $H_0-K\ge10^5C_*$. Then $\eta\le3\cdot10^{-5}$ and $\eta'\le10^{-5}$ for
$H\ge H_0$, and
\[
a_t\le\vartheta\max_{1\le j\le m_1}a_{t-j}\quad(t\ge t_0+m_1),\qquad\vartheta:=\frac{\theta+2\eta+\eta'}{1-\eta}<0.9992 .
\]
Since $a_t<1$, this gives $a_t\le\vartheta^{\lfloor(t-t_0)/m_1\rfloor}$ for $t\ge t_0$.

\begin{proof}[Proof of Theorem~\ref{thm:rows}]
Let $t^*=t_0+m_1\bigl(1+\lceil(\lambda+1)\ln(n+1)/\ln(1/\vartheta)\rceil\bigr)$. Then $a_{t^*}\le(n+1)^{-\lambda-1}$ and
$\sum_{i\le n}u_{t^*}(i)\le n^\lambda a_{t^*}\le1/n$. The
events $\{\mathsf T\le n\}$ and the complement of $\Gamma^\circ_n$ have probability at most $1/n$
each, and outside them $\Gamma^*_i$ holds for every $i\le n$. So with probability at least $1-3/n$
no position is open at level $t^*$, and the relevance-pruned iteration has finished after $t^*$
rounds. The constants $K_0,K_1,K_d,m_0,m_1,\lambda,\vartheta$ depend on $\eps$ only, and
$t_0\le(K_0+K_d)\ln(n+1)+K_1+K_d$. The prefix--Jacobi iteration is covered by
Proposition~\ref{prop:jacobi}.
\end{proof}

\begin{remark}[scope]\label{rem:rows}
\begin{enumerate}[label=(\roman*),leftmargin=2.4em]
\item \emph{Positions, not chains.} On these bases the expected number of dependency chains of a
given length grows below noise $0.69$ (the table after Theorem~\ref{thm:fixed}), while the number
of open positions decreases from level to level. The reason is the second case of $f$: a gate
that follows a one and reads a one stays open only if it was open and its read is open. A count of
chains takes the worst case at each gate and cannot use this. The crude count of Step~1 already
shows the difference: the positions that depend on one position through $m$ kept pointers form
an interval of expected length at most $1+mK_r$, while the number of chains to them grows
exponentially in $m$.
\item \emph{What is not covered.} The theorem is about the depth $D$. It says nothing about the
Kleene--prefix iteration, whose depth $d$ has no such rule. It needs $H\ge H_0$, because the two ends of a window of $H-K$ positions cost a bounded
number of positions; for a fixed $H$,
Theorem~\ref{thm:short} gives $O(\log^3n)$ rounds at every positive noise level. And it is about
one family of bases: the proof uses that the gates of a row read the row below in order, so that
two rows form a Markov chain.
\item \emph{The level $0.455$.} The number comes from Lemma~\ref{lem:rowscert}, with potentials
that depend on one symbol and a look back of $K=8$ steps. With these choices the search finds
nothing at noise $0.45$. In the ideal model, with a fixed density $d$ of ones among the resampled
reads, the largest ratio of the expected numbers of open positions in two consecutive rows, over
all periodic rows below of period up to $9$, is $0.958$ at noise $0.455$ for $d=1/3$; it is $0.988$ at noise $0.44$ and $1.029$ at noise
$0.42$. The worst row has two open zeros around a closed one on an
alternating background. So an argument with two rows cannot go far below $0.44$; below that one would have to follow more rows.
In a simulation of the ideal model the open positions decrease by a factor $0.74$ from row to row
at noise $0.46$, and by $0.89$ at noise $0.30$ (Section~\ref{sec:checks}).
\item \emph{Constants.} The theorem is qualitative. The constant $K_d$ of
Lemma~\ref{lem:fddens} for $\tau=\delta/40$ is astronomically large, and $A$ is larger. On
instances with $H=2000$ the largest depth is $42$, $63$, $92$ and $111$ for $40$, $80$, $160$ and
$320$ rows at noise $0.46$ (Section~\ref{sec:checks}).
\end{enumerate}
\end{remark}

\subsection{Three rows: fixed-distance bases below the level of the hardness theorem}\label{sec:rows3}

Theorem~\ref{thm:rows} stops at noise $0.455$ because it follows two rows. With three rows the
same method reaches noise $0.30$. This is below $\eps_S=0.4468$: at noise levels at which the
smoothing of the worst base is as hard as Theorem~\ref{thm:hards24} says, the smoothing of every
fixed-distance base with a large distance is solved in $O(\log n)$ rounds.

\begin{definition}[the three-row chain]\label{def:threerow}
As Definition~\ref{def:tworow}, with a third pair. A state is $(x^1,o^1,x^2,o^2,x^3,o^3)$; it is
closed if $o^1=o^2=o^3=0$. In a step with the symbol $y$ and the densities $d_1,d_2,d_3$, the pair
$(x^1,o^1)$ reads $y$ with probability $1-\eps$, $(1,0)$ with probability $\eps d_1$ and $(0,0)$
with probability $\eps(1-d_1)$; then $(x^2,o^2)$ reads the new pair $(x^1,o^1)$ or $(1,0)$ or
$(0,0)$, with the probabilities $1-\eps$, $\eps d_2$, $\eps(1-d_2)$; then $(x^3,o^3)$ reads the
new pair $(x^2,o^2)$ or $(1,0)$ or $(0,0)$, with the probabilities $1-\eps$, $\eps d_3$,
$\eps(1-d_3)$. In $T_y$ and $T^-_y$ the densities range over $[d_a,d_b]$, and $d_2$ and $d_3$ may
depend on what the pairs before them have read. Here $\phi(s)=o^3-\theta o^2$.
\end{definition}

\begin{lemma}[potentials for three rows]\label{lem:rows3cert}
Let $0.30\le\eps\le0.455$, and let $\delta$, $\theta$, $K$, $\ell_4$, $\mu_4$ be as in
Lemma~\ref{lem:rowscert}. For every $d$ with $\ell_4-\delta\le d\le\mu_4+\delta$ there are
functions $c_y$ on the $64$ states, one for each symbol $y$, with values in $[-4,4]$ and equal to
$0$ on the closed states, such that~\eqref{eq:rowsC1} and~\eqref{eq:rowsC2} hold for the operators
of Definition~\ref{def:threerow} with $[d_a,d_b]=[d-\delta,d+\delta]$ and for every admissible
word $z_1\dots z_Ky$.
\end{lemma}

\begin{proof}
By computation, as for Lemma~\ref{lem:rowscert}. The program \texttt{rows3\_check.c} is
\texttt{rows\_check.c} with the chain of three rows; its check, its arithmetic and its cover are
the same, and its search is different. In its operators the density and the noise level are
maximised separately for each of the three reads and may depend on what the pairs before have
read. The runs
\texttt{rows3\_check 0.30 0.34 0.999 8 0.004 0.016},
\texttt{rows3\_check 0.34 0.38 0.999 8 0.004 0.016} and
\texttt{rows3\_check 0.38 0.455 0.999 8 0.004 0.016} print \texttt{PROVED}, with $218$, $106$ and
$179$ boxes (Section~\ref{sec:checks}).
\end{proof}

\begin{theorem}[fixed-distance bases, three rows]\label{thm:rows3}
Let $0.30\le\eps<1$. There are constants $H_0$, $A$ and $B$ such that for every $H\ge H_0$ and
every $n\ge2$, with probability at least $1-3/n$, the relevance-pruned iteration on
$S_\eps(j^H)$ finishes within $A\ln(n+1)+B$ rounds, and the prefix--Jacobi iteration within one
round more.
\end{theorem}

\noindent For $\eps\ge0.455$ this is Theorem~\ref{thm:rows}. Let $0.30\le\eps\le0.455$. The proof
follows that of Theorem~\ref{thm:rows}, and we say what changes. The new point is that the middle
row also needs an ideal copy, so the comparison of Step~3 is made three times.

\paragraph{Step 0.} The root $\varpi(\eps)$ is at least $\varpi(0.30)>0.19$, so we use the pair
$(0.19,0)$ of Lemma~\ref{fact:densall}: $\mathsf T$ is defined with $0.19$ in place of $1/4$, and
$r_0=1-0.19\,\eps$ in~\eqref{eq:rowsstop}. Three rows span $3H$ positions, so we take
$m_0\ge12/\delta$, $A_0=g_1+(m_0+3)H$, $A_1=A_0+H-K$, $m_1=K_d+m_0+5$, and
$\Gamma^*_i=\{\mathsf T>i\}\cap\Gamma^\circ_{i-3H}$; statement~\eqref{eq:rowslate} holds with
$4H$ in place of $3H$, since $4/m_0\le\delta/3$ and $1/3-\delta-1/13>0.19$. The numbers $u_t$,
$\nu$, $a_t$, $q_t$, the sources, and~\eqref{eq:rowsq},~\eqref{eq:rowssource} are as before.
Step~1 and~\eqref{eq:rowsearly} are unchanged.

\paragraph{Step 2: a window of three rows.}
Fix $t\ge t_0$ and $A_0\le a<n$, and put $s=a-2H-K$, $J=\min(H-K,n-a)$ and
$g^r_j=a-(3-r)H+j$ for $r=1,2,3$ and $-K\le j\le J$. The gates $g^r_j$ with $j>-K$ are
\emph{row $r$}; the rows are disjoint and lie in this order, and gate $g^r_j$ has the base
pointer $g^{r-1}_j$, where $g^0_j=g^1_j-H\le s$. Let
\[
z_j=\bigl(x_{g^0_j},I^t_{g^0_j}\bigr),\qquad(x^1_j,o^1_j)=\bigl(x_{g^1_j},I^{t+1}_{g^1_j}\bigr),\qquad
m_j=\ind\bigl[\Gamma^\circ_{a+j-3H}\cap\{\mathsf T>g_1\}\bigr]\quad(1\le j\le J).
\]
As before these are $\cF_s$-measurable, the words are admissible, and on $\{m_j=1\}$ the
statements~\eqref{eq:rowswin} hold, so that $m_j=\ind[\Gamma^*_{a+j}]$ and
$m_j\le\ind[\Gamma^*_p]$ for $p\le a+j$.

\emph{Ideal rows.} For $r=2,3$ and $-K<j\le J$ let $\beta^r_j$ be bits that, given $\cF_s$, are
independent of each other and of the noise and equal to $1$ with probability $\psi_s$. Call
$(r,j)$ a \emph{near step} if $\pi_{g^r_j}=1$ and $\rho_{g^r_j}>g^1_j$; given the past this has
probability at most $2\eps H/g^r_j\le2\eps/m_0$. Put $(y^r_{-K},\hat o^r_{-K})=(0,0)$ and let
$(y^r_j,\hat o^r_j)$ be the pair $(y^r_{j-1},\hat o^r_{j-1})$ after it has read
\[
\begin{cases}
(x^1_j,o^1_j)\text{ if }r=2,\quad(y^2_j,\hat o^2_j)\text{ if }r=3,&\text{when }\pi_{g^r_j}=0,\\
(x_\rho,I^{t+r-1}_\rho)&\text{when }\pi_{g^r_j}=1\text{ and }\rho:=\rho_{g^r_j}\le g^1_j,\\
(\beta^r_j,0)&\text{at a near step.}
\end{cases}
\]
So the ideal row $3$ reads the ideal row $2$. Let
$S_j=(x^1_j,o^1_j,y^2_j,\hat o^2_j,y^3_j,\hat o^3_j)$, and let $\mathcal G_j$ be generated by
$\cF_{g^1_j}$ and by the noise of $g^2_i,g^3_i$ and the bits $\beta^2_i,\beta^3_i$ with
$-K<i\le j$. Then $S_j$ is $\mathcal G_j$-measurable, and given $\mathcal G_{j-1}$ the noise of
$g^1_j,g^2_j,g^3_j$ and the bits $\beta^2_j,\beta^3_j$ are independent with their unconditional
laws. Call $j$ a source step if $g^1_j$ is a source at level $t$, or if for $r=2$ or $r=3$ gate
$g^r_j$ is resampled to a position $\rho\le g^1_j$ that is open at level $t+r-1$; then $g^r_j$ is
a source at level $t+r-1$. The one-step inequality~\eqref{eq:rowsone} holds for the operators of
Definition~\ref{def:threerow}, with the same proof: the densities are
$d_1=\psi_{g^1_j-1}$ and $d_r=\bigl((g^1_j+1)\psi_{g^1_j}+(g^r_j-g^1_j-1)\psi_s\bigr)/g^r_j$ for
$r=2,3$, and $d_r$ depends only on what the pairs before it have read.

With the functions $c_y$ of Lemma~\ref{lem:rows3cert} for $d=\psi_s$ the computation
before~\eqref{eq:rowsideal0} gives~\eqref{eq:rowsideal0} again, with $\phi=\hat o^3-\theta\hat o^2$
and $\Phi_0\le4\,\EE[o^1_0+\hat o^2_0+\hat o^3_0\mid\cF_s]$. The ideal rows start closed. So
$\hat o^2_i=1$ only if $o^1_{i'}=1$ or $i'$ is a source step for some $-K<i'\le i$, and
$\hat o^3_0=1$ only if $\hat o^2_i=1$ or $i$ is a source step for some $-K<i\le0$. Hence
$\hat o^2_0+\hat o^3_0\le(K+1)\sum_{-K<i\le0}(o^1_i+\ind[i\text{ is a source step}])$, and
\begin{multline}\label{eq:rows3ideal}
\sum_{j=1}^{J}\nu_{a+j}\,\EE[m_j\hat o^3_j]\le\theta\sum_{j=1}^{J}\nu_{a+j}\,\EE[m_j\hat o^2_j]
+4(K+2)\sum_{i=1-K}^{0}\nu_{g^1_i}u_{t+1}(g^1_i)\\
+C_1\,\eps\sum_{i=1-K}^{J}\ \sum_{r=1}^{3}\nu_{g^r_i}\,q_{t+r-1}(g^r_i),\qquad C_1=170+4(K+1).
\end{multline}

\paragraph{Step 3: three comparisons.}
We keep $\gamma=(\eps/4)^2$, $\kappa=(1-\gamma)^{1/2}$, $C_\gamma=\kappa^{-1}(1-\kappa)^{-2}$, and
put $r_1=1-\eps/4$. For a row $r\in\{2,3\}$ call a step $i>-K$ a \emph{far one} of row $r$ if
gate $g^r_i$ is resampled to a position $\rho\le g^1_i$ that holds a one. As in the proof of
Theorem~\ref{thm:rows}, on $\{m_j=1\}$ and for $i\le j$ this has conditional probability at least
$\eps/4$ given the noise of all gates below $g^r_i$ and the bits $\beta$; now
$g^1_i+1\ge(1-2/m_0)\,g^r_i$, and $0.33\,(1-2/m_0)\ge1/4$. We compare two pairs that belong to gates of row $r$, call some steps
\emph{disturbed}, and suppose that the two pairs read the same pair at every step that is not
disturbed. In (A) and (B) one of the pairs is a true pair and the disturbed steps are the near
steps of row $r$; at the other steps a true pair changes as in Definition~\ref{def:tworow}, by
the case of equality in Lemma~\ref{lem:rowstatus}. In (C) the disturbed steps are named there. If the
steps $i-1$ and $i$ are far ones of row $r$ and are not disturbed, then both pairs have the value
$0$ after step $i-1$, both read a pair $(1,w)$ at step $i$ and become $(0,w)$, and they stay equal
until the next disturbed step. So if the two pairs differ after step $j$, then there are no two
consecutive far ones of row $r$ in $(j^*,j]$, where $j^*$ is the last disturbed step up to $j$,
or $j^*=-K$ if there is none. As for~\eqref{eq:rowsregen}, for $-K\le i\le i'\le j$ and $j\ge1$,
\begin{equation}\label{eq:rows3regen}
m_j\,\PP[\text{no two consecutive far ones of row }r\text{ in }(i,i']\mid\mathcal H]\le\kappa^{\,i'-i-1},
\end{equation}
where $\mathcal H$ is generated by $\cF_{g^r_i}$ and the bits $\beta$. A pair that is closed stays
closed as long as it reads pairs of status $0$, and a pair that reads a pair of value $1$ and
status $0$ becomes closed. So if a pair that starts closed is open after step $i'$, and $l\le i'$ is the last step
at which it read a pair of status $1$, then no step in $(l,i']$ is a far one of its row, because
there it would read a one of status $0$; and for $-K<l\le i'\le j$, $j\ge1$,
\begin{equation}\label{eq:rows3alive}
m_j\,\PP[\text{no far one of row }r\text{ in }(l,i']\mid\mathcal H]\le r_1^{\,i'-l},
\end{equation}
with $\mathcal H$ generated by $\cF_{g^r_l}$ and the bits $\beta$. The step $l$ is not a stopping
time, so we use this through the bound: the status after step $i'$ is at most the sum over
$l\le i'$ of the indicator that the pair read a pair of status $1$ at step $l$ and that there is
no far one in $(l,i']$. We write $\Sigma^r_i$ for the
event that $g^r_i$ is a source at level $t+r-1$, and $N^r_i$ for the event that $(r,i)$ is a near
step. In (A), (B) and (C) every expectation that belongs to a step $j$ of a sum over $1\le j\le J$
carries the factor $\nu_{a+j}m_j$, which we do not write. A term that belongs to an earlier step
$i\le j$ is then bounded with $\nu_{a+j}m_j\le\nu_{a+\max(i,1)}m_{\max(i,1)}$, and
$\nu_{a+\max(i,1)}\le\nu_{g^r_i}$ for $r\le3$; so on the right sides $u$ and $q$ at a position
stand for $\nu u$ and $\nu q$ at that position, and $\EE[\hat o^2_i]$ for
$\nu_{a+\max(i,1)}\EE[m_{\max(i,1)}\hat o^2_i]$.

\emph{(A) The ideal row $2$ and the true row $2$.} The true pair of gate $g^2_j$ is
$(x^2_j,o^2_j)=(x_{g^2_j},I^{t+2}_{g^2_j})$. It reads the same pair as the ideal one except at the
near steps of row $2$, and it starts with the pair of position $a-H-K$. Let $\Delta_j$ be the
event that the two pairs differ after step $j$. By the above,
\begin{equation}\label{eq:rows3delta}
\ind[\Delta_j]\le\sum_{-K\le c\le j}\ind[c=-K\text{ or }N^2_c]\cdot\ind[\text{no two consecutive far ones of row }2\text{ in }(c,j]],
\end{equation}
and on the event of the term $c$ the pair $(y^2,\hat o^2)$, which reads $(\beta,0)$ at near steps,
satisfies
\[
\hat o^2_j\le\hat o^2_{c-1}\ind[c>-K]+\sum_{e=\max(c,1-K)}^{j}\bigl(o^1_e+\ind[\Sigma^2_e]\bigr).
\]
The event $N^2_c$ is independent of $\mathcal G_{c-1}$ and of row $1$, and $\hat o^2_{c-1}$ is
$\mathcal G_{c-1}$-measurable. For the sources we use that at most one $c$ is the last disturbed
step, so that these terms are bounded by $\sum_{e\le j}\ind[\Sigma^2_e]\,\ind[\text{no two
consecutive far ones in }(e,j]]$. With~\eqref{eq:rows3regen} this gives, for the sum over
$1\le j\le J$,
\begin{multline}\label{eq:rows3A}
\sum_j\EE[\hat o^2_j\ind[\Delta_j]]\le\frac{2\eps C_\gamma}{m_0}\Bigl(\sum_{i=-K}^{J-1}\EE[\hat o^2_i]+\sum_{i=1-K}^{J}u_{t+1}(g^1_i)\Bigr)
+C_\gamma\sum_{i=1-K}^{J}\kappa^{\max(i,1)}u_{t+1}(g^1_i)\\
+C_\gamma\,\eps\sum_{i=1-K}^{J}q_{t+1}(g^2_i).
\end{multline}
Since $\hat o^2_j\le o^2_j+\hat o^2_j\ind[\Delta_j]$, and $\hat o^2_i\le\sum_{-K<i'\le i}(o^1_{i'}+\ind[\Sigma^2_{i'}])$
for $i\le0$, the number $Y=\sum_{j=1}^{J}\EE[\hat o^2_j]$ satisfies
$Y\le\sum_{j=1}^Ju_{t+2}(g^2_j)+\frac{2\eps C_\gamma}{m_0}Y+E_A$, where $E_A$ is the right side
of~\eqref{eq:rows3A} without the term $Y$, plus
$\frac{2\eps C_\gamma}{m_0}K\sum_{i\le0}(u_{t+1}(g^1_i)+\eps q_{t+1}(g^2_i))$. As
$2\eps C_\gamma/m_0\le1/2$ by the choice of $m_0$ below,
\begin{equation}\label{eq:rows3Y}
\sum_{j=1}^{J}\EE[\hat o^2_j]\le\Bigl(1+\frac{4\eps C_\gamma}{m_0}\Bigr)\sum_{j=1}^{J}u_{t+2}(g^2_j)+2E_A .
\end{equation}
In the same way, with the true pair in place of the ideal one in the bound for the status, and
with the term $o^2_{a-H-K}$ for $c=-K$, where $o^2_{a-H-K}:=I^{t+2}_{a-H-K}$,
\begin{multline}\label{eq:rows3A2}
\EE\bigl[\ind[\Delta_i]\,o^2_l\bigr]\le\frac{2\eps}{m_0}\sum_{c\le i}\kappa^{\,i-c-1}\Bigl(u_{t+2}(g^2_{\min(c-1,l)})+\sum_{e=c}^{l}u_{t+1}(g^1_e)\Bigr)\\
+\kappa^{\,i+K-1}\Bigl(u_{t+2}(a-H-K)+\sum_{e=1-K}^{l}u_{t+1}(g^1_e)\Bigr)
+\eps\sum_{-K<e\le l}\kappa^{\max(i-e-1,0)}q_{t+1}(g^2_e)
\end{multline}
for $-K<l$ and $-K<i$, where the first sum is over $-K<c\le i$. For $c>l$ the event $N^2_c$ is independent of
$\cF_{g^2_c-1}$, which contains $o^2_l$. For $c\le l$ we used
$o^2_l\le o^2_{c-1}+\sum_{e=c}^{l}(o^1_e+\ind[\Sigma^2_e])$, the independence of $N^2_c$ from
$\cF_{g^2_c-1}$, which contains $o^2_{c-1}$ and row $1$, and for the sources that the last
disturbed step is unique, which gives the factor $\kappa^{\,i-e-1}$ for $e\le i$ and no factor
for $i<e\le l$.

\emph{(B) The true row $3$ and a row $3$ that reads the true row $2$.} Let
$(\tilde y_j,\tilde o_j)$ be defined as the ideal row $3$, except that at a step with
$\pi_{g^3_j}=0$ it reads the true pair $(x^2_j,o^2_j)$. The true pair of gate $g^3_j$ is
$(x_{a+j},I^{t+3}_{a+j})$; it reads the same pairs as $(\tilde y,\tilde o)$ except at the near
steps of row $3$, and row $2$ is $\cF_{a-K}$-measurable. This is the situation of Step~3 of the
proof of Theorem~\ref{thm:rows}, with the true row $2$ as the lower row, and that step gives
\begin{multline}\label{eq:rows3B}
\sum_{j=1}^{J}u_{t+3}(a+j)\le\sum_{j=1}^{J}\EE[\tilde o_j]+C_\gamma\,\eps\sum_{i=1-K}^{J}q_{t+2}(g^3_i)
+\frac{2\eps C_\gamma}{m_0}\Bigl(\sum_{i=-K}^{J-1}u_{t+3}(g^3_i)+\sum_{i=1-K}^{J}u_{t+2}(g^2_i)\Bigr)\\
+C_\gamma\,u_{t+3}(a-K)+C_\gamma\sum_{i=1-K}^{J}\kappa^{\max(i,1)}u_{t+2}(g^2_i).
\end{multline}

\emph{(C) The two rows $3$.} The pairs $(\tilde y,\tilde o)$ and $(y^3,\hat o^3)$ start equal and
read the same pairs except at the steps $i$ with $\pi_{g^3_i}=0$ at which $\Delta_i$ holds; call
these disturbed. Let $\mathcal F'$ be generated by $\cF_{a-K}$ and the bits $\beta^2$; the events
$\Delta_i$ are $\mathcal F'$-measurable, and given $\mathcal F'$ the noise of row $3$ and the bits
$\beta^3$ have their unconditional laws. If the two pairs differ after step $j$, let $i$ be the
last disturbed step up to $j$; then there are no two consecutive far ones of row $3$ in $(i,j]$,
and
\[
\tilde o_j\le\tilde o_{i-1}+\sum_{l=i}^{j}\bigl(o^2_l+\ind[\Sigma^3_l]\bigr).
\]
By~\eqref{eq:rows3alive} and the bound after it, applied to the pair $(\tilde y,\tilde o)$, with
$\mathcal H\supseteq\mathcal F'$,
$\EE[\tilde o_{i-1}\mid\mathcal F']\le\sum_{-K<l<i}r_1^{\,i-1-l}\bigl(o^2_l+\PP[\Sigma^3_l\mid\mathcal F']\bigr)$.
For the sources at steps $l<i$ we drop the indicator of the disturbed step and add
$r_1^{\,i-1-l}\kappa^{\,j-i-1}$ over $i$ and $j$; for the sources at steps $l\in[i,j]$ we use that
the last disturbed step is unique. So
by~\eqref{eq:rows3regen}, with $\tilde o_j\le\hat o^3_j+\tilde o_j\ind[\text{the pairs differ after }j]$,
\begin{multline}\label{eq:rows3C}
\sum_{j=1}^{J}\EE[\tilde o_j]\le\sum_{j=1}^{J}\EE[\hat o^3_j]
+\sum_{j=1}^{J}\ \sum_{-K<i\le j}\kappa^{\,j-i-1}\Bigl(\sum_{-K<l<i}r_1^{\,i-1-l}\,\EE[\ind[\Delta_i]o^2_l]+\sum_{l=i}^{j}\EE[\ind[\Delta_i]o^2_l]\Bigr)\\
+\frac{2C_\gamma}{1-r_1}\,\eps\sum_{i=1-K}^{J}q_{t+2}(g^3_i).
\end{multline}
Insert~\eqref{eq:rows3A2}. Every term is a number $u_{t+2}$ or $u_{t+1}$ at a position of row $2$
or row $1$, or a number $\eps q_{t+1}$ at a position of row $2$, times a product of powers of
$\kappa$ and $r_1$ whose exponents add up, over the other indices, to a convergent multiple
geometric sum; in some of them a factor that is linear in a difference of indices appears, from
the number of $l$ between $i$ and $j$. Let $C_2=C_\gamma^3(1-r_1)^{-2}$. The terms with a near step $c>-K$ carry the
factor $2\eps/m_0$ and add up to at most
$\frac{2\eps C_2}{m_0}\bigl(\sum_{i=-K}^{J}u_{t+2}(g^2_i)+\sum_{i=1-K}^{J}u_{t+1}(g^1_i)\bigr)$.
The terms with $c=-K$ carry $\kappa^{\,i+K-1}$ with $i\le j$ and $j\ge1$, so they add up to at
most $C_2\,u_{t+2}(a-H-K)+C_2\sum_{e=1-K}^{J}\kappa^{\max(e,1)/2}u_{t+1}(g^1_e)$. The terms with
a source add up to at most $C_2\,\eps\sum_{e}q_{t+1}(g^2_e)$.

\paragraph{Step 4.} Combine~\eqref{eq:rows3B},~\eqref{eq:rows3C},~\eqref{eq:rows3ideal}
and~\eqref{eq:rows3Y}, add over the windows of one shift $\varsigma$, and average over
$\varsigma$, as in Step~4 of the proof of Theorem~\ref{thm:rows}. Every position is $g^r_i$ for
at most three windows of one $\varsigma$; the terms with one position per window, with the $K$
positions $i\le0$, or with weights $\kappa^{\max(i,1)}$ or $\kappa^{\max(i,1)/2}$, get the factor
$1/H'$. This gives, with a constant $C_*$ that depends on $\eps$ only,
\[
\sum_{A_1<g\le n}\nu_gu_{t+3}(g)\le\theta\Bigl(1+\frac{4\eps C_\gamma}{m_0}\Bigr)a_{t+2}
+\eta\,(a_{t+3}+a_{t+2}+a_{t+1}+a_t),\qquad\eta=C_*\Bigl(\frac1{H'}+\frac1{m_0}+\frac1\lambda\Bigr).
\]
Use~\eqref{eq:rowsearly} with $m=m_1$ at the level $t+3\ge t_0+m_1$, which implies $t\ge t_0$,
and choose $m_0$, $\lambda$ and $H_0$ as before, with
$m_0\ge\max(12/\delta,10^5C_*,10^5\eps C_\gamma)$. Since $m_1\ge3$, the numbers $a_{t+2}$,
$a_{t+1}$, $a_t$ are among the $a_{t+3-j}$ with $1\le j\le m_1$, and
\[
(1-\eta)\,a_{t+3}\le\bigl(\theta(1+4\cdot10^{-5})+3\eta+\eta'\bigr)\max_{1\le j\le m_1}a_{t+3-j},
\]
so $a_t\le\vartheta\max_{1\le j\le m_1}a_{t-j}$ for $t\ge t_0+m_1$ with $\vartheta<0.9993$. The
proof ends as that of Theorem~\ref{thm:rows}.

\begin{remark}[what the three rows give]\label{rem:rows3}
\begin{enumerate}[label=(\roman*),leftmargin=2.4em]
\item \emph{Below the hardness level.} For $0.30\le\eps\le\eps_S$ a randomized NC algorithm for
the $\eps$-smoothing of every base would give $\mathrm P\subseteq\mathrm{BPNC}$
(Theorem~\ref{thm:hards24}), while the smoothing of every fixed-distance base with $H\ge H_0$ is
solved in $O(\log n)$ rounds by an iteration that is never wrong. So the hardness at these noise
levels is a property of particular bases, such as the layouts of
Sections~\ref{sec:latch} and~\ref{sec:s24}, and not of long pointers as such.
\item \emph{More rows.} In the ideal model with $d=0.3345$ the largest ratio of the open positions
of the top row to those of the row below, over the periodic rows of period up to $8$, is
$0.946$, $0.984$, $1.019$ and $1.075$ at noise $0.30$, $0.25$, $0.20$ and $0.10$ for three rows,
and $0.966$ and $1.001$ at noise $0.20$ and $0.10$ for four rows (Section~\ref{sec:checks}). So
each further row lowers the level that the method can reach. Near $0.30$ the margin of three
rows is thin, and the run for $[0.30,0.34]$ needs twice as many boxes as that for $[0.34,0.38]$.
We have not tried to prove a bound
for every positive noise level in this way; for a fixed $H$ that is Theorem~\ref{thm:short}, with
$O(\log^3n)$ rounds.
\item \emph{Several distances.} The same computation for a base with two distances, in which the
upper row reads two lower rows in turn, gives ratios above $1$ at noise $0.46$ for two rows
(Section~\ref{sec:checks}): an open run of a lower row is split between two upper rows, and
positions that are read with gaps each start a run of their own. So the method of this subsection
does not extend to several distances as it stands.
\end{enumerate}
\end{remark}

\section{All bases: at noise \texorpdfstring{$0.72$}{0.72} the expected number of chains can grow}\label{sec:barrier}

The bounds at the upper end control an \emph{expected number} of dependency chains of length
$\ell$. Lemma~\ref{lem:kernel}, Lemma~\ref{lem:typed} and the chain count of Theorem~\ref{kp:thm34}
bound the expected number of all dependency chains by $\theta^\ell$ times a polynomial in $n$ and
need $\theta<1$; Theorem~\ref{thm:ordered} counts fewer chains, the top-reading ones. This section
shows that at noise $0.70$ and at noise $0.72$ no bound of the first kind holds for all bases:
there is a base fixed in advance on which the expected number of all chains grows exponentially
(Proposition~\ref{prop:barrier}). On a base with the same pattern the iterations still finish in
few rounds in simulation. So what fails there is this count, not the iterations. The section then
computes, in a model that can be solved exactly on repeating bases, where three counts reach $1$:
the count of all chains, a count of the chains that satisfy a local consequence of being canonical,
and a count that uses only relevant gates and the last reset.

\subsection{The density of ones at every noise level}

\begin{lemma}\label{lem:dens2any}
Lemma~\ref{lem:dens2} holds for every $\eps\in(0,1)$ and every integer $k\ge3$, provided
$1/(1+\rho)+\tau\le1/4$, with $\eps^5/2$ replaced by $\eps^{2k-1}/2$ and the condition on
$\mathrm{pre}$ replaced by $\mathrm{pre}\ge(5/\tau^2)^{2k+1}W$.
\end{lemma}

\begin{proof}
The proof of Lemma~\ref{lem:dens2} uses $\eps\le1/4$ to make $\eps^5/2$ small and to keep the
arguments of $\Lambda$ at most $1$. Run the stages up to $u_{2k}$ instead of $u_6$. The comparison with the starred sequence, which approaches
$\varpi$ by a factor $\eps$ per step, gives $\psi\ge\ell_{2k-1}\ge\varpi-\eps^{2k-1}/2-\eta/(1-\eps)$
on $[Q_{2k-1},n]$ and $\psi\le u_{2k}\le\varpi+\eps^{2k}/2+\eta/(1-\eps)$ on $[Q_{2k},n]$, and the
same induction gives $Q_{2k}\le(5/\tau^2)^{2k+1}W$. Since $\eta=1/(1+\rho)+\tau\le1/4$, the
arguments of $\Lambda$ stay in $[-\eta,1]$: every $\ell_k$ is at least $\Lambda(1)-\eta=-\eta$, and
every $u_k$ is at most $\Lambda(-1/4)+1/4\le5/9+1/4$. On $(-\infty,1]$ we have $|\Lambda'|\le\eps$.
\end{proof}

\subsection{A base on which the count grows}

Let $n$ be even, $A=\lfloor n^{1/4}\rfloor$, $n_0=n-A$ and $b=n_0+8$. The base $j^{\mathrm{bar}}$ has
$j_i=0$ for $i\le n_0+3$ and, for $n_0+3<i\le n$,
\[
j_i=i-2\quad(i\text{ even}),\qquad j_i=i-3\quad(i\text{ odd}).
\]
So every gate of the \emph{block} $(n_0+3,n]$ reads an even position two or three below it, and
all other gates are zero gates. For a gate $q$ of the block write $d_q=q-j_q\in\{2,3\}$. Let
$\jt=S_\eps(j^{\mathrm{bar}})$, and let $N_t(q)$ be the number of dependency chains
$(P_0,k_1,P_1,\dots,k_t,P_t)$ of $\jt$ with $P_0=q$, every $k_s\ge b$, and $\pi_{k_s}=0$ for every
$s$: every step keeps its base pointer. Put $N_0(q)=1$.

\begin{proposition}[a base on which the count grows]\label{prop:barrier}
Let $\delta=1/2000$, and let $(\eps,\mu)$ be one of the pairs $(7/10,113/100)$ and
$(18/25,101/100)$. There are $\kappa>0$ and $n_1$ such that for every even $n\ge n_1$ there is an
event $H$ with
$\PP[H^c]\le(n+1)^{-A}$, on which $|\psi_g-\varpi(\eps)|\le\delta$ for all
$g\in[n_0+3,n]$, and
\[
\EE\bigl[N_t(n)\,\mathbf 1_H\bigr]\ \ge\ \kappa\,\mu^t\qquad\text{for every }t\le A/80-1 .
\]
\end{proposition}

\begin{proof}
\emph{Density.} Apply Lemma~\ref{lem:dens2any} with the block as the set of cell gates and
$\mathrm{pre}=n_0+3$: every interval contains at most $A$ cell gates, so the hypothesis holds with
$c=A$ and any $\rho$. Take $\rho=10^6$, $\tau=1/20000$, $k=13$ and $K=A$. Then
$W=A+3+(A+2)\ln(n+1)$, so $\mathrm{pre}\ge(5/\tau^2)^{27}W$ for large $n$, and the error is at most
$(10^{-6}+1/20000)/(1-\eps)+\eps^{25}/2<3.2\cdot10^{-4}<\delta$ for $\eps\le18/25$. For
$q\in[n_0+4,n+1]$ let $H_q$ be the event that $|\psi_g-\varpi|\le\delta$ for all
$g\in[n_0+3,q-1]$, and let $H=H_{n+1}$. Then $H_q\in\cF_{q-1}$, $H_q\subseteq H_{q-1}$, and $H$
contains the event of the lemma, so $\PP[H^c]\le(n+1)^{-A}$.

\emph{One gate.} Let $q\ge b$, let $S_q=(x_{q-2},x_{q-1},x_q)$, and $p_q=\psi_{q-1}=O_{q-1}/q$.
Given $\cF_{q-1}$, gate $q$ keeps its pointer with probability $1-\eps$; it then reads
$v=x_{q-d_q}$, a coordinate of $S_{q-1}$, is a reset if and only if $v=1$, and outputs
$\NOR(x_{q-1},v)$. If it is resampled and $x_{q-1}=1$, it outputs $0$ and is a reset with probability
$(O_{q-1}-1)/q\le p_q$. If it is resampled and $x_{q-1}=0$, it is a reset and outputs $0$ with
probability $p_q$, and otherwise outputs $1$ and is not a reset. A gate that is not a reset
outputs $1-x_{q-1}$. On $H_q$ we have $p_q\in[p_-,p_+]$ with $p_\pm=\varpi\pm\delta$. For a state
$s'=(x_{q-3},x_{q-2},x_{q-1})$ and a value $y$ write $s'y=(x_{q-2},x_{q-1},y)$, and let $v$ be the
coordinate $x_{q-d_q}$ of $s'$. Define $8\times8$ matrices by
\begin{align*}
\mathrm{NR}_q(s',s'(1-x_{q-1}))&=(1-\eps)\mathbf 1[v=0]+\eps(1-p_+),\\
\mathrm{KP}_q(s',s'y)&=1-\eps\quad\text{for }y=\NOR(x_{q-1},v),\\
\mathrm{FU}_q(s',s'0)&=1\quad\text{if }x_{q-1}=1,\\
\mathrm{FU}_q(s',s'0)&=(1-\eps)\mathbf 1[v=1]+\eps p_-\quad\text{if }x_{q-1}=0,\\
\mathrm{FU}_q(s',s'1)&=(1-\eps)\mathbf 1[v=0]+\eps(1-p_+)\quad\text{if }x_{q-1}=0,
\end{align*}
and all other entries $0$. On $H_q$ and on $\{S_{q-1}=s'\}$, the conditional probability given
$\cF_{q-1}$ that $S_q=s$ and $q$ is not a reset is at least $\mathrm{NR}_q(s',s)$, that $S_q=s$ is
at least $\mathrm{FU}_q(s',s)$, and that $S_q=s$ and $\pi_q=0$ equals $\mathrm{KP}_q(s',s)$. The
matrices depend only on the parity of $q$, and their rows have sums at most $1$.

\emph{Recursion.} For $q\ge b$ and $t\ge1$, the first gate of a counted chain from $q$ is either $q$
itself, with $\pi_q=0$, or lies below $q$, and then $q$ is not a reset. Hence, with $N_t(b-1)=0$,
\[
N_t(q)=\mathbf 1[\pi_q=0]\,N_{t-1}(q-d_q)+\mathbf 1[q\text{ is not a reset}]\,N_t(q-1),
\]
and $N_t(q)\le A^t$, since each step chooses one of at most $A$ gates. For each state $s$ let
\[
V_t(q,s)=\EE\bigl[N_t(q)\,\mathbf 1_{H_q}\mathbf 1[S_q=s]\bigr],
\]
a row vector in $s$. Below, ``$V\ge U-\xi$
with error $E$'' means that $\xi\ge0$ and $\|\xi\|_1\le E$; multiplying by a matrix whose rows sum
to at most $1$ does not increase $\|\xi\|_1$. Condition on $\cF_{q-1}$. Since $H_q\in\cF_{q-1}$ and
$H_{q-1}\setminus H_q\subseteq H^c$, replacing $\mathbf 1_{H_q}$ by $\mathbf 1_{H_{q-1}}$ costs at
most $A^t\PP[H^c]$ per state. So $V_t(q)\ge V_t(q-1)\,\mathrm{NR}_q-\xi$ with error $8e_t$,
where $e_t=A^t(n+1)^{-A}$. In the same way the first term, conditioned successively on
$\cF_{q-1},\dots,\cF_{q-d_q}$, contributes at least
$V_{t-1}(q-d_q)\,\mathrm{FU}_{q-d_q+1}\cdots\mathrm{FU}_{q-1}\mathrm{KP}_q$ with error $24e_t$.
Unfolding the window $70$ times and dropping the remainder, which is nonnegative,
\[
V_t(q)\ \ge\ \sum_{m=0}^{70}V_{t-1}(k-d_k)\,\mathrm{Push}_k\,\mathrm{NR}_{k+1}\cdots\mathrm{NR}_q
\;-\;\xi\quad\text{with error }2300\,e_t,\qquad k=q-m,
\]
where $\mathrm{Push}_k=\mathrm{FU}_{k-d_k+1}\cdots\mathrm{FU}_{k-1}\mathrm{KP}_k$, provided all the
gates involved are at least $b$.

\emph{Certificate.} Every target $k-d_k$ is even. The script \texttt{barrier\_check.py} computes a
row vector $v\ge0$, positive exactly on the states $001,010,100,101$ (written
$x_{q-2}x_{q-1}x_q$), and checks in exact rational arithmetic that for even $q$
\[
\sum_{m=0}^{70}v\,\mathrm{Push}_{q-m}\,\mathrm{NR}_{q-m+1}\cdots\mathrm{NR}_q\ \ge\ \mu\,v
\]
entrywise. The smallest ratio is $1.1355$ at $\eps=7/10$ and $1.0173$ at $\eps=18/25$.

\emph{Induction.} Let $F_q=\mathrm{FU}_{q-3}\mathrm{FU}_{q-2}\mathrm{FU}_{q-1}\mathrm{FU}_q$. For $s$
in the support of $v$ and every state $s'$, the entry $F_q(s',s)$ is positive, since every string
without two adjacent ones is reachable in four steps. So, with $\kappa_0>0$ the smallest ratio of
these entries to $v(s)$, $V_0(q)\ge\kappa_0v-\xi$ with error $E_0=32(n+1)^{-A}$, for $q\ge b+4$. By induction on $t$, for even $q\ge b+80t+4$,
\[
V_t(q)\ \ge\ \kappa_0\,\mu^t\,v-\xi\quad\text{with error }E_t,\qquad E_t\le71\,E_{t-1}+2300\,e_t .
\]
Indeed the targets $k-d_k\ge q-73$ are even and at least $b+80(t-1)+4$, and each of the $71$
matrices $\mathrm{Push}_k\mathrm{NR}_{k+1}\cdots\mathrm{NR}_q$ has rows summing to at most $1$. Hence
$E_t\le71^t(32+2300t)A^t(n+1)^{-A}$, which tends to $0$ uniformly in $t\le A/80$, because
$(A/80)\ln(71A)\le A\ln(n+1)/2$ for large $n$. At $q=n$, where $n\ge b+80t+4$ for $t\le A/80-1$,
summing over $s$ gives $\EE[N_t(n)\mathbf 1_{H_n}]\ge\kappa_0\mu^t\|v\|_1-E_t$. Finally
$\EE[N_t(n)\mathbf 1_H]\ge\EE[N_t(n)\mathbf 1_{H_n}]-A^t\PP[H^c]$, and $\kappa=\kappa_0\|v\|_1/2$
works for large $n$.
\end{proof}

\begin{corollary}[no bound on the number of all chains, for all bases, at noise $0.72$]\label{cor:barrier}
Let $\eps\in\{7/10,18/25\}$, $\theta<1$ and $\lambda\ge0$. For large even $n$, the base
$j^{\mathrm{bar}}$ of length $n$ violates
\[
\sum_\sigma\PP[E_\sigma\cap H]\ \le\ \theta^\ell(n+1)^\lambda ,
\]
where the sum is over the skeletons of length $\ell$ from $n$ with all positions at least
$n_0+5$, and $\ell=\lceil(\lambda+1)\ln(n+1)/\ln(\mu/\theta)\rceil$. On $H$ the conditions $G_g$
of Definition~\ref{def:admissible} hold for all $g>n_0+5$ when $c<\varpi-\delta$ and
$\beta<1-\varpi-\delta$ (and $n$ is large), so Lemma~\ref{lem:onechain}, whose proof uses only the
events $G_g$ at the gates of the chain, holds with $H$ in place of $G$. So for this base the
conclusion of Lemma~\ref{lem:kernel} with $H$ in place of $G$ is false for every $\theta<1$.
Consequently, at these two noise levels, no bound on the probability of single chains that is
valid on every event contained in the events $G_g$, as that of Lemma~\ref{lem:onechain} is, can be
summed over \emph{all} chains to give $O(\log n)$ rounds for all bases. This holds whatever the
bound for one chain uses, including the values that a chain reads. It does not exclude a sum over
a smaller class of chains, such as the top-reading chains of Theorem~\ref{thm:ordered}, or a bound
on a smaller event.
\end{corollary}

\begin{proof}
Every chain counted by $N_\ell(n)$ is a dependency chain whose positions are at least $n_0+5$,
because $P_s=k_s-d_{k_s}\ge b-3$. So the sum is at least
$\EE[N_\ell(n)\mathbf 1_H]\ge\kappa\mu^\ell$, and $\kappa\mu^\ell>\theta^\ell(n+1)^\lambda$ by the
choice of $\ell$ once $n$ is large; also $\ell\le A/80-1$ for large $n$.
\end{proof}

\noindent The pattern of $j^{\mathrm{bar}}$ is the worst repeating pattern found by the
computations below. In simulation, on a base in which this pattern fills the last tenth of the
positions and the other gates are zero gates, the depths grow like $\log n$, as for a harmless
base. (In $j^{\mathrm{bar}}$ itself the block has only $n^{1/4}$ gates.) There is one smoothing per
entry (script \texttt{parity\_depths.c}). Each entry gives $\max_id(i)$, which is the number of rounds of the
Kleene--prefix iteration, and then the number of rounds of the prefix--Jacobi iteration:

\begin{center}
\begin{tabular}{lrrr}
\toprule
active part, noise & $n=10^5$ & $10^6$ & $10^7$\\
\midrule
$j^{\mathrm{bar}}$, $0.50$ & 55 / 25 & 72 / 25 & 88 / 35\\
$j^{\mathrm{bar}}$, $0.60$ & 48 / 20 & 53 / 25 & 64 / 29\\
$j^{\mathrm{bar}}$, $0.70$ & 39 / 18 & 44 / 26 & 57 / 31\\
NOT gates ($j_i=i-1$), $0.70$ & 28 / 18 & 34 / 22 & 42 / 27\\
\bottomrule
\end{tabular}
\end{center}

\subsection{Exact counts on repeating bases}

Let the active part repeat with period $p$: gate $q$ reads $q-d_{q\bmod p}$, where $d=1$ is a NOT
gate. In the model in which a resampled gate reads a one with probability exactly $\varpi(\eps)$,
independently of everything else, the last $\max d$ values form a Markov chain. The recursion of
the proof above then holds with equalities, and the growth rate of the expected number of chains
is the spectral radius of a finite operator. Proposition~\ref{prop:barrier} is the rigorous form of
one such computation, with the probability known only to lie in a band. The script
\texttt{first\_moment.c} computes three counts.
\begin{itemize}[leftmargin=2.2em]
\item The \emph{Kleene--prefix count}: all dependency chains whose steps keep the base pointer,
as in Proposition~\ref{prop:barrier}. (Lemma~\ref{lem:kernel} also counts the chains with resampled
steps.)
\item The \emph{locally canonical count} (script \texttt{canonical.c}). At every step choose,
among the gates of the run whose targets start a chain of the required length, the lowest target
and then the highest gate. This gives one \emph{canonical} chain from every position that starts a
chain. A kept gate of the step's window that read the step's target, or a position of the next
window, would itself be such a choice: from a position $\tau'$ in the next window
$[k_{s+1},P_s]$ the chain can continue with $k_{s+1}$. So in a canonical chain such gates are
resampled. The count is of all chains with this local property. It is an upper bound for the
number of canonical chains, which is at most one per position, and it can grow. On the base
$(2,3)$ the local property is the same as being top-reading in the sense of
Theorem~\ref{thm:ordered}.
\item The \emph{$D_{\mathrm{last}}$ count}: a step from a position $i$ is taken only at a relevant
gate in $(\varrho(i),i]$ that is not a NOT gate, or at the last reset $\varrho(i)$. Taking the anchor
$r=\varrho(i)$ in the formula for $D$ of Section~\ref{sec:upper} gives
$D(i)\le1+\max\bigl(\{D(j_{\varrho(i)}):\varrho(i)\ge1\}\cup\{D(j_k):\varrho(i)<k\le i,\ k\text{
relevant and not a NOT gate}\}\bigr)$. So $\max_iD(i)$ is at most one more than the length of the longest chain
of this form, and this count bounds the rounds of the relevance-pruned iteration and of the
prefix--Jacobi iteration (Lemma~\ref{fact:depths} and Proposition~\ref{prop:jacobi}).
\end{itemize}
In the table, $(a,b)$ is the base with $j_q=q-a$ for even $q$ and $j_q=q-b$ for odd $q$, so
$j^{\mathrm{bar}}$ is $(2,3)$. The last column gives the noise level at which the growth rate
equals $1$.

\begin{center}
\begin{tabular}{llrrrr}
\toprule
base & count & $\eps=0.66$ & $0.70$ & $0.74$ & equals $1$ at\\
\midrule
$(2,3)$ & Kleene--prefix & 1.411 & 1.138 & 0.911 & 0.723\\
$(2,3)$ & locally canonical & 1.221 & 1.009 & 0.826 & 0.702\\
$(2,3)$ & $D_{\mathrm{last}}$ & 0.764 & 0.643 & 0.537 & 0.599\\
$(3,17)$ & $D_{\mathrm{last}}$ & 0.947 & 0.826 & 0.705 & 0.642\\
$j_i=i-3$ & Kleene--prefix & 0.921 & 0.819 & 0.712 & 0.628\\
$j_i=i-3$ & $D_{\mathrm{last}}$ & 0.878 & 0.774 & 0.668 & 0.611\\
bound of Thm.~\ref{kp:thm34} & & 1.650 & 1.331 & 1.061 & 0.750\\
\bottomrule
\end{tabular}
\end{center}

\noindent The search covered all bases of period at most $4$ with pointers at most $7$ back,
$1500$ random bases of period $6$ or $8$ with pointers at most $9$ back, and the families
$(2,b)$ and $(3,b)$ for $b\le19$. For the Kleene--prefix count, $(2,3)$ was the largest at noise
$0.74$. With pointers at most $5$ back, it was also the largest for the locally canonical count. The
$D_{\mathrm{last}}$ count was searched with period $2$ up to $12$ back and period $3$ up to $8$
back. Its largest values come from the family $(3,b)$, odd $b$, and increase with $b$: the rate
equals $1$ at $0.611$, $0.630$ and $0.642$ for $b=3,11,17$. At noise $0.66$ it is
$0.947,0.954,0.960,0.965$ for $b=17,19,21,23$, with increments shrinking by a factor of about
$0.82$, which suggests a limit near $0.99$. The sampled rates of
Remark~\ref{rem:needs}(iii) agree with these exact values where the sample is reliable: for
$j_i=i-3$, $0.82$ against $0.819$ at noise $0.70$. On $(2,3)$ the same kind of estimate, taken
over lengths $4$ to $16$ at $n=4\cdot10^6$ with three smoothings, gave $0.91$, $0.94$ and $1.10$ at
noise $0.66$, where the exact value is $1.41$, and $0.71$, $0.85$ and $0.97$ at noise $0.70$, where it
is $1.14$. It underestimates because the expectation is carried by rare long runs.

\begin{remark}[what this means for the upper end]\label{rem:barrier}
\begin{enumerate}[label=(\roman*),leftmargin=2.4em]
\item At noise $0.70$ and $0.72$, a proof that covers all bases cannot bound the expected number
of all chains on an event like $H$. It has to count fewer chains, or work on a smaller event, or
bound directly the probability that a long chain \emph{exists}. On the pattern $(2,3)$ the chains
overlap: from a typical position many chains share their lower parts. So the count multiplies,
while the longest chain stays logarithmic in the simulation above; Theorem~\ref{thm:short} proves
that it is polylogarithmic.
\item Two ways of counting fewer chains are available. One is the locally canonical count, which
on $(2,3)$ is the count of top-reading chains used in Theorem~\ref{thm:ordered}. In the model it
moves the crossing on $(2,3)$ only from $0.723$ to $0.702$: at noise $0.72$ it is $0.915$, so
$j^{\mathrm{bar}}$ does not block that method there, and at $0.70$ it is $1.009$, which no
certificate confirms. The other is the $D_{\mathrm{last}}$ count, which suffices for the
relevance-pruned and prefix--Jacobi iterations. On every repeating base tried it stays below $1$
down to about $0.66$, against $0.723$ for the Kleene--prefix count. Among the first-moment routes
below $3/4$ for all bases, a bound on the $D_{\mathrm{last}}$ count has the most room;
Section~\ref{sec:allbases} carries out a bound of this kind for all bases, down to noise $0.7288$.
\item Such a proof still meets the obstruction of Remark~\ref{rem:needs}(ii): the window gates of a
$D_{\mathrm{last}}$ chain can read positions below the next target, and the
$(A,B)$-inequalities of that remark, written for this count, still allow the growth rate
$(1-\eps)/(\eps c)$. The computations only show that no repeating base realizes the obstruction
strongly enough to push this count above $1$ at noise above $0.66$.
\item Proposition~\ref{prop:barrier} is a statement about a method, like
Proposition~\ref{prop:adaptive}. It says nothing about the critical noise level. The computations
of this subsection are exact only in the model of the resampled reads, and they are evidence only
about the bases searched.
\end{enumerate}
\end{remark}

\subsection{Short pointers: the iterations are fast at every positive noise level}\label{sec:short}

The base $j^{\mathrm{bar}}$ has only short pointers and zero gates. For every such base the
iterations need polylogarithmically many rounds at every positive noise level, although the count
of chains may grow. The proof does
not count chains. It uses cuts: a few consecutive resampled gates, the lowest of them a reset,
that no step with a base pointer can cross.

\begin{theorem}[short pointers]\label{thm:short}
Let $0<\eps\le1$ and $R\ge2$. There is a constant $C$ such that for every $n\ge2$ and every base
of length $n$ in which every gate $i$ has $j_i\ge i-R$ or $j_i=0$, the Kleene--prefix iteration on
the $\eps$-smoothing finishes within $C\ln^3(n+1)$ rounds with probability at least $1-3/n$. The
same holds for the relevance-pruned iteration and, with one more round, for the prefix--Jacobi
iteration.
\end{theorem}

\begin{proof}
Fix $c\in(0,\varpi(\eps))$ and let $g_0$ and $G\subseteq\bigcap_{g>g_0}G_g$ be as in
Lemma~\ref{fact:densall}, so that $\PP[G^c]\le1/n$, $g_0=O(\ln n)$, and $(O_{g-1}-1)/g\ge c$ on
$G_g\in\cF_{g-1}$.

\emph{Cuts.} Call a position $a$ with $g_0<a\le n-R+1$ a \emph{cut} if the gates
$a,\dots,a+R-1$ are all resampled and gate $a$ is a reset. Given $\cF_{a-1}$, gate $a$ is resampled
to a one below $a-1$ with probability $\eps(O_{a-1}-x_{a-1})/a\ge\eps c$ on $G_a$, and the other
$R-1$ gates are resampled with probability $\eps^{R-1}$, independently. So $a$ is a cut with
conditional probability at least $\kappa=\eps^Dc$ on $G_a$, and this event is
$\cF_{a+R-1}$-measurable. Let $\Gamma=R\lceil3\ln(n+1)/\kappa\rceil$. An interval of $\Gamma$
consecutive positions in $(g_0,n-R+1]$ contains $\Gamma/R$ positions that are $R$ apart, and by
successive conditioning the probability that $G$ occurs and none of them is a cut is at most
$(1-\kappa)^{\Gamma/R}\le(n+1)^{-3}$. Let $\mathcal C$ be the event that every such interval
contains a cut; then $\PP[\mathcal C^c\cap G]\le(n+1)^{-2}$. Put $\Gamma'=\Gamma+R$ and
$g_1=g_0+\Gamma'$. On $\mathcal C$, for every position $q\ge g_1$ there is a cut $a$ with
$q-\Gamma'<a\le q-R$.

\emph{No step with a base pointer crosses a cut.} Let $a$ be a cut and let step $s$ of a dependency
chain have $P_{s-1}\ge a$. Gate $a$ is a reset and no gate in $(k_s,P_{s-1}]$ is a reset, so
$k_s\ge a$. If moreover $g_0\le P_s<a$, then $\pi_{k_s}=1$: this is clear for $k_s\le a+R-1$, and
if $k_s\ge a+R$ and $\pi_{k_s}=0$ then $P_s=j_{k_s}$ is $0$ or at least $k_s-R\ge a$.

\emph{Segments.} Work on $\mathcal C$ and consider a dependency chain with all positions at least
$g_1$. Its steps with $\pi_{k_s}=1$ split it into segments of steps with $\pi_{k_s}=0$. Let $q$ be
the position at which a segment starts, and let $a$ be a cut with $q-\Gamma'<a\le q-R$. By the
previous paragraph all positions of the segment are at least $a$, and so is the gate of the
resampled step that ends it. Positions decrease by at least $2$ per step, so the segment has at
most $\Gamma'/2$ steps.

\emph{Resampled steps.} A \emph{jump skeleton} of length $r$ from $P_0$ is a tuple
$(k_1,q_1,\dots,k_r,q_r)$ with $P_0-\Gamma'<k_1\le P_0$, $g_1\le q_i\le k_i-2$ and
$q_i-\Gamma'<k_{i+1}\le q_i$. The event that $\pi_{k_i}=1$ and $\rho_{k_i}=q_i$ for all $i$ has
probability $\prod_i\eps/k_i$, because the gates $k_i$ are distinct. On $\mathcal C$, a dependency
chain from $P_0$ with all positions at least $g_1$ and with $r$ resampled steps gives a jump
skeleton whose event occurs: the gates and targets of its resampled steps. With
$h(q)=(q+1)^\lambda$,
\[
\sum_{P-\Gamma'<k\le P}\frac{\eps}{k}\sum_{q\le k-2}h(q)\ \le\ \frac{\Gamma'\eps}{\lambda+1}\,h(P),
\]
as in Lemma~\ref{lem:kernel}. Take $\lambda=\lceil2\eps\Gamma'\rceil$, so that the factor is at
most $1/2$. Then the expected number of jump skeletons of length $r$ from $P_0$ whose event occurs
is at most $2^{-r}(n+1)^\lambda$, and with $r^*=\lceil(\lambda+2)\log_2(n+1)\rceil$ and a union
bound over $P_0$, with probability at least $1-1/(n+1)$ no jump skeleton of length $r^*$ occurs.

\emph{Rounds.} On $G$, on $\mathcal C$, and when no jump skeleton of length $r^*$ occurs, every
dependency chain with all positions at least $g_1$ has fewer than $r^*$ resampled steps and at
most $r^*$ segments, hence fewer than $\ell^*=r^*(\Gamma'/2+1)$ steps. Then the Kleene--prefix
iteration finishes within $\ell^*+g_1-1$ rounds, as in the proof of Theorem~\ref{thm:fixed}. Here
$\Gamma'=O(\ln n)$, $\lambda=O(\ln n)$ and $r^*=O(\ln^2n)$, so $\ell^*+g_1=O(\ln^3n)$. The other two
iterations are covered by Lemma~\ref{kp:pruned}(b) and Proposition~\ref{prop:jacobi}.
\end{proof}

\begin{remark}\label{rem:short}
\begin{enumerate}[label=(\roman*),leftmargin=2.4em]
\item Theorem~\ref{thm:short} covers $j^{\mathrm{bar}}$ and every repeating base of the table
above. So on $j^{\mathrm{bar}}$ the expected number of chains grows at noise $0.72$
(Proposition~\ref{prop:barrier}) while the iteration needs only $O(\log^3n)$ rounds, at that and
at every other positive noise level. The simulations above suggest $O(\log n)$.
\item It also covers the fixed-distance bases $j^H$ of Section~\ref{sec:fixed}, at every positive
noise level, with $O(\log^3n)$ rounds in place of $O(\log n)$. The constant is large: the
distance between cuts is of order $\eps^{-R}/c$, and the proof gives
$\ell^*\approx\eps\,\Gamma'^2\log_2n$, a constant of order $R^2\eps^{1-2R}c^{-2}$ in front of
$\ln^3n$, besides the constants hidden in $g_0$ (Remark~\ref{rem:fixedscope}(iv)).
Theorem~\ref{thm:fixed} keeps its interest for large $H$, where that constant is useless, and its
method extends to bases with long pointers (Theorem~\ref{thm:ordered}), about which
Theorem~\ref{thm:short} says nothing.
\item So, for a fixed reach $R$ and a fixed noise level $\eps>0$, short pointers cannot make
smoothed instances asymptotically hard for these iterations. The bases on which the iterations fail
(Section~\ref{sec:block}) and the bases of the hardness proofs (Sections~\ref{sec:tree}
and~\ref{sec:tree2}) use pointers that are long compared with the distance between resets. The
statement is asymptotic: at $\eps=0$ the base $j_i=i-2$ needs linearly many rounds, and at small
noise the bound exceeds $n$ at every size that can be run. On $j_i=i-2$ at noise $0.1$ a simulation
by the referee gave $\max_id(i)=10{,}784$, $28{,}771$ and $55{,}805$ at $n=10^5,10^6,10^7$, with
increments still growing; at noise $0.5$ it gave $430$, $619$ and $858$.
\end{enumerate}
\end{remark}

\section{Lower bounds for the two iterations: block bases and the golden ratio}\label{sec:block}

This section proves that both iterations of Section~\ref{sec:iterations} fail on bases fixed in
advance: on random block bases they need $\Omega(n/\log n)$ rounds up to explicit noise levels,
and for the Kleene--prefix iteration it finds the exact level on such bases. Let
\[
d^*=\frac{3-\sqrt5}{2}\approx0.382,\qquad 1-d^*=\frac{\sqrt5-1}{2}\approx0.618 ,
\]
so that $d^*=(1-d^*)^2$; $d^*$ is the density of ones of uniform instances (Section~\ref{sec:newton}). On block bases
the Kleene--prefix iteration needs $\Omega(n/\log n)$ rounds for every $\eps<1-d^*$ and
$O(\log n)$ rounds for every $\eps>1-d^*$, and the relevance-pruned iteration needs
$\Omega(n/\log n)$ rounds for every $\eps<d^*$ (Theorem~\ref{thm:block}). When the blocks sit on
top of zero gates, the two lower bounds hold up to $\eps_K=0.65896\ldots$ and up to $4/9$
(Theorem~\ref{thm:blockzero}). So the threshold of the Kleene--prefix iteration lies in
$[\eps_K,3/4]$ and that of the relevance-pruned iteration in $[4/9,3/4]$
(Corollary~\ref{cor:thresholds}). For the relevance-pruned iteration over zero gates the level
$4/9$ is exact: above it the iteration needs $O(\log n)$ rounds on these bases
(Theorem~\ref{thm:blockabove}). These statements are unconditional. They are
about two algorithms and say nothing about $\eps_c$.

The reason is that on a block base the reads of a window are independent samples of the block
below. The dependency chains then form a branching process, and the first-moment count of Theorem~\ref{kp:thm34} is
sharp instead of only an upper bound. This is the opposite of the base of Section~\ref{sec:barrier},
where the chains overlap.

\subsection{Block bases and their densities}

\begin{definition}[block bases]\label{def:block}
Let $m\ge2$, $n_0\ge0$ and $T\ge1$ with $n_0+(T+1)m\le n$. For $0\le t\le T$, \emph{block} $t$ is
the set of positions $\mathcal B_t=(n_0+tm,\,n_0+(t+1)m]$. In a \emph{block base} every gate
outside blocks $1,\dots,T$ is a zero gate, $j_i=0$, and every gate $i$ of a block $t\ge1$ has
$j_i\in\mathcal B_{t-1}$. In a \emph{random block base} these pointers, written $J_i$, are
independent and uniform on $\mathcal B_{t-1}$. The base is drawn before the noise and independently
of it.
\end{definition}

\noindent A base drawn at random is still fixed before the noise: if a statement holds with probability
$1-\delta$ over base and noise, then by Markov's inequality, for all bases outside a set of
probability $\sqrt\delta$, it holds with probability at least $1-\sqrt\delta$ over the noise.

Throughout, $\mathcal H_g$ is the $\sigma$-field of the pointers $J_i$ and the noise
$(\pi_i,\rho_i)$ of the gates $i\le g$; the trace up to $g$ is $\mathcal H_g$-measurable. Let
$\delta_t$ be the fraction of ones among the positions of block $t$, let $p_g=O_{g-1}/g$, and, for a
gate $g$ of a block $t\ge1$,
\[
q_g=1-(1-\eps)\,\delta_{t-1}-\eps\,p_g .
\]
Ones are never adjacent, so $\delta_t\le\frac12+\frac1{2m}$ and $p_g\le\frac12+\frac1{2g}$.

\begin{lemma}[one gate]\label{lem:blockgate}
Let $g$ be a gate of a block $t\ge1$. Then $(J_g,\pi_g,\rho_g)$ is independent of
$\mathcal H_{g-1}$, and
\begin{enumerate}[label=(\roman*),leftmargin=2.4em]
\item $\PP[x_g=1\mid\mathcal H_{g-1}]=(1-x_{g-1})\,q_g$;
\item if $g\ge n_0+tm+2$, then $q_g\le\PP[g\text{ is not a reset}\mid\mathcal H_{g-1}]\le q_g+\eps/g$;
\item if $S\subseteq\mathcal B_{t-1}$ is determined by $\mathcal H_{n_0+tm}$, then
$\PP[\pi_g=0,\ J_g\in S\mid\mathcal H_{g-1}]=(1-\eps)\,|S|/m$.
\end{enumerate}
\end{lemma}

\begin{proof}
A gate that keeps its pointer reads $x_{J_g}$, and $J_g$ is uniform on block $t-1$, whose values
are $\mathcal H_{n_0+tm}$-measurable; so it reads a one with probability $\delta_{t-1}$. A resampled
gate reads $x_{\rho_g}$ with $\rho_g$ uniform on $\{0,\dots,g-1\}$, a one with probability $p_g$.
So the value read is a zero with conditional probability $q_g$, and $x_g=1$ exactly when
$x_{g-1}=0$ and the value read is a zero. This is (i). For (ii), $J_g\le n_0+tm\le g-2$, so a gate
that keeps its pointer is not a NOT gate, and it is a reset exactly when it reads a one. A
resampled gate is a reset exactly when $\rho_g\le g-2$ and $x_{\rho_g}=1$, which has probability
$(O_{g-1}-x_{g-1})/g\in[p_g-1/g,\,p_g]$. (iii) is the independence.
\end{proof}

Let $\Lambda(q)=q/(1+q)$ for $q\ge0$; it is increasing with $0<\Lambda'\le1$. As in
Lemma~\ref{lem:fddens}, let $a_i=\PP[x_i=1\mid\mathcal H_{i-1}]$ and let $G^*$ be the event that
$\bigl|\sum_{i=a+1}^{b}(x_i-a_i)\bigr|\le\lambda(b-a)$ for all $0\le a<b\le n$, where
$\lambda(M)=\sqrt{2M\ln(n+1)}$. By the Azuma--Hoeffding inequality $\PP[(G^*)^c]\le(n+1)^{-2}$. Put
$\xi=(2+\lambda(m))/m$; if $m\ge C\ln(n+1)$ then $\xi\le2/m+\sqrt{2/C}$.

\begin{lemma}[intervals]\label{lem:blockint}
Work on $G^*$. Let $(a,b]$ be an interval of positions inside the blocks $t\ge1$, of length
$M=b-a$, and let $X=O_b-O_a$. If $q_i\le q^+$ for all $i\in(a,b]$, then
$X\le\Lambda(q^+)M+1+\lambda(M)$. If $q_i\ge q^-\ge0$ for all $i\in(a,b]$, then
$X\ge\Lambda(q^-)M-1-\lambda(M)$.
\end{lemma}

\begin{proof}
By Lemma~\ref{lem:blockgate}(i), $a_i=(1-x_{i-1})q_i$, and positions $a,\dots,b-1$ hold between
$M-X-1$ and $M-X+1$ zeros. So $X\le q^+(M-X+1)+\lambda(M)$, that is,
$X\le(q^+M+q^++\lambda(M))/(1+q^+)$; and $X\ge q^-(M-X-1)-\lambda(M)$, that is,
$X\ge(q^-M-q^--\lambda(M))/(1+q^-)$.
\end{proof}

\begin{lemma}[densities; no zero gates below]\label{lem:blockdens}
Let $n_0=0$. For every $\eta\in(0,1/100]$ there are $C_\eta$ and $t_\eta$ such that, if
$m\ge C_\eta\ln(n+1)$, then on $G^*$: $|\delta_t-d^*|\le\eta$ for every block $t\ge t_\eta$, and
$|p_g-d^*|\le\eta$ for every gate $g$ of a block $t\ge t_\eta$.
\end{lemma}

\begin{proof}
Let $h(x)=\Lambda(1-x)=(1-x)/(2-x)$ for $x\le1$. It is decreasing, $h(d^*)=d^*$ because
$d^{*2}-3d^*+1=0$, and $|h'(x)|=(2-x)^{-2}<0.46$ for $x\le0.52$. Fix $\eta_1=\eta/16$ and let
$C_\eta$ be so large that $m\ge100$ and $\xi\le\eta/32$; put $e=\eta_1+2\xi\le\eta/8$.

We show by induction on $k$ that there are numbers $\ell_k\le d^*\le u_k$ and $t_k$ with: for all
blocks $t\ge t_k$ and all gates $g$ of such blocks, $\ell_k\le\delta_t\le u_k$ and
$\ell_k\le p_g\le u_k$. For $k=0$ take $\ell_0=0$, $u_0=\frac12+\frac1m$ and $t_0=1$. Suppose the
claim holds for $k$, with $0\le\ell_k$ and $u_k\le0.52$. For a gate $g$ of a block $t\ge t_k+1$,
both $\delta_{t-1}$ and $p_g$ lie in $[\ell_k,u_k]$, so $q_g\in[1-u_k,\,1-\ell_k]$. By
Lemma~\ref{lem:blockint}, every interval of length $M\le m$ inside one block $t\ge t_k+1$ has
\[
h(u_k)M-1-\lambda(m)\ \le\ X\ \le\ h(\ell_k)M+1+\lambda(m).
\]
With $M=m$ this gives $h(u_k)-\xi\le\delta_t\le h(\ell_k)+\xi$. Let $t_{k+1}=\lceil(t_k+2)/\eta_1\rceil$
and let $g$ be a gate of a block $t\ge t_{k+1}$. The positions in $((t_k+1)m,\,g-1]$ number
$M'\ge(1-\eta_1)g$, and they meet at most $M'/m+1$ blocks, all with index at least $t_k+1$. Summing
the display over these blocks, and using $O_{(t_k+1)m}\le(t_k+1)m+1\le\eta_1g$,
\[
h(u_k)(1-\eta_1)-2\xi\ \le\ p_g\ \le\ \eta_1+h(\ell_k)+2\xi ,
\]
because $(M'/m+1)(1+\lambda(m))\le2\xi g$. As $h\le1$, the claim holds for $k+1$ with
\[
\ell_{k+1}=h(u_k)-e,\qquad u_{k+1}=h(\ell_k)+e .
\]
Here $h(u_k)\le d^*\le h(\ell_k)$, so $\ell_{k+1}\le d^*\le u_{k+1}$; also
$u_{k+1}\le h(0)+e\le0.52$ and $\ell_{k+1}\ge h(0.52)-e\ge0$. The widths $w_k=u_k-\ell_k$ satisfy
$w_{k+1}\le0.46\,w_k+2e$, hence $w_k\le0.46^k+4e\le0.46^k+\eta/2$. Take $k$ with $0.46^k\le\eta/2$
and $t_\eta=t_k$.
\end{proof}

For blocks on top of zero gates the prefix density is that of the zero gates. For $\eps\in(0,1)$
let $q(\eps)\in(0,1)$ be the solution of
\begin{equation}\label{eq:qeps}
q=1-(1-\eps)\,\frac{q}{1+q}-\eps\,\varpi(\eps),
\end{equation}
which is unique because the right-hand side decreases in $q$, exceeds $0$ at $q=0$ and is below
$1$ at $q=1$. Put $d_\eps=\Lambda(q(\eps))$, so that $q(\eps)=1-(1-\eps)d_\eps-\eps\varpi(\eps)$.

\begin{lemma}[densities; blocks over zero gates]\label{lem:blockdens0}
Let $\eps\in(0,1)$ and $\eta>0$. There are $C_\eta$, $t_\eta$ and $c_\eta>0$ such that for every
large $n$ and every random block base with $m\ge C_\eta\ln(n+1)$ and $(T+1)m\le c_\eta n$, where
$n_0=n-(T+1)m$, the following holds with probability at least $1-2(n+1)^{-2}$ over the base and the
noise:
$|p_g-\varpi(\eps)|\le\eta$ for every gate $g>n_0$, and $|\delta_t-d_\eps|\le\eta$ for every block
$t\ge t_\eta$.
\end{lemma}

\begin{proof}
\emph{Prefix densities.} Apply Lemma~\ref{lem:dens2any} with the gates of blocks $1,\dots,T$ as
cell gates, $\mathrm{pre}=n_0$ and $c=(T+1)m$, for each choice of the pointers: every interval
contains at most $c$ cell gates.
Choose $\rho\ge1$, $\tau\le1/8$ and $k\ge3$, depending on $\eps$ and $\eta$, with $1/(1+\rho)+\tau\le1/4$ and
$(1/(1+\rho)+\tau)/(1-\eps)+\eps^{2k-1}/2\le\eta'$, where $\eta'=\min(\eta,\eps\eta)/4$, and
$K=2$. Then $W=c+3+4\ln(n+1)$, and $\mathrm{pre}\ge(5/\tau^2)^{2k+1}W$ holds for large $n$ if
$c_\eta\le\frac12(5/\tau^2)^{-(2k+1)}$. So with probability at least $1-(n+1)^{-2}$,
$|\psi_g-\varpi|\le\eta'$ for all $g\in[n_0,n]$; and $p_g=\psi_{g-1}$.

\emph{Blocks.} Work also on $G^*$. Let $t\ge1$ and suppose $|\delta_{t-1}-d_\eps|\le e$. For a gate
$g$ of block $t$, $|q_g-q(\eps)|\le(1-\eps)e+\eps\eta'$, and $q_g\ge0$. By Lemma~\ref{lem:blockint}
with $M=m$, and since $\Lambda$ is $1$-Lipschitz with $\Lambda(q(\eps))=d_\eps$,
\[
|\delta_t-d_\eps|\le(1-\eps)\,e+\eps\eta'+\xi .
\]
Starting from $|\delta_0-d_\eps|\le1$, this gives
$|\delta_t-d_\eps|\le(1-\eps)^t+\eta'+\xi/\eps$. Choose $C_\eta$ with $\xi\le\eps\eta/4$ and $t_\eta$
with $(1-\eps)^{t_\eta}\le\eta/4$.
\end{proof}

\subsection{One block}

The next lemma uses only the formula for $D$ recalled in Section~\ref{sec:upper}.

\begin{lemma}[two facts about $D$]\label{lem:Dfacts}
Let $j$ be a parent array. (i) If gate $k$ is relevant and not a NOT gate, then
$D(k)\ge1+D(j_k)$. (ii) If gate $g\ge1$ is not a reset, then $D(g)\ge D(g-1)$.
\end{lemma}

\begin{proof}
(i) If $k$ is a reset, it is a reset at a relevant gate, so $\varrho'(k)=k$, the only anchor of $k$
is $k$, and $D(k)=1+D(j_k)$. If $k$ is not a reset, every anchor $r$ of $k$ satisfies $r<k$, and
$k$ is a relevant gate of $(r,k]$ that is not a NOT gate; so the maximum in the formula is at least
$D(j_k)$ for every anchor. (ii) Since $g$ is not a reset, $\varrho'(g)=\varrho'(g-1)$ and $g$ and
$g-1$ have the same anchors. For each anchor the set over which the maximum is taken for $g$
contains the one for $g-1$. For $g=1$ the claim is $D(1)\ge0$.
\end{proof}

\begin{lemma}[one block]\label{lem:oneblock}
Let $t\ge1$, let $S\subseteq\mathcal B_{t-1}$ be determined by $\mathcal H_{n_0+tm}$, and put
$\sigma=(1-\eps)|S|/m$. Let $L,J\ge1$ and let $i\in\mathcal B_t$ with $g_0:=i-L-J\ge n_0+tm+1$. For
$0\le l<L$ and $k=i-l$ let
\[
E^{d}_l=\{\pi_k=0,\ J_k\in S,\ \text{no gate in }(k,i]\text{ is a reset}\},\qquad
E^{D}_l=E^{d}_l\cap\{x_{k-1}=0\}.
\]
Let $0<q^-\le q^+<1$ and let $\mathcal D\in\mathcal H_{g_0}$ be an event on which
$q^-\le q_g\le q^+$ for every gate $g\in(g_0,i]$. Then, on $\mathcal D$,
\[
\PP\Bigl[\bigcup_{l<L}E^{d}_l\Bigm|\mathcal H_{g_0}\Bigr]\ \ge\ \sigma\,\frac{1-(q^-)^L}{1-q^-}-\frac{L^2\sigma^2}{2},
\qquad
\PP\Bigl[\bigcup_{l<L}E^{D}_l\Bigm|\mathcal H_{g_0}\Bigr]\ \ge\ \zeta\,\sigma\,\frac{1-(q^-)^L}{1-q^-}-\frac{L^2\sigma^2}{2},
\]
where $\zeta=(1-q^+)/(1-q^+q^-)-(q^+q^-)^{\lfloor J/2\rfloor}$. The first inequality needs only
$q_g\ge q^-$ on $\mathcal D$.
\end{lemma}

\begin{proof}
All gates $g\in(g_0,i]$ satisfy $g\ge n_0+tm+2$. Work on $\mathcal D$.

\emph{The value below the step's gate.} For $g\in(g_0,i]$ and $0\le j\le g-g_0$ we claim
$v_j\le\PP[x_g=0\mid\mathcal H_{g-j}]\le w_j$, where $v_0=0$, $w_0=1$, $v_j=1-q^+w_{j-1}$ and
$w_j=1-q^-v_{j-1}$. For $j\ge1$, Lemma~\ref{lem:blockgate}(i) gives
$\PP[x_g=0\mid\mathcal H_{g-j}]=1-\EE[(1-x_{g-1})q_g\mid\mathcal H_{g-j}]$, where
$q_g\in[q^-,q^+]$ and $\PP[x_{g-1}=0\mid\mathcal H_{g-j}]\in[v_{j-1},w_{j-1}]$ by induction (for
$g-1=g_0$ and $j=1$ this is the case $j-1=0$). Now $v_{j+2}=(1-q^+)+q^+q^-v_j$, with fixed point
$v_\infty=(1-q^+)/(1-q^+q^-)$, and $0\le v_\infty-v_j\le1$ for $j\in\{0,1\}$. Hence
$v_j\ge v_\infty-(q^+q^-)^{\lfloor j/2\rfloor}$. Since $k-1-g_0=L+J-l-1\ge J$, we get
$\PP[x_{k-1}=0\mid\mathcal H_{g_0}]\ge\zeta$.

\emph{One event.} Condition successively on $\mathcal H_{i-1},\dots,\mathcal H_k$: by
Lemma~\ref{lem:blockgate}(ii), each gate of $(k,i]$ is not a reset with conditional probability at
least $q^-$. Then condition on $\mathcal H_{k-1}$: by Lemma~\ref{lem:blockgate}(iii),
$\pi_k=0$ and $J_k\in S$ have conditional probability $\sigma$. So
$\PP[E^{d}_l\mid\mathcal H_{g_0}]\ge\sigma(q^-)^l$ and, with the previous paragraph,
$\PP[E^{D}_l\mid\mathcal H_{g_0}]\ge\zeta\sigma(q^-)^l$.

\emph{The union.} For $l\ne l'$, the event $E^{d}_l\cap E^{d}_{l'}$ requires $\pi_k=0$, $J_k\in S$
at two different gates, which has conditional probability $\sigma^2$. By the Bonferroni
inequality the probability of the union is at least the sum of the probabilities minus the sum
over the at most $L^2/2$ pairs.
\end{proof}

\subsection{The theorems}

By Lemma~\ref{fact:depths}, the Kleene--prefix iteration needs exactly $\max_id(i)$ rounds and the
relevance-pruned iteration exactly $\max_iD(i)$ rounds.

\begin{theorem}[block bases and the golden ratio]\label{thm:block}
Let $\eps\in(0,1)$ with $\eps\ne1-d^*$. There are $C,C',c>0$ and $n_1$, depending only on $\eps$, such
that for every $n\ge n_1$ the following holds for the $\eps$-smoothing of a random block base with
$n_0=0$, $m=\lceil C\ln(n+1)\rceil$ and $T=\lfloor n/m\rfloor-1$, with probability at least
$1-3/n$ over the base and the noise.
\begin{enumerate}[label=(\alph*),leftmargin=2.4em]
\item If $\eps<1-d^*$, then $\max_id(i)\ge cn/\ln n$: the Kleene--prefix iteration needs at least
$cn/\ln n$ rounds.
\item If $\eps<d^*$, then $\max_iD(i)\ge cn/\ln n$: the relevance-pruned iteration needs at least
$cn/\ln n$ rounds.
\item If $\eps>1-d^*$, then the Kleene--prefix iteration finishes within $C'\ln n$ rounds.
\end{enumerate}
For $\eps<1/2$, part (a) has a proof that uses no density bound.
\end{theorem}

\begin{proof}
\emph{Parameters for (a) and (b).} Let $\eta\in(0,1/100]$, $q^\pm=1-d^*\pm2\eta$, and let $\zeta$ be
as in Lemma~\ref{lem:oneblock} in case (b) and $\zeta=1$ in case (a). Put
$\varrho_0=(1-\eps)\zeta/(1-q^-)$. In case (a), $\varrho_0=(1-\eps)/(d^*+2\eta)$, which exceeds $1$
for small $\eta$ because $\eps<1-d^*$. In case (b), as $\eta\to0$ and $J\to\infty$ we have
$\zeta\to1/(2-d^*)=1-d^*$ and $\varrho_0\to(1-\eps)(1-d^*)/d^*=(1-\eps)/(1-d^*)$, which exceeds
$1$ because $\eps<d^*$; here we used $d^*=(1-d^*)^2$ twice. Fix $\eta$, $J$ and $\kappa>0$ with
$\varrho_0\ge1+4\kappa$, then $L$ with $\varrho_0(1-(q^-)^L)\ge1+3\kappa$, and put
$\pi_0=\kappa/(4L^2)$, $s_0=\lceil\pi_0m\rceil$ and $\eta_2=\kappa\pi_0/2$. Let $C$ be so large that
$m\ge C_\eta\ln(n+1)$, $m\ge\max(1/\pi_0,(L+J)/\eta)$, and the two estimates marked $(\ast)$ below
hold.

\emph{The sets.} Let $t_\eta$ be as in Lemma~\ref{lem:blockdens} and $S_{t_\eta}=\mathcal B_{t_\eta}$.
For $t>t_\eta$: if $|S_{t-1}|\ge s_0$, let $S'_{t-1}$ consist of the $s_0$ smallest elements of
$S_{t-1}$, and let $S_t$ be the set of $i\in\mathcal B_t$ with $i-L-J\ge tm+1$ for which
$\bigcup_{l<L}E_l$ occurs with $S=S'_{t-1}$, where $E_l=E^{d}_l$ in case (a) and $E_l=E^{D}_l$ in
case (b); otherwise $S_t=\emptyset$. We claim that $d(i)\ge t-t_\eta$, respectively
$D(i)\ge t-t_\eta$, for $i\in S_t$. Indeed let $E_l$ occur and $k=i-l$. Then $\jt_k=J_k\in S_{t-1}$,
gate $k$ is not a NOT gate because $J_k\le tm\le k-2$, and no gate in $(k,i]$ is a reset. In case
(a), $d(i)\ge1+d(J_k)$ by the recursion for $d$ in Lemma~\ref{fact:depths}. In case (b), gate $k$ is
also relevant, so $D(i)\ge D(k)\ge1+D(J_k)$ by Lemma~\ref{lem:Dfacts}. Induction on $t$ gives the
claim.

\emph{One block.} Let $t>t_\eta$ and $i\in\mathcal B_t$ with $g_0=i-L-J\ge tm+1$, and let
$\mathcal D_i\in\mathcal H_{g_0}$ be the event that $|\delta_{t-1}-d^*|\le\eta$ and
$|p_{g_0+1}-d^*|\le\eta$. For $g\in(g_0,i]$ we have $|p_g-p_{g_0+1}|\le(L+J)/g\le\eta$, so on
$\mathcal D_i$ every such gate has $q_g\in[q^-,q^+]$. With $\sigma=(1-\eps)s_0/m$ we have
$(1-\eps)\pi_0\le\sigma\le2\pi_0$, so by Lemma~\ref{lem:oneblock}, on $\mathcal D_i$ and on
$\{|S_{t-1}|\ge s_0\}$,
\[
\PP[i\in S_t\mid\mathcal H_{g_0}]\ \ge\ \pi_0\,\varrho_0\bigl(1-(q^-)^L\bigr)-2L^2\pi_0^2
\ \ge\ (1+3\kappa)\pi_0-\tfrac12\kappa\pi_0\ \ge\ \gamma:=(1+2\kappa)\pi_0 .
\]
Let $Y_i=\mathbf 1[i\in S_t]$ and $Y'_i=\max(Y_i,\,1-\mathbf 1_{\mathcal D_i})$, so that
$\EE[Y'_i\mid\mathcal H_{g_0}]\ge\gamma$ on $\{|S_{t-1}|\ge s_0\}$. The indices $i$ in question
number $m-L-J$. For each residue class modulo $L+J$, the variables $Y'_i-\EE[Y'_i\mid\mathcal H_{i-L-J}]$
with $i$ in the class are martingale differences with values in $[-1,1]$. By the Azuma--Hoeffding
inequality, on $\{|S_{t-1}|\ge s_0\}$, except with conditional probability at most
\[
(L+J)\exp\Bigl(-\tfrac12\eta_2^2\bigl(\tfrac{m-L-J}{L+J}-1\bigr)\Bigr)\ \le\ (n+1)^{-3}\qquad(\ast)
\]
given $\mathcal H_{tm}$, we have
$\sum_iY'_i\ge(\gamma-\eta_2)(m-L-J)\ge s_0$ $(\ast)$.

\emph{Conclusion.} By Lemma~\ref{lem:blockdens}, on $G^*$ every event $\mathcal D_i$ with
$i\in\mathcal B_t$, $t>t_\eta$, occurs, so there $Y'_i=Y_i$. Hence, except with probability
$(n+1)^{-2}+n(n+1)^{-3}\le2/n$, we get $|S_t|\ge s_0$ for all $t_\eta\le t\le T$ by induction on
$t$. Then $S_T\ne\emptyset$, and the maximal depth is at least $T-t_\eta\ge cn/\ln n$.

\emph{No density bound.} Since $\delta_{t-1}\le\frac12+\frac1{2m}$ and $p_g\le\frac12+\frac1{2g}$,
every gate of a block $t\ge1$ has $q_g\ge\frac12-\frac1m$. For $\eps<1/2$ the argument for (a)
goes through with $q^-=\frac12-\frac1m$, with $\mathcal D_i$ the sure event and with $t_\eta=0$,
because $(1-\eps)/(1-q^-)>1$ for large $m$.

\emph{(c).} Let $\eta$ be so small that $\theta_0:=(1-\eps)/(d^*-3\eta)<1$, and let $G$ be the
event $G^*$, on which the bounds of Lemma~\ref{lem:blockdens} hold. Put $g_0=(t_\eta+1)m$ and
$y=1-d^*+3\eta$. For a gate $g>g_0$ let $G_g\in\mathcal H_{g-1}$ be the event that the density
bounds of Lemma~\ref{lem:blockdens} hold for $p_g$ and, if $g$ lies in a block $t$, for
$\delta_{t-1}$. On $G_g$, gate $g$ is not a reset with conditional probability at most $y$, for
large $m$: for $g\ge tm+2$ in a block this is Lemma~\ref{lem:blockgate}(ii); the first gate of a
block is a reset when it keeps a pointer other than $g-1$ to a one, or is resampled to a one below
$g-1$, which has probability at least $(1-\eps)(\delta_{t-1}-1/m)+\eps(p_g-1/g)$; and a zero gate
above the blocks, of which there are fewer than $m$, is a reset when it keeps its pointer or is
resampled to a one below $g-1$, which has probability at least $(1-\eps)+\eps(p_g-1/g)$, where
$p_g$ differs from its value at the last gate of block $T$ by at most $m/g$ (for these gates let
$G_g$ be the event that the bound holds at that last gate). Also, for $q\le k-2$, the event
$\{\jt_k=q\}$ is independent of $\mathcal H_{k-1}$ and has probability
$w(k,q)=\eps/k+(1-\eps)\PP[J_k=q]$, where $\PP[J_k=q]$ is $1/m$ for $q$ in the block below that of
$k$ and $0$ otherwise (for a zero gate, $\PP[J_k=q]=\mathbf 1[q=0]$). With these two facts the
proof of Lemma~\ref{lem:onechain} gives, for every skeleton with $P_\ell\ge g_0$,
$\PP[E_\sigma\cap G]\le\prod_s y^{P_{s-1}-k_s}w(k_s,P_s)$. Since
$\sum_{q=g_0}^{k-2}w(k,q)(q+1)^\lambda\le(\eps/(\lambda+1)+1-\eps)k^\lambda$, the proof of
Lemma~\ref{lem:kernel} gives the bound $\theta^\ell(n+1)^\lambda$ for the expected number of
dependency chains of length $\ell$ from a position, with
$\theta=(\eps/(\lambda+1)+1-\eps)/(d^*-3\eta)$, which is below $1$ for large $\lambda$. As in the
proof of Theorem~\ref{thm:fixed}, with probability at least $1-2/n$ there is no dependency chain of
length $\ell^*=\lceil((\lambda+1)\ln(n+1)+\ln n)/\ln(1/\theta)\rceil$ with all positions at least
$g_0$, and then the Kleene--prefix iteration finishes within $\ell^*+g_0-1\le C'\ln n$ rounds.
\end{proof}

\begin{theorem}[blocks over zero gates]\label{thm:blockzero}
Let $\eps\in(0,1)$ and let $q(\eps)$ be given by~\eqref{eq:qeps}. There are $C,c,c'>0$ and $n_1$,
depending only on $\eps$, such that for every $n\ge n_1$ the following holds for the
$\eps$-smoothing of a random block base with $m=\lceil C\ln(n+1)\rceil$,
$T=\lfloor c'n/m\rfloor-1$ and $n_0=n-(T+1)m$, with probability at least $1-3/n$ over the base and
the noise.
\begin{enumerate}[label=(\alph*),leftmargin=2.4em]
\item If $q(\eps)>\eps$, then $\max_id(i)\ge cn/\ln n$.
\item If $q(\eps)^2>\eps$, then $\max_iD(i)\ge cn/\ln n$.
\end{enumerate}
The condition of (a) holds exactly for $\eps<\eps_K$, where $\eps_K=0.65896\ldots$ is the root of
$\eps^3+4\eps^2+3\eps-4$ in $(0,1)$. The condition of (b) holds exactly for $\eps<4/9$.
\end{theorem}

\begin{proof}
Repeat the proof of Theorem~\ref{thm:block}(a),(b) with Lemma~\ref{lem:blockdens0} in place of
Lemma~\ref{lem:blockdens}, with $q^\pm=q(\eps)\pm2\eta$, and with $\mathcal D_i$ the event that
$|\delta_{t-1}-d_\eps|\le\eta$ and $|p_{g_0+1}-\varpi(\eps)|\le\eta$; on it,
$|q_g-q(\eps)|\le(1-\eps)\eta+2\eps\eta\le2\eta$ for the gates of the window. As $\eta\to0$ and
$J\to\infty$, the number $\varrho_0=(1-\eps)\zeta/(1-q^-)$ tends to $(1-\eps)/(1-q(\eps))$ in case
(a) and, since $\zeta\to1/(1+q(\eps))$, to $(1-\eps)/(1-q(\eps)^2)$ in case (b). These limits
exceed $1$ exactly under the two conditions. The blocks occupy $(T+1)m\le c'n$ positions, with
$c'=c_\eta$, and $T\ge cn/\ln n$.

\emph{The two conditions.} Write $\varpi=\varpi(\eps)$ and $q=q(\eps)$. If $q=\eps$,
then~\eqref{eq:qeps} gives $\eps\varpi=(1-\eps)/(1+\eps)$, and $\varpi=\eps(1-\varpi)^2$ becomes
$1-\eps^2=(\eps^2+2\eps-1)^2$, that is, $\eps(\eps^3+4\eps^2+3\eps-4)=0$. The cubic is increasing
with one root $\eps_K$ in $(0,1)$. Since $q(\eps)-\eps$ is continuous, positive near $0$ and
negative near $1$, it is positive exactly on $(0,\eps_K)$. Now let $q^2=\eps$ and put
$s=\sqrt\eps$, so $q=s$. Since $(1-\eps)/(1+s)=1-s$, equation~\eqref{eq:qeps} becomes
$\eps\varpi=(1-s)^2$, and then $\varpi=\eps(1-\varpi)^2$ becomes $(1-s)^2=(2s-1)^2$, whose only
solution in $(0,1)$ is $s=2/3$. At $\eps=4/9$ indeed $\varpi=1/4$ and $q=2/3$. So
$q(\eps)^2-\eps$ vanishes only at $4/9$; it is continuous, positive near $0$ and negative near
$1$, hence positive exactly on $(0,4/9)$. The script \texttt{block\_thresholds.py} confirms both
sign patterns on a grid of step $10^{-4}$.
\end{proof}

\begin{corollary}[thresholds of the two iterations]\label{cor:thresholds}
Call the threshold of an iteration the infimum of the noise levels at which it finishes in
polylogarithmically many rounds, with probability tending to $1$, on the smoothing of every base. The threshold of the Kleene--prefix iteration lies in
$[\eps_K,3/4]\subset[0.6589,\,0.75]$, and that of the relevance-pruned iteration in $[4/9,\,3/4]$.
\end{corollary}

\begin{proof}
The upper bound $3/4$ is Theorem~\ref{kp:thm39}. For $\eps<\eps_K$, respectively $\eps<4/9$,
Theorem~\ref{thm:blockzero} and Markov's inequality give, for every large $n$, a base on whose
smoothing the iteration needs at least $cn/\ln n$ rounds with probability at least $1-2/\sqrt n$.
\end{proof}

\noindent For the relevance-pruned iteration Theorem~\ref{thm:allbases} lowers the upper end to
$\eps_1<0.7288$.

\subsection{Above \texorpdfstring{$4/9$}{4/9} the relevance-pruned iteration is fast on block bases}\label{sec:above}

Theorem~\ref{thm:blockzero}(b) is a lower bound below $4/9$. This subsection proves the matching
upper bound on the same kind of bases: for every $\eps>4/9$ the relevance-pruned iteration needs
$O(\log n)$ rounds on a random block base over zero gates (Theorem~\ref{thm:blockabove}). So on
these bases $4/9$ is the exact level of that iteration; it is also the level of the hardness
theorem of Section~\ref{sec:latch}, which Section~\ref{sec:s24} extends to $0.4468$.

A count of chains with a fixed rule for the anchor does not give this just above $4/9$
(Remark~\ref{rem:above}). What makes the iteration fast is a second-order effect: a gate whose
predecessor holds $1$ and which reads a $1$ is determined as soon as \emph{one} of the two is
determined, so an undetermined value passes such a gate only if both are undetermined. The proof
has two parts. It first follows the fraction of positions that are still undetermined after $s$
rounds: in a block that fraction is bounded, up to small errors, by the image of the fractions
in the block below under an explicit map $\Psi$ of the plane, and $\Psi$ has a fixed point other
than $0$ exactly for $\eps<4/9$ (Proposition~\ref{prop:psi}). Once the fraction is small, a
count of chains that uses relevance and the resets with determined targets finishes.

For a parent array and $s\ge0$ call a position $i$ \emph{$s$-open} if $D(i)>s$ and
\emph{$s$-determined} otherwise. By Lemma~\ref{fact:depths} the relevance-pruned iteration has
finished after $s$ rounds exactly when no position is $s$-open.

\begin{lemma}[the recursion for $D$]\label{lem:status}
Let $j$ be a parent array with trace $x$, and let $k\ge2$ be a gate with pointer $p=j_k$.
\begin{enumerate}[label=(\roman*),leftmargin=2.4em]
\item If $k$ is a reset and relevant, then $D(k)=1+D(p)$.
\item If $k$ is a reset and not relevant, then $D(k)=\min\bigl(D(k-1),\,1+D(p)\bigr)$.
\item If $k$ is relevant, not a reset and not a NOT gate, then $D(k)=\max\bigl(D(k-1),\,1+D(p)\bigr)$.
\item In all other cases $D(k)=D(k-1)$.
\end{enumerate}
Consequently, for every $s\ge0$:
\begin{itemize}[leftmargin=3.2em]
\item[(O1)] if $k$ is $(s+1)$-open and $k-1$ is $(s+1)$-determined, then $x_{k-1}=0$, gate $k$ is
not a NOT gate, and $p$ is $s$-open;
\item[(O2)] if $k$ is $(s+1)$-open and a reset, then $p$ is $s$-open.
\end{itemize}
\end{lemma}

\begin{proof}
(i) is shown in the proof of Lemma~\ref{lem:Dfacts}(i), and (iii) and the case of (iv) in which
$k$ is not relevant are in Lemma~\ref{fact:depths}. The cases left are (ii) and a relevant NOT
gate. In both, $k$ is not a reset at a relevant gate, so $\varrho'(k)=\varrho'(k-1)$, and the
anchors of $k$ are those of $k-1$, together with $k$ itself in case (ii). Gate $k$ is not at the
same time relevant and not a NOT gate, so for an anchor $r\le k-1$ the maximum in the formula for
$D$ is taken over the same set for $k$ and for $k-1$. For the anchor $k$ of case (ii) the set is
$\{D(p)\}$. Since $k-1\ge1$, the formula applies to $k-1$, and (ii) and (iv) follow.

(O1) Here $D(k-1)\le s+1<D(k)$, which excludes (ii) and (iv). In (i) and (iii) gate $k$ is
relevant and not a NOT gate, and $1+D(p)>s+1$. (O2) A reset is in case (i) or (ii), where
$D(k)\le1+D(p)$.
\end{proof}

\noindent The script \texttt{status\_check.py} compares (i)--(iv), (O1) and (O2) with the
definition of $D$ on all parent arrays with $n\le7$ (Section~\ref{sec:checks}).

\medskip\noindent\emph{The status chain.} Fix $\eps\in(0,1)$ and write
\[
A=\eps\varpi(\eps),\qquad q=q(\eps),\qquad\sigma=1-q,\qquad
\lambda=\frac{1-\eps}{1-q^2};
\]
recall $d_\eps=q/(1+q)$, so that $\sigma=A+(1-\eps)d_\eps$ by~\eqref{eq:qeps}: a gate of a block
reads a one with probability $\sigma$ in the limit. (In this subsection $\sigma$ is this number
and not a skeleton.) For numbers $0\le a<\sigma$ and $0\le b\le q$ the \emph{status chain} with
parameters $(a,b)$ is the Markov chain on the four states $K_1,K_0,U_1,U_0$ that moves as follows.
At each step a symbol is drawn, independently of the past: a \emph{determined one} with
probability $\sigma-a$, an \emph{open one} with probability $a$, a \emph{determined zero} with
probability $q-b$, an \emph{open zero} with probability $b$. From $K_1$ the chain moves to $K_0$.
From $K_0$ it moves to $K_0$ on a determined one, to $K_1$ on a determined zero, to $U_0$ on an
open one and to $U_1$ on an open zero. From $U_1$ it moves to $K_0$ on a determined one and to
$U_0$ otherwise. From $U_0$ it moves to $K_0$ on a determined one, to $U_0$ on an open one and to
$U_1$ on a zero. The index is the value the position holds, and $U$ stands for open: the value
becomes $1$ exactly when it was $0$ and a zero is read; a determined state stays determined unless
the value is $0$ and the symbol is open; an open state becomes determined exactly on a determined
one. Put
\[
\Delta(a,b)=\sigma(1+q)+q(b-a)+ab,\qquad
\psi_1(a,b)=\frac{b+qa-ab}{(1+q)\,\Delta(a,b)},\qquad
\psi_0(a,b)=\frac{a+qb+ab}{(1+q)\,\Delta(a,b)},
\]
and, for $u=(u_1,u_0)$ with $0\le u_1<\sigma/(1-\eps)$ and $0\le u_0\le q/(1-\eps)$,
\[
\Psi(u)=\bigl(\psi_1((1-\eps)u_1,(1-\eps)u_0),\ \psi_0((1-\eps)u_1,(1-\eps)u_0)\bigr).
\]

\begin{proposition}[the map of the open fractions]\label{prop:psi}
Let $\eps\in(0,1)$.
\begin{enumerate}[label=(\alph*),leftmargin=2.4em]
\item For $0\le a<\sigma$ and $0\le b\le q$ the status chain has exactly one stationary law. It
gives the states $U_1$ and $U_0$ the probabilities $\psi_1(a,b)$ and $\psi_0(a,b)$, and the states
$K_1,U_1$ together the probability $d_\eps$. From every initial state, the law of the chain after
$L$ steps is within $(1-\sigma+a)^L$ of the stationary law in total variation.
\item $\psi_1$ and $\psi_0$ are continuous and nondecreasing in $a$ and in $b$.
\item The equation $\Psi(u)=u$ has a solution $u\ne0$ with nonnegative coordinates if and only if
$\lambda>1$, that is, if and only if $\eps<4/9$. It is then unique: with
$\omega=\sigma(\lambda-1)$, it is given by $(1-\eps)u_1=a_*$ and $(1-\eps)u_0=b_*$, where
$a_*=\frac12\bigl(\sqrt{\omega^2+4\omega}-\omega\bigr)$ and $b_*=a_*+\omega=a_*+a_*b_*$; and
$a_*<\sigma$, $b_*\le q$, so it lies in the domain of $\Psi$.
\item Let $\eps>4/9$. If $u_1+u_0\le r$ and $qr<\sigma(1+q)$, then the two coordinates of
$\Psi(u)$ add up to at most $\lambda\,(u_1+u_0)\big/\bigl(1-\frac{qr}{\sigma(1+q)}\bigr)$. And
for every $r>0$ there is $\bar\gamma>0$ such that for every $\gamma_1\in(0,\bar\gamma]$ there is
$k_0$ such that the sequence $u^{(0)}=(d_\eps+\gamma_1,\,1-d_\eps+\gamma_1)$,
$u^{(k+1)}=\Psi(u^{(k)})+(\gamma_1,\gamma_1)$ is defined and nonincreasing in both coordinates, and
$u^{(k)}_1+u^{(k)}_0\le r$ for all $k\ge k_0$.
\end{enumerate}
\end{proposition}

\begin{proof}
(a) From every state the chain moves to $K_0$ with probability at least $\sigma-a>0$: from $K_1$
with probability $1$, from the others on a determined one. So there is exactly one stationary law,
and two copies of the chain can be coupled to meet with probability at least $\sigma-a$ at each
step, which gives the bound on the total variation. Let $p_{K_1},p_{K_0},p_{U_1},p_{U_0}$ be the
stationary probabilities, $p_U=p_{U_1}+p_{U_0}$, $\upsilon=a+b$, $\alpha_1=\sigma-a$ and $\alpha_0=q-b$. The
open states are left only on a determined one and entered only from $K_0$ on an open symbol, and
$K_1$ is entered only from $K_0$ on a determined zero:
\[
p_U\,\alpha_1=p_{K_0}\,\upsilon,\qquad p_{K_1}=p_{K_0}\,\alpha_0 .
\]
With $p_{K_1}+p_{K_0}+p_U=1$ this gives $p_U=\upsilon/(\upsilon+\alpha_1(1+\alpha_0))$, and
$\upsilon+\alpha_1(1+\alpha_0)=\Delta(a,b)$ because $1-\sigma=q$; in particular $\Delta>0$. The state $U_1$
is entered from $K_0$ on an open zero and from $U_0$ on a zero, so
$p_{U_1}=p_{K_0}\,b+(p_U-p_{U_1})\,q$, that is
$p_{U_1}(1+q)=p_U\,(\alpha_1b/\upsilon+q)$, which is $\psi_1$ after inserting $p_U$ and $\sigma+q=1$;
and $p_{U_0}=p_U-p_{U_1}=\psi_0$. (For $\upsilon=0$ the open states have probability $0$, as the
formulas say.) The value held follows the rule $x'=1$ if and only if $x=0$ and a zero is read,
whatever the status, so its stationary law has $\PP[x=1]=q\,(1-\PP[x=1])$, which gives $d_\eps$.

(b) Continuity is clear from the formulas, as $\Delta>0$. Let $a\le a'$ and $b\le b'$, and run
the chains with parameters $(a,b)$ and $(a',b')$ on the same sequence of ones and zeros, with the
statuses of the symbols coupled so that a symbol that is open for $(a,b)$ is open for $(a',b')$;
this is possible because a one is open with probability $a/\sigma\le a'/\sigma$ and a zero with
probability $b/q\le b'/q$. The values held agree. Write $\hat U=1$ for an open state. The rule for
the status is
\begin{equation}\label{eq:statusrule}
\hat U'=\bigl(\hat U\wedge\neg\,\mathrm{dOne}\bigr)\vee\bigl(\neg\hat U\wedge[x=0]\wedge\mathrm{open}\bigr),
\end{equation}
where $\mathrm{dOne}$ says that the symbol is a determined one and $\mathrm{open}$ that it is
open. The right-hand side is nondecreasing in $\hat U$, because an open symbol is not a
determined one; it is nondecreasing in $\mathrm{open}$ and nonincreasing in $\mathrm{dOne}$. So,
started from the same state, the chain for $(a',b')$ is open whenever the chain for $(a,b)$ is,
and the stationary probabilities of $U_1$ and of $U_0$ are ordered.

(c) Write $a=(1-\eps)u_1$, $b=(1-\eps)u_0$ and $e=1-\eps=\lambda\sigma(1+q)$. Then
$\Psi(u)=u$ reads $a(1+q)\Delta=e\,(b+qa-ab)$ and $b(1+q)\Delta=e\,(a+qb+ab)$. Adding the two,
$(a+b)\Delta=e\,(a+b)$; subtracting, $(b-a)(1+q)\Delta=e\,\bigl((a-b)\sigma+2ab\bigr)$. For
$(a,b)\ne0$ the first says $\Delta=e$, and then the second says $(b-a)(1+q+\sigma)=2ab$, that is
$b-a=ab$. Now $\Delta=\sigma(1+q)+(1+q)ab=e$ gives $ab=e/(1+q)-\sigma=\omega$. If $\lambda<1$
this is impossible; if $\lambda=1$ it forces $ab=0=b-a$, so $(a,b)=0$; if $\lambda>1$ it leaves
$b=a+\omega$ and $a^2+\omega a-\omega=0$, which has one positive root. Conversely that pair solves
both equations. It lies in the domain: $a_*<\sigma$ means $\sigma^2+\omega\sigma-\omega>0$, that
is $\lambda<1/q$, or $(1-\eps)q<1-q^2$, and~\eqref{eq:qeps} gives $(1-\eps)q=1-q^2-A(1+q)$; and
$b_*=a_*/(1-a_*)\le q$ means $a_*\le d_\eps$, that is $d_\eps^2+\omega d_\eps-\omega\ge0$, or
$q^2/(1+q)\ge\omega$, which holds because $\omega=(q^2-\eps)/(1+q)$. Finally $\lambda>1$ means
$q^2>\eps$, which holds exactly for $\eps<4/9$ (Theorem~\ref{thm:blockzero}).

(d) The coordinates of $\Psi(u)$ add up to $(a+b)/\Delta$, and
$\Delta\ge\sigma(1+q)-qa\ge\sigma(1+q)-q(1-\eps)r$, so the sum is at most
$\lambda(u_1+u_0)/(1-qr/(\sigma(1+q)))$, since $1-\eps\le1$. For the sequence, note first that
$\psi_1<d_\eps$ and $\psi_0<1-d_\eps$ at $a=(1-\eps)d_\eps$, $b=(1-\eps)(1-d_\eps)$: by (a) the difference
$d_\eps-\psi_1$ is the stationary probability of $K_1$, and $1-d_\eps-\psi_0$ that of $K_0$, and both are
positive there because $\sigma-a=A>0$ and $q-b=\eps(1-\varpi)>0$. By continuity, for all small
$\gamma_1$ the point $u^{(0)}$ lies in the domain of $\Psi$ and $u^{(1)}\le u^{(0)}$ in both
coordinates. By (b), $\Psi$ is nondecreasing, so induction gives $u^{(k+1)}\le u^{(k)}$ for all
$k$, and the same induction shows that $u^{(k)}$ is nondecreasing in $\gamma_1$ for each $k$. So
the limit $u^{(\infty)}(\gamma_1)$ exists, solves $u=\Psi(u)+(\gamma_1,\gamma_1)$, and is
nondecreasing in $\gamma_1$; its limit as $\gamma_1\to0$ is a solution of $\Psi(u)=u$ with
nonnegative coordinates, hence $0$ by (c). Let $\bar\gamma$ be so small that
$u^{(\infty)}_1+u^{(\infty)}_0\le r/2$ at $\gamma_1=\bar\gamma$, hence for all smaller
$\gamma_1$; then $k_0$ exists for each such $\gamma_1$.
\end{proof}

\noindent For $\eps<4/9$ the iterates of $\Psi$ from $(d_\eps,1-d_\eps)$ converge to the solution of
(c) in the numerical evaluation of \texttt{status\_check.py}; at noise $0.40$, for example, to
$(0.2496,0.2936)$. So, according to this map, $54\%$ of the positions of a high block stay open
until the values of the block below are determined. We do not prove this. In a simulation with
blocks of $2000$ positions the fraction is between $54\%$ and $57\%$ after $50$ to $300$ rounds
(Section~\ref{sec:checks}).

\medskip
We return to random block bases. Let $t\ge1$, let $g_1=n_0+tm+1$ be the first gate of block $t$,
and let $W_1,W_0\subseteq\mathcal B_{t-1}$ be disjoint sets determined by $\mathcal H_{g_1-1}$
such that the positions of $W_1$ hold $1$ and those of $W_0$ hold $0$. For a gate $g$ of block
$t$ say that $g$ \emph{reads an open symbol} if $\pi_g=0$ and $J_g\in W_1\cup W_0$, or $\pi_g=1$
and $\rho_g>n_0$; and that it \emph{reads a determined one} if it reads a one and either $\pi_g=0$
and $J_g\notin W_1\cup W_0$, or $\pi_g=1$ and $\rho_g\le n_0$. Define $\hat U_{g_1-1}=1$ and, for
the gates $g$ of block $t$ in increasing order, $\hat U_g$ by the rule~\eqref{eq:statusrule} with
$\hat U=\hat U_{g-1}$, $x=x_{g-1}$ and the symbol read by $g$.

\begin{lemma}[one block of statuses]\label{lem:statusblock}
With these definitions put $a=(1-\eps)|W_1|/m$ and $b=(1-\eps)|W_0|/m$.
\begin{enumerate}[label=(\alph*),leftmargin=2.4em]
\item Let $s\ge0$. If every $s$-open position of block $t-1$ lies in $W_1\cup W_0$ and every
position up to $n_0$ is $s$-determined, then every $(s+1)$-open gate $g$ of block $t$ has
$\hat U_g=1$.
\item Let $L\ge1$ and $\nu\in(0,1)$, and let $i$ be a gate of block $t$ with $i-L\ge g_1$. Let
$\mathcal D_i\in\mathcal H_{i-L}$ be an event on which $|\delta_{t-1}-d_\eps|\le\nu$,
$|p_g-\varpi(\eps)|\le\nu$ and $(g-n_0)/g\le\nu$ for every gate $g\in(i-L,i]$. If $a<\sigma$ and
$b\le q$, then on $\mathcal D_i$, for $v\in\{0,1\}$,
\[
\PP\bigl[\hat U_i=1,\ x_i=v\bigm|\mathcal H_{i-L}\bigr]\ \le\ \psi_v(a,b)+(1-\sigma+a)^L+6L\nu .
\]
\end{enumerate}
\end{lemma}

\begin{proof}
(a) Induction on $g$. Let $g$ be $(s+1)$-open. Suppose first that $g-1$ is $(s+1)$-open; then
$\hat U_{g-1}=1$, by induction or, for $g=g_1$, by definition. If $g$ read a determined one, its
target would hold $1$ and be $s$-determined, by the two hypotheses. If $g$ is a reset this
contradicts (O2). If not, $g$ is a NOT gate whose target $g-1$ holds $1$; but $g-1$ is
$(s+1)$-open, hence $s$-open, a contradiction. So $g$ does not read a determined one, and
$\hat U_g=1$ by~\eqref{eq:statusrule}. Suppose now that $g-1$ is $(s+1)$-determined. By (O1),
$x_{g-1}=0$ and the target of $g$ is $s$-open; so the target lies in $W_1\cup W_0$ or above $n_0$,
and $g$ reads an open symbol. Then both alternatives of~\eqref{eq:statusrule} give $\hat U_g=1$.

(b) Given $\mathcal H_{g-1}$, a gate $g$ of block $t$ keeps its pointer with probability $1-\eps$,
and then $J_g$ is uniform on $\mathcal B_{t-1}$, whose values and the sets $W_1,W_0$ are
$\mathcal H_{g_1-1}$-measurable; otherwise $\rho_g$ is uniform on $\{0,\dots,g-1\}$. So the four
symbols have the conditional probabilities
\begin{align*}
\text{determined one: }&(1-\eps)\bigl(\delta_{t-1}-|W_1|/m\bigr)+\eps\,O_{n_0}/g,\\
\text{open one: }&a+\eps\,(O_{g-1}-O_{n_0})/g,\\
\text{determined zero: }&(1-\eps)\bigl(1-\delta_{t-1}-|W_0|/m\bigr)+\eps\,(n_0+1-O_{n_0})/g,\\
\text{open zero: }&b+\eps\,(g-1-n_0-O_{g-1}+O_{n_0})/g .
\end{align*}
On $\mathcal D_i$ and for $g\in(i-L,i]$ these differ from $\sigma-a$, $a$, $q-b$, $b$ by at most
$3\nu$, $\nu$, $3\nu$, $\nu$: use $\sigma=\eps\varpi+(1-\eps)d_\eps$,
$|O_{n_0}/g-p_g|\le(g-n_0)/g$ and $q=\eps(1-\varpi)+(1-\eps)(1-d_\eps)$. The pair
$(x_g,\hat U_g)$ is obtained from $(x_{g-1},\hat U_{g-1})$ and the symbol by the rule of the
status chain. So, given $\mathcal H_{i-L}$ and on $\mathcal D_i$, the pairs for
$g=i-L,\dots,i$ can be coupled with a status chain with parameters $(a,b)$, started in the same
state, so that the two disagree at any given step with conditional probability at most $6\nu$,
because the laws of the symbols are within $6\nu$ in total variation at every step. With (a) of
Proposition~\ref{prop:psi} this gives the bound.
\end{proof}

\begin{theorem}[above $4/9$]\label{thm:blockabove}
Let $\eps\in(4/9,1)$. There are $C,C',c'>0$ and $n_1$, depending only on $\eps$, such that for
every $n\ge n_1$ the following holds for the $\eps$-smoothing of a random block base with
$m=\lceil C\ln(n+1)\rceil$, $T=\lfloor c'n/m\rfloor-1$ and $n_0=n-(T+1)m$, with probability at
least $1-5/n$ over the base and the noise: $\max_iD(i)\le C'\ln n$. So the relevance-pruned
iteration finishes within $C'\ln n$ rounds, and the prefix--Jacobi iteration within
$C'\ln n+1$ applications of $G$.
\end{theorem}

\begin{proof}
The last sentence is Lemma~\ref{fact:depths} and Proposition~\ref{prop:jacobi}. Throughout,
$\lambda<1$, and ``for large $n$'' means for $n\ge n_1$ after $C$ is chosen.

\emph{Parameters.} Fix $r_0,\eta,\chi_0>0$ and integers $J,J_2\ge1$ such that, with
\[
q^+=q+r_0+2\eta+\chi_0,\qquad
\bar w=\frac{1-q+2\eta}{1-(q+2\eta)(q-2\eta)}+(q+2\eta)^{J-1},
\]
we have $q^+<1$, $\bar w\le1$ and
\begin{equation}\label{eq:theta2}
\frac{(1-\eps)\bigl(\bar w+(q^+)^{J_2}\bigr)}{1-q^+}<\theta<1
\end{equation}
for some $\theta$; this is possible because $\bar w\to1/(1+q)$ and the left-hand side tends to
$(1-\eps)/((1+q)(1-q))=\lambda$ as $r_0,\eta,\chi_0\to0$ and $J,J_2\to\infty$. Put
$\kappa=1/(\theta(1-q^+))$, and fix $\chi\le\chi_0$ and $\beta>0$ so small that
\begin{equation}\label{eq:theta3}
\frac{(1-\eps)(1+\beta)\bigl(\bar w+(q^+)^{J_2}\bigr)+\eps\chi(1+\kappa)}{1-q^+}\le\theta
\qquad\text{and}\qquad(1-\eps)(1+\beta)+\eps\chi(1+\kappa)\le1 .
\end{equation}
Let $\bar\gamma$ be as in Proposition~\ref{prop:psi}(d) for $r=r_0$, fix $\gamma_1\le\bar\gamma$
with $\gamma_1\le\min(\eta/2,\varpi/4)$ and $(1-\eps)\gamma_1\le\min(A,\eps(1-\varpi))/2$, and let
$k_0$ be as there for this $\gamma_1$. Fix $L\ge1$ and $\nu\in(0,\gamma_1]$ with
$(1-A/2)^L+12L\nu\le\gamma_1/3$. Let $t_\eta$, $C_\eta$,
$c_\eta$ be the constants of Lemma~\ref{lem:blockdens0} for $\nu$ in place of $\eta$, and let
$c'\le c_\eta$ be so small that $c'/(1-c')\le\min(\nu,\chi)$; then $(g-n_0)/g\le\min(\nu,\chi)$
for every $g\le n$. The constant $C\ge C_\eta$ is chosen large enough for the finitely many
estimates below that ask for a large $m/\ln(n+1)$. Put $t_1=t_\eta+1$ and let
$n_z=n_0+(t_1+1)m$ be the last position of block $t_1$.

\emph{The density event.} By Lemma~\ref{lem:blockdens0}, with probability at least
$1-2(n+1)^{-2}$ we have $|p_g-\varpi|\le\nu$ for every gate $g>n_0$ and
$|\delta_t-d_\eps|\le\nu$ for every block $t\ge t_\eta$. Call this event $\mathcal E$. On it, every
gate $g$ of a block $t>t_\eta$ has $|q_g-q|\le2\nu\le2\eta$.

\emph{Step 1: the positions up to $n_z$.} We claim that there is $s_z\le C_z\ln n$ such that, on
$\mathcal E$ and outside an event of probability at most $2/n$, every position up to $n_z$ is
$s_z$-determined. First, if a position $i$ is $s$-open with $s\ge1$, then there is a gate
$g\le i$ that is not a NOT gate, with no reset in $(g,i]$, whose target is $(s-1)$-open. Indeed,
let $k\le i$ be the largest gate such that $k-1$ is $s$-determined; it exists because position
$0$ is, and $k\ge2$ because $D(1)=1$. All of $k,\dots,i$ are $s$-open. If $(k,i]$ contains a
reset, let $g$ be the last one; its target is $(s-1)$-open by (O2), and a reset is not a NOT gate.
Otherwise let $g=k$ and use (O1). Iterating, an $s$-open position is the start of a dependency
chain of length $s$ in the sense of Section~\ref{sec:setting}.

Next, a lower bound for the density of ones at every gate. The first $n_0$ gates are zero gates,
and the trace on them is the trace of the base of length $n_0$ that consists of zero gates. Apply
Lemma~\ref{lem:dens2any} to that base, with no cell gates, $c=1$, $K=2$, and with its parameters
$\rho,\tau,k$ chosen so that its error is at most $\varpi/4$: there is $g_0=O(\ln n)$, with
$g_0\ge4/\varpi$, such that with probability at least $1-(n_0+1)^{-2}$ we have
$\psi_g\ge3\varpi/4$ for all $g_0\le g\le n_0$. Let $G_0$ be the intersection of this event with
$\mathcal E$. Since $\nu\le\varpi/4$, on $G_0$ we have $(O_{g-1}-1)/g\ge c_0:=\varpi/2$ for every
$g>g_0$. Condition on the base. Lemma~\ref{lem:onechain} and its proof hold with $c_0$, $G_0$ and
$w(k,P)=\eps/k+(1-\eps)\mathbf 1[J_k=P]$. Split the kernel $W$ of Lemma~\ref{lem:kernel}
accordingly, $W=W_\eps+W_b$, and let $h(P)=(P+1)^{\lambda_0}$ with $\lambda_0+1=4/c_0$. As in the
proof of Lemma~\ref{lem:kernel}, $W_\eps h\le h/4$ and $W_bh\le\Theta_bh$ with
$\Theta_b=\max(1,(1-\eps)/(\eps c_0))$. A step of a skeleton that keeps a pointer $J_k\ge g_0$
starts at a gate of a block $t'\ge1$ and leads to block $t'-1$, and no step leads to a higher
block. So a skeleton from a position up to $n_z$ has at most $t_1$ such steps, and the expected
number of dependency chains of length $\ell$ from such a position, with all positions at least
$g_0$, that occur together with $G_0$ is at most
\[
(\ell+1)^{t_1}\,\Theta_b^{t_1}\,4^{t_1-\ell}\,(n+1)^{\lambda_0} .
\]
Let $\ell_z=\lceil C_z'\ln n\rceil$ with $C_z'$ so large that this is at most $n^{-2}$ for large
$n$. Then, on $G_0$ and outside an event of probability $1/n$, no such chain exists. As the
positions of a chain decrease by at least $2$ per step, every dependency chain from a position up
to $n_z$ then has fewer than $\ell_z+g_0$ steps, and the claim holds with $s_z=\ell_z+g_0$. Let
$\mathcal Z$ be the event that every position up to $n_z$ is $s_z$-determined; so
$\PP[\mathcal E\setminus\mathcal Z]\le2/n$.

\emph{Step 2: the open fractions become small.} Let $u^{(k)}$ be the sequence of
Proposition~\ref{prop:psi}(d), let $s_1=s_z+k_0$ and $s_2=s_1+\ell_2$ with
$\ell_2=\lceil(3\ln n+\ln(1+\kappa))/\ln(1/\theta)\rceil$. For $s_z\le s\le s_2$ and a block
$t>t_1$ let $\mathcal F(s,t)$ be the event that at most $m\,u^{(k)}_1$ positions of block $t$ are
$s$-open and hold $1$, and at most $m\,u^{(k)}_0$ are $s$-open and hold $0$, where
$k=\min(s-s_z,k_0)$. We claim that on $\mathcal E\cap\mathcal Z$, outside an event of probability
at most $1/n$, all these events hold. For $s=s_z$ they hold on $\mathcal E$, because block $t$
has at most $(d_\eps+\nu)m$ ones and at most $(1-d_\eps+\nu)m$ zeros and $\nu\le\gamma_1$. Let
$s_z\le s<s_2$ and $t>t_1$, let $g_1$ be the first gate of block $t$, and let $W_1$ and $W_0$ be
the sets of $s$-open positions of block $t-1$ that hold $1$ and $0$. They are determined by
$\mathcal H_{g_1-1}$. Let $\mathcal F'\in\mathcal H_{g_1-1}$ be the event that
$|W_v|\le m\,u^{(k)}_v$ for $v=0,1$; it contains $\mathcal F(s,t-1)$ if $t-1>t_1$, and it
contains $\mathcal Z$ if $t-1=t_1$, because then the sets are empty. On $\mathcal Z$, by
Lemma~\ref{lem:statusblock}(a), every $(s+1)$-open gate $g$ of block $t$ has $\hat U_g=1$. For a
gate $i$ of block $t$ with $i-L\ge g_1$ let $\mathcal D_i\in\mathcal H_{i-L}$ be the event that
$|\delta_{t-1}-d_\eps|\le\nu$ and $|p_{i-L+1}-\varpi|\le\nu$. On it, for large $n$, the
hypotheses of Lemma~\ref{lem:statusblock}(b) hold with $2\nu$ in place of $\nu$, because $p_g$
changes by at most $2L/g$ on $(i-L,i]$. On $\mathcal F'$ we have
$a=(1-\eps)|W_1|/m\le(1-\eps)(d_\eps+\gamma_1)\le\sigma-A/2$ and
$b=(1-\eps)|W_0|/m\le(1-\eps)(1-d_\eps+\gamma_1)\le q$. So that lemma and
Proposition~\ref{prop:psi}(b) give, on $\mathcal D_i\cap\mathcal F'$ and for $v\in\{0,1\}$,
\[
\PP\bigl[\hat U_i=1,\ x_i=v\bigm|\mathcal H_{i-L}\bigr]\ \le\ \Psi_v\bigl(u^{(k)}\bigr)+(1-A/2)^L+12L\nu
\ \le\ \Psi_v\bigl(u^{(k)}\bigr)+\frac{\gamma_1}3 ,
\]
where $\Psi_v$ is the coordinate of $\Psi$ for the value $v$. Let
$Y_i=\mathbf 1[\hat U_i=1,\,x_i=v]\,\mathbf 1_{\mathcal D_i}$. For each residue class modulo $L$,
the variables $Y_i-\EE[Y_i\mid\mathcal H_{i-L}]$ with $i$ in the class are martingale
differences with values in $[-1,1]$. By the Azuma--Hoeffding inequality, applied to each class
as in the proof of Theorem~\ref{thm:block}, given $\mathcal H_{g_1-1}$ and on $\mathcal F'$,
\[
\sum_iY_i\ \le\ m\Bigl(\Psi_v\bigl(u^{(k)}\bigr)+\frac{2\gamma_1}3\Bigr)
\]
except with probability at most $L\exp\bigl(-\gamma_1^2m/(20L)\bigr)\le(n+1)^{-3}$, for $C$
large. On $\mathcal E$ every $\mathcal D_i$ occurs. Adding the at most $L\le\gamma_1m/3$ gates
with $i-L<g_1$, the numbers of $(s+1)$-open positions of block $t$ that hold $1$ and $0$ are then
at most $m(\Psi_v(u^{(k)})+\gamma_1)=m\,u^{(k+1)}_v$, which is at most $m\,u^{(\min(k+1,k_0))}_v$
because the sequence is nonincreasing. This is $\mathcal F(s+1,t)$. There are fewer than $n^2$
triples $(s,t,v)$ for large $n$, which gives the claim. Let $\mathcal F$ be the event that all
$\mathcal F(s,t)$ hold. On $\mathcal F$, for $s_1\le s\le s_2$ every block $t>t_1$ has at most
$r_0m$ positions that are $s$-open.

\emph{Step 3: no open position after $s_2$ rounds.} Let $s_1\le s<s_2$ and let $i>n_z$ be
$(s+1)$-open. As in Step~1 let $k\le i$ be the largest gate such that $k-1$ is
$(s+1)$-determined. By (O1), $x_{k-1}=0$, gate $k$ is not a NOT gate and its target is $s$-open;
by (O2), no gate of $(k,i]$ is a reset with an $s$-determined target. On $\mathcal Z$ an $s$-open
position exceeds $n_z$. Iterating from the round $s_2$ down to the round $s_1$: on $\mathcal Z$,
a position $P_0>n_z$ that is $s_2$-open is the start of a \emph{witness}
$\varsigma=(P_0,k_1,P_1,\dots,k_{\ell_2},P_{\ell_2})$, with $n_z<P_r\le k_r-2$ and $k_r\le P_{r-1}$,
for which the following event $E'_\varsigma$ occurs: for every $r$, $\jt_{k_r}=P_r$, $x_{k_r-1}=0$,
and no gate of $(k_r,P_{r-1}]$ is a reset with an $(s_2-r)$-determined target.

Three bounds for single gates. Let $g>n_z$ be a gate of a block $t$, so $t>t_1$, and let
$s\ge s_1$. (1) Let $G'_g\in\mathcal H_{g-1}$ be the event that $|\delta_{t-1}-d_\eps|\le\eta$,
$|p_g-\varpi|\le\eta$, at most $r_0m$ positions of block $t-1$ are $s$-open, and all positions up
to $n_0$ are $s$-determined. On $G'_g$, given $\mathcal H_{g-1}$, gate $g$ is a reset with an
$s$-determined target with probability at least
\[
(1-\eps)\bigl(\delta_{t-1}-r_0-\tfrac1m\bigr)+\eps\,\frac{O_{n_0}}g\ \ge\
(1-\eps)\bigl(d_\eps-\eta-r_0-\tfrac1m\bigr)+\eps\,(\varpi-\eta-\chi)\ \ge\ 1-q^+
\]
for large $m$: it keeps its pointer and reads one of the at least $(\delta_{t-1}-r_0)m-1$
positions of block $t-1$ below $g-1$ that hold $1$ and are $s$-determined, or it is resampled to
a one among the positions up to $n_0$. (2) For $P\le g-2$ the event $\{\jt_g=P\}$ is independent
of $\mathcal H_{g-1}$ and has probability
$w(g,P)=\eps/g+(1-\eps)\,\mathbf 1[P\in\mathcal B_{t-1}]/m$. (3) Call $g$ \emph{low} if
$g\le n_0+tm+J$. If $g$ is not low, then by the recursion in the proof of
Lemma~\ref{lem:oneblock}, applied with $q\pm2\eta$, we have
$\PP[x_{g-1}=0\mid\mathcal H_{g-1-J}]\le\bar w$ on the event $G''_g\in\mathcal H_{g-1-J}$ that
$|\delta_{t-1}-d_\eps|\le\eta$ and $|p_{g-J}-\varpi|\le\eta/2$, for large $n$.

For a witness $\varsigma$ let $c_r=\bar w$ if $P_r$ lies in the block below that of $k_r$ and $k_r$
is not low, and $c_r=1$ otherwise. We claim that
\begin{equation}\label{eq:witness}
\PP\bigl[E'_\varsigma\cap\mathcal E\cap\mathcal Z\cap\mathcal F\bigr]\ \le\ \prod_{r=1}^{\ell_2}(q^+)^{P_{r-1}-k_r}\,c_r\,w(k_r,P_r).
\end{equation}
The proof is that of Lemma~\ref{lem:onechain}. The gates that are constrained are those of the
intervals $(k_r,P_{r-1}]$, the gates $k_r$, and, when $c_r=\bar w$, the $J$ gates
$k_r-J,\dots,k_r-1$; these lie in the block of $k_r$, hence above $P_r$, so all these groups are
disjoint. Conditioning from the top, a gate $g$ of an interval contributes the factor $q^+$ by
(1), with the event $G'_g$ for the round $s=s_2-r$, which contains
$\mathcal E\cap\mathcal Z\cap\mathcal F$; the gate $k_r$ contributes $w(k_r,P_r)$ by (2); and when
$c_r=\bar w$ the event $x_{k_r-1}=0$ contributes $\bar w$ by (3), with
$G''_{k_r}\supseteq\mathcal E$.

\emph{The sum over witnesses.} Call a position $P>n_z$ \emph{near} if it is among the first
$J+J_2$ positions of its block, and let $f(P)=1+\kappa$ for near $P$ and $f(P)=1$ otherwise. Let
\[
(\mathcal Kf)(P)=\sum_{n_z<k\le P}(q^+)^{P-k}\sum_{n_z<P'\le k-2}c(k,P')\,w(k,P')\,f(P'),
\]
with $c(k,P')$ defined as $c_r$. For a gate $k$ of a block $t$ the inner sum is at most
$(1-\eps)\,\bar c\,\bar f+\eps\chi(1+\kappa)$, where $\bar f\le1+\kappa(J+J_2)/m\le1+\beta$ is the
mean of $f$ over block $t-1$ (for large $m$), $\bar c=\bar w$ unless $k$ is low, in which case
$\bar c=1$, and we used $\bar w\le1$ and $(k-n_z)/k\le\chi$. If $P$ is
not near, the low gates $k\le P$ are at least $J_2$ below $P$, so by~\eqref{eq:theta3}
\[
(\mathcal Kf)(P)\ \le\ \frac{(1-\eps)(1+\beta)\bigl(\bar w+(q^+)^{J_2}\bigr)+\eps\chi(1+\kappa)}{1-q^+}\ \le\ \theta=\theta f(P).
\]
If $P$ is near, then by~\eqref{eq:theta3}
$(\mathcal Kf)(P)\le\bigl((1-\eps)(1+\beta)+\eps\chi(1+\kappa)\bigr)/(1-q^+)\le1/(1-q^+)=\theta\kappa\le\theta f(P)$.
So $\mathcal Kf\le\theta f$. By~\eqref{eq:witness}, the expected number of witnesses from a
given $P_0$ for which $E'_\varsigma\cap\mathcal E\cap\mathcal Z\cap\mathcal F$ occurs is at most
$(\mathcal K^{\ell_2}\mathbf 1)(P_0)\le\theta^{\ell_2}(1+\kappa)\le n^{-3}$.

\emph{Conclusion.} By a union bound over $P_0$, on $\mathcal E\cap\mathcal Z\cap\mathcal F$ and
outside an event of probability $n^{-2}$, no position above $n_z$ is $s_2$-open; and on
$\mathcal Z$ no position up to $n_z$ is $s_z$-open. The events $\mathcal E$, $\mathcal Z$ and
$\mathcal F$ fail with probability at most $2(n+1)^{-2}+2/n+1/n$. So with probability at least
$1-5/n$ we have $\max_iD(i)\le s_2\le C'\ln n$.
\end{proof}

\begin{corollary}[the exact level on block bases]\label{cor:exact}
On random block bases over zero gates, with the parameters of
Theorems~\ref{thm:blockzero} and~\ref{thm:blockabove}, the relevance-pruned iteration needs at
least $cn/\ln n$ rounds for every $\eps<4/9$ and at most $C'\ln n$ rounds for every $\eps>4/9$,
with probability at least $1-5/n$ over the base and the noise. The constants $C$, $c'$, $c$,
$C'$ of the two theorems depend on $\eps$, so the block length and the fraction of positions in
blocks do; nothing is claimed at $\eps=4/9$.
\end{corollary}

\begin{remark}[the range between $4/9$ and $3/4$]\label{rem:above}
\begin{enumerate}[label=(\roman*),leftmargin=2.4em]
\item \emph{Two descriptions of one level.} Below $4/9$ the relevance-pruned iteration needs
$cn/\ln n$ rounds on these bases, the map $\Psi$ has a fixed point with a positive fraction of
open positions, and blocks of this kind can hold a bit (Section~\ref{sec:latch}). Above $4/9$
the open positions die out; once they are few, the proof bounds the expected number of witnesses
by a factor $\theta$ per round, for any $\theta>\lambda$. The number
$\lambda=(1-\eps)/(1-q^2)$ is the multiplier of Remark~\ref{rem:latch}(i) at the diagonal fixed
point.
\item \emph{A count with a fixed anchor does not give Theorem~\ref{thm:blockabove}.} With the
last reset as the anchor, a chain continues through every relevant gate above that reset and
through the reset itself, whether or not the reset is relevant. Let $\alpha=(1-\eps)q/(1-q^2)$
and $\beta'=(1-\eps)/(q(1+q))$ be the probabilities, in the limit, that a gate that reads a one,
respectively a zero, has kept its pointer. A computation that we do not include gives, for the
mean numbers of continuations from a position that holds $1$ and from one that holds $0$, to
positions that hold $1$ and $0$, the matrix with rows $(\alpha,\ \beta'/(1-q^2))$ and
$(\alpha,\ \beta'q^2/(1-q^2))$. Its largest eigenvalue is $1.32$ at $\eps=4/9$ and $1.19$ at
$\eps=0.50$, and it reaches $1$ only at about $\eps=0.576$; the referee's simulation of that
count agrees with these values (Section~\ref{sec:checks}). So this count grows between $4/9$ and
about $0.576$, where the iteration is fast. The gain to $\lambda$ comes from case (ii) of
Lemma~\ref{lem:status}: at a reset that is not relevant the depth is a minimum, so an open value
needs two open inputs there, and in the proof this appears as the rule that an open state ends on
every determined one.
\item \emph{Simulation.} The script \texttt{block\_unknown.c} measures the fractions of $s$-open
positions on one smoothing of a random block base; the numbers are in Section~\ref{sec:checks}.
At noise $0.46$ and $0.50$ they fall to $0$, or almost, within $200$ rounds; at noise $0.40$ more
than half of the positions of the upper blocks are still open after $300$ rounds.
\item \emph{What is not shown.} The theorem is about one iteration on random block bases. It does
not make the smoothed instance of an arbitrary base easy above $4/9$. The count of all chains
cannot do that at noise $0.72$ (Corollary~\ref{cor:barrier}); counts of fewer chains are not
excluded there, but in the model of Section~\ref{sec:barrier} the one for the last reset grows on
some repeating bases below about $0.66$. In the other direction, a search for a restoring
organ that works above $4/9$ found none: in the ideal model of Section~\ref{sec:latch}, among
all periodic chains of period at most $7$ in which every gate reads its own rail, the other rail,
position $0$, position $1$, or the position $2$, $3$, $4$ or $5$ before it in the chain, and in
which each rail is read at least once, the plain latch has the largest multiplier at noise
$0.46$, namely $0.972$ (script \texttt{latch\_local\_search.c}). Longer designs do better: in the
ideal model a design of period $24$ has two attracting states for pairs of layers up to noise
about $0.447$, numerically (Remark~\ref{rem:beyond}), and it gives hardness up to $0.4468$
(Section~\ref{sec:s24}). So $4/9$ is not the critical noise level, unless
$\mathrm P\subseteq\mathrm{BPNC}$.
To prove that some level $\eps_0$ is critical one would have to show that some iteration is fast
above $\eps_0$ on every base. By Remark~\ref{rem:mixed} it would not be the relevance-pruned
or the prefix--Jacobi iteration, and pointers of small reach would have to be treated exactly,
with a reach that the base cannot exceed by a little.
And the dependence between the reads of one window has to be controlled; the present proof
avoids this by taking the base pointers at random.
\end{enumerate}
\end{remark}

\begin{remark}[one short pointer makes both iterations slow above $4/9$; numerical]\label{rem:mixed}
Theorem~\ref{thm:blockabove} is special to block bases. Consider the following \emph{mixed base}:
zero gates, then layers of $m$ positions, $m$ even; layer $0$ consists of zero gates; in a layer
$t\ge1$ a gate at an even offset has $j_i=i-2$, and a gate at an odd offset points to a
uniformly random position at an even offset of layer $t-1$. Like a block base it is fixed before
the noise. This remark reports simulations; nothing in it is proved except the two facts in (ii).
\begin{enumerate}[label=(\roman*),leftmargin=2.4em]
\item \emph{Measurements.} One smoothing each, with $m=1000$, the layers in the last twentieth of
the positions, $124$, $249$ and $499$ layers, and seed $1$ (script \texttt{mixed\_base.c}). The
entries are $\max_iD(i)$, which is the number of rounds of the relevance-pruned iteration, the
number of applications of $G$ in the prefix--Jacobi iteration, and the same for the table--Jacobi
iteration of Section~\ref{sec:table} with reach $2$:
\begin{center}
\begin{tabular}{lccc}
\toprule
noise & $\max_iD(i)$ & prefix--Jacobi & table--Jacobi, $R=2$\\
\midrule
$0.46$ & $426$, $766$, $1622$ & $426$, $765$, $1618$ & $22$, $24$, $25$\\
$0.50$ & $313$, $630$, $1180$ & $312$, $625$, $1140$ & $20$, $20$, $22$\\
$0.52$ & $266$, $368$, $841$ & $264$, $358$, $831$ & $24$, $22$, $23$\\
$0.55$ & $156$, $158$, $161$ & $155$, $158$, $161$ & $22$, $25$, $24$\\
$0.58$ & $63$, $55$, $68$ & $58$, $52$, $67$ & $23$, $26$, $25$\\
$0.62$ & $55$, $41$, $45$ & $53$, $38$, $45$ & $25$, $26$, $27$\\
\bottomrule
\end{tabular}
\end{center}
At noise $0.46$ and $0.50$ the first two numbers are proportional to the number of layers. The
referee's simulator confirmed this with three seeds per entry, with $999$ and $1999$ layers at
noise $0.50$ ($2434$ and $2331$, then $4811$ and $4854$), with layers of $5000$ positions and
with the layers in the last half, tenth, fiftieth and two-hundredth of the positions. The rate
depends on $m$: about $2.5$ per layer for $m=1000$ and $3.1$ for $m=20000$ at noise $0.50$. On a
random block base of the same size ($499$ layers) the three numbers are $232$, $91$, $80$ at
noise $0.46$, then $91$, $79$, $79$ at $0.50$ and $52$, $45$, $45$ at $0.55$.

With $m=1000$ the growth at noise $0.52$ is irregular ($\max_iD(i)$ between $542$ and $1057$
over fifteen seeds with $499$ layers) and stops near $0.55$. That is an effect of the layer
length: the few deep positions of a short layer die out by chance. With layers of $20000$
positions the referee found, for $49$, $99$, $199$ and $399$ layers and two seeds:
\begin{center}
\begin{tabular}{lcccc}
\toprule
noise & $49$ & $99$ & $199$ & $399$ layers\\
\midrule
$0.50$ & $166$, $168$ & $311$, $323$ & $620$, $639$ & $1223$, $1209$\\
$0.52$ & $142$, $152$ & $265$, $268$ & $538$, $542$ & $1046$, $1093$\\
$0.54$ & $132$, $128$ & $241$, $232$ & $451$, $436$ & $859$, $834$\\
$0.55$ & $133$, $121$ & $201$, $212$ & $326$, $342$ & $582$, $592$\\
$0.56$ & $106$, $103$ & $143$, $148$ & $235$, $215$ & $275$, $241$\\
$0.58$ & $65$, $71$ & $90$, $88$ & $101$, $90$ & $116$, $135$\\
\bottomrule
\end{tabular}
\end{center}
So on this base $\max_iD(i)$ grows in proportion to the number of layers up to noise about
$0.55$ and stalls from about $0.56$ on. The level depends a little on the fraction of positions in
layers, because a resampled pointer that lands in a layer can pass on an open status.
\item \emph{Two facts.} Let $j$ be a parent array and $i\ge3$ a gate with $j_i=i-2$. Then
$D(i)\ge D(i-2)$. If moreover gate $i-1$ is relevant and not a NOT gate, then
$D(i)\ge1+D(j_{i-1})$. Indeed, if $x_{i-1}=0$, gate $i$ is relevant and not a NOT gate, so
$D(i)\ge1+D(i-2)$ by Lemma~\ref{lem:Dfacts}(i). If $x_{i-1}=1$, then $x_{i-2}=0$, so neither
gate $i-1$ nor gate $i$ is a reset, and $D(i)\ge D(i-1)\ge D(i-2)$ by Lemma~\ref{lem:Dfacts}(ii).
For the second claim, $x_{i-2}=0$, so gate $i$ is not a reset and $D(i)\ge D(i-1)$, and
$D(i-1)\ge1+D(j_{i-1})$ by Lemma~\ref{lem:Dfacts}(i).
\item \emph{The mechanism.} In the smoothed mixed base a gate at an even offset that keeps its
pointer inherits the depth of the previous even position, whatever the gate between them reads:
a reset at the gate between, which on a block base cuts the depth, does not cut it here. It also
inherits one more than the depth of the target of the gate between, when that gate is relevant.
The two facts alone do not account for the growth: the gates they cover pass on the depth of
$0.37$ positions of the layer below on average at noise $0.50$, and the depth also survives two
thirds of the resampled even gates. The referee measured the multiplier directly, by marking a
sparse set of even positions of one layer as infinitely deep and counting the positions of the
layers above that inherit this: it is $1.31$, $1.20$, $1.03$, $1.00$, $0.94$, $0.87$ at noise
$0.50$, $0.52$, $0.54$, $0.55$, $0.56$, $0.58$, where the same measurement on block bases returns
$\lambda$ ($0.900$ against $0.899$ at noise $0.50$). When the kept pointers $j_i=i-2$ are not
allowed to pass on an open status, the multiplier drops to $0.98$, $0.88$, $0.78$ at noise $0.46$,
$0.50$, $0.52$. So it is the depth passed through the pointers of reach $2$ that puts this base
above $1$. The prefix--Jacobi iteration reads the previous iterate through these pointers as
through all others, and needs as many rounds. The table--Jacobi iteration with reach $2$ reads
the current iterate through them; it needs about as many rounds as the zero gates below the
layers need by themselves.
\item \emph{What it means.} At noise up to about $0.55$ the relevance-pruned and the
prefix--Jacobi iteration need a number of rounds proportional to the number of layers on a base
fixed in advance, as far as simulation can show. These are the first slow bases we know for
either iteration above $4/9$; for the prefix--Jacobi iteration no lower bound is proved at any
noise level, and below $4/9$ it is slow on block bases in simulation
(Remark~\ref{rem:block}). So neither is a candidate for an algorithm that is fast above $4/9$
on every base. It does not make the smoothed mixed base hard: the table--Jacobi iteration with
reach $2$ needs about $25$ rounds on it.
\item \emph{Longer short pointers, and searches.} The same happens to the table--Jacobi
iteration when the short pointer is longer than its reach. With $j_i=i-4$ at the even offsets
and reach $2$ it needs about $340$ and $690$ rounds for $120$ and $240$ layers at noise $0.46$.
With $j_i=i-6$ and reach $5$ it needs $259$, $586$ and $1088$ rounds for $120$, $240$ and $480$
layers of $4000$ positions at noise $0.46$ (with layers of $1000$ positions the growth collapses,
which misled the first version of this remark). The effect fades as the pointer gets longer: at
noise $0.46$ a pointer of reach $8$ gave $130$, $211$, $258$ rounds for these sizes, and pointers
of reach $12$ and $20$ gave no growth. The mixed base was found by a search over layered bases
whose layers are chains with a periodic program, in a model in which a resampled read returns a
one with probability $\varpi(\eps)$ (scripts \texttt{layered\_search.c},
\texttt{layered\_design.c}); every gate of a period reads a random position with a given offset
in the layer below, or position $0$ or $1$, or the position $2$, $3$, $4$ or $5$ before it. At
noise $0.50$ and period at most $5$ the designs with the most positions of large depth all
contain gates that read $2$ or $4$ back. With the table--Jacobi iteration of reach $5$ in place
of the depth, no design of period at most $5$ needed more than $80$ rounds at noise $0.50$ with
$300$ layers of $240$ gates ($153{,}275$ designs of period $5$; script
\texttt{layered\_table\_search.c}); but this search allows local pointers only up to reach $5$,
and with layers of $240$ gates it has little power.
\end{enumerate}
\end{remark}

\begin{remark}[simulations and meaning]\label{rem:block}
\begin{enumerate}[label=(\roman*),leftmargin=2.4em]
\item \emph{Simulations.} Each entry is one smoothing of one random block base; the script is
\texttt{block\_depths.c}. Without zero gates below and with blocks of $m=1000$ positions, the
sizes $n=10^5,10^6,10^7$ give $99$, $999$ and $9999$ blocks. The pairs $\max_id(i)$ /
$\max_iD(i)$ are:

\begin{center}
\begin{tabular}{lrrr}
\toprule
noise & $n=10^5$ & $10^6$ & $10^7$\\
\midrule
$0.30$ & 122 / 107 & 1082 / 1024 & 10193 / 10037\\
$0.45$ & 129 / 98 & 1065 / 388 & 10150 / 817\\
$0.55$ & 121 / 67 & 1039 / 154 & 10067 / 238\\
$0.66$ & 102 / 60 & 428 / 84 & 1062 / 127\\
$0.70$ & 95 / 48 & 289 / 76 & 549 / 116\\
\bottomrule
\end{tabular}
\end{center}

\noindent Below $d^*\approx0.382$ both depths are close to the number of blocks; between $d^*$ and
$1-d^*$ only $d$ is. Above $1-d^*$ the depth $d$ grows more slowly than the number of blocks. With
blocks of $100$ positions, where the three sizes give $999$, $9999$ and $99999$ blocks,
$\max_id(i)$ is $286$, $675$, $1081$ at noise $0.66$ and $178$, $433$, $645$ at noise $0.70$: a
roughly constant amount per factor $10$, with a large constant because $\theta$ is close to $1$.
With the blocks in
the last tenth of $n=10^7$ positions, on top of zero gates ($m=2000$, $500$ blocks), $\max_id(i)$
is $534$, $525$, $383$, $197$, $123$ at noise $0.60$, $0.63$, $0.65$, $0.67$, $0.70$, and
$\max_iD(i)$ is $508$, $505$, $402$, $212$, $101$ at noise $0.38$, $0.42$, $0.44$, $0.46$, $0.50$,
consistent with $\eps_K\approx0.659$ and $4/9$. These runs do not locate the levels to two
digits: with the blocks in the last tenth, the prefix density is about $0.01$ above $\varpi$, and
the entries just below the levels are already under the number of blocks. All entries of this
remark use seed $1$.
\item \emph{The Kleene--prefix iteration.} Its threshold is within $0.1$ of the bound $3/4$ of
Theorem~\ref{kp:thm39}. No analysis of this iteration can go below $\eps_K$. Together with
Corollary~\ref{cor:barrier}, this leaves little room: the count of all chains cannot be bounded
for all bases at noise $0.72$, and the iteration itself fails on some base at every noise level
below $\eps_K\approx0.659$.
\item \emph{The relevance-pruned and prefix--Jacobi iterations.} The relevance-pruned iteration
fails on some base at every noise level below $4/9$. The prefix--Jacobi iteration is never slower
(Proposition~\ref{prop:jacobi}), and we have no lower bound for it. In simulation it needs about
as many rounds: on the block bases of the table with $n=10^6$ it takes $1010$, $293$, $126$, $81$
and $72$ rounds at the five noise levels, against $\max_iD(i)=1024$, $388$, $154$, $84$, $76$.
\item \emph{The interval for $\eps_c$.} Below $4/9$ no randomized NC algorithm succeeds on all
bases unless $\mathrm P\subseteq\mathrm{BPNC}$ (Theorem~\ref{thm:hard4}). The bases of that proof
are layered like block bases, and its condition is the condition $q(\eps)^2>\eps$ of
Theorem~\ref{thm:blockzero}(b) (Remark~\ref{rem:latch}). So the lower bound for the
relevance-pruned iteration on block bases and that hardness proof stop at the same level. Above
$4/9$ the relevance-pruned iteration is fast on random block bases
(Theorem~\ref{thm:blockabove}); the construction of Section~\ref{sec:latch} is not proved to fail
there, and below the level $\eps_1<0.7288$ of Section~\ref{sec:allbases} nothing is proved to work
for all bases. The construction of
Section~\ref{sec:s24} goes past $4/9$: up to $0.4468$ no such algorithm exists either, unless
$\mathrm P\subseteq\mathrm{BPNC}$, while the relevance-pruned iteration is fast on random block
bases above $4/9$ (Remark~\ref{rem:s24}(i)).
\item \emph{What is sharp.} Two statements have both sides. On block bases without zero gates
below, the Kleene--prefix iteration needs at least $cn/\ln n$ rounds below $1-d^*$ and
$O(\log n)$ above (Theorem~\ref{thm:block}(a),(c)). On block bases over zero gates, the
relevance-pruned iteration needs at least $cn/\ln n$ rounds below $4/9$ and $O(\log n)$ above
(Corollary~\ref{cor:exact}). For the relevance-pruned iteration without zero gates below and for
the Kleene--prefix iteration over zero gates only the lower bounds are proved; the simulations
suggest that $d^*$ and $\eps_K$ are the levels there. The first-moment argument of part (c) should give $O(\log n)$
rounds above $\eps_K$ for blocks over zero gates; we did not write it out.
\item \emph{Constants.} The theorems are asymptotic. The constants $C$ and $t_\eta$ of the proofs
are astronomical (the proof of Lemma~\ref{lem:blockdens} gives $t_\eta$ of order
$(16/\eta)^k$), so $n_1$ is far beyond any size that can be run, and the simulations, with blocks
of at most $2000$ positions, illustrate the statements without testing their regime.
\end{enumerate}
\end{remark}

\section{All bases: the relevance-pruned iteration below \texorpdfstring{$3/4$}{3/4}}\label{sec:allbases}

This section lowers the upper end for every base. Let $\varpi=\varpi(\eps)$, $r=1-\eps\varpi$,
$\eta_2=\eps^2\varpi(1-\varpi)$, and
\[
\chi(\eps)=(1-\eps)(1-\eta_2)-\eps\varpi,\qquad
\kappa_2(\eps)=r^2+(1-\eps)\bigl(1-\tfrac\eps2+r\bigr).
\]
Both functions are decreasing on $(0,1)$: $\varpi$ increases with $\eps$, so $r$ decreases, and
$\eta_2$ increases because $\varpi<1/2$. Let $\eps_1$ be the root of $\chi$. Then
\[
0.7287<\eps_1<0.7288\qquad\text{and}\qquad\kappa_2(\eps)<0.96\ \text{ for }\eps\ge0.7287
\]
(script \texttt{allbases\_levels.py}, in exact rational arithmetic). Theorem~\ref{thm:allbases}
says that for every $\eps>\eps_1$ the relevance-pruned iteration finishes in $O(\log n)$ rounds on
the $\eps$-smoothing of \emph{every} base. So the critical noise level is at most $\eps_1<0.7288$
(Corollary~\ref{cor:upperend}).

\emph{Where the gain comes from.} The count behind the level $3/4$ charges a step of a chain the
probability $1-\eps$ that its gate keeps its pointer, and a gate of a window the probability
$1-\eps\varpi$ that it is not a resampled reset. Remark~\ref{rem:needs} looked for more in the
\emph{reads} of the window: a gate that keeps its pointer must read a zero. A base can aim these
reads at positions whose values are tied to the chain below, and that is where those attempts
stopped. Here the gain is taken from the \emph{gate of the step} instead. In the relevance-pruned
iteration a step is taken only at a relevant gate $k$, one with $x_{k-1}=0$. The value $x_{k-1}$
depends on the noise of the gates $k-2$ and $k-1$: whatever happened before, with probability
nearly $\eta_2$ or more gate $k-2$ is resampled to a one and gate $k-1$ to a zero, and then
$x_{k-1}=1$. The rest of the chain, below the pointer of $k$, is decided by the noise at least
three positions below $k$ when that pointer has reach at least $3$. So each such step costs a
factor close to $1-\eta_2$, for every base;
nothing the base does below can remove it. A pointer of reach $2$ escapes this when it lists a
one, because then $x_{k-1}=0$ for certain. But such a gate is tied to the position two below it in
every other respect too, and a direct count over the two gates $k-1$ and $k$ shows that it does
better than the level requires (inequality~\eqref{eq:reachtwo}).

\begin{lemma}[one gate, two kinds of open positions]\label{lem:twotype}
Let $j$ be a parent array with trace $x$, let $s\ge0$, and let $k\ge2$ be a gate with pointer
$p=j_k$. Call a position \emph{$s$-open} if $D>s$ there.
\begin{enumerate}[label=(\alph*),leftmargin=2.4em]
\item If $k$ is $(s+1)$-open and $x_k=1$, then $x_{k-1}=0$, gate $k$ reads a zero, and either
$k-1$ is $(s+1)$-open, or $k$ is not a NOT gate and $p$ is $s$-open with $x_p=0$.
\item If $k$ is $(s+1)$-open and $x_k=0$, then either $x_{k-1}=0$, $k$ is not a NOT gate and $p$
is $s$-open with $x_p=1$; or $x_{k-1}=1$, $k-1$ is $(s+1)$-open, and gate $k$ is a NOT gate, or
reads a zero, or $p$ is $s$-open with $x_p=1$.
\item $D(k)\le k$, also for $k=0$ and $k=1$.
\end{enumerate}
\end{lemma}

\begin{proof}
Lemma~\ref{lem:status}. (a) $x_k=1$ means that $x_{k-1}=0$ and the read is $0$: case (iii),
where $D(k)=\max(D(k-1),1+D(p))$, or a NOT gate, case (iv), where $D(k)=D(k-1)$. (b) If
$x_{k-1}=0$ then $k$ reads a one and is not a NOT gate, case (i), and $D(k)=1+D(p)$. If
$x_{k-1}=1$: for a NOT gate or a read of zero, case (iv) gives $D(k)=D(k-1)$; for a reset, case
(ii) gives $D(k)=\min(D(k-1),1+D(p))$, so both $k-1$ and $p$ are open. (c) Induction on $k$ with
the four cases, from $D(0)=0$ and $D(1)=1$: a pointer of a gate that is not a NOT gate is at most
$k-2$.
\end{proof}

\noindent The script \texttt{allbases\_exact.py} checks (a), (b) and (c) on all parent arrays with
$2\le n\le7$. It also checks the inequalities~\eqref{eq:twoa} and~\eqref{eq:twob} below, and the
three inequalities from which~\eqref{eq:reachtwo} is put together, in the form in which the
densities are those of the instance, by summing over all outcomes of the noise for random bases
of length $7$. The referee of this section checked~\eqref{eq:reachtwo} itself in the same way, in
exact arithmetic, for all bases of length up to $8$ (Section~\ref{sec:checks}).

\begin{theorem}[every base, below $3/4$]\label{thm:allbases}
Let $\eps_1<\eps\le1$. There is a constant $C$, depending only on $\eps$, such that for every
$n\ge2$ and every base of length $n$, with probability at least $1-2/n$ over the smoothing,
$\max_iD(i)\le C\ln(n+1)$. So the relevance-pruned iteration finishes within $C\ln(n+1)$ rounds
on the $\eps$-smoothing of every base, and the prefix--Jacobi iteration within $C\ln(n+1)+1$
applications of $G$.
\end{theorem}

\begin{proof}
\emph{Constants.} For $0<c<\varpi$, $\frac14\le\beta<\frac12$ and $\lambda\ge1$ put
\begin{gather*}
r_c=1-\eps c,\qquad\eta=\eps^2c(1-c),\\
\theta_3=\frac{(1-\eps)(1-\eta)+\frac{2\eps}{\lambda+1}}{1-r_c},\qquad
\theta_2=r_c^2+(1-\eps)(1-\eps\beta+r_c)+\frac{4\eps}{\lambda+1}.
\end{gather*}
As $c\to\varpi$, $\beta\to\frac12$ and $\lambda\to\infty$, $\theta_3$ tends to
$(1-\eps)(1-\eta_2)/(\eps\varpi)$, which is less than $1$ because $\chi(\eps)<0$, and $\theta_2$
tends to $\kappa_2(\eps)<1$. Fix $c$, $\beta$ and $\lambda$ with $\theta:=\max(\theta_2,\theta_3)<1$.
By Lemma~\ref{fact:densall} the pair $(c,\beta)$ is admissible: there are $g_0=O(\ln n)$ and an
event $G\subseteq\bigcap_{g>g_0}G_g$ with $\PP[G^c]\le1/n$, where $G_g\in\cF_{g-1}$ is the event that
$(O_{g-1}-1)/g\ge c$ and $(g-O_{g-1})/g\ge\beta$. Let $G^*_k=\bigcap_{g_0<g\le k}G_g\in\cF_{k-1}$,
the sure event for $k\le g_0$, and for $v\in\{0,1\}$ let
\[
I^v_s(k)=\PP\bigl[D(k)>s,\ x_k=v,\ G^*_k\bigr],\qquad I_s(k)=I^0_s(k)+I^1_s(k).
\]
These numbers do not increase with $s$. Let $k_0$ be a constant with
$\frac14-\frac1{4(k-2)^2}\ge c(1-c)$ for $k\ge k_0$, and $k_1=\max(g_0,k_0)+4$. By
Lemma~\ref{lem:twotype}(c), $I_s(k)=0$ whenever $s\ge k$.

\emph{Three facts about single gates.} Let $k>k_1-1$; $\cF_g$ is the $\sigma$-field of the noise
of the gates up to $g$, and $D(g)$, $x_g$ and $O_g$ are $\cF_g$-measurable.

(F1) On $G_k$, given $\cF_{k-1}$, gate $k$ reads a zero with probability at most
$(1-\eps)+\eps(k-O_{k-1})/k\le r_c$; the probability that it keeps its pointer, or is resampled
to a zero, or is resampled to $k-1$, is at most $1-\eps(O_{k-1}-1)/k\le r_c$; and the probability
that it is resampled to a zero or to $k-1$ is at most $\eps(k-O_{k-1}+1)/k\le\eps(1-c)$.

(F2) On $G_{k-2}\in\cF_{k-3}$, $\PP[x_{k-1}=1\mid\cF_{k-3}]\ge\eta$. Indeed, with probability
$\eps\,O_{k-3}/(k-2)$ gate $k-2$ is resampled to a one, and then $x_{k-2}=0$ and
$O_{k-2}=O_{k-3}$; then with probability $\eps(k-1-O_{k-3})/(k-1)$ gate $k-1$ is resampled to a
zero, and $x_{k-1}=1$. With $t=O_{k-3}/(k-2)$ the product is at least $\eps^2t(1-t)$. On $G_{k-2}$
we have $t\ge c$, and $t\le\frac12+\frac1{2(k-2)}$ because ones are never adjacent; on this
interval $t(1-t)\ge c(1-c)$, by the choice of $k_0$ and $c<\frac12$.

(F3) On $G_{k-1}\cap\{x_{k-2}=0\}\in\cF_{k-2}$,
$\PP[x_{k-1}=1\mid\cF_{k-2}]\ge\eps(k-1-O_{k-2})/(k-1)\ge\eps\beta\ge\eta$: gate $k-1$ is resampled
to a zero.

\emph{Two inequalities for every gate.} For $k>k_1-1$ write $d_k=k-j_k$ for the reach of the base
pointer. We claim that for all $s\ge0$
\begin{align}
I^1_{s+1}(k)&\le r_c\,I^0_{s+1}(k-1)+(1-\eps)(1-\eta)\,\mathbf 1[d_k\ge2]\,I^0_s(j_k)
+\frac\eps k\sum_{p\le k-2}I^0_s(p),\label{eq:twoa}\\
I^0_{s+1}(k)&\le r_c\,I^1_{s+1}(k-1)+(1-\eps)\,\gamma'_k\,\mathbf 1[d_k\ge2]\,I^1_s(j_k)
+\frac{2\eps}k\sum_{p\le k-2}I^1_s(p),\label{eq:twob}
\end{align}
where $\gamma'_k=1-\eta$ if $d_k\ge3$ and $\gamma'_k=1$ if $d_k=2$.

\emph{Proof of~\eqref{eq:twoa}.} Use Lemma~\ref{lem:twotype}(a) for the smoothed instance. In the
first alternative the event $\{D(k-1)>s+1,\ x_{k-1}=0\}\cap G^*_k$ is $\cF_{k-1}$-measurable and
gate $k$ reads a zero, so by (F1) this alternative has probability at most
$r_c\,I^0_{s+1}(k-1)$, as $G^*_k\subseteq G^*_{k-1}$. In the second alternative gate $k$ is not a
NOT gate, $x_{k-1}=0$, and its pointer $p\le k-2$ is $s$-open with $x_p=0$. If gate $k$ is
resampled to $p$, this has probability at most $(\eps/k)\,I^0_s(p)$, because $\rho_k$ is independent
of $\cF_{k-1}$ and $G^*_k\subseteq G^*_p$. If it keeps its pointer, then $d_k\ge2$ and $p=j_k$; let
$E=\{x_{j_k}=0,\ D(j_k)>s\}\in\cF_{j_k}$. The probability is $(1-\eps)\,\PP[x_{k-1}=0,\ E,\ G^*_k]$,
since $\pi_k$ is independent of $\cF_{k-1}$. If $d_k\ge3$, then $E\cap G^*_{k-2}\in\cF_{k-3}$, and
(F2) gives $\PP[x_{k-1}=0,\ E,\ G^*_{k-2}]\le(1-\eta)\,\PP[E\cap G^*_{k-2}]\le(1-\eta)\,I^0_s(j_k)$. If
$d_k=2$, then $E\cap G^*_{k-1}\in\cF_{k-2}$ and $x_{k-2}=0$ on $E$, and (F3) gives the same bound.

\emph{Proof of~\eqref{eq:twob}.} Use Lemma~\ref{lem:twotype}(b). Suppose first that $x_{k-1}=1$
and $k-1$ is $(s+1)$-open. Either gate $k$ is resampled to a position $p\le k-2$ with $x_p=1$,
and then $p$ is $s$-open, which has probability at most $(\eps/k)\,I^1_s(p)$; or it keeps its
pointer, or is resampled to a zero or to $k-1$, which by (F1) has probability at most $r_c$ given
$\cF_{k-1}$, and contributes at most $r_c\,I^1_{s+1}(k-1)$. Suppose now that $x_{k-1}=0$, gate $k$
is not a NOT gate and its pointer $p$ is $s$-open with $x_p=1$. For a resampled pointer this has
probability at most $(\eps/k)\,I^1_s(p)$. For a kept pointer with $d_k\ge3$ the argument
for~\eqref{eq:twoa}, with $E=\{x_{j_k}=1,\ D(j_k)>s\}$ and (F2), gives
$(1-\eps)(1-\eta)\,I^1_s(j_k)$. For $d_k=2$ we only use $\PP[\pi_k=0]=1-\eps$.

\emph{Gates whose pointer does not have reach $2$.} Adding~\eqref{eq:twoa} and~\eqref{eq:twob}: if
$d_k\ne2$, then
\begin{equation}\label{eq:generic}
I_{s+1}(k)\ \le\ r_c\,I_{s+1}(k-1)+(1-\eps)(1-\eta)\,\mathbf 1[d_k\ge3]\,I_s(j_k)
+\frac{2\eps}k\sum_{p\le k-2}I_s(p).
\end{equation}

\emph{Gates whose pointer has reach $2$.} Let $k>k_1$ and $j_k=k-2$. We claim
\begin{equation}\label{eq:reachtwo}
I_{s+1}(k)\ \le\ c_0\,I^0_s(k-2)+c_1\,I^1_s(k-2)+r_c(1-\eps)\,\mathbf 1[d_{k-1}\ge2]\,I_s(j_{k-1})
+\frac{3\eps}{k-1}\sum_{p\le k-2}I_s(p),
\end{equation}
with $c_0=r_c^2+(1-\eps)(1-\eps\beta)$ and $c_1=\eps(1-c)\,r_c+(1-\eps)$. Note
$c_0-c_1=(1-\eps)(r_c-\eps\beta)\ge0$.

Split the event that $k$ is $(s+1)$-open according to the noise of gate $k$.

\emph{Gate $k$ is resampled}, to a position $p$. If $p=k-1$ it is a NOT gate, and it is open only
if $k-1$ is. If $p\le k-2$, the four cases of Lemma~\ref{lem:status} show that $k$ is open only
if $p$ is $s$-open, or $x_p=0$ and $k-1$ is $(s+1)$-open. By (F1) the resampled case has
probability at most
\[
\eps(1-c)\,I_{s+1}(k-1)+\frac\eps k\sum_{p\le k-2}I_s(p).
\]

\emph{Gate $k$ keeps its pointer} $k-2$; this has probability $1-\eps$ and is independent of
$\cF_{k-1}$. If $x_{k-2}=1$, then $x_{k-1}=0$, gate $k$ is a relevant reset, and it is open
exactly when $k-2$ is $s$-open: probability at most $(1-\eps)\,I^1_s(k-2)$. If $x_{k-2}=0$ and
$x_{k-1}=1$, gate $k$ reads a zero and is not relevant, case (iv): it is open exactly when $k-1$
is: probability at most $(1-\eps)\,I^1_{s+1}(k-1)$. If $x_{k-2}=0$ and $x_{k-1}=0$, case (iii): $k$ is open
exactly when $k-1$ is $(s+1)$-open or $k-2$ is $s$-open. Let
\begin{align*}
Z&=\PP\bigl[x_{k-2}=0,\ x_{k-1}=0,\ D(k-1)>s+1,\ G^*_{k-1}\bigr],\\
Y&=\PP\bigl[x_{k-2}=1,\ D(k-1)>s+1,\ G^*_{k-1}\bigr],
\end{align*}
so that $I_{s+1}(k-1)=I^1_{s+1}(k-1)+Z+Y$. By (F3), applied as in the proof of~\eqref{eq:twoa},
$\PP[x_{k-2}=0,\ D(k-2)>s,\ x_{k-1}=0,\ G^*_{k-1}]\le(1-\eps\beta)\,I^0_s(k-2)$. Altogether, with
$(1-\eps)+\eps(1-c)=r_c$,
\begin{align*}
I_{s+1}(k)\ \le\ {}&(1-\eps)\,I^1_s(k-2)+(1-\eps)(1-\eps\beta)\,I^0_s(k-2)\\
&+r_c\bigl(I^1_{s+1}(k-1)+Z\bigr)+\eps(1-c)\,Y+\frac\eps k\sum_{p\le k-2}I_s(p).
\end{align*}
Now look at gate $k-1$, whatever its pointer. By~\eqref{eq:twoa},
$I^1_{s+1}(k-1)\le r_c\,I^0_{s+1}(k-2)+(1-\eps)\mathbf 1[d_{k-1}\ge2]I^0_s(j_{k-1})
+\frac\eps{k-1}\sum_{p\le k-3}I^0_s(p)$. The event of $Z$ is the second case in the proof
of~\eqref{eq:twob} for gate $k-1$: it is a relevant reset whose pointer is an $s$-open one, so
$Z\le(1-\eps)\mathbf 1[d_{k-1}\ge2]I^1_s(j_{k-1})+\frac\eps{k-1}\sum_{p\le k-3}I^1_s(p)$. The event of
$Y$ is the first case there: $x_{k-2}=1$ and $k-1$ open, so
$Y\le r_c\,I^1_{s+1}(k-2)+\frac\eps{k-1}\sum_{p\le k-3}I^1_s(p)$. Inserting these three bounds,
and then $I^v_{s+1}(k-2)\le I^v_s(k-2)$, gives~\eqref{eq:reachtwo}; the sums add up to at most
$\frac\eps k+\frac{2\eps}{k-1}$ times $\sum_{p\le k-2}I_s(p)$.

\emph{The recursion.} Let $w_k=(k+1)^\lambda$ and, for $s\ge k_1$, $Q_s=\max_kI_s(k)/w_k$. For
$k\le k_1$ we have $I_s(k)=I_{s+1}(k)=0$. Since
$\sum_{p\le k-2}(p+1)^\lambda\le k^{\lambda+1}/(\lambda+1)$ and $w_{j_k},w_{j_{k-1}},w_{k-2}\le w_k$,
\eqref{eq:generic} gives $I_{s+1}(k)\le r_c\,I_{s+1}(k-1)+(1-r_c)\,\theta_3\,Q_s\,w_k$ for $k>k_1$
with $d_k\ne2$, and~\eqref{eq:reachtwo} gives $I_{s+1}(k)\le\theta_2\,Q_s\,w_k$ for $k>k_1$ with
$d_k=2$, because $c_1\le c_0$, $I^0_s+I^1_s=I_s$ and $k/(k-1)\le4/3$. By induction on $k$, with
$w_{k-1}\le w_k$, $I_{s+1}(k)\le\theta\,Q_s\,w_k$ for every $k$: in the first case
$r_c\theta+(1-r_c)\theta_3\le\theta$. So $Q_{s+1}\le\theta\,Q_s$, and $Q_{k_1}\le1$ gives
$Q_s\le\theta^{\,s-k_1}$ for $s\ge k_1$.

\emph{Conclusion.} For every $k$ and $s\ge k_1$,
$\PP[D(k)>s,\ G]\le I_s(k)\le\theta^{\,s-k_1}(n+1)^\lambda$. With
$s^*=k_1+\lceil(\lambda+2)\ln(n+1)/\ln(1/\theta)\rceil$ this is at most $(n+1)^{-2}$, so outside
$G^c$ and an event of probability at most $1/n$ every position has $D\le s^*$. As
$k_1=O(\ln n)$, $s^*\le C\ln(n+1)$. The statements about the iterations are
Lemma~\ref{fact:depths} and Proposition~\ref{prop:jacobi}.
\end{proof}

\begin{corollary}[the upper end]\label{cor:upperend}
For every $\eps>\eps_1$ the prefix--Jacobi iteration, stopped after $\lceil C\ln(n+1)\rceil+1$
applications of $G$ and with the output $\bot$ unless it has reached a fixed point, is a
deterministic errorless algorithm of depth $O(\log n)$ and polynomial size with unbounded fan-in
that outputs the trace with probability at least $1-2/n$ on the $\eps$-smoothing of every base. So
the critical noise level satisfies $\eps_c\le\eps_1<0.7288$; with Corollary~\ref{cor:gaps24},
\[
0.4468<\eps_S\le\eps_c\le\eps_1<0.7288\qquad\text{unless }\mathrm P\subseteq\mathrm{BPNC},
\]
and $\inf E_K\in\{0\}\cup[\eps_S,\eps_1]$ for every $K$. The threshold of the relevance-pruned
iteration (Corollary~\ref{cor:thresholds}) lies in $[4/9,\eps_1]$.
\end{corollary}

\begin{proof}
One application of $G$ is in $\mathrm{AC}^0$ (Section~\ref{sec:upper}), a fixed point of $G$ is the
trace (Proposition~\ref{prop:jacobi}(a)), and by Theorem~\ref{thm:allbases} the iteration reaches
it in time with probability at least $1-2/n$. The lower bound $4/9$ for the relevance-pruned
iteration is Corollary~\ref{cor:thresholds}.
\end{proof}

\begin{remark}[what this does and does not give]\label{rem:allbases}
\begin{enumerate}[label=(\roman*),leftmargin=2.4em]
\item \emph{Against Remark~\ref{rem:needs}.} The inequalities written there for the counts
$A_t,B_t$ drop the indicator of $x_{q-1}=0$ at the gate of a step. Keeping it, and conditioning
three positions below the gate instead of one, gives~\eqref{eq:twoa} and~\eqref{eq:twob}. The
assignment that makes the old inequalities tight, in which the depth of a position is tied to its
value, does not help against (F2): the value $x_{k-1}$ in question lies above the target, and its
last two gates have fresh noise. The proof uses that the base is fixed in advance exactly here:
the target $j_k$ is known before the noise, so the event $E$ is decided by the noise of the gates
up to $k-3$. Against the
adaptive bases of Proposition~\ref{prop:adaptive} this fails, as it must.
\item \emph{Pointers of reach $2$.} Without~\eqref{eq:reachtwo}, treating the two kinds of open
positions separately with the worst factor for each, the same proof gives the level $0.7394$, the
root of $\bigl(r+(1-\eps)(1-\eta_2)\bigr)\bigl(r+1-\eps\bigr)=1$: a pointer of reach $2$ that lists
a one, at a gate with output $0$, is then charged nothing. That bound takes the worst pointer
separately for the gates with output $1$ (reach at least $3$) and for the gates with output $0$
(reach $2$), although a base must choose the pointer of a gate before its output is known.
Inequality~\eqref{eq:reachtwo} looks at the gate and its predecessor together.
\item \emph{Longer pointers.} For an arbitrary base only two facts about densities are
available: ones have density at least $c<\varpi$, and zeros at least $\beta<1/2$. The bound (F2)
uses both for one density $t$, through $t(1-t)\ge c(1-c)$. It conditions three positions below
the gate, which is all that a pointer of reach $3$ allows, and there it is attained: if gate
$k-2$ is a NOT gate, gate $k-1$ has pointer $0$ and $x_{k-3}=0$, then $x_{k-1}=1$ exactly when
gate $k-2$ is resampled to a one and gate $k-1$ to a zero. For a pointer of reach $4$ the same bound
remains when the target holds a one, because then $x_{k-3}=0$. For a pointer of reach $d\ge5$ one
can condition $d$ positions below the gate. Iterating the two one-gate bounds $\eps(1-t)$ and $1-\eps t$
for the probability of a one after a zero gives a larger $\eta$ there, and with it the level
$0.7208$ for bases all of whose pointers have reach $1$, $2$ or at least $5$, and about $0.7185$
in the limit of long pointers (floating point; we have not written this out; the referee of this
section found the same numbers). To use it for every base one would need a count for the pointers
of reach $3$ and $4$ in the manner of~\eqref{eq:reachtwo}. So this route ends near $0.72$.
\item \emph{Not sharp.} The slowest bases we know for this iteration are the mixed bases of
Remark~\ref{rem:mixed}, up to noise about $0.55$ in simulation. The proof charges a gate of a
window only the resampled resets. In simulation at noise $0.74$ the base $j_i=i-2$ is the slowest
we found: $\max_iD(i)$ grows by about $13$ per unit of $\ln n$ (means $58$, $87$, $113$, $145$
over five smoothings for $n=10^4,\dots,10^7$, and $174$ over two for $n=10^8$); at noise $0.73$
the means are $54$, $84$, $117$ for $n=10^4,10^5,10^6$, against $45$, $63$, $84$ for $j_i=i-3$ and
$51$, $77$, $103$ for the base that alternates reach $2$ and $3$.
\item \emph{Constants.} The constant $C$ is astronomical near $\eps_1$: the proof needs the density
of ones within a small distance of $\varpi$, and the constants of Lemma~\ref{fact:densall} grow
like a high power of the inverse of that distance. The theorem is asymptotic.
\item \emph{Which iteration.} The proof needs relevance, so it says nothing about the
Kleene--prefix iteration, whose threshold stays in $[\eps_K,3/4]$
(Corollary~\ref{cor:thresholds}); and the count of all chains cannot work at $0.72$
(Corollary~\ref{cor:barrier}). The logspace algorithm of Section~\ref{sec:logspace} follows
Kleene's derivations and is proved for $\eps>3/4$ only. Between $\eps_1$ and $3/4$ we have depth
$O(\log n)$ with unbounded fan-in, hence $\mathrm{NC}^2$, not logspace.
\item \emph{The interval.} With Section~\ref{sec:s24} the critical noise level lies in
$[0.4468,0.7288]$, unless $\mathrm P\subseteq\mathrm{BPNC}$. Section~\ref{sec:layered} shows that the
hardness side cannot pass $0.455$ by layered embeddings; whether the algorithmic side can come
down to that level for every base is open.
\end{enumerate}
\end{remark}

\section{A table iteration: short pointers exactly, long pointers by iteration}\label{sec:table}

The block bases of Section~\ref{sec:block} defeat the two iterations of Section~\ref{sec:iterations} at every noise
level below $\eps_K$ and $4/9$. Their base pointers are short: they reach at most twice the block
length, which is $O(\log n)$. This section gives an iteration that treats all pointers of reach at
most $R$ exactly inside each round, by composing tables with $2^R$ entries, and iterates only
over the longer pointers. It is never slower than the bound $\max_iD(i)+1$ of the prefix--Jacobi
iteration, it needs at most $\lceil n/(R+1)\rceil+1$ rounds on every instance, and on the
smoothing of every base of reach $R$ it needs $O(\sqrt{\eps n}+\log n)$ rounds
(Theorem~\ref{thm:table}). On the block bases of Theorems~\ref{thm:block}
and~\ref{thm:blockzero} this is $O(\sqrt n)$ rounds of polynomial work, where the iterations of
Section~\ref{sec:iterations} need $\Omega(n/\log n)$ rounds below the noise levels of those theorems. We prove no lower
bound for the new iteration. In simulation at low noise its number of rounds grows like $\sqrt n$,
so it does not appear to be an algorithm of polylogarithmic parallel time for those bases, and we
found none (Remark~\ref{rem:table}). For all bases there is none up to noise $0.4468$ unless
$\mathrm P\subseteq\mathrm{BPNC}$ (Theorems~\ref{thm:hard4} and~\ref{thm:hards24}); random block bases do not encode
circuits, and for them the question is open. On the mixed base of Remark~\ref{rem:mixed}, which
adds one pointer of reach $2$ to a block base, the relevance-pruned and the prefix--Jacobi
iteration are slow up to noise about $0.55$ in simulation, and this iteration with reach $2$
needs about $25$ rounds; a short pointer that is longer than the reach makes this iteration slow
in the same way, at least up to reach $6$ at noise $0.46$. The two ingredients are exact maps across a cut for instances of bounded
width and a Jacobi iteration over long reads.

\begin{definition}[table--Jacobi iteration]\label{def:table}
Let $R\ge1$ and let $j$ be a parent array of length $n$. Call gate $i$ \emph{short} if
$j_i\ge i-R$, and \emph{long} otherwise. For $y\in\{0,1\}^{n+1}$ with $y_0=1$ let $G_R(y)=z$, where
$z_0=1$ and, for $i\ge1$, $z_i=\NOR(z_{i-1},r_i)$ with $r_i=z_{j_i}$ if gate $i$ is short and
$r_i=y_{j_i}$ if it is long. Put $y^{(0)}=(1,0,\dots,0)$ and $y^{(s+1)}=G_R(y^{(s)})$. The iteration
stops at the first $s$ with $y^{(s+1)}=y^{(s)}$ and outputs $y^{(s)}$.
\end{definition}

\noindent For $R=1$ the short gates are the NOT gates, and $G_1$ is the map $G$ of
Definition~\ref{def:jacobi}.

\emph{One application is a prefix computation.} Put $z_p=0$ for $p<0$ and
$\sigma_i=(z_{i-R+1},\dots,z_i)\in\{0,1\}^R$. Given $y$, gate $i$ defines a map $\tau_i$ on
$\{0,1\}^R$ with $\sigma_i=\tau_i(\sigma_{i-1})$: it drops the first coordinate and appends
$\NOR(z_{i-1},r_i)$, where $z_{i-1}$ is the last coordinate of $\sigma_{i-1}$ and $r_i$ is either
the coordinate $z_{j_i}$ of $\sigma_{i-1}$ (a short gate) or the constant $y_{j_i}$ (a long gate).
So $\sigma_i=\tau_i\circ\dots\circ\tau_1(\sigma_0)$ with $\sigma_0=(0,\dots,0,1)$. All these
prefixes are computed by a parallel prefix circuit over the monoid of maps on $\{0,1\}^R$, each map
stored as a table with $2^R$ entries: $O(n)$ compositions in $O(\log n)$ layers, a composition
being $2^R$ table look-ups. For $R=O(\log n)$ one application has polynomial size and depth
$O(\log^2n)$.

\begin{theorem}[table--Jacobi iteration]\label{thm:table}
Let $j$ be a parent array of length $n$ with trace $x$, and let $R\ge1$.
\begin{enumerate}[label=(\alph*),leftmargin=2.4em]
\item If $G_R(y)=y$ then $y=x$.
\item If $D(i)\le t\le s$ then $y^{(s)}_i=x_i$. Hence the iteration applies $G_R$ at most
$\max_iD(i)+1$ times.
\item If $y$ agrees with $x$ on $[0,c]$, then $G_R(y)$ agrees with $x$ on $[0,g-1]$, where $g$ is the
least long gate with $j_g>c$, and $g=n+1$ if there is none. Hence the iteration applies $G_R$ at
most $\lceil n/(R+1)\rceil+1$ times, on every parent array.
\item Let $\eps\in(0,1]$ and let $j$ be a base in which every gate has $j_i\ge i-R$ or $j_i=0$. On
the $\eps$-smoothing of $j$, with probability at least $1-n^{-4}$, the iteration applies $G_R$ at
most $17\sqrt{\eps n}+35\ln n+6$ times.
\end{enumerate}
\end{theorem}

\begin{proof}
(a) If $G_R(y)=y$ then $y_0=1$ and $y_i=\NOR(y_{i-1},y_{j_i})$ for every $i\ge1$, whether gate $i$
is short or long. The trace is the only solution.

(b) Repeat the proof of Proposition~\ref{prop:jacobi}(b). The only change is that a short gate
reads $z=y^{(s)}$ at its pointer and not $y^{(s-1)}$. In that proof every gate whose read matters,
namely the anchor $r$ and the relevant gates $p\in(r,i]$ that are not NOT gates, has a pointer $q$
with $D(q)\le t$, where $D(i)\le t+1\le s$. The induction hypothesis for $t$ holds for both
$s-1\ge t$ and $s\ge t$, so $y^{(s-1)}_q=y^{(s)}_q=x_q$, and the gate reads the right value in
either case.

(c) By induction on $i<g$: $z_0=1=x_0$, and for $1\le i<g$ a short gate reads $z_{j_i}=x_{j_i}$ by
induction, while a long gate has $j_i\le c$ and reads $y_{j_i}=x_{j_i}$. So
$z_i=\NOR(x_{i-1},x_{j_i})=x_i$. A long gate $g$ has $j_g\le g-R-1$, so $j_g>c$ implies
$g\ge c+R+2$. Since $y^{(0)}$ agrees with $x$ on $[0,0]$, it follows that $y^{(s)}$ agrees with $x$
on $[0,s(R+1)]$. So $y^{(s)}=x$ for $s=\lceil n/(R+1)\rceil$, and the next application stops the
iteration.

(d) In the smoothed instance the long gates are zero gates that keep their pointer $0$, and
resampled gates $g$ with $\rho_g<g-R$. Let $c_0=0$ and $c_{s+1}=g_s-1$, where $g_s$ is the least
gate $g$ with $\pi_g=1$ and $c_s<\rho_g<g-R$, and $g_s=n+1$ if there is none. By (c) and
induction, $y^{(s)}$ agrees with the trace on $[0,c_s]$; note $c_1\ge R+1\ge2$. The gate
$g_{s-1}=c_s+1$ has $\rho\le c_s$, so $g_s$ is determined by the noise of the gates above
$c_s+1$, which is independent of the noise that determines $c_s$. For $\delta\ge0$ and $c_s=c\ge2$,
\[
\PP[g_s\le c+\delta+1\mid c_0,\dots,c_s]\ \le\ \sum_{u=1}^{\delta}\frac{\eps\,u}{c+1+u}
\ \le\ \frac{\eps\,\delta(\delta+1)}{2c},
\]
because gate $c+1+u$ is resampled to one of the $u$ positions in $(c,c+u]$ with probability
$\eps u/(c+1+u)$. With $\delta=\lfloor\sqrt{c/(2\eps)}\rfloor$ the bound is at most
$\frac14+\frac12\sqrt{\eps/(2c)}\le\frac12$. Call step $s$ \emph{good} if
$c_{s+1}-c_s>\sqrt{c_s/(2\eps)}$ or $c_{s+1}=n$. A step that is not good has a gate that qualifies
among the gates $c_s+2,\dots,c_s+\delta+1$, so each step with $2\le c_s<n$ is good with
conditional probability at least $1/2$. After a good step from $c$ to $c'<n$, either $c'>4c$, or
$\sqrt{c'}-\sqrt c=(c'-c)/(\sqrt{c'}+\sqrt c)\ge\sqrt{c/(2\eps)}/(3\sqrt c)=1/(3\sqrt{2\eps})$. So
at most $3\sqrt{2\eps n}+\log_4n+2$ good steps can occur while $c_s<n$, the last of them reaching
$n$. Let $S=\lceil12\sqrt{2\eps n}+4\log_4n+32\ln n\rceil$. The number of good steps among steps
$1,\dots,S$, counting a step with $c_s\ge n$ as good, dominates a binomial variable with
parameters $S$ and $1/2$, so by Hoeffding's inequality it is at least $S/4$ except with
probability $e^{-S/8}\le n^{-4}$; and $S/4>3\sqrt{2\eps n}+\log_4n+2$. Then $c_{S+1}\ge n$, the
iterate $y^{(S+1)}$ is the trace, and the iteration stops after at most $S+2$ applications. Finally
$12\sqrt2\le17$ and $4\log_4n+32\ln n\le35\ln n$.
\end{proof}

\begin{corollary}[block bases]\label{cor:tableblock}
Consider a block base as in Definition~\ref{def:block}, with blocks of $m$ positions, and
$R=2m$. For every $\eps\in(0,1]$, on its $\eps$-smoothing the table--Jacobi iteration finds the
trace within $17\sqrt{\eps n}+35\ln n+6$ applications of $G_{2m}$, with probability at least
$1-n^{-4}$. For $m=\lceil C\ln(n+1)\rceil$, as in Theorems~\ref{thm:block}
and~\ref{thm:blockzero}, each application is a parallel prefix over tables with at most
$4(n+1)^{2C}$ entries. For the random block bases of those theorems, with their parameters, the
Kleene--prefix and relevance-pruned iterations need at least $cn/\ln n$ rounds below the noise
levels stated there ($1-d^*$ and $d^*$ without zero gates below, $\eps_K$ and $4/9$ over zero
gates), with probability at least $1-3/n$ over the base and the noise. The work is polynomial
only formally: $C$ is the constant of those theorems (Remark~\ref{rem:block}(vi)).
\end{corollary}

\begin{proof}
A gate $i$ of a block $t\ge1$ has $j_i\in\mathcal B_{t-1}$, so $i-j_i\le2m-1$, and every other gate
has $j_i=0$. So Theorem~\ref{thm:table}(d) applies with $R=2m$, and $2^{2m}\le4e^{2C\ln(n+1)}$.
\end{proof}

\begin{remark}[what this gives, and what it does not]\label{rem:table}
\begin{enumerate}[label=(\roman*),leftmargin=2.4em]
\item \emph{Simulations.} The bases are random block bases without zero gates below, with blocks
of $100$ positions; $R=200$ and the seed is $1$. The script simulates $G_R$ sequentially and
counts applications (\texttt{table\_jacobi.c}). Applications of $G_{200}$ until the fixed point:

\begin{center}
\begin{tabular}{lrrrrr}
\toprule
noise & $n=10^5$ & $4\cdot10^5$ & $1.6\cdot10^6$ & $6.4\cdot10^6$ & $2.56\cdot10^7$\\
\midrule
$0.15$ & 38 & 80 & 159 & 329 & 555\\
$0.30$ & 37 & 57 & 96 & 131 & 175\\
$0.50$ & 27 & 41 & 53 & 62 & \\
$0.70$ & 23 & 32 & 39 & 48 & \\
\bottomrule
\end{tabular}
\end{center}

\noindent On the same instances the prefix--Jacobi iteration ($R=1$) takes $988$ and $3979$
applications at noise $0.15$ and $592$ and $1294$ at noise $0.30$, for $n=10^5$ and $4\cdot10^5$,
where there are $1000$ and $4000$ blocks. The chain $c_s$ of the proof of (d) has $152$, $394$,
$897$, $1939$ and $4139$ steps at noise $0.30$, so the iteration is $4$ to $24$ times faster than
that chain at these sizes, more at the larger ones. At noise $0.30$ and $n=1.6\cdot10^6$, blocks of $25$, $50$, $100$, $200$,
$400$ and $800$ positions (with $R$ twice the block length) give $66$, $75$, $96$, $109$, $94$ and
$83$ applications.
\item \emph{It does not appear to be a polylogarithmic algorithm.} We prove no lower bound for
this iteration; (d) is an upper bound. The evidence is one block length at one noise level: at
noise $0.15$ the count doubles with every factor $4$ in $n$ up to $6.4\cdot10^6$, as $\sqrt n$
does, and grows by $1.69$ at the last step (seeds $2$ and $3$ give $530$ and $549$ there, so this
slowdown is systematic). The regime of Section~\ref{sec:block}, with $m=C\ln n$, is not tested.
Heuristically, the reason is the one behind (d): a long read into the part that is not yet correct
returns a wrong value about half the time, and below the level of the relevance-pruned iteration
the damage of a wrong read spreads through the blocks above it with positive probability (it
follows the chains counted by $D$). So the correct part advances by about $\sqrt{c/\eps}$ per
round from $c$, up to a constant. At noise $0.30$ with blocks of $100$ positions the growth is slower at these
sizes; the theorems of Section~\ref{sec:block} concern much longer blocks. At noise $0.50$ and
$0.70$, above the level $d^*$, the count grows by a roughly constant amount per factor $4$.
\item \emph{Exact maps across a cut.} Exact maps across a cut alone run out of memory on wide
instances, because smoothing puts long reads into the live set, and a windowed Jacobi iteration
alone is slowed by wrong long reads. The table--Jacobi iteration keeps the long reads out of the
state and iterates over them. Blocks of $\Theta(\log^2n)$ positions would need tables of
quasi-polynomial size; the bases of Section~\ref{sec:block} have blocks of $\Theta(\log n)$
positions.
\item \emph{Below $4/9$.} This iteration was an attempt at the range between $1/4$ and $4/9$,
made before Section~\ref{sec:latch} was found. By Theorems~\ref{thm:hard4} and~\ref{thm:hards24},
for every $\eps\le0.4468$ no randomized NC algorithm succeeds on the smoothing of every base unless
$\mathrm P\subseteq\mathrm{BPNC}$. The bases of these proofs have listed pointers of every length,
so Theorem~\ref{thm:table}(d) does not apply to them. For bases of reach $O(\log n)$, such as the
random block bases of Section~\ref{sec:block}, no hardness is known, and we found no algorithm of
polylogarithmic parallel time. On block bases the obstacle, as far as we understand it, has two
parts, neither of which is proved: below the levels of Section~\ref{sec:block} the blocks do not
seem to forget wrong inputs, and the resampled pointers form stacked chains of length of order
$\sqrt{\eps n}$, which is how far the argument for (d) goes. The iteration is faster than that
chain by a constant factor in the simulations and shows the same growth at noise $0.15$. By a
rough estimate, keeping in the state the $O(\log n)$ long reads just above the correct part,
instead of iterating over them, would save a factor of order $\sqrt{\log n}$ only. So
Section~\ref{sec:table} shows only that the lower bound $n/\log n$ for the iterations of Section~\ref{sec:iterations} is
not a lower bound for parallel time on these bases.
\item \emph{Consequences of (b) and (c).} By (b) every round bound for the prefix--Jacobi iteration
in this paper holds for the table--Jacobi iteration with any $R$. By (c), for
$R=\lceil\log_2n\rceil$ every parent array with $n>64$ is evaluated in $n/\log_2n$ rounds of polynomial work;
this is the classical saving of a logarithmic factor by tables and says nothing new about
worst-case instances.
\end{enumerate}
\end{remark}

\subsection{Sets of windows: a sound variant}\label{sec:setval}

The table--Jacobi iteration guesses the long reads and is right only at its fixed point. The
following variant never guesses. It keeps, for every position, the \emph{set} of windows that are
possible when the long reads into positions that are not yet determined are left free,
separately for each gate that reads them. It is sound in every round, its rounds are bounded by
the depth $D$ and by a depth $D_R$ in which short pointers cost nothing, and it bounds the
table--Jacobi iteration from above (Proposition~\ref{prop:setval}). On the mixed base with reach
$2$ it is as fast as the table--Jacobi iteration, while $D_2$ grows with the number of layers
(Remark~\ref{rem:setval}); so it is a monotone object that sees the healing of wrong reads.

\begin{definition}[set-valued table iteration]\label{def:setval}
Let $R\ge1$, let $j$ be a parent array of length $n$, and call gates short and long as in
Definition~\ref{def:table}. A \emph{window} is a vector $\sigma\in\{0,1\}^R$; at position $i$ it
stands for $(z_{i-R+1},\dots,z_i)$, with $z_p=0$ for $p<0$. Let $K\subseteq\{0,\dots,n\}$ with
$0\in K$, and let $v\in\{0,1\}^K$. Put $S_0=\{\sigma_0\}$ with $\sigma_0=(0,\dots,0,1)$, and for
$i\ge1$
\[
S_i=\bigl\{\tau(\sigma,b):\ \sigma\in S_{i-1},\ b\in B_i(\sigma)\bigr\},
\]
where $\tau(\sigma,b)$ drops the first coordinate of $\sigma$ and appends $\NOR(a,b)$, $a$ being
the last coordinate of $\sigma$, and where the set of reads $B_i(\sigma)$ consists of the
coordinate $z_{j_i}$ of $\sigma$ if gate $i$ is short, is $\{v_{j_i}\}$ if gate $i$ is long and
$j_i\in K$, and is $\{0,1\}$ if gate $i$ is long and $j_i\notin K$. Position $i$ is
\emph{determined} if all windows of $S_i$ have the same last coordinate $w_i$. Let
$\Phi_R(K,v)=(K',w)$, where $K'$ is the set of determined positions. The iteration starts from
$K^{(0)}=\{0\}$, $v^{(0)}_0=1$, puts $(K^{(s+1)},v^{(s+1)})=\Phi_R(K^{(s)},v^{(s)})$, and stops at
the first $s$ with $K^{(s)}=\{0,\dots,n\}$, with output $v^{(s)}$.

Let $D_R(0)=0$, $D_R(1)=1$, and for $k\ge2$ let $D_R(k)$ be given by the four cases of
Lemma~\ref{lem:status} with $D_R$ in place of $D$ and with $c_k+D_R(p)$ in place of $1+D(p)$,
where $c_k=0$ if gate $k$ is short and $c_k=1$ if it is long.
\end{definition}

\begin{proposition}[set-valued table iteration]\label{prop:setval}
Let $j$ be a parent array of length $n$ with trace $x$, and $R\ge1$.
\begin{enumerate}[label=(\alph*),leftmargin=2.4em]
\item \emph{Sound and monotone.} For every $s$, $v^{(s)}_p=x_p$ for all $p\in K^{(s)}$, and
$K^{(s)}\subseteq K^{(s+1)}$.
\item \emph{Depth.} $D_R\le D$; every $i$ with $D(i)\le s$ lies in $K^{(s)}$, and every $i$ with
$D_R(i)\le s$ lies in $K^{(s+1)}$. So the iteration stops after at most
$\min\bigl(\max_iD(i),\ \max_iD_R(i)+1\bigr)$ rounds.
\item \emph{It bounds the table--Jacobi iteration.} For the iterates $y^{(s)}$ of
Definition~\ref{def:table}, $y^{(s)}_i=x_i$ for all $i\in K^{(s)}$. So the table--Jacobi iteration
applies $G_R$ at most once more than the set-valued iteration has rounds.
\item \emph{Cost.} One round is a prefix computation over Boolean matrices of size
$2^R\times2^R$. For $R=O(\log n)$ it has polynomial size and depth $O(\log^2n)$.
\item \emph{Monotone in the reach.} If $R\le R'$, then for every $s$ the set $K^{(s)}$ for the
reach $R$ is contained in the set $K^{(s)}$ for the reach $R'$. So the number of rounds does not
increase with the reach.
\item \emph{The bounds of Theorem~\ref{thm:table}.} If $[0,c]\subseteq K^{(s)}$, then
$[0,g-1]\subseteq K^{(s+1)}$, where $g$ is the least long gate with $j_g>c$, and $g=n+1$ if there
is none. So the iteration stops after at most $\lceil n/(R+1)\rceil$ rounds on every parent
array, and on the $\eps$-smoothing of a base as in Theorem~\ref{thm:table}(d) after at most
$17\sqrt{\eps n}+35\ln n+6$ rounds, with probability at least $1-n^{-4}$.
\end{enumerate}
\end{proposition}

\begin{proof}
Write $\omega_i=(x_{i-R+1},\dots,x_i)$ for the window of the trace, with $x_p=0$ for $p<0$, and
$S^{(s)}_i$ for the sets computed from $(K^{(s-1)},v^{(s-1)})$ in round $s\ge1$.

(a) Suppose $v^{(s-1)}$ agrees with $x$ on $K^{(s-1)}$; this holds for $s=1$. Then
$\omega_i\in S^{(s)}_i$ for all $i$, by induction on $i$: $\omega_0=\sigma_0$, and
$\omega_i=\tau(\omega_{i-1},x_{j_i})$ with $x_{j_i}\in B_i(\omega_{i-1})$ in each of the three
cases. So a determined position has $w_i=x_i$, and $v^{(s)}$ agrees with $x$ on $K^{(s)}$.
For the inclusion, $0\in K^{(1)}$ because $S_0$ has one window. If $K^{(s-1)}\subseteq K^{(s)}$,
then the sets of reads of round $s+1$ are contained in those of round $s$, since the values agree
with $x$; so $S^{(s+1)}_i\subseteq S^{(s)}_i$ by induction on $i$, and these sets are not empty.
Hence a position that is determined in round $s$ is determined in round $s+1$.

(b) $D_R\le D$ follows by induction on $k$ from Lemma~\ref{lem:status}, because the four
expressions are nondecreasing in their arguments, $c_k\le1$, and $D(0)=0$, $D(1)=1$. For the
second claim we show by induction on $s\ge0$, and for fixed $s$ on $i$, that $D_R(i)\le s$
implies $i\in K^{(s+1)}$. For $i=0$ this is (a). Gate $1$ is a NOT gate, hence short, and
$S^{(s+1)}_1=\{\tau(\sigma_0,1)\}$, so $1\in K^{(s+1)}$. Let $i\ge2$, $p=j_i$, and consider round
$s+1$. Call a read \emph{settled} if $c_i+D_R(p)\le s$. A settled read has the value $x_p$ for
every window of $S^{(s+1)}_{i-1}$: if gate $i$ is long, then $D_R(p)\le s-1$, so $s\ge1$ and
$p\in K^{(s)}$ by the hypothesis for $s-1$, and the read is $v^{(s)}_p=x_p$; if it is short, then
$D_R(p)\le s$ and $p<i$, so $p\in K^{(s+1)}$ by the hypothesis for $p$, that is, all windows of
$S^{(s+1)}_p$ end with $x_p$, and every window of $S^{(s+1)}_{i-1}$ arises from one of them by
fewer than $R$ steps and still has this coordinate. Likewise, if $D_R(i-1)\le s$, then all
windows of $S^{(s+1)}_{i-1}$ end with $x_{i-1}$. Now go through the cases for a gate with
$D_R(i)\le s$. (i) A relevant reset: the read is settled and equals $1$, so every window gets the
new coordinate $0=x_i$. (ii) A reset that is not relevant: the read is settled and equals $1$, or
all windows end with $x_{i-1}=1$; in both cases the new coordinate is $0$. (iii) Relevant, not a
reset, not a NOT gate: the read is settled and equals $0$, and all windows end with $0$, so the
new coordinate is $1=x_i$. (iv) Otherwise $D_R(i-1)\le s$ and all windows end with $x_{i-1}$; if
$x_{i-1}=1$ the new coordinate is $0$, and if $x_{i-1}=0$ the gate is a NOT gate, whose read is
that coordinate, and the new coordinate is $1$. So $i\in K^{(s+1)}$.

The claim for $D$ is the same induction without the shift. For $i\ge1$ we have $D(i)\ge1$, by
induction on $i$ from Lemma~\ref{lem:status} and $D(1)=1$, so only position $0$ has depth $0$,
and it lies in $K^{(0)}$. Let $s\ge1$ and $D(i)\le s$, and consider round $s$. Now every gate
costs a round: a read is settled if $1+D(p)\le s$, and then $p\in K^{(s-1)}$ by the hypothesis
for $s-1$. For a long gate the read is $v^{(s-1)}_p=x_p$; for a short gate
$p\in K^{(s-1)}\subseteq K^{(s)}$ by (a), so all windows of $S^{(s)}_p$ end with $x_p$. The four
cases are as before, and $i\in K^{(s)}$.

(c) By induction on $s$; for $s=0$, $y^{(0)}_0=1$. Let $z=G_R(y^{(s)})$ and
$\zeta_i=(z_{i-R+1},\dots,z_i)$. Then $\zeta_i\in S^{(s+1)}_i$ for all $i$, by induction on $i$:
$\zeta_0=\sigma_0$ and $\zeta_i=\tau(\zeta_{i-1},r_i)$, where $r_i$ is the coordinate $z_{j_i}$ of
$\zeta_{i-1}$ for a short gate, and $r_i=y^{(s)}_{j_i}$ for a long gate, which equals
$x_{j_i}=v^{(s)}_{j_i}$ if $j_i\in K^{(s)}$, by the hypothesis and (a). So $z_i=w_i=x_i$ for
$i\in K^{(s+1)}$. When the set-valued iteration stops at $s$, $y^{(s)}=x$, and one more
application of $G_R$ confirms the fixed point.

(d) Gate $i$ defines a relation $\varrho_i$ on $\{0,1\}^R$, with $\sigma\,\varrho_i\,\sigma'$ if
$\sigma'=\tau(\sigma,b)$ for some $b\in B_i(\sigma)$; it is a Boolean matrix with one or two
ones in each row, and $S_i$ is the row $\sigma_0$ of the product of $\varrho_1,\dots,\varrho_i$.
A product of two Boolean matrices of this size has $8^R$ conjunctions and depth $O(R)$ with
bounded fan-in, and
a parallel prefix circuit uses $O(n)$ products in $O(\log n)$ layers. Whether a position is
determined is read off from its row.

(e) Mark the objects for the reach $R'$ with a prime, and for a window $\sigma'\in\{0,1\}^{R'}$
let $\pi(\sigma')$ be its last $R$ coordinates. We show by induction on $s$ that
$K^{(s)}\subseteq K'^{(s)}$, and for this, given the inclusion for $s$, by induction on $i$ that
$\pi(\sigma')\in S^{(s+1)}_i$ for every $\sigma'\in S'^{(s+1)}_i$. This holds for $i=0$. Let
$\sigma'\in S'^{(s+1)}_{i-1}$ and $b\in B'_i(\sigma')$; since $\pi$ commutes with $\tau$, it is
enough to see that $b\in B_i(\pi(\sigma'))$. If gate $i$ is short for $R$, both reads are the
coordinate $z_{j_i}$, which $\pi$ keeps. If it is long for $R'$, then $B'_i\subseteq B_i$ because
$K^{(s)}\subseteq K'^{(s)}$ and the values agree with $x$. If it is short for $R'$ and long for
$R$, then $b$ is the coordinate $z_{j_i}$ of $\sigma'$; this is allowed if $j_i\notin K^{(s)}$,
and if $j_i\in K^{(s)}\subseteq K'^{(s)}\subseteq K'^{(s+1)}$, then all windows of
$S'^{(s+1)}_{j_i}$ end with $x_{j_i}$ and $\sigma'$ still has this coordinate, so
$b=x_{j_i}=v^{(s)}_{j_i}$. Hence a position that is determined in round $s+1$ for $R$ is
determined for $R'$.

(f) If $[0,c]\subseteq K^{(s)}$, every gate $i<g$ has exactly one read for every window: a short
gate by definition, and a long gate because $j_i\le c$. So $S^{(s+1)}_i=\{\omega_i\}$ for $i<g$, by
induction on $i$ and (a). This is the statement of Theorem~\ref{thm:table}(c) with $K^{(s)}$ in
place of the set on which $y^{(s)}$ is right, and the two bounds follow as in the proofs of
Theorem~\ref{thm:table}(c) and (d): $[0,s(R+1)]\cap[0,n]\subseteq K^{(s)}$, and
$[0,c_s]\subseteq K^{(s)}$ for the chain $c_s$ of the proof of (d).
\end{proof}

\begin{remark}[the mixed base; numerical]\label{rem:setval}
On the mixed base of Remark~\ref{rem:mixed}, here with layers of $600$ positions over a zero
preamble of half the positions and three smoothings each (script \texttt{setval\_table.c}; the
depths $D$ and $D_2$ are from \texttt{design\_search.c}, on the same instances), at noise $0.50$:
\begin{center}
\begin{tabular}{lcccc}
\toprule
layers & $40$ & $80$ & $160$ & $320$\\
\midrule
$\max_iD(i)$ & $99$--$106$ & $186$--$205$ & $314$--$354$ & $597$--$629$\\
$\max_iD_2(i)$ & $48$--$52$ & $85$--$91$ & $149$--$170$ & $284$--$296$\\
rounds, sets of windows, $R=2$ & $19$--$25$ & $24$--$25$ & $27$--$31$ & $27$--$32$\\
least $s$ with $y^{(s)}=x$, table--Jacobi, $R=2$ & $18$--$24$ & $21$--$23$ & $24$--$27$ & $24$--$31$\\
\bottomrule
\end{tabular}
\end{center}
The last row is one less than the number of applications of $G_2$ in
Definition~\ref{def:table}, and by Proposition~\ref{prop:setval}(c) it is at most the row above.
At noise $0.46$ the last two rows are $24$, $31$, $29$ and $23$, $28$, $26$ for $40$, $160$,
$320$ layers (one smoothing each). Below $4/9$ the mixed base stays easy for reach $2$: with
$40$, $160$ and $640$ layers the sets of windows need $26$, $32$, $40$ rounds at noise $0.40$ and
$37$, $46$, $60$ at $0.30$; at $0.20$ and $0.10$ the growth is faster ($37$, $64$, $97$ and
$42$, $103$, $165$). So on this base derivations in which the pointers $i-2$ are free are still
long, by a count proportional to the number of layers (at $1280$ layers the referee of this
subsection found $\max_iD_2(i)$ to be $0.74$ per layer at noise $0.50$ and $1.0$ at $0.46$), and
yet the sets of windows collapse within a few dozen rounds ($29$ to $36$ at $1280$ layers). What
three-valued logic misses is that a gate with pointer $i-2$ soon reads two undetermined values
that are complements of each other, and then its output is $0$.

For longer pointers the picture at these sizes is less clear. With pointers $i-6$ in place of
$i-2$ and $40$ and $160$ layers, reach $5$ gives $83$ and $245$ rounds at noise $0.46$ and $79$
and $115$ at $0.50$ (the table--Jacobi iteration reaches the trace at $s=78$, $195$, $73$,
$110$), and reach $6$ gives $34$ and $61$, then $29$ and $37$ (and $s=30$, $45$, $23$, $30$).
Over six smoothings the referee found, for reach $6$ at noise $0.46$, $27$--$38$, $56$--$69$ and
$67$--$95$ rounds for $40$, $160$ and $320$ layers, and $41$--$43$, $67$--$69$, $83$--$91$ with
layers of $4000$ positions; at noise $0.50$ and $640$ layers, $48$--$60$ rounds against
$\max_iD_6(i)$ between $87$ and $106$. So with reach $6$ the rounds still grow with the number
of layers at noise $0.46$, much more slowly than with reach $5$, and $D_6$ is not proportional
to the number of layers.

For $R=1$ the sets $K^{(s)}$ were equal to $\{i:D(i)\le s\}$ on every parent array with $n\le9$
(the referee's check), so the iteration with reach $1$ is then the relevance-pruned iteration.
The only bounds we have for larger reach are those of Proposition~\ref{prop:setval}.
Heuristically, once few positions are undetermined, a position stays undetermined for one more
round only if its value is sensitive, in the dynamics with exact short pointers and with all
other long reads fixed, to the value of some undetermined position; so what has to be bounded is
the expected number of positions to whose value a given position is sensitive in this sense. For
random block bases over zero gates, where the long reads are independent and there are no short
pointers, a short computation gives for this number the quantity $\lambda=(1-\eps)/(1-q^2)$ of
Section~\ref{sec:above}, which is $1$ at noise $4/9$. The
question whether, for a reach of order $\log n$, it needs $O(\log n)$ rounds on the smoothing of
every base from noise $0.455$ on is the form in which we would now ask for an algorithm on all
bases: it concerns a monotone iteration, as in Section~\ref{sec:allbases}, but numerically the
count $D_R$ does not answer it.
\end{remark}

\section{Other layered designs}\label{sec:designs}

The latch of Section~\ref{sec:latch} is one way to put chains over a stage. This section considers
a family of such arrangements in the ideal model. For each of them it looks at two things: how
the densities of a layer react to the densities of the layer below, and how a flipped position of
the layer below spreads into the layer above. There are seven results.
\begin{itemize}[leftmargin=2em]
\item Theorem~\ref{thm:indep}: if every gate lists a position that is drawn independently of the
chain (a position of some class of the layer below, or position $0$ or $1$, at random), or is a
NOT gate, then for every noise level $\eps\ge0.455$ the density map is a contraction in a suitable
metric. It has one fixed point and the densities of the layers converge to it from every start;
when the design changes from layer to layer, the densities still forget the first layer, as two
starts approach each other geometrically; and the mean spread of a flipped position shrinks from
layer to layer. The proof is a potential argument with one inequality in two variables, which is
verified by interval arithmetic (Lemma~\ref{lem:potential}).
\item Theorem~\ref{thm:two}: the same holds when gates may also list the position two before
them, in any arrangement. A difference between two copies of a chain then has three states, one
of them hidden, and the same potential works, with one more inequality for every length of a
run of such gates (Lemma~\ref{lem:runs}, again by interval arithmetic).
\item Remark~\ref{rem:beyond} (numerical): the latch is not the best design. Designs of this kind
beat the plain chain above $4/9$ (with one rail from period $12$ on, and with three rails already
at period $6$): in some of them a flipped
position keeps spreading up to noise about $0.4487$, and in one of them, of period $24$, the
density map has an attracting cycle of length two up to noise about $0.447$. So pairs of layers
of that design have two attracting states, as a stage of the latch does below $4/9$.
\item Remark~\ref{rem:designs}: an exhaustive computation over about $120$ million designs of
period at most $7$, which also have local pointers and repeated reads. At noise $0.46$ none of
them spreads a flipped position faster than the plain chain. This search was too short to see the
designs of Remark~\ref{rem:beyond}.
\item Theorem~\ref{thm:layered}: the consequence of Theorem~\ref{thm:indep} for smoothed instances.
From noise $0.455$ on, the relevance-pruned iteration finishes in $O(\log n)$ rounds on the
smoothing of every layered base with wide bundles (at least $C_1\ln n$ chains, reading random copies)
and long chains over a long preamble; the circuit embeddings of Sections~\ref{sec:latch}
and~\ref{sec:s24} can be built in this form and still prove their hardness. So embeddings of this
kind stop between $0.4468$ and $0.455$.
\item Theorem~\ref{thm:twolayered}: the consequence of Theorem~\ref{thm:two}. From noise
$0.455$ on, the table iteration with reach $2$ finishes in $O(\log n)$ rounds on the smoothing
of every layered base with wide bundles whose programs may also list the position two back.
Two identities turn such a gate into one whose value depends on one earlier value of the chain
and one read, which gives a lower and an upper trace again. On these bases the relevance-pruned
iteration is slow in simulation.
\item Theorem~\ref{thm:scan}: a first bound for reads that depend on each other. In towers
that alternate plain layers with layers that read runs of consecutive positions of the layer
below, in any arrangement, a flipped position spreads to fewer than one position two layers up
from noise $0.446$ on, as it does for independent reads. An exact computation for a wider
family shows that dependence can nevertheless help, by about $0.002$
(Remark~\ref{rem:scan}).
\end{itemize}
The first four and the last are statements about the ideal model, and none of the seven is a statement about all bases.
The design of period $24$ of the second item is turned into an embedding in
Section~\ref{sec:s24}, which gives hardness up to noise $0.4468$. In the ideal model, the supremum of the noise
levels at which the density map of some design with independent reads has more than one
attracting set is at least about $0.447$, numerically, and at most $0.455$ by the theorem. It is
not $4/9$.

\begin{definition}[periodic designs]\label{def:design}
Let $0<\eps<1$ and let $R,T,K\ge1$ be integers. A \emph{class} is a pair $(r,\varphi)$ of a \emph{rail}
$r\in\{1,\dots,R\}$ and a \emph{phase} $\varphi\in\{0,\dots,T-1\}$. A \emph{design} gives every class a
\emph{read type}: $\mathrm E(c)$ for a class $c$, or $\mathrm U$, $\mathrm S$, $\mathrm Z_0$,
$\mathrm Z_1$, or $\mathrm L_k$ with $1\le k\le K$. Some class has type $\mathrm E(c)$ or
$\mathrm U$, and every rail with a class of type $\mathrm S$ has a class of type $\mathrm E(c)$ or
$\mathrm U$.

Let $d=(d_c)$ be a vector of numbers in $[0,1]$, one for every class, and $\bar d$ their mean. A
\emph{chain of rail $r$ over $d$} is a sequence of bits $(x_g)_{g\in\mathbb Z}$ with
$x_g=\NOR(x_{g-1},y_g)$, in which gate $g$ has the read type of the class $(r,g\bmod T)$ and its
read $y_g$ is as follows. With probability $\eps$ the read is \emph{resampled}, and then $y_g=1$
with probability $\varpi$. Otherwise $y_g$ is the \emph{target} of the gate: for type
$\mathrm E(c)$ a new bit that is $1$ with probability $d_c$; for type $\mathrm U$ a new bit that
is $1$ with probability $\bar d$; the bit $1$ for $\mathrm Z_0$ and $0$ for $\mathrm Z_1$; the bit
$x_{g-k}$ for $\mathrm L_k$; and for type $\mathrm S$ the target of the last gate before $g$ of
type $\mathrm E(c)$ or $\mathrm U$. All these choices are independent, and the law of the
sequence is invariant under shifts by $T$. This law is unique: with probability $(\eps\varpi)^{K+T}$ the gates $g-K-T+1,\dots,g$ are all resampled and
read a one, and then the values after $g$ do not depend on what happened before gate $g-K-T+1$.
The \emph{density map} $\Phi$ sends $d$ to the vector $d'$ with $d'_{(r,\varphi)}=\PP[x_g=1]$ for a
gate $g$ of phase $\varphi$ in a chain of rail $r$ over $d$.
\end{definition}

\noindent This is the ideal model of Section~\ref{sec:latch} for a layer that consists of long chains of
every rail over a layer of the same kind. The layer below has density $d_c$ of ones in class $c$.
A gate of type $\mathrm E(c)$ lists a uniformly random position of class $c$ of the layer below,
a gate of type $\mathrm U$ a uniformly random position of that layer, a gate of type
$\mathrm Z_0$ or $\mathrm Z_1$ position $0$ or $1$, a gate of type $\mathrm L_k$ the position $k$
before it (for $k=1$ it is a NOT gate), and a gate of type $\mathrm S$ the position listed by the
last gate of type $\mathrm E$ or $\mathrm U$ before it. The density of ones outside the layers is
$\varpi$, and the positions listed by different gates of type $\mathrm E$ or $\mathrm U$ are
independent. A chain of plain gates over a stage, as in a block base, is the design with $R=T=1$
and type $\mathrm U$. The \emph{latch with a copy after every pair} is the design with $R=T=2$
whose rail $1$ has the types $\mathrm E(1,1)$, $\mathrm E(2,1)$ at the phases $0$, $1$, and whose
rail $2$ has the types $\mathrm E(2,1)$, $\mathrm E(1,1)$: by Lemma~\ref{lem:latch} the gate of
phase $1$ is set by its own rail and reset by the other one.

\begin{definition}[two matrices]\label{def:multipliers}
Let $d^*$ be a fixed point of $\Phi$ and let $g$ be a gate of type $\mathrm E(c)$ or $\mathrm U$ in
a chain over $d^*$. Run two copies $x$ and $x'$ of that chain with the same resampling choices,
the same resampled values and the same targets, except that the target of $g$ is $0$ in $x$ and
$1$ in $x'$. For a class $k=(r,\varphi)$ of the rail of that chain let $N^+_g(k)$ and $N^-_g(k)$ be the
expected numbers of gates $h\ge g$ of phase $\varphi$ with $(x_h,x'_h)=(0,1)$ and with
$(x_h,x'_h)=(1,0)$. For classes $k$ and $c$ put
\[
\overline M_{k,c}=\sum_{g}w_{g,c}\bigl(N^+_g(k)+N^-_g(k)\bigr),\qquad
J_{k,c}=\sum_{g}w_{g,c}\bigl(N^+_g(k)-N^-_g(k)\bigr),
\]
where the sums are over the $T$ gates $g$ of one period of the rail of $k$, and $w_{g,c}=1$ if $g$
has type $\mathrm E(c)$, $w_{g,c}=1/(RT)$ if $g$ has type $\mathrm U$, and $w_{g,c}=0$ otherwise.
Let $\rho_u$ and $\rho_s$ be the spectral radii of $\overline M$ and of $J$.
\end{definition}

\noindent The numbers $N^\pm_g(k)$ are finite, by the argument for uniqueness. If the roles of
$x$ and $x'$ are exchanged, $N^+_g+N^-_g$ does not change; so $\overline M_{k,c}$ is the mean
number of positions of class $k$ that change when one position of class $c$ of the layer below is
flipped and both layers are evaluated with the same noise, whatever that position holds. Here
every rail has the same number of periods in a layer, so that a position of class $c$ is listed by
one gate on average for every class of type $\mathrm E(c)$ and by $1/(RT)$ gates for every class
of type $\mathrm U$; the probability $1-\eps$ that the read is kept is part of $N^\pm_g$. If the
design has no gate of type $\mathrm U$, other proportions of the rails conjugate $\overline M$ by a
diagonal matrix and leave $J$, $\rho_u$ and $\rho_s$ unchanged; with gates of type $\mathrm U$
they change the density map itself, and the numbers below are for equal proportions. When few positions of a layer are flipped, the changes spread from layer
to layer like a branching process with mean matrix $\overline M$: they die out if $\rho_u<1$, and
survive with positive probability if $\rho_u>1$. We use this only as the meaning of the number
$\rho_u$ and prove nothing about it.

\begin{proposition}[three facts]\label{prop:designs}
Let $\lambda=(1-\eps)/(1-q^2)$ with $q=q(\eps)$, and $\pi_0=q/(1+q)$, as in
Proposition~\ref{prop:ideal}.
\begin{enumerate}[label=(\alph*),leftmargin=2.4em]
\item At a fixed point $d^*$ of the density map of a design, $J$ is the Jacobian of $\Phi$, and
$|J_{k,c}|\le\overline M_{k,c}$ for all classes $k$ and $c$. Hence $\rho_s\le\rho_u$, and if
$\rho_u<1$ then $d^*$ is an attracting fixed point of $\Phi$.
\item If every class has type $\mathrm E(c)$ or $\mathrm U$, the vector with all entries $\pi_0$
is a fixed point of $\Phi$, and at this fixed point $\rho_u=\lambda$.
\item For the latch with a copy after every pair, $\rho_s=\rho_u=\lambda$ at the fixed point with
all entries $\pi_0$.
\end{enumerate}
\end{proposition}

\begin{proof}
(a) Let $h$ be a gate of class $k$. Given all other choices, $x_h$ is a function of the targets of
the gates $g\le h$ of type $\mathrm E$ or $\mathrm U$, which are independent bits; the target of a
gate of type $\mathrm E(c)$ is $1$ with probability $d_c$, and that of a gate of type $\mathrm U$
with probability $\bar d$, whose derivative in $d_c$ is $1/(RT)$. For the chain started $n$ gates
before $h$ from a fixed state, $\PP[x_h=1]$ is a polynomial in $d$, and its derivative in $d_c$
is $\sum_gw_{g,c}\bigl(\EE[x_h\mid t_g=1]-\EE[x_h\mid t_g=0]\bigr)$, over the gates $g$ among
these $n$, with $t_g$ the target of $g$. The term of $g$ is at most
$2(1-(\eps\varpi)^{K+T})^{\lfloor(h-g)/(K+T)\rfloor}$ in absolute value, uniformly in $d$ and in
the starting state. By the same regeneration argument the truncated probabilities converge to
$d'_k=\PP[x_h=1]$, uniformly in $d$. So the derivatives converge uniformly as $n\to\infty$, $\Phi$
is continuously differentiable, and the limit of the derivatives is the derivative of $d'_k$. With the two copies of Definition~\ref{def:multipliers},
$\EE[x_h\mid t_g=1]-\EE[x_h\mid t_g=0]=\PP[(x_h,x'_h)=(0,1)]-\PP[(x_h,x'_h)=(1,0)]$. Summing over
$g\le h$ with $h$ fixed is the same, by stationarity, as summing over the gates $g$ of one period
and all $h\ge g$ of the phase of $k$. This gives $\partial d'_k/\partial d_c=J_{k,c}$ at $d^*$. The
inequality is clear from the definitions. For matrices with $|J|\le\overline M$ entrywise the
spectral radius of $J$ is at most that of $\overline M$, and a fixed point at which the Jacobian
has spectral radius less than $1$ is attracting.

(b) Let $d$ have all entries $\pi_0$. Every gate reads a one with probability
$A+(1-\eps)\pi_0=1-q$ (Proposition~\ref{prop:ideal}(b)), independently of the earlier values of
the chain. So $p_g=\PP[x_g=1]$ satisfies $p_g=(1-p_{g-1})q$, and the stationary solution is
$p_g=q/(1+q)=\pi_0$ at every phase: $d$ is a fixed point. In Definition~\ref{def:multipliers},
the copies differ at $g$ exactly if the read of $g$ is kept and $x_{g-1}=0$, which has probability
$(1-\eps)/(1+q)$. If they differ at a gate $h\ge g$, they differ at $h+1$ exactly if gate $h+1$
reads a zero, because that read is the same in both copies; this has probability $q$ and is
independent of what happened before. So the copies differ at $g+j$ with probability
$(1-\eps)q^j/(1+q)$. For a class $k=(r,\varphi)$ every $j\ge0$ has exactly one gate $g$ of the period of
rail $r$ with $g+j\equiv \varphi\pmod T$, and $\sum_cw_{g,c}=1$ for every gate. Hence
\[
\sum_c\overline M_{k,c}=\sum_{j\ge0}\frac{(1-\eps)q^j}{1+q}=\frac{1-\eps}{(1+q)(1-q)}=\lambda
\]
for every $k$. A nonnegative matrix all of whose row sums are $\lambda$ has spectral radius
$\lambda$.

(c) By (b), $\rho_u=\lambda$. In the computation of (b), $x_g=1$ and $x'_g=0$ if the copies
differ at $g$, and the two values alternate along the chain while the copies differ. So the pair
is $(1,0)$ at even distances $j$ from $g$ and $(0,1)$ at odd ones. Only the classes $(1,1)$ and
$(2,1)$ are listed. On rail $1$ the gate of phase $0$ lists class $(1,1)$ and is at odd distances
from the gates of phase $1$, so
$J_{(1,1),(1,1)}=\sum_{j\text{ odd}}(1-\eps)q^j/(1+q)=a$ with $a=(1-\eps)q/((1+q)(1-q^2))$; the
gate of phase $1$ lists class $(2,1)$ and is at even distances, so $J_{(1,1),(2,1)}=-b$ with
$b=(1-\eps)/((1+q)(1-q^2))$. By symmetry $J_{(2,1),(2,1)}=a$ and $J_{(2,1),(1,1)}=-b$, and the
columns of the classes $(1,0)$ and $(2,0)$ are zero. The nonzero eigenvalues of $J$ are therefore
$a+b=\lambda$, for the difference of the two rails, and $a-b=-(1-\eps)/(1+q)^2$. This is the
multiplier of~\eqref{eq:uv} at the diagonal fixed point.
\end{proof}

\subsection{Independent reads: a contraction for noise at least \texorpdfstring{$0.455$}{0.455}}\label{sec:indep}

\begin{definition}[independent reads]\label{def:indep}
A \emph{mixed type} is a probability vector $\mu=(\mu_0,\mu_1,(\mu_c)_c)$, with one entry for
position $0$, one for position $1$ and one for every class. A gate of mixed type $\mu$ lists
position $0$ with probability $\mu_0$, position $1$ with probability $\mu_1$, and a uniformly
random position of class $c$ of the layer below with probability $\mu_c$. In a chain over $d$ its
target is a new bit that is $1$ with probability $\mu_0+\sum_c\mu_cd_c$, and its read is resampled
with probability $\eps$ as before. The types $\mathrm E(c)$, $\mathrm U$, $\mathrm Z_0$ and
$\mathrm Z_1$ are mixed types; so is a gate that is thinned to position $1$ with some
probability. A design \emph{has independent reads} if every class has a mixed type or the type
$\mathrm L_1$. For such a design and $d\in[0,1]^{\mathcal C}$, where $\mathcal C$ is the set of
classes, and for a gate $g$ of a chain over $d$, put $y_g=\mu_1+\sum_c\mu_c(1-d_c)$ if $g$ has
mixed type $\mu$, put $y_g=1$ if $g$ has type $\mathrm L_1$, and in both cases
\[
q_g=B+(1-\eps)y_g .
\]
The mixed types may also refer to the classes $\mathcal C'$ of another design, that of the layer
below; then the density map goes from $[0,1]^{\mathcal C'}$ to $[0,1]^{\mathcal C}$. For designs with
independent reads, $\rho_s$ denotes the spectral radius of the Jacobian of $\Phi$ at a fixed point,
and $\rho_u$ that of the matrix $\overline M$ of Lemma~\ref{lem:onebit}(d).
\end{definition}

\noindent So $y_g\in[0,1]$ and $B\le q_g\le B+(1-\eps)=1-A$. Throughout this subsection
\[
hi=\frac{1-B}{1-B(1-A)},\qquad lo=1-(1-A)\,hi,
\]
so that $0<lo\le hi<1$, and $\mathcal K$ is the set of all $d\in[0,1]^{\mathcal C}$ with $1-d_c\in[lo,hi]$
for every class $c$.

\begin{lemma}[chains with independent reads]\label{lem:onebit}
Consider a design with independent reads, a vector $d\in[0,1]^{\mathcal C}$ and a chain over $d$.
\begin{enumerate}[label=(\alph*),leftmargin=2.4em]
\item For every gate $g$ there is an event $Z_g$ of probability $q_g$, independent of the choices
made at the gates before $g$, with $x_g=1$ exactly if $Z_g$ occurs and $x_{g-1}=0$.
\item The numbers $\alpha_g=\PP[x_g=0]$ depend only on the phase of $g$, satisfy
$\alpha_g=1-q_g\alpha_{g-1}$, and lie in $[lo,hi]$. So $\Phi(d)\in\mathcal K$, and $\Phi$ is continuously
differentiable.
\item Let $g$ have a mixed type and run two copies of the chain with the same choices, except
that the target of $g$ is $0$ in one copy and $1$ in the other. For $h\ge g$ the copies differ at
$h$ with probability $(1-\eps)\alpha_{g-1}\prod_{g<i\le h}q_i$.
\item For classes $k$ and $c$ let
$\overline M_{k,c}(d)=\sum_{g\le h}(1-\eps)\mu_c(g)\,\alpha_{g-1}\prod_{g<i\le h}q_i$, where $h$ is
a gate of class $k$, the sum is over the gates $g\le h$ of mixed type in its chain, and $\mu(g)$ is
the type of $g$. Then $|\partial\Phi(d)_k/\partial d_c|\le\overline M_{k,c}(d)$.
\end{enumerate}
\end{lemma}

\begin{proof}
(a) For a mixed type let $Z_g$ be the event that the read of $g$ is $0$. The read is resampled
and $0$ with probability $\eps(1-\varpi)=B$, and kept and $0$ with probability $(1-\eps)y_g$,
since the target is $0$ with probability $y_g$. For type $\mathrm L_1$ let $Z_g$ be the event that
the read is not both resampled and $1$; it has probability $1-\eps\varpi=1-A$. If the read of an
$\mathrm L_1$ gate is kept, $x_g=\NOR(x_{g-1},x_{g-1})=1-x_{g-1}$; if it is resampled, $x_g=0$
when it is $1$ and $x_g=1-x_{g-1}$ when it is $0$. In both cases $x_g=1$ exactly if $Z_g$ occurs
and $x_{g-1}=0$, and $Z_g$ depends only on the choices made at gate $g$.

(b) By (a), $\PP[x_g=1]=q_g(1-\PP[x_{g-1}=1])$, that is $\alpha_g=1-q_g\alpha_{g-1}$. Over one
period this is an affine map with slope of absolute value $\prod q_i<1$, so there is exactly one
periodic solution, it is the stationary one, and it is a rational function of $d$ without poles
on $[0,1]^{\mathcal C}$. Let $m$ and $M$ be the least and the largest of the $\alpha_g$. From
$B\le q_g\le1-A$ we get $M\le1-Bm$ and $m\ge1-(1-A)M$, hence $M\le1-B+B(1-A)M$, so $M\le hi$, and
then $m\ge lo$.

(c) The copies agree before $g$. They differ at $g$ exactly if the read of $g$ is kept and
$x_{g-1}=0$, which has probability $(1-\eps)\alpha_{g-1}$. For $i>g$ the event $Z_i$ is the same
in both copies. By (a), if the copies differ at $i-1$ they differ at $i$ exactly if $Z_i$ occurs,
and if they agree at $i-1$ they agree at $i$.

(d) Let $h$ be a gate of class $k$, so that $\Phi(d)_k=1-\alpha_h$. Differentiating
$\alpha_h=1-q_h\alpha_{h-1}$, with $\partial q_h/\partial d_c=-(1-\eps)\mu_c(h)$ for a mixed type
and $0$ for $\mathrm L_1$, gives
$\partial\alpha_h/\partial d_c=-q_h\,\partial\alpha_{h-1}/\partial d_c+(1-\eps)\mu_c(h)\alpha_{h-1}$.
The numbers $\overline M_h=\overline M_{k,c}(d)$ satisfy
$\overline M_h=q_h\overline M_{h-1}+(1-\eps)\mu_c(h)\alpha_{h-1}$. So
$E_h=|\partial\alpha_h/\partial d_c|-\overline M_h$ satisfies $E_h\le q_hE_{h-1}$; as $E_h$ is
periodic in $h$, this gives $E_h\le(\prod q_i)E_h$ and $E_h\le0$.
\end{proof}

\noindent By (c), for the types of Definition~\ref{def:design} the matrix $\overline M(d^*)$ of
(d) is the matrix $\overline M$ of Definition~\ref{def:multipliers}.

\begin{lemma}[the inequality]\label{lem:potential}
Let $f(x)=x^{43/20}(1-x)^{33/20}$ on $[0,1]$, and let $\hat f$ be the function on $[0,1]$ with
$\hat f(y)=S_1y$ for $y\le\frac{23}{56}$, $\hat f(y)=f(y)$ for
$\frac{23}{56}\le y\le\frac{43}{56}$ and $\hat f(y)=S_2(1-y)$ for $y\ge\frac{43}{56}$, where
$S_1=f(\frac{23}{56})/\frac{23}{56}$ and $S_2=f(\frac{43}{56})/\frac{13}{56}$.
\begin{enumerate}[label=(\alph*),leftmargin=2.4em]
\item $\hat f$ is concave, $\hat f\ge f$, and $\hat f(0)=\hat f(1)=0$.
\item For every $\eps\in[0.455,1)$, every $x\in[lo,hi]$ and every $y\in[0,1]$,
\[
\bigl(B+(1-\eps)y\bigr)f(x)+(1-\eps)\,x\,\hat f(y)\ \le\ \theta\,f\bigl(1-x(B+(1-\eps)y)\bigr),
\qquad\theta=0.9998 .
\]
(Here $\theta$ is a contraction factor, not the thinning of Section~\ref{sec:latch}.)
\end{enumerate}
\end{lemma}

\begin{proof}
(a) Write $a=43/20$ and $b=33/20$. The function $f(t)/t=t^{a-1}(1-t)^b$ on $(0,1)$ has its maximum
at $t=(a-1)/(a+b-1)=23/56$, so $S_1y\ge f(y)$ for all $y$, and the line $S_1y$ touches $f$ at
$23/56$. In the same way $f(t)/(1-t)=t^a(1-t)^{b-1}$ has its maximum at $t=a/(a+b-1)=43/56$, so
$S_2(1-y)\ge f(y)$, and this line touches $f$ at $43/56$. Finally
$t^2(1-t)^2f''(t)/f(t)=(a-(a+b)t)^2-a(1-t)^2-bt^2$ is a quadratic in $t$ with positive leading
coefficient and with the negative values $-759/1120$ and $-559/1120$ at $23/56$ and at $43/56$, so
$f''<0$ between these points. Hence $\hat f$ is continuously differentiable with a nonincreasing
derivative.

(b) This is verified by the program \texttt{potential\_check.c}, with interval arithmetic. For an
interval $[e_0,e_1]$ of noise levels it computes enclosures of $\varpi$, $A$, $B$, $1-\eps$, $lo$
and $hi$; here it uses that $\varpi(\eps)$ is increasing, which follows from
$\varpi=\eps(1-\varpi)^2$, and it accepts a bound $v$ for $\varpi(e_0)$ or $\varpi(e_1)$ only
after checking the sign of $e(1-v)^2-v$. On a box $[x_0,x_1]\times[y_0,y_1]$ it bounds the left
side from above and the right side from below, using that $f$ and $\hat f$ increase up to $43/76$
and decrease after it. A box on which this fails is halved, and so is an interval of noise levels
on which a box of the smallest size fails. All arithmetic is in IEEE double precision with every
result moved one unit in the last place in the safe direction; a number $r$ is used as an upper
(lower) bound for a twentieth root of $g$ only after $r^{20}\ge g$ ($r^{20}\le g$) has been
checked in this arithmetic, so apart from \texttt{nextafter} and the conversion of the decimal
arguments no library function is trusted. Because those arguments are rounded, the runs start and
end a little outside the intervals they cover and use a factor a little below $\theta$: one run
covers $[0.45499,0.46001]$ with the factor $0.99979$ in $51$ intervals and $778{,}769$ boxes, and
three runs cover $[0.45999,1]$ with the factor $0.9995$ in $108$ intervals and $351{,}696$ boxes
(Section~\ref{sec:checks}).
\end{proof}

\begin{theorem}[independent reads]\label{thm:indep}
Let $0.455\le\eps<1$ and $\theta=0.9998$, and let $f$ be the function of
Lemma~\ref{lem:potential}. For $d,d'\in\mathcal K$ let $\delta(d,d')$ be the infimum of
$\int_0^1\max_c|\gamma_c'(t)|/f(1-\gamma_c(t))\,dt$ over the piecewise continuously
differentiable paths $\gamma$ in $\mathcal K$ from $d$ to $d'$. Consider a design with independent reads
and its density map $\Phi$.
\begin{enumerate}[label=(\alph*),leftmargin=2.4em]
\item $\Phi$ maps $[0,1]^{\mathcal C}$ into $\mathcal K$, and
$\delta(\Phi(d),\Phi(d'))\le\theta\,\delta(d,d')$ for all $d,d'\in\mathcal K$.
\item $\Phi$ has exactly one fixed point $d^*$, and
$\|\Phi^{t}(d)-d^*\|_\infty\le C\theta^{t}$ for every $d\in[0,1]^{\mathcal C}$ and $t\ge1$, with
a constant $C$ that depends only on $\eps$.
\item At $d^*$ the matrix $\overline M(d^*)$ of Lemma~\ref{lem:onebit}(d) and the Jacobian of $\Phi$
have spectral radius at most $\theta$: $\rho_s\le\rho_u\le\theta$.
\item Let $\Phi_1,\Phi_2,\dots$ be the density maps of designs with independent reads, where the
classes of the design of $\Phi_{t}$ are the classes listed by the design of $\Phi_{t+1}$. Then
$\|\Phi_t\circ\dots\circ\Phi_1(d)-\Phi_t\circ\dots\circ\Phi_1(d')\|_\infty\le C\theta^{t}$ for all
$d$, $d'$ and $t\ge1$.
\end{enumerate}
\end{theorem}

\begin{proof}
Let $d\in[0,1]^{\mathcal C}$, write $\beta_c=1-d_c$, and let $h$ be a gate of a chain over $d$,
with $\alpha$ and $q$ as in Lemma~\ref{lem:onebit}. We claim
\begin{equation}\label{eq:gatepot}
q_hf(\alpha_{h-1})+(1-\eps)\alpha_{h-1}\sum_c\mu_c(h)f(\beta_c)\ \le\ \theta f(\alpha_h),
\end{equation}
where the sum is absent if $h$ has type $\mathrm L_1$. Indeed $\alpha_{h-1}\in[lo,hi]$ and
$\alpha_h=1-\alpha_{h-1}(B+(1-\eps)y_h)$, so by Lemma~\ref{lem:potential}(b) the right side is at
least $q_hf(\alpha_{h-1})+(1-\eps)\alpha_{h-1}\hat f(y_h)$. For a mixed type,
$\sum_c\mu_cf(\beta_c)\le\mu_0\hat f(0)+\mu_1\hat f(1)+\sum_c\mu_c\hat f(\beta_c)\le\hat f(y_h)$ by
Lemma~\ref{lem:potential}(a) and Jensen's inequality, because
$y_h=\mu_0\cdot0+\mu_1\cdot1+\sum_c\mu_c\beta_c$. For type $\mathrm L_1$, $y_h=1$ and
$\hat f(1)=0$.

Now let $N_h=\sum_c\overline M_{k,c}(d)f(\beta_c)$ for a gate $h$ of class $k$. By the definition
of $\overline M$, $N_h=q_hN_{h-1}+(1-\eps)\alpha_{h-1}\sum_c\mu_c(h)f(\beta_c)$. With
\eqref{eq:gatepot} and $\theta\le1$, the numbers $D_h=N_h-\theta f(\alpha_h)$ satisfy
$D_h\le q_h(N_{h-1}-f(\alpha_{h-1}))\le q_hD_{h-1}$, and they are periodic, so $D_h\le0$:
\begin{equation}\label{eq:Mpot}
\sum_c\overline M_{k,c}(d)\,f(1-d_c)\ \le\ \theta\,f\bigl(1-\Phi(d)_k\bigr)\qquad\text{for every class }k .
\end{equation}

(a) $\Phi$ maps into $\mathcal K$ by Lemma~\ref{lem:onebit}(b). For $d\in\mathcal K$ and a vector $\xi$ put
$\|\xi\|_d=\max_c|\xi_c|/f(1-d_c)$. By Lemma~\ref{lem:onebit}(d) and~\eqref{eq:Mpot}, the
derivative $D\Phi(d)$ satisfies
\[
\frac{|(D\Phi(d)\xi)_k|}{f(1-\Phi(d)_k)}\le\frac{\sum_c\overline M_{k,c}(d)f(1-d_c)}{f(1-\Phi(d)_k)}\,\|\xi\|_d\le\theta\|\xi\|_d ,
\]
so $\|D\Phi(d)\xi\|_{\Phi(d)}\le\theta\|\xi\|_d$. If $\gamma$ is a path in $\mathcal K$ from $d$ to $d'$,
then $\Phi\circ\gamma$ is a path in $\mathcal K$ from $\Phi(d)$ to $\Phi(d')$ whose length, in the sense of
the definition of $\delta$, is at most $\theta$ times that of $\gamma$.

(b) Let $m_f$ and $M_f$ be the minimum and the maximum of $f$ on $[lo,hi]$; both are positive.
Every path in $\mathcal K$ from $d$ to $d'$ has length at least $\|d-d'\|_\infty/M_f$, and the straight
path has length at most $\|d-d'\|_\infty/m_f$. So $\delta$ is a metric on $\mathcal K$ that is equivalent
to the maximum norm, $(\mathcal K,\delta)$ is complete, and $\delta\le(hi-lo)/m_f$. By (a) and the
contraction principle $\Phi$ has exactly one fixed point $d^*$ in $\mathcal K$, and
$\delta(\Phi^t(d),d^*)\le\theta^{t-1}\delta(\Phi(d),d^*)$ for every $d\in[0,1]^{\mathcal C}$; every
fixed point lies in $\mathcal K$. This gives (b) with $C=M_f(hi-lo)/(\theta m_f)$.

(c) At $d=d^*$, \eqref{eq:Mpot} says $\overline Mv\le\theta v$ for the positive vector
$v_c=f(1-d^*_c)$. So the spectral radius of $\overline M$ is at most $\theta$, and that of the
Jacobian, which is bounded entrywise by $\overline M$ in absolute value
(Lemma~\ref{lem:onebit}(d)), is at most that.

(d) The claim \eqref{eq:gatepot} and its consequences did not use that the layer below has the
same design. So (a) holds for every $\Phi_t$, with the metric $\delta$ of its own set of classes
on each side, and
$\delta(\Phi_t\circ\dots\circ\Phi_1(d),\Phi_t\circ\dots\circ\Phi_1(d'))\le\theta^{t-1}\delta(\Phi_1(d),\Phi_1(d'))\le\theta^{t-1}(hi-lo)/m_f$.
\end{proof}

\begin{remark}[what the theorem covers]\label{rem:indep}
\begin{enumerate}[label=(\roman*),leftmargin=2.4em]
\item \emph{The latch.} The ideal map $T$ of Section~\ref{sec:latch} is the density map of the
latch with a copy after every pair, restricted to the two classes that are listed, and a stage
with thinning is a design with a mixed type. So for $\eps\ge0.455$ the stage densities of every
such construction, with any thinning and any assignment of rails to the two checks, forget the
first stage: for a repeated design they converge to the fixed point, and for stages that differ
from one another two starts approach each other like $\theta^t$. The same holds when a stage reads
several earlier stages, as a gate stage of $\mathrm{Latch}(C)$ reads two wires: row $k$
of~\eqref{eq:Mpot} involves only the classes that the chain of $k$ lists, so the argument of (d)
applies to every arrangement of stages without cycles, with $t$ the length of the shortest path
of stages from the start. For $T$ itself Proposition~\ref{prop:ideal}(c) gives the uniqueness of
the fixed point already for $\eps\ge4/9$; the theorem adds the convergence, the other designs, and
stages that differ from one another, from $0.455$ on.
\item \emph{Between $4/9$ and $0.455$.} Part (c) fails below about $0.4487$ and part (b) below
about $0.447$: by Remark~\ref{rem:beyond} there are designs with independent reads with $\rho_u>1$
up to the first level, and one whose density map has an attracting cycle of length two up to the
second. The function $f$ was found by a search over $x^a(1-x)^b$; the inequality of
Lemma~\ref{lem:potential} holds with a factor below $1$ down to about $0.4538$ for the exponents
chosen, and the best exponents gain very little (about $0.4537$). A value iteration for the best potential that depends on $\alpha$ alone (script
\texttt{potential\_bellman.py}, not rigorous) gives a growth rate that crosses $1$ at about
$0.449$. So for single changes the level of this family appears to be about $0.449$, and a better
potential of the same kind could probably move the theorem close to it. For more than one
attracting set the supremum of the noise levels lies between about $0.447$ (numerically, one
design) and $0.455$.
\item \emph{Chains that are not stationary.} The proof uses periodicity in two places only: to
get $D_h\le0$ from $D_h\le q_hD_{h-1}$, and $E_h\le0$ from $E_h\le q_hE_{h-1}$ in
Lemma~\ref{lem:onebit}(d). Consider instead a chain that begins at a gate $h_0$ whose predecessor
holds $0$ with a probability $\alpha_{h_0-1}\in[lo,hi]$ that does not depend on $d$. Define
$\overline M$ and $N$ by sums over the gates of the chain from $h_0$ on. Then
$E_{h_0-1}=0$ and $D_{h_0-1}=-\theta f(\alpha_{h_0-1})\le0$, and every $\alpha_h$ lies in
$[lo,hi]$, because $\alpha\mapsto1-q\alpha$ maps this interval into itself for $q\in[B,1-A]$. So
Lemma~\ref{lem:onebit}(d) and~\eqref{eq:Mpot} hold by induction along the chain, with a class for
every position of a chain, and parts (a) and (d) of the theorem hold for layers of finite chains
of this kind. In $\mathrm{Latch}(C)$ every chain follows $n=2h\rho$ zero gates, after which the
predecessor holds $0$ with probability $1/(1+B)+(-B)^n(\alpha'-1/(1+B))$, where $\alpha'$ is the
probability before the stretch, which depends on the densities. So the hypothesis holds there only
up to an error of order $B^n$; we have not extended the argument to such errors. Chains that start
from a fixed bit are not covered, since $f$ vanishes at $0$ and $1$.
\item \emph{What it does not cover.} Local pointers of reach $2$ or more and repeated reads of one
position are not independent reads. Pointers two back are added in Theorem~\ref{thm:two}; for
the others see Remark~\ref{rem:designs}. A gate of type
$\mathrm U$ in a layer whose rails have unequal proportions has a mixed type with unequal
weights, so such layers are covered. The theorem concerns
the ideal model: it says that in one kind of construction the class densities carry no
information from layer to layer at noise $0.455$ or above. It does not give an algorithm, and it says nothing about bases that are
not of this kind.
\end{enumerate}
\end{remark}

\subsection{Pointers two back}\label{sec:two}

Theorem~\ref{thm:indep} allows one local pointer, the NOT gate. This subsection adds the type
$\mathrm L_2$, a gate that lists the position two before it, in any arrangement. A flipped
position can then stay in a chain without being visible: the two copies agree at a position and
differ at the one before it, and the next gate of type $\mathrm L_2$ makes them differ again. So
the state of a difference has three values, and the inequality of Lemma~\ref{lem:potential} is
replaced, at such gates, by one inequality for every length of a run of them. The potential $f$
stays the same.

\begin{definition}[pointers two back]\label{def:two}
A design \emph{has independent reads and pointers two back} if every class has a mixed type, the
type $\mathrm L_1$ or the type $\mathrm L_2$. For a chain over $d$ of such a design let
$\alpha_g=\PP[x_g=0]$, and let $y_g$ and $q_g$ be as in Definition~\ref{def:indep} for the gates
of mixed type and of type $\mathrm L_1$. For a gate $g$ of mixed type and a gate $h\ge g$ of the
same chain let $p_g(h)$ be the probability that two copies of the chain differ at $h$ when they
are run with the same choices, except that the target of $g$ is $0$ in one copy and $1$ in the
other. For classes $k$ and $c$ put
\[
\overline M_{k,c}(d)=\sum_{g\le h}\mu_c(g)\,p_g(h),
\]
where $h$ is a gate of class $k$, the sum is over the gates $g\le h$ of mixed type in its chain,
and $\mu(g)$ is the type of $g$. We write $r=1-A$.
\end{definition}

\noindent For a design with independent reads this is the matrix of Lemma~\ref{lem:onebit}(d),
by part (c) of that lemma. At a fixed point, and for the types of Definition~\ref{def:design},
it is the matrix $\overline M$ of Definition~\ref{def:multipliers}.

\begin{lemma}[chains with pointers two back]\label{lem:twochain}
Consider a design with independent reads and pointers two back, a vector
$d\in[0,1]^{\mathcal C}$ and a chain over $d$.
\begin{enumerate}[label=(\alph*),leftmargin=2.4em]
\item Every $\alpha_g$ lies in $[lo,hi]$. If $g$ has a mixed type or the type $\mathrm L_1$,
then $\alpha_g=1-q_g\alpha_{g-1}$. If $g$ has type $\mathrm L_2$, then
$\alpha_g=(2-\eps)-r\,\alpha_{g-1}-(1-\eps)\,\alpha_{g-2}$.
\item $\Phi$ is continuously differentiable, and
$|\partial\Phi(d)_k/\partial d_c|\le\overline M_{k,c}(d)$ for all classes $k$ and $c$.
\item Let $g$ have a mixed type and take the two copies of Definition~\ref{def:two}. For
$h\ge g$ let $a_h$ be the probability that the copies differ at $h$ and agree at $h-1$, $b_h$
that they differ at both, and $c_h$ that they agree at $h$ and differ at $h-1$. Then
$p_g(h)=a_h+b_h$, and $a_g=(1-\eps)\alpha_{g-1}$, $b_g=c_g=0$. For $h>g$,
\[
(a_h,b_h,c_h)=\bigl(0,\ q_h(a_{h-1}+b_{h-1}),\ (1-q_h)(a_{h-1}+b_{h-1})\bigr)
\]
if $h$ has a mixed type or the type $\mathrm L_1$, and
\[
(a_h,b_h,c_h)=\bigl((1-\eps)c_{h-1},\ r\,a_{h-1}+B\,b_{h-1},\ A\,a_{h-1}+(1-B)\,b_{h-1}\bigr)
\]
if $h$ has type $\mathrm L_2$.
\end{enumerate}
\end{lemma}

\begin{proof}
A chain has no two consecutive ones, since $x_g=1$ needs $x_{g-1}=0$.

(a) For a gate of mixed type or of type $\mathrm L_1$, Lemma~\ref{lem:onebit}(a) and its proof
apply, because they concern the choices at gate $g$ only; this gives $\alpha_g=1-q_g\alpha_{g-1}$.
Let $g$ have type $\mathrm L_2$. Then $x_g=1$ exactly if $x_{g-1}=0$ and the read is $0$, and the
read is $0$ exactly if it is resampled to $0$, which has probability $B$ and is independent of
the earlier choices, or it is kept and $x_{g-2}=0$. So
$\PP[x_g=1]=B\alpha_{g-1}+(1-\eps)\PP[x_{g-2}=0,\,x_{g-1}=0]$. As $x_{g-2}=1$ forces $x_{g-1}=0$,
$\PP[x_{g-2}=0,\,x_{g-1}=0]=\alpha_{g-1}-(1-\alpha_{g-2})$. With $B+(1-\eps)=r$ this is the
formula. For every gate, of either kind, the read is $0$ with a conditional probability between
$B$ and $r$ given the earlier choices, so $B\alpha_{g-1}\le1-\alpha_g\le r\alpha_{g-1}$, and
$\alpha_g\in[lo,hi]$ follows as in the proof of Lemma~\ref{lem:onebit}(b).

(b) The proof of Proposition~\ref{prop:designs}(a) applies to every $d$ and to mixed types: the
target of a gate $g$ of mixed type is $1$ with probability $\mu_0(g)+\sum_c\mu_c(g)d_c$, whose
derivative in $d_c$ is $\mu_c(g)$, and the regeneration argument does not depend on the types.
It gives that $\Phi$ is continuously differentiable and that $\partial\Phi(d)_k/\partial d_c$ is
the sum over the gates $g\le h$ of mixed type of $\mu_c(g)$ times
$\PP[(x_h,x'_h)=(0,1)]-\PP[(x_h,x'_h)=(1,0)]$ for the two copies $x$, $x'$ of
Definition~\ref{def:two}. The absolute value of this difference is at most $p_g(h)$.

(c) The copies agree before $g$, and they differ at $g$ exactly if the read of $g$ is kept and
$x_{g-1}=0$. Call a gate $h\ge g$ in state $\mathsf a$, $\mathsf b$ or $\mathsf c$ according to
the three events of the statement. In state $\mathsf a$ both copies hold $0$ at $h-1$: the copy
that holds $1$ at $h$ does, and they agree there. In state $\mathsf b$ the pairs
$(x_{h-1},x_h)$ of the two copies are $(0,1)$ and $(1,0)$. In state $\mathsf c$ both copies hold
$0$ at $h$: the copy that holds $1$ at $h-1$ does, and they agree there. If the copies agree at
$h-1$ and at $h\ge g$, they agree from then on, because no gate looks further back than two
positions and all later choices are the same. Now let $h>g$. If gate $h$ reads the same bit $v$
in both copies, then for $v=1$ both hold $0$ at $h$, so the states $\mathsf a$ and $\mathsf b$ of
$h-1$ become $\mathsf c$ and the state $\mathsf c$ ends; and for $v=0$ both copies take the
complement of their value at $h-1$, so $\mathsf a$ and $\mathsf b$ become $\mathsf b$, and
$\mathsf c$ ends, as both copies hold $1$ at $h$. A gate of mixed type reads the same bit in
both copies, and it is $0$ with probability $q_h$, independently of the earlier choices; a gate
of type $\mathrm L_1$ behaves in the same way with the event $Z_h$ of
Lemma~\ref{lem:onebit}(a) in place of the event that the read is $0$. This gives the first
formula. Let $h$ have type $\mathrm L_2$. Its read is resampled to $0$ with probability $B$ and
to $1$ with probability $A$, with the effects just described. With probability $1-\eps$ it is
kept, and then $x_h=\NOR(x_{h-1},x_{h-2})$ in both copies. From state $\mathsf a$ at $h-1$ both
copies hold $0$ at $h-2$, so both take the complement of their value at $h-1$: state
$\mathsf b$. From state $\mathsf b$ each copy holds a one at $h-2$ or at $h-1$, so both hold $0$
at $h$: state $\mathsf c$. From state $\mathsf c$ both copies hold $0$ at $h-1$ and they differ
at $h-2$, so they differ at $h$ and agree at $h-1$: state $\mathsf a$. So
$a_h=(1-\eps)c_{h-1}$, $b_h=(B+1-\eps)a_{h-1}+Bb_{h-1}$ and
$c_h=Aa_{h-1}+(A+1-\eps)b_{h-1}$, and $A+1-\eps=1-B$.
\end{proof}

\begin{lemma}[runs of pointers two back]\label{lem:runs}
Let $0.455\le\eps<1$, $x\in[lo,hi]$, $y\in[0,1]$ and $q=B+(1-\eps)y$. Put $\alpha_{-1}=x$,
$\alpha_0=1-xq$ and $\alpha_j=(2-\eps)-r\alpha_{j-1}-(1-\eps)\alpha_{j-2}$ for $j\ge1$. Let
$\Lambda$ be the linear map
$(v_1,v_2,v_3)\mapsto\bigl((1-\eps)v_3,\ rv_1+Bv_2,\ Av_1+(1-B)v_2\bigr)$, and
\[
v^0=\bigl((1-\eps)\,x\,\hat f(y),\ q\,f(x),\ (1-q)\,f(x)\bigr),\qquad v^j=\Lambda v^{j-1}\quad(j\ge1),
\]
with $f$ and $\hat f$ as in Lemma~\ref{lem:potential}.
\begin{enumerate}[label=(\alph*),leftmargin=2.4em]
\item $\alpha_j\in[lo,hi]$ for every $j\ge-1$.
\item $v^k_1+v^k_2\le\frac9{10}f(\alpha_k)$ for $1\le k\le17$.
\item $v^{18}_1+v^{18}_2+v^{18}_3\le\frac9{10}\min\{f(lo),f(hi)\}$.
\item $v^k_1+v^k_2\le\frac9{10}f(\alpha_k)$ for every $k\ge1$.
\end{enumerate}
\end{lemma}

\begin{proof}
(a) The interval $[lo,hi]$ is mapped into itself by $\alpha\mapsto1-q'\alpha$ for every
$q'\in[B,r]$, since $1-B\,lo=hi$ and $1-r\,hi=lo$. We show by induction that
$\alpha_j=1-q'_j\alpha_{j-1}$ with some $q'_j\in[B,r]$, for $j\ge0$. This holds for $j=0$ with
$q'_0=q$. If it holds for $j$, then $\alpha_j>0$ and
$\alpha_j+\alpha_{j-1}-1=(1-q'_j)\alpha_{j-1}$ lies between $0$ and $\alpha_j$, the latter
because $\alpha_{j-1}\le1$. So
$\alpha_{j+1}=1-B\alpha_j-(1-\eps)(\alpha_j+\alpha_{j-1}-1)=1-q'_{j+1}\alpha_j$ with
$q'_{j+1}=B+(1-\eps)(\alpha_j+\alpha_{j-1}-1)/\alpha_j\in[B,r]$.

(b), (c) These are verified by the program \texttt{runs\_check.c}, with the interval arithmetic
of \texttt{potential\_check.c} (Lemma~\ref{lem:potential}). The components of $v^j$ are sums and
products of nonnegative numbers, so upper bounds for $\eps$-dependent constants, for $f$ on an
interval of $x$ and for $\hat f$ on an interval of $y$ give upper bounds for $v^j$. For
$\alpha_k$ the program uses $\alpha_k=C_k+x(P_k+Q_kq)$, where $C_k$, $P_k$ and $Q_k$ depend on
$\eps$ only and follow the recursion of $\alpha$, with the term $2-\eps$ for $C$ and without it
for $P$ and $Q$, from $(C,P,Q)_{-1}=(0,1,0)$ and $(C,P,Q)_0=(1,0,-1)$; it encloses these
numbers for an interval of noise levels, intersects the enclosure of $\alpha_k$ with $[lo,hi]$,
which is allowed by (a), and bounds $f$ from below on the result by its values at the two ends.
The runs \texttt{runs\_check 0.45499 0.50001 0.8999 18 0.0005} and
\texttt{runs\_check 0.49999 1.0 0.8999 18 0.005} cover $[0.455,1)$ in $192$ intervals of noise
levels and $9142$ boxes; as in Lemma~\ref{lem:potential} they use a factor a little below the
stated one, because the decimal arguments are rounded.

(d) For $k\le17$ this is (b). The coordinates of $\Lambda v$ add up to
$v_1+v_2+(1-\eps)v_3$, so the sum of the coordinates of $v^j$ does not increase with $j$. For
$k\ge18$ therefore $v^k_1+v^k_2\le v^{18}_1+v^{18}_2+v^{18}_3\le\frac9{10}\min\{f(lo),f(hi)\}$
by (c), and this is at most $\frac9{10}f(\alpha_k)$ by (a), because $f$ increases up to
$43/76$ and decreases after it.
\end{proof}

\begin{theorem}[pointers two back]\label{thm:two}
Let $0.455\le\eps<1$ and $\theta=0.9998$, and consider a design with independent reads and
pointers two back. Then inequality~\eqref{eq:Mpot} holds for every $d\in[0,1]^{\mathcal C}$, with
the matrix $\overline M(d)$ of Definition~\ref{def:two}:
\[
\sum_c\overline M_{k,c}(d)\,f(1-d_c)\ \le\ \theta\,f\bigl(1-\Phi(d)_k\bigr)\qquad\text{for every class }k .
\]
Consequently parts (a) to (d) of Theorem~\ref{thm:indep} hold for these designs: the density map
is a contraction of $\mathcal K$ in the metric $\delta$ with factor $\theta$, it has exactly one
fixed point, to which the densities converge from every start, $\rho_s\le\rho_u\le\theta$ there,
and the densities of layers with different designs of this kind forget the first layer. Here
$\rho_u$ is the spectral radius of the matrix $\overline M(d^*)$ of Definition~\ref{def:two}.
\end{theorem}

\begin{proof}
Let $d\in[0,1]^{\mathcal C}$ and $\beta_c=1-d_c$. For a gate $g$ of mixed type put
$s_g=\sum_c\mu_c(g)f(\beta_c)$; as in the proof of Theorem~\ref{thm:indep},
$s_g\le\hat f(y_g)$. For a gate $h$ of class $k$ let
$N_h=\sum_c\overline M_{k,c}(d)f(\beta_c)=\sum_{g\le h}s_g\,p_g(h)$, and let $N^a_h$, $N^b_h$
and $N^c_h$ be the same sums with the numbers $a_h$, $b_h$ and $c_h$ of
Lemma~\ref{lem:twochain}(c) for $g$ in place of $p_g(h)$, so that $N_h=N^a_h+N^b_h$. These sums
are finite by the regeneration argument of Definition~\ref{def:design}, and they are periodic in
$h$. Put $I_h=(1-\eps)\alpha_{h-1}s_h$ if $h$ has a mixed type and $I_h=0$ otherwise, and
$D_h=N_h-\theta f(\alpha_h)$. We show that for every gate $h$ of a chain that has a gate which
is not of type $\mathrm L_2$ there are a gate $h'<h$ and a number $\kappa\in[0,1)$ with
\begin{equation}\label{eq:Dtwo}
D_h\le\kappa\,\bigl(N_{h'}-f(\alpha_{h'})\bigr)\le\kappa\,D_{h'} .
\end{equation}
Then $D_h\le0$ for every such $h$: if the largest of these finitely many numbers were positive
and attained at $h$, \eqref{eq:Dtwo} would give $D_h\le\kappa D_{h'}\le\kappa D_h<D_h$. In a chain all of whose
gates have type $\mathrm L_2$ no gate has a mixed type, and $N_h=0$. So the inequality of the
theorem holds, because $\Phi(d)_k=1-\alpha_h$.

Let $h$ have a mixed type or the type $\mathrm L_1$. By Lemma~\ref{lem:twochain}(c), summed over
$g$ with the weights $s_g$,
\[
N^a_h=I_h,\qquad N^b_h=q_hN_{h-1},\qquad N^c_h=(1-q_h)N_{h-1} ,
\]
so $N_h=q_hN_{h-1}+I_h$. By Lemma~\ref{lem:twochain}(a), $\alpha_{h-1}\in[lo,hi]$ and
$\alpha_h=1-q_h\alpha_{h-1}$, so~\eqref{eq:gatepot} holds for $h$ as in the proof of
Theorem~\ref{thm:indep}, and \eqref{eq:Dtwo} holds with $h'=h-1$ and $\kappa=q_h\le r<1$.

Let $h$ have type $\mathrm L_2$, let $h_0<h$ be the last gate before it that is not of type
$\mathrm L_2$, and $k=h-h_0$. Write $x=\alpha_{h_0-1}$, $y=y_{h_0}$ and $q=q_{h_0}$. By the
display above and by Lemma~\ref{lem:twochain}(c) for the gates $h_0+1,\dots,h$, all of type
$\mathrm L_2$,
\[
(N^a_h,N^b_h,N^c_h)=\Lambda^k\bigl(I_{h_0},\ qN_{h_0-1},\ (1-q)N_{h_0-1}\bigr),
\]
with the map $\Lambda$ of Lemma~\ref{lem:runs}. Hence $N_h=\kappa N_{h_0-1}+\iota I_{h_0}$, where
$\kappa$ and $\iota$ are the sums of the first two coordinates of $\Lambda^k(0,q,1-q)$ and of
$\Lambda^k(1,0,0)$; both are nonnegative. By Lemma~\ref{lem:twochain}(a) the numbers
$\alpha_{h_0-1},\alpha_{h_0},\dots,\alpha_h$ are the numbers $\alpha_{-1},\alpha_0,\dots,\alpha_k$
of Lemma~\ref{lem:runs}, and the vector $v^0$ of that lemma is
$(1-\eps)x\hat f(y)\,(1,0,0)+f(x)\,(0,q,1-q)$. So Lemma~\ref{lem:runs}(d) says
\[
\kappa\,f(x)+\iota\,(1-\eps)\,x\,\hat f(y)\ \le\ \tfrac9{10}f(\alpha_h)\ \le\ \theta f(\alpha_h).
\]
As $I_{h_0}\le(1-\eps)x\hat f(y)$, this gives
$D_h\le\kappa\bigl(N_{h_0-1}-f(\alpha_{h_0-1})\bigr)$, which is \eqref{eq:Dtwo} with
$h'=h_0-1$. Finally $\kappa<1$: the coordinates of $\Lambda(0,q,1-q)$ add up to
$q+(1-\eps)(1-q)=1-\eps(1-q)<1$, and $\Lambda$ does not increase the sum of the coordinates.

For the consequences, the proofs of parts (a) to (d) of Theorem~\ref{thm:indep} use
\eqref{eq:Mpot}, that $\Phi$ maps into $\mathcal K$ and is continuously differentiable, and the
bound of the derivative by $\overline M$. The last three are Lemma~\ref{lem:twochain}(a)
and~(b), and the argument above, like that of Theorem~\ref{thm:indep}, does not use that the
layer below has the same design.
\end{proof}

\begin{remark}[pointers two back: numbers, meaning and limits]\label{rem:two}
\begin{enumerate}[label=(\roman*),leftmargin=2.4em]
\item \emph{The margin.} The constant $\frac9{10}$ of Lemma~\ref{lem:runs} is not close.
Numerically, on a grid, the largest value of $(v^k_1+v^k_2)/f(\alpha_k)$ at noise $0.455$ is
$0.798$, $0.604$, $0.557$, $0.456$ and $0.364$ for $k=1,\dots,5$ and decreases further; for
$k=1$ it is $0.754$ at noise $0.50$ and $0.514$ at $0.80$. For a NOT gate before the run, which
is the case $y=1$, it is $0.639$ at noise $0.455$. So at a gate of type $\mathrm L_2$ the
potential has room, and the factor $\theta$ of the theorem is needed only at the gates with
far reads. For the design with the types $\mathrm L_2$ and $\mathrm E$ at the two phases, the
design of the mixed base of Remark~\ref{rem:mixed}, $\rho_u=0.585$ at noise $0.46$, against
$\lambda=0.972$ for the plain chain.
\item \emph{A test.} The script \texttt{two\_back\_test.py} takes random designs of period $2$
to $7$ with the types $\mathrm E(c)$, $\mathrm Z_0$, $\mathrm Z_1$, $\mathrm L_1$ and
$\mathrm L_2$ and random $d$, and computes $\overline M(d)$ twice: from the joint law of the two
copies, and from the recursion of Lemma~\ref{lem:twochain}(c). In $700$ designs the two agreed
to $10^{-15}$, the formula of Lemma~\ref{lem:twochain}(a) to $10^{-12}$, and the largest ratio
of the two sides of the inequality of the theorem was $0.981$, in a design without a gate of
type $\mathrm L_2$. This is numerical and tests the statement, not the proof.
\item \emph{What it means for the iterations.} For the set-valued table iteration of
Definition~\ref{def:setval} with reach $2$, take a chain all of whose reads are determined
except one. The set of windows that the iteration keeps at a later position consists of the
windows of the two copies above, so the position is undetermined exactly if the copies differ
there, and the state $\mathsf c$ is a position that is determined while its window is not. So in
the ideal model $\overline M$ is the mean spread, for that iteration, of a single undetermined
read when all other reads are determined; with two or more undetermined reads in a chain the
sets of windows are not unions of such pairs. For the relevance-pruned iteration even this is
not so: on the mixed base it needs many rounds up to noise about $0.55$
(Remark~\ref{rem:mixed}), because it cannot use that a kept gate of type $\mathrm L_2$ turns
a difference in state $\mathsf b$ into the hidden state $\mathsf c$, which ends at the next
gate if that gate has a far read.
\item \emph{What is missing.} This is a statement about the ideal model, like
Theorem~\ref{thm:indep}. Theorem~\ref{thm:layered} turns Theorem~\ref{thm:indep} into a
statement about smoothed instances with the help of a lower and an upper trace
(Lemma~\ref{lem:lowup}), which exist for the relevance-pruned iteration. We know no such pair
for the table iteration with reach $2$, and so we have no theorem of that kind for the sets of windows themselves.
Subsection~\ref{sec:twolayered} goes another way: it rewrites the gates of type $\mathrm L_2$
so that a lower and an upper trace exist, and proves the theorem on smoothed instances for the
iteration that these traces define, which the table iteration with reach $2$ matches or beats.
Pointers three or more back
are not covered either: the state of a difference then has more values, and the clean chain is
no longer described by the numbers $\alpha_g$ alone.
\end{enumerate}
\end{remark}

\subsection{Above \texorpdfstring{$4/9$}{4/9} the latch is not the best design}\label{sec:beyond}

\begin{remark}[designs above $4/9$; numerical]\label{rem:beyond}
The following designs have one rail and independent reads: every gate has a type $\mathrm E(c)$,
where $c$ is a phase of the rail, except one gate of type $\mathrm Z_1$. A list gives the types of
the phases $0,1,\dots,T-1$, and $\mathrm E(c)$ is written $c$.
\begin{align*}
D_{12}&:\ 2,4,2,6,2,\mathrm Z_1,1,2,2,4,2,4\\
D_{16}&:\ 11,12,12,14,12,14,12,14,12,14,12,14,12,0,11,\mathrm Z_1\\
S_{24}&:\ 22,9,22,9,22,9,22,9,22,9,\mathrm Z_1,3,8,9,8,9,8,9,8,11,8,11,8,9
\end{align*}
and $D_{40}$, of period $40$, is listed in the script \texttt{organ\_long\_designs.py}. They were
found by the referee of this section, by a local search that excludes the designs with types
$\mathrm E$ and $\mathrm U$ only. Those designs all have $\rho_u=\lambda$
(Proposition~\ref{prop:designs}(b)), and this plateau had absorbed our own local searches. We
recomputed all numbers below with a separate program written from
Definition~\ref{def:design} and Lemma~\ref{lem:onebit} (script \texttt{organ\_long\_designs.py});
the referee's exact program and ours agree to six digits on all entries, the program of
Remark~\ref{rem:designs} agrees where it was run (all four designs at noise $0.46$; at $0.445$ it
does not converge for $S_{24}$), and the referee's Monte Carlo simulation of coupled chains gave
$1.0062$ for $D_{16}$ at noise $0.445$, with error about $10^{-3}$.
\begin{center}
\begin{tabular}{lccccc}
\toprule
noise & $0.445$ & $0.448$ & $0.45$ & $0.455$ & $0.46$\\
\midrule
$\lambda$ & $0.99900$ & $0.99359$ & $0.98999$ & $0.98095$ & $0.97190$\\
$\rho_u(D_{12})$ & $1.00186$ & $0.99607$ & $0.99221$ & $0.98255$ & $0.97290$\\
$\rho_u(D_{16})$ & $1.00555$ & $0.99968$ & $0.99578$ & $0.98601$ & $0.97626$\\
$\rho_u(D_{40})$ & $1.00733$ & $1.00138$ & $0.99741$ & $0.98750$ & $0.97761$\\
$\rho_s=\rho_u(S_{24})$ & $1.00430$ & $0.99846$ & $0.99457$ & $0.98484$ & $0.97511$\\
\bottomrule
\end{tabular}
\end{center}
The values of $\rho_s$ for $D_{12}$, $D_{16}$ and $D_{40}$ are below $0.6$.
\begin{enumerate}[label=(\roman*),leftmargin=2.4em]
\item \emph{Changes that keep spreading.} All four designs have $\rho_u>\lambda$ at every noise
level in the table, and at noise $0.445$ all have $\rho_u>1$. For $D_{40}$, $\rho_u=1.00039$ at
$0.4485$ and $0.99939$ at $0.449$, so it crosses $1$ at about $0.4487$. The largest ratio
$\rho_u/\lambda$ the referee found among designs that are not made of types $\mathrm E$ and
$\mathrm U$ only is $1.0083$ at noise $0.445$, $1.0062$ at $0.46$ and $1.0017$ at $0.50$, for
periods $32$ to $48$. The best such designs it found with period $8$ to $11$ stay below
$\lambda$ at noise $0.46$, and the exhaustive searches of Remark~\ref{rem:designs} stop at period
$7$. With several rails or mixed types shorter designs suffice. The referee of
Theorem~\ref{thm:indep} found the following design with three rails of period $6$, where
$\mathrm E(r,\varphi)$ reads class $(r,\varphi)$:
\begin{align*}
&\mathrm E(3,0),\ \mathrm Z_1,\ \mathrm E(3,2),\ \mathrm E(1,4),\ \mathrm E(3,4),\ \mathrm E(1,0);\\
&\mathrm E(3,0),\ \mathrm E(1,2),\ \mathrm E(3,2),\ \mathrm E(1,4),\ \mathrm E(3,4),\ \mathrm E(1,0);\\
&\mathrm E(2,0),\ \mathrm E(3,1),\ \mathrm E(2,2),\ \mathrm E(3,3),\ \mathrm E(2,4),\ \mathrm E(3,5).
\end{align*}
It has $\rho_u=0.97309$ at noise $0.46$ and $1.00002$ at $0.446$ (by the program of
Remark~\ref{rem:designs}). With two rails of period $2$ and one gate
thinned to position $1$ with probability $0.067$, it found $\rho_u=1.0017$ at noise $0.445$.
With one gate of type $\mathrm Z_1$ the densities are no longer constant along the period, the
classes differ in how much a flipped position of them spreads, and the gates list the classes
that spread most. With type $\mathrm U$ in place of the types $\mathrm E(c)$ the same gate lowers
$\rho_u$ (Remark~\ref{rem:designs}(iii)). On smoothed instances, with layers of $320000$
positions over zero gates, the referee measured for $D_{16}$ at noise $0.445$ a one-layer matrix
with spectral radius $1.0045\pm0.0007$ over seven seeds, against $0.9991\pm0.0011$ for the plain
chain over six seeds.
\item \emph{A bit that is kept.} At noise $0.445$ Newton's method from $40$ random starts finds
one fixed point of the density map of $S_{24}$, and its Jacobian there has the eigenvalue
$-1.0043$, minus the Perron root of $\overline M$. Iterated from the vector with all
entries $0.35$, the density map converges to a cycle $\{a,\Phi(a)\}$ of length two with
$\|\Phi(a)-a\|_\infty=0.130$, and the Jacobian of $\Phi\circ\Phi$ at $a$ has spectral radius
$0.983$, so the cycle is attracting. Then $\Phi\circ\Phi$ has the two attracting fixed points $a$
and $\Phi(a)$, as the ideal map of the latch has $L$ and its mirror image below $4/9$, and a pair
of layers keeps a bit: the phase of the oscillation. The amplitude $\|\Phi(a)-a\|_\infty$ is
$0.146$ at noise $0.4444445$, $0.130$ at $0.445$, $0.097$ at $0.446$ and $0.040$ at $0.447$; at
$0.448$ the iteration from the same start converges to the fixed point, where $\rho_s=0.998$. On smoothed instances, $60$ layers of $240000$ positions with the
first layer prepared in one phase, the referee found the alternating pattern still present in
layers $50$ to $59$ at noise $0.445$, with $0.83\pm0.09$ times the ideal amplitude over eight
seeds, and $0.23\pm0.07$ at noise $0.46$, where the ideal value after the same number of layers is
about $0.25$; at this size the run cannot tell a multiplier of $1.004$ from one of $0.999$.
\item \emph{Where it stops.} By Theorem~\ref{thm:indep}, no design with independent reads has
$\rho_u>0.9998$, a cycle, or a second fixed point at noise $0.455$ or above. For single changes
the designs above and the value iteration of Remark~\ref{rem:indep}(ii) meet at about $0.449$.
So the supremum of the noise levels at which some design with independent reads has more than one
attracting set lies between about $0.447$, numerically and for one design, and $0.455$.
\item \emph{What it means.} In the ideal model the level $4/9$ belongs to the latch, not to
layered arrangements of chains. Section~\ref{sec:s24} turns $S_{24}$ into an embedding and moves
the lower end of the critical interval to $0.4468$. The proof of Theorem~\ref{thm:latch} uses that
the ideal map of the latch is monotone. The period of $S_{24}$ is even, and every listed class is
listed only by gates of its own parity. By the alternation argument in the proof of
Proposition~\ref{prop:designs}(c), $\partial\Phi_k/\partial d_c$ then has the sign of
$-(-1)^{k+c}$ at every $d$, so $\Phi$ reverses the order in which the even phases go up and the odd
ones down, and $\Phi\circ\Phi$ preserves it; Lemma~\ref{lem:envelope} is the version of this for
finite chains with slack, and with it the sandwich of Theorem~\ref{thm:latch} works stage by stage.
The gate stage reads both inputs with a bias toward the top state. Designs with $\rho_u>1$ are also
natural candidates for bases on which the table--Jacobi iteration of Section~\ref{sec:table} is
slow above $4/9$, because a wrong value read from the layer below spreads like a flipped position.
We have not tested this.
\end{enumerate}
\end{remark}

\subsection{Short periods: an exhaustive search}\label{sec:search}

\begin{remark}[the search]\label{rem:designs}
The script \texttt{organ\_multiplier.c} looks for fixed points of $\Phi$ by a lagged iteration from
five starting vectors (one with all entries $0.35$, four with entries $0.03$ and $0.9$) and computes
$\rho_u$ and $\rho_s$ at every fixed point it finds. It propagates the law of the last $K$ values
of one chain, and of two coupled chains, gate by gate, so the two matrices are exact up to
rounding. The lagged iteration has the same fixed points as $\Phi$ but not the same stability: it
can converge to a fixed point that repels under $\Phi$, it can miss a fixed point, and it does not
see cycles of $\Phi$. So ``found'' below means found by this iteration. The spectral radius was
taken as $(\|M^{256}\|/\|M^{128}\|)^{1/128}$. This ratio converges only like $1/n$ when the
dominant eigenvalue is defective or not the only one of its modulus. For one rail of period at
most $5$ it gives $\rho_u$ exactly to rounding, but $\rho_s$ only to $6.3\cdot10^{-3}$; for
several rails the referee found errors up to $4.6\cdot10^{-3}$ in $\rho_u$, up to
$1.0\cdot10^{-3}$ downwards where $\rho_u>0.9\lambda$, and up to $8.2\cdot10^{-3}$ in
$\rho_s$. The bounds below for several rails are therefore uncertain by about $10^{-3}$; the
referee recomputed random samples of $150{,}000$ designs with two rails of period $3$, $150{,}000$
with three rails of period $2$ and $100{,}000$ mirrored designs of period $5$ with an exact
spectral radius, and found none above $\lambda$.
\begin{enumerate}[label=(\roman*),leftmargin=2.4em]
\item \emph{Noise $0.46$.} Here $\lambda=0.97190$. The families are the following; ``all''
means the types $\mathrm E$, $\mathrm U$, $\mathrm S$, $\mathrm Z_0$, $\mathrm Z_1$ and
$\mathrm L_1,\dots,\mathrm L_K$, and ``no $\mathrm U$, $\mathrm S$'' the same without those two.
In a mirrored design there are two rails, and the second is the first with the two rails
exchanged in its types $\mathrm E$.
\begin{center}
\begin{tabular}{llr}
\toprule
rails and periods & types & designs\\
\midrule
one rail, $T=7$, $K=4$ & no $\mathrm U$, $\mathrm S$ & $62{,}468{,}581$\\
one rail, $T\le6$, $K=4$ & all & $7{,}785{,}749$\\
one rail, $T=5$, $K=5$ & all & $505{,}056$\\
two rails, $T=3$, $K=4$ & all & $6{,}802{,}033$\\
three rails, $T=2$, $K=4$ & all & $6{,}081{,}831$\\
mirrored, $2\le T\le5$, $K=4$ & all & $1{,}938{,}392$\\
mirrored, $T=6$, $K=4$ & no $\mathrm U$, $\mathrm S$ & $33{,}965{,}568$\\
\bottomrule
\end{tabular}
\end{center}
No design has a computed $\rho_u$ or $\rho_s$ above $0.9720$, and the iteration found no second
fixed point and converged from every start. The largest values, $\rho_u=\rho_s=\lambda$, belong
to designs with types $\mathrm E$ and $\mathrm U$ only (Proposition~\ref{prop:designs}(b)). The
referee added three checks: local pointers up to $\mathrm L_5$ for periods $2$ to $4$, up to
$\mathrm L_6$ for periods $2$ and $3$ and up to $\mathrm L_8$ for period $2$, run with the program
and a more accurate spectral radius ($(\|M^{4096}\|/\|M^{2048}\|)^{1/2048}$), with none above
$0.97195$; with its own code, Newton's method from random starts, which found a single fixed point
for every design with one rail and period at most $5$ and for $30{,}000$ random designs with three
rails; and the true density map, which converged from all five starts in $100{,}000$ random
designs of several families. As a check of Theorem~\ref{thm:indep}, the $214{,}352{,}320$ designs with one
rail of period $8$ and the types $\mathrm E$, $\mathrm Z_0$, $\mathrm Z_1$ and $\mathrm L_1$ also
have $\rho_u\le\lambda$ at this noise level. With period $12$ and more there are designs above
$\lambda$ (Remark~\ref{rem:beyond}).
\item \emph{Just above $4/9$.} At noise $0.445$, where $\lambda=0.99900$, no design with one rail,
period at most $6$ and all types has a computed $\rho_u$ above $0.99900$. For $746$ of them,
all with types $\mathrm E$ only, the iteration had not converged after $6000$ steps and computed
nothing; for these designs Proposition~\ref{prop:designs}(b) gives $\rho_u=\lambda$ at the
constant fixed point.
\item \emph{Rare changes.} In a rail of period $32$ whose gates have type $\mathrm U$, one gate of
another type lowers $\rho_u$ at noise $0.46$ from $0.97190$ to $0.92130$ ($\mathrm Z_0$),
$0.95475$ ($\mathrm Z_1$ and $\mathrm L_1$), $0.94900$, $0.96381$, $0.96127$, $0.96783$,
$0.96671$, $0.96961$ and $0.96912$ ($\mathrm L_2,\dots,\mathrm L_8$). With period $64$ the loss
is about half as large. Among the designs of period $8$, $16$, $32$ and $64$ in which a block of at
most three gates has arbitrary types ($\mathrm E$ of one of up to five phases, $\mathrm Z_0$,
$\mathrm Z_1$, $\mathrm L_1,\dots,\mathrm L_4$) and all other gates have type $\mathrm U$, none is
above $0.97190$. In a background of type $\mathrm U$ a rare change does not help; in a background
of types $\mathrm E(c)$ it does (Remark~\ref{rem:beyond}).
\item \emph{Below $4/9$.} At noise $0.30$ the design with one rail and the types $\mathrm E(2)$,
$\mathrm L_3$, $\mathrm E(2)$, $\mathrm E(2)$ has $\rho_u=1.3015$, where $\lambda=1.2480$; at
noise $0.38$ the design $\mathrm E(4)$, $\mathrm E(4)$, $\mathrm E(4)$, $\mathrm L_5$,
$\mathrm E(4)$ has $\rho_u=1.1222$, where $\lambda=1.1139$. Local pointers are not needed: at
noise $0.10$ the design $\mathrm E(0)$, $\mathrm E(0)$, $\mathrm E(0)$, $\mathrm Z_1$ has
$\rho_u=1.7474$, where $\lambda=1.5254$, and $D_{16}$ of Remark~\ref{rem:beyond} has
$\rho_u/\lambda=1.0138$ at noise $0.40$ and $1.0073$ at $0.44$. Below $4/9$ a design with types
$\mathrm E$ only can have other fixed points than the constant one, as the latch does. Among the
short designs the advantage over $\lambda$ disappears near noise $0.39$: at noise $0.40$ no design
with one rail, period $5$, $K=5$ and all types ($505{,}056$ designs) is above $\lambda$, nor is any
of period at most $4$ with $K=4$; at noise $0.42$ and $0.44$ none with one rail, period $6$ and no
$\mathrm U$, $\mathrm S$ is ($2{,}939{,}328$ designs), and at both levels every design for which
the iteration found two fixed points has types $\mathrm E$ only. A rough reason why far local
pointers help only below $4/9$: a gate that lists a far position of its own chain works like a
layer inside the layer. If a fraction $\theta$ of the gates did this without any other effect, a
change would produce
$(1-\theta)\lambda(1+\theta\lambda+\theta^2\lambda^2+\dots)=(1-\theta)\lambda/(1-\theta\lambda)$
changes in the next layer, which is more than $\lambda$ exactly when $\lambda>1$. This is a
heuristic, and it does not apply to the designs of Remark~\ref{rem:beyond}.
\item \emph{Dependent reads, on real instances.} In a design the positions listed by different
gates of the layer below are independent. The script \texttt{flipspread.c} measures the simplest
dependent arrangement on smoothed instances: zero gates, then $24$ layers of $200000$ positions
in the last fiftieth of the instance; the gates of a layer come in groups of $m$, and a group
lists $m$ consecutive positions of the layer below, starting at a uniformly random one. Two
copies are evaluated with the same pointers, and in the fourth layer one position in $200$ of the
second copy is flipped. The factor per layer by which the number of positions that differ
changes, from the layer $2$ to the layer $10$ after that one and from the layer $10$ to the layer
$20$, summed over six seeds:
\begin{center}
\begin{tabular}{lccccc}
\toprule
noise & $m=1$ & $m=2$ & $m=3$ & $m=8$ & $m=200000$\\
\midrule
$0.40$ & $1.062$, $1.052$ & $0.963$, $0.975$ & $0.932$, $0.949$ & $0.897$, $0.912$ & $0.877$, $0.895$\\
$0.46$ & $0.967$, $0.980$ & $0.902$, $0.917$ & $0.858$, $0.921$ & $0.836$, $0.864$ & $0.824$, $0.869$\\
\bottomrule
\end{tabular}
\end{center}
The entries are uncertain by about $0.01$; the referee's own simulator gave $0.961$, $0.974$ and
$0.889$, $0.919$ at noise $0.46$ for $m=1$ and $m=2$, and $1.059$, $1.048$ and $0.977$, $0.966$ at
$0.40$. For $m=1$ this is a random block base, and the ideal value is $\lambda=1.079$ and $0.972$;
at noise $0.40$ the differences reach $4\%$ of a layer, where they no longer spread
independently. For $m=200000$ a layer is a shifted copy of the layer below. In this arrangement
gates that list neighbouring positions spread a change less, not more.
\item \emph{Real instances and the ideal model.} The referee compared the column sums of
$\overline M$ with the mean numbers of changed positions on smoothed instances, with layers of
$40000$ positions over zero gates: $0.9725\pm0.0039$ against $0.97190$ for the plain chain at
noise $0.46$, $5.401\pm0.013$ against $5.391$ for $\mathrm E(2),\mathrm L_3,\mathrm E(2),\mathrm E(2)$ at
$0.30$, $3.743\pm0.010$ against $3.746$ for the same design at $0.46$, and similar agreement for
designs with types $\mathrm S$, $\mathrm U$ and $\mathrm Z_1$ and for the latch. At this size some
runs deviate by $1$ to $2\%$ (the design $\mathrm E(0),\mathrm S,\mathrm S$, the two rails of the
latch, one short run); for $\mathrm E(0),\mathrm S,\mathrm S$ the deviation disappears with layers of
$300000$ positions ($2.358$ to $2.369$ against $2.372$).
\item \emph{What it does not cover.} The searches are finite; the lagged iteration need not find
every fixed point and does not see cycles; $\rho_u<1$ concerns single changes. Not covered are
periods above $7$ (where Remark~\ref{rem:beyond} finds better designs), local pointers of reach
above $4$ except as stated in (i) and (iii), several rails with different programs beyond
period $3$, unequal proportions of rails with gates of type $\mathrm U$, dependent reads other than
those of (v), designs that are not periodic, and bases that are not layered.
\end{enumerate}
\end{remark}

\begin{remark}[dependent reads; numerical]\label{rem:dependent}
In Theorem~\ref{thm:indep} the reads of a chain are independent of each other and of the chain.
A base need not respect this: the gates of a chain can list neighbouring positions of one chain
below, or a position together with the position that its gate lists. The script
\texttt{dep\_design.c} measures, in the ideal model, how a flipped position spreads then. A layer
has $M$ positions and continues the chain of the layer below, and a design of period $T$ gives
every phase a rule: $\mathrm U$, a uniformly random position of the layer below; $\mathrm I(c)$,
one of phase $c$; $\mathrm F(d)$, the position of the layer below at the gate's own offset plus
$d$; $\mathrm H(d)$, the same two layers below; $\mathrm R(d)$, the position at distance $d$ from
the last one listed by a rule $\mathrm U$, $\mathrm I$, $\mathrm F$ or $\mathrm R$; a local
pointer $\mathrm L_k$; a NOT gate; position $0$ or
position $1$. A read is resampled with probability $\eps$ and is then a new bit that is $1$ with
probability $\varpi$, the same in both copies. Two copies are evaluated with the same choices, a
fraction $\delta$ of one layer of the second copy is flipped, and the \emph{growth} is
$(S_2/S_1)^{1/K}$, where $S_1$ and $S_2$ are the numbers of positions that differ, summed over
$K$ consecutive layers and over the $K$ layers after them. For independent reads this estimates
$\rho_u$. At noise $0.46$, with $M=2\cdot10^6$, $K=15$, $\delta=0.003$ and $100$ runs (mean
and standard error):
\begin{center}
\begin{tabular}{llc}
\toprule
design & reads & growth\\
\midrule
$\mathrm U$ & independent (the plain chain) & $0.9698\pm0.0008$\\
$D_{16}$ of Remark~\ref{rem:beyond} & independent & $0.9742\pm0.0011$\\
$\mathrm U,\mathrm R(8)$ & pairs at distance $8$ & $0.9640\pm0.0008$\\
$\mathrm U,\mathrm R(3)$ & pairs at distance $3$ & $0.9335\pm0.0010$\\
$\mathrm U,\mathrm U,\mathrm R(1)$ & one pair of neighbours in three gates & $0.9103\pm0.0011$\\
$\mathrm U,\mathrm R(2)$ & pairs at distance $2$ & $0.9078\pm0.0011$\\
$\mathrm U,\mathrm R(-1)$ & a position and the one before it & $0.8859\pm0.0013$\\
$\mathrm U,\mathrm R(1)$ & a position and the one after it & $0.8845\pm0.0013$\\
$\mathrm F(0)$ & a shifted copy of the layer below & $0.8199\pm0.0015$\\
$\mathrm F(0),\mathrm F(0),\mathrm H(-2)$ & two neighbours and what the first one lists & $0.7776\pm0.0016$\\
$\mathrm F(1),\mathrm F(-1)$ & neighbours in reversed order & $0.7836\pm0.0023$\\
\bottomrule
\end{tabular}
\end{center}
The exact values of the first two rows are $\lambda=0.97190$ and $\rho_u(D_{16})=0.97626$: with
this $\delta$ the estimator is about $0.002$ too low, and more for larger $\delta$ ($0.967$ for
the plain chain at $\delta=0.01$). In all other rows the reads depend on each other, and a
flipped position spreads less; the loss fades as the related positions move apart.

\emph{Searches.} A screen of $59{,}000$ random designs of period $2$ to $6$ used the rules
$\mathrm U$, $\mathrm I(c)$, $\mathrm F(d)$ and $\mathrm H(d)$ with $|d|\le5$,
$\mathrm L_2,\dots,\mathrm L_5$, NOT gates and the positions $0$ and $1$, at noise $0.46$ with
$M=2\cdot10^5$, $K=8$ and $\delta=0.02$. With so many flipped positions the plain chain measures
$0.960$, with a standard deviation of $0.012$ per run; a design whose first run was above
$0.945$ was measured twice more, and $2226$ designs had a mean above it. None of these contains a
NOT gate or a constant read, so the screen compares, in effect, the roughly one design in seven
without them. The largest means were $0.982$ for
$\mathrm H(4),\mathrm I(3),\mathrm U,\mathrm U$ and $0.981$ for
$\mathrm F(-4),\mathrm I(0),\mathrm H(5),\mathrm U$ and
$\mathrm F(5),\mathrm I(3),\mathrm H(3),\mathrm I(2)$; the largest mean of a design with two
rules $\mathrm F$ was $0.967$, and of a design with a local pointer $0.962$. These are maxima of
noisy means. Measured again as in the table, with $100$ runs, the three designs gave
$0.9748\pm0.0006$, $0.9705\pm0.0006$ and $0.9747\pm0.0007$. Hill climbing from the best of
$3000$ random designs, now also with the rules $\mathrm R(d)$, $|d|\le4$, for the periods $3$,
$4$, $6$ and $8$ and $250$ steps each, accepted no step for the periods $3$ and $4$ and one step
each for $6$ and $8$; its last two designs contain rules $\mathrm F$ and $\mathrm R$ and
measured $0.9706\pm0.0007$ and $0.9626\pm0.0009$ as in the table. So it adds little to the
screen.

\emph{The unit.} Some of these values are above the plain chain's. They are growths per layer,
and a rule $\mathrm H$ raises them without making a flipped position change more positions: it
makes the changes skip a layer. If one gate in four reads a uniformly random position two layers
down and the reads are independent, a flipped position still changes $\lambda$ positions on
average, and the growth $g$ per layer solves $g^2=\tfrac34\lambda g+\tfrac14\lambda$, which gives
$g=0.9775$. A rule $\mathrm H$ reads one phase, and then the exact value for independent reads
depends on the design (script \texttt{exact\_indep.py} of the referee): it is $0.9778$ for
$\mathrm H(4),\mathrm I(3),\mathrm U,\mathrm U$, $0.9770$ for
$\mathrm U,\mathrm U,\mathrm U,\mathrm H(-1)$ and $0.9794$ for
$\mathrm I(1),\mathrm H(-2),\mathrm U$, while the number of positions changed per flipped
position, the growth \emph{per generation}, is $\lambda$ in all three. The measurements are
$0.9748\pm0.0006$, $0.9752\pm0.0007$ and $0.9683\pm0.0006$. They are below the values for
independent reads, by the bias of the estimator or by more: the gates with rule $\mathrm H$ of
consecutive periods read positions $T$ apart in one chain, which is a dependent read when the
period is short.

\emph{A second simulator.} An automated referee wrote its own simulator (\texttt{refsim.c}),
which also measures the growth per generation, and added rules that follow pointers, among them
$\mathrm G(d)$: the position listed by the gate at distance $d$ from the last position listed.
It reproduced the table. A position together with the position that its gate lists,
$\mathrm U,\mathrm G(0)$, gave $0.942\pm0.002$ per generation, against $0.971$ to $0.976$ for
controls with the same depths and independent reads. Its searches were: $165$ changes of one
rule of $D_{16}$; $3000$ random designs of period $2$ to $8$ with all rules, the best $24$
measured with $60$ runs; and $2000$ random designs of period $2$ to $6$ whose reads go one layer
down, the best $40$ measured with $30$ runs. Per generation, the best designs of the two random
searches gave $0.9727\pm0.0016$ and $0.9713\pm0.0013$, against $0.9697\pm0.0002$ for the plain
chain, and they are designs whose dependent reads are far apart. The best change of $D_{16}$
gave $0.9758\pm0.0009$ against $0.9743\pm0.0009$ for $D_{16}$ in the same runs; it is a read at
distance $16$, which is nearly independent, and every rule that lists a neighbour lowered the
value by about $0.02$. A spot check at noise $0.45$ gave the same order.

So no design that was measured spreads a flipped position faster, per generation, than the
designs with independent reads do, within a resolution of $0.001$ to $0.003$. A gain of a few
thousandths could have been missed, in particular in the first stages of the searches, which use
few runs. The searches are random, at one noise level apart from the spot check, for periods of
at most $8$ apart from the changes of $D_{16}$, and they are in the ideal model.

\emph{One fact behind the question, and what is missing.} Take two assignments of values to the
positions of an instance, and call a gate $k$ \emph{regular} if both satisfy
$x_k=\NOR(x_{k-1},x_{j_k})$; in two evaluations of the table--Jacobi kind the short gates are
regular, and so is every gate above the flipped layer in the experiment above. If the two
assignments differ at $p$, then position $p+1$, if its gate is regular, and every regular gate
that lists $p$ do not hold $1$ in both, and if gate $p$ is regular, neither $p-1$ nor the
position that it lists holds $1$
in both: a one at an input of a regular gate, in both assignments, forces a zero at its output in
both, and a one at its output needs zeros at both inputs. So a gate whose read is such a
neighbour of the position that caused a difference in its chain reads a zero or a differing
value, and never ends that difference by reading a one in both. In this way a base can make single
differences last longer than independent reads allow, and the count of Section~\ref{sec:allbases}
charges every gate that keeps its pointer as if it did. In the measurements above this gain is
more than offset. A plausible reason is that the differences of a chain come in runs, so that
gates listing related positions share one cause and produce one run where independent reads
would produce several. The referee's runs support this for neighbours: the mean length of a run
is $3.0$ for $\mathrm U$, $3.1$ for $\mathrm U,\mathrm R(1)$ and $4.4$ for $\mathrm F(0)$, and
the number of runs started per differing position of the layer below is $0.32$, $0.29$ and
$0.18$. For $\mathrm U,\mathrm R(2)$ the runs are not longer ($2.9$ and $0.31$), so this is not
the whole reason. We have no proof for the designs above, and none for all bases.
This is the point at which the level $0.455$ stops being a theorem. The argument of
Theorem~\ref{thm:indep} treats a layered base as a formula, in which no value is used twice, and
a general base is a circuit. For circuits of noisy gates the corresponding question is open as
well: Evans and Schulman~\cite{evansschulman} give an upper bound on the noise that circuits
tolerate, Unger~\cite{unger} determined the threshold for formulas of two-input gates, and that
the same threshold holds for circuits is his conjecture.

\emph{Fixed bases that are formulas.} One might hope to prove the level $0.455$ at least for
fixed bases in which no value is used twice. Such bases are nearly paths. Suppose that every
position depends on two earlier positions through pointers that are kept, and these on two each,
and so on. Following the pointers for $t$ levels gives $2^t$ sequences, so for $t>\log_2n$ two of
them end at the same position. But at noise $0.46$ a flipped position of the plain chain changes
$\lambda=0.972$ positions one level up, so the probability that a position is still open falls
by about this factor per level, and a union bound over the positions with an error of $1/n$
needs about $2\ln n/\ln(1/\lambda)\approx70\ln n$, that is $49\log_2n$, levels; at noise $0.455$
it needs $72\log_2n$. A tree of that depth with at most $n$ positions has a branching number
below $1.02$ (below $1.03$ if the union bound is allowed a constant error). In
Theorem~\ref{thm:layered} the
independence does not come from a tree. It comes from the coins of the base, which make every
read a new draw from a class of many copies, so that the chains above see only the densities of
the classes. In one sense that theorem is about fixed bases too: its probability $5/n$ is over
the coins and the smoothing, so by Markov's inequality all but a fraction $\sqrt{5/n}$ of the
outcomes of the coins give a fixed base on which the iteration finishes in time with probability
at least $1-\sqrt{5/n}$ over the smoothing. For a base that is not drawn in this way one would
have to show instead that values which are reused only some levels down are nearly independent,
with an error below the margin of the contraction. By a rough count for the plain chain at noise
$0.46$ (the influence of everything $d$ levels down is of order $\lambda^d$, and the margin is
$1-\lambda\approx0.03$) this asks for more than a hundred levels without reuse; for designs that
spread more it asks for many more, since the contraction that is proved has
$\theta=0.9998$. And it asks for a proof of the decay of correlations that we do not have.
\end{remark}

\subsection{From noise \texorpdfstring{$0.455$}{0.455} on, layered bases are easy}\label{sec:layered}

Theorem~\ref{thm:indep} is a statement about the ideal model. This subsection turns it into a
statement about smoothed instances. From noise $0.455$ on, the relevance-pruned iteration
finishes in $O(\log n)$ rounds on the smoothing of every \emph{layered base} with wide bundles and
long chains over a long preamble: a base built from bundles of at least $C_1\ln n$ chains, each
running a periodic program for many periods, whose gates read uniformly random copies of earlier
bundles, a constant, or the position before them (Theorem~\ref{thm:layered}). The bases of
Sections~\ref{sec:latch} and~\ref{sec:s24} are of this kind for every circuit when they are built
with long chains, many copies and a long preamble, and built in this way they still prove the
hardness of those sections (Lemma~\ref{lem:longer}). So, in this form, neither construction, nor
any other layered one, gives hardness at noise $0.455$ or above (Corollary~\ref{cor:layered}),
while Section~\ref{sec:s24} reaches $\eps_S>0.4468$: for this method the level lies in
$[0.4468,0.455]$. The width matters: with one chain per bundle the definition covers nearly every
base with a long preamble, and the statement is false there (Remark~\ref{rem:layered}(ii)).

Two facts make the transfer work. The positions that are still open after $s$ rounds of the
relevance-pruned iteration are exactly those where a lower and an upper trace differ, and these
two traces satisfy a crossed NOR recursion (Lemma~\ref{lem:lowup}); so the densities of the two
traces of a bundle follow the density map of a design with twice as many classes, and
Theorem~\ref{thm:indep}(a) contracts it. Once few positions are open,
inequality~\eqref{eq:Mpot} makes the expected number of open positions shrink by a fixed factor
per round.

\begin{definition}[layered bases]\label{def:layered}
Fix integers $T,h,m\ge1$ and $\rho\ge0$. A \emph{bundle} $\mathcal Q$ is given by a
\emph{program}, a list of $T_{\mathcal Q}\le T$ types, each a mixed type
(Definition~\ref{def:indep}) whose classes are classes of earlier bundles, or the type
$\mathrm L_1$; and by a nonempty set $\Lambda_{\mathcal Q}\subseteq\{0,\dots,T_{\mathcal Q}-1\}$ of
\emph{output phases}. The bundle consists of $m$ \emph{chains}, each preceded by
$\rho T_{\mathcal Q}h$ zero gates; a chain is a sequence of $T_{\mathcal Q}h$ consecutive gates,
and its gate of number $r$ has the type of the phase $(r-1)\bmod T_{\mathcal Q}$. A gate of mixed
type $\mu$ lists position $0$ with probability $\mu_0$, position $1$ with probability $\mu_1$, and a
uniformly random copy of class $c$ with probability $\mu_c$, with fresh coins; a gate of type
$\mathrm L_1$ lists the position before it. For $\varphi\in\Lambda_{\mathcal Q}$ the class
$(\mathcal Q,\varphi)$ has as its \emph{copies} the gates of phase $\varphi$ in the last period of
the $m$ chains. A \emph{layered base} consists of zero gates at the positions $1,\dots,n_0$,
followed by bundles $\mathcal Q_1,\dots,\mathcal Q_N$ in this order, where the classes listed by
$\mathcal Q_i$ belong to $\mathcal Q_1,\dots,\mathcal Q_{i-1}$.
\end{definition}

\noindent With $m=h=1$ every gate can be its own bundle, and a gate that lists a position $p$ then
reads the class of the bundle of $p$; so the definition alone allows nearly every base with a
long preamble. The theorem below needs $m\ge C_1\ln(n+1)$, so that a gate reads a uniformly random
copy among many, and $h\ge h_0$. With $\mathrm{pre}=n_0$, the base $\mathrm{Latch}(C)$ of
Definition~\ref{def:latch} is a layered base with $T=2$: every rail is a bundle with the program (set source, reset source) and
the output phase $1$, and the sources $\theta S$ and $\frac12S+\frac12S'$ and the constant rails are
mixed types. The base $\mathrm{Prog}(C)$ of Definition~\ref{def:prog} is one with $T=24$, the
program $S_{24}$ and the output phases $\mathcal L$.

\begin{theorem}[layered bases above $0.455$]\label{thm:layered}
Let $0.455\le\eps<1$ and $T\ge1$. There are $c'>0$ and $h_0$, $C_1$, $C_2$, $n_1$, depending only on
$\eps$ and $T$, such that for every $n\ge n_1$ the following holds for every layered base of length
$n$ with $h\ge h_0$, $m\ge C_1\ln(n+1)$ and $n_0\ge(1-c')n$: with probability at least $1-5/n$ over
the coins and the smoothing, $\max_iD(i)\le C_2\ln n$. So the relevance-pruned iteration finishes
within $C_2\ln n$ rounds, and the prefix--Jacobi iteration within $C_2\ln n+1$ applications of $G$;
the latter checks its output, so it is errorless.
\end{theorem}

\begin{lemma}[longer chains, more copies, longer preambles]\label{lem:longer}
\begin{enumerate}[label=(\alph*),leftmargin=2.4em]
\item If $(\theta,h,J,\eta,\gamma)$ is admissible in the sense of Definition~\ref{def:adm}, or
$(\beta,h,J,\eta,\gamma)$ in the sense of Definition~\ref{def:adms24}, then so is the same tuple with
any $h'\ge h$ in place of $h$, with the same numbers.
\item Theorems~\ref{thm:hard4} and~\ref{thm:hards24} remain true when, in
Definitions~\ref{def:latcha} and~\ref{def:proga}, $h$ is replaced by a fixed $h'\ge h$, $m$ by the
least power of two at least $K_mM(a+1)\log_2s$ for a fixed $K_m\ge1$, and $K_p(n-\mathrm{pre})$
zero gates, for a fixed $K_p\ge0$, are added to the preamble.
\end{enumerate}
\end{lemma}

\begin{proof}
(a) For the latch, $F^-_h=F(1-\tau^h)$ is nondecreasing and $F^+_h=F+(1-F)\tau^h$ nonincreasing in
$h$, so $\mathrm{lo}$ does not decrease and $\mathrm{up}$ does not increase, and (R), (G), (V) keep
holding. For $S_{24}$: a chain of $h'$ periods with the worst start is a chain of $h$ periods
started from the value after its first $h'-h$ periods, and $p_o$ is affine in the probability of
the start, so $P^-_c$ does not decrease and $P^+_c$ does not increase with $h$; so neither does
$\Psi^-$, nor increases $\Psi^+$, and (R), (G), (V) keep holding.
(b) Theorems~\ref{thm:latch} and~\ref{thm:progembed} hold for all $m$, $\rho$ and $\mathrm{pre}$,
and a larger $m$ only lowers the bound $e^{-2\gamma^2m}$. The length grows by at most the constant
factor $(1+K_p)K_m\cdot2$, so parts (i)--(iii) of Lemmas~\ref{lem:param4} and~\ref{lem:paramprog} hold
with $Z_0$ enlarged by that factor, and condition (ii) on the preamble only becomes easier. The proofs
of Corollaries~\ref{cor:latch} and~\ref{cor:prog} and of the two theorems are then unchanged.
\end{proof}

\begin{corollary}[the method stops at $0.455$]\label{cor:layered}
Let $0.455\le\eps<1$, and let $h_0$, $C_1$, $C_2$, $c'$ be as in Theorem~\ref{thm:layered} for
$T=24$. For every NOR circuit $C$, the relevance-pruned iteration finishes within $C_2\ln n$ rounds,
with probability at least $1-5/n$ over the coins and the smoothing, on the $\eps$-smoothing of
$\mathrm{Latch}(C)$ and of $\mathrm{Prog}(C)$ whenever these are built with $h\ge h_0$,
$m\ge C_1\ln(n+1)$ and a preamble of at least $(1-c')n$ zero gates. By Lemma~\ref{lem:longer}, the
bases built in this way still give the hardness of Theorems~\ref{thm:hard4} (below $4/9$)
and~\ref{thm:hards24} (up to $\eps_S$). So, in this form, neither embedding, nor any embedding by a
layered base with $m\ge C_1\ln(n+1)$ and $h\ge h_0$, proves hardness at a noise level of $0.455$ or
more: the prefix--Jacobi iteration, an errorless NC algorithm, solves the smoothed instances it
produces with high probability.
\end{corollary}

\begin{proof}
Theorem~\ref{thm:layered} with $T=24$, and the remarks after Definition~\ref{def:layered}; in
Lemma~\ref{lem:longer}(b) take $K_m$ with $K_mM\ge2C_1$, so that $m\ge C_1\ln(n+1)$ because
$n+1\le s^2$, and $K_p$ with $K_p/(1+K_p)\ge1-c'$.
\end{proof}

For the proof, for $s\ge0$ and a position $i$ let $\sigma^{(s)}_i=x_i$ if $D(i)\le s$ and
$\sigma^{(s)}_i=\bot$ otherwise, and let $L^{(s)}_i$ and $U^{(s)}_i$ be $\sigma^{(s)}_i$ with $\bot$
replaced by $0$ and by $1$. So $L^{(s)}\le x\le U^{(s)}$, and $i$ is $s$-open exactly when
$L^{(s)}_i<U^{(s)}_i$. Since $D(0)=0$ and $D(i)\ge1$ for $i\ge1$, $L^{(0)}_i=0$ and $U^{(0)}_i=1$ for
$i\ge1$.

\begin{lemma}[lower and upper traces]\label{lem:lowup}
Let $s\ge1$, let $k\ge1$ be a gate and $p=\jt_k$ its pointer. If gate $k$ is not a NOT gate, then
\[
L^{(s)}_k=\NOR\bigl(U^{(s)}_{k-1},U^{(s-1)}_p\bigr),\qquad U^{(s)}_k=\NOR\bigl(L^{(s)}_{k-1},L^{(s-1)}_p\bigr).
\]
If it is a NOT gate, the same holds with $U^{(s-1)}_p$ and $L^{(s-1)}_p$ replaced by $0$.
\end{lemma}

\begin{proof}
Gate $1$ is a NOT gate with $D(1)=1$, and both sides give $L^{(s)}_1=U^{(s)}_1=0$. Let $k\ge2$; we
use Lemma~\ref{lem:status}. First, $L^{(s)}_k=1$ means $x_k=1$ and $D(k)\le s$. Then $x_{k-1}=0$
and gate $k$ reads a zero: case (iii), where $D(k)\le s$ means $D(k-1)\le s$ and $D(p)\le s-1$,
that is $U^{(s)}_{k-1}=U^{(s-1)}_p=0$; or a NOT gate, case (iv), where it means
$U^{(s)}_{k-1}=0$. Conversely, $U^{(s)}_{k-1}=0$ gives $x_{k-1}=0$ and $D(k-1)\le s$, and for a gate
that is not a NOT gate $U^{(s-1)}_p=0$ gives $x_p=0$ and $D(p)\le s-1$; so we are in these cases
with $D(k)\le s$. Second, $U^{(s)}_k=0$ means $x_k=0$ and $D(k)\le s$. If $x_{k-1}=0$, then $k$ is
a relevant reset, case (i), and $D(k)\le s$ means $D(p)\le s-1$, that is $L^{(s-1)}_p=1$; here
$L^{(s)}_{k-1}=0$. If $x_{k-1}=1$ and $k$ is a reset, case (ii) gives
$D(k)=\min(D(k-1),1+D(p))$, so $D(k)\le s$ means $L^{(s)}_{k-1}=1$ or $L^{(s-1)}_p=1$. If
$x_{k-1}=1$ and $k$ reads a zero or is a NOT gate, case (iv) gives $D(k)=D(k-1)$, and
$L^{(s-1)}_p=0$ for a gate that is not a NOT gate, as $x_p=0$. In all cases
$U^{(s)}_k=0$ exactly when $L^{(s)}_{k-1}=1$ or (for a gate that is not a NOT gate)
$L^{(s-1)}_p=1$.
\end{proof}

\noindent A NOT gate thus acts like a gate that reads a determined zero, and a gate of type
$\mathrm L_1$ that keeps its pointer is a NOT gate. So below we treat $\mathrm L_1$ as the mixed
type with $\mu_1=1$; in the ideal model the two have the same density map
(Lemma~\ref{lem:onebit}(a)). Applying the lemma repeatedly along a chain with gates
$g_1<\dots<g_N$ and $g_0=g_1-1$: for $o\le N$, $L^{(s)}_{g_o}$ is the value $x_o$ of
Lemma~\ref{lem:progchain} with the start $L^{(s)}_{g_0}$ if $o$ is even and $U^{(s)}_{g_0}$ if $o$
is odd, and with the read of $g_r$ equal to the $L^{(s-1)}$-value of its target if $o-r$ is odd and
to the $U^{(s-1)}$-value if $o-r$ is even; for $U^{(s)}_{g_o}$ exchange $L$ and $U$ everywhere.
These are exactly the choices that make $x_o$ smallest and largest
(Lemma~\ref{lem:progchain}(a)).

\emph{The doubled design.} For a bundle $\mathcal Q$, a class $(\mathcal Q,k)$ and a sign
$\sigma\in\{-,+\}$, let the \emph{doubled rail} $(k,\sigma)$ be the following design of period
$2T_{\mathcal Q}$, over the classes $(c,-)$ and $(c,+)$, $c$ a class listed by $\mathcal Q$: its
gate of phase $\varphi$ has the type of phase $\varphi\bmod T_{\mathcal Q}$ of the program, with
every class $c$ replaced by $(c,\sigma)$ if $k-\varphi$ is odd and by $(c,-\sigma)$ if $k-\varphi$ is
even. Let $\Phi^{(2)}_{\mathcal Q}$ be the map that sends a vector $v$ of densities of the classes
$(c,\pm)$ to the vector of the densities, at phase $k$, of the doubled rails $(k,\sigma)$. It is
the density map of a design with independent reads, restricted to some classes. If
$v_{(c,-)}=v_{(c,+)}=d_c$ for every $c$, then $\Phi^{(2)}_{\mathcal Q}(v)_{(k,\pm)}=\Phi_{\mathcal Q}(d)_k$,
where $\Phi_{\mathcal Q}$ is the density map of the program. And $\Phi^{(2)}_{\mathcal Q}$ is
monotone for the order in which $v\sqsubseteq v'$ means $v_{(c,-)}\ge v'_{(c,-)}$ and
$v_{(c,+)}\le v'_{(c,+)}$ for all $c$ (a narrower vector): by Lemma~\ref{lem:progchain}(b), in the
rail $(k,-)$ the classes $(c,-)$ are read at gates where the output rises with the read and the
classes $(c,+)$ where it falls, and the other way round in $(k,+)$.

\begin{proof}[Proof of Theorem~\ref{thm:layered}]
Write $A=\eps\varpi$, $B=\eps(1-\varpi)$, and let $\theta=0.9998$, $f$, $lo$, $hi$, $\mathcal K$ and
$\delta$ be as in Theorem~\ref{thm:indep}; let $m_f$ and $M_f$ be the minimum and the maximum of
$f$ on $[lo,hi]$. By Theorem~\ref{thm:indep}(a) and the proof of its part (d), every density map
$\Phi$ of a design with independent reads, between two sets of classes, maps $[0,1]^{\mathcal C'}$
into $\mathcal K$ and satisfies $\delta(\Phi(d),\Phi(d'))\le\theta\,\delta(d,d')$ on $\mathcal K$. On
$\mathcal K$, $\|d-d'\|_\infty/M_f\le\delta(d,d')\le\|d-d'\|_\infty/m_f$. If the classes are the
union of several sets, the distance $\delta$ on the union is at most the largest of the distances
on the parts: reparametrize nearly optimal paths in the parts proportionally to their lengths, so
that each has constant speed on $[0,1]$ (the infimum defining $\delta$ is the same over Lipschitz
paths, by approximation with piecewise linear ones). Every density map is Lipschitz in the maximum norm with a
constant $L_\Phi$ that depends only on $\eps$, by Lemma~\ref{lem:onebit}(d), because the rows of
$\overline M$ add up to at most $(1-\eps)/(1-(1-A))$. Let $\pi$ denote the coordinatewise
projection onto $\mathcal K$.

\emph{Step 0: the prefix and the density.} The first $n_0$ positions form the smoothing of the base
of length $n_0$ of zero gates, and their trace and depths do not depend on the rest. Step~1 of the
proof of Theorem~\ref{thm:blockabove} applies to them with no blocks ($t_1=0$): every pointer that
is kept there is $0$, so a skeleton with all positions at least $g_0$ takes resampled steps only.
So there is $s_z\le C_z\ln n$ such that, outside an event of probability at most $2/n$, every
position up to $n_0$ is $s_z$-determined; call this event $\mathcal Z$. By
Lemma~\ref{lem:dens2any} applied to that base with $K=3$, for any $\eta_0>0$ and large $n$, with
probability at least $1-(n_0+1)^{-3}\ge1-(n+1)^{-2}$ we have $|\psi_g-\varpi|\le\eta_0$ for all $g\in[n_0/2,n_0]$; then, as
$(g-n_0)/g\le c'/(1-c')\le2c'$ for $g\le n$ and $c'\le1/2$, $|\psi_g-\varpi|\le\eta_0+3c'$ for all
$g\in[n_0/2,n]$. Call this event $\mathcal E$, and put $\eta=\eta_0+3c'$ and $\eta'=\eta+2c'$. On
$\mathcal Z$, for a round $s>s_z$ and a gate $g>n_0$, the fractions $\psi^L_{g-1}$ and
$\psi^U_{g-1}$ of positions $p<g$ with $L^{(s-1)}_p=1$ and with $U^{(s-1)}_p=1$ differ from
$\psi_{g-1}$ by at most $(g-n_0)/g\le2c'$; so on $\mathcal E\cap\mathcal Z$ they lie in
$[\varpi-\eta',\varpi+\eta']$. The events $\mathcal Z$ and $\mathcal E$ depend only on the noise of
the first $n_0$ gates, so they are independent of the coins and of the noise of the bundles. From
now on we work conditionally on $\mathcal E\cap\mathcal Z$; this changes none of the conditional
laws used below.

\emph{Step 1: one bundle in one round.} Let $\mathcal Q$ be a bundle, $s>s_z$, and let
$\mathcal G_{\mathcal Q}$ be generated by the coins of the bundles below $\mathcal Q$ and by
$(\pi_p,\rho_p,U_p)$ for the positions $p$ below the first position of $\mathcal Q$. For a gate $g$ of
$\mathcal Q$ let $\mathcal H_{g-1}$ be generated by $\mathcal G_{\mathcal Q}$, the coins of the gates of
$\mathcal Q$ below $g$, and $(\pi_p,\rho_p,U_p)$ for all positions $p<g$. The values $L^{(s-1)}$ and $U^{(s-1)}$ of positions below
$\mathcal Q$ are $\mathcal G_{\mathcal Q}$-measurable. For a listed class $c$ let $\lambda_c$ and
$\upsilon_c$ be the fractions of its copies with $L^{(s-1)}=1$ and with $U^{(s-1)}=1$, and let
$v$ be the vector with $v_{(c,-)}=\lambda_c$, $v_{(c,+)}=\upsilon_c$. For a chain gate $g$ of
$\mathcal Q$ of mixed type $\mu$, the $\sigma^{(s-1)}$-value $R_g$ of its read has, given
$\mathcal H_{g-1}$, the probabilities
\[
\PP[R_g=1]=(1-\eps)\Bigl(\mu_0+\sum_c\mu_c\lambda_c\Bigr)+\eps\psi^L_{g-1},\qquad
\PP[R_g\ne0]=(1-\eps)\Bigl(\mu_0+\sum_c\mu_c\upsilon_c\Bigr)+\eps\psi^U_{g-1},
\]
both $\mathcal H_{g-1}$-measurable up to an error of at most $\eps/g$; position $1$ is
$(s-1)$-determined because $s-1\ge1$. Here, if the gate is resampled to $g-1$ it becomes a NOT
gate, and by Lemma~\ref{lem:lowup} its effective read is a determined zero; we take $R_g$ to be this
effective read, which changes the probabilities by at most $\eps/g$, and we absorb this into
$\eta'$ (as $g>n_0\ge n/2$). As in
Lemma~\ref{lem:reads}, a uniform variable $\Xi_g$, independent of $\mathcal H_{g-1}$, can be built from
$R_g$ and $U_g$ with $R_g=1$ if and only if $\Xi_g<\PP[R_g=1]$ and $R_g\ne0$ if and only if
$\Xi_g<\PP[R_g\ne0]$, almost surely; the $\Xi_g$ of the gates of $\mathcal Q$ are independent and
independent of $\mathcal G_{\mathcal Q}$. Put $y^-_g=\mathbf 1[\Xi_g<(1-\eps)(\mu_0+\sum_c\mu_c\lambda_c)+\eps(\varpi-\eta')]$
and $y^+_g=\mathbf 1[\Xi_g<(1-\eps)(\mu_0+\sum_c\mu_c\upsilon_c)+\eps(\varpi+\eta')]$. On
$\mathcal E\cap\mathcal Z$ the $L^{(s-1)}$-value of the read of $g$ is at least $y^-_g$ and its
$U^{(s-1)}$-value at most $y^+_g$. For $k\in\Lambda_{\mathcal Q}$ and a chain, define
$X^-(k)$ and $X^+(k)$ as in the proof of Theorem~\ref{thm:progembed}, with these $y^\pm_g$ and the
worst start. By the remark after Lemma~\ref{lem:lowup}, on $\mathcal E\cap\mathcal Z$ the copy of
phase $k$ of the chain has $L^{(s)}\ge X^-(k)$ and $U^{(s)}\le X^+(k)$. Given $\mathcal G_{\mathcal Q}$,
the variables $X^-(k)$ of the $m$ chains are independent, and so are the $X^+(k)$, with
$\PP[X^\pm(k)=1\mid\mathcal G_{\mathcal Q}]=P^\pm_k(v)$, where $P^\pm_k(v)$ is $p_o$ of
Lemma~\ref{lem:progchain}(b) for these read probabilities. Compare with the stationary chain of the
doubled rail $(k,\pm)$, which reads the same classes with resampled reads exactly $\varpi$: every
read probability differs by at most $\eps\eta'$, the map $p\mapsto(1-p)(1-\ell)$ has slope at most
$1-A/2$ when $\eta'\le\varpi/2$, and the chain has $T_{\mathcal Q}(h-1)$ gates before the last
period. So
\begin{equation}\label{eq:Pphi}
\bigl|P^\pm_k(v)-\Phi^{(2)}_{\mathcal Q}(v)_{(k,\pm)}\bigr|\ \le\ \frac{2\eps\eta'}{A}+(1-A/2)^{h-1}=:\kappa_0 .
\end{equation}
The same argument with the true reads (Lemma~\ref{lem:reads}, which needs only $\mathcal E$) shows
that, given $\mathcal G_{\mathcal Q}$, the fraction $t_{\mathcal Q,k}$ of copies of class
$(\mathcal Q,k)$ that hold $1$ lies between the means of $m$ independent indicators with means
$P^-_k(t,t)$ and $P^+_k(t,t)$, where $t$ is the vector of the true densities of the listed classes
and both means are within $\kappa_0$ of $\Phi_{\mathcal Q}(t)_k$. By Hoeffding's inequality, for
$\gamma>0$, outside an event of probability at most $4Te^{-2\gamma^2m}$ given
$\mathcal G_{\mathcal Q}$, on $\mathcal E\cap\mathcal Z$ and for all $k\in\Lambda_{\mathcal Q}$:
\begin{equation}\label{eq:onebundle}
\lambda^{(s)}_{\mathcal Q,k}\ge\Phi^{(2)}_{\mathcal Q}(v)_{(k,-)}-\kappa_1,\quad
\upsilon^{(s)}_{\mathcal Q,k}\le\Phi^{(2)}_{\mathcal Q}(v)_{(k,+)}+\kappa_1,\quad
\bigl|t_{\mathcal Q,k}-\Phi_{\mathcal Q}(t)_k\bigr|\le\kappa_1,
\end{equation}
with $\kappa_1=\kappa_0+\gamma$, where $\lambda^{(s)}_{\mathcal Q,k}$ and $\upsilon^{(s)}_{\mathcal Q,k}$ are
the fractions of copies of $(\mathcal Q,k)$ with $L^{(s)}=1$ and $U^{(s)}=1$. Let $\mathcal F$ be the
event that~\eqref{eq:onebundle} holds for every bundle and every round $s_z<s\le s_2$, where
$s_2\le n$ is fixed below, and $\mathcal F_{\mathcal Q}\in\mathcal G_{\mathcal Q}$ the same for the
bundles below $\mathcal Q$. There are at most $n^2$ pairs of a bundle and a round, so for
$m\ge C_1\ln(n+1)$ with $C_1$ large, $\PP[\mathcal F^c\mid\mathcal E\cap\mathcal Z]\le(n+1)^{-3}$.

\emph{Step 2: the open fractions become small.} The ideal densities $d$ are defined along the
layout: positions $0$ and $1$ have densities $1$ and $0$, and $d_{\mathcal Q}=\Phi_{\mathcal Q}(d)$
for every bundle, applied to the ideal densities of the classes it lists. On $\mathcal F$, by
induction along the layout, $\delta(\pi t_{\mathcal Q},d_{\mathcal Q})\le\theta\max\delta(\pi t,d)+2(1+L_\Phi)\kappa_1/m_f$,
the maximum over the bundles listed by $\mathcal Q$, so
$\|t_{\mathcal Q}-d_{\mathcal Q}\|_\infty\le\kappa_2:=4M_f(1+L_\Phi)\kappa_1/(m_f(1-\theta))$ for every
bundle. Next define deterministic interval vectors $V^{(s)}_{\mathcal Q}$, with entries
$V_{(k,-)}\le V_{(k,+)}$: for $s=s_z$ all entries are $0$ and $1$, and
$V^{(s)}_{\mathcal Q}$ is $\Phi^{(2)}_{\mathcal Q}(V^{(s-1)})$ widened by $\kappa_1$ (subtract $\kappa_1$
from the entries $(k,-)$, add it to the entries $(k,+)$, clip to $[0,1]$), where
$V^{(s-1)}$ collects the vectors of the listed bundles, and the constant rails get the exact
values. On $\mathcal F$, by~\eqref{eq:onebundle} and the monotonicity of $\Phi^{(2)}$ for
$\sqsubseteq$, induction on $s$ gives
$\lambda^{(s)}_{\mathcal Q,k}\ge V^{(s)}_{(k,-)}$ and $\upsilon^{(s)}_{\mathcal Q,k}\le V^{(s)}_{(k,+)}$
for all $s_z\le s\le s_2$. Let $W_{\mathcal Q}$ be the vector with both entries $(k,\pm)$ equal to
$t_{\mathcal Q,k}$. By~\eqref{eq:onebundle} and $\Phi^{(2)}(W)=(\Phi(t),\Phi(t))$, both
$V^{(s)}_{\mathcal Q}$ and $W_{\mathcal Q}$ are within $\kappa_1$ of the images under
$\Phi^{(2)}_{\mathcal Q}$ of $V^{(s-1)}$ and $W$, and these images are within $L_\Phi\kappa_1$ of the
images of $\pi V^{(s-1)}$ and $\pi W$. With the contraction, $\Delta_s=\max_{\mathcal Q}\delta(\pi V^{(s)}_{\mathcal Q},\pi W_{\mathcal Q})$
satisfies $\Delta_s\le\theta\Delta_{s-1}+2(1+L_\Phi)\kappa_1/m_f$ for $s\ge s_z+2$, and
$\Delta_{s_z+1}\le(hi-lo)/m_f+2\kappa_1/m_f$. So $\Delta_s\le\theta^{s-s_z-1}(hi-lo+2)/m_f+\kappa_2/(2M_f)$.
Every $s$-open copy has $L^{(s)}<U^{(s)}$, so the open fraction of $(\mathcal Q,k)$ is
$\upsilon^{(s)}_{\mathcal Q,k}-\lambda^{(s)}_{\mathcal Q,k}\le V_{(k,+)}-V_{(k,-)}\le2M_f\Delta_s+4\kappa_1$,
since $V$ and $W$ are within $\kappa_1$ of $\pi V$ and $\pi W$; the same bound holds for
$t-\lambda^{(s)}$ and $\upsilon^{(s)}-t$.
Given $\kappa_3>0$, if $\kappa_1$ is small enough there is a constant $k_1$ with: on $\mathcal F$,
for every $s_1:=s_z+k_1\le s\le s_2$, every bundle and class has
$\lambda^{(s)}\ge t-\kappa_3$, $\upsilon^{(s)}\le t+\kappa_3$, and $|t-d|\le\kappa_3$; for the classes
listed by a bundle $\mathcal Q$ this holds on $\mathcal F_{\mathcal Q}$.

\emph{Step 3: the expected numbers of open positions shrink.} Let $s_1\le s<s_2$. A position
$i>n_0$ that is $(s+1)$-open lies in a chain or in the zero gates before one. In the chain,
let $g$ be the largest position up to $i$ such that $g$ is open and $g-1$ is not; it exists
because the positions up to $n_0$ are determined. If $g$ lies before the first gate of the chain,
then the gate before the chain, the \emph{start}, is open, and so is every gate of the chain up
to $i$. By Lemma~\ref{lem:lowup}, an open
gate whose predecessor is determined reads an $s$-open position and its predecessor is not a
determined one; and an open gate whose predecessor is open reads a position that is not a
determined one, in round $s$. Zero gates are the mixed type $\mu_0=1$, so the same holds before the
chain. Hence $i$ is open only if, for some $g\le i$: $g$ reads an $s$-open position and $x_{g-1}$
is not a determined one, or $g$ is the start and is open; and every gate of $(g,i]$ reads a
position that is not $s$-determined with value $1$.

Now bound probabilities. Fix $\kappa>0$ and let $\overline M^\kappa_{k,c}(d)$ be defined as
$\overline M_{k,c}(d)$ in Lemma~\ref{lem:onebit}(d), with every $\alpha_{g-1}$ replaced by
$\alpha_{g-1}+\kappa$ and every $q_i$ by $q_i+\kappa$. As $\kappa\to0$, $\overline M^\kappa\to\overline M$
uniformly in $d$ and in the program, because $q_i\le1-A$ and the sums are geometric; so
by~\eqref{eq:Mpot}, for small $\kappa$,
\begin{equation}\label{eq:Mpotkappa}
\sum_c\overline M^\kappa_{k,c}(d)\,f(1-d_c)\ \le\ \theta_1f\bigl(1-\Phi(d)_k\bigr),\qquad\theta_1:=\tfrac{1+\theta}2<1 .
\end{equation}
Fix a constant $J_0$ with $(1-A/2+\kappa)^{J_0}\le\kappa/4$. On $\mathcal F_{\mathcal Q}$, given
$\mathcal H_{g-1}$, a chain gate $g$ of $\mathcal Q$ reads a position that is not an $s$-determined
one with probability $1-(1-\eps)(\mu_0+\sum_c\mu_c\lambda_c)-\eps\psi^L_{g-1}\le q_g(d)+2\kappa_3+\eta'$,
by Step 2 and Step 0, where $q_g(d)$ is the number of Definition~\ref{def:indep}; this is at most
$q_g(d)+\kappa$ if $\kappa_3$, $\eta_0$ and $c'$ are small. It reads an $s$-open position with
probability at most $(1-\eps)\sum_c\mu_c(g)o_c+\eps\Omega_{g-1}$, where $o_c=\upsilon_c-\lambda_c$ is the
open fraction of class $c$ in round $s$, which is $\mathcal G_{\mathcal Q}$-measurable, and
$\Omega_{g-1}$ is the fraction of $s$-open positions below $g$. If $g$ is at least $J_0$ gates above the start of its
chain, then $\PP[x_{g-1}\text{ is not a determined one in round }s+1\mid\mathcal G_{\mathcal Q}]\le\alpha_{g-1}(d)+\kappa$
on $\mathcal F_{\mathcal Q}$: the probability that $x_{g-1}=0$ is within $(1-A/2)^{J_0}+K\kappa_3+K\eta'$
of $\alpha_{g-1}(d)$: by the sandwich of Step 1 applied to the true reads (Lemma~\ref{lem:reads}),
$x_{g-1}$ lies between the values of two chains with independent reads whose probabilities are
within $2\kappa_3+\eta'$ of those at $d$, and Lemma~\ref{lem:onebit}(a),(b) applies to them; and by the cascade argument below,
the probability that $g-1$ is open is at most $K(\kappa_3+c')+(1-A/2+\kappa)^{J_0}$. Here $K$
depends only on $\eps$. Now fix a copy $i$ of $(\mathcal Q,k)$ and $g\le i$ in its chain.
Conditioning successively on $\mathcal H_{i-1},\dots,\mathcal H_g$, as in
Lemma~\ref{lem:onechain}, and then on $\mathcal H_{g-1}$ and $\mathcal G_{\mathcal Q}$, the
probability that $\mathcal F_{\mathcal Q}$ holds, that $g$ reads an open position, that $x_{g-1}$ is
not a determined one, and that every gate of $(g,i]$ reads a position that is not a determined one,
is at most
\[
\Bigl((1-\eps)\sum_c\mu_c(g)\,\EE[o_c;\mathcal F_{\mathcal Q}]\,(\alpha_{g-1}(d)+\kappa)+2\eps c'e^*_s\Bigr)\prod_{g<r\le i}(q_r(d)+\kappa),
\]
for $g$ at least $J_0$ gates above the start; here
$e^*_s$ is the largest probability that a position above $n_0$ is $s$-open, and we used that
positions up to $n_0$ are $s$-determined, so $\EE[\Omega_{g-1}]\le2c'e^*_s$. The complement of
$\mathcal F_{\mathcal Q}$ is counted once, below. For the at most $J_0$ gates
next to the start, and for the start itself, which is open in round $s+1$ with probability at most
$e^*_{s+1}$, replace $\alpha_{g-1}(d)+\kappa$ by $1$; their products have at least
$T_{\mathcal Q}(h-1)-J_0$ factors. Adding over $g$, and using
$\EE[o_c;\mathcal F_{\mathcal Q}]\le e_s(c)$, the mean over the copies $p$ of $c$ of
$\PP[p\text{ is }s\text{-open}]$:
\[
\PP[i\text{ is }(s+1)\text{-open}]\ \le\ \sum_c\overline M^\kappa_{k,c}(d)\,e_s(c)+K_0\,\zeta\,(e^*_s+e^*_{s+1})+K_0\,\beta,
\]
because the stationary sums defining $\overline M^\kappa$ dominate the sums over the gates of the
chain. Here $\zeta=2c'+(1-A/2+\kappa)^{h-1-J_0}$, $\beta=(n+1)^{-3}$ bounds
$\PP[\mathcal F^c_{\mathcal Q}]$, and $K_0$ depends only on $\eps$ and $T$. In the same way, every
position above $n_0$, also one that is not a copy, is $(s+1)$-open with probability at most
$K_0\max_ce_s(c)+K_0\zeta(e^*_s+e^*_{s+1})+K_0\beta$; for the zero gates before a chain, and for
the gates near its start, this uses only that the factors of the cascade are at most $1-A/2+\kappa$.
Taking the maximum over the positions, for $K_0\zeta\le1/2$,
\[
e^*_{s+1}\ \le\ 2K_0\max_ce_s(c)+2K_0\zeta e^*_s+2K_0\beta .
\]
Put $w_c=f(1-d_c)$ and $R_s=\max_ce_s(c)/w_c$, over all classes. By~\eqref{eq:Mpotkappa} and the last
two displays,
\[
R_{s+1}\le\theta_1R_s+\frac{K_0\zeta}{m_f}(e^*_s+e^*_{s+1})+\frac{K_0\beta}{m_f},\qquad
e^*_{s+1}\le2K_0M_fR_s+2K_0\zeta e^*_s+2K_0\beta .
\]
Insert the second inequality into the first. For $\zeta\le\zeta_0$, where $\zeta_0$ depends only on
$\eps$ and $T$, the resulting linear recursion for $(R_s,e^*_s)$ has a matrix with nonnegative
entries and spectral radius at most $\theta_2:=(1+\theta_1)/2<1$; choose $c'$ small and $h_0$ large
so that $\zeta\le\zeta_0$. Then, from $R_{s_1}\le1/m_f$ and $e^*_{s_1}\le1$,
$e^*_s\le K_1\theta_2^{s-s_1}+K_1\beta$ for $s_1\le s\le s_2$, with $K_1$ depending only on $\eps$
and $T$. Let $s_2=s_1+\lceil\ln(K_1n^3)/\ln(1/\theta_2)\rceil$. Then
$e^*_{s_2}\le n^{-3}+K_1(n+1)^{-3}$, and the expected number of $s_2$-open positions above $n_0$ is
at most $n\,e^*_{s_2}\le1/n$ for large $n$.

\emph{Conclusion.} All probabilities in Steps 1--3 are conditional on $\mathcal E\cap\mathcal Z$.
Outside $\mathcal Z^c$ (probability at most $2/n$), $\mathcal E^c$ (at most $(n+1)^{-2}$) and, given
$\mathcal E\cap\mathcal Z$, the event that some position above $n_0$ is $s_2$-open (at most $1/n$),
every position is $s_2$-determined, so $\max_iD(i)\le s_2\le C_2\ln n$, as $s_z\le C_z\ln n$ and
$k_1$ is a constant.
The last sentence of the theorem is Lemma~\ref{fact:depths} and Proposition~\ref{prop:jacobi}.
\end{proof}

\begin{remark}[what this says, and what it does not]\label{rem:layered}
\begin{enumerate}[label=(\roman*),leftmargin=2.4em]
\item \emph{The method.} The hardness proofs by embedding in this paper build layered bases:
chains over stages, reading random copies. Built with long chains, many copies and a long
preamble, they still prove hardness below $4/9$ and up to $\eps_S$ (Lemma~\ref{lem:longer}), and
Corollary~\ref{cor:layered} shows that from noise $0.455$ on the smoothed instances they produce
are solved by an errorless NC algorithm in $O(\log n)$ rounds. So the level of this method, in this
form, is in $[\eps_S,0.455]\subset[0.4468,0.455]$; numerically it is near $0.447$ for the designs
found (Remark~\ref{rem:beyond}). For the bases of Definitions~\ref{def:latcha}
and~\ref{def:proga} as written, with a preamble of logarithmic length, the theorem does not apply,
since its proof uses the long preamble; with heavy zero padding inside the bundles the same
argument should go through, but we have not written it.
\item \emph{Wide bundles, not all bases.} The theorem needs bundles of at least $C_1\ln n$ chains,
so that a gate reads a uniformly random copy among many. With one chain per bundle the definition
covers nearly every base with a long preamble, the mixed base of Remark~\ref{rem:mixed} among
them, and there the relevance-pruned iteration is slow up to noise about $0.55$ in simulation; the
referee of this subsection found $\max_iD(i)=159$, $276$, $509$ for $49$, $99$, $199$ layers of
$2000$ positions at noise $0.50$, over a zero preamble of $99\%$ of the positions. A base fixed by
an adversary can also aim many pointers at one position or read its own recent positions.
On every base the relevance-pruned iteration is fast above $\eps_1<0.7288$
(Theorem~\ref{thm:allbases}). Whether some NC algorithm works on every base at a level near
$0.455$ is open. The candidate we know is the table--Jacobi iteration of Section~\ref{sec:table} with reach of order $\log n$. On
mixed bases whose even gates read the position $H$ before them ($H$ even, layers of $4000$
positions in the last twentieth, one smoothing each), with reach $H-1$, so that these pointers are
just too long for it, it needed $102$ and $336$ rounds for $31$ and $125$ layers at noise $0.46$
when $H=6$, but $71$ and $94$ when $H=10$ and $54$ and $79$ when $H=16$; at $0.50$ the counts were
$52$ and $113$, $43$ and $70$, $57$ and $53$. So pointers just beyond its reach slow it down, and the
effect fades as they get longer. A search for designs that defeat it when its reach covers their
short pointers found none (script \texttt{design\_search.c}; numerical, and of little power). The
designs have layers of $600$ positions over a zero preamble of half the positions and a period of
at most $8$; each offset class has a pointer of a fixed reach at most $6$, or reads a random
position of a given class in the layer below, or reads at a fixed distance into the layer below,
or has pointer $0$; the iteration has reach $6$. Among $6000$ random designs at noise $0.50$ and
$6000$ at $0.55$, and in six hill-climbing runs at $0.50$ (period up to $12$, reach up to $8$),
the slowest designs had far reads only, or nearly so; with random far reads these are block bases
with offset classes. In the ideal model, no design whose far reads depend on each other spread a
flipped position faster than the designs with independent reads, within the resolution of the
measurements (Remark~\ref{rem:dependent}).
With $160$, $320$, $640$ and $1280$ layers the slowest of them reached the trace after about
$70$, $90$, $100$ and $130$ iterates at noise $0.50$ and $50$, $60$, $65$ and $80$ at $0.55$
(means of three smoothings); these sizes do not separate logarithmic growth from a small power. The speed of this
iteration is not a matter of short derivations: on the mixed base, derivations in which pointers
of reach $2$ are free are still long, and yet the iteration with reach $2$ is fast
(Remark~\ref{rem:setval}). Section~\ref{sec:setval} gives a sound variant with the same speed
there.
\item \emph{Constants.} The proof uses the contraction factor $\theta=0.9998$ of
Theorem~\ref{thm:indep}, so the constant $C_2$ is of order $10^4$ times the constants of
Lemma~\ref{lem:dens2any}; the real rate is the flip spread $\rho_u$ of the design, about $0.975$
for $S_{24}$ at noise $0.46$. In simulation (script \texttt{s24\_embed.c} with the variable
\texttt{DEPTH}), on $\mathrm{Prog}(C)$ for random circuits with $h=1$, a short preamble and one
smoothing each, $\max_iD(i)$ grows like the number of stages at noise $0.40$ ($77$, $135$, $260$ for
$J=16,32,64$, with $4$ gates, $\rho=10$, $m=500$), and slowly at $0.50$ and $0.55$ ($86$, $119$,
$129$, $173$ and $60$, $82$, $90$, $103$ for $J=64,128,256,512$, with $2$ gates, $\rho=5$, $m=200$),
while the length grows eightfold. At $0.46$ it still grows with $J$ at the sizes we ran ($71$, $99$,
$216$ for $J=16,32,64$), as expected when the contraction per stage is only about $0.975$.
\end{enumerate}
\end{remark}

\subsection{Layered bases with pointers two back}\label{sec:twolayered}

Theorem~\ref{thm:layered} is about the relevance-pruned iteration, and the programs of its
bundles have no gate that lists the position two before it. With such gates that iteration is
slow in simulation (Remark~\ref{rem:mixed}, and Remark~\ref{rem:twolayered} below), although
the ideal model contracts (Theorem~\ref{thm:two}). This subsection closes the gap between the
two for the table iteration. From noise $0.455$ on, the set-valued table iteration and the
table--Jacobi iteration with reach $2$ finish in $O(\log n)$ rounds on the smoothing of every
layered base with wide bundles whose programs may contain gates of type $\mathrm L_2$
(Theorem~\ref{thm:twolayered}).

The proof of Theorem~\ref{thm:layered} rests on a lower and an upper trace. A gate with pointer
two back seems to destroy them, because its value is not monotone along the chain. Two
identities restore them. Let $z$ be a sequence of bits with $z_k=\NOR(z_{k-1},b_k)$ and
$z_{k-1}=\NOR(z_{k-2},b_{k-1})$. If $b_k=z_{k-2}$, then
\begin{equation}\label{eq:idtwo}
z_k=(1-z_{k-2})\,b_{k-1},
\end{equation}
because $z_k=(1-z_{k-1})(1-z_{k-2})$ and $1-z_{k-1}$ is $1$ exactly if $z_{k-2}=1$ or
$b_{k-1}=1$. If moreover $b_{k-1}=z_{k-3}$ and $z_{k-2}=\NOR(z_{k-3},b_{k-2})$, then
\begin{equation}\label{eq:idthree}
z_k=z_{k-3},
\end{equation}
because $z_k=(1-z_{k-2})z_{k-3}$ by~\eqref{eq:idtwo}, and $z_{k-3}=1$ forces $z_{k-2}=0$. In
both identities the value at $k$ is a monotone or antitone function of one earlier value of the
chain and of at most one read, and that read is the read of the gate before.

\begin{definition}[traces for pointers two back]\label{def:twotraces}
Let $j$ be a parent array of length $n$. Put $\ell^{(s)}_0=u^{(s)}_0=1$ for all $s\ge0$, and
$\ell^{(0)}_i=0$, $u^{(0)}_i=1$ for $i\ge1$. For $s\ge1$ and $k=1,\dots,n$, with $p=j_k$, let
\[
\ell'=\NOR\bigl(u^{(s)}_{k-1},u^{(s-1)}_p\bigr),\qquad u'=\NOR\bigl(\ell^{(s)}_{k-1},\ell^{(s-1)}_p\bigr),
\]
with $u^{(s-1)}_p$ and $\ell^{(s-1)}_p$ replaced by $0$ if $p=k-1$. If $p\ne k-2$, put
$\ell^{(s)}_k=\ell'$ and $u^{(s)}_k=u'$. If $p=k-2$, let $q=j_{k-1}$ and put
$\ell^{(s)}_k=\max(\ell',\ell'')$ and $u^{(s)}_k=\min(u',u'')$, where
\begin{itemize}[leftmargin=2em]
\item $\ell''=u''=0$ if $q=k-2$;
\item $\ell''=\bigl(1-u^{(s)}_{k-2}\bigr)\,\ell^{(s-1)}_q$ and
$u''=\bigl(1-\ell^{(s)}_{k-2}\bigr)\,u^{(s-1)}_q$ if $q<k-3$;
\item $\ell''=\max\bigl\{\ell^{(s)}_{k-3},\ (1-u^{(s)}_{k-2})\,\ell^{(s-1)}_{k-3}\bigr\}$ and
$u''=\min\bigl\{u^{(s)}_{k-3},\ (1-\ell^{(s)}_{k-2})\,u^{(s-1)}_{k-3}\bigr\}$ if $q=k-3$.
\end{itemize}
Position $i$ is \emph{$s$-settled} if $\ell^{(s)}_i=u^{(s)}_i$.
\end{definition}

\noindent The pair $(\ell',u')$ is the rule of Lemma~\ref{lem:lowup}. The pair $(\ell'',u'')$
uses~\eqref{eq:idtwo} with the read of gate $k-1$ taken from the round before,
\eqref{eq:idthree} when gate $k-1$ also has its pointer two back, and that $x_k=0$ when gate
$k-1$ is a NOT gate.

\begin{lemma}[traces for pointers two back]\label{lem:twotraces}
Let $j$ be a parent array of length $n$ with trace $x$.
\begin{enumerate}[label=(\alph*),leftmargin=2.4em]
\item $\ell^{(s)}\le x\le u^{(s)}$, $\ell^{(s)}\le\ell^{(s+1)}$ and $u^{(s+1)}\le u^{(s)}$, for
every $s\ge0$.
\item $L^{(s)}\le\ell^{(s)}$ and $u^{(s)}\le U^{(s)}$ for the traces of
Lemma~\ref{lem:lowup}. So every position $i$ with $D(i)\le s$ is $s$-settled.
\item Every $s$-settled position lies in the set $K^{(s)}$ of the set-valued table iteration
with reach $2$ (Definition~\ref{def:setval}).
\item If every position is $s$-settled, then for every reach $R\ge2$ the set-valued table
iteration has stopped after at most $s$ rounds, and the table--Jacobi iteration applies $G_R$ at
most $s+1$ times.
\end{enumerate}
\end{lemma}

\begin{proof}
All the rules of Definition~\ref{def:twotraces} give a lower value that is nondecreasing in the
lower values it uses and nonincreasing in the upper ones, and an upper value with the opposite
behaviour.

(a) The bounds hold for $s=0$. Let $s\ge1$ and suppose they hold for $s-1$ and, in round $s$,
below $k$. The trace satisfies $x_k=\NOR(x_{k-1},x_p)$, so $\ell'\le x_k\le u'$; for $p=k-1$,
$x_k=1-x_{k-1}$. If $p=k-2$, then $x_k=(1-x_{k-2})x_q$ by~\eqref{eq:idtwo}, which gives
$\ell''\le x_k\le u''$ in the second case and for the second entries of the third; for $q=k-2$,
$x_{k-1}=1-x_{k-2}$ and $x_k=0$; and for $q=k-3$, $x_k=x_{k-3}$ by~\eqref{eq:idthree}, where
$k-2\ge1$ is a gate. The monotonicity in $s$ follows by induction on $s$ and $k$ from the remark
on the rules, starting from $\ell^{(0)}\le\ell^{(1)}$ and $u^{(1)}\le u^{(0)}$, which hold
because $\ell^{(0)}$ and $u^{(0)}$ are the least and the largest vectors with the entry $1$ at
$0$.

(b) For $s=0$ the traces are equal. By Lemma~\ref{lem:lowup}, $L^{(s)}_k$ and $U^{(s)}_k$ are
given by the rule $(\ell',u')$ applied to $L$ and $U$. By induction on $s$ and $k$ and the
remark on the rules, $\ell'\ge L^{(s)}_k$ and $u'\le U^{(s)}_k$, and
$\ell^{(s)}_k\ge\ell'$, $u^{(s)}_k\le u'$. If $D(i)\le s$ then $L^{(s)}_i=U^{(s)}_i$.

(c) By induction on $s$; for $s=0$ both sets are $\{0\}$. Suppose the claim holds for $s-1$.
In round $s$ of the set-valued iteration with reach $2$, call a sequence $z_0,\dots,z_i$
\emph{admissible} if $z_0=1$ and $z_t=\NOR(z_{t-1},b_t)$ for $1\le t\le i$, where $b_t=z_{j_t}$
if gate $t$ is short, $b_t=x_{j_t}$ if it is long and $j_t\in K^{(s-1)}$, and $b_t$ is $0$ or
$1$ otherwise. By the definition, and by Proposition~\ref{prop:setval}(a) for the values
$v^{(s-1)}$, the set $S_i$ of round $s$ is the set of the windows $(z_{i-1},z_i)$ of the
admissible sequences. We show $\ell^{(s)}_t\le z_t\le u^{(s)}_t$ for every admissible sequence,
by induction on $t$. A long read $b_t$ lies between $\ell^{(s-1)}_{j_t}$ and $u^{(s-1)}_{j_t}$:
if $j_t\in K^{(s-1)}$ by (a), and otherwise because $j_t$ is not $(s-1)$-settled, by the
induction hypothesis on $s$, so that these two numbers are $0$ and $1$. A short read is
$z_{t-1}$ or $z_{t-2}$, which lies between the traces of round $s$ and, by (a), of round $s-1$.
So $\ell'\le z_t\le u'$. If $j_t=t-2$, the identities~\eqref{eq:idtwo}
and~\eqref{eq:idthree} hold for $z$ with $b_{t-1}$ the read of gate $t-1$, and give
$\ell''\le z_t\le u''$ as in (a). Hence, if $i$ is $s$-settled, all admissible sequences have the
same $z_i$, and $i\in K^{(s)}$.

(d) By (c), $K^{(s)}$ is the set of all positions for the reach $2$, and then for every reach
$R\ge2$ by Proposition~\ref{prop:setval}(e). By Proposition~\ref{prop:setval}(c), $y^{(s)}=x$,
so $G_R(y^{(s)})=y^{(s)}$.
\end{proof}

\begin{definition}[the trace map]\label{def:tracemap}
Consider a design with independent reads and pointers two back (Definition~\ref{def:two}) with
the set of classes $\mathcal C$, and a vector $v$ with entries $v_{(c,-)},v_{(c,+)}\in[0,1]$ for
$c\in\mathcal C$. For a gate $g$ of mixed type $\mu$ put
$y^\mp_g=\mu_1+\sum_c\mu_c(1-v_{(c,\mp)})$, for a gate of type $\mathrm L_1$ put $y^\mp_g=1$,
and in both cases $q^\mp_g=B+(1-\eps)y^\mp_g$. The \emph{trace recursion} asks for numbers
$a^-_g$ and $a^+_g$, for the gates $g$ of every rail, periodic with the period of the design,
with:
\begin{itemize}[leftmargin=3.2em]
\item[(T1)] if $g$ has a mixed type or the type $\mathrm L_1$:
$a^-_g=1-q^+_g\,a^+_{g-1}$ and $a^+_g=1-q^-_g\,a^-_{g-1}$;
\item[(T2)] if $g$ has type $\mathrm L_2$ and $g-1$ has a mixed type or the type $\mathrm L_1$:
$1-a^-_g=B\,a^+_{g-1}+(1-\eps)(1-q^-_{g-1})\,a^+_{g-2}$ and
$1-a^+_g=B\,a^-_{g-1}+(1-\eps)(1-q^+_{g-1})\,a^-_{g-2}$;
\item[(T3)] if $g$ and $g-1$ have type $\mathrm L_2$:
$1-a^-_g=B\,a^+_{g-1}+(1-\eps)A\,a^+_{g-2}+(1-\eps)^2(1-a^-_{g-3})$ and
$1-a^+_g=B\,a^-_{g-1}+(1-\eps)A\,a^-_{g-2}+(1-\eps)^2(1-a^+_{g-3})$.
\end{itemize}
The \emph{trace map} $\Psi$ sends $v$ to the vector with the entries $1-a^-_g$ at $(k,-)$ and
$1-a^+_g$ at $(k,+)$, for a gate $g$ of class $k$.
\end{definition}

\noindent In the ideal model, $a^-_g$ and $a^+_g$ are the probabilities that the lower and the
upper trace of Definition~\ref{def:twotraces} hold $0$ at $g$, when the lower and the upper
values of a read of class $c$ are $1$ with the probabilities $v_{(c,-)}$ and $v_{(c,+)}$, a
gate of type $\mathrm L_2$ keeps its pointer with probability $1-\eps$, and the rule
$(\ell'',u'')$ is used at the gates that do. For example, a kept gate $g$ after a gate of mixed
type has $\ell_g=(1-u_{g-2})\ell_q$, a product of two independent bits, and a resampled one has
$\ell_g=\NOR(u_{g-1},\cdot)$; this is (T2). We use the recursion itself, and derive it for the
chains of a smoothed instance in the proof of Theorem~\ref{thm:twolayered}.

\begin{lemma}[two inequalities]\label{lem:traceineq}
Let $0.455\le\eps<1$, and let $f$ and $\hat f$ be as in Lemma~\ref{lem:potential}.
\begin{enumerate}[label=(\alph*),leftmargin=2.4em]
\item For all $X,X'\in[lo,hi]$ and $y\in[0,1]$, with $q=B+(1-\eps)y$,
\begin{multline*}
Bq\,f(X)+(1-\eps)(1-q)\,f(X')+(1-\eps)\bigl(BX+(1-\eps)X'\bigr)\hat f(y)\\
\le\ \tfrac9{10}\,f\bigl(1-B(1-Xq)-(1-\eps)(1-q)X'\bigr).
\end{multline*}
\item For all $a_1,a_2,a_3\in[lo,hi]$,
\[
B\,f(a_1)+(1-\eps)A\,f(a_2)+(1-\eps)^2f(a_3)\ \le\ \tfrac9{10}\,f\bigl(1-Ba_1-(1-\eps)Aa_2-(1-\eps)^2(1-a_3)\bigr).
\]
\end{enumerate}
In both, the argument of $f$ on the right lies in $[lo,hi]$.
\end{lemma}

\begin{proof}
For the last statement recall $1-B\,lo=hi$ and $1-r\,hi=lo$. In (a), $1-Xq\in[lo,hi]$, so
$B(1-Xq)+(1-\eps)(1-q)X'$ lies between $B\,lo$ and $(B+1-\eps)hi=r\,hi$. In (b),
$Ba_1+(1-\eps)Aa_2+(1-\eps)^2(1-a_3)$ is at least $B\,lo$ and at most
$\bigl(B+(1-\eps)A+(1-\eps)^2r\bigr)hi\le r\,hi$, because $1-lo=r\,hi$ and
$A+(1-\eps)r\le1$. The inequalities are verified by the program \texttt{traces\_check.c}, with
the interval arithmetic of \texttt{potential\_check.c} on boxes of the three variables. The runs
\texttt{traces\_check 0.45499 0.50001 0.8999 0.0005} and
\texttt{traces\_check 0.49999 1.0 0.8999 0.005} cover $[0.455,1)$ in $192$ intervals of noise
levels, with $419{,}506$ boxes for (a) and $195{,}092$ for (b). Numerically, on a grid, the
largest ratios of the two sides at noise $0.455$ are $0.842$ and $0.791$.
\end{proof}

\begin{proposition}[the trace map]\label{prop:tracemap}
Consider a design with independent reads and pointers two back, with density map $\Phi$ and
trace map $\Psi$.
\begin{enumerate}[label=(\alph*),leftmargin=2.4em]
\item For every $v$ the trace recursion has exactly one solution. All its numbers lie in
$[lo,hi]$, and $\Psi$ is continuously differentiable.
\item If $v_{(c,-)}=v_{(c,+)}=d_c$ for all $c$, then $\Psi(v)_{(k,-)}=\Psi(v)_{(k,+)}=\Phi(d)_k$.
\item $\Psi$ is monotone for the order $\sqsubseteq$ of Subsection~\ref{sec:layered}: if
$v'\sqsubseteq v$, that is $v_{(c,-)}\le v'_{(c,-)}$ and $v_{(c,+)}\ge v'_{(c,+)}$ for all $c$,
then $\Psi(v')\sqsubseteq\Psi(v)$.
\item Let $0.455\le\eps<1$. For every $v$ and every vector $\xi$ with
$|\xi_{(c,\tau)}|\le f(1-v_{(c,\tau)})$ for all $(c,\tau)$, the derivative of $\Psi$ at $v$
satisfies $|(D\Psi(v)\xi)_{(k,\sigma)}|\le\theta\,f(1-\Psi(v)_{(k,\sigma)})$ for all
$(k,\sigma)$, with $\theta=0.9998$. So $\Psi$ maps into the set $\mathcal K$ of the classes
$(c,\pm)$ and satisfies $\delta(\Psi(v),\Psi(v'))\le\theta\,\delta(v,v')$ there, and it is
Lipschitz in the maximum norm with a constant that depends only on $\eps$.
\item Let $0.455\le\eps<1$, $d\in[0,1]^{\mathcal C}$, and let $\alpha_g$ and $q_g$ be the
numbers of the chain over $d$ (Definition~\ref{def:two}). For $0\le\kappa\le A/8$ and numbers
$e_c\ge0$ let $w_g$, for the gates of every rail, be the periodic solution of
\begin{itemize}[leftmargin=3.2em]
\item[(U1)] $w_g=(q_g+\kappa)w_{g-1}+\kappa(w_{g-2}+w_{g-3})+(1-\eps)(\alpha_{g-1}+\kappa)\sum_c\mu_c(g)e_c$
if $g$ has a mixed type or the type $\mathrm L_1$ (the sum is absent for $\mathrm L_1$);
\item[(U2)] $w_g=(B+\kappa)w_{g-1}+\bigl((1-\eps)(1-q_{g-1})+\kappa\bigr)w_{g-2}+\kappa w_{g-3}
+(1-\eps)(1-\eps+\kappa)(\alpha_{g-2}+\kappa)\sum_c\mu_c(g-1)e_c$ if $g$ has type $\mathrm L_2$ and $g-1$ a
mixed type or the type $\mathrm L_1$;
\item[(U3)] $w_g=(B+\kappa)w_{g-1}+\bigl((1-\eps)A+\kappa\bigr)w_{g-2}+\bigl((1-\eps)^2+\kappa\bigr)w_{g-3}$
if $g$ and $g-1$ have type $\mathrm L_2$.
\end{itemize}
It is unique and linear in $e$; write $w_g=\sum_c\overline M^\kappa_{k,c}(d)e_c$ for a gate $g$
of class $k$. There is $\kappa_*>0$, depending only on $\eps$, such that for $\kappa\le\kappa_*$
\begin{equation}\label{eq:Mpottwo}
\sum_c\overline M^\kappa_{k,c}(d)\,f(1-d_c)\ \le\ \theta_1\,f\bigl(1-\Phi(d)_k\bigr),\qquad\theta_1=\tfrac{1+\theta}2 .
\end{equation}
\end{enumerate}
\end{proposition}

\begin{proof}
(a) In each of (T1), (T2), (T3) the two numbers at $g$ are affine functions of at most three
earlier numbers, with coefficients whose absolute values add up to at most $r$: in (T1) the
coefficient is at most $r$; in (T2) the sum is $B+(1-\eps)(1-q)\le B+(1-\eps)=r$; in (T3) it is
$B+(1-\eps)(A+1-\eps)\le r$, as $A\le\eps$. So on bounded sequences, periodic or not, the
recursion is a fixed point equation for an affine map that contracts the supremum norm by the
factor $r<1$: there is exactly one bounded solution, it is periodic, and over one period it is
the solution of a linear system whose matrix is invertible and whose coefficients are
polynomials in $v$. The solution is the limit of the iterates of that map from any start, and
the map sends sequences with entries in $[0,1]$ to such sequences, since the right sides of
(T2) and (T3) lie between $0$ and $r$; so the solution has entries in $[0,1]$. Let $m$ and $M$
be its least and largest numbers. In all three cases $1-a_g\ge Bm$, so $M\le1-Bm$. In (T1) and
(T2), $1-a_g\le rM$, and in (T3), $1-a_g\le(B+(1-\eps)A)M+(1-\eps)^2(1-m)$. If $m$ is attained
at a gate of case (T3), then $(1-m)\bigl(1-(1-\eps)^2\bigr)\le(B+(1-\eps)A)M$, and
$r\bigl(1-(1-\eps)^2\bigr)-B-(1-\eps)A=\eps(1-\eps)r\ge0$ gives $1-m\le rM$ there too. So
$1-m\le rM$ in every case, and with $M\le1-Bm$ this gives $M\le hi$ and $m\ge lo$ as in the
proof of Lemma~\ref{lem:onebit}(b).

(b) On the diagonal the recursion is symmetric in the two signs, so by uniqueness $a^-=a^+$. It
is enough to show that the numbers $\alpha_g$ of the chain over $d$ satisfy the recursion with
$q^\mp_g=q_g$. (T1) is Lemma~\ref{lem:twochain}(a). For a gate $g$ of type $\mathrm L_2$ that
lemma gives $1-\alpha_g=B\alpha_{g-1}+(1-\eps)(\alpha_{g-1}+\alpha_{g-2}-1)$. If $g-1$ has a
mixed type or the type $\mathrm L_1$, then
$\alpha_{g-1}+\alpha_{g-2}-1=(1-q_{g-1})\alpha_{g-2}$, which is (T2). If $g-1$ has type
$\mathrm L_2$, then
$\alpha_{g-1}+\alpha_{g-2}-1=(1-\eps)(1-\alpha_{g-3})+(1-r)\alpha_{g-2}$ by the same lemma for
$g-1$, and $1-r=A$; this is (T3).

(c) Call a pair of sequences $(a^-,a^+)$ narrower than another if its $a^-$ is smaller and its
$a^+$ larger at every gate. If $v'\sqsubseteq v$, then
$q^-_g(v)\ge q^-_g(v')$ and $q^+_g(v)\le q^+_g(v')$. In (T1) to (T3), with nonnegative entries,
each right side for $1-a^-_g$ increases, and each right side for $1-a^+_g$ decreases, when the
earlier pair is made narrower and when $v$ is replaced by $v'$. Iterate the recursion from the
sequences that are constant $1$ for $a^-$ and $0$ for $a^+$, for $v$ and for $v'$: all iterates
have entries in $[0,1]$, the iterates for $v'$ are narrower than those for $v$, and both converge
to the solutions by (a).

(d) Let $v$ and $\xi$ be as stated and let $D^\mp_g$ be the absolute value of the derivative of
$a^\mp_g$ in the direction $\xi$. For a gate $g$ of mixed type, the derivative of $q^\mp_g$ in
the direction $\xi$ is $-(1-\eps)\sum_c\mu_c(g)\xi_{(c,\mp)}$, of absolute value at most
$(1-\eps)\sum_c\mu_c(g)f(1-v_{(c,\mp)})\le(1-\eps)\hat f(y^\mp_g)$, by
Lemma~\ref{lem:potential}(a) and Jensen's inequality as in the proof of
Theorem~\ref{thm:indep}; for type $\mathrm L_1$ it is $0=\hat f(1)$. Differentiating (T1),
\[
D^-_g\le q^+_gD^+_{g-1}+(1-\eps)a^+_{g-1}\hat f(y^+_g),
\]
and by Lemma~\ref{lem:potential}(b), with $x=a^+_{g-1}\in[lo,hi]$ and $y=y^+_g$,
\[
q^+_gf(a^+_{g-1})+(1-\eps)a^+_{g-1}\hat f(y^+_g)\le\theta f(a^-_g).
\]
In (T2) replace $a^+_{g-1}$ by $1-q^-_{g-1}a^-_{g-2}$, which is (T1) at $g-1$, and differentiate:
with $q=q^-_{g-1}$, $X=a^-_{g-2}$, $X'=a^+_{g-2}$ and $y=y^-_{g-1}$,
\[
D^-_g\le Bq\,D^-_{g-2}+(1-\eps)(1-q)\,D^+_{g-2}+(1-\eps)\bigl(BX+(1-\eps)X'\bigr)\hat f(y),
\]
and $a^-_g=1-B(1-Xq)-(1-\eps)(1-q)X'$; Lemma~\ref{lem:traceineq}(a) bounds the right side, with
$f(X)$ and $f(X')$ in place of $D^-_{g-2}$ and $D^+_{g-2}$, by $\frac9{10}f(a^-_g)$. In (T3),
\[
D^-_g\le B\,D^+_{g-1}+(1-\eps)A\,D^+_{g-2}+(1-\eps)^2D^-_{g-3},
\]
and Lemma~\ref{lem:traceineq}(b) with $a_1=a^+_{g-1}$, $a_2=a^+_{g-2}$, $a_3=a^-_{g-3}$ bounds
the same expression in $f$ by $\frac9{10}f(a^-_g)$. The same holds with the signs exchanged. So
at every gate and for each sign there are earlier gates and signs, indexed by $t$, and
coefficients $w_t\ge0$ with $\sum_tw_t\le r$, such that
$D_g-\theta f(a_g)\le\sum_tw_t\bigl(D_t-f(a_t)\bigr)\le\sum_tw_t\bigl(D_t-\theta f(a_t)\bigr)$.
If the largest of the periodic numbers $D^\mp_g-\theta f(a^\mp_g)$ were positive, this would
bound it by $r$ times itself. So $D^\mp_g\le\theta f(a^\mp_g)$, which is the claim. The
consequences follow as in the proof of Theorem~\ref{thm:indep}(a), with (a) for the range; for
the Lipschitz constant, bound $\hat f$ by $1$ in the three inequalities for $D$, which gives
$D\le(1-\eps)/(1-r)$ when $|\xi_{(c,\tau)}|\le1$.

(e) For $\kappa\le A/8$ the coefficients of (U1) to (U3) add up to at most
$r+3\kappa\le1-A/2$, so the periodic solution exists, is unique and nonnegative, and is linear
in $e$. Let $e_c=f(1-d_c)$ and $s_g=\sum_c\mu_c(g)f(1-d_c)\le\hat f(y_g)$. For $\kappa=0$:
(U1) with~\eqref{eq:gatepot} gives $w_g-\theta f(\alpha_g)\le q_g(w_{g-1}-f(\alpha_{g-1}))$. In
(U2) insert (U1) for $w_{g-1}$:
$w_g=\bigl(Bq+(1-\eps)(1-q)\bigr)w_{g-2}+(1-\eps)\bigl(B+(1-\eps)\bigr)\alpha_{g-2}s_{g-1}$ with
$q=q_{g-1}$; by Lemma~\ref{lem:traceineq}(a) with $X=X'=\alpha_{g-2}$ and by (T2) on the
diagonal, $w_g-\theta f(\alpha_g)\le\bigl(Bq+(1-\eps)(1-q)\bigr)(w_{g-2}-f(\alpha_{g-2}))$. In
(U3), Lemma~\ref{lem:traceineq}(b) and (T3) on the diagonal give
$w_g-\theta f(\alpha_g)\le B(w_{g-1}-f(\alpha_{g-1}))+(1-\eps)A(w_{g-2}-f(\alpha_{g-2}))
+(1-\eps)^2(w_{g-3}-f(\alpha_{g-3}))$. As in (d), $w_g\le\theta f(\alpha_g)$ for $\kappa=0$. For
$\kappa>0$, the solution is at most $2/A$ times the largest source, so it is bounded by a
constant that depends only on $\eps$; the difference between the solutions for $\kappa$ and for
$0$ solves the recursion for $\kappa=0$ with sources of size at most a constant times $\kappa$,
so it is at most $C\kappa$ with $C$ depending only on $\eps$. As $f(\alpha_g)\ge m_f>0$, the
claim holds for $\kappa\le\kappa_*=(\theta_1-\theta)m_f/C$.
\end{proof}

\noindent For $\kappa=0$ the number $w_g$ of (e) is the probability that the two copies of
Definition~\ref{def:two} differ at $g$, added over the sources with the weights $e_c$: the
identities~\eqref{eq:idtwo} and~\eqref{eq:idthree} express a difference at $g$ by a difference
at one earlier gate and one read. So $\overline M^0(d)$ is the matrix of
Definition~\ref{def:two}, and (e) with $\kappa=0$ is another proof of Theorem~\ref{thm:two},
without the runs of Lemma~\ref{lem:runs}. We do not use this.

\begin{theorem}[layered bases with pointers two back]\label{thm:twolayered}
Define layered bases as in Definition~\ref{def:layered}, except that the program of a bundle
may also contain the type $\mathrm L_2$, whose gates list the position two before them. Let
$0.455\le\eps<1$ and $T\ge1$. There are $c'>0$ and $h_0$, $C_1$, $C_2$, $n_1$, depending only on
$\eps$ and $T$, such that for every $n\ge n_1$ the following holds for every such base of length
$n$ with $h\ge h_0$, $m\ge C_1\ln(n+1)$ and $n_0\ge(1-c')n$: with probability at least $1-5/n$
over the coins and the smoothing, every position is $s$-settled for $s=\lfloor C_2\ln n\rfloor$.
So, for every reach $R\ge2$, the set-valued table iteration stops within $C_2\ln n$ rounds, and
the table--Jacobi iteration applies $G_R$ at most $C_2\ln n+1$ times; the latter stops only at
the trace (Theorem~\ref{thm:table}(a)), so it is errorless.
\end{theorem}

\begin{proof}
The last sentence is Lemma~\ref{lem:twotraces}(d). For the rest we follow the proof of
Theorem~\ref{thm:layered}, with the traces $\ell^{(s)}$, $u^{(s)}$ in place of $L^{(s)}$,
$U^{(s)}$, with ``settled'' in place of ``determined'' and ``open'' meaning ``not settled'', and
with the trace map $\Psi_{\mathcal Q}$ of the program of $\mathcal Q$ in place of
$\Phi^{(2)}_{\mathcal Q}$. The $\sigma$-value of a position in round $s$ is $x$ if it is
$s$-settled and $\bot$ otherwise. As there, a gate of type $\mathrm L_1$ is treated as the mixed
type with $\mu_1=1$. We say what changes in each step.

\emph{Step 0.} By Lemma~\ref{lem:twotraces}(b) every position up to $n_0$ is $s_z$-settled on
$\mathcal Z$. Nothing else changes.

\emph{Step 1.} For a gate $g$ of mixed type, $R_g$, $\Xi_g$ and $y^\mp_g$ are as before. A gate
$g$ of type $\mathrm L_2$ is \emph{kept} if its pointer after the smoothing is $g-2$, which has
probability $1-\eps+\eps/g$ and is independent of $\mathcal H_{g-1}$. Let $R_g$ be its read
when it is not kept. Given $\mathcal H_{g-1}$, the event that $g$ is not kept and $R_g=1$ has
probability $\eps\psi^L_{g-1}$, and the event that it is not kept and $R_g\ne0$ has probability
$\eps\psi^U_{g-1}$, both up to $2\eps/g$. In this proof we enlarge $\eta'$ by $2/n_0$, which
absorbs these errors, and take $\eta'$ small and $n$ large, so that
$\eps(\varpi+\eta')\le\eps(1-1/g)$. Build $\Xi_g$, uniform and independent of $\mathcal H_{g-1}$, with
$R_g=1$ and $g$ not kept exactly if $\Xi_g$ is below the first of these probabilities, $R_g\ne0$ and
$g$ not kept exactly if it is below the second, and $g$ kept
exactly if $\Xi_g\ge\eps(1-1/g)$; put $y^-_g=\mathbf 1[\Xi_g<\eps(\varpi-\eta')]$ and
$y^+_g=\mathbf 1[\Xi_g<\eps(\varpi+\eta')]$. So for every gate of a chain the lower value of the
read is at least $y^-_g$ and the upper value at most $y^+_g$, when the gate is not a kept gate
of type $\mathrm L_2$.

For a chain with the gates $g_1<\dots<g_N$ define bits $X^-_g\le X^+_g$: for $g<g_1$,
$X^-_g=0$ and $X^+_g=1$, and $y^-_{g_1-1}=0$, $y^+_{g_1-1}=1$. For $g\ge g_1$:
if $g$ is not a kept gate of type $\mathrm L_2$, $X^-_g=\NOR(X^+_{g-1},y^+_g)$ and
$X^+_g=\NOR(X^-_{g-1},y^-_g)$; if $g$ is one and $g-1$ is not, $X^-_g=(1-X^+_{g-2})y^-_{g-1}$ and
$X^+_g=(1-X^-_{g-2})y^+_{g-1}$; if both are, $X^\mp_g=X^\mp_{g-3}$. On
$\mathcal E\cap\mathcal Z$, $\ell^{(s)}_g\ge X^-_g$ and $u^{(s)}_g\le X^+_g$ for every gate of
the chain, by induction on $g$: in each case the rule used for $X$ is one of the rules whose
maximum, respectively minimum, defines the traces, applied to bounds of its entries; in the
second case, if gate $g-1$ has become a NOT gate, the traces are $0$ and $X^-_g=0$, because the
effective read of a NOT gate is a determined zero, and if $g-1$ has pointer $g-3$, the entries
with the read of the round before are used. Given
$\mathcal G_{\mathcal Q}$ these pairs are independent for the $m$ chains, with a law that
depends only on $v$ and on the positions of the gates.

Let $a^\mp_g=\PP[X^\mp_g=0\mid\mathcal G_{\mathcal Q}]$. In the first case $X^+_{g-1}$ and
$y^+_g$ are independent, and for a gate of type $\mathrm L_2$ the event that it is not kept and
$y^+_g=0$ has probability $\eps(1-1/g)-\eps(\varpi+\eta')$. In the second case, $g$ kept,
$X^+_{g-2}$ and $y^-_{g-1}$ are independent, and for a gate $g-1$ of type $\mathrm L_2$ the
event that it is not kept and $y^-_{g-1}=1$ has probability $\eps(\varpi-\eta')$. In the third
case the two kept events are independent of $X^-_{g-3}$. Adding the cases, the numbers
$a^\mp_g$ satisfy (T1), (T2), (T3) along the chain, with $q^\mp$ computed from $v$, with every
coefficient changed by at most $\varsigma=2\eps\eta'+4/n_0$, and with the start $a^-=1$,
$a^+=0$ before $g_1$. For small $\eta'$ and large $n$ the changed coefficients still add up to
at most $1-A/2$ in every case, and then two solutions of such recursions whose coefficients
differ by at most $\varsigma$ differ at the $t$-th gate by at most
$(1-A/2)^{\lfloor t/3\rfloor}+6\varsigma/A$, whatever their starts in $[0,1]$. The chain has $T_{\mathcal Q}(h-1)$ gates before
its last period. So, in place of~\eqref{eq:Pphi},
\[
\bigl|P^\pm_k(v)-\Psi_{\mathcal Q}(v)_{(k,\pm)}\bigr|\ \le\ \frac{6(2\eps\eta'+4/n_0)}{A}+(1-A/2)^{\lfloor(h-1)/3\rfloor}=:\kappa_0,
\]
where $P^\pm_k(v)=\PP[X^\pm_g=1\mid\mathcal G_{\mathcal Q}]$ for the copy $g$ of phase $k$. The
true values are treated in the same way: with thresholds for the true reads
(Lemma~\ref{lem:reads}), the same rules give bits $X^-_g\le x_g\le X^+_g$ whose probabilities
satisfy the recursion for $v=(t,t)$ up to the same errors, and
$\Psi_{\mathcal Q}(t,t)=(\Phi_{\mathcal Q}(t),\Phi_{\mathcal Q}(t))$ by
Proposition~\ref{prop:tracemap}(b). So~\eqref{eq:onebundle} holds with $\Psi_{\mathcal Q}$, and
the event $\mathcal F$ has the same bound.

\emph{Step 2.} This step uses of $\Phi^{(2)}_{\mathcal Q}$ only that it maps into $\mathcal K$,
contracts $\delta$ by $\theta$, is monotone for $\sqsubseteq$, is Lipschitz with a constant
depending on $\eps$, and equals $(\Phi_{\mathcal Q},\Phi_{\mathcal Q})$ on the diagonal; and of
$\Phi_{\mathcal Q}$ that it contracts $\delta$. These are Proposition~\ref{prop:tracemap}(a)
to~(d) and Theorem~\ref{thm:two}. So the step holds as written.

\emph{Step 3.} Let $s_1\le s<s_2$ and let $g>n_0$ be a gate that is open in round $s+1$. The
rules of Definition~\ref{def:twotraces} give, for the gate $g$:
\begin{itemize}[leftmargin=3.2em]
\item[(O1)] in every case, by the rule $(\ell',u')$: $g-1$ is open and the read of $g$ is not a
settled one in round $s$, or the read of $g$ is $s$-open and $x_{g-1}$ is not a settled one;
\item[(O2)] if $g$ has pointer $g-2$ and $g-1$ is not a NOT gate, by the second rule
$(\ell'',u'')$: $g-2$ is open and the read of $g-1$ is not a settled zero in round $s$, or the
read of $g-1$ is $s$-open and $x_{g-2}$ is not a settled one;
\item[(O3)] if $g$ has pointer $g-2$ and $g-1$ has pointer $g-3$: $g-3$ is open.
\end{itemize}
Here (O1) is as in the proof of Theorem~\ref{thm:layered}; for (O2), $\ell''<u''$ needs
$\ell_{g-2}=0$ and an upper value $1$ of the read, and then $u_{g-2}=1$ or a lower value $0$ of
the read. Let $E_g$ be the probability that $\mathcal F_{\mathcal Q}$ holds and $g$ is open in
round $s+1$. For a gate of mixed type use (O1); for a gate of type $\mathrm L_2$ use (O1) if it
is not kept, (O2) if it is kept and $g-1$ is not a kept gate of type $\mathrm L_2$, and (O3) if
both are kept; (O2) covers a gate $g-1$ whose pointer is $g-3$ by resampling. Conditioning as
in that proof on $\mathcal H_{g-1}$, or on $\mathcal H_{g-2}$ for the read of $g-1$ in (O2), and
using Step 2 and Step 0: a read of a gate of mixed type is not a settled one with probability at
most $q_g(d)+\kappa$ and is not a settled zero with probability at most $1-q_g(d)+\kappa$; a
gate of type $\mathrm L_2$ is not kept and its read is not a settled one with probability at
most $B+\kappa$, is not kept and its read is not a settled zero with probability at most
$A+\kappa$, and is kept with probability at most $1-\eps+\kappa$; and a read is $s$-open with
probability at most $(1-\eps)\sum_c\mu_c o_c+\eps\Omega$ for a mixed type and $\eps\Omega$ for a
gate that is not kept. Here $\kappa$ is fixed with $3\kappa\le\min(\kappa_*,A/8)$, and $J_0$ with
$(1-A/2)^{\lfloor J_0/3\rfloor}\le\kappa/4$, since a contribution of the start now decays once
in three gates. For gates at least $J_0$ above the start of the chain, the predecessor
used in (O1) or (O2) is not a settled one with probability at most $\alpha(d)+\kappa$: as
before, its value is $0$ with a probability close to $\alpha(d)$ by the sandwich of Step~1 for
the true values, and it is open with probability at most
$K(\kappa_3+c')+(1-A/2)^{\lfloor J_0/3\rfloor}$, by the recursion that follows, run given
$\mathcal G_{\mathcal Q}$ with the bound $1$ in place of $\alpha(d)+\kappa$.
So the numbers $E_g$ satisfy, with inequality, the recursions (U1), (U2), (U3) of
Proposition~\ref{prop:tracemap}(e) at the ideal densities $d$, with $3\kappa$ in place of
$\kappa$ and $e_c=\EE[o_c;\mathcal F_{\mathcal Q}]$, up to additional sources of size at most
$K_0c'e^*_s$ and, for the gates within $J_0$ of the start and for the start, with sources of
size at most $K_0$ and $e^*_{s+1}$. The periodic solution dominates the solution along the
chain, and the contributions of the start decay by the factor $1-A/2$ every three gates. Hence
\[
\PP[i\text{ is }(s+1)\text{-open}]\ \le\ \sum_c\overline M^{3\kappa}_{k,c}(d)\,e_s(c)+K_0\,\zeta\,(e^*_s+e^*_{s+1})+K_0\,\beta
\]
for a copy $i$ of $(\mathcal Q,k)$, with
$\zeta=2c'+(1-A/2)^{\lfloor(T_{\mathcal Q}(h-1)-J_0)/3\rfloor}$, and every position above $n_0$
is $(s+1)$-open with probability at most
$K_0\max_ce_s(c)+K_0\zeta(e^*_s+e^*_{s+1})+K_0\beta$. With~\eqref{eq:Mpottwo} in place
of~\eqref{eq:Mpotkappa}, for $3\kappa\le\kappa_*$, the rest of the step and the conclusion are
unchanged.
\end{proof}

\begin{remark}[what this adds, and what it does not]\label{rem:twolayered}
\begin{enumerate}[label=(\roman*),leftmargin=2.4em]
\item \emph{The two iterations on these bases; numerical.} The script
\texttt{twoback\_layered.c} builds a layered base with $m=200$ chains of four periods per
bundle, in which the output class of a bundle is read by the bundle above, over a preamble of
half the positions, and counts $\max_iD(i)$ and the rounds until every position is settled. For
the program $(\mathrm E,\mathrm L_2)$ and $25$, $100$ and $400$ bundles, one smoothing each:
at noise $0.46$, $\max_iD(i)=90$, $331$, $1282$ and $20$, $30$, $30$ rounds; at noise $0.50$,
$71$, $253$, $885$ and $18$, $30$, $32$ rounds. For the program
$(\mathrm E,\mathrm L_2,\mathrm L_2)$ at noise $0.46$ with $100$ and $400$ bundles,
$\max_iD(i)=268$ and $718$, against $19$ and $21$ rounds. So on these bases the
relevance-pruned iteration needs a number of rounds that grows with the number of bundles, and
the traces settle in a number of rounds that hardly grows. Both statements are simulations: the
bases have $h=4$ and a preamble of half the positions, and the theorem asks for $h\ge h_0$ and
a preamble of all but a small fraction.
\item \emph{A check of Lemma~\ref{lem:twotraces}.} The script \texttt{twoback\_check.py}
computes the traces of Definition~\ref{def:twotraces} on all $46{,}232$ parent arrays with
$2\le n\le8$ and confirms parts (a), (b) and (c) in every round. On these arrays the number of
rounds until every position is settled equals the number of rounds of the set-valued iteration
with reach $2$ in every case; the referee of this subsection found that on all arrays with
$n\le10$ the settled positions are exactly the set $K^{(s)}$ of that iteration, in every
round. We have not proved this.
\item \emph{What is still missing.} The theorem needs wide bundles, like
Theorem~\ref{thm:layered}, and for the same reason: the reads of a chain must be independent.
It covers pointers one and two back. For pointers $k\ge3$ back there are identities of the same
kind only for some arrangements, and we have not found traces. Nothing here is a statement
about all bases.
\end{enumerate}
\end{remark}

\subsection{Dependent reads: a family with a proved bound}\label{sec:scan}

Every theorem of this section so far needs reads that are independent of each other, and
Remark~\ref{rem:dependent} is numerical. This subsection proves a bound for an infinite family of
designs whose reads depend on each other: chains that read runs of consecutive positions of a
plain chain, in every arrangement and with every period. The bound is the one for independent
reads, up to a small constant, and it holds from noise $0.446$ on. An exact computation for a
second family then shows that dependence can help, by a small amount
(Remark~\ref{rem:scan}).

\begin{definition}[scanning chains]\label{def:scan}
Let $0<\eps<1$ and $d\in[0,1]$. A \emph{plain chain over $d$} is a chain of the design with one
class and the type $\mathrm E$ over the density $d$ (Definition~\ref{def:design}): its targets
are independent bits that are $1$ with probability $d$. Let $W=\ell_1\dots\ell_T$ be a word over
the letters $\mathrm U$ and $\mathrm R$ with $\ell_1=\mathrm U$. A \emph{scanning chain of $W$
over $d$} is a sequence of bits with $x_g=\NOR(x_{g-1},y_g)$, in which gate $g$ has the letter
$\ell_{((g-1)\bmod T)+1}$, its read is resampled with probability $\eps$ to a bit that is $1$
with probability $\varpi$, and otherwise is its target: for a gate with the letter $\mathrm U$,
the value at position $0$ of a new stationary plain chain over $d$, its \emph{window}, which is
independent of everything before; for a gate with the letter $\mathrm R$, the value at the next
position of the window of the gate before it. A \emph{block} is a gate with the letter
$\mathrm U$ together with the gates $\mathrm R$ that follow it; a block with $k$ gates
$\mathrm R$ reads the positions $0,\dots,k$ of its window. The chain is stationary with period
$T$. A \emph{density of $W$} is a number $d$ such that the mean density of ones of the scanning
chain of $W$ over $d$, over a period, is $d$.

For a block and a position $j\le k$ of its window, run two copies of the chain with the same
choices, except that the target of the gate at position $j$ of the window is $0$ in one copy and
$1$ in the other, and let $N(j)$ be the expected number of gates of the scanning chain at which
the copies differ. The \emph{two-layer spread} is
\[
G_2(W,d)=\frac1T\sum_{\text{blocks of a period}}\ \sum_{j\le k}N(j),
\]
where the inner sum is over all positions $j\le k$ of the window of the block, also the negative
ones.
\end{definition}

\noindent This is the ideal model of a tower in which plain layers and scanning layers
alternate. A gate of a plain layer lists a uniformly random position of the scanning layer
below, so its targets are independent with the density of that layer, and a density of $W$ is a
density that such a tower can keep. A gate $\mathrm U$ of a scanning layer lists a uniformly
random position of the plain layer below, and a gate $\mathrm R$ the position after the one
listed by the gate before it; these are the rules $\mathrm U$ and $\mathrm R(1)$ of
Remark~\ref{rem:dependent}. If a position of a scanning layer is flipped, it is the target of
one gate of the plain layer above on average, and that gate is at a given place relative to the
windows of the next scanning layer with a frequency of $1/T$ per place and block. So $G_2(W,d)$
is the mean number of positions of a scanning layer that change when one position of the
scanning layer two below is flipped; as for $\overline M$, we use this only as the meaning of
the number. For the word $\mathrm U$, all reads are independent, the density is $\pi_0$, and
$G_2=\lambda^2$ by Proposition~\ref{prop:designs}(b).

\begin{theorem}[scanning programs]\label{thm:scan}
Let $W$ be a word over $\mathrm U$ and $\mathrm R$ that begins with $\mathrm U$ and has at most
$120$ consecutive letters $\mathrm R$, cyclically, and let $d$ be a density of $W$. If
$0.446\le\eps<1$, then $G_2(W,d)\le0.999$. If $0.455\le\eps<1$, then $G_2(W,d)\le0.99$.
\end{theorem}

\begin{proof}
\emph{Blocks.} Let a block have $k$ gates $\mathrm R$, and let $p$ be the probability that the
chain holds $0$ at the gate before the block. The window of the block is independent of the
chain before it. So the joint law of the two copies inside the block, for every $j$, is
determined by $p$, and every probability of an event inside the block is an affine function of
$p$. Let $\mathrm{cost}_k(p)$ be the expected number of gates of the block at which the copies
differ, added over all $j\le k$; let $X_k(p)$ be the probability that they differ at the last
gate of the block, added over all $j\le k$; let $\mathrm{out}_k(p)$ be the probability that the
chain holds $0$ at the last gate of the block, and $\mathrm{ones}_k(p)$ the expected number of
ones in the block. The sums over $j\le k$ are finite: the two copies of the window differ at
position $0$ only if they differ at all positions from $j$ on, which has probability
$(1-\eps)\alpha_dq_d^{-j}$ for $j\le0$, where $q_d=B+(1-\eps)(1-d)$ is the probability that a
gate of the plain chain reads a zero and $\alpha_d=1/(1+q_d)$ the probability that it holds one.

\emph{Later blocks.} After the block the windows are new, so both copies read the same bits.
If they differ at a gate, they differ at the next one exactly if its read is $0$. Let $S_k$ be
the expected number of gates of a block with $k$ gates $\mathrm R$ up to which all reads of the
block are $0$, and $Q_k$ the probability that all its reads are $0$; these do not depend on the
chain before the block. Let $V$ be the supremum, over all sequences of blocks with at most $120$
gates $\mathrm R$ each, of the expected
number of gates at which two copies differ that differ at the gate before the first of these
blocks. It is finite, as every read is $0$ with conditional probability at most $r$, and
$V\le\max_k(S_k+Q_kV)$, so
\[
V\le C:=\max_{k\le120}\frac{S_k}{1-Q_k}.
\]

\emph{The bound.} Number the blocks of a period $1,\dots,b$, let $k_i$ be the number of gates
$\mathrm R$ of block $i$ and $p_i$ the probability that the chain holds $0$ before it; then
$p_{i+1}=\mathrm{out}_{k_i}(p_i)$, with $p_{b+1}=p_1$. For every real $\gamma$,
\[
T\,G_2(W,d)\le\sum_{i=1}^b\Bigl(\mathrm{cost}_{k_i}(p_i)+C\,X_{k_i}(p_i)\Bigr)
=\sum_{i=1}^b\Bigl(\mathrm{cost}_{k_i}(p_i)+C\,X_{k_i}(p_i)+\gamma\bigl(\mathrm{out}_{k_i}(p_i)-p_i\bigr)\Bigr),
\]
because the added terms cancel around the period. So if, for some $\gamma$ and some interval $I$
that contains every $p_i$,
\begin{equation}\label{eq:scanB}
\mathrm{cost}_k(p)+C\,X_k(p)+\gamma\bigl(\mathrm{out}_k(p)-p\bigr)\le\theta\,(k+1)
\qquad\text{for all }k\le120\text{ and }p\in I,
\end{equation}
then $G_2(W,d)\le\theta$, since $\sum_i(k_i+1)=T$. The left side is affine in $p$, so it is
enough to check the two ends of $I$. In the same way, for every real $g$ the density of ones of
the chain lies between the least and the largest value of
$\bigl(\mathrm{ones}_k(p)+g(\mathrm{out}_k(p)-p)\bigr)/(k+1)$ over $k\le120$ and $p\in I$. If $d$
is below the least value for some $g$, or above the largest for some $g$, then $d$ is not a
density of any word.

\emph{The interval.} Put $I_0=[0,1]$ and let $I_{n+1}$ be the least interval that contains
$\mathrm{out}_k(I_n)$ for all $k\le120$. By induction on $n$, the probability $p$ before a
block lies in $I_n$ for every $n$, because the chain is stationary and the block is preceded by
$n$ blocks. We use $I=I_{60}$.

\emph{The computation.} The program \texttt{scan\_check.c} computes enclosures of
$\mathrm{cost}_k$, $X_k$, $\mathrm{out}_k$, $\mathrm{ones}_k$ at $p=0$ and $p=1$, and of $S_k$
and $Q_k$, for $k\le120$, by following the joint law of the window and of the two copies of the
chain gate by gate; a state consists of the two values of the window and the two values of the
chain. It does this for an interval of noise levels and an interval of $d$, with the arithmetic
of \texttt{potential\_check.c}. For such a pair of intervals it accepts if $d$ is excluded as
above for one of $101$ values of $g$, or if~\eqref{eq:scanB} holds for one of $81$ values of
$\gamma$; otherwise it halves the interval of $d$, and if that becomes too short, the interval
of noise levels. It starts from sixteen intervals that cover all $d\in[0,1]$. The run
\texttt{scan\_check 0.44599 0.45501 0.999 120 0.0005} covers $[0.446,0.455]$ with
$\theta=0.999$, and the runs \texttt{scan\_check 0.45499 0.50001 0.99 120 0.001} and
\texttt{scan\_check 0.49999 1.0 0.99 120 0.01} cover $[0.455,1)$ with $\theta=0.99$; they print
\texttt{PROVED}, with $182$ intervals of noise levels in all, $6718$ boxes in which $d$ is excluded and $7162$ in which the bound is checked (Section~\ref{sec:checks}).
\end{proof}

\begin{remark}[what this shows, and an exact computation]\label{rem:scan}
\begin{enumerate}[label=(\roman*),leftmargin=2.4em]
\item \emph{The bound is that of independent reads.} For the word $\mathrm U$,
$G_2=\lambda^2$, which is $0.9623$ at noise $0.455$ and reaches $1$ at $4/9$. The same run with
$\theta=0.96$ on $[0.455,0.4552]$ is not proved, and with $\theta=0.965$ it is. So near noise
$0.455$ the proved bound is within $0.003$ of the value for independent reads. Away from it the
theorem proves less than is true: numerically, no word of period up to $16$ exceeds $\lambda^2$
at the noise levels $0.446$, $0.455$, $0.46$, $0.5$ and $0.7$ (the referee's enumeration of all
$65{,}535$ such words), and the largest value belongs to the words without $\mathrm R$; so
reading runs of consecutive positions does not seem to help at all. For the words $\mathrm{UR}$ and $\mathrm{URR}$ at noise
$0.46$, $G_2=0.8647$ and $0.8484$, against $\lambda^2=0.9446$. Remark~\ref{rem:dependent} shows a loss of the same kind for pairs
of neighbours, in a tower in which every layer scans.
\item \emph{What the family is.} Every second layer of the tower is plain. The proof uses this:
the window is then a Markov chain with a known law, and the joint law of a block is a finite
computation. A tower in which every layer scans the layer below is not covered. The number
$120$ is a limit of the computation, not of the method.
\item \emph{Dependence can help; an exact computation.} For a wider family the number $G_2$
is computed in the same way, in floating point and without a proof for all periods, by the
program \texttt{window\_designs.c}. A gate of the scanning layer may read the value at any position of a window,
or the target of the plain gate at that position, which is a position two layers down; it may
do so in any order, and with a second window of one position. At noise $0.46$, among all such
programs of period at most $5$ with a window of at most three positions, and of period at most
$3$ with four, the largest value is $\lambda=0.9719$, attained by programs of period at most
$4$ that read targets only and so skip the plain layers; nothing comes near $1$. But the program that reads position
$1$ of a window, then the target at a second window, then the target at position $0$ of the
first, has $G_2=0.964489$, and the same program with three independent windows has
$G_2=0.962795$. So a design with dependent reads can spread a flipped position more than the
same design with independent reads. The gain is $0.0017$ here, the largest we found is of this
size, and it is below the resolution of the simulations of Remark~\ref{rem:dependent}. The
order matters: with the two reads of the first window next to each other the same three reads
give $G_2=0.9519$, or $0.9369$ if the target comes first, against $0.9628$; and the program of
period $2$ that reads position $1$ of a window and then the target at position $0$ has
$G_2=0.9397$, against $0.9582$ for independent reads. These are floating-point values in the
ideal model, not proved statements; the referee confirmed the gain with a second exact program
and with a simulation of single flips ($0.0018\pm0.0001$).
\item \emph{What this means for all bases.} By (iii), a program with dependent reads can
spread more than the same program with independent reads, so an argument that compares the two
read by read cannot work. What we have not seen is a program with dependent reads above the
best program with independent reads: in the family of (iii) every value is at most $\lambda$.
The method of the theorem consists of blocks that start from an independent window, a bound for
the differences that leave a block, and a term that cancels around the period. We expect it to
apply to other finite sets of block types, but we have carried it out only for the runs.
\end{enumerate}
\end{remark}

\section{Levels of NC and polynomial speedups}\label{sec:levels}

The results of Sections~\ref{sec:logspace}--\ref{sec:s24} are stated for NC and BPNC. Their
proofs give more. The reductions of Theorems~\ref{thm:hard4} and~\ref{thm:hards24} add depth
$O(\log^2n)$ to the algorithm they use, and they map a circuit of $s$ gates to an instance of
length $O(s\log s)$. The
iterations above $\eps_1$ do only a logarithmic factor more work than sequential evaluation. This
section records the consequences for the levels $\mathrm{NC}^k$ and for polynomial parallel
speedups in the sense of Kruskal, Rudolph and Snir~\cite{krs}: parallel time $t^{1-\kappa}$ for a
problem of sequential time $t$, with work within a constant factor of $t$ (their class EP), a
polylogarithmic factor (AP) or a polynomial factor (SP); their classes ENC and ANC ask for
polylogarithmic time with the same work. We use these descriptions, not the exact definitions
of~\cite{krs}, which are stated for deterministic algorithms and the worst case. Whether P
has such speedups is usually studied through circuit evaluation. Nothing in
this section bears on $\mathrm P$ versus $\mathrm{NC}$, $\mathrm{BPNC}$ or SP: each statement says
that smoothed instances are easy for some class, or that they are as hard for it as circuit
evaluation.

\subsection{Each level of NC}

For $k\ge1$, a \emph{randomized $\mathrm{NC}^k$ family} is a randomized NC family (the definition
before Lemma~\ref{lem:errorless}) of depth $O(\log^kn)$, and $\mathrm{BPNC}^k$ is the class of languages
decided by such families with two-sided error at most $1/3$.

\begin{theorem}[the gap at each level]\label{thm:levels}
Let $k\ge2$ and $K\ge1$.
\begin{enumerate}[label=(\alph*),leftmargin=2.4em]
\item If for some $\eps\le\eps_S$ a randomized $\mathrm{NC}^k$ family outputs the trace with
probability at least $n^{-K}$ on the $\eps$-smoothing of every base of every length $n\ge n_0$,
then $\mathrm P\subseteq\mathrm{BPNC}^k$.
\item If $\mathrm P\subseteq\mathrm{BPNC}^k$, there is an errorless randomized $\mathrm{NC}^k$
family that outputs the trace with probability at least $1-2^{-n}$ on every parent array.
\item For every $\eps>\eps_1$, with $\eps_1<0.7288$ the level of Theorem~\ref{thm:allbases}, there
is a logspace-uniform family of deterministic circuits of polynomial size and depth $O(\log^2n)$
that outputs the trace or $\bot$ on every parent array, and $\bot$ with probability at most $3/n$
on the $\eps$-smoothing of every base.
\end{enumerate}
Consequently, let $E^{(k)}_K$ be the set of $\eps\in[0,1]$ at which some randomized
$\mathrm{NC}^k$ family outputs the trace with probability at least $n^{-K}$ on the
$\eps$-smoothing of every base, for all large $n$. Then $E^{(k)}_K$ is an up-set that contains
$(\eps_1,1]$, and $\inf E^{(k)}_K\in\{0\}\cup[\eps_S,\eps_1]$, with $\inf E^{(k)}_K=0$ if and only if
$\mathrm P\subseteq\mathrm{BPNC}^k$.
\end{theorem}

\begin{proof}
(a) Follow the proof of Theorem~\ref{thm:hards24}, which is that of Theorem~\ref{thm:hard} with
$\mathrm{Prog}_a$ at the noise level $\eps_S$, and count depth. The check of Lemma~\ref{lem:errorless} has depth
$O(\log n)$. The compiler is computable in logspace, so by logspace-uniform circuits of depth
$O(\log^2n)$, because $\mathrm L\subseteq\mathrm{AC}^1\subseteq\mathrm{NC}^2$ uniformly. The
sampling of the noise, the votes and the final majority of polynomially many votes are computable
in logspace too, hence in the same depth. A trial runs the family on an instance of length $n\le s^2$
(Lemma~\ref{lem:paramprog}(iii)), in depth $O(\log^kn)=O(\log^ks)$, and the trials run in parallel.
So evaluating NOR circuits is in $\mathrm{BPNC}^k$. That problem is P-complete under logspace
reductions, and for $k\ge2$ the class $\mathrm{BPNC}^k$ is closed under them: a logspace
reduction followed by a circuit of depth $O(\log^kn)$ is a circuit of depth
$O(\log^2n+\log^kn)$.

(b) As in the proof of Corollary~\ref{cor:gap}(a): the majority of $36(n+1)$ runs of a
$\mathrm{BPNC}^k$ algorithm for each bit, followed by the check of Lemma~\ref{lem:errorless}, adds
depth $O(\log n)$.

(c) For $\eps>3/4$ the machine of Theorem~\ref{thm:logspace} works in logspace; its output is
computed by logspace-uniform circuits of depth $O(\log^2n)$, again because
$\mathrm L\subseteq\mathrm{AC}^1\subseteq\mathrm{NC}^2$. For $\eps_1<\eps\le3/4$ use
Corollary~\ref{cor:upperend}: $O(\log n)$ applications of $G$, each in $\mathrm{AC}^0$, and the test
$G(y)=y$; the output is $\bot$ with probability at most $2/n$.

For the last claims, $E^{(k)}_K$ is an up-set by Lemma~\ref{lem:errorless} and
Lemma~\ref{gb:comp}; it contains $(\eps_1,1]$ by (c), as $k\ge2$; it is $[0,1]$ if
$\mathrm P\subseteq\mathrm{BPNC}^k$, by (b); and otherwise it contains no $\eps\le\eps_S$, by (a).
\end{proof}

\begin{remark}[the level $k=1$, and the short pointers]\label{rem:levels}
(i) For $k=1$, part (a) would need the reduction in $\mathrm{NC}^1$. Circuit evaluation is
P-complete under $\mathrm{NC}^1$ reductions~\cite{ghr}, and the NOR version follows by replacing
gates locally; the sampling, the votes, the majority and the check are in $\mathrm{NC}^1$. What
remains is to write out the compiler in $\mathrm{NC}^1$, which is routine: all positions of
$\mathrm{Latch}_a(C)$ are explicit arithmetic expressions in the wire, the stage, the rail, the
block, the gate and the coins, and the same holds for $\mathrm{Prog}_a(C)$. We have not done it. Even then the picture at $k=1$ would be
incomplete, because the algorithms above $\eps_1$ are in $\mathrm{AC}^1$, and above $3/4$ in logspace,
neither of which is known to be in $\mathrm{NC}^1$.

(ii) For bases with short pointers and zero gates, Theorem~\ref{thm:short} gives $O(\log^3n)$
applications of $G$ with probability at least $1-3/n$, each a prefix computation of constant depth
with unbounded fan-in. With the round count capped and the output checked, these instances are
solved errorlessly in $\mathrm{AC}^3$, and so in $\mathrm{NC}^4$, at every positive noise level;
a bound of $O(\log^2n)$ rounds would put them in $\mathrm{NC}^3$.
\end{remark}

\subsection{Work}

For statements about work we use a CREW PRAM with words of $O(\log n)$ bits, unit-cost addition,
subtraction, multiplication, comparison and shifts, and a unit-cost random word. The \emph{time}
of an algorithm is the number of parallel steps and its \emph{work} the total number of
operations. A NOR circuit with $s$ gates is given by the arrays $(u_t)$ and $(v_t)$; it is
evaluated sequentially in time $\Theta(s)$, and a parent array of length $n$ in time $\Theta(n)$.

\begin{theorem}[work]\label{thm:work}
\begin{enumerate}[label=(\alph*),leftmargin=2.4em]
\item Let $\eps>\eps_1$. There is a deterministic algorithm that outputs the trace on every parent
array of length $n$, with work $O(n\log n)$ and time $O(n)$ in the worst case, and that on the
$\eps$-smoothing of every base of length $n$ runs in time $O(\log^2n)$ with probability at least
$1-2/n$. Its expected time is $O(\log^2n)$.
\item Let $\eps\le\eps_S$ and $0<p\le1$, and let $T(N)$ and $W(N)$ be nondecreasing. Suppose a
randomized algorithm outputs the trace with probability at least $p$ on the $\eps$-smoothing of
every base of every length $N\ge N_0$, in time at most $T(N)$ and work at most $W(N)$. Then there
is a constant $c_\eps$ such that NOR circuits with $s$ gates are evaluated with error probability
at most $1/s$, for all large $s$, by a randomized algorithm with time $T(N_s)+O(\log s)$ and work
$O\bigl((W(N_s)+N_s)\log(s)/p\bigr)$, where $N_s\le c_\eps s\log s$.
\end{enumerate}
\end{theorem}

\begin{proof}
(a) Let $C$ be the constant of Theorem~\ref{thm:allbases} for $\eps$: on the $\eps$-smoothing of
every base the relevance-pruned iteration finishes within $C\ln(n+1)$ rounds with probability at
least $1-2/n$. Apply $G$ of Definition~\ref{def:jacobi} at most $\lceil C\ln(n+1)\rceil+1$
times, starting from $y^{(0)}$, and stop at the first $y$ with $G(y)=y$; then output $y$, which is
the trace by Proposition~\ref{prop:jacobi}(a). If no fixed point is reached, evaluate the trace
sequentially (time and work $O(n)$). One application of $G$ is done as follows: mark the gates $i$
that are not NOT gates and have $y_{j_i}=1$ (time $O(1)$, work $O(n)$, with concurrent reads);
compute for every $i$ the last marked gate $\ell_i\le i$, a prefix maximum (time $O(\log n)$, work
$O(n)$); put $z_i=(i-\ell_i)\bmod2$, or $(i+1)\bmod2$ if there is no marked gate up to $i$; and
test $z=y$, an OR of $n$ bits (time $O(\log n)$, work $O(n)$). So the work is $O(n\log n)$ on every
parent array. By Proposition~\ref{prop:jacobi} the iteration applies $G$ at most once more than
the relevance-pruned iteration needs rounds, so with probability at least $1-2/n$ it reaches the
fixed point within the cap, in time $O(\log^2n)$. The expected time is at most
$O(\log^2n)+\frac2n\,O(n)$.

(b) Make the algorithm errorless by the check of Lemma~\ref{lem:errorless}, which costs time
$O(\log N)$ and work $O(N)$. Let first $\eps<4/9$, and fix a dyadic rational $\eps^*$ with $\eps\le\eps^*<4/9$ and $\eps^*>0$. By
Lemma~\ref{gb:comp} the errorless algorithm succeeds with probability at least $p$ also at noise
$\eps^*$, on every base. Take $\mathrm{Latch}_a$ for $\eps^*$ with $a=3$, and let $N_s$ be its
length $n$; then $N_s+1\le Z<2Z_0M(a+2)^3s\log_2s$ by Lemma~\ref{lem:param4}(i) and the proof of
Lemma~\ref{lem:param4}(iii).

\emph{A trial} builds $\mathrm{Latch}_3(C)$ with fresh coins, draws an approximate
$\eps^*$-smoothing of it, runs the errorless algorithm and votes as in the proof of
Theorem~\ref{thm:hard4}. The global parameters ($m$, $Z$, $b$, $\mathrm{pre}$) are computed once,
in time $O(\log s)$. The preamble is written by one processor per position. Every other position is
written by one processor for its tuple (wire, stage, rail, block, position in the block), in the
order of Definition~\ref{def:latch}: the constants have the stages $1,\dots,J$ and come first, and
within a wire the order is $P_0,N_0,P_1,N_1,\dots$. Each processor computes its position and
pointer with $O(1)$ word operations from the arrays $(u_t),(v_t)$ and its coins, so this costs
time $O(1)$ and work $O(N_s)$. For the smoothing, $\pi_i$ compares a random integer of $w$ bits with
the threshold $\eps^*2^w$, which is exact because $\eps^*$ is dyadic and $w$ is large, and
$\rho_i=\lfloor ri/2^w\rfloor$ for a random integer $r$ of $w=\lceil 4\log_2(N_s+1)\rceil$ bits
(two random words, if a word is shorter), which is within $i2^{-w}$ of uniform; so the total
sampling error is $\nu\le N_s^{-2}$, at time $O(1)$ and work $O(N_s)$. The vote counts the copies of
$P_J(g_s)$ that hold $1$, of which there are $O(\log s)$, at time $O(\log s)$ and work $O(\log s)$.
As in that proof, a trial abstains with probability at most $1-p+\nu$, and votes wrongly with
probability at most $\delta+\nu$, where $\delta=s^{-3}+(n+1)^{-4}$ by Corollary~\ref{cor:latch}.

\emph{The algorithm} runs $L=\lceil2\ln(s)/p\rceil+1$ independent trials in parallel and outputs
the vote of the first trial that does not abstain, found by a prefix computation over the trials.
It errs only if all trials abstain, with probability at most $(1-p+\nu)^L\le e^{-(p-\nu)L}$, or
if some trial votes wrongly, with probability at most $L(\delta+\nu)$. For large $s$ both are at
most $1/(2s)$. The time is $T(N_s)+O(\log N_s)$ and the work
$L\cdot O(W(N_s)+N_s)$, and $\log N_s=O(\log s)$. The constant $c_\eps$ is at least $\Lambda_0$,
which is astronomical (Remark~\ref{rem:latch}(v)).

For $4/9\le\eps\le\eps_S$ the same proof works with $\eps^*=\eps_S$, which is dyadic, with
$\mathrm{Prog}_3$ of Section~\ref{sec:s24} in place of $\mathrm{Latch}_3$, and with
Lemma~\ref{lem:paramprog} and Corollary~\ref{cor:prog} in place of Lemma~\ref{lem:param4} and
Corollary~\ref{cor:latch}. The positions of $\mathrm{Prog}_3(C)$ are written by one processor for
each tuple (wire, stage, block, position in the block), in the order of Definition~\ref{def:prog},
and the vote counts the copies in $R_J(g_s,8)$ that hold $1$.
\end{proof}

\begin{corollary}[polynomial speedups]\label{cor:speedup}
Let $\eps\le\eps_S$, $\kappa\in(0,1)$ and $c,\gamma\ge0$, and suppose that the $\eps$-smoothed
instances of every base are solved with probability at least $1/2$ by a randomized algorithm in
time $O(N^{1-\kappa})$.
\begin{enumerate}[label=(\alph*),leftmargin=2.4em]
\item If its work is $O(N\log^cN)$, then NOR circuits with $s$ gates are evaluated, with error
probability at most $1/s$, in time $O(s^{1-\kappa}\log s)$ and work $O(s\log^{c+2}s)$.
\item If its work is $O(N^{1+\gamma})$, the time is the same and the work is
$O(s^{1+\gamma}\log^{2+\gamma}s)$.
\end{enumerate}
\end{corollary}

\begin{proof}
Theorem~\ref{thm:work}(b) with $p=1/2$ and $N_s=O(s\log s)$.
\end{proof}

\begin{remark}[what this says about speedups]\label{rem:speedup}
\begin{enumerate}[label=(\roman*),leftmargin=2.4em]
\item \emph{Up to $0.4468$.} In the terms of~\cite{krs}: a randomized algorithm of the kind of EP or
AP that succeeds with constant probability on the smoothed instances of every base at some noise
level $\eps\le\eps_S$ gives a randomized algorithm of the kind of AP, with error $1/s$, for NOR circuit
evaluation, and one of the kind of SP gives one of the kind of SP. So up to $\eps_S$ the noise does
not make Sequential NOR easier for polynomial speedups than circuit evaluation is, up to
polylogarithmic factors in time and work, randomization, and constants that depend on $\eps$ and
are astronomical. This uses that the embedding has length $O(s\log s)$.
\item \emph{Above $\eps_1$.} Theorem~\ref{thm:work}(a) is of the kind of ANC in~\cite{krs}:
polylogarithmic time with a logarithmic factor of extra work, errorless, for every base; the work
bound holds in the worst case and the time bound with probability $1-2/n$ over the noise. It
implies speedups of the kinds of AP and SP, but not of EP, because the work is $\Theta(n\log n)$
on typical instances. The same holds on the random block bases over zero gates of
Theorem~\ref{thm:blockabove} for every noise level above $4/9$, with probability at least $1-5/n$
over the base and the noise.
\item \emph{Block bases: a speedup without NC.} On every block base with blocks of $m$ positions,
at every noise level, the table--Jacobi iteration with reach $2m$ applies $G_{2m}$ at most
$17\sqrt{\eps n}+35\ln n+6$ times with probability at least $1-n^{-4}$
(Corollary~\ref{cor:tableblock}), and at most $\lceil n/(2m+1)\rceil+1$ times on every parent array
(Theorem~\ref{thm:table}(c)). One application is a parallel prefix over $n$ maps stored as tables
with $2^{2m}$ entries; a composition of two maps is $2^{2m}$ look-ups. So these instances are
solved in time $O(\sqrt n\log n)$ with probability $1-n^{-4}$ over the noise, a polynomial speedup,
with work $O(n^{3/2}2^{2m})$. With $m=\lceil C\ln(n+1)\rceil$, as in Theorems~\ref{thm:block}
and~\ref{thm:blockzero}, the exponent of the work is $3/2+2C\ln2$, polynomial only formally since
$C$ is large. On the random block bases of those theorems, the Kleene--prefix iteration needs
$cn/\ln n$ rounds below $1-d^*$ (and below $\eps_K$ over zero gates), and the relevance-pruned
iteration below $d^*$ (and below $4/9$ over zero gates). There no algorithm of polylogarithmic
time is known, and no hardness either, because random block bases do not encode circuits. So,
among smoothed instances of Sequential NOR, random block bases at low noise have a proved
polynomial speedup and no known NC algorithm. Whether they are in NC, in randomized NC, or hard in
some sense is open.
\end{enumerate}
\end{remark}

\subsection{Two restatements: greedy independent sets and miswired circuits}\label{sec:restate}

The results on smoothed instances can be read as results on two older questions. Nothing new is
proved here beyond Lemma~\ref{lem:greedysteps}.

\paragraph{Greedy independent sets.} For a parent array $j$ of length $n$ let $\Gamma(j)$ be
the graph on $\{0,\dots,n\}$ with the edges $\{i-1,i\}$ and $\{j_i,i\}$, $1\le i\le n$: a
Hamiltonian path, and one more edge from each vertex to an earlier one. As in
Corollary~\ref{nw:cor-lfmis}, the lexicographically first maximal independent set of $\Gamma(j)$
in the natural order is $\{i:x_i=1\}$, where $x$ is the trace of $j$. For a general graph and
order this problem is P-complete~\cite{ghr}, and so is its restriction to the graphs
$\Gamma(j)$~\cite{okhotin}. The $\eps$-smoothing moves the second edge of each vertex,
independently with probability $\eps$, to a uniform earlier vertex. The \emph{parallel greedy
algorithm}~\cite{nw-crt87,nw-bfs12} works in steps: in each step every remaining vertex without
a remaining earlier neighbour joins the set, and these vertices and their neighbours are
removed. Let $s(i)$ be the step in which vertex $i$ joins or is removed. An earlier neighbour of
a vertex that joins is removed in an earlier step, and a vertex is removed in the first step in
which one of its earlier neighbours joins; a later neighbour cannot join before the vertex is
removed. So $s(0)=1$ and, with $N(i)=\{i-1,j_i\}$,
\[
s(i)=1+\max_{u\in N(i)}s(u)\ \text{ if }x_i=1,\qquad s(i)=\min\{s(u):u\in N(i),\ x_u=1\}\ \text{ if }x_i=0 .
\]

\begin{lemma}[greedy steps and the pruned depth]\label{lem:greedysteps}
For every parent array and every $i$, $D(i)\le2s(i)$. So the parallel greedy algorithm on
$\Gamma(j)$ needs at least $\frac12\max_iD(i)$ steps.
\end{lemma}

\begin{proof}
Put $\tau(0)=0$, and for $i\ge1$ let $\tau(i)=1+\max_{u\in N(i)}\tau(u)$ if $x_i=1$ and
$\tau(i)=1+\min\{\tau(u):u\in N(i),\ x_u=1\}$ if $x_i=0$. By induction on $i$,
$\tau(i)\le2s(i)-x_i$: this holds for $i=0$; if $x_i=1$ then
$\tau(i)\le1+\max_u2s(u)\le1+2(s(i)-1)$; and if $x_i=0$ and $u$ attains the minimum for $s$, then
$\tau(i)\le1+\tau(u)\le2s(u)=2s(i)$. Also by induction, $D(i)\le\tau(i)$. For $i=1$ both are $1$.
Let $i\ge2$ and use the cases of Lemma~\ref{kp:switch}(a). At a NOT gate
$\tau(i)\ge1+\tau(i-1)$, both when $x_{i-1}=1$ and when $x_{i-1}=0$, and $D(i)=D(i-1)$. At a
relevant reset $\tau(i)=1+\tau(j_i)\ge1+D(j_i)=D(i)$. At an irrelevant reset
$\tau(i)=1+\min(\tau(i-1),\tau(j_i))\ge\min(D(i-1),1+D(j_i))=D(i)$. At a relevant gate that is
neither a reset nor a NOT gate, $x_i=1$ and
$\tau(i)=1+\max(\tau(i-1),\tau(j_i))\ge\max(D(i-1),1+D(j_i))=D(i)$. In the remaining case
$x_{i-1}=1$ and $x_{j_i}=0$, so $\tau(i)=1+\tau(i-1)\ge D(i-1)=D(i)$.
\end{proof}

\begin{corollary}[smoothed greedy independent sets]\label{cor:smoothedmis}
Consider the lexicographically first maximal independent set of $\Gamma(\jt)$,
$\jt=S_\eps(j)$, in the natural order.
\begin{enumerate}[label=(\alph*),leftmargin=2.4em]
\item \emph{Hardness.} Let $\eps\le\eps_S=3661/8192$. If a randomized NC family outputs this set
with probability at least $n^{-K}$ for every base $j$ of every length $n\ge n_0$, then
$\mathrm P\subseteq\mathrm{BPNC}$.
\item \emph{Easiness.} For every $\eps>\eps_1=0.7288$ the set is computed for every base in
$O(\log n)$ rounds, each in $\mathrm{AC}^0$, with probability at least $1-2/n$, without error.
For $j=j^H$, where the second edges join the vertices at distance $H$, this holds from
$\eps=0.455$ on if $H\ge H_0$.
\item \emph{The parallel greedy algorithm.} For every $\eps<d^*=(3-\sqrt5)/2$ there is $c>0$ such
that on the smoothing of the random block bases of Theorem~\ref{thm:block} the parallel greedy
algorithm needs at least $cn/\ln n$ steps, with probability at least $1-3/n$.
\end{enumerate}
\end{corollary}

\begin{proof}
(a) is Theorem~\ref{thm:hards24}, (b) is Theorem~\ref{thm:allbases} and Theorem~\ref{thm:rows};
the relevance-pruned iteration records only correct values (Lemma~\ref{kp:pruned}(a)). (c) is
Theorem~\ref{thm:block}(b) and Lemma~\ref{lem:greedysteps}.
\end{proof}

\noindent For an arbitrary graph in a uniformly random order the parallel greedy algorithm needs
$O(\log n)$ steps~\cite{nw-bfs12,nw-fn18}, and for $G_{n,p}$ in the natural order polylogarithmic
expected time suffices~\cite{nw-crt87,nw-cf90}. The corollary describes what lies between the
worst case and these two: when a fraction $\eps$ of the edges that are not on the path is random
and the order is the natural one, the problem stays as hard as in the worst case up to
$\eps=0.4468$, is in NC for every graph of this form from $\eps=0.7288$ on, and the parallel greedy
algorithm itself is slow on some of these graphs for every $\eps<0.3819$. We do not know of
earlier results on this problem with a fixed order and partly random edges.

\begin{remark}[miswired circuits]\label{rem:miswired}
Reliable computation with unreliable components~\cite{gb-vonneumann,gb-pippenger} is usually
studied for gates whose outputs are flipped at random; for formulas of two-input gates the exact
noise threshold is known~\cite{evansschulman,unger}. The smoothing is a different fault: with
probability $\eps$ the second input wire of a gate is attached to a uniform earlier gate. In this
language the sections on hardness and Section~\ref{sec:designs} say the following.
\begin{enumerate}[label=(\roman*),leftmargin=2.4em]
\item \emph{Reliable computation is possible up to $\eps_S$.} For every $\eps\le\eps_S$ every NOR
circuit with $s$ gates has a layout as a parent array of length $O(s\log s)$, such that after the miswiring the value of every gate of the
circuit is the majority, in a fixed sense, of a bundle of positions of the trace, with high
probability (Theorems~\ref{thm:latch} and~\ref{thm:progembed}).
\item \emph{Layouts of the known kind fail from $0.455$ on.} A layered design whose reads are
independent of each other forgets its input at a geometric rate at every noise level from
$0.455$ on: its density map is a contraction (Theorem~\ref{thm:indep}), and the trace of the
smoothing of every layered base with wide bundles is then computed in $O(\log n)$ rounds
(Theorem~\ref{thm:layered}). This is unconditional. It concerns these layouts only, and the
ideal model for the first statement.
\item \emph{What is open.} Whether some layout computes reliably above $0.455$ is the question
whether the critical noise level is above $0.455$; by Theorem~\ref{thm:allbases} no layout can
make the evaluation of its own trace hard above $0.7288$. Unlike a flipped output, a miswired
input is not symmetric: it reads a one with the density of ones of the trace, which the layout
influences.
\end{enumerate}
\end{remark}

\begin{remark}[a fixed distribution of hard instances: what is missing]\label{rem:fixeddist}
The hardness theorems are about the worst base: the distribution $S_\eps(j)$ depends on the
circuit that $j$ encodes. A stronger statement would name one samplable distribution $\mathcal D_n$
for each $n$, not depending on any circuit, on which every NC algorithm fails unless
$\mathrm P\subseteq\mathrm{BPNC}$. We have no such statement, and we record why the two obvious
routes do not give it.
\begin{enumerate}[label=(\roman*),leftmargin=2.4em]
\item \emph{A fixed base.} Let $j^{(n)}$ be one base for each $n$, for instance the layout of a
universal circuit with a fixed input, and $\mathcal D_n=S_\eps(j^{(n)})$. For uniform algorithms
this reduces to a statement about one instance per length, which the assumption
$\mathrm P\not\subseteq\mathrm{BPNC}$ does not give. Against nonuniform algorithms the values
of the encoded circuit are advice, and the hardness proof says nothing.
\item \emph{A random self-reduction.} We know no reduction that maps a worst-case base to
instances of a fixed distribution. The usual reductions of this kind use polynomials of low
degree, and a P-complete problem with such a structure is not known.
\item \emph{What the noise could give.} For a fixed base the noise is the only randomness, and a
layout can use it: a gate whose pointer is aimed at a position that holds a one outputs a biased
coin, and bundles that read this gate turn the coin into a value that the layout computes with
reliably. This suggests, but we have not proved, that for $\eps\le\eps_S$ there is a fixed
sequence of bases whose smoothing is hard on average for NC if and only if circuit evaluation is
hard on average for NC under some samplable distribution of inputs. That would make the
smoothing of one explicit base a complete problem for this average-case question; it would not
derive the hardness from $\mathrm P\not\subseteq\mathrm{BPNC}$.
\end{enumerate}
\end{remark}

\section{Computations and checks}\label{sec:checks}

The scripts are in the folder \texttt{scripts/} of this paper's verification folder. The Python scripts use only
the standard library; \texttt{embed.c} is plain C.

\medskip\noindent\textbf{\texttt{logspace\_eval.py}.} Implements the evaluator of
Section~\ref{sec:logspace} with $\ghat=0$ and an explicit stack of offsets, checks the result
against the trace (all positions for $n=256$; four positions per instance otherwise), and measures
the largest stack in gamma-coded bits. Three instances per row.

\begin{center}
\begin{tabular}{lrrrrr}
\toprule
base, noise & $n$ & calls for one position & stack bits & frames & bits$/\log_2n$\\
\midrule
uniform ($\eps=1$) & 256 & 618 & 32 & 12 & 4.00\\
 & 1024 & 10{,}102 & 58 & 22 & 5.80\\
 & 4096 & 78{,}781 & 74 & 26 & 6.17\\
$j_i=i-2$, $\eps=0.9$ & 4096 & 5{,}017{,}011 & 112 & 43 & 9.33\\
$j_i=0$, $\eps=0.9$ & 4096 & 3{,}451 & 59 & 21 & 4.92\\
$j_i=\lfloor i/2\rfloor$, $\eps=0.9$ & 4096 & 69{,}128 & 76 & 26 & 6.33\\
$j_i=i-2$, $\eps=0.8$ & 4096 & 12{,}667{,}908 & 137 & 51 & 11.42\\
$j_i=\lfloor i/2\rfloor$, $\eps=0.8$ & 4096 & 11{,}584 & 52 & 20 & 4.33\\
\bottomrule
\end{tabular}
\end{center}

\noindent Each row reports the worst of the measured positions. The number of bits per frame stays
near $3$, as the weighted count predicts. These runs are a sanity check of the algorithm and of
the order of magnitude; they are not a test of the asymptotic statement.

\medskip\noindent\textbf{\texttt{check\_note\_constants.py}.} Asserts the numerical inequalities
used above: $0.0111\cdot65/\ln2>1$, $D(1/6\,\|\,0.1161)\ge0.0111$, $D(1/30\,\|\,1/64)\ge0.0077$,
$0.0077\cdot130\ge1$, $92/\ln2>130$, the gate count $413{,}340$, the three parts of
Lemma~\ref{lem:param} on a grid of $(a,s)$ up to $a=10^6$ and $s=10^{100}$, the inequality
$2\delta+3\nu\le n^{-K}/2$, the recursion $q_{10}\le0.1161$ at $\eps=1/64$, the two gadgets of
Lemma~\ref{gb:gadgets} on all inputs, and the constants of Remark~\ref{rem:constants}.

\medskip\noindent\textbf{\texttt{check\_embedding.py}.} Asserts the numerical claims of
Section~\ref{sec:tree}. The recursion for $Q_0,Q_1$ is evaluated in exact rational arithmetic,
rounded upwards to $40$ digits after every step, so the values $0.2311$, $0.0900$ (noise $1/10$)
and $0.2099$, $0.0493$ (noise $3/20$) are rigorous upper bounds. The script also asserts the four
divergences, the block sizes, the constants of Lemma~\ref{lem:zeropad}, and
Lemma~\ref{lem:param2} on the same grid of $(a,s)$ as before.

\medskip\noindent\textbf{\texttt{domination\_test.py}.} Tests Lemma~\ref{lem:dom} pathwise. It
compiles random circuits with $d\in\{1,3\}$, smooths them at noise between $0$ and $1$, and checks
at every node that a wrong output implies the event $D$, with every flag replaced by the smallest
event it is known to contain (a misread for $B_g$; a misread of a $0$ for $A_g$). In $1500$
instances it checked $404{,}200$ nodes, of which $98{,}061$ were wrong, and found no violation. At
noise $0$ no node is wrong.

\medskip\noindent\textbf{\texttt{embed.c}.} Builds $\mathrm{Tree}(C)$ for a random circuit, smooths
it exactly as in Section~\ref{sec:setting}, evaluates the trace and reports the bad cells per
bundle. Cells of depth $11$, no padding, $m=31$, $s=30$ ($n\approx1.1\cdot10^7$), two smoothings
per row. The columns give the fraction of bad cells per bundle (its mean over the gates of value
$0$, over the gates of value $1$, and its largest value) and the number of bundles without a
correct majority:

\begin{center}
\begin{tabular}{lrrrr}
\toprule
noise & mean, value $0$ & mean, value $1$ & largest & majority lost\\
\midrule
0.05 & 0.027 & 0.018 & 0.129 & 0 of 60\\
0.10 & 0.058 & 0.040 & 0.194 & 0 of 60\\
0.125 & 0.077 & 0.060 & 0.226 & 0 of 60\\
0.15 & 0.106 & 0.089 & 0.290 & 0 of 60\\
0.175 & 0.131 & 0.095 & 0.355 & 0 of 60\\
0.20 & 0.152 & 0.122 & 0.323 & 0 of 60\\
0.25 & 0.317 & 0.265 & 0.645 & 5 of 60\\
\bottomrule
\end{tabular}
\end{center}

\noindent With the cells of Corollary~\ref{cor:tenth} (depth $15$; $m=31$, $s=8$, four smoothings
at noise $0.10$) the mean fraction of bad cells was at most $0.05$ and the largest $0.097$. With
zero padding (ratio $20$, depth $13$, $m=15$, $s=4$, three smoothings at noise $0.15$) the density
of ones was $0.128$ and the largest fraction of bad cells $0.133$. The bundles here are much
smaller than the corollaries require, so these runs test the construction, not the tail bounds.

\medskip\noindent\textbf{\texttt{lean/SeqNORTree.lean}.} A Lean~4 file (core library only, no
Mathlib) that proves Lemmas~\ref{lem:exact} and~\ref{lem:dom} in an abstract form. A block is a
leaf or a node together with what happened to it: for a leaf the value of the copy and whether its
two gates misread, for a node whether its gate misreads. The value of a block is computed from the
value before its first gate exactly as in the smoothed instance, with
$x_g=\NOR(x_{g-1},y_g)$ at every gate, and positions are forgotten. The theorem
\texttt{domination} states that a block whose last gate differs from its intended value has its
event $D$, where a flag $B_g$ is replaced by ``gate $g$ misreads'' and a flag $A_g$ by ``gate $g$
misreads a $0$''; it covers nodes whose two blocks have different intended values, hence the top
gate of a cell. The theorems \texttt{exact}, \texttt{want\_tree} and \texttt{cell\_exact} are
Lemma~\ref{lem:exact}, and \texttt{dom\_clean} shows that $D$ fails on every block without
misreads and bad copies. The file compiles with Lean~4.35.0 without \texttt{sorry}; the axioms
used are \texttt{propext}, \texttt{Quot.sound} and \texttt{Classical.choice}. What is not
formalized: the layout by positions, Lemma~\ref{lem:misread}, and every probabilistic step.

\medskip\noindent\textbf{Section~\ref{sec:tree2}.} \texttt{lean/SeqNORTree2.lean} proves the leaf,
block and cell cases of Lemma~\ref{lem:dom2} and the exactness of Lemma~\ref{lem:exact2} in the
abstract model of \texttt{SeqNORTree.lean}; it compiles with Lean~4.35.0 without \texttt{sorry},
using at most the axioms \texttt{propext} and \texttt{Quot.sound}. \texttt{check\_tree2.py} asserts,
in exact rational arithmetic with upward rounding, the bound on the error of
Lemma~\ref{lem:dens2}, the stage bound $Q_k\le(5/\tau^2)^{k+1}W$ on examples, the enclosure of
$\varpi(1/4)$, the values $Q_0\le0.13848$ and $Q_1\le0.15904$ at depth $51$ with the rounded $A$ and
$B$ of Corollary~\ref{cor:twentythree}, the Chernoff
constants, the monotonicity of $A(\eps)$ and $B(\eps)$ on a grid up to $1/4$, the recursion at five
smaller noise levels, the inequalities of Lemma~\ref{lem:param3}, and the cell size; it also
certifies the parameters for noise $23/100$ quoted in Remark~\ref{rem:limits2}(ii). \texttt{tree2\_explore.py}, run with the argument \texttt{note}, computes every threshold quoted in
Remarks~\ref{rem:limits} and~\ref{rem:limits2}. \texttt{embed2.c} simulates $\mathrm{Tree}'(C)$ for a random NOR
circuit ($s=24$ gates, depth $7$, padding ratio $99$):

\begin{center}
\begin{tabular}{rrrrr}
\toprule
noise & bundle & density of ones over cell gates & largest bad fraction (0 / 1) & output\\
\midrule
0.15 & 15 & 0.1175 -- 0.1207 & 0.40 / 0.33 & correct\\
0.20 & 15 & 0.1463 -- 0.1490 & 0.20 / 0.27 & correct\\
0.23 & 31 & 0.1611 -- 0.1639 & 0.19 / 0.23 & correct\\
0.23 & 31 & 0.1617 -- 0.1639 & 0.19 / 0.16 & correct\\
0.25 & 31 & 0.1719 -- 0.1739 & 0.48 / 0.58 & correct\\
0.28 & 15 & 0.1862 -- 0.1884 & 0.40 / 0.33 & correct\\
0.32 & 15 & 0.2041 -- 0.2061 & 0.33 / 0.53 & correct\\
\bottomrule
\end{tabular}
\end{center}

\noindent For comparison, $\varpi(\eps)=0.117$, $0.146$, $0.162$, $0.172$, $0.186$ and $0.203$ at these
noise levels. At depth $7$ and noise $0.23$, inputs with $10\%$ bad cells give $Q_0\approx0.153$ and
$Q_1\approx0.192$ by the recursion (the second figure depends on the density error used, in the
third digit), against measured means of $0.06$ to $0.09$ and $0.09$ to $0.15$.

\medskip\noindent\textbf{Referee pass on Section~\ref{sec:tree2}.} The automated referee of
Section~\ref{sec:fixed} read Section~\ref{sec:tree2} in its version for noise $23/100$ and reported
no error in the proofs. It checked Lemma~\ref{lem:dens2} line by line and by simulation on twelve
padded bases ($n=2\cdot10^6$, padding ratios $3$ to $999$, noise $0.01$ to $0.25$), compiled
\texttt{SeqNORTree2.lean} and compared its model with Definition~\ref{def:treeprime} gate by gate,
checked domination by brute force over all flag patterns and pathwise on $200$ small instances,
recomputed the recursion in its own rational arithmetic, and checked Lemma~\ref{lem:param3}. It
noted that the recursion closes at noise $1/4$, which led to the present statement, and it pointed
out one rounded constant and four sentences of the remarks, which are corrected here. A second
pass checked the version for noise $1/4$: the recursion and Chernoff constants, Lemma~\ref{lem:param3}
for $d=51$, the error of Lemma~\ref{lem:dens2} and the monotonicity of $A$, $B$ up to $1/4$, and every
threshold quoted in the remarks; the corrections it pointed out (the recursion must be evaluated at
the rounded $A$, $B$ that are printed, and four figures in the remarks) are made here.

\medskip\noindent\textbf{\texttt{adaptive\_adversary.py}.} Simulates the adaptive bases of
Proposition~\ref{prop:adaptive}, with the depths computed by the recursions of
Lemma~\ref{fact:depths}:

\begin{center}
\begin{tabular}{lrrrr}
\toprule
noise & $n$ & $\max_id(i)/n$ & $\max_iD(i)/n$ & bound of part (a), $/n$\\
\midrule
0.50 & $10^6$ & 0.2614 & 0.1575 & 0.1000\\
0.80 & $10^6$ & 0.1256 & 0.0712 & 0.0571\\
0.90 & $10^6$ & 0.0688 & 0.0370 & 0.0310\\
0.95 & $10^6$ & 0.0365 & 0.0187 & 0.0161\\
\bottomrule
\end{tabular}
\end{center}

\noindent The ratios at $n=10^4$ and $10^5$ agree with these to two digits. The first column of
depths is for the base built against the Kleene--prefix iteration, the second for the base built
against the relevance-pruned one. For the oblivious base $j_i=i-2$ at noise $0.9$ the Kleene--prefix
depth was $48$, $71$ and $98$ at $n=10^4,10^5,10^6$.

\medskip\noindent\textbf{\texttt{depths.c}, \texttt{theta.c}, \texttt{jacobi\_check.py}.}
The first computes, on one smoothing, the depths $d$, $D_{\mathrm{last}}$ and $D$ by the recursions
of Section~\ref{sec:iterations} and the number of applications of $G$ after which the prefix--Jacobi iterate equals the
trace. For the base with distances uniform in $[64,128)$ and no filler:

\begin{center}
\begin{tabular}{lrrrrr}
\toprule
noise & $n$ & $\max d$ & $\max D_{\mathrm{last}}$ & $\max D$ & prefix--Jacobi\\
\midrule
0.60 & $2^{16}$ & 403 & 179 & 89 & 75\\
 & $2^{20}$ & 1756 & 388 & 165 & 122\\
0.66 & $2^{16}$ & 217 & 125 & 74 & 61\\
 & $2^{20}$ & 674 & 253 & 110 & 104\\
0.70 & $2^{16}$ & 176 & 105 & 71 & 56\\
 & $2^{20}$ & 386 & 192 & 103 & 88\\
0.73 & $2^{16}$ & 163 & 92 & 56 & 55\\
 & $2^{20}$ & 284 & 159 & 82 & 77\\
\bottomrule
\end{tabular}
\end{center}

\noindent Two sizes do not give a growth rate, and near its threshold the Kleene--prefix depth
has a large constant in front of $\log n$. The second script produces the table of
Remark~\ref{rem:needs}; the third is the exhaustive check of Proposition~\ref{prop:jacobi}.

\medskip\noindent\textbf{\texttt{fd\_bound.py}, \texttt{fd\_check.py}.} The first evaluates
$\Theta_H$ of Definition~\ref{def:Tm} (the sum over $m$ is cut where $r^{m}<10^{-16}$), prints the
tables of Section~\ref{sec:fixed}, asserts the six numerical claims of
Theorem~\ref{thm:fixed}, and spot-checks that $\Theta_H$ is nonincreasing in $H$. For the
constants of Corollary~\ref{cor:fdadm} it iterates $\ell_s,\mu_s$ forty times. The second
simulates $S_\eps(j^H)$ at $n=10^6$ and measures the mean number of gates $k\le P$ that keep their
base pointer with no reset in $(k,P]$, split by the value at $k-H$; this is the quantity that
$(1-\eps)\sum_mT_m$ bounds. A third script, \texttt{fd\_step.py}, tests Lemma~\ref{lem:onestep}
for each $m<10$ separately at $n=1.5\cdot10^6$: the measured probability is between $0.3$ and
$0.95$ of the bound for $1\le m<10$, and equals it within sampling error at $m=0$, where the
bound is an equality.

\begin{center}
\begin{tabular}{rrrrrr}
\toprule
$H$ & noise & density of ones & measured, $x_{k-H}=0$ & measured, $x_{k-H}=1$ & bounds\\
\midrule
3 & 0.60 & 0.408 & 0.835 & 1.031 & 1.822, 2.025\\
3 & 0.70 & 0.399 & 0.672 & 0.787 & 1.121, 1.223\\
3 & 0.74 & 0.396 & 0.593 & 0.681 & 0.911, 0.984\\
8 & 0.60 & 0.390 & 0.906 & 1.095 & 1.486, 1.647\\
8 & 0.70 & 0.386 & 0.714 & 0.820 & 0.978, 1.067\\
8 & 0.74 & 0.385 & 0.630 & 0.709 & 0.813, 0.880\\
\bottomrule
\end{tabular}
\end{center}

\noindent The bounds in this table use the constants of Lemma~\ref{fact:densall}.

\medskip\noindent\textbf{\texttt{fd\_density.py}.} Computes $\ell_{40}$ and $\mu_{40}$ and
simulates $S_\eps(j^H)$ at $n=4\cdot10^5$ for $H\in\{2,3,8,32\}$ and
$\eps\in\{0.1,0.3,0.5,0.7,0.9,1\}$, recording the smallest and largest prefix density $\psi_g$
over $g\ge20000$. In all $24$ cases these lie between the two bounds of
Lemma~\ref{lem:fddens}, up to $0.002$ at $\eps=1$, where the bounds coincide. Some rows:

\begin{center}
\begin{tabular}{rrrrr}
\toprule
$H$ & noise & $\ell_{40}$ & $\mu_{40}$ & measured $\psi_g$\\
\midrule
2 & 0.1 & 0.3335 & 0.4915 & 0.3370 to 0.3375\\
2 & 0.7 & 0.3496 & 0.4303 & 0.3633 to 0.3650\\
3 & 0.1 & 0.3335 & 0.4915 & 0.4758 to 0.4773\\
3 & 0.7 & 0.3496 & 0.4303 & 0.3992 to 0.4007\\
8 & 0.7 & 0.3496 & 0.4303 & 0.3841 to 0.3871\\
32 & 0.5 & 0.3401 & 0.4535 & 0.3961 to 0.3995\\
\bottomrule
\end{tabular}
\end{center}

\medskip\noindent\textbf{\texttt{ord\_bound.py}.} For Section~\ref{sec:ordered}. It evaluates
$\Theta^{\mathrm{ord}}_D$, asserts the two numerical claims of Theorem~\ref{thm:ordered}, and
checks the inequality $\eps c\le r(1-\gamma)$ at $200{,}000$ random parameter values with
$\eps\le3/4$, $c\le1/2$ and $\beta<1-c$. It tests Lemma~\ref{lem:sumspacing} against the
operators $B_d$ themselves, at six parameter triples: all sequences of at most $8$ distances from
$\{1,2,3,4,9\}$, and $3000$ random sequences of length $300$ per triple; the largest partial sum
found is $0.9815$ of the bound. Finally it simulates three ordered bases at $n=6\cdot10^5$ and
noise $0.60$ and $0.74$ (even gates reading $32$ below with NOT gates between; $j_{2t}=t$,
$j_{2t+1}=2t$; and random increments $1$ to $3$ at distance at least $32$): the measured number
of base gates in a window is between $0.38$ and $0.67$, far below $\Theta^{\mathrm{ord}}_{32}$,
which is $0.968$ at noise $0.74$ and $2.01$ at $0.60$.

\medskip\noindent\textbf{\texttt{ties\_check.py}.} For the parts of Section~\ref{sec:ordered}
that concern gates reading the same position. It checks the inequalities (b) and (c) in the proof
of Lemma~\ref{lem:sumspacing} at $300{,}000$ random parameter values with $\eps\le3/4$,
$c\le1/2$ and $\beta\le\min(2/3,1-c)$. It tests the lemma itself with the exact recursion
$y_{u+1}=M^{d_u}(\varphi^{\mu_u}y_u)$ at six parameter triples: all sequences of at most $6$ pairs
$(d,\mu)$ from a list of seven, and $4000$ random sequences of length $200$ per triple; the
largest sum found is $0.968$ of the bound. It also simulates two bases with ties at
$n=6\cdot10^5$, $j_i=\lfloor i/2\rfloor$ and blocks of $64$ gates reading one position, and
counts for a sample of positions $P_0$ the gates of the window that keep their pointer, and
those among them that are the highest with that pointer. At noise $0.74$ the two counts are
$0.70$ and $0.65$ for the first base and $0.77$ and $0.51$ for the second; at noise $0.60$ they
are $1.09$ and $0.97$, and $1.35$ and $0.67$.

\medskip\noindent\textbf{\texttt{typed\_model.py}.} Exploration for
Remark~\ref{rem:below}(iii), not a check of a proof. It runs a value iteration on a grid of $96$
directions of the pair of weights, with up to $8$ pending offsets and window lengths up to $28$,
and prints the growth factor of the one-gate model ($1.046$ at noise $0.74$; with no credit the
same code gives $1.064$, against the exact $1.061$, so its discretization error is about
$0.003$). It also evaluates the $2\times2$ recursion of the remark: $1.071$, $1.012$, $0.955$
and $0.900$ at noise $0.72$, $0.73$, $0.74$ and $0.75$.

\medskip\noindent\textbf{\texttt{newton.c}.} Runs the GF(2) Newton iteration of Section~\ref{sec:newton} from the
standard start on smoothed bases and counts rounds until the iterate is the trace. It was run to
see whether Newton's method, whose round solves a whole linear system, fails at a lower noise
level than the iterations above. It does not appear to: it needs fewer rounds than $\max_iD(i)$
on every instance tried, by a factor between about $1.6$ and $7$, but the counts grow with $n$ in
the same way. For the base with distances uniform in $[64,128)$, at $n=2^{14},2^{16},\dots,2^{22}$:
at noise $0.5$, $\max D=82,161,223,320,444$ and Newton $27,43,49,63,80$; at noise $0.3$,
$\max D=180,717,2785,10439,35891$ and Newton $66,229,749,2572,8631$. One smoothing each.

\medskip\noindent\textbf{Referee pass on Section~\ref{sec:fixed}.} A separate automated referee,
given the section and no summary of it, read Lemmas~\ref{lem:onegate}, \ref{lem:onestep}
and~\ref{lem:typed} and Theorem~\ref{thm:fixed} and reported no error. It checked
Lemmas~\ref{lem:onegate} and~\ref{lem:onestep} by exact rational enumeration over all trace
prefixes for $H\le6$ and $m\le7$, Lemma~\ref{lem:typed} for all typed skeletons of length at
most $3$ with $n\le14$, the three numerical claims in exact arithmetic, and the one-step bound by
its own simulation. It found four imprecise statements, which are corrected in this version:
the measurability sentence in the proof of Lemma~\ref{lem:typed}, the missing condition
$\beta\ge0$, and two sentences of Remark~\ref{rem:fixedscope}. It did not check
Lemma~\ref{fact:densall}, the round bound of Theorem~\ref{kp:thm39}, or Proposition~\ref{prop:jacobi}.
A second pass by the same referee covered Definition~\ref{def:admissible},
Lemma~\ref{lem:fddens}, Corollary~\ref{cor:fdadm} and the changed statements, and again
reported no error. It checked the formula for the conditional mean by exact enumeration
($n=13$, $H\le5$), the six numerical claims and both tables, and, by simulation at
$n=2\cdot10^6$ with $H$ from $1$ to $1000$ and noise from $0.01$ to $1$, that the prefix
densities stay between the two bounds and that neither interval inequality of the proof is
violated on about $1.35$ million intervals per run. Three small points it raised are corrected
here (the inequality $\eps\beta<r$, the case $H=1$, one rounded figure), and
Remark~\ref{rem:fixedscope}(iv) records its observation about the size of the constants.
A third pass covered Section~\ref{sec:ordered}. It found one wrong figure, a bound stated as
$0.985$ for a value of $0.985027$, now $0.986$, and one inequality used without derivation
($c<1/2$), now stated; otherwise no error. It checked Lemma~\ref{lem:spacing} in exact
rational arithmetic at $300$ random parameter sets for $d\le60$; Lemma~\ref{lem:sumspacing} on
$2.4$ million partial sums with random factors $\rho_u$, by a value iteration over distances up
to $12$ at $20$ parameter sets (the supremum is attained by the constant sequence $2$ and stays
below the bound), and by hill-climbing at $400$ parameter sets; Lemma~\ref{lem:ordstep} by exact
enumeration on $11$ small ordered bases; and Lemma~\ref{lem:ordsum} on $7500$ windows of random
ordered bases.
A fourth pass covered the rewriting of Section~\ref{sec:ordered} for nondecreasing pointers. It
reported no error in a statement or proof, and four inaccurate sentences in the remarks and
examples, corrected here. It checked facts (b) and (c) of Lemma~\ref{lem:sumspacing} in exact
rational arithmetic, and found that $\beta\le2/3$ is really used ((b) fails from about
$\beta=0.77$); the lemma itself with the exact recursion on $559{,}240$ prefixes and by a value
iteration over distances up to $12$ and group sizes up to $20$; Lemma~\ref{lem:ordstep} with the
event $C^\circ$ by exact enumeration on $26$ small bases, including arbitrary and decreasing
pointers, where it also found that the bound fails without the condition $\pi_g=1$; and the
replacement by top-reading chains exhaustively for all bases and all noise with $n\le6$.
A fifth pass covered Section~\ref{sec:below}, which claims no theorem. It reported no error in
the two lemmas and one missing clause in Remark~\ref{rem:below}(i), added here. It checked
Lemma~\ref{lem:typebit} in exact rational arithmetic on $15$ bases, for every trace prefix and
eight functions $\Psi$ ($111{,}192$ cases), and Lemma~\ref{lem:highest} exhaustively for all
$409{,}110$ parent arrays with $3\le n\le9$ and $g_0\le3$. It reproduced the growth factors of the
recursion in Remark~\ref{rem:below}(iii); it did not check the one-gate value iteration.

\medskip\noindent\textbf{Section~\ref{sec:barrier}.} \texttt{barrier\_check.py} computes the vector
$v$ of Proposition~\ref{prop:barrier} by power iteration and checks the certificate inequality in
exact rational arithmetic, rounding every intermediate vector down to a multiple of $10^{-40}$.
It reports the smallest ratios $1.13549$ at $\eps=7/10$ and $1.01726$ at $\eps=18/25$, and
support $\{001,010,100,101\}$. \texttt{first\_moment.c} computes the exact growth rates of the
Kleene--prefix and $D_{\mathrm{last}}$ counts in the model of Section~\ref{sec:barrier}, and, with a
band for the resampled reads, the lower bounds that the certificate makes rigorous ($1.1355$ and
$1.0173$ for the band of the proposition). \texttt{canonical.c} adds the canonical constraint;
with it switched off it reproduces \texttt{first\_moment.c}. As a check of the model against the
instances themselves, the expected numbers of chains of lengths $1,2,3$ from a position of
$(2,3)$ at noise $0.70$ are $1.060$, $1.178$, $1.331$ in the model, and the sampled means at
$n=8\cdot10^6$ (last tenth active, three smoothings) are $1.042$ to $1.050$, $1.129$ to $1.161$, and
$1.212$ to $1.318$. The active part raises the density of ones from $\varpi=0.3219$ to $0.3292$;
with $0.3292$ in place of $\varpi$ the model gives $1.037$, $1.126$, $1.241$. The sampled means
are close to both sets of values; at length $3$ they scatter too much to say more
(\texttt{parity\_counts.c}). The searches are the loops of \texttt{search\_bases.sh}. \texttt{parity\_depths.c} produced the table of depths.

\medskip\noindent\textbf{Section~\ref{sec:block}.} \texttt{block\_thresholds.py} checks the
identities for $d^*$, the exact values at $\eps=4/9$, that $q(\eps)-\eps$ changes sign exactly at
the root $\eps_K=0.658967$ of the cubic, and the sign of $q(\eps)^2-\eps$ on a grid.
\texttt{block\_depths.c} produced the simulations of Remark~\ref{rem:block}.

\medskip\noindent\textbf{Section~\ref{sec:table}.} \texttt{table\_jacobi.c} simulates the
table--Jacobi iteration on random block bases by evaluating $G_R$ sequentially (it does not build
tables), checks that the fixed point is the trace, and prints the number of applications and the
length of the chain of Theorem~\ref{thm:table}(c), which is the chain $c_s$ of (d) when the base
has reach at most $R$, as in the runs with $R$ twice the block length; with $R=1$ it is the
prefix--Jacobi iteration. It produced the numbers of Remark~\ref{rem:table}. A third pass of the
referee covered Section~\ref{sec:table}. It found Theorem~\ref{thm:table} true in all four parts
and reproduced every number of the remark with the script and with its own implementation. It
checked (a), (b) and (c) by brute force on all $5913$ parent arrays with $n\le7$ for five values
of $R$ and on $3000$ random arrays, and the table description of one application against the
sequential definition on $1500$ random cases. It pointed out two gaps in the proof of (d), a step
that ends at $n$ and a ceiling in the last line, and that the first version stated without proof
that the iteration is not polylogarithmic and that no method of this kind can do better; these are
changed.

\medskip\noindent\textbf{Section~\ref{sec:latch}.} \texttt{latch\_check.py} is the certificate
of Corollary~\ref{cor:twofifths}. It works in exact rational arithmetic, rounds every lower bound
down and every upper bound up to a multiple of $10^{-18}$, and encloses $\varpi(2/5)$ between
rationals; it asserts (R), (G) and (V) of Definition~\ref{def:adm} for $\theta=1/2$, $h=40$,
$J=75$, $\eta=\gamma=1/500$, and prints the constants of Remark~\ref{rem:latch}(v). The least
value of $a_k-b_k$ along the sequence is $0.1997$, at $k=1$. With the argument \texttt{table} it
prints the table of Remark~\ref{rem:latch}(i) from the closed forms of
Proposition~\ref{prop:ideal} and asserts that they are fixed points of the ideal map.
\texttt{latch\_patterns.py} is the search of Remark~\ref{rem:latch}(iii). The file
\texttt{lean/SeqNORLatch.lean} proves Lemma~\ref{lem:latch}(a) in Lean: the identity for a pair of
gates, the monotonicity of a chain, and the sandwich used in the proof of
Theorem~\ref{thm:latch}.

\texttt{latch\_embed.c} builds $\mathrm{Latch}(C)$ as in Definition~\ref{def:latch}, with a
preamble of $20{,}000$ zero gates, for a random circuit of $s$ gates in which most gates read the
previous gate, so that the circuit is deep; it draws the smoothing and evaluates it. It reports
the number of gates $g$ for which $P_J(g)$ and $N_J(g)$ are in the wrong order, and the mean
densities of the high and of the low rails of stage $J$. With $\theta=1/2$, and $\theta=19/32$
in the last row:
\begin{center}
\begin{tabular}{ccccccccc}
\toprule
$\eps$ & $s$ & $m$ & $h$ & $J$ & $\rho$ & seeds & wrong gates & densities (ideal)\\
\midrule
$0.30$ & $30$ & $1024$ & $12$ & $30$ & $20$ & $1$--$5$ & none & $0.869$, $0.044$ \ ($0.869$, $0.043$)\\
$0.35$ & $30$ & $1024$ & $12$ & $30$ & $20$ & $1$--$5$ & none & $0.796$, $0.082$ \ ($0.798$, $0.081$)\\
$0.40$ & $30$ & $1024$ & $12$ & $30$ & $20$ & $1$--$5$ & none & $0.674$, $0.163$ \ ($0.682$, $0.157$)\\
$0.40$ & $20$ & $2048$ & $14$ & $50$ & $40$ & $1$--$5$ & none & $0.687$, $0.154$ \ ($0.687$, $0.154$)\\
$0.42$ & $4$ & $4096$ & $16$ & $80$ & $40$ & $1$--$3$ & none & $0.616$, $0.209$ \ ($0.613$, $0.209$)\\
$0.46$ & $30$ & $1024$ & $12$ & $30$ & $20$ & $1$ & $12$ of $30$ & $0.407$, $0.382$ \ (none)\\
$0.50$ & $30$ & $1024$ & $12$ & $30$ & $20$ & $1$ & $13$ of $30$ & $0.404$, $0.392$ \ (none)\\
\bottomrule
\end{tabular}
\end{center}
The densities are means over the gates and the seeds. The measured density of ones was $0.2048$,
$0.2240$, $0.2414$, $0.2381$ and $0.2449$ in the first five rows; it exceeds $\varpi(\eps)$
because the padding ratio is small. The numbers in brackets are the fixed point $L$ of the ideal
map with the measured density in place of $\varpi$; in the last two rows that map has no fixed
point off the diagonal. In the third row $30$ stages are too few to reach the fixed point. So the
runs agree with the analysis up to noise $0.42$ and fail above $4/9$, as the ideal map predicts.
A run at noise $0.43$ with $m=2048$, $J=50$, $\rho=40$ and $\theta=11/16$, made with the first
version of the script, had $7$ wrong gates of $20$; there the multiplier is $1.026$ and these
sizes are too small. All runs are far below the sizes of Corollary~\ref{cor:twofifths}, so they
test the mechanism and not the theorem.

\medskip\noindent\textbf{Referee pass on Section~\ref{sec:latch}.} A separate automated referee
read the section with the instruction to find an error. It reported none in the chain from
Lemma~\ref{lem:latch} to Corollary~\ref{cor:gap4}. It redid the algebra of
Proposition~\ref{prop:ideal} by hand and checked every identity with its own code in $60$-digit
arithmetic on a grid of noise levels from $0.001$ to $0.999$, counted the fixed points by brute
force (three below $4/9$, one above), and verified Corollary~\ref{cor:twofifths} with its own
implementation of Definition~\ref{def:adm} in exact arithmetic (the same bounds to six digits;
$J=64$ would suffice and $J=60$ would not). With its own simulator, which follows
Definition~\ref{def:latch}, it obtained no wrong gate in $13$ runs at noise $0.40$ with $m=1024$,
none at $0.30$ and $0.35$, none at $0.42$ with $m=8192$, and failure at $0.46$ and $0.50$; it also
compared the density of every rail with the conditional law used in the proof of
Theorem~\ref{thm:latch} and found the deviations as large as the sampling error and no larger. It
found three errors outside the proofs. The first version of \texttt{latch\_embed.c} did not build
the rails of the constants and read constants directly at a gate, which Definition~\ref{def:adm}
does not cover; this, and not the size of the bundles, had caused the wrong gates of that version
at noise $0.40$ (one or two in two of five runs), and the script and the table above are now the
corrected ones. One digit of the table of Remark~\ref{rem:latch}(i) was wrong. It asked for the
converse half of Proposition~\ref{prop:adm} and supplied its proof, for the sentence about the
coins in the proof of Corollary~\ref{cor:latch}, and for a rational threshold in the vote. It
objected to ``$4/9$ is exact for this organ'', which is proved only for the ideal map and for the
conditions of the proof, to a statement about success probabilities in the abstract that
holds for infinitely many $n$ only, and to several clashes of notation. All of this is changed.
It did not check Sections~\ref{sec:iterations} and~\ref{sec:dense}, the proofs of Lemmas~\ref{lem:dens2}
and~\ref{lem:dens2any} at the scale of the theorem, or the Lean file.

\medskip\noindent\textbf{Section~\ref{sec:s24}.} \texttt{s24\_check.py} is the certificate of
Proposition~\ref{prop:s24cert} (usage in the file; \texttt{python3 -I s24\_check.py 3661 8192 1
1000000 1 1000000 4 3 128 0.004 6000}, under a minute). In the coordinates $t$, with entries for
the phases $3,8,9,11,22$, it uses the boxes
$Q^0=[(0.5736,0.3920,0.5819,0.6692,0.3311),(0.5816,0.4000,0.5899,0.6772,0.3391)]$ and
$Q^1=[(0.6220,0.4409,0.6302,0.7102,0.3771),(0.6300,0.4489,0.6382,0.7182,0.3851)]$ (to four
digits; the script prints them to six), and finds
$L^1_0\ge(0.6100,0.4268,0.6180,0.7013,0.3654)$ and $U^0_0\le(0.5945,0.4113,0.6025,0.6882,0.3508)$
for the gate stage, and
$[L^1_J,U^1_J]\subseteq[(0.6238,0.4426,0.6321,0.7123,0.3789),(0.6274,0.4464,0.6357,0.7154,0.3826)]$,
$[L^0_J,U^0_J]\subseteq[(0.5761,0.3945,0.5844,0.6719,0.3337),(0.5798,0.3982,0.5881,0.6752,0.3371)]$.
The arguments \texttt{1117 2500 \dots\ 4000} and \texttt{89 200 1 100000 1 100000 4 1 16 0.01 1000}
give the checks at $0.4468$ and $0.445$ quoted after the proposition. The script asserts that every
read phase of $S_{24}$ has the parity of the gates that read it. \texttt{s24\_embed.c} is the
simulator of Remark~\ref{rem:s24}(v) (\texttt{./s24\_embed 8 0.40 1 32 20 1000 0.25 0.43 200000
seed}, seeds $1$ to $8$; set \texttt{VERBOSE} to list the gates). With the variable
\texttt{DEPTH} set it also computes $\max_iD(i)$ by the recursion of Lemma~\ref{lem:status}; these
are the numbers in Remark~\ref{rem:layered}(iii) (\texttt{./s24\_embed 4 0.40 1 J 10 500 0.25 0.43
200000 1} and \texttt{./s24\_embed 2 eps 1 J 5 200 0.25 0.43 100000 1}).

\medskip\noindent\textbf{Referee pass on Section~\ref{sec:s24}.} A separate automated referee read
the section, with Section~\ref{sec:latch} as far as it is used, with the instruction to find an
error, and reported none. It checked Lemma~\ref{lem:progchain}(a) by brute force for all chains of
length at most $9$, the parities behind Lemma~\ref{lem:envelope} against the program, and the
antitonicity of the envelopes on random comparable pairs. It wrote its own implementation of
Definition~\ref{def:adms24} in $50$-digit decimal arithmetic with directed rounding, with its own
placement of the boxes, and found the parameters of Proposition~\ref{prop:s24cert} admissible,
with $(L^1_J)_8=0.44265\ldots$ and $(U^0_J)_8=0.39815\ldots$, the same as \texttt{s24\_check.py} to
six digits; the check closes from $J\approx4200$ on and fails at $J=4000$, it also closes with the
four-digit boxes printed above, and the checks at $0.445$ and $0.4468$ and the failure at $0.447$
were reproduced. It found one gap, now closed: the proof of Proposition~\ref{prop:s24cert} did not
say why $L^v_k\le U^v_k$. Its other points (a leftover symbol, the description of the rounding in
the script, the constant $m\approx6.6\cdot10^{12}(a+1)\log_2s$, a sentence of
Remark~\ref{rem:latch}(iii) that gave $0.4468$ where $\eps_S$ is meant, the count in the comparison
of the embeddings) are corrected. It did not check the simulation of Remark~\ref{rem:s24}(v), and it
checked the results reused from earlier sections only as far as they are applied here.

\medskip\noindent\textbf{Section~\ref{sec:above}.} \texttt{status\_check.py} computes $D$ from its
definition with anchors and asserts the four cases of Lemma~\ref{lem:status} and the statements
(O1), (O2), for every round, on all $5913$ parent arrays with $n\le7$ and on $4000$ random ones
with $8\le n\le40$. For Proposition~\ref{prop:psi} it compares the closed form of the stationary
law with the law obtained by iterating the four-state chain on a grid of parameters at seven
noise levels (largest difference $2\cdot10^{-15}$), asserts monotonicity on the grid and that
the stationary probability of the value $1$ is $d_\eps$, checks the solution of (c) below $4/9$
and that it lies in the domain, and iterates $\Psi$ from $(d_\eps,1-d_\eps)$: the limit is that
solution at noise $0.20$, $0.40$, $0.44$ and $0$ at noise $0.46$, $0.50$, $0.60$, $0.70$.
\texttt{block\_unknown.c} computes $D$ on one smoothing of a random block base and reports the
fraction of $s$-open positions in the upper blocks. With $n=5\cdot10^7$ and blocks of $2000$
positions in the last fiftieth of the positions ($500$ blocks, of which the upper $150$ are
counted), the smallest and the largest value over the seeds $1,2,3$ are:
\begin{center}
\begin{tabular}{lccccc}
\toprule
noise, rounds $s$ & $10$ & $50$ & $200$ & $300$ & $\max_iD(i)$\\
\midrule
$0.40$, measured & $0.627$--$0.637$ & $0.561$--$0.571$ & $0.547$--$0.563$ & $0.539$--$0.556$ & $505$--$508$\\
$0.40$, iterates & $0.600$ & $0.569$ & $0.569$ & $0.569$ & \\
$0.46$, measured & $0.462$--$0.478$ & $0.106$--$0.134$ & $0.000$--$0.002$ & $0.000$ & $148$--$214$\\
$0.46$, iterates & $0.370$ & $0.122$ & $0.009$ & $0.002$ & \\
$0.50$, measured & $0.370$--$0.382$ & $0.003$--$0.008$ & $0.000$ & $0.000$ & $69$--$94$\\
$0.50$, iterates & $0.222$ & $0.005$ & $0.000$ & $0.000$ & \\
\bottomrule
\end{tabular}
\end{center}
The rows ``iterates'' are the iterates of $\Psi$ from the top, computed by the mode
\texttt{iterates} of \texttt{status\_check.py} with the measured density of ones ($0.2377$, $0.2581$,
$0.2706$) in place of $\varpi$ and with the reads that are resampled into the blocks, a fraction
$0.017$ of the resampled reads on average over the counted blocks, counted like reads of the block
below. In the first rounds the measured fractions lag behind, because at the start no position
is determined, also below the blocks. From $50$ rounds on they agree with the iterates to about
$0.02$ at noise $0.40$, where the fractions of the upper blocks decrease slowly because the
determined part comes closer; at noise $0.46$ the last open positions go faster than the
iterates. The regime of Theorem~\ref{thm:blockabove}, blocks of $C\ln n$ positions with a large
$C$, is not reached. With blocks of $400$ positions the fractions at noise $0.40$ are lower and
vary more from seed to seed. \texttt{latch\_local\_search.c} is the search of
Remark~\ref{rem:above}(iv); with period at most $7$ it examines $811{,}068$ patterns at noise
$0.46$, and the largest multipliers are $0.9719$ (the plain latch and its repetitions), $0.9676$
and $0.9642$.

\medskip\noindent\textbf{Remark~\ref{rem:mixed}.} \texttt{mixed\_base.c} builds the mixed base
with real zero gates below the layers, draws the smoothing, computes $D$ by the recursion of
Lemma~\ref{lem:status} and runs the prefix--Jacobi and the table--Jacobi iteration until they
stop, checking that they stop at the trace; it produced the table of the remark with
$n=2.5\cdot10^6$, $5\cdot10^6$, $10^7$, $m=1000$, the layers in the last twentieth and seed $1$.
\texttt{layered\_search.c} and \texttt{layered\_table\_search.c} are the two searches, and
\texttt{layered\_design.c} evaluates one design; these three use the model with a fixed density
for resampled reads and are explorations. A separate automated referee checked the remark with
its own simulator, which it validated against the definition of $D$ on $400$ small instances of
this base, and with a brute-force check of the two facts on all $46{,}233$ parent arrays with
$n\le8$. It confirmed the growth at noise $0.46$ and $0.50$ and could not break it (more layers,
other layer lengths and padding fractions, three seeds; the controls with a filler or with
$j_i=i-3$ at the even offsets show no growth). It corrected three statements of the first
version: the growth stops near $0.55$, not $0.52$, once the layers are long; a pointer of reach
$6$ does make the table--Jacobi iteration with reach $5$ slow; and the prefix--Jacobi iteration
is slow in simulation on block bases below $4/9$, so the mixed base is the first slow one above
$4/9$ only. The second table of the remark, the measured multipliers and the numbers on longer
pointers are the referee's. It also noted that the script counts applications of $G$, one more
than the index of Definition~\ref{def:jacobi}, and that the searches have little power at their
sizes.

\medskip\noindent\textbf{Referee pass on Section~\ref{sec:above}.} A separate automated referee
read the subsection with the instruction to find an error. It reported none in the chain from
Lemma~\ref{lem:status} to Corollary~\ref{cor:exact}. It derived Lemma~\ref{lem:status} by hand and
checked it with its own implementation of the definition of $D$ on all $5913$ parent arrays with
$n\le7$ and on $6000$ random ones, and the two passages from an open position to a step of a
chain or of a witness on $241{,}650$ cases each. It derived the stationary law again and checked
it numerically, the monotonicity on $320{,}000$ random pairs, and the fixed points on a grid. It
checked Lemma~\ref{lem:statusblock}(a) by brute force on $400$ small smoothed block instances
($303{,}774$ gate-rounds) and the four symbol probabilities by exact enumeration. It read the
proof of Theorem~\ref{thm:blockabove} line by line and found the order of the parameters
consistent, the groups of gates in~\eqref{eq:witness} disjoint and the events measurable as
used. With its own simulator, blocks of $400$ positions and $250$, $1000$, $4000$, $16000$
blocks, $\max_iD(i)$ was $57$--$58$, $91$--$94$, $94$--$108$, $109$--$116$ at noise $0.50$ and
$253$--$255$, $898$--$1005$, $3929$--$3934$, $14733$--$14848$ at noise $0.40$; and one
application of the status chain to the measured fractions of the block below reproduced the
fractions of the block above to $0.001$ (for example $(0.2052,0.2327)$ measured against
$(0.2049,0.2326)$ at noise $0.46$ after $10$ rounds). It found two errors outside the proof of
the theorem. Remark~\ref{rem:above}(ii) gave the rate $(1-\eps)/(1-q)$ and the level $\eps_K$ for
the count with the last reset as the anchor; that is the rate of the count of all chains. The
matrix that is now there is the referee's, and its simulation of the count gave $1.28$, $1.18$,
$1.05$, $0.93$ at noise $0.46$, $0.50$, $0.55$, $0.60$ against the eigenvalues $1.282$, $1.186$,
$1.063$, $0.939$. And the first sentence of Proposition~\ref{prop:psi}(d) lacked the condition
$qr<\sigma(1+q)$. It asked for the proof that the solution of (c) lies in the domain of $\Psi$,
for the condition $\bar w\le1$, and for a density bound in Step~1 that does not rest on a result
proved elsewhere. It objected to the sentence that the proof does not count chains, to the claim
that the open fraction is the image under $\Psi$ where only a bound is proved, to a table with one
seed, and to clashes of notation. All of this is changed. It did not check Section~\ref{sec:dense},
and no simulation is near the regime of the theorem, whose constants are astronomical.

\medskip\noindent\textbf{Referee passes on Sections~\ref{sec:barrier} and~\ref{sec:block}.} The
automated referee read Section~\ref{sec:barrier} without the subsection on short pointers. It
found Lemma~\ref{lem:dens2any}, Proposition~\ref{prop:barrier} and Corollary~\ref{cor:barrier}
correct, rebuilt the certificate with its own code from the matrices of the proof (smallest ratios
$1.135493$ and $1.017261$), and reproduced the tables. It pointed out that the first version
claimed too much: Theorem~\ref{thm:ordered} counts only top-reading chains, which the base
$j^{\mathrm{bar}}$ does not block at noise $0.72$; the abstract, the section and the corollary now
say that the barrier concerns the count of all chains. It also pointed out a missing hypothesis in
Lemma~\ref{lem:dens2any} ($1/(1+\rho)+\tau\le1/4$), a false comparison of sampled and model
counts, and several loose phrases; all are changed. For Section~\ref{sec:block} it found the
lemmas and theorems correct. It checked Lemma~\ref{lem:blockgate} by exact enumeration on
$46{,}800$ gate and history pairs, Lemma~\ref{lem:Dfacts} and the depth recursions by brute force
on all $5913$ parent arrays with $n\le7$ and $4000$ random ones, Lemma~\ref{lem:oneblock} by exact
enumeration on small histories and by measured offspring rates, and the four levels with its own
simulator on both sides; these runs are at sizes far below the regime of the theorems, and one
entry went against the pattern ($\max_id(i)=250$ with $499$ blocks at noise $0.645$ over zero
gates). It pointed out one error, a lower bound $\eps_K=0.65896\ldots$ that was
rounded up to $0.659$, and several overstatements (only the Kleene--prefix iteration without zero
gates has both sides; $\Theta$ where only $\Omega$ is proved; the constant in
Theorem~\ref{thm:block}(c)); it supplied the two-line proof that $q(\eps)^2=\eps$ only at $4/9$.
All are changed here. It tested
Lemma~\ref{fact:depths} by brute force, which showed that its last clause needs $k\ge2$, now
stated. It did not check the proofs of Theorem~\ref{kp:thm39} and Lemma~\ref{kp:pruned}(b). In a second pass it read the subsection on short pointers and reported no error in
Theorem~\ref{thm:short}. It checked the claim about cuts by brute force on all dependency chains
of $3000$ small smoothed instances ($718{,}214$ steps, $3027$ cuts), the deterministic round count
for every threshold on the same instances, and the segments at sizes up to $10^7$. It corrected the
size of the constant quoted in Remark~\ref{rem:short} and asked for ``every positive noise
level''; both are changed. The reports on Section~\ref{sec:barrier} and on
Section~\ref{sec:block} come from two separate referee runs.

\medskip\noindent\textbf{Section~\ref{sec:designs}.} \texttt{potential\_check.c} proves
Lemma~\ref{lem:potential}(b). Compiled with \texttt{-ffp-contract=off}, the run
\texttt{potential\_check 0.45499 0.46001 0.99979 0.0001} covers $[0.45499,0.46001]$ in $51$
intervals of noise levels with $778{,}769$ boxes. Three runs with the factor $0.9995$, from
$0.45999$ to $0.47001$, from $0.46999$ to $0.50001$ and from $0.49999$ to $1$, cover
$[0.45999,1]$ in $108$ intervals with $351{,}696$ boxes. Each run prints \texttt{PROVED}. As controls, the factor $0.999$ on
$[0.455,0.4551]$ and the factor $1$ on $[0.4535,0.4536]$ are not proved. \texttt{runs\_check.c}
proves parts (b) and (c) of Lemma~\ref{lem:runs} with the same arithmetic: the runs
\texttt{runs\_check 0.45499 0.50001 0.8999 18 0.0005} and
\texttt{runs\_check 0.49999 1.0 0.8999 18 0.005} print \texttt{PROVED}, with $91$ intervals and
$7741$ boxes and with $101$ intervals and $1401$ boxes. As controls, on $[0.45499,0.4552]$ the
factor $0.79$ with $18$ steps and the factor $0.8999$ with $12$ steps are not proved.
\texttt{two\_back\_test.py} is the numerical test of Remark~\ref{rem:two}(ii)
(\texttt{python3 -I two\_back\_test.py 1 300} and the same with the arguments $7$ and $400$).
A separate automated referee read Subsection~\ref{sec:two} with the instruction to find an
error and found none. It checked Lemma~\ref{lem:twochain} against its own joint chain of two
copies, in exact rational arithmetic on $420$ designs and in floating point on $1500$, compared
the derivative of $\Phi$ with finite differences, read \texttt{runs\_check.c} line by line and
reran it, and searched about $110{,}000$ designs of period up to $30$ for a violation of the
inequality of Theorem~\ref{thm:two}; the largest ratio of its two sides was $0.9972$, in a chain
of far reads only, and $0.7814$ at a gate of type $\mathrm L_2$. Its four remarks on wording
are taken into the text, and its programs are in \texttt{referee-scripts/two/}.
\texttt{traces\_check.c} proves Lemma~\ref{lem:traceineq} with the same arithmetic: the runs
\texttt{traces\_check 0.45499 0.50001 0.8999 0.0005} and
\texttt{traces\_check 0.49999 1.0 0.8999 0.005} print \texttt{PROVED}, with $91$ intervals and
$394{,}241+175{,}379$ boxes and with $101$ intervals and $25{,}265+19{,}713$ boxes; as a
control, the factor $0.83$ on $[0.45499,0.4552]$ is not proved. \texttt{twoback\_check.py}
(\texttt{python3 -I twoback\_check.py 8}) is the check of Remark~\ref{rem:twolayered}(ii),
\texttt{twoback\_layered.c} the simulation of part (i) of that remark
(\texttt{twoback\_layered 0.46 EL 4 200 400 0.5 1} prints its largest entry), and
\texttt{tracemap\_test.py} compares the trace recursion of Definition~\ref{def:tracemap} with a
simulation of the traces in the ideal model for random programs
(\texttt{python3 -I tracemap\_test.py 3 18 400000}: largest deviation $0.006$, within the
sampling error). A separate automated referee read Subsection~\ref{sec:twolayered} with the
instruction to find an error and found no mathematical error. With its own programs it checked
Lemma~\ref{lem:twotraces} and the conditions (O1) to (O3) on all $4{,}037{,}912$ parent arrays
with $n\le10$, the sandwich of Step~1 on all arrays with $n\le8$ with every choice of kept
gates and every start of a chain, the trace recursion against simulations of the ideal model
and of smoothed layered instances, part (d) of Proposition~\ref{prop:tracemap} by
finite-difference Jacobians on $3600$ random programs (largest ratio $0.9972$, in a chain of far
reads only), and part (e) against the matrix of Definition~\ref{def:two}; it read
\texttt{traces\_check.c} line by line and reran it. It judged Steps 0 and 2 of the proof of
Theorem~\ref{thm:twolayered} proved, and Steps 1 and 3 proved up to routine details, which are
now written out (the choice of $J_0$ and of $\kappa$, the error $2\eps/g$ in Step~1); it also
corrected the direction of the order in Proposition~\ref{prop:tracemap}(c) and the wording of
Remark~\ref{rem:twolayered}(i). It did not recheck the older results that the proof cites, nor
the unchanged end of Step~3. It noted one point in the proof of Theorem~\ref{thm:layered}
itself: without zero gates before a chain, the first gate can list the gate before it through
its own coin, with probability $1/m$, and become a NOT gate; this changes one threshold at one
gate by at most $(1-\eps)/m$ and is covered by the error terms. Its programs are in
\texttt{referee-scripts/twolayered/}.
\texttt{scan\_check.c} proves Theorem~\ref{thm:scan}: the three runs quoted in its proof print
\texttt{PROVED}, with $19$, $48$ and $115$ intervals of noise levels. As controls,
\texttt{scan\_check 0.45499 0.4552 0.96 40 0.0001} and
\texttt{scan\_check 0.44 0.441 0.99 40 0.001} are not proved, as they must be, since
$\lambda^2$ is $0.9623$ at noise $0.455$ and exceeds $1$ below $4/9$; and
\texttt{scan\_check 0.4445 0.445 0.9999 120 0.0005} is proved, so the theorem also holds a
little above $4/9$. \texttt{window\_designs.c} is the exact computation of
Remark~\ref{rem:scan}(iii), in floating point: \texttt{window\_designs 0.46 2,1:P01,T10,T00}
and \texttt{window\_designs 0.46 2,1,1:P01,T10,T20} print the two values of that remark, and
\texttt{window\_designs 0.46 search 5 3} runs through the programs of period $5$. For the
words $\mathrm{UR}$ and $\mathrm{URR}$ the two programs agree to six digits. A separate automated
referee read Subsection~\ref{sec:scan} with the instruction to find an error. It found none in
Definition~\ref{def:scan}, in Theorem~\ref{thm:scan} and its proof, or in
\texttt{scan\_check.c}, which it read line by line, reran, and compared with its own computation
of the block quantities for all $k\le120$ at eight parameter points (agreement to $2\cdot10^{-12}$).
It computed $G_2$ with its own program for all $65{,}535$ words of period up to $16$ at five
noise levels, and simulated real towers of $10^6$ positions per layer for the words
$\mathrm U$, $\mathrm{UR}$ and $\mathrm{URR}$ ($0.9460\pm0.0020$, $0.8649\pm0.0018$ and
$0.8498\pm0.0018$ at noise $0.46$). It found one wrong sentence in Remark~\ref{rem:scan}(iii),
which attributed the values of a program of period $2$ to a reordering of the program of period
$3$, and several statements of that remark that said more than was proved; these are corrected.
Its programs are in \texttt{referee-scripts/scan/}. The script \texttt{potential\_scan.c} was used to choose the exponents,
and the script \texttt{potential\_bellman.py} is the value iteration of Remark~\ref{rem:indep}(ii);
neither is part of the proof. \texttt{organ\_multiplier.c} computes the
fixed points and the two matrices of Definition~\ref{def:multipliers} for the searches of
Remark~\ref{rem:designs}; \texttt{organ\_onebit.py} and \texttt{organ\_long\_designs.py} compute
them by the recursions of Lemma~\ref{lem:onebit} for designs with independent reads, and the second
also iterates the density map (the cycle of $S_{24}$); \texttt{flipspread.c} is the simulation of
Remark~\ref{rem:designs}(v).

\medskip\noindent\textbf{\texttt{rows\_check.c}.} It proves Lemma~\ref{lem:rowscert}: the three
runs quoted in its proof print \texttt{PROVED}, with $24$, $67$ and $113$ intervals of noise
levels and $747$, $712$ and $535$ boxes; the largest absolute value of a potential is $3.29$. The
left end of every interval of noise levels is moved down by one unit in the last place, so the
decimal number $0.455$ is covered. Only the check of~\eqref{eq:rowsC1} and~\eqref{eq:rowsC2} is
part of the proof; the search that proposes the functions $c_y$ is a relaxation method in floating
point. As controls, \texttt{rows\_check 0.455 0.457 0.9 8 0.002 0.004} is not proved (run by the
second referee below), and on the noise levels $[0.45,0.455]$ an earlier version of the search
found no potentials for $\theta=0.999$; this says nothing about potentials of another form. The other programs of Subsection~\ref{sec:rows} are
simulations and explorations in floating point. \texttt{rows\_periodic 0.455 0.3333333333 9}
gives the largest ratio $0.9476$ over the periodic rows of period up to $9$, and
\texttt{rows\_word.py} with the word \texttt{o1oA} (an open zero, a closed one, an open zero and
thirty times a closed one and a closed zero) gives $0.9583$, $0.9883$ and $1.0290$ at noise
$0.455$, $0.44$ and $0.42$, for $d=1/3$. \texttt{rows\_ideal 0.46 200000 12 0.4} gives ratios
between $0.74$ and $0.77$ from row to row, and \texttt{rows\_ideal 0.30 200000 12 0.4} between
$0.88$ and $0.89$. \texttt{rows\_depths 0.46 2000 320 3 1} gives the largest depth $41.3$,
$60.3$, $83.0$ and $109.7$ for $40$, $80$, $160$ and $320$ rows, as averages of three runs, and
$40.7$, $55.7$, $78.0$ and $104.3$ rounds of the prefix--Jacobi iteration.

\medskip\noindent\textbf{\texttt{rows3\_check.c}.} It proves Lemma~\ref{lem:rows3cert}. It is
\texttt{rows\_check.c} with a chain of \texttt{NR} rows, $\texttt{NR}=3$; compiled with
\texttt{-DNR=2} its check is that of \texttt{rows\_check.c}. The three runs quoted in the proof print
\texttt{PROVED}, with $11$, $10$ and $19$ intervals of noise levels and $218$, $106$ and $179$
boxes; the largest absolute value of a potential is $2.40$. The run for $[0.30,0.34]$ was made
after the referee passes below, with the same program. Its search
takes linear inequalities from the words at which the check fails. The numbers of
Remark~\ref{rem:rows3}(ii) are from \texttt{exploration/more\_rows/per3.c}, which computes the
ratio for a chain of $R$ rows over periodic rows (\texttt{per3 0.30 0.3345 3 8 0} and so on; with
$R=2$ it agrees with \texttt{rows\_periodic.c}), and those of Remark~\ref{rem:rows3}(iii) from
\texttt{exploration/several\_distances/per2.c} and \texttt{per2b.c}; these are explorations in
floating point. A separate automated referee read the proof of Theorem~\ref{thm:rows3} with the
instruction to find an error and found none. It asked for six clarifications, which are made:
the factor $1-2/m_0$ for a far one of row $3$, the intervals in~\eqref{eq:rows3regen}
and~\eqref{eq:rows3alive}, the union bound behind the use of~\eqref{eq:rows3alive}, which steps
are disturbed in (C), three windows instead of two, and the assembly of Step~4. It built the true
rows, the ideal rows and the row $(\tilde y,\tilde o)$ of $4800$ windows of smoothed instances
with $H=40$ and tested every deterministic claim of Steps~2 and~3 on them, and it computed the
multiple sums of (C) by brute force for three pairs $(\kappa,r_1)$. A second referee read
\texttt{rows3\_check.c} against \texttt{rows\_check.c}, found the proof part correct and, for
two rows, identical in its outputs to \texttt{rows\_check.c}; it checked that the search cannot
change a verdict; and with its own exact program for the chain of three rows it
verified~\eqref{eq:rowsC1} and~\eqref{eq:rowsC2} for the potentials of one box at noise
$[0.42,0.422]$ over all $45{,}568$ admissible words without leaving branches. Its search over
periodic rows found the largest ratio $0.947$ at noise $0.30$. Their programs are in
\texttt{referee-scripts/rows3/}.

\medskip\noindent\textbf{\texttt{greedy\_steps\_check.py}.} For Lemma~\ref{lem:greedysteps}. On
$3000$ random parent arrays of length up to $60$, of three kinds, it simulates the parallel greedy
algorithm on the graph $\Gamma(j)$ and checks that the set it finds is $\{i:x_i=1\}$, the recursion
for $s$, and $D(i)\le2s(i)$; the largest ratio of $\max_iD(i)$ to $\max_is(i)$ it meets is $1.67$.

\medskip\noindent\textbf{Referee pass on Subsection~\ref{sec:rows}.} Two separate automated
referees read the subsection with the instruction to find an error. The first read
Lemma~\ref{lem:rowstatus} and Steps~0 to~4 line by line and recomputed the constants. It found no
error. It found four places where the text did not say what the computation uses, which are
corrected: the fields $\mathcal G_j$ must not contain the noise of the gate $a-K$; the event that
a given step is the last near step is not measurable where the first version used it, and the
proof now uses the larger event that it is a near step; the constants of
Lemma~\ref{fact:densall} do not depend on the base; and the lower bound on $\Phi_j$ is needed for
$j\ge1$ only. It tested Lemma~\ref{lem:rowstatus} on $4000$ random parent arrays, the claim of
Step~1 on about $12{,}600$ cases on smoothed instances, and the agreement of the true and the
ideal upper row after a regeneration, with~\eqref{eq:rowssplit}, on $2400$ windows. The second
read \texttt{rows\_check.c} line by line and found the operators, the rounding, the walk over the
words and the cover sound. It found that the first version did not cover the decimal number
$0.455$ itself, because the nearest double is larger, and that two ends of the cover were not
rounded outward; both are corrected, and the runs were repeated. With its own program in exact
rational arithmetic, written from Definition~\ref{def:tworow}, it checked~\eqref{eq:rowsC1}
and~\eqref{eq:rowsC2} for the potentials of one box at noise $[0.455,0.456]$ over all $45{,}568$
admissible words, without leaving branches, and for one box at noise $[0.7,0.704]$; it compared
the operators of the program with exact ones at $400$ random boxes and functions; and it
searched the periodic rows of period up to $9$, at eleven noise levels from $0.455$ on and three
densities, for a ratio above $\theta$: the largest was $0.944$. Neither referee checked the
older results that the proof cites. Their programs are in \texttt{referee-scripts/rows/}.

\medskip\noindent\textbf{Referee pass on the search.} A separate automated referee checked
Proposition~\ref{prop:designs}, the program and the claims of the first version of the search
remark, with its own implementation (exact propagation, Monte Carlo of coupled chains and Newton's
method), $101$ comparisons on $33$ designs that cover all read types, and simulations on smoothed
instances. It found Proposition~\ref{prop:designs} correct and the program right, with one
weakness: its spectral radius, taken from the ratio of the norms of $M^{256}$ and $M^{128}$, can be
off by up to $5\cdot10^{-3}$ for several rails. It found that the main claim of the first version
was false: with periods of $12$ and more there are designs with independent reads that beat the
plain chain above $4/9$, and one whose density map has an attracting cycle of length two; these
are the designs of Remark~\ref{rem:beyond}, which we recomputed with a separate program. It also
found that our local searches had been absorbed by the plateau $\rho_u=\lambda$, that the lagged
iteration of the program can converge to a fixed point that repels under $\Phi$ and cannot see
cycles, that unequal proportions of the rails change the density map when there are gates of type
$\mathrm U$, and several overstatements and wrong counts. The section was rewritten accordingly,
and the referee read the new version. It confirmed the counts and the sign of the eigenvalue of
$S_{24}$, corrected three digits of the table of Remark~\ref{rem:beyond}, the sizes of the layers
in Remark~\ref{rem:designs}(vi) and the accuracy of the old spectral radius for $\rho_s$, and
pointed out that $\Phi\circ\Phi$ is monotone for $S_{24}$ (Remark~\ref{rem:beyond}(iv)); the
first version had said the opposite. It asked for the cautious formulation of the levels in the
introduction and in Remarks~\ref{rem:indep} and~\ref{rem:beyond}. All of this is changed. It did not
check Lemma~\ref{lem:onebit}, Lemma~\ref{lem:potential} or Theorem~\ref{thm:indep}.

\medskip\noindent\textbf{Referee pass on Section~\ref{sec:layered}.} A separate automated referee
read the subsection with the instruction to find an error. It checked Lemma~\ref{lem:lowup} case by
case against Lemma~\ref{lem:status}, the doubled design and its monotonicity, the use of
Theorem~\ref{thm:indep}(a) between class sets and on unions, Step~1 of the proof of
Theorem~\ref{thm:blockabove} for a zero preamble, the constants of~\eqref{eq:Pphi}, Step 2, the
cascade and the recursion of Step 3, and the Corollary, and found the mechanism sound. It found no
broken argument, eight gaps, all fillable and now filled (the filtrations, gates resampled to the
position before them, the law of $x_{g-1}$, one failure term per gate instead of one in all, the
final count, the constant $J_0$, Lipschitz paths, the exponent of the density event), and two
errors in the claims around the theorem, now corrected: the definition alone allows bundles of one
chain and so covers nearly every base with a long preamble, where the conclusion is false, so
every statement now carries the width $m\ge C_1\ln n$ (it found $\max_iD(i)=159$, $276$, $509$ on
the mixed base, written as such a base, at noise $0.50$); and the bases of
Definitions~\ref{def:latcha} and~\ref{def:proga} as written have a short preamble, so the
Corollary now refers to the form of Lemma~\ref{lem:longer}, which still proves the hardness. In the
ideal model it computed the doubled map of $S_{24}$ from the widest intervals: the widths shrink
to $10^{-14}$ at noise $0.455$, $0.46$ and $0.50$, at rates $0.985$, $0.975$ and $0.898$ per stage,
while at $0.44$ the width stays at $0.233$ (the two-cycle) and at $0.447$ it is still $0.041$ after
$4000$ stages; for the latch the width at $0.455$ is $10^{-9}$ after $1000$ stages and stays at
$0.33$ at $0.43$. Its simulations of $\mathrm{Prog}(C)$ with a long preamble gave $\max_iD(i)=34$,
$52$, $64$, $91$ for $J=8,32,128,512$ at noise $0.50$, against linear growth at $0.40$.

\medskip\noindent\textbf{Section~\ref{sec:allbases}.} \texttt{allbases\_levels.py} encloses
$\varpi(\eps)$ between rationals and checks in exact rational arithmetic that
$\chi(0.7287)>0>\chi(0.7288)$, that $\kappa_2(0.7287)<0.96$ (it is $0.95730\ldots$), and that the
level of Remark~\ref{rem:allbases}(ii) lies between $0.7393$ and $0.7394$; in floating point
$\eps_1=0.728776$, and $\kappa_2=1$ at $0.7112$. \texttt{allbases\_exact.py} (\texttt{python3 -I
allbases\_exact.py 7 7 30}) checks Lemma~\ref{lem:twotype}(a), (b), (c) on all $5912$ parent
arrays with $2\le n\le7$, and then, for $240$ pairs of a base of length $7$ and a noise level
($0.3$, $0.5$, $0.7$, $0.73$; random bases with many pointers of reach $2$ and $3$), sums over all
$8!$ outcomes of the noise and checks five inequalities for every gate $k\ge5$ and every $s$:
\eqref{eq:twoa} and~\eqref{eq:twob}, and for the gates with base pointer $k-2$ the inequality
that collects the cases of gate $k$ in the derivation of~\eqref{eq:reachtwo}, and the bounds on
$Z$ and $Y$. In these the constants
$c$, $\beta$, $\eta$ are replaced by the densities of the instance (for example $r_c$ by
$1-\eps\,O_{k-1}/k$ under the expectation), which the constants bound from one side, and $G^*$ is
the sure event; the file states them. No violation; the largest difference is $1.8\cdot10^{-13}$,
which is rounding. The inequalities~\eqref{eq:generic} and~\eqref{eq:reachtwo} follow from these
by the algebra in the proof. \texttt{design\_search.c} is the search of
Remark~\ref{rem:layered}; its usage and the runs of the remark are in the file, and with the
variable \texttt{DR} it prints the depth $D_R$ of Definition~\ref{def:setval}.
\texttt{setval\_table.c} runs the set-valued table iteration of Definition~\ref{def:setval} and
the table--Jacobi iteration on one smoothing of a mixed base and prints the number of rounds of
the first and the least $s$ with $y^{(s)}=x$ for the second (Remark~\ref{rem:setval}); in every
run the second number was at most the first, as Proposition~\ref{prop:setval}(c) requires.
\texttt{design\_search.c} prints the same least $s$. \texttt{dep\_design.c} is the simulation of
Remark~\ref{rem:dependent}: \texttt{DELTA=0.003 dep\_design eval 0.46 U,R1 2000000 40 100} prints
the entry of the table for the pairs of neighbours, and the commands of the two searches are in
the file. It is a Monte Carlo estimate in the ideal model; its two checks are the first rows of
the table, which lie about $0.002$ below $\lambda$ and $\rho_u(D_{16})$, the bias of the
estimator. A separate automated referee read the remark with the instruction to refute it. Its
own simulator \texttt{refsim.c}, its exact values for independent reads
(\texttt{exact\_indep.py}), its check of the fact on regular gates on all parent arrays with
$n\le5$ (\texttt{fact\_bruteforce.py}) and its search driver are in
\texttt{referee-scripts/dependent/}. It reproduced the table, found that the first version of
the paragraph on searches could not be reproduced from the commands in the file and that its
conclusion was stated too strongly, and corrected the statement of the fact on regular gates;
the remark is rewritten accordingly, with its numbers. \texttt{setval\_check.py}
(\texttt{python3 -I setval\_check.py 8 4}) runs both iterations from their definitions on all
$46{,}232$ parent arrays with $2\le n\le8$ and every reach $R\le4$ and checks parts (a), (b), (c),
(e) and the first bound of (f) of Proposition~\ref{prop:setval} in every round, and for $R=1$ that
$K^{(s)}=\{i:D(i)\le s\}$; the additive terms $1$ in (b), for $D_R$, and in (c) are attained.

\medskip\noindent\textbf{Referee pass on Section~\ref{sec:setval}.} A separate automated referee
read the subsection with the instruction to find an error. It found none in
Definition~\ref{def:setval} or in parts (a) to (e) of Proposition~\ref{prop:setval}, and checked
them with its own programs on all $409{,}112$ parent arrays with $2\le n\le9$ and $R\le4$, in
every round; two deliberately wrong variants of the iteration failed these checks at $n=4$. Its
findings are all taken up. The first version said that there is no bound beyond (b), although
Theorem~\ref{thm:table}(c) and (d) carry over; this is now part (f). Part (b) had only the bound
$\max_iD_R(i)+1\le\max_iD(i)+1$; the bound $\max_iD(i)$, and the identity of the iteration with
reach $1$ and the relevance-pruned iteration on these arrays, are the referee's. The scripts
print the least $s$ with $y^{(s)}=x$, which the first version called the number of applications
of $G_R$; it is one less. The sentence on the collapse of the sets of windows was supported for
the pointer $i-2$ only, and the numbers for the pointer $i-6$ with reach $6$ in
Remark~\ref{rem:setval} are the referee's, as are those at $1280$ layers. It also gave three
examples: the term $1$ in (b) is needed for $n=5$, $j=(0,0,0,0,2)$, $R=2$, and the one in (c) for
$n=2$, $j=(0,1)$; the sets of windows are strictly ahead of three-valued logic with exact short
reads first at $n=6$, $j=(0,0,0,1,3,4)$, $R=2$; and the number of applications of $G_R$ is not
monotone in $R$ ($n=8$, $j=(0,1,0,2,2,0,0,5)$: three applications for one reach, four for the
next), unlike the rounds of the set-valued iteration. It noted a variant of $D_R$, with cost
$\max(1,D(p))$ at a short gate, whose level sets lie in $K^{(s)}$ without a shift; we did not
include it. With its own simulator on the mixed base at noise $0.50$ (six smoothings, $40$ and
$160$ layers) it found $\max_iD(i)$ in $91$--$121$ and $302$--$370$, $\max_iD_2(i)$ in $50$--$56$
and $151$--$170$, and $20$--$28$ and $24$--$29$ rounds, and it reproduced the numbers of the
remark with our programs. Its programs are kept in \texttt{verification/referee-scripts/setval}.

\medskip\noindent\textbf{Referee passes on Section~\ref{sec:allbases}.} A separate automated
referee read a first version of the section, which proved the level $0.7394$ of
Remark~\ref{rem:allbases}(ii) by a recursion for the two kinds of open positions, with the
instruction to find an error. It found none in Lemma~\ref{lem:twotype}, in (F1)--(F3), in
\eqref{eq:twoa} and~\eqref{eq:twob} or in that recursion, and corrected statements in the remarks
that had used a density of zeros which is not available for every base. A second referee read
the present version with the same instruction and found no error in the theorem or its proof: the
facts (F1)--(F3) for $\beta\in[\frac14,\frac12)$, the indices and $\sigma$-fields, \eqref{eq:generic},
the derivation of~\eqref{eq:reachtwo} for every pointer of gate $k-1$, the recursion and the
analytic claims ($\eps_1=0.728776$, $\kappa_2(0.7287)=0.95730$, $\kappa_2=1$ at $0.7112$). With its
own program, in exact rational arithmetic, it checked every step of the derivation
of~\eqref{eq:reachtwo} and the inequality itself, and~\eqref{eq:twoa}, \eqref{eq:twob}
and~\eqref{eq:generic}, with the densities of the instance and $G^*$ the sure event: for all $720$
bases of length $6$ and all $5040$ of length $7$ at four or five noise levels, all $40{,}320$ of
length $8$ at noise $1/2$ and $0.73$, $400$ random bases of length $9$ and structured bases of
length $10$. There was no violation; several intermediate steps hold with equality in many
cases, and in~\eqref{eq:reachtwo} the left side was at most $0.632$ times the right side. It also
checked the form with $G^*_k$ and constants $c$, $\beta$, $g_0$ on bases of length up to $11$, and
it computed $D$ in three ways (the iteration, the definition, the four cases of
Lemma~\ref{lem:status}) on all parent arrays with $n\le7$. Its corrections are all made: the
remark on longer pointers was wrong in saying that an upper bound on the density of ones is
needed (the numbers $0.7208$ and $0.7185$ of Remark~\ref{rem:allbases}(iii) are its own as well);
the growth on the base $j_i=i-2$ had been written as $13\ln n$ where $13$ is only the slope; the
sentence on the script had named~\eqref{eq:reachtwo} where the script checks its three
ingredients; and two strict inequalities fail at $\eps=1$. It noted that $2\eps/k$
in~\eqref{eq:twob} and~\eqref{eq:generic} can be $\eps/k$, which we did not change, and that it
read the source and not the typeset pages. The simulations quoted in
Remark~\ref{rem:allbases}(iv) are its own, and its programs are kept in
\texttt{verification/referee-scripts/allbases} (\texttt{ref\_check.py} is the exact check).

\medskip\noindent\textbf{Referee pass on Section~\ref{sec:levels}.} A separate automated referee
read the section with the instruction to find errors. It found none in Theorem~\ref{thm:levels},
Theorem~\ref{thm:work} or Corollary~\ref{cor:speedup}; it checked the depth of every step of the
reduction, the closure of $\mathrm{BPNC}^k$ under logspace reductions for $k\ge2$, the cost of each
step of the iteration and of the construction, the use of Theorem~\ref{kp:thm39}, the error analysis
and the arithmetic of the corollary. It found that the first version overstated what
Theorem~\ref{thm:work}(a) gives in the terms of~\cite{krs} (the work is $\Theta(n\log n)$, so it is
not of the kind of EP), that the remark on block bases read as if it concerned noise between
$4/9$ and $3/4$ and did not say that its time bound holds with high probability only, that the
bibliography lacked~\cite{krs}, that the completeness of circuit evaluation under $\mathrm{NC}^1$
reductions is known, that the choice of $\eps^*$ in the proofs of Theorems~\ref{thm:hard4}
and~\ref{thm:work}(b) allowed $\eps^*=0$, and several smaller points in the construction and the
notation. All are changed; it also suggested the cap in the proof of Theorem~\ref{thm:work}(a),
which makes its work bound hold in the worst case. It could not see the exact definitions
of~\cite{krs}, and neither could we; the section uses descriptions of the classes.

\medskip\noindent\textbf{Referee pass on Theorem~\ref{thm:indep}.} A separate automated referee
read Definition~\ref{def:indep}, Lemmas~\ref{lem:onebit} and~\ref{lem:potential},
Theorem~\ref{thm:indep}, Remark~\ref{rem:indep} and the verifier with the instruction to find an
error. It found none in the mathematics, which it checked by hand and against brute force on $480$
random designs (density map, coupled copies and finite differences, all to $10^{-10}$ or better).
It found the verifier sound but its quoted runs not quite covering: the loop over noise levels
stopped two units in the last place short of the end of each run of width $10^{-4}$, and the
decimal arguments $0.455$ and $0.9998$ are rounded upwards. The loop is corrected, and the runs
quoted above start below $0.455$ and use a factor below $0.9998$. It reimplemented the check from
the statement in $30$-digit decimal arithmetic, with the same box counts, evaluated the inequality
in $50$-digit arithmetic (largest ratio $0.99931$ at noise $0.455$, crossing $1$ near $0.4538$),
and tried to break the theorem on several thousand designs at noise $0.455$, $0.46$ and $0.5$,
including those of Remark~\ref{rem:beyond}, without success. It asked for the wording on changing
designs (they forget the first layer, without converging to a point), for the statement of
part (d) for stages that read several stages, for the error term in the remark on finite chains,
for two definitions in part (d) and (c), and for a few changes of notation; all are made. It also
found the shorter designs quoted in Remark~\ref{rem:beyond}(i).

\medskip\noindent\textbf{Sections~\ref{sec:iterations}, \ref{sec:dense}, \ref{sec:newton}
and~\ref{sec:gapbase}.} These sections were read by independent automated referees in an earlier
form; their findings are resolved there, and their computations (the constants of
Theorem~\ref{kp:thm34}, the recursion of Theorem~\ref{gb:embed}, the arithmetic of the window
lemma of Section~\ref{sec:newton}) were rechecked when they were set into this paper. The
deterministic lemmas listed in the Lean library of the collection's repository are formalized; the
probabilistic steps are not.

\medskip\noindent\textbf{Not checked.} No human referee has read any of this paper. Apart from
the lemmas checked in Lean and the referee passes described above, the proofs have been checked
only by their authors.

\begin{thebibliography}{99}
\bibitem{adleman} L.~Adleman. Two theorems on random polynomial time. \emph{FOCS}, 1978.
\bibitem{elias} P.~Elias. Universal codeword sets and representations of the integers.
\emph{IEEE Trans.\ Inform.\ Theory} 21(2), 1975.
\bibitem{evansschulman} W.~S.~Evans and L.~J.~Schulman. Signal propagation and noisy circuits.
\emph{IEEE Trans.\ Inform.\ Theory} 45(7), 1999.
\bibitem{ghr} R.~Greenlaw, H.~J.~Hoover and W.~L.~Ruzzo. \emph{Limits to Parallel Computation:
P-Completeness Theory}. Oxford University Press, 1995.
\bibitem{hoeffding} W.~Hoeffding. Probability inequalities for sums of bounded random variables.
\emph{J.\ Amer.\ Statist.\ Assoc.} 58, 1963.
\bibitem{krs} C.~P.~Kruskal, L.~Rudolph and M.~Snir. A complexity theory of efficient parallel
algorithms. \emph{Theoret.\ Comput.\ Sci.} 71(1), 1990.
\bibitem{okhotin} A.~Okhotin. A simple P-complete problem and its language-theoretic
representations. \emph{Theoret.\ Comput.\ Sci.} 412, 2011.
\bibitem{unger} F.~Unger. Noise threshold for universality of two-input gates.
\emph{IEEE Trans.\ Inform.\ Theory} 54(8), 2008; arXiv:0711.0351.
% Bibliography entries for sec_kleene.tex.
% The fragment also cites the existing key "hoeffding"
% (W. Hoeffding, Probability inequalities for sums of bounded random
% variables, J. Amer. Statist. Assoc. 58 (1963) 13--30); it is not repeated here.

\bibitem{kp-borodin}
A.~Borodin, \emph{On relating time and space to size and depth}, SIAM J.
Comput. \textbf{6} (1977), no.~4, 733--744.

\bibitem{kp-devroye}
L.~Devroye, \emph{Branching processes in the analysis of the heights of
trees}, Acta Inform. \textbf{24} (1987), no.~3, 277--298.

\bibitem{kp-mcdiarmid}
C.~McDiarmid, \emph{On the method of bounded differences}, Surveys in
Combinatorics, 1989 (Norwich, 1989), London Math. Soc. Lecture Note Ser.,
vol.~141, Cambridge Univ. Press, Cambridge, 1989, pp.~148--188.

\bibitem{kp-pittel}
B.~Pittel, \emph{Note on the heights of random recursive trees and random
$m$-ary search trees}, Random Structures Algorithms \textbf{5} (1994), no.~2,
337--347.
% Bibliography entries for sec_gapbase.tex (existing keys used there:
% okhotin, ghr, hoeffding)
\bibitem{gb-blasermanthey}
M.~Bl\"aser and B.~Manthey, \emph{Smoothed complexity theory}, ACM
Transactions on Computation Theory \textbf{7} (2015), no.~2, Article~6.

\bibitem{gb-ladner}
R.~E. Ladner, \emph{The circuit value problem is log space complete for
{P}}, SIGACT News \textbf{7} (1975), no.~1, 18--20.

\bibitem{gb-pippenger}
N.~Pippenger, \emph{On networks of noisy gates}, Proceedings of the 26th
Annual Symposium on Foundations of Computer Science, 1985, pp.~30--38.

\bibitem{gb-spielmanteng}
D.~A. Spielman and S.-H. Teng, \emph{Smoothed analysis of algorithms: why the
simplex algorithm usually takes polynomial time}, Journal of the ACM
\textbf{51} (2004), no.~3, 385--463.

\bibitem{gb-vonneumann}
J.~von Neumann, \emph{Probabilistic logics and the synthesis of reliable
organisms from unreliable components}, Automata Studies (C.~E. Shannon and
J.~McCarthy, eds.), Princeton University Press, 1956, pp.~43--98.
\bibitem{nw-azuma67}
K.~Azuma, \emph{Weighted sums of certain dependent random variables},
T\^ohoku Math. J. \textbf{19} (1967), no.~3, 357--367.

\bibitem{nw-bfs12}
G.~E. Blelloch, J.~T. Fineman, and J.~Shun, \emph{Greedy sequential maximal
independent set and matching are parallel on average}, Proc. 24th ACM
Symposium on Parallelism in Algorithms and Architectures (SPAA), 2012.

\bibitem{nw-cf90}
N.~Calkin and A.~Frieze, \emph{Probabilistic analysis of a parallel algorithm
for finding maximal independent sets}, Random Structures Algorithms
\textbf{1} (1990), no.~1, 39--50.

\bibitem{nw-crt87}
D.~Coppersmith, P.~Raghavan, and M.~Tompa, \emph{Parallel graph algorithms
that are efficient on average}, Proc. 28th IEEE Symposium on Foundations of
Computer Science (FOCS), 1987.

\bibitem{nw-ekl10}
J.~Esparza, S.~Kiefer, and M.~Luttenberger, \emph{Newtonian program analysis},
J. ACM \textbf{57} (2010), no.~6.

\bibitem{nw-fn18}
M.~Fischer and A.~Noever, \emph{Tight analysis of parallel randomized greedy
MIS}, Proc. 29th ACM-SIAM Symposium on Discrete Algorithms (SODA), 2018.

\bibitem{nw-freedman75}
D.~A. Freedman, \emph{On tail probabilities for martingales}, Ann. Probab.
\textbf{3} (1975), no.~1, 100--118.

\end{thebibliography}

\end{document}
